% Satellite article to the corrected edition of James P. Jones,
% "Undecidable Diophantine Equations", Bull. Amer. Math. Soc. (N.S.) 3 (1980), 859--862.
% 14 September 2026.
% Companion program: ../verification/round4_1980_operation_count.py
% The certificate table jones1980_theorem5_schedule.tex is generated by
% ../verification/round4_1980_schedule_table.py from the verified schedule.
% Compile twice with: pdflatex jones1980_theorem5_operations.tex
\documentclass[11pt,letterpaper]{article}
\usepackage[T1]{fontenc}
\usepackage[utf8]{inputenc}
\usepackage{lmodern}
\usepackage[margin=0.9in]{geometry}
\usepackage{amsmath,amssymb,amsthm}
\usepackage{microtype}
\usepackage{booktabs,longtable,array}
\usepackage{enumitem}
\usepackage{alphalph}
\usepackage[hidelinks]{hyperref}
\hypersetup{pdftitle={Straight-line certificates for Jones's Theorem 5: from 129 to 89 operations},
 pdfauthor={Editors of the corrected edition},pdfsubject={Satellite article to the corrected 1980 edition}}
\newtheorem{theorem}{Theorem}[section]
\newtheorem{proposition}[theorem]{Proposition}
\newtheorem{lemma}[theorem]{Lemma}
\newtheorem{corollary}[theorem]{Corollary}
\theoremstyle{definition}
\newtheorem{definition}[theorem]{Definition}
\theoremstyle{remark}
\newtheorem{remark}[theorem]{Remark}
\numberwithin{equation}{section}
\newcommand{\Matijasevic}{Matijasevi\v{c}}
\newcommand{\pow}{\ \mathrm{pow}\ }
\newcommand{\file}[1]{\texttt{#1}}
\setlength{\parskip}{2pt}
\setlength{\jot}{4pt}

\title{\Large\bfseries Straight-line certificates for Jones's Theorem 5 (1980):\\
 \large the reported count $100$, and a reduction from $129$ to $89$}
\author{\normalsize Satellite article to the corrected edition of
 \textit{Undecidable Diophantine Equations}\\
 \normalsize Satellite to the corrected edition, 14 September 2026}
\date{}

\begin{document}
\maketitle

\begin{abstract}
\noindent
Jones's 1980 announcement states that for the equations of its Theorem~3,
``modified to do away with $q=b^{5^{60}}$'', the total number of indicated
arithmetical operations is $o=100$, and derives from this Theorem~5: every
theorem of an effectively axiomatizable theory has ``another proof consisting
of 100 additions and multiplications of integers''. Earlier revisions
of the corrected edition retained 100 as the author's reported count. This
satellite article settles what the number is and proves Theorem~5 in a
precise form.

(i)~\emph{Where 100 comes from.} The seventeen polynomial equations of
Theorem~3 (all equations except $q=b^{5^{60}}$) contain exactly 100
addition, subtraction and multiplication signs when each sign is counted
once as printed and exponentiation is not counted. The same rule applied to
the fourteen equations of Theorem~2.12 of Jones, Sato, Wada and Wiens (1976)
gives exactly their reported count~87, and applied to the thirty-six
equations (1.3) of Jones (1978) gives 241 against the reported 243. This is
the only counting rule we found that reproduces the published numbers.

(ii)~\emph{An explicit certificate.} We write down the modified system:
the equation $q=b^{5^{60}}$ is replaced, by Lemma~2.26 of Jones's 1982
paper, by three polynomial equations in four new unknowns, the exponent
$5^{60}$ entering only as an integer constant. We prove that the modified
system still defines $x\in W_{\langle z,u,y\rangle}$, and we give a
straight-line program, checked symbolically, that verifies all twenty
equations with 129 additions and multiplications of integers (75
multiplications, 54 additions, subtractions being checked as additions).
Taking certificate length as the complexity measure, we reduce this to
a \textbf{89-instruction certificate} using 47
multiplications and 42 additions for 22 equations in 34 positive
unknowns, with three fixed index parameters besides the input $x$.
The reductions combine shared arithmetic, exact Pell index recovery,
and a coefficient compiler with an affine radix, signed carry resets,
and two copied inputs. An odd-index polynomial identity reduces the
auxiliary Pell parameter; a positive product bound and binary codes
remove two coding operations and a redundant supplied witness.
Narrower packing and a carry-exclusion argument remove one more
multiplication while preserving the same fixed index and raw input.
Every power is built by multiplication; fixed numerals have no cost.
A finite support layout encodes $L$ by the fixed parameter $\psi_2(L)$.
We give complete instruction tables and positive-domain proofs for
the successive constructions, with no lower bound or optimality claim.

(iii)~\emph{Theorem 5.} With these certificate bounds, Theorem~5 follows from
Theorem~3 by the standard passage from an effectively axiomatizable theory
to a recursively enumerable set of G\"odel numbers; we spell this out and
state exactly what the ``proof'' of Theorem~5 is and what it is not.

(iv)~\emph{How the count arose.} Quotations from the 1976, 1978, 1980 and
1982 papers show that $o$ was defined as the number of operations
``necessary to write the system'', ``very closely associated with the
number of symbols'', that is, as typographical complexity; the
proof-theoretic reading was attached to that number afterwards and the
convention was inherited from the 1976 primality bound, where the two
counts coincide.
\end{abstract}

\section{The claim and its context}

The 1980 announcement \cite{J80} presents, without proof, three universal
systems of Diophantine equations and reports several size measures for them.
On its last page it introduces a further measure:

\begin{quote}\small
A different measure of size and complexity of a system of diophantine
equations is the total number of indicated arithmetical operations, $o$, i.e.\
the number of additions, subtractions and multiplications necessary to
evaluate or determine the correctness of a proposed solution. For the
equations of Theorem~3, modified to do away with $q=b^{5^{60}}$, it can be
shown that $o=100$.
\end{quote}

\noindent and states

\begin{quote}\small
\textsc{Theorem 5.} For any effectively axiomatizable theory $T$ and any
proposition $P$, if $P$ has a proof in $T$, then $P$ has another proof
consisting of 100 additions and multiplications of integers.
\end{quote}

\noindent The full paper \cite{J82} repeats both statements verbatim
(``By an arithmetical operation we shall mean an addition or multiplication.
We will consider subtraction as merely a signed form of addition.\ \ldots\ it
can be seen that $o=100$''), and adds that the corresponding count for the
36-equation system of \cite{J78} was $o=243$. Neither paper explains the
count, and neither writes down the ``modified'' system. The corrected
edition therefore carried the footnote that the exact count 100 is the
author's reported count, not independently reconstructed. This article
removes that caveat.

The system in question is Theorem~3 of \cite{J80} (identical to Theorem~3 of
\cite{J82}), in the corrected form of the present edition:
\begin{align}
 e\ell g^2+\alpha&=(b-xy)q^2, \tag{E1}\\
 q&=b^{5^{60}}, \tag{E0}\\
 \lambda+q^4&=1+\lambda b^5, \tag{E2}\\
 \theta+2z&=b^5, \tag{E3}\\
 \ell&=u+t\theta, \tag{E4}\\
 e&=y+m\theta, \tag{E5}\\
 n&=q^{16}, \tag{E6}\\
 r&=\Bigl[g+eq^3+\ell q^5+\bigl(2(e-z\lambda)(1+xb^5+g)^4
     +\lambda b^5+\lambda b^5q^4\bigr)q^7\Bigr](n^2-n) \notag\\
  &\qquad+\Bigl[q^3-b\ell+\ell+\theta\lambda q^3+(b^5-2)q^8\Bigr](n^2-1), \tag{E7}\\
 p&=2ws^2r^2n^6, \tag{E8}\\
 p^2k^2-k^2+1&=\tau^2, \tag{E9}\\
 4(c-ksn^2)^2+\eta&=k^2, \tag{E10}\\
 k&=r+1+hp-h, \tag{E11}\\
 a&=(wn^2+1)rsn^2, \tag{E12}\\
 c&=2r+1+\varphi, \tag{E13}\\
 d&=bw+ca-2c+4a\gamma-5\gamma, \tag{E14}\\
 d^2&=(a^2-1)c^2+1, \tag{E15}\\
 f^2&=(a^2-1)i^2c^4+1, \tag{E16}\\
 (d+of)^2&=\Bigl(\bigl(a+f^2(d^2-a)\bigr)^2-1\Bigr)(2r+1+jc)^2+1. \tag{E17}
\end{align}
The unknowns are the 28 positive integers
\[
 a,b,c,d,e,f,g,h,i,j,k,\ell,m,n,o,p,q,r,s,t,w,\alpha,\gamma,\eta,\theta,\lambda,\tau,\varphi;
\]
the parameters are $x$ and the admissible coding triple $\langle z,u,y\rangle$
of \cite[(4.1)]{J82}, with $\delta=4$ and $\nu=58$. We have numbered the
equations in the printed order, giving the exponential equation the label
(E0) so that the seventeen polynomial equations keep the labels (E1)--(E17)
throughout.

\medskip\noindent\textbf{Where $5^{60}$ comes from.} In \cite[\S4]{J82} the
base is $B=b^{\delta+1}=b^5$ and the code length is
$q=B^{(\delta+1)^{\nu+1}}=B^{5^{59}}$, hence $q=b^{5\cdot5^{59}}=b^{5^{60}}$.
The exponent is a fixed integer; it is not built up inside the system.

\section{Two counting conventions}

The phrase ``number of additions, subtractions and multiplications necessary
to evaluate or determine the correctness of a proposed solution'' admits two
readings, and the two give different numbers for Theorem~3. We define both.

\begin{definition}[indicated operations]\label{def:indicated}
The \emph{indicated operation count} of a system of polynomial equations is
the number of addition, subtraction and multiplication signs in the printed
system, each occurrence counted once: a binary $+$ or $-$ counts one, a
multiplication sign or a juxtaposition of two factors counts one, a numeral
coefficient such as $2z$ or $4a\gamma$ counts one multiplication per factor,
and an exponent $x^k$ is not counted. Repeated subexpressions are counted each
time they are printed.
\end{definition}

\begin{definition}[straight-line certificate]\label{def:certificate}
A \emph{straight-line certificate} for a system of polynomial equations in
the variables $v_1,\dots,v_m$ consists of integer values for the variables,
integer values for finitely many auxiliary names, and a finite list of
instructions, each of the form $t=u+v$ or $t=u\cdot v$ where $u,v$ are
variables, auxiliary names, earlier instruction targets, or integer numerals,
followed by a finite list of equality tests between names. It is
\emph{valid} if every instruction and every equality test holds for the
supplied values. Its \emph{length} is the number of instructions. A
subtraction $t=u-v$ is recorded as the instruction $t+v=u$, so subtraction
costs one operation and no subtraction primitive is needed. Equality tests,
reading the supplied integers, and (in the basic form) writing down numerals
are free. A certificate \emph{verifies} the system if the validity of the
certificate implies that the supplied values satisfy every equation of the
system, and every solution of the system extends to a valid certificate.
\end{definition}

\begin{definition}[certificate complexity]\label{def:complexity}
The complexity $\mathcal C(\Sigma)$ is the least length of a
\emph{uniform certificate scheme} verifying $\Sigma$ in the preceding
sense. Its finite list of names, instruction operands and equality tests
is fixed before the values of the input and witnesses are supplied.
Witness domains are part of the scheme. The integer constants and any
fixed admissible index are specified when the represented set is chosen;
an existential witness cannot stand in for an unchecked encoding parameter.
\end{definition}

Definition~\ref{def:certificate} is the ``calculator'' reading: a calculator
that can only add and multiply integers, variables with fixed values, any
auxiliary integers typed for free, intermediate results written down and
reused at will. It is the convention under which the corrected 1976 edition
reproduces the 87-operation primality certificate of \cite[Theorem~5]{JSWW},
and it is the convention that makes ``a proof consisting of $o$ additions and
multiplications of integers'' literally true: the proof is the list of
instructions. Definition~\ref{def:indicated} is the ``as written'' reading of
``indicated operations'': it charges nothing for the exponentiations that a
calculator would have to perform by repeated multiplication, and it charges
twice for a subexpression printed twice.

The two conventions coincide for a system in which every power is a square of
a variable that is not reused and no subexpression repeats. They diverge for
Theorem~3, which contains 44 exponentiations, among them $q^{16}$, $q^8$,
$q^7$, $q^5$, $q^3$, six occurrences of $b^5$, $n^6$, $c^4$ and a fourth
power of a trinomial.

\section{Where the number 100 comes from}

\begin{proposition}\label{prop:100}
The seventeen polynomial equations \textup{(E1)--(E17)} of Theorem~3 contain
exactly 100 indicated operations in the sense of
Definition~\ref{def:indicated}: 33 additions, 15 subtractions and 52
multiplications. The 44 exponentiations are not counted, and the equation
\textup{(E0)}, $q=b^{5^{60}}$, contributes nothing.
\end{proposition}

\begin{proof}
Count each equation as printed. The per-equation totals are
\[
\begin{array}{@{}l*{17}{r}@{}}
\text{equation} & \text{E1}&\text{E2}&\text{E3}&\text{E4}&\text{E5}&\text{E6}&\text{E7}&\text{E8}&\text{E9}&\text{E10}&\text{E11}&\text{E12}&\text{E13}&\text{E14}&\text{E15}&\text{E16}&\text{E17}\\
\text{signs} & 6&3&2&2&2&0&32&4&3&5&4&5&3&10&3&4&12
\end{array}
\]
with sum 100. For instance (E1) has $e\ell g^2$ (two multiplications),
$+\alpha$, $xy$, $b-xy$ and $(\cdot)q^2$: six signs; (E7) has 32 signs, of
which the first bracket contributes $g+eq^3$ (2), $+\ell q^5$ (2),
$+2(e-z\lambda)(1+xb^5+g)^4$ (8), $+\lambda b^5$ (2), $+\lambda b^5q^4$ (3),
$(\cdot)q^7$ (1), and the rest is $(n^2-n)$ (1), the product (1), the second
bracket $q^3-b\ell+\ell+\theta\lambda q^3+(b^5-2)q^8$ (9), $(n^2-1)$ (1), the
product (1) and the final sum (1). The companion program
\file{round4\_1980\_operation\_count.py} parses the transcribed equations and
returns the same totals.
\end{proof}

\begin{proposition}[calibration]\label{prop:calibration}
Under Definition~\ref{def:indicated}:
\begin{enumerate}[nosep]
\item the fourteen equations of Theorem~2.12 of \cite{JSWW}, from which the
 authors ``calculated'' the bound 87 of their Theorem~5, contain exactly 87
 indicated operations (35 additions, 17 subtractions, 35 multiplications;
 24 exponentiations not counted);
\item the thirty-six equations \textup{(1.3)} of \cite{J78}, for which the
 author reports 243, contain 241 indicated operations (111 additions,
 31 subtractions, 99 multiplications; 54 exponentiations not counted);
\item the seventeen polynomial equations of Theorem~3 of \cite{J80} contain
 100 indicated operations, the reported value.
\end{enumerate}
\end{proposition}

\begin{proof}
Direct counts, reproduced by the companion program from transcriptions of
the three systems, which were checked against the scans.
\end{proof}

\begin{remark}[reading the evidence]\label{rem:evidence}
Two of the three published counts are reproduced exactly and the third to
within two signs; we could not locate the two missing signs in (1.3) and
record the discrepancy. No other rule we tried comes close on all three
systems. In particular, under Definition~\ref{def:certificate} with every
power built by repeated multiplication and every repeated subexpression
computed once, a direct shared-expression evaluation of the printed
Theorem~3 without (E0) uses 127 operations before anything replaces the
exponential equation, and the initial certificate of
Section~\ref{sec:certificate} uses 129 with the
replacement; for \cite{J78} the same convention gives 242, and for
\cite{JSWW} the naive count is 99, which the corrected 1976 edition brings down to 87
only by algebraic rewriting. The natural conclusion is that Jones's $o$ is
the count of indicated operation signs of Definition~\ref{def:indicated},
with exponentiation regarded as notation rather than as an operation, and
that ``modified to do away with $q=b^{5^{60}}$'' means that the exponential
equation is set aside and not that its polynomial replacement has been
counted. One small inconsistency should be recorded: \cite{J78} remarks that
combining its 36 equations into one by transposing and summing squares
``would increase this number to 350'', and $350-243=107=36+36+35$ counts
one subtraction, one squaring and (for 35 joins) one addition per equation;
there the squaring is charged, whereas the powers inside the printed
equations are not.
\end{remark}

\section{The modified system}\label{sec:modified}

To make Theorem~5 a statement about a polynomial system we must actually do
away with (E0). The 1982 paper provides the tool: Lemma~2.26 of \cite{J82}
defines an exponential $Q=B^{B_1}$ on top of the Pell machinery that
Theorem~3 already contains, at the cost of three conditions and four new
unknowns, with the exponent $B_1$ entering only as an integer constant. We
recall it in the notation of Theorem~3.

\begin{lemma}[{\cite[Lemma 2.26, primed form]{J82}}]\label{lem:226}
Let $R,N,b$ satisfy the hypotheses of \textup{\cite[Lemma 2.25]{J82}}
($R\ge8$, $N\ge8$, $R\ge b$, $N\ge b>0$), and let $B,Q,B_1$ be positive
integers with $3<3B_1\le B\le Q\le N\le R$. Then
$N^2\mid\binom{2R}{R}$, $b\pow2$ and $Q=B^{B_1}$ hold together if and only if
the equations \textup{(B1)--(B14)} of Lemma~2.25 hold with \textup{(B8)}
replaced by $C=C_1+B'+\varphi$ and the following three conditions adjoined,
for some positive integers $C_1,D_1,\Delta$:
\begin{align*}
 D_1&\equiv Q+C_1(A-B)\pmod{2AB-B^2-1},\tag{C1$'$}\\
 (A^2-1)C_1^2+1&=D_1^2,\tag{C2$'$}\\
 C_1&=B_1+\Delta(A-1).\tag{C3$'$}
\end{align*}
\end{lemma}

In Theorem~3 the quantities of Lemma~2.25 appear as $R=r$, $N=n$, $A=a$,
$C=c$, $K=k$, $P=p$, $M=rsn^2$, $U=wn^2$, $Y=n^2s$, $W=bw$, $V=2$,
$B'=2r+1$; this dictionary was verified equation by equation. We now take $B:=b$, $Q:=q$ and $B_1:=5^{60}$ in Lemma~\ref{lem:226},
write $\kappa,\mu,\rho,\Delta$ for $C_1$, $D_1$, the multiplier of the
modulus in (C1$'$), and $\Delta$, and obtain:

\begin{definition}[the modified system $\Sigma$]\label{def:sigma}
$\Sigma$ consists of the seventeen equations \textup{(E1)--(E17)} with
\textup{(E13)} replaced by
\begin{equation}
 c=2r+1+\kappa+\varphi \tag{E13$'$}
\end{equation}
and the three equations
\begin{align}
 \mu&=q+\kappa(a-b)+\rho\,(2ab-b^2-1), \tag{E18}\\
 (a^2-1)\kappa^2+1&=\mu^2, \tag{E19}\\
 \kappa&=5^{60}+\Delta(a-1), \tag{E20}
\end{align}
in the 32 unknowns $a,\dots,\varphi,\kappa,\mu,\rho,\Delta$, all ranging over
positive integers, and the parameters $x,z,u,y$.
\end{definition}

The constant $5^{60}$ in (E20) is a numeral, exactly as $2$, $4$ and $5$ are
numerals in (E3), (E10) and (E14). Nothing in $\Sigma$ has degree larger than
that of (E17); the degree of $\Sigma$ in the unknowns is $\max(16,\,\deg
\text{(E7)},\,\deg\text{(E17)})$ and no longer $5^{60}$.

\begin{theorem}\label{thm:sigma}
For every admissible coding triple $\langle z,u,y\rangle$ and every positive
integer $x$: $x\in W_{\langle z,u,y\rangle}$ if and only if $\Sigma$ has a
solution in positive integers.
\end{theorem}

\begin{proof}
The proof of Theorem~3 in \cite[\S4]{J82} shows that (E1)--(E17), including
(E0), are equivalent to the conditions (U1)--(U24) of \S4 together with
$b\pow2$ and $n^2\mid\binom{2r}{r}$, the latter two being expressed by
(B1)--(B14) of Lemma~2.25 with (B14) replaced by (Q1)--(Q3) and (B13) in the
$\chi$-form. In that proof the hypotheses of Lemma~2.25 are established from
the other equations: $8\le r$, $8\le n$, $b\le r$, $0<b\le n$; the two bounds
on $r$ are obtained there through the block bounds $0\le S_i,T_i$, whose
derivation uses (E0). We check these two bounds and the additional
hypotheses of Lemma~\ref{lem:226} with $B=b$, $Q=q$, $B_1=5^{60}$, in both
directions, using only the equations of $\Sigma$ or the parameters.

\emph{$3<3\cdot5^{60}\le b$.} From (E1), $(b-xy)q^2=e\ell g^2+\alpha\ge2$,
hence $b>xy\ge y$. For an admissible triple, by \cite[(4.1)]{J82} with
$\delta=4$, $\nu=58$, $z\ge2$ and $z+P_{0,\dots,0}\ge1$,
\[
 y\ \ge\ (z+P_{0,\dots,0})(2z)^{5^{59}}\ \ge\ 2^{5^{59}}\ >\ 3\cdot5^{60},
\]
the last inequality because $\log_2(3\cdot5^{60})<141<5^{59}$.

\emph{$b\le q$.} From (E2), $\lambda(b^5-1)=q^4-1$ with $\lambda\ge1$, so
$q^4\ge b^5>b^4$ and $q>b$.

\emph{$q\le n\le r$, with $8\le r$ and $b\le r$.} $n=q^{16}\ge q$ by (E6). The
block bounds $0\le S_i,T_i$ of \cite[\S4]{J82} cannot be invoked here: the
derivation of $0\le S_3$ there uses $xb^5\le q$, which comes from (E0).
Instead write $S$ and $T$ for the two brackets of (E7), so that
$r=S(n^2-n)+(T+1)(n^2-1)$ factors as
\[
 r=(n-1)\bigl(Sn+(T+1)(n+1)\bigr).
\]
Since $r\ge1$ and $n-1\ge1$, the second factor is a positive integer, and it
equals $1$ only if $(T+1)(n+1)\equiv1\pmod n$, that is, only if $n\mid T$.
But $0<T<n$: by (E5), (E3), $m\ge1$ and $y\ge2z$ (from the bound on $y$
above) we have $e\ge y+\theta\ge b^5$, so (E1) gives
$b^5\ell\le e\ell g^2<(b-xy)q^2<bq^2$, hence $\ell<q^2$ and
$\ell(b-1)<bq^2<q^3$, which makes $T>q^3-\ell(b-1)>0$; and
$\theta\lambda<q^4$ because $\theta<b^5-1$, so
$T<q^3+q^7+b^5q^8<q^{16}=n$. Therefore the second factor is at least $2$,
and $r\ge2(n-1)\ge n$; in particular $r\ge8$ and $r>b$.

\emph{Sufficiency.} Let the 32 unknowns satisfy $\Sigma$. Lemma~2.25 and
Lemma~\ref{lem:226} apply, the unprimed conditions (B1)--(B14) being
(E8)--(E17) with $c=2r+1+\kappa+\varphi$ in the role of $C=C_1+B'+\varphi$,
and (E18)--(E20) being (C1$'$)--(C3$'$): (E18) with $\rho\ge1$ implies the
congruence (C1$'$). Hence $b\pow2$, $n^2\mid\binom{2r}{r}$ and $q=b^{5^{60}}$.
So (E0) holds, and the original proof of Theorem~3 gives
$x\in W_{\langle z,u,y\rangle}$.

\emph{Necessity.} Let $x\in W_{\langle z,u,y\rangle}$. The necessity proof of
Theorem~3 supplies positive values of the 28 original unknowns satisfying
(E0)--(E17), with $q=b^{5^{60}}$, $a=A$ the Pell modulus of Lemma~2.25 and
$c=\psi_a(2r+1)$. Put $\kappa=\psi_a(5^{60})$ and $\mu=\chi_a(5^{60})$, the
solutions of $\mu^2-(a^2-1)\kappa^2=1$; this is (E19). Since
$\psi_a(k)\equiv k\pmod{a-1}$ and $\psi_a(k)>k$ for $a\ge2$, $k\ge2$, we have
$\kappa=5^{60}+\Delta(a-1)$ with an integer $\Delta\ge1$; this is (E20).
Lemma~2.22 of \cite{J82} (the $\chi$-form of \cite[Remark after Lemma
2.28]{J82}) gives $\mu\equiv q+\kappa(a-b)\pmod{2ab-b^2-1}$, so
$\mu-q-\kappa(a-b)=\rho(2ab-b^2-1)$ for an integer $\rho$. To see $\rho\ge1$
note that $2ab-b^2-1=(a^2-1)-(a-b)^2>0$ since $1<b<a$, that
$\mu>\sqrt{a^2-1}\,\kappa>(a-1)\kappa=(a-b)\kappa+(b-1)\kappa$, and that
$(b-1)\kappa\ge\kappa=\psi_a(5^{60})\ge(2a-1)^{5^{60}-1}>a>q$, the last step
by the estimate $Q^3<A$ in \cite[(13)]{J82}; hence the left side is positive
and $\rho\ge1$. Finally $\varphi$ must be re-chosen for (E13$'$):
$c-2r-1-\kappa=\psi_a(2r+1)-(2r+1)-\psi_a(5^{60})>0$ by
\cite[(14)]{J82}, which is exactly the inequality
$\psi_A(B_1)+2R+1+2^{2R+1}<\psi_A(2R+1)$ proved there for $3B_1\le R<A$.
Thus all 32 values are positive integers satisfying $\Sigma$.
\end{proof}

\begin{remark}
One may equally take $B:=b^5$ and $B_1:=5^{59}$, which is the literal reading
of \cite[\S4]{J82}; then (E18) reads
$\mu=q+\kappa(a-b^5)+\rho\,(2ab^5-b^{10}-1)$ and (E20) has the constant
$5^{59}$. Since $b^5$ is already computed, the modulus identity used below
gives the same basic instruction count for both choices. The hypothesis
$3\cdot5^{59}\le b^5$ is checked in the same way. The unprimed
group (C1)--(C3), which uses the congruence
$(B^2-1)QC_1\equiv B(Q^2-1)$, costs more operations and is not used.
\end{remark}

\section{An explicit straight-line certificate}\label{sec:certificate}

\begin{theorem}\label{thm:129}
The system $\Sigma$ of Definition~\ref{def:sigma} is verified, in the sense
of Definition~\ref{def:certificate}, by the straight-line certificate of the
Appendix~\ref{app:129}, which has 129 instructions: 75 multiplications and 54 additions
(37 additions and 17 subtractions, each subtraction recorded as an
addition), followed by 20 equality tests. If the numerals $2,4,5$ and
$5^{60}$ must be generated from $1$ by additions and multiplications, ten
further instructions suffice ($2=1+1$, $4=2+2$, $5=4+1$ and the addition
chain $1,2,3,6,12,15,30,60$ for $5^{60}$), for a total of 139.
\end{theorem}

\begin{proof}
For each of the 20 equations of $\Sigma$ the certificate computes both sides
(or one side, when the other is a single variable) and the program
\file{round4\_1980\_operation\_count.py} verifies, with SymPy, that the
difference of the two compared names is, as a polynomial in the 36 variables
of $\Sigma$, equal to the left side minus the right side of the source
equation, up to sign. Hence the equality tests hold exactly when the
equations hold, so a valid certificate proves the equations, and every
solution extends to a valid certificate by evaluating the instructions. The
counts are read off the list. The residual signs, the operation histogram
and the count of the numeral chain are written to
\file{round4\_1980\_operation\_count.json}.
\end{proof}

The economies used, all of them plain integer identities, are the following.
Powers are shared: $q^2,q^3,q^4,q^8,q^{16}$ (five multiplications, the
shortest chain containing $2,3,4,8,16$) and $b^2,b^4,b^5$ (three). The first
bracket of (E7) is evaluated by Horner's rule as
$g+q^3\bigl(e+\ell q^2+q^4(2(e-z\lambda)C^4+\Lambda(1+q^4))\bigr)$ with
$C=1+xb^5+g$ and $\Lambda=\lambda b^5$ taken from (E2), which removes the
need for $q^5$ and $q^7$; $C^4$ is two squarings. In (E8)--(E12) the products
$wn^2$, $sn^2$ and $rsn^2$ are computed once and
$p=2ws^2r^2n^6=2\,(wn^2)(rsn^2)^2$. (E9) is evaluated as $(p^2-1)k^2+1$,
(E11) as $r+1+h(p-1)$, (E14) as $bw+c(a-2)+\gamma(4a-5)$, (E16) as
$(a^2-1)(ic^2)^2+1$, and $a^2-1$, $d^2$, $f^2$, $c^2$, $r+1$, $2r+1$ are
shared between (E15), (E16), (E17), (E19). The modulus of (E18) is
$2ab-b^2-1=(a^2-1)-(a-b)^2$, the identity that the corrected 1976 edition uses for the
1976 count. All intermediate integers, including the signed ones such as
$e-z\lambda$ and $b^5-2$, are supplied with the certificate.

\begin{remark}[the count is an upper bound]
\label{rem:whynot}
Neither the printed sign count nor a particular evaluation schedule gives
a lower bound for all certificates. In particular, computations may be
shared across equations, previously checked equalities may be reused, and
the underlying Diophantine encoding may be changed. The 129 instructions
above are a reproducible starting point for these optimizations. The
historical interpretation of 100 rests on the evidence in
Remark~\ref{rem:evidence}, not on an assertion that smaller certificates
are impossible.
\end{remark}

\begin{theorem}[sharing a product]\label{thm:128}
The same system $\Sigma$ has a certificate with 128 instructions:
74 multiplications and 54 additions, followed by the same 20 equality
tests. Generating the numerals from $1$ raises this bound to 138.
\end{theorem}

\begin{proof}
Put $U=wn^2$, $Y=sn^2$, $M=rY$ and $V=UM$. Then the two equations
\[
 a=(U+1)M,\qquad p=2UM^2
\]
can be checked as $a=V+M$ and $p=2(VM)$. Computing $U,Y,M,V,V+M,VM,2(VM)$
takes seven instructions. These replace the eight instructions 58--65
of Appendix~\ref{app:129} and preserve both residual polynomials exactly. The
companion program checks all twenty residuals again for this modified
schedule and records the result as \texttt{shared\_product\_certificate}.
\end{proof}

% Input section for jones1980_theorem5_operations.tex.
% The archived 129-step certificate is preserved; this section supplies
% independently verified improvements, culminating in the packed 120 steps.
\section{Reducing certificate length to 120}\label{sec:optimization}

We now take straight-line certificate length, rather than indicated
operation count, as the complexity measure to minimize. The preceding
129-instruction schedule is a useful starting point, but it is not a
minimum. Two shared computations reduce the certificate for exactly the
same system to 127. A positive Pell-solution gap removes a redundant
inequality and gives an equivalent system with a 126-instruction
certificate. Finally, changing the admissible index and packing its two
coefficient congruences gives a universal system with a certificate of
length 120. The last construction still has 32 positive unknowns, twenty
equations, and three fixed index parameters besides the input~$x$.

\subsection{Two savings without changing the system}

Put $U=wn^2$ and $M=rsn^2$, quantities already computed in the baseline
schedule. Equations (E8) and (E12) require
\[
 p=2UM^2,\qquad a=(U+1)M.
\]
Compute $J=UM$ once and instead check
\[
 a=J+M,\qquad p=2JM.
\]
This replaces five instructions, after $U,M$ are available, by four: one
multiplication is saved.

For the second saving, temporarily write $B=b^5$. Equations (E2), (E3)
give the useful identity
\begin{equation}
 1+\theta\lambda=\lambda+q^4-2z\lambda.\label{eq:reuse-theta}
\end{equation}
Compute $Z_\lambda=(2z)\lambda$, and write the coefficient
$2(e-z\lambda)$ in the first bracket of (E7) as $2e-Z_\lambda$.
This costs the same three instructions as before. Its second bracket
then becomes
\[
 q^3(\lambda+q^4-Z_\lambda)-(b-1)\ell+(B-2)q^8.
\]
The quantity $\lambda+q^4$ is already available from (E2), and
$Z_\lambda$ is now shared between the two brackets. Eight instructions
replace the previous nine.

\begin{theorem}\label{thm:127}
The original modified system $\Sigma$ has a straight-line certificate
with 127 instructions: 73 multiplications and 54 additions, with
subtractions recorded as additions, followed by twenty equality tests.
\end{theorem}

\begin{proof}
Apply the two independent savings to the archived 129-instruction
schedule. The shared-product identities are polynomial identities. The
second saving uses equations already checked, so its exact verification
must record that dependence. If $F_i$ denotes the source residual, left
side minus right side, of (Ei), the residual evaluated by the new (E7)
schedule is
\[
 F_{7,\mathrm{new}}
   =F_7+(\lambda F_3-F_2)q^3(n^2-1).
\]
Thus it vanishes exactly when $F_7$ does, once (E2) and (E3) have been
checked. The companion program verifies this identity symbolically,
along with the other nineteen residuals and the instruction count.
\end{proof}

\subsection{Removing a redundant Pell gap}

The strengthened equation (E13$'$) asks for
$c>\kappa+2r+1$. The two Pell equations already ensure that $c>\kappa$
implies this stronger inequality. The following elementary argument does
not need the general classification of Pell solutions.

\begin{lemma}\label{lem:pellgap}
Let $a\ge2$ and suppose positive integers $c,d,\kappa,\mu$ satisfy
\[
 d^2-(a^2-1)c^2=\mu^2-(a^2-1)\kappa^2=1,
 \qquad c>\kappa.
\]
Then $c-\kappa\ge2a-1$.
\end{lemma}

\begin{proof}
Write $D=a^2-1$, $t=\mu c-d\kappa$ and
$v=d\mu-Dc\kappa$. Since
\[
 \frac dc=\sqrt{D+c^{-2}}<\sqrt{D+\kappa^{-2}}
 =\frac\mu\kappa,
\]
the integer $t$ is positive. Also $v>0$, because
$d>\sqrt D\,c$ and $\mu>\sqrt D\,\kappa$. Direct multiplication gives
\[
 v^2-Dt^2=1,\qquad c=\mu t+\kappa v.
\]
As $t\ge1$, the first identity implies $v\ge a$; likewise $\mu\ge a$.
Consequently
\[
 c-\kappa=\mu t+\kappa(v-1)
       \ge a+\kappa(a-1)\ge2a-1.\qedhere
\]
\end{proof}

Let $\Sigma_{\rm flat}$ replace (E13$'$) by $c=\kappa+\varphi$,
with $\varphi$ still positive, and retain the other nineteen equations.

\begin{theorem}\label{thm:126}
The system $\Sigma_{\rm flat}$ defines the same set of inputs as
$\Sigma$, with an explicit bijection of their slack witnesses, and has a
126-instruction certificate: 73 multiplications and 53 additions.
\end{theorem}

\begin{proof}
The other equations give $b\ge2$, $q\ge2$, $n\ge2$, and
\[
 a=(wn^2+1)rsn^2\ge20r.
\]
By Lemma~\ref{lem:pellgap}, the new equation gives
$c-\kappa\ge2a-1>2r+1$. Therefore
\[
 \varphi_{\rm old}=\varphi_{\rm new}-(2r+1)>0,
 \qquad
 \varphi_{\rm new}=\varphi_{\rm old}+2r+1
\]
are inverse maps between solutions. Since $2r+1$ is already computed for
(E17), checking $c=\kappa+\varphi$ saves one addition over checking
$c=2r+1+\kappa+\varphi$. Combine this with Theorem~\ref{thm:127}.
\end{proof}

\subsection{A larger radix supplied by the index}

Put
\[
 L=5^{59},\qquad M=5^{58},\qquad L=5M.
\]
Let $(z,u,y)$ be an admissible index of the original construction and
write $Z=2z$. Choose a power of two $H$ so large that
\begin{equation}
 H>\max\{2Z^{4L+1},\ Z\,60^4,\ 3L\}.
 \label{eq:index-H}
\end{equation}
The new radix will be $B=Hb^4$. Constructing $H$ is part of constructing
the index of the represented r.e.\ set. It is an explicitly declared
fixed parameter; it is not an unrestricted existential witness supplied
by a purported certificate. This distinction is necessary for soundness.

First retain both $u$ and $y$, so the admissible index has four
components $(Z,u,y,H)$. Define $\Sigma_{\rm rec}$ by the following
replacements in $\Sigma$:
\begin{align*}
 b&=x+\beta, & e\ell g^2+\alpha&=q^2,
       &&\text{in place of (E1)},\\
 \theta+Z&=B, &&&&\text{in place of (E3)},\\
 c&=\kappa+\varphi, &&&&\text{in place of (E13$'$)},\\
 \mu&=q+\kappa(a-B)+\rho(2aB-B^2-1),
       &&&&\text{in place of (E18)},\\
 \kappa&=L+\Delta(a-1), &&&&\text{in place of (E20)}.
\end{align*}
Replace every occurrence of $b^5$ by $B=Hb^4$, and replace
$2(e-z\lambda)$ in (E7) by $2e-Z\lambda$. All remaining equations and
positive domains are retained, and $\beta$ is an additional positive
unknown. Thus $\Sigma_{\rm rec}$ has 33 unknowns and 21 equations.

\begin{theorem}\label{thm:123}
For every admissible $(Z,u,y,H)$, the system $\Sigma_{\rm rec}$ defines
the same r.e.\ set as the original index $(Z/2,u,y)$. Its certificate
uses 123 instructions: 70 multiplications and 53 additions.
\end{theorem}

\begin{proof}
The radix $B=Hb^4$ is computed in three multiplications, the same number
as $b^5$. The split (E1) costs five instructions instead of seven.
Using $Z$ directly saves the separate computation of $2z$; its occurrence
in $2e-Z\lambda$ costs the same three instructions as the old coefficient.
The remaining three savings are those of Theorems~\ref{thm:127}
and~\ref{thm:126}. Hence the count is $129-2-1-2-1=123$.

For the mathematical equivalence, the original proof uses $b>xy$ to
ensure $b>x$, a sufficiently large radix for Lemma~2.9, and the coefficient
bound needed to prevent carries. Here $b>x$ is explicit, while
\eqref{eq:index-H} gives
\[
 B>2Z^{2L+1},\qquad z\,60^4b^4<B/2.
\]
These are the required radix and coefficient bounds. The original
uncombined inequality $e\ell g^2<q^2$ gives $e,\ell<q^2$ and $g<q$.
The retained congruence equations and $u,y\ge Z$ also give
$e,\ell\ge B$, hence $q^2>e\ell g^2\ge B^2$ and $q>B$ before using the
intended exponential relation. Therefore $\lambda=(q^4-1)/(B-1)>q^3>e$, and
$0<D_0=z\lambda-e<z\lambda$.

For completeness, the positivity of the third block can also be checked
before invoking the exponent lemma. With $C=1+xB+g$,
\[
 C\le xB+q<(b+1)q,
 \qquad (b+1)^4<6b^4\quad(b\ge2).
\]
Since $H>12z$, this gives
$2C^4D_0<B\lambda q^4$, so
\[
 0<S_3=-2C^4D_0+B\lambda(1+q^4)<2q^8<q^9.
\]
The other block bounds follow from $e,\ell<q^2$, $g<q$ and $b<B<q$.
Thus $n=q^{16}\le r$, and $3<3L\le B<q\le n\le r$.
The Pell gap restores (E13$'$); Lemma~\ref{lem:226} applies with base
$B$ and exponent $L$, yielding $b\pow2$ and $q=B^L$. Since $H$ is a
power of two, all required radices are powers of two. The original
digit-extraction and no-carry proof now applies with the new radix and
the same coefficient index.

Conversely, choose a power-of-two bound $b$ for an actual solution and
use its intended codes. The source estimates
\[
 e<2zB^L,\qquad \ell<2B^M,\qquad g<bB^M
\]
give $e\ell g^2<4zb^2B^{8M}<B^{10M}=q^2$, so the new
$\alpha$ is positive; $\beta=b-x>0$. The remaining digit and Pell
witnesses are supplied by the original necessity construction, now using
base $B$ and exponent $L$. Positivity of its congruence multipliers and
of the new $\varphi=c-\kappa$ follows as in Theorem~\ref{thm:sigma}.
\end{proof}

\subsection{Packing the coefficient congruences and sharing a square}

There is a further useful change to the encoding. Put
\begin{equation}
 V=y+uZ^{2L}.\label{eq:packed-index}
\end{equation}
Call $(Z,V,H)$ admissible if it is obtained this way from an admissible
$(Z/2,u,y)$ and satisfies \eqref{eq:index-H}. This restores three index
parameters, so the complete parameter list is $x,Z,V,H$.

Define $\Sigma_{\rm pack}$ from $\Sigma_{\rm rec}$ by writing
$C=1+xB+g$, replacing its inequality equation by
\begin{equation}
 e\ell C^2+\alpha=q^2,\label{eq:packed-bound}
\end{equation}
and replacing (E4), (E5) together by
\begin{equation}
 e+\ell q^2=V+t\theta.\label{eq:packed-congruence}
\end{equation}
Delete the now-unused unknown $m$. The quantities $B,C$ are abbreviations
whose arithmetic is charged, not free auxiliary definitions. The resulting
system again has 32 positive unknowns and twenty equations.

\begin{theorem}\label{thm:120}
For every admissible packed index $(Z,V,H)$, the system
$\Sigma_{\rm pack}$ defines the same recursively enumerable set as its
original Jones index. It has a straight-line certificate with 120
instructions: 68 multiplications and 52 additions, followed by twenty
equality tests.
\end{theorem}

\begin{proof}
We first establish the bounds needed to use the Pell lemmas, without
assuming $q=B^L$. Equation~\eqref{eq:packed-bound} and positivity give
$e,\ell<q^2$ and $C<q$. Since $b>x\ge1$ and $C>B$, we have
$q>B>b\ge2$ and $g<q$. As before,
$\lambda>q^3>e$ and $0<D_0=z\lambda-e<z\lambda$. Now the stronger
bound $C<q$ directly gives
\[
 2C^4D_0<2z\lambda q^4<B\lambda q^4,
 \qquad 0<S_3<2q^8<q^9.
\]
The remaining block bounds are
\begin{gather*}
 0<g<q^3,\qquad
 0<q^3-1-(b-1)\ell<q^3,\\
 0<e+\ell q^2<q^4,\qquad
 0<(B-Z)\lambda<(B-1)\lambda=q^4-1,\\
 0<(B-2)q<q^2<q^9.
\end{gather*}
For the second use $\ell<q^2$ and $b<q$; for the third use integer
$e,\ell\le q^2-1$. Thus the three blocks lie within their widths
$q^3,q^4,q^9$, and the definition of $r$ gives $r\ge n^2-1\ge n$.
We obtain $3<3L\le B<q\le n\le r$ entirely from the polynomial
equations. Lemma~\ref{lem:pellgap} restores the stronger gap needed by
Lemma~\ref{lem:226}. The Pell lemmas therefore give $b\pow2$,
$q=B^L$ and the central-binomial divisibility. The resulting powers of
two allow the three no-carry conditions to be unpacked.

It remains to justify that the single congruence loses none of the
coefficient information. Let $e_0(B),\ell_0(B)$ denote the intended
digit polynomials of (4.5)--(4.6) in \cite{J82}, and put
\[
 F(B)=e_0(B)+\ell_0(B)B^{2L}.
\]
The first summand has degree at most $L$. The second occupies positions
$2L+5^i$ for $1\le i\le58$, all exceeding $2L$ and at most
$2L+M<4L$. Their digit positions are disjoint, all digits lie in
$[0,Z)$, and $F(Z)=V$.

Apply \cite[Lemma~2.9]{J82} to $Y=e+\ell q^2$, digit base $Z$, and
$m=k=4L$. Its congruence hypothesis is
\eqref{eq:packed-congruence}; its size hypothesis follows from
$Y<q^4=B^{4L}<ZB^{4L}$; and its mask hypothesis is precisely the second
no-carry condition
\[
 \tau_2\bigl(Y,(B-Z)\lambda\bigr)=0.
\]
The radix hypothesis $2Z^{4L+1}\le B$ follows from
\eqref{eq:index-H}. The lemma gives $Y=F(B)$. Both the supplied and
the intended $e,\ell$ lie in $[0,q^2)$, so uniqueness of quotient and
remainder upon division by $q^2$ gives $e=e_0(B)$ and
$\ell=\ell_0(B)$. The first no-carry condition now identifies the
witness code $g$, and the third gives the vanishing coefficient, with
the bound $z\,60^4b^4<B/2$ preventing carries. This proves soundness.

For necessity, choose a power-of-two $b$ bounding a solution, take
$B=Hb^4$, $q=B^L$, and use the intended codes. Since $x\le b-1$,
$1+xB\le bB\le bB^M$, and hence $C<2bB^M$. Therefore
\[
 e\ell C^2<16zb^2B^{L+3M}
       =16zb^2B^{8M}<B^{8M+1}<B^{10M}=q^2.
\]
Here $H>16z$ and $2M>1$ justify the last two inequalities. Thus the
new $\alpha$ is positive. The polynomial $F$ has nonnegative digits
and positive degree, so $F(B)>F(Z)=V$; its difference is divisible by
$B-Z$, giving the positive integer $t$ in
\eqref{eq:packed-congruence}. All other witnesses follow from the same
digit and Pell necessity construction.

Finally, $C^2$ was already computed to obtain $C^4$ in (E7), so
\eqref{eq:packed-bound} removes the separate squaring of $g$. Equation
\eqref{eq:packed-congruence} needs one multiplication and one addition
instead of the two of each in (E4), (E5); its left side was already
computed as the second S-block. These three savings reduce 123 to 120.
The program schedules the inequality check after $C^2$ has been
computed, verifies all twenty residuals, and records the explicit
dependence on (E2), (E3) for the transformed (E7).
\end{proof}

\medskip
\noindent The progression distinguishes a shorter certificate for the
same system from a change of slack variable and from a change of encoding:
\begin{center}
\begin{tabular}{@{}lrrr@{}}
\toprule
Construction & Multiplications & Additions & Total\\
\midrule
Archived certificate for $\Sigma$ &75&54&129\\
Shared product $UM$ &74&54&128\\
Also reuse (E2), (E3): same $\Sigma$ &73&54&127\\
Pell gap removed: $\Sigma_{\rm flat}$ &73&53&126\\
Radix index and split bound: $\Sigma_{\rm rec}$ &70&53&123\\
Packed index and shared square: $\Sigma_{\rm pack}$ &68&52&120\\
\bottomrule
\end{tabular}
\end{center}
These are proved upper bounds under Definition~\ref{def:certificate};
no minimum is asserted. In particular, the successful reductions replace
the earlier suggestion that algebraic rewriting cannot materially close
the gap. The 120-operation result uses free integer numerals and
admissible index parameters. It does not include the computation of the
index $(Z,V,H)$, nor claim a 120-operation bound when every such integer
must itself be generated from~$1$. If the parameters $x,Z,V,H$ remain
supplied but the numerical constants $2,4,5,L$ must be generated from~$1$,
the checked addition chain for $L=5^{59}$ gives eleven additional
instructions and a total of 131.

The independently executable companions are
\begin{center}\small
\begin{tabular}{@{}l@{}}
\file{round4\_1980\_optimized\_certificate.py},\\
\file{round4\_1980\_recoded\_certificate.py},\\
\file{round4\_1980\_packed\_certificate.py}.
\end{tabular}
\end{center}
The receipt of the second of these records the 126-instruction count of
the recoding alone; the 123-instruction count of Theorem~\ref{thm:123}
is the entry \texttt{recoded\_system\_123} of the first receipt.
They preserve the original 129-step schedule, verify the new schedules
by exact symbolic identities, and write their operation histograms and
residual checks as JSON receipts.
The positive-domain arguments above supply the mathematical implications
that polynomial expansion alone cannot certify.

% Independent Pell improvements; input by the satellite article.
\section{Further reductions in the Pell block}\label{sec:further-pell}

Two changes to the Pell block save four operations independently of the
digit encoding: a shorter parameter in the final Pell equation saves two,
and a pair of linear interval equations saves another two. All unknowns
remain positive integers. Signed intermediate expressions retain the
auxiliary domain of Definition~\ref{def:certificate}.

\subsection{Factoring the final Pell parameter}

Write $J=2r+1$. The remark after Lemma~2.28 of~\cite{J82} explicitly
permits replacing the parameter $a+f^2(d^2-a)$ by
\[
 G_0=a+f^2(f^2-a).
\]
This already saves one operation: if $X=(a^2-1)(ic^2)^2$, then (E16)
asserts $f^2=X+1$, and
\[
 G_0-1=(f^2-1)(f^2-a+1).
\]
The four instructions
\[
 T=f^2-(a-1),\quad G_-=XT,\quad G_+=G_-+2,
 \quad D_G=G_-G_+
\]
replace five instructions for $G_0^2-1$; both $X$ and $a-1$ are already
needed elsewhere.

This replacement preserves positive witnesses constructively. Once
$c=\psi_a(J)$ and $d=\chi_a(J)$, keep positive $f,i$ satisfying (E16).
For either the original parameter or $G_0$, one has
$G\equiv a\pmod f$, $G\equiv1\pmod c$, and $G\ge2a+1$.
Set $I=\chi_G(J)$ and $H=\psi_G(J)$. The Pell recurrence polynomials give
$I\equiv d\pmod f$, $H\equiv J\pmod c$, while growth gives $I>d$ and
$H>J$. Thus $o=(I-d)/f$ and $j=(H-J)/c$ are positive integers.
This explicitly supplies the required positive first coordinate; it does
not assume that an arbitrary signed square root in Lemma~2.28 is positive.

A signed congruence permits a further reduction. Replace (E17) by
\begin{equation}\label{eq:signed-pell-new}
 (of-d)^2=(G^2-1)(J+jc)^2+1,
 \qquad G=1+(a+1)(f^2-1).
\end{equation}
The variable $o$ is positive; only the expression $of-d$ may have either
sign.

\begin{lemma}\label{lem:signed-chi-new}
For $a>1$ and $0<k\le m$, the congruence
$\chi_a(n)\equiv\pm\chi_a(k)\pmod{\chi_a(m)}$ implies
$n\equiv\pm k\pmod{2m}$.
\end{lemma}
\begin{proof}
The positive-sign case is the usual $\chi$ step-down lemma quoted after
Lemma~2.28 of~\cite{J82}, with its stronger modulus $4m$. For the negative
sign, the Pell addition identities give
\[
 \chi_a(n+2m)\equiv-\chi_a(n)\pmod{\chi_a(m)},
\]
because $\chi_a(2m)=2\chi_a(m)^2-1$ and
$\psi_a(2m)=2\chi_a(m)\psi_a(m)$. Apply the positive-sign case at
$n+2m$ and reduce modulo $2m$.
\end{proof}

\begin{proposition}\label{prop:signed-pell-new}\label{lem:signed-pell}
Under $a>1$ and $1<J<c$ with $J$ odd, equations (E15), (E16), and
\eqref{eq:signed-pell-new} are existentially equivalent to
$c=\psi_a(J)$, with every system unknown positive.
\end{proposition}
\begin{proof}
For sufficiency, write
\[
 d=\chi_a(k),\quad c=\psi_a(k),\qquad
 f=\chi_a(m),\quad ic^2=\psi_a(m).
\]
The standard Pell divisibility fact used in the proof of Lemma~2.28,
$\psi_a(k)^2\mid\psi_a(m)\Rightarrow\psi_a(k)\mid m$, gives $c\mid m$;
in particular $0<k<c\le m$. Our parameter satisfies
\[
 G>1,\qquad G\equiv-a\pmod f,\qquad G\equiv1\pmod c.
\]
Put $I=|of-d|$ and $H=J+jc$. The last Pell equation gives a positive
index $n$ with $I=\chi_G(n)$ and $H=\psi_G(n)$. Since
$I\equiv\pm d\pmod f$ and
$\chi_{-a}(n)=(-1)^n\chi_a(n)$, Lemma~\ref{lem:signed-chi-new} gives
$n\equiv\pm k\pmod{2m}$, hence modulo $c$. On the other hand,
$G\equiv1\pmod c$ gives $H\equiv n\pmod c$, so $J\equiv\pm k\pmod c$.
The positive-sign case yields $J=k$. The other case yields $J+k=c$,
impossible because $c=\psi_a(k)\equiv k\pmod2$ and $J$ is odd.

For necessity, choose a positive Pell denominator divisible by $c^2$,
as in the construction of Lemma~2.28 of~\cite{J82}; this supplies positive
$f,i$ satisfying (E16). Put
$I=\chi_G(J)$, $H=\psi_G(J)$. Since $J$ is odd,
$I\equiv-d\pmod f$ and $H\equiv J\pmod c$. Therefore
$o=(I+d)/f>0$ and $j=(H-J)/c>0$ supply the desired witnesses.
\end{proof}

For its certificate, compute $a_-=a-1$, $a_+=a+1$, and
$A=a_-a_+$ in place of the three existing instructions jointly needed
for $a^2-1$ and $a-1$. The parameter coefficient then takes only
\[
 G_-=a_+X,\qquad G_+=G_-+2,\qquad D_G=G_-G_+.
\]
Thus \eqref{eq:signed-pell-new} saves two operations relative to the
original five-instruction parameter computation. Substituting $X=f^2-1$
uses (E16): the verifier records the exact triangular correction
\[
 F_{17}^{\rm calc}=F_{17}^{\rm source}
  +F_{16}(a+1)(G_{\rm source}+G_{\rm calc})(J+jc)^2.
\]

\subsection{Replacing the quadratic approximation by an interval}

The factor $4$ in (E10) is unnecessary for these encodings, and the
square can subsequently be removed too. The proof must establish the
evenness used in rounding before appealing to binomial divisibility.
We give that argument explicitly.

Use the notation
\[
 \begin{gathered}
 N=n,\quad R=r,\quad Y=sN^2,\quad U=wN^2,\quad M=RY,\\
 A=a=M(U+1),\quad P=p=2UM^2,\quad K=k,\quad C=c,\quad J=2R+1.
 \end{gathered}
\]
The encoding, independently of the Pell block, establishes
\[
 3<3L\le B<q\le N\le R,\qquad R,N\ge8,\qquad 0<b\le N.
\]
Moreover $B=Hb^4$ is even because admissibility fixes an even $H$.
This does not yet assert that $q$ or $N$ is even. Both $N=q^{16}$ and
the shorter encoding $N=q^8$ satisfy the argument below.

\begin{lemma}\label{lem:weak-pell-rounding-new}
Replace (E10) temporarily by
$(C-KY)^2+\eta=K^2$ with $\eta>0$, retaining the other equations.
Then all conclusions of the original Pell block still hold, and in fact
$|C/K-Y|<1/2$.
\end{lemma}
\begin{proof}
We have $U,Y\ge64$, $M\ge512$, $A\ge33280$, and
\begin{equation}\label{eq:weak-error-new}
 |C/K-Y|<1.
\end{equation}
Equation (E11) gives $K\ge R+2$, so
$C>(Y-1)K\ge63(R+2)>J$. Therefore either the source Pell construction
or Proposition~\ref{prop:signed-pell-new} gives $C=\psi_A(J)$.
By (E9), (E11), and the elementary Pell congruence,
\[
 K=\psi_P(R+1+v(P-1)),\qquad v\ge0.
\]
If $v\ge1$, Lemma~2.19 of~\cite{J82} gives
\[
 \frac CK\le
 \frac{(2A)^{2R}}{(2P-1)^{R+P-1}}
 \le\left(\frac{2A}{2P-1}\right)^{2R}<\frac12,
\]
contradicting $C/K>Y-1\ge63$. Hence $K=\psi_P(R+1)$.

Set $\xi=(U+1)^{2R}/U^R$. The elementary estimates in the proof of
Lemma~2.25 now apply to these exact Pell indices and give
\begin{equation}\label{eq:pell-ratio-new}
 \xi\left(1-\frac{R}{M(U+1)}\right)<\frac CK
 <\xi\left(1+\frac{R}{2M^2U}\right).
\end{equation}
In particular
$C/K>\xi/2>U^R/2>U^{R-1}+1$. Combining this with
\eqref{eq:weak-error-new} yields $Y>U^{R-1}$ and therefore
\begin{equation}\label{eq:a-growth-new}
 A=RY(U+1)>RU^R.
\end{equation}

Put $W=bw$. We have
$2^{3J}\le8N^{2R}\le RU^R<A$ and $W^3\le U^R<A$.
Equation (E14), with $C=\psi_A(J)$ and $d=\chi_A(J)$, is exactly the
$\chi$ congruence of Lemma~2.22. Thus $W=2^J$, which implies that
$b$ is a power of two and
\begin{equation}\label{eq:u-growth-new}
 U=N^2w\ge Nbw=NW\ge8\cdot2^J>4\cdot2^{2R}.
\end{equation}
Next, $\kappa<C$, (E19), and (E20), together with $0<L<J<A$, give
$\kappa=\psi_A(L)$ by Lemma~2.23. By \eqref{eq:a-growth-new},
\[
 B^{3L}\le B^B\le N^N\le U^R<A,
 \qquad q^3\le N^N\le U^R<A.
\]
Equation (E18) and the $\chi$ form of Lemma~2.22 therefore give $q=B^L$.
Since $B$ is even, $q$ and $N$ are now even. These two exponential
arguments precede the rounding step; they do not invoke Lemma~2.26
with a weakened hypothesis.

From $\xi/2<C/K<Y+1$ we get $\xi<2Y+2<3Y$. Equation
\eqref{eq:pell-ratio-new} therefore gives
\[
 \left|\frac CK-\xi\right|
 <\frac{R\xi}{M(U+1)}
 =\frac{\xi}{Y(U+1)}<\frac3{65}<\frac14.
\]
The binomial expansion of Lemma~2.24 and \eqref{eq:u-growth-new} give
\[
 0\le\xi-\lfloor\xi\rfloor<\frac14,
 \qquad \lfloor\xi\rfloor\equiv\binom{2R}{R}\pmod U.
\]
Hence $|Y-\lfloor\xi\rfloor|<3/2<2$. Both integers are even:
$Y=sN^2$, $U=wN^2$, and
$\binom{2R}{R}=2\binom{2R-1}{R-1}$. Thus $Y=\lfloor\xi\rfloor$.
This proves $N^2\mid\binom{2R}{R}$, and the last two error estimates
also give $|C/K-Y|<1/2$, restoring the original stronger bound.
\end{proof}

\begin{proposition}\label{prop:linear-pell-interval-new}\label{lem:linear-ratio}
In the admissible encodings considered here, (E10) may be replaced by
\begin{equation}\label{eq:linear-pell-interval-new}
 c=ksn^2+\eta,\qquad k=\eta+\zeta,
 \qquad \eta,\zeta>0.
\end{equation}
This replacement saves two operations and adds one positive unknown
and one free equality test.
\end{proposition}
\begin{proof}
The two equations give $0<C/K-Y<1$, so Lemma~\ref{lem:weak-pell-rounding-new}
proves sufficiency. For necessity, an original solution satisfies the
weaker bound too: $K^2-(C-KY)^2$ is a positive integer and supplies its
slack. Applying that lemma gives
$Y=\lfloor\xi\rfloor$ and $\xi<Y+1/4<2Y$. The lower estimate in
\eqref{eq:pell-ratio-new} implies
\[
 \xi-\frac CK<\frac{\xi}{Y(U+1)}<\frac2{U+1}<\frac2U.
\]
The positive term of order $1/U$ in the binomial expansion gives
\[
 \xi-Y\ge\frac{\binom{2R}{R-1}}{U}
 \ge\frac{2R}{U}>\frac2U.
\]
Consequently $C/K>Y$. Together with $|C/K-Y|<1/2$, this proves that
$\eta=C-KY$ and $\zeta=K-\eta$ are positive integers.
The certificate computes $ksn^2$, adds $\eta$, and computes
$\eta+\zeta$: three operations in place of the previous five.
\end{proof}

The proofs above preserve the other positive witnesses; $j,o$ may be
re-chosen for the new Pell parameter, and $\eta,\zeta$ supply the new
interval. The residual verifier checks all changed equations and the
two triangular identities exactly. Positivity and the equivalence of
the encodings depend on the mathematical arguments, not on polynomial
expansion alone.

% New input section; independent successor to the 120-step construction.
\section{A coefficient normalization allowing shorter masks}
\label{sec:short-masks}

The 120-step system still pays for two powers that can be removed by
changing the coefficient encoding. Its first coefficient code has a
nonzero digit in position $L$, and the packing shifts the second code by
$q^2$. We construct a quartic representation whose digit in position $L$
is zero. Both coefficient codes can then be bounded by $q$, their packing
uses a shift by $q$, and the block widths can be reduced to
$q,q^2,q^5$. Their product is $q^8$; only $q^2,q^4,q^8$ need be computed.

The normalization is part of the admissible index, not an assumption
about an arbitrary polynomial. In particular, simply deleting the leading
digit of the old code would change the coefficient being tested. The
construction below makes that digit zero while preserving the represented
set.

\subsection{Making the leading coefficient digit vanish}

Start with the quadratic-system representation underlying the universal
pair $(58,4)$, and shift its witnesses to nonnegative coordinates. Write
its residuals as $F_i(X,z_1,\ldots,z_{58})$, with constant coefficients
$c_i$. Introduce a nonnegative witness $h$ and put
\[
 F_i^h=F_i-c_i+c_ih,
 \qquad
 S=\sum_i(F_i^h)^2+(Xh-X)^2.
\]
The polynomial $S$ has degree at most four and zero constant coefficient.
For $X>0$, $S=0$ forces $h=1$ and hence is equivalent to the original
system. Introduce a further nonnegative witness $\sigma$ and choose a
power of two $z\ge4$. The polynomial to be encoded is
\begin{equation}
 R=24S+z(\sigma-1).\label{eq:short-normalization}
\end{equation}
For nonnegative integral witnesses, $R=0$ is equivalent to $S=0$ and
$\sigma=1$: values $\sigma\ge2$ make $R$ positive, while $\sigma=0$
would require $24S=z$, impossible because $3$ does not divide a power of
two. Thus this remains a universal quartic representation for positive
inputs. It now uses $\nu=60$ witness coordinates, including $h,\sigma$.
These are coordinates encoded in $g$, not additional unknowns of the
final twenty-equation system.
The forced coordinate $\sigma=1$ ensures that the actual witness code
$g$ is positive.

For multiindices $I=(i_0,\ldots,i_\nu)$ of total degree at most four,
put
\[
 c_I=\frac{4!}{i_0!\cdots i_\nu!(4-\sum_j i_j)!},
 \qquad R=\sum_I c_IP_I X^{i_0}z_1^{i_1}\cdots z_\nu^{i_\nu}.
\]
Every $c_I$ divides 24. Consequently the contribution of $24S$ to every
$P_I$ is integral; the new linear coefficient of $\sigma$ contributes
$z/4$, also integral. Choose $z$ so large that the fixed contributions
from $24S$ have absolute value less than $z$. The variable $\sigma$ is
fresh and hence these contributions do not depend on $z$. We obtain
\begin{equation}
 P_{\boldsymbol0}=-z,
 \qquad -z<P_I<z\quad(I\ne\boldsymbol0).
 \label{eq:short-coefficients}
\end{equation}
The convention here encodes $R$ directly as $\sum c_IP_I$ times the
monomials. There is no further factor of 24 in the extracted coefficient.

Set
\[
 Z=2z,\qquad M=5^{60},\qquad L=5^{64}=625M,
 \qquad w(I)=i_0+i_1 5+\cdots+i_\nu5^\nu.
\]
Then $w(I)\le4M<L$, and distinct $I$ have distinct weights by uniqueness
of base-five digits. Define
\begin{align*}
 \ell_0(B)&=\sum_{j=1}^{\nu}B^{5^j},\\
 e_0(B)&=\sum_I(z+P_I)B^{L-w(I)},\\
 V&=e_0(Z)+\ell_0(Z)Z^L.
\end{align*}
The leading digit of $e_0$ vanishes by
\eqref{eq:short-coefficients}, so $\deg e_0\le L-1$. All its other
digits lie in $[0,Z)$; the two parts of the packed polynomial
$e_0(B)+\ell_0(B)B^L$ have disjoint positions.

Choose a power of two $H$ with
\begin{equation}
 H>\max\{2Z^{2L+1},\ Z(\nu+2)^4,\ 3L,\ 64\}.
 \label{eq:short-H}
\end{equation}
An index $(Z,V,H)$ is now called admissible if constructed by these rules
from a quadratic-system representation. Its three components remain
fixed parameters. The complete parameter list is $x,Z,V,H$, as in the
preceding packed system, although this is a different class of admissible
indices.

\subsection{The shortened system and its certificate}

Keep the 32 positive unknowns of $\Sigma_{\rm pack}$ and the
abbreviations $B=Hb^4$, $C=1+xB+g$. Replace the coding equations by
\begin{align}
 b&=x+\beta,& e+\ell C^4+\alpha&=q,\label{eq:short-bound}\\
 \lambda+q^2&=1+\lambda B,&\theta+Z&=B,\label{eq:short-geometric}\\
 e+\ell q&=V+t\theta,&n&=q^8.\label{eq:short-packed}
\end{align}
Use the following abbreviations in the equation defining $r$:
\begin{align*}
 S_3&=(2e-Z\lambda)C^4+B\lambda(1+q^2),\\
 S&=g+q(e+\ell q+q^2S_3),\\
 T+1&=q-(b-1)\ell+\theta\lambda q+(B-2)q^4,\\
 r&=S(n^2-n)+(T+1)(n^2-1).
\end{align*}
Retain the remaining Pell equations, including $c=\kappa+\varphi$,
and use the new fixed exponent $L$ in
$\kappa=L+\Delta(a-1)$. The exponent-replacement congruence still uses
base $B$. There are twenty equations; all abbreviations are charged in
the certificate.

\begin{theorem}\label{thm:short118}
For every admissible index just defined, the shortened system represents
the same recursively enumerable set as the quadratic system from which
the index was constructed. It has a certificate of 118 instructions:
65 multiplications and 53 additions, followed by twenty equality tests.
\end{theorem}

\begin{proof}
We first verify the hypotheses needed to apply the Pell lemmas without
assuming the exponential equation. Equation~\eqref{eq:short-bound}
gives
\[
 e<q,\qquad \ell<q,\qquad C^4<q,
 \qquad (b-1)\ell<\ell C^4<q.
\]
Since $C>B>b\ge2$, it follows that $q>B$ and $g<q$. Equation
\eqref{eq:short-geometric} gives $\lambda>q$, so, writing $z=Z/2$,
\[
 0<D_0=z\lambda-e<z\lambda,
 \qquad 2C^4D_0<2z\lambda q<B\lambda q.
\]
The blocks therefore satisfy
\begin{gather*}
 0<g<q,\qquad 0\le q-1-(b-1)\ell<q,\\
 0<e+\ell q<q^2,\qquad
 0<(B-Z)\lambda<(B-1)\lambda=q^2-1,\\
 0<S_3<B\lambda(1+q^2)
      =\frac{B(q^4-1)}{B-1}<2q^4<q^5,\\
 0<(B-2)q<q^2<q^5.
\end{gather*}
Thus their respective widths are $N_1=q$, $N_2=q^2$, $N_3=q^5$;
their product is $n=q^8$. The equation for $r$ gives $r\ge n^2-1\ge n$,
and \eqref{eq:short-H} yields
$3<3L\le B<q\le n\le r$. The elementary Pell gap restores the
stronger inequality required by Lemma~\ref{lem:226}. The Pell lemmas
now give $b\pow2$, $q=B^L$, and the central-binomial divisibility.
Since $H$ is a power of two, the block radices are powers of two, and the
three separate no-carry conditions follow.

Apply \cite[Lemma~2.9]{J82} to $Y=e+\ell q$ with digit base $Z$ and
$m=k=2L$. The polynomial
\[
 F(B)=e_0(B)+\ell_0(B)B^L
\]
has degree at most $L+M<2L$, digits less than $Z$, and $F(Z)=V$.
The congruence in \eqref{eq:short-packed}, the bound $Y<q^2=B^{2L}$,
the mask $(B-Z)\lambda$, and \eqref{eq:short-H} supply the lemma's
hypotheses. Hence $Y=F(B)$. The supplied and intended $e,\ell$ all lie
in $[0,q)$, so quotient-and-remainder uniqueness gives $e=e_0(B)$ and
$\ell=\ell_0(B)$. The first mask identifies the witness digits of $g$
and bounds them by $b$.

To identify what the last mask tests, regard $B$ temporarily as a formal
base. The coefficient of $B^L$ in
\[
 C^4(e_0(B)-z\lambda)
\]
is exactly $R$ evaluated at the encoded tuple. Indeed, the monomial
exponent $w(I)$ in $C^4$ has coefficient $c_I$ times its monomial value;
the complementary digit of $e_0-z\lambda$ is $P_I$. This includes
$I=\boldsymbol0$, where that digit is $-z$, precisely the constant
coefficient in \eqref{eq:short-normalization}.

Every coefficient of $e_0-z\lambda$ has absolute value at most $z$.
Thus every coefficient of the product has absolute value less than
$z((\nu+2)b)^4<B/2$. Its degree is at most $2L-1+4M<3L$.
The offset $(B/2)\lambda(1+q^2)$ supplies the digit $B/2$ at every
position below $4L$, preventing carries throughout. The third no-carry
condition therefore tests, by \cite[Lemma~2.8]{J82}, whether the
coefficient $R$ is zero. The equivalence following
\eqref{eq:short-normalization} proves soundness.

Conversely, take a solution of the original quadratic system, set
$h=\sigma=1$, and choose a power-of-two $b$ greater than the input and
all sixty witnesses. Form the intended codes using $B=Hb^4$ and
$q=B^L$. Their bounds are
\[
 e_0(B)<ZB^{L-1},\qquad \ell_0(B)<2B^M,
 \qquad C<2bB^M.
\]
Because $B>2Z$, the first is less than $q/2$. Also
\[
 \ell C^4<32b^4B^{5M}
       <\tfrac12 B^{5M+1}\le\tfrac12q,
\]
using $H>64$ and $L=625M\ge5M+1$. Consequently the new
$\alpha=q-e-\ell C^4$ is positive. The packed polynomial $F$ has
nonnegative digits and positive degree, so $F(B)>F(Z)$ and
$t=(F(B)-F(Z))/(B-Z)>0$. The coordinate $\sigma=1$ ensures $g>0$.
All remaining witnesses follow from the same digit and Pell necessity
constructions with the new base and exponent.

For the count, replace the three instructions computing $e\ell$,
$(e\ell)C^2$ and addition of $\alpha$ by the three instructions
computing $\ell C^4$, addition of $e$, and addition of $\alpha$.
This exchanges one multiplication for one addition but preserves the
length. All packing formulas retain their previous instruction counts.
The power chain now consists only of $q^2,q^4,q^8$; the instructions
for $q^3$ and $q^{16}$ disappear. This reduces 120 to 118.

The exact verifier checks all twenty residuals. Its shared (E7)
calculation records the adapted identity
\[
 F_{7,\mathrm{new}}=F_7+(\lambda F_3-F_2)q(n^2-1),
\]
where $F_2,F_3$ now refer to
\eqref{eq:short-geometric}. Thus the symbolic check accounts for the
dependence on the preceding equations as well as the operation count.
\end{proof}

\begin{corollary}\label{thm:best114}
With the same admissible index $(Z,V,H)$, the shortened system together
with the signed Pell replacement of Proposition~\ref{lem:signed-pell} and the
linear ratio replacement of Proposition~\ref{lem:linear-ratio} has a
114-instruction certificate: 63 multiplications and 51 additions,
followed by 21 equality tests. It has 33 positive unknowns and the four
parameters $x,Z,V,H$, and defines the same recursively enumerable set.
\end{corollary}

\begin{proof}
Both Pell replacements apply to the shortened system: the relevant
equations retain their form, and the bounds $q>B>b\ge2$, $n=q^8\ge2$
and $a\ge20r$ supply their hypotheses. Their positive-domain equivalences
therefore compose with Theorem~\ref{thm:short118}. The final certificate
is independently checked by
\file{round5\_1980\_certificate.py}; its exact residual checks and
instruction histogram give the asserted 114 instructions and 21 tests.
The linear ratio replacement adds one positive unknown and one equation.
\end{proof}

For comparison, applying the three independent changes in this order gives:
\begin{center}
\begin{tabular}{@{}lrrr@{}}
\toprule
Construction & Multiplications & Additions & Total\\
\midrule
Packed encoding of Theorem~\ref{thm:120} & 68 & 52 & 120\\
Shorter masks & 65 & 53 & 118\\
Also the signed Pell parameter & 65 & 51 & 116\\
Also the positive interval & 63 & 51 & 114\\
\bottomrule
\end{tabular}
\end{center}

This is a further verified upper bound, not an assertion of minimality.
The choice $L=5^{64}$ also makes the fixed numeral inexpensive to
generate: after $2=1+1$, $4=2+2$, $5=4+1$, six successive squarings
give $5^2,5^4,\ldots,5^{64}$. Thus, if $x,Z,V,H$ remain supplied
parameters while $2,4,5,L$ must be generated from $1$, nine additional
instructions suffice, giving a total of 123.

% Direct quadratic coefficient tests; preserves the earlier 114-operation milestone.
\section{Testing quadratic residuals with one mask}\label{sec:quadratic-masks}

The sum of squares used to obtain a quartic polynomial is not needed for
coefficient testing. We can test the original quadratic residuals at
separate digit positions, using the already decoded variable mask as the
test mask. Extra digit coordinates make this possible without extra
unknowns in the final polynomial system.

\subsection{A fixed layout for all quadratic systems}

Start from the quadratic-system representation in 58 existential
variables underlying the universal pair $(58,4)$ in~\cite[\S5]{J82}.
Shift its positive variables to nonnegative variables
$X_1,\ldots,X_{58}$, and put $X_0=x$. The integer quadratic residuals
have fixed coefficients for the represented r.e.\ set. Their rational
linear span has dimension at most
\[
 m=\binom{61}{2}=1830,
\]
the number of monomials of degree at most two in 59 variables.
Choose a basis from the original integer residuals and pad it with zero
polynomials to obtain $F_0,\ldots,F_{m-1}$. Simultaneous vanishing is
unchanged: every original residual is a rational linear combination of
the chosen basis, and every chosen residual is original.

Use the following fixed positions:
\begin{align*}
 v_i&=1+3(6729i+i^2)\quad(0\le i\le58),& v_*&=1180939,\\
 D&=4v_*+1,& T&=(m+1)D+2v_*,\\
 t_j&=T+jD\quad(0\le j<m),& t_*&=T+(m-1)D.
\end{align*}
Let $d_*=T-2v_*$. The inequalities
\begin{equation}\label{eq:quad-layout-bounds}
 D>2v_*,\qquad
 d_*>\max(v_*,t_*-T),\qquad
 L=5^{16}=152587890625>3t_*+1=51873937495
\end{equation}
are exact integer inequalities; the companion program checks them.
The positions $v_i$ are for the input and the 58 actual variables.
The positions $t_j$ are for 1830 additional, unconstrained digit
coordinates and will also be the equation-testing positions.

For a multiindex $I=(I_0,\ldots,I_{58})$ with $|I|\le2$, put
\[
 w(I)=\sum_{i=0}^{58}I_i v_i,
 \qquad c_I=\frac{2!}{I_0!\cdots I_{58}!(2-|I|)!}\in\{1,2\}.
\]
Write $F_j=\sum_I a_{j,I}X^I$, and define
$P_{j,I}=2a_{j,I}/c_I\in\mathbb Z$. Choose a power of two $z\ge2$
strictly larger than every $|P_{j,I}|$, and put $Z=2z$.
Define the polynomial with possibly negative coefficients
\[
 \mathcal D(B)=\sum_{j=0}^{m-1}\sum_{|I|\le2}
        P_{j,I}B^{t_j-w(I)}.
\]
Its exponents are nonnegative. Within one row, distinct multiindices
have distinct weights. Indeed $v_i\equiv1\pmod3$, so the weight
determines the total degree, which is at most two. Degree-one weights
are strictly increasing. For a degree-two weight $v_i+v_j$, subtracting
$2$ and dividing by $3$ gives $6729(i+j)+i^2+j^2$. Since
$0\le i^2+j^2\le2\cdot58^2<6729$, division by $6729$ recovers both
$i+j$ and $i^2+j^2$. These determine the unordered pair $\{i,j\}$.
The supports of different rows lie in disjoint intervals
$[t_j-2v_*,t_j]$, by \eqref{eq:quad-layout-bounds}. Thus every
coefficient of $\mathcal D$ has absolute value less than $z$, and its
support lies in $[d_*,t_*]$.

The two nonnegative digit polynomials are
\begin{align*}
 \ell_0(B)&=\sum_{i=1}^{58}B^{v_i}+\sum_{j=0}^{m-1}B^{t_j},\\
 e_0(B)&=z\sum_{h=0}^{L-1}B^h+\mathcal D(B).
\end{align*}
Every coefficient of $e_0$ lies in $[0,Z)$, and its degree is $L-1$.
The mask $\ell_0$ has coefficients zero or one and degree $t_*<L$.
The admissible index is $(Z,V,H)$, where
\begin{equation}\label{eq:quad-index}
 V=e_0(Z)+\ell_0(Z)Z^L,\qquad
 H>\max\{2Z^{2L+1},\ Z(1890)^2,\ 3L,\ 16\}
\end{equation}
and $H$ is a power of two. There are $58+1830=1888$ non-input digit
coordinates; their number explains the factor $1890=1888+2$.
They are encoded in one integer $g$ and are not new unknowns of the
final polynomial system. All these choices are effective index
construction. The certificate parameters remain $x,Z,V,H$.

\subsection{The changed equations and the bounds before decoding}

Retain the signed Pell parameter and the linear interval equations
proved above. Put $B=Hb^2$ and $\mathcal C=1+xB+g$, with both
abbreviations computed by charged instructions. Replace the bound and
the third testing block by
\begin{align*}
 e+\ell\mathcal C^2+\alpha&=q,\\
 S_3&=(2e-Z\lambda)\mathcal C^2+B\lambda(1+q^2),\\
 T_3&=(B-2)\ell.
\end{align*}
The other coding equations remain
\[
 b=x+\beta,\quad \lambda+q^2=1+\lambda B,\quad
 \theta+Z=B,\quad e+\ell q=V+t\theta,\quad n=q^8.
\]
The first two pairs of blocks are
\[
 S_1=g,\quad T_1=q-1-(b-1)\ell,\qquad
 S_2=e+\ell q,\quad T_2=\theta\lambda,
\]
with widths $N_1=q$, $N_2=q^2$, $N_3=q^5$. Accordingly,
\begin{align*}
 S&=g+q(e+\ell q+q^2S_3),\\
 T+1&=q-(b-1)\ell+\theta\lambda q+(B-2)\ell q^3,\\
 r&=S(n^2-n)+(T+1)(n^2-1).
\end{align*}

These blocks have the required signs and widths before either
$b$ being a power of two or $q=B^L$ is known. Indeed, positivity gives
$e,\ell,\mathcal C^2<q$, while $\mathcal C>B>b\ge2$ gives $q>B$.
The geometric equation yields $\lambda>q$, so
$0<z\lambda-e<z\lambda$. Moreover
\[
 2\mathcal C^2(z\lambda-e)<2z\lambda q<B\lambda q,
 \qquad 0<S_3<B\lambda(1+q^2)<2q^4<q^5.
\]
Since $(b-1)\ell<\ell\mathcal C^2<q$, we have $0\le T_1<q$.
Likewise $0<T_3=(B-2)\ell<\ell\mathcal C^2<q$.
The familiar integer bounds give $0<S_2<q^2$ and
$0<T_2<(B-1)\lambda=q^2-1$, while $0<g<q$.
Consequently $r\ge n^2-1\ge n$, and admissibility gives
\[
 3<3L\le B<q\le n\le r,
 \qquad r,n\ge8,\qquad 0<b\le n.
\]
As $B$ is even already, the noncircular Pell and interval arguments of
Section~\ref{sec:further-pell} apply. They give $b$ a power of two,
$q=B^L$, and $n^2\mid\binom{2r}{r}$. Thus the three separate no-carry
conditions follow from Lemmas~2.16 and~2.11 of~\cite{J82}.

\subsection{Decoding and excluding unwanted coefficients}

Apply Lemma~2.9 of~\cite{J82} to $Y=e+\ell q$, with digit base $Z$
and $m=k=2L$. Its fixed digit polynomial is
$e_0(B)+\ell_0(B)B^L$, of degree $L+t_*<2L$, with coefficients
in $[0,Z)$. Equation~\eqref{eq:quad-index} supplies its value at $Z$;
the packed congruence, $Y<q^2=B^{2L}$, the mask $\theta\lambda$, and
$B>2Z^{2L+1}$ supply the other hypotheses. Therefore
$e+\ell q=e_0(B)+\ell_0(B)q$. Since all four components lie in
$[0,q)$, quotient-and-remainder uniqueness gives
$e=e_0(B)$ and $\ell=\ell_0(B)$.

The first no-carry condition now decodes
\[
 \mathcal C(B)=1+xB+
       \sum_{i=1}^{58}X_i B^{v_i}
       +\sum_{j=0}^{m-1}Y_j B^{t_j},
 \qquad 0\le X_i,Y_j<b.
\]
The $Y_j$ are the unconstrained extra coordinates. They do not affect
the equation tests. To see this, regard $B$ as a formal variable and
write
\[
 e_0-z\lambda
   =\mathcal D-z\sum_{h=L}^{2L-1}B^h.
\]
The tail starts above every testing position, since $t_*<L$.
The remaining factor $\mathcal D$ starts at degree $d_*$.
Any monomial of $\mathcal C^2$ involving a $Y_j$ has degree at least
$T$, so its product with $\mathcal D$ has degree at least
$T+d_*>t_*$. Thus no extra coordinate contributes to any testing
position. Also $d_*>v_*$, so the coefficients at the low positions
$v_i$ are automatically zero.

At a high position $t_j$, a contribution from row $k$ and true-variable
monomials $I,J$ requires
\[
 t_k-w(I)+w(J)=t_j.
\]
Both weights lie in $[0,2v_*]$, so the row spacing forces $k=j$.
Uniqueness of the quadratic monomial weights then forces $I=J$. Therefore
\begin{equation}\label{eq:quad-target-coefficient}
 [B^{t_j}]\,\mathcal C(B)^2(e_0(B)-z\lambda(B))
       =\sum_I c_I P_{j,I}X^I=2F_j(x,X_1,\ldots,X_{58}).
\end{equation}
This proves both kinds of separation: other rows cannot interfere,
and the extra coordinates at the row positions cannot interfere.

Every coefficient of $e_0-z\lambda$ has absolute value at most $z$.
The sum of coefficients of $\mathcal C^2$ is less than
$(1890b)^2$. Hence every coefficient of their product has absolute
value less than
\[
 z(1890b)^2<B/2.
\]
Its degree is at most $2L-1+2t_*<3L$. The offset
$(B/2)\lambda(1+q^2)$ supplies the digit $B/2$ at every position
below $4L$, so all resulting digits lie strictly between zero and $B$.
Now $S_3,T_3$ are even, and their third no-carry condition is equivalent
to that for
\[
 \frac{S_3}{2}\quad\hbox{and}\quad
 \frac{T_3}{2}=(B/2-1)\ell_0(B).
\]
The latter mask tests exactly the low positions $v_i$ and the high
positions $t_j$. At each such position it requires the digit of
$S_3/2$ to equal $B/2$. The low tests hold automatically; by
\eqref{eq:quad-target-coefficient}, the high tests are equivalent to
$F_j=0$ for every $j$. This proves soundness.

\subsection{Positive witnesses and the certificate count}

Conversely, choose a nonnegative solution of the quadratic system and
assign the extra coordinates arbitrarily in $[0,b)$, with one of them
equal to $1$. Choose the power of two $b$ larger than $x$ and all these
coordinates. The extra value $1$ guarantees $g>0$ even if every true
existential coordinate is zero. Take the intended $B,q,e,\ell,g$.
Their bounds are
\[
 e<ZB^{L-1}<q/2,\qquad
 \ell<2B^{t_*},\qquad
 \mathcal C<2bB^{t_*}.
\]
Using $H>16$ and $L>3t_*+1$, we obtain
\[
 \ell\mathcal C^2<8b^2B^{3t_*}
      <\tfrac12B^{3t_*+1}\le q/2.
\]
Thus $\alpha=q-e-\ell\mathcal C^2$ is positive, as is $\beta=b-x$.
The packed quotient $t$ is a positive integer because its fixed
nonnegative digit polynomial has positive degree and $B>Z$.
All decoded no-carry conditions hold by the preceding coefficient
calculation. The remaining positive witnesses follow from the Pell
necessity constructions and the positive interval proof. This proves
necessity with all domains preserved.

\begin{theorem}\label{thm:best113}
For each r.e.\ set, an effective admissible index $(Z,V,H)$ as above
gives a system in 33 positive unknowns and 21 equations such that, for
positive input $x$, membership is equivalent to a valid certificate with
\textbf{113 arithmetic instructions}: 62 multiplications and 51 additions.
Equality tests and supplied parameters are free. Generating the fixed
numerals $2,4,5,L$ from $1$ adds seven instructions, giving 120 under
that convention.
\end{theorem}
\begin{proof}
Relative to the 114-instruction certificate, $B=Hb^2$ saves the
multiplication constructing $b^4$, and $\mathcal C^2$ saves the one
constructing $\mathcal C^4$. Computing the changed third mask in the
packing expression costs one extra multiplication; the other counts
are unchanged. The $T+1$ computation may be organized as
\[
 T+1=q(\lambda+q^2-Z\lambda)
      +(B-2)\bigl((\ell q)q^2\bigr)-(b-1)\ell,
\]
using the geometric and $\theta$ equations and the already computed
product $\ell q$. The exact schedule and
its triangular residual corrections are verified by
\file{round6\_1980\_quadratic\_certificate.py}.
The source system has the same 33 positive unknowns and 21 equations;
the additional digit coordinates are encoded inside $g$.
The equivalence is proved above. Starting from $1$, three additions
produce $2,4,5$ and four successive squarings of $5$ produce $L=5^{16}$,
giving the seven-instruction numeral construction.
\end{proof}

% A further Pell shift, independent of the direct quadratic mask improvement.
\section{Shifting the first Pell parameter}\label{sec:shifted-pell}

One more operation can be removed by shifting the first Pell parameter.
Use the notation of Section~\ref{sec:further-pell}, and put
\[
 X=2UM^2=2ws^2r^2n^6,\qquad P=X+1.
\]
Delete the unknown $p$ and its defining equation (E8), and replace
(E9) and (E11) by
\begin{equation}\label{eq:shifted-pell-system}
 X(X+2)k^2+1=\tau^2,\qquad k=r+1+hX.
\end{equation}
Every other equation is retained, including the positive interval
$c=ksn^2+\eta$, $k=\eta+\zeta$.
The mathematical parameter $P$ need not be computed: its coefficient
is $P^2-1=X(X+2)$, and its index modulus is $P-1=X$.

\begin{lemma}\label{lem:shifted-pell}
Under the initial encoding bounds
\[
 3<3L\le B<q\le N\le R,\qquad R,N\ge8,\qquad b\le N,
\]
with $B$ even and $N=q^8$, the replacement
\eqref{eq:shifted-pell-system} preserves existential equivalence with
positive witnesses. It re-chooses only $k,\tau,h,\eta,\zeta$ and removes
or restores $p$. It applies both to $B=Hb^4$, $L=5^{64}$ and to the
direct quadratic encoding $B=Hb^2$, $L=5^{16}$.
\end{lemma}
\begin{proof}
The initial bounds give $U,Y\ge64$, $M=RY\ge512$, and
$A=M(U+1)\ge33280$. Write $C=c$, $K=k$, $J=2R+1$ as before.
The positive interval says $0<C/K-Y<1$. Equation
\eqref{eq:shifted-pell-system} gives $K\ge R+2$, so
$C>(Y-1)K>J$. Proposition~\ref{prop:signed-pell-new} therefore gives
$C=\psi_A(J)$ and $d=\chi_A(J)$, before any approximation is used.
The other Pell equation and its elementary index congruence give
\[
 K=\psi_P(R+1+vX),\qquad v\ge0.
\]
If $v\ge1$, the Pell growth bounds imply
\[
 \frac CK\le\frac{(2A)^{2R}}{(2P-1)^{R+X}}
 \le\left(\frac{2A}{2P-1}\right)^{2R}<\frac12,
\]
contrary to $C/K>Y-1\ge63$. Here $X\ge R$ and
$2A/(2P-1)=2M(U+1)/(4UM^2+1)<1/2$. Thus
$K=\psi_P(R+1)$.

Put $\xi=(U+1)^{2R}/U^R$. At these exact indices the same growth bounds
give
\begin{align}
 \frac CK
 &\le\frac{(2A)^{2R}}{(2P-1)^R}
   =\xi\left(1+\frac1{4UM^2}\right)^{-R}<\xi,
       \label{eq:shifted-upper}\\
 \frac CK
 &\ge\xi\left(1-\frac1{2A}\right)^{2R}
             \left(1+\frac1X\right)^{-R}
  >\xi\left(1-\frac RA-\frac RX\right).
       \label{eq:shifted-lower}
\end{align}
The last inequality follows from Bernoulli's inequality, with both
factors positive. Since $A/X=(U+1)/(2UM)<1$, the error coefficient is
\[
 \varepsilon=\frac RA+\frac RX
 <\frac{2}{Y(U+1)}<\frac12.
\]
Consequently $C/K>\xi/2>U^R/2>U^{R-1}+1$, while $C/K<Y+1$.
Hence $Y>U^{R-1}$ and $A>RU^R$.

This last bound permits the two exponent arguments in their original
order. First,
\[
 2^{3J}\le8N^{2R}\le RU^R<A,\qquad (bw)^3\le U^R<A,
\]
so (E14) and the $\chi$ form of Lemma~2.22 of~\cite{J82} give
$bw=2^J$, hence
\begin{equation}\label{eq:shifted-u-bound}
 U\ge Nbw\ge8\cdot2^J>4\cdot2^{2R}.
\end{equation}
Next, $\kappa<C$, (E19), (E20), and $0<L<J<A$ give
$\kappa=\psi_A(L)$. The bounds
$B^{3L},q^3\le N^N\le U^R<A$ and (E18) give $q=B^L$ by the same
exponent lemma. Since $B$ is even already, $q$ and $N$ are now even.
Neither argument has used parity rounding or central-binomial
divisibility.

From $\xi/2<C/K<Y+1$ we have $\xi<2Y+2<3Y$, so
\eqref{eq:shifted-upper}--\eqref{eq:shifted-lower} imply
\[
 0<\xi-C/K<\xi\varepsilon
 <\frac6{U+1}\le\frac6{65}<\frac14.
\]
The binomial expansion, using \eqref{eq:shifted-u-bound}, gives
\[
 0<\xi-\lfloor\xi\rfloor<\frac14,
 \qquad \lfloor\xi\rfloor\equiv\binom{2R}{R}\pmod U.
\]
Thus $|Y-\lfloor\xi\rfloor|<3/2<2$. Both integers are even, since
$Y,U$ are even and $\binom{2R}{R}$ is even. Therefore
$Y=\lfloor\xi\rfloor$, so $N^2\mid\binom{2R}{R}$ as required.

It remains to justify strict positivity when changing witnesses.
Start with an old solution and set
\[
 K_{\rm new}=\psi_{X+1}(R+1),\quad
 \tau_{\rm new}=\chi_{X+1}(R+1),\quad
 h_{\rm new}=\frac{K_{\rm new}-R-1}{X}>0.
\]
The quotient is integral by the Pell congruence and positive by strict
growth. The old solution has $Y=\lfloor\xi\rfloor$ and $\xi<2Y$.
Thus the new ratio estimates give
$0<\xi-C/K_{\rm new}<4/U$, whereas the positive fractional tail gives
\[
 \xi-Y\ge\frac{\binom{2R}{R-1}}U\ge\frac{2R}U>\frac4U.
\]
Consequently $Y<C/K_{\rm new}<\xi<Y+1/4$, and
$\eta_{\rm new}=C-K_{\rm new}Y$,
$\zeta_{\rm new}=K_{\rm new}-\eta_{\rm new}$ are positive integers.

Conversely, restore $p=X$ and choose
$K_{\rm old}=\psi_X(R+1)$, $\tau_{\rm old}=\chi_X(R+1)$, and
$h_{\rm old}=(K_{\rm old}-R-1)/(X-1)>0$. The elementary old-base
estimates at these exact indices give
\[
 \xi(1-R/A)<C/K_{\rm old}<\xi(1+R/X),
 \qquad |C/K_{\rm old}-\xi|<2/U.
\]
The same positive tail yields $C/K_{\rm old}>Y$, while its upper
bound gives $C/K_{\rm old}-Y<1/4+2/U<1$. The corresponding old
$\eta,\zeta$ are therefore positive as well. Every other witness is
preserved. The proof uses only the stated initial bounds and evenness
of $B$, which were established for both encodings independently of
the Pell block.
\end{proof}

\begin{theorem}\label{thm:best112}
For each r.e.\ set, the effective admissible index $(Z,V,H)$ of
Theorem~\ref{thm:best113} gives a system in 32 positive unknowns and
20 equations such that, for positive input $x$, membership is equivalent
to a valid straight-line certificate with \textbf{112 arithmetic
instructions}: 62 multiplications and 50 additions. Equality tests and
supplied parameters are free. Constructing $2,4,5,L=5^{16}$ from $1$
adds seven instructions, giving 119 under that convention.
\end{theorem}
\begin{proof}
Apply Lemma~\ref{lem:shifted-pell} to the 113-operation direct quadratic
system of Theorem~\ref{thm:best113}. The shared value $X$ was already
computed by the old (E8) block. Replace the two instructions for
$p^2-1$ by
\[
 X_+=X+2,\qquad D_P=XX_+,
\]
and use $hX$ in place of $h(p-1)$. This deletes one subtraction,
the positive unknown $p$, and its free defining equality. All other
counts are unchanged. The complete primitive schedule, equality list,
declared domains, and exact residual checks are in
\file{round6\_1980\_certificate.py} and its JSON receipt. The standalone
shift and its full positive-witness proof are also recorded in
\file{round5\_1980\_shifted\_pell\_certificate.py} and
\file{PELL\_SHIFTED\_BASE\_PROOF.md}.
\end{proof}

% Reversed code order and high-mask quotient test: 111 arithmetic instructions.
\section{Reversing the packed codes}\label{sec:reversed-packing}

The inequality $e+\ell\mathcal C^2<q$ costs a multiplication and two
additions, including its positive slack. The already computed packed
quantity can replace the product, but its two components must still be
recovered uniquely. Reversing their order and shortening the coefficient
code makes the third mask enforce that uniqueness. The resulting bound
costs two additions.

\subsection{The revised fixed index and system}

Retain the quadratic layout, the polynomial $\mathcal D$, and the mask
$\ell_0$ from Section~\ref{sec:quadratic-masks}. Put $K=t_*+1$ and keep
$L=5^{16}$. The fixed numerical layout satisfies
\begin{equation}\label{eq:reverse-margin}
 L>3K+2.
\end{equation}
Shorten the dense part of the coefficient polynomial and reverse the
order of the two codes:
\begin{equation}\label{eq:reverse-code}
 e_0(B)=z\sum_{j=0}^{K-1}B^j+\mathcal D(B),\qquad
 V=\ell_0(Z)+e_0(Z)Z^L.
\end{equation}
Both $e_0$ and $\ell_0$ now have degree at most $t_*$; their coefficients
are still in $[0,Z)$, with the coefficients of $e_0$ positive.
Choose the fixed power of two $H$ so that
\begin{equation}\label{eq:reverse-index}
 H>\max\{2Z^{2L+1},\ 2^{t_*+4}Z(1890)^2,\ 3L,\ 16\}.
\end{equation}
This includes $H>4Z(1890)^2$ and all earlier radix bounds. These index
choices depend only on the represented r.e.\ set.

Start from the system of Theorem~\ref{thm:best112}, retaining
$B=Hb^2$, $\mathcal C=1+xB+g$, and all its Pell equations. The changed
coding equations and packing expressions are
\begin{align}
 Y&=\ell+eq,&
 Y+\mathcal C^2+\alpha&=q^2,&
 Y&=V+t\theta,\label{eq:reverse-bound}\\
 S_3&=(2e-Z\lambda)\mathcal C^2+B\lambda(1+q),&
 \sigma&=S_3,\label{eq:reverse-sigma}\\
 S&=g+q^2(Y+q^2S_3),\label{eq:reverse-S}\\
 T+1&=q^2(1+\theta\lambda)-(b-1)\ell
                    +(B-2)\ell q^4,\label{eq:reverse-T}\\
 r&=S(n^2-n)+(T+1)(n^2-1).\label{eq:reverse-r}
\end{align}
Here $Y,S_3,S,T$ remain computed abbreviations. The only new unknown is
$\sigma$, required to be positive. Its equality to the already computed
$S_3$ is a free equality test. The unchanged coding equations include
\[
 b=x+\beta,\qquad \lambda+q^2=1+\lambda B,\qquad
 \theta+Z=B,\qquad n=q^8.
\]

\subsection{Bounds before decoding and a possible borrow}

Equation~\eqref{eq:reverse-bound} implies
\[
 \ell<Y<q^2,\qquad e<q,\qquad \mathcal C<q,\qquad g<q.
\]
In particular $q>B$, since $x>0$ and $\mathcal C>B$. With
$N=\theta\lambda$, the geometric equation and the radix bounds give
\[
 q<\lambda<2q^2/B,\qquad b<N<q^2-1.
\]
The new unknown supplies $S_3>0$. Before any decoding we also have
\[
 S_3<2q^3+B\lambda(1+q)<5q^3<q^4.
\]
The three blocks of $S$ therefore have widths $q^2,q^2,q^4$ and give
$0<S<n=q^8$.

The first raw mask $q^2-1-(b-1)\ell$ may be negative at this stage.
Its complete packed value is positive, because
\[
 q^2(1+N)-(b-1)\ell>q^2,
 \qquad (B-2)\ell<q^3.
\]
In addition $T+1<q^7+q^4+q^2<q^8$. Thus $0\le T<n$ and
$r\ge n^2-1\ge n$. The hypotheses used by the unchanged Pell
implications hold without assuming the sign of that first raw mask.
Those implications give $q=B^L$, the required powers of two, and
$\tau_2(S,T)=0$.

To read the masks, write
\[
 d=\left\lfloor\frac{(b-1)\ell}{q^2}\right\rfloor,
 \qquad u=(b-1)\ell-dq^2.
\]
Then $0\le d<b$, $0\le u<q^2$, and the actual block decomposition is
\[
 T=(q^2-1-u)+q^2(N-d)+q^4(B-2)\ell.
\]
All three blocks are in their stated ranges. Every base-$B$ digit of
$N=(B-Z)\lambda$ equals $B-Z$, in positions $0,\ldots,2L-1$.
Since $d<b<B-Z$, subtracting $d$ changes only the unit digit. The middle
no-carry condition consequently makes every nonunit digit of $Y$ less
than $Z$, and makes its unit digit at most $Z+d-1$. Hence
\[
 Y<\frac{Zq^2}{B-1}+b,
 \qquad
 (b-1)\ell<\frac{(b-1)Zq^2}{B-1}+b(b-1)<q^2.
\]
For the last inequality, \eqref{eq:reverse-index} gives
$B-1>4(b-1)Z$ and $q>4b$. Therefore $d=0$ after all.

All digits of $Y$ are now less than $Z$. The digit-evaluation lemma
applied to \eqref{eq:reverse-code} and the packed congruence gives
\begin{equation}\label{eq:reverse-quotient}
 Y=\ell_0(B)+e_0(B)q,\qquad
 \ell=\ell_0(B)+mq,\qquad e=e_0(B)-m
\end{equation}
for an integer $m$ with $0\le m<e_0(B)$. Positivity of $\ell$ and
$\ell_0(B)<q$ exclude negative $m$. The following argument rules out
positive $m$ as well.

\subsection{The high mask checks the quotient}

The fixed support gives $(b-1)\ell_0(B)<q$. The low $q$ block of the
first mask is therefore $q-1-(b-1)\ell_0(B)$. Since $g<q$, its no-carry
condition decodes exactly the earlier low coordinate positions, with
every digit less than $b$. Extra coordinates at positions belonging to
$mq$ cannot occur in $g$, because those positions are at least $L$.
Thus $\mathcal C$ has degree at most $t_*$, constant digit $1$, and input
digit $x$. Writing $A$ for the sum of the coefficients of
$\mathcal C^2$, we obtain
\begin{equation}\label{eq:reverse-A}
 A=\mathcal C(1)^2<(1890b)^2,
 \qquad 4\le ZA<B/4.
\end{equation}
Here $\mathcal C(1)$ denotes evaluation of its digit polynomial, not of
the integer expression $1+xB+g$ with $g$ held constant.

By \eqref{eq:reverse-quotient}, both $e$ and $m$ are smaller than
$e_0(B)<B^K$. The decoded $\mathcal C$ also satisfies $\mathcal C<B^K$.
The explicit margin \eqref{eq:reverse-margin} consequently gives an
integer bound which includes all carries:
\begin{equation}\label{eq:reverse-small-product}
 2e\mathcal C^2<2B^{3K}<B^L=q.
\end{equation}

Put
\[
 X=(B-2)m=\sum_{j=0}^{t_*+1}X_jB^j,
 \qquad 0\le X_j<B.
\]
Since $(B-2)\ell_0(B)<q$, the digits of the third mask at positions
$L,\ldots,L+t_*+1$ are exactly the $X_j$. In this window the coefficient
of $\lambda\mathcal C^2$ is always $A$: the window is above $2t_*$
and below $2L$. By \eqref{eq:reverse-small-product}, the positive integer
$2e\mathcal C^2$ has no digit there. The offset
$B\lambda(1+q)$ has raw coefficient $1$ at position $L$ and raw
coefficient $2$ at the subsequent positions of the window.

The incoming carry at $L$ lies in $\{-1,0,1\}$. Indeed the raw negative
coefficients below $L$ have absolute value at most $ZA<B/4$, while the
positive contribution below $L$ is less than $q+q/(B-1)<2q$.
Thus the normalized digit at $L$ lies between $B-ZA$ and $B-ZA+2$;
because $ZA\ge4$, its outgoing carry is $-1$. The later digits in the
window are $B-ZA+1$ and their outgoing carry remains $-1$. In particular
every digit in this window is at least $B-ZA$.

No carry with the third mask now implies $X_j\le ZA-1$. Consider the
ordinary polynomial $F(U)=\sum_jX_jU^j$. Its value
$F(B)=(B-2)m$ is divisible by $B-2$, so $F(2)$ is divisible by $B-2$.
But \eqref{eq:reverse-index} and \eqref{eq:reverse-A} give
\[
 0\le F(2)<ZA\,2^{t_*+2}
       <Z(1890)^2b^2 2^{t_*+2}<B/4<B-2.
\]
Therefore $F(2)=0$. All coefficients of $F$ are nonnegative, so $X=0$
and $m=0$. We have recovered both canonical codes
$e=e_0(B)$ and $\ell=\ell_0(B)$.

Every target position is below $K$. Below that cutoff,
$e_0-z\lambda=\mathcal D$, and the offset in $S_3/2$ still supplies
the digit $B/2$ at every target. The coefficient bound and the support
separation of Section~\ref{sec:quadratic-masks} apply unchanged.
Carries from overlapping offset blocks above $L$ cannot affect these
lower positions. Thus all original quadratic residuals vanish.

\subsection{Positive witnesses and the verified count}

Conversely, encode a genuine quadratic-system solution, assign one dummy
coordinate the value $1$, and choose a power of two $b$ larger than $x$
and all coordinate values. Use the canonical $e_0(B)$ and $\ell_0(B)$.
Their degrees are at most $t_*$, and \eqref{eq:reverse-margin} gives
$Y+\mathcal C^2<q^2$, with a strictly positive slack.
It also gives $Z\mathcal C^2<q$, so
\[
 S_3\ge\lambda\bigl(B(1+q)-Z\mathcal C^2\bigr)>0.
\]
Set $\sigma=S_3$. The quotient in the packed congruence is positive
because $B>Z$ and its fixed digit polynomial has positive nonconstant
coefficients. The three no-carry conditions hold, and the preceding
Pell necessity constructions supply all remaining positive unknowns.

\begin{theorem}\label{thm:best111}
For each r.e.\ set, an effective admissible index $(Z,V,H)$ defined by
\eqref{eq:reverse-code}--\eqref{eq:reverse-index} gives a system in
33 positive unknowns and 21 equations such that, for positive input $x$,
membership is equivalent to a valid certificate with
\textbf{111 arithmetic instructions}: 61 multiplications and 50 additions.
Equality tests and supplied parameters are free. Generating the fixed
numerals $2,4,5,L$ from $1$ adds seven instructions, giving 118 under
that convention.
\end{theorem}
\begin{proof}
The equivalence, including both positive-witness constructions, was
proved above. Relative to Theorem~\ref{thm:best112}, the product
$\ell\mathcal C^2$, addition of $e$, and addition of the positive slack
are replaced by $Y+\mathcal C^2$ and addition of that slack, saving one
multiplication. The reversed code $\ell+eq$ has the same cost as
$e+\ell q$. Computing $S$ still uses two multiplications and two
additions. The $T+1$ calculation uses
\[
 q^2(\lambda+q^2-Z\lambda)
       -(b-1)\ell+(B-2)(\ell q^4),
\]
which has the same cost as before. Replacing $1+q^2$ by $1+q$ also has
equal cost, and $\sigma=S_3$ is an equality test on an existing computed
register. The exact primitive instructions, equality tests, and
triangular polynomial residual corrections are checked by
\file{round7\_1980\_certificate.py}. The numeral construction is unchanged.
\end{proof}

% Decoded unit coordinate: a successor to the frozen 111-operation system.
\section{Decoding the unit coordinate}\label{sec:unit-digit}

The addition inserting the constant digit in
$\mathcal C=1+xB+g$ can also be removed. We encode a coordinate
$\delta$ in that position, homogenize the quadratic residuals using
$\delta$, and use a slightly less restrictive mask to force
$\delta=1$. The resulting system has 110 arithmetic instructions.
The encoded coordinates introduced in this construction remain digits
of $g$; they do not increase the number of system unknowns.

\subsection{Homogeneous rows and the unit test}

Take the 1830 integer quadratic residuals $F_i(x,X_1,\ldots,X_{58})$
obtained by the basis selection in Section~\ref{sec:quadratic-masks}.
Put $X_0=x$ and homogenize each residual:
\[
 F_i^{\mathrm h}(\delta,X)
   =\sum_{|I|\le2}a_{i,I}\delta^{2-|I|}X^I.
\]
Introduce one more nonnegative encoded coordinate $u=X_{59}$ and
one ordinary row $u\delta-x^2$. There are now 1831 ordinary homogeneous
quadratic rows. Their target coefficients will be twice their values.
Add a final special target whose coefficient is $\delta^2$.
The new mask will permit a target coefficient to be either zero or one.
Thus each ordinary target, being even, must be zero, while the special
target forces $\delta\in\{0,1\}$. The guard $u\delta-x^2=0$ and the
positive input $x$ then force $\delta=1$.

The distinction in scaling is essential: the special target encodes
$\delta^2$, rather than $2\delta^2$. It is not an additional equation
requiring the special value to vanish.

\subsection{The uniform support and admissible index}

Give $\delta$ weight zero. For the input, the 58 original coordinates,
and $u$, use the 60 positive weights
\[
 v_i=1+3(6963i+i^2)\quad(0\le i\le59),
 \qquad v_*=v_{59}=1242895.
\]
There are 1891 homogeneous quadratic monomials in these coordinates
and $\delta$, and all their weights are distinct. Reduction modulo
three distinguishes the number of positive-weight factors. For two
such factors, division by $6963=2\cdot59^2+1$ recovers the sum and
sum of squares of their indices, hence their unordered pair. This is
the same elementary construction used in the preceding sections.

There are $m=1832$ target positions. Define
\begin{align*}
 D_*&=4v_*+1=4971581,\\
 T_*&=(m+1)D_*+2v_*=9115393763,\\
 t_j&=T_*+jD_*\quad(0\le j<m),\\
 t_*&=18218358574,\qquad K=t_*+1=18218358575.
\end{align*}
The exact checker verifies
\begin{equation}\label{eq:unit-layout}
\begin{gathered}
 D_*>2v_*,\qquad T_*-2v_*>v_*,\qquad
 2T_*-2v_*>t_*,\\
 L=5^{16}=152587890625>3K+2=54655075727.
\end{gathered}
\end{equation}
The 1830 original basis rows, 1832 target positions, and 1891 monomial
weights count different objects.

For each ordinary homogeneous quadratic monomial, let $a$ be its
coefficient in its row and $c\in\{1,2\}$ its multinomial coefficient
in a square. Insert the integer $2a/c$ at the complementary exponent
$t_j-w$, where $w$ is the monomial weight. For the special row insert
the coefficient $1$ at $t_{m-1}$, complementary to the weight zero
of $\delta^2$. Let $\mathcal D(B)$ be the resulting signed coefficient
polynomial. Its row intervals are disjoint by
\eqref{eq:unit-layout}; within a row, weight uniqueness prevents
coefficient collisions.

Choose a power of two $z\ge2$ larger than the absolute values of all
these coefficients, and put $Z=2z$. Define
\begin{align*}
 \ell_0(B)&=1+\sum_{i=1}^{59}B^{v_i}
                  +\sum_{j=0}^{m-1}B^{t_j},\\
 e_0(B)&=z\sum_{h=0}^{K-1}B^h+\mathcal D(B),\\
 V&=\ell_0(Z)+e_0(Z)Z^L.
\end{align*}
The coefficients of $\ell_0$ are zero or one, and those of $e_0$ lie
in $[1,Z-1]$. Both codes have degree $t_*$. The mask includes 1892
non-input coordinates: the unit, 59 other genuine coordinates, and
1832 dummy row coordinates. Including the input gives 1893 coordinates;
we use the convenient strict bound $1900b$ for their coefficient sum.

Choose the fixed power of two $H$ so that
\begin{equation}\label{eq:unit-H}
 H>\max\{2Z^{2L+1},\ 4^{t_*+3}Z\,1900^2,\ 3L,\ 16\}.
\end{equation}
This also gives $H>4Z\,1900^2$ and $H>4Z$.
All index construction depends only on the represented r.e.\ set,
independently of the queried input $x$.

\subsection{The changed system and recovery of the codes}

Retain the reversed-packing system of
Section~\ref{sec:reversed-packing}, including its positive unknown
$\sigma=S_3$. Make only the replacements
\[
 \mathcal C=xB+g,\qquad T_3=(B-4)\ell.
\]
Thus its coding and packing expressions are
\begin{align*}
 B&=Hb^2,& Y&=\ell+eq,\\
 Y+\mathcal C^2+\alpha&=q^2,&Y&=V+t\theta,\\
 S_3&=(2e-Z\lambda)\mathcal C^2+B\lambda(1+q),&\sigma&=S_3,\\
 S&=g+q^2(Y+q^2S_3),\\
 T+1&=q^2(1+\theta\lambda)-(b-1)\ell+(B-4)\ell q^4,\\
 r&=S(n^2-n)+(T+1)(n^2-1).
\end{align*}
All the preceding geometric and Pell equations are retained.

The bootstrap remains valid before any unit coordinate has been
decoded. Indeed $\mathcal C=xB+g>B$, since $x,g>0$. The bound still
gives $\mathcal C<q$ and $q>B$. Replacing $B-2$ by the smaller positive
number $B-4$ preserves the complete bounds $0<S<n$ and $0\le T<n$,
including the treatment of a possibly negative first raw mask.
The preceding borrow argument therefore applies unchanged: it first
recovers the Pell conclusions and then proves that the first-block
borrow is zero. The middle mask decodes
\[
 Y=\ell_0(B)+e_0(B)q,
 \qquad \ell=\ell_0(B)+mq,
 \qquad e=e_0(B)-m,
 \qquad 0\le m<e_0(B).
\]

The low first mask restricts the digits of $g$ to the support of
$\ell_0$, each less than $b$. Its unit digit is some
$0\le\delta<b$. Position one is excluded from $\ell_0$, so the input
digit of $\mathcal C=xB+g$ is exactly $x$. The digit polynomial
$\mathcal C(B)$ has degree at most $t_*$. Write
\[
 \mathcal A=\mathcal C(1)^2,
 \qquad 1\le\mathcal A<1900^2b^2,
 \qquad 4\le Z\mathcal A<B/4.
\]
The lower bound follows from $x>0$ and does not assume $\delta=1$.
Also $e,\mathcal C<B^K$, so $2e\mathcal C^2<2B^{3K}<q$.

To remove the remaining quotient ambiguity, put
\[
 X=(B-4)m=\sum_{j=0}^{t_*+1}X_jB^j,
 \qquad 0\le X_j<B.
\]
The term $(B-4)\ell_0$ is less than $q$. Thus the digits $X_j$
occupy exactly the high third-mask positions $L$ through $L+t_*+1$.
The raw-coefficient and carry calculation in
Section~\ref{sec:reversed-packing} gives digits of $S_3$ at least
$B-Z\mathcal A$ throughout this window. It applies unchanged, because
the formula for $S_3$ is unchanged and its hypotheses were just proved.
The third no-carry condition consequently implies $X_j\le Z\mathcal A-1$.

Consider the digit polynomial $F(U)=\sum_jX_jU^j$.
Since $F(B)=(B-4)m$, the integer $F(4)$ is divisible by $B-4$.
But \eqref{eq:unit-H} gives
\[
 0\le F(4)<Z\mathcal A\,4^{t_*+2}
       <Z\,1900^2b^2\,4^{t_*+2}<B/4<B-4.
\]
Hence $F(4)=0$. Nonnegativity of its coefficients forces every $X_j$
to be zero, so $m=0$. We have recovered both canonical codes.
The evaluation point four is required here by the changed mask;
the earlier evaluation-at-two argument is not used for this step.

\subsection{The relaxed mask forces the unit}

At every tested position, which is below $K$, the residual factor
$e_0-z\lambda$ agrees with $\mathcal D$. The support separation
in \eqref{eq:unit-layout} proves that dummy row coordinates cannot
contribute to any target. Distinct rows cannot interfere, and
quadratic weight uniqueness isolates the intended monomial within
each row. The coefficient of
$\mathcal C^2(e_0-z\lambda)$ is therefore zero at every low mask
position, twice the homogeneous residual at each ordinary target,
and $\delta^2$ at the special target.

Every coefficient of this product has absolute value at most
$z\mathcal A<B/4$. Below position $L$, the offset in $S_3/2$
supplies the digit $B/2$ at every position. The resulting digits
lie strictly between $B/4$ and $3B/4$, without carries between them.
The even third block and mask can both be divided by two. The mask
then has digit
\[
 B/2-2
\]
at each selected position. Its only zero bits within a base-$B$
digit are the unit bit and the highest bit. A digit in
$(B/4,3B/4)$ with no overlap with this mask must consequently be
$B/2$ or $B/2+1$. Its centered coefficient is zero or one.

Every ordinary target coefficient is even, so every ordinary row
vanishes. At the special target, $\delta^2\in\{0,1\}$ gives
$\delta\in\{0,1\}$. The ordinary guard $u\delta-x^2=0$ excludes
$\delta=0$ because $x>0$. Thus $\delta=1$, and the original quadratic
residuals all vanish. This proves soundness without an additional
system equation imposing the unit coordinate.

Conversely, given genuine witnesses, set $\delta=1$ and $u=x^2$,
and set all dummy row coordinates to zero. Choose the power of two
$b$ greater than $x$, $u$, and every original coordinate. The unit
digit guarantees $g>0$. Use the canonical code values. The degree
margin and \eqref{eq:unit-H} give
$Y+\mathcal C^2<q^2$ and $Z\mathcal C^2<q$, hence positive
$\alpha$ and $\sigma=S_3$. The packed quotient is positive, as
before. Ordinary target coefficients are zero and the special
coefficient is one, so all masks hold. The preceding Pell necessity
constructions furnish the remaining positive unknowns.

\begin{theorem}\label{thm:best110}
For each r.e.\ set, the effective admissible index $(Z,V,H)$ above
gives a system in 33 positive unknowns and 21 equations equivalent
to membership of the positive input $x$. Its straight-line certificate
has \textbf{110 arithmetic instructions}: 61 multiplications and
49 additions. Charging the seven fixed-numeral construction instructions
gives 117 operations.
\end{theorem}
\begin{proof}
The equivalence has just been proved. Relative to the 111-instruction
schedule, replacing $1+xB+g$ by $xB+g$ removes one addition.
Replacing $B-2$ by $B-4$ has equal cost, and the numeral $4$ is
already used elsewhere. No positive system unknown or equality is
added. The complete primitive list and all 21 residual identities,
including the existing triangular corrections, are verified by
\file{round8\_1980\_certificate.py}. The fixed exponent remains
$L=5^{16}$, so numeral construction still requires seven instructions.
\end{proof}

% Exchange the input and encoded unit positions; preserves the 110 milestone.
\section{Putting the input in the unit position}\label{sec:input-unit}

The supplied input can occupy the constant coefficient of the code.
Exchanging its position with that of the encoded unit changes
$\mathcal C=xB+g$ to $\mathcal C=x+g$, removing one multiplication.
The main proof obligation is to recover the Pell conclusions without
assuming $\mathcal C>B$. The geometric equation provides the weaker
bound $B\le q^2$, which suffices.

\subsection{Remapping the homogeneous code}

Retain the homogeneous rows, guard $u\delta-x^2$, special target
$\delta^2$, and all numerical layout constants of
Section~\ref{sec:unit-digit}. Give $x$ weight zero and $\delta$
weight $v_0=1$. The original coordinates retain weights
$v_1,\ldots,v_{58}$, and $u$ retains weight $v_{59}$.
The mask becomes
\[
 \ell_0(B)=\sum_{i=0}^{59}B^{v_i}
                    +\sum_{j=0}^{1831}B^{t_j}.
\]
It excludes position zero, which now belongs to the input.
There are still 1892 non-input encoded coordinates and 1893 coordinates
including the input. The same bound $1900b$ applies to their sum.

Recompute each complementary coefficient position using the new
monomial weights. The ordinary coefficient rule is still $2a/c$.
The special row has coefficient $1$ at $t_{1831}-2$, since
$\delta^2$ now has weight two. The guard's two monomial weights are
$v_{59}+1$ and zero. All 1891 homogeneous monomial weights remain
distinct: exchanging which coordinate has weight zero does not affect
the number of positive-weight factors or the Sidon argument.
All support-separation bounds consequently remain valid.

With the remapped signed polynomial $\mathcal D$, define
\[
 e_0(B)=z\sum_{h=0}^{K-1}B^h+\mathcal D(B),
 \qquad V=\ell_0(Z)+e_0(Z)Z^L,
\]
and use the same admissibility bound \eqref{eq:unit-H}.
These indices are again fixed independently of the input value.
Make the sole arithmetic replacement
\[
 \mathcal C=x+g
\]
in the 110-operation system. In particular $Y=\ell+eq$,
$Y+\mathcal C^2+\alpha=q^2$, $\sigma=S_3>0$, the mask
$(B-4)\ell$, and all Pell equations remain unchanged.

\subsection{The weaker bounds before the Pell step}

The positive bound gives
\[
 0<\ell<Y<q^2,\qquad e<q,\qquad
 \mathcal C<q,\qquad g<q.
\]
The geometric equation and $\lambda\ge1$ imply
\begin{equation}\label{eq:input-weaker-radix}
 q^2-1=\lambda(B-1)\ge B-1,
 \qquad B\le q^2\le n=q^8.
\end{equation}
Since $B=Hb^2$ and $H>16$, we already have $q>4b\ge8$.
Put $\mathcal N=(B-Z)\lambda$. Then
\[
 b<B-Z\le\mathcal N<q^2-1,
 \qquad B\lambda<2q^2.
\]
No inequality $\lambda>q$ is needed. The imposed positive register
and its elementary upper bound give
\[
 0<S_3<2q^3+2q^2(1+q)<5q^3<q^4.
\]
Thus $S=g+q^2(Y+q^2S_3)$ satisfies $0<S<n$.

For completeness, the possibly negative first raw mask still has
an exact normalized packing. Set
\[
 d_0=\left\lfloor\frac{(b-1)\ell}{q^2}\right\rfloor,
 \qquad \varepsilon_0=(b-1)\ell-d_0q^2.
\]
Then $0\le d_0<b$, $0\le\varepsilon_0<q^2$, and
\[
 T=(q^2-1-\varepsilon_0)
       +q^2(\mathcal N-d_0)+q^4(B-4)\ell.
\]
The first two blocks belong to $[0,q^2)$, since
$\mathcal N>b>d_0$. By \eqref{eq:input-weaker-radix},
the last block satisfies
$0<(B-4)\ell<q^4$. Hence $0\le T<n$, and
\[
 r=S(n^2-n)+(T+1)(n^2-1)\ge n^2-1\ge n.
\]
This establishes all packing bounds before decoding or any conclusion
that $B<q$.

\subsection{The Pell bootstrap with \texorpdfstring{$B\le n$}{B at most n}}

Here is the dependency check for the relaxed radix comparison.
The information now available is
\begin{equation}\label{eq:input-pell-initial}
 3<3L\le B\le n\le r,\qquad q,b\le n,
 \qquad n,r\ge8.
\end{equation}
Moreover $B$ is even because the fixed parameter $H$ is an even
power of two, independently of any assertion about $b$.

Write $N=n$, $R=r$, $U=wN^2$, $Y_P=sN^2$,
$M=RY_P$, $a=M(U+1)$, $X_P=2UM^2$, $P=X_P+1$, and
$J=2R+1$. We have $U,Y_P\ge64$ and $M\ge512$.
The index congruence gives $k\ge R+2$, and the positive interval
$Y_P<c/k<Y_P+1$ gives $c>J$. The signed Pell block therefore
gives $c=\psi_a(J)$ and $d=\chi_a(J)$.
The growth and index-congruence argument of
Section~\ref{sec:shifted-pell} gives $k=\psi_P(R+1)$ and
\[
 \xi\left(1-\frac{R}{a}-\frac{R}{X_P}\right)
       <\frac ck<\xi,
 \qquad \xi=\frac{(U+1)^{2R}}{U^R},
 \qquad \frac{R}{a}+\frac{R}{X_P}
       <\frac{2}{Y_P(U+1)}<\frac12.
\]
These arguments use the numerical lower bounds just stated, not
a comparison between $B$ and $q$. The interval consequently gives
$Y_P>U^{R-1}$ and $a>RU^R$.

The base-two exponent argument has the same strict size bounds:
\[
 2^{3J}=8\cdot8^{2R}\le8N^{2R}\le RU^R<a,
 \qquad (bw)^3\le U^R<a.
\]
Here $b\le N$ gives $bw\le U$.
Equation $E14$ and the exponent lemma yield $bw=2^J$;
thus $b$ is a power of two and $U>4\cdot2^{2R}$.

Next \eqref{eq:input-pell-initial} gives $0<L<J<a$.
The positive gap $c>\kappa$ and equations $E19,E20$ give
$\kappa=\psi_a(L)$ and $\mu=\chi_a(L)$. The second exponent
argument now uses
\begin{equation}\label{eq:input-exponent-bounds}
 B^{3L}\le B^B\le N^N\le U^R<a,
 \qquad q^3\le N^N\le U^R<a.
\end{equation}
These follow from $3L\le B\le N$, $q\le N$, and $R\ge N$.
Also $a>N\ge B$, so the modulus in $E18$ has its required
positive sign. Equation $E18$ therefore gives $q=B^L$.
Only now do we infer $B<q$.

As $B$ was already even, $q$ and $N$ are even. To make the final
rounding dependency explicit, the exact-index estimates give
\[
 0<\xi-c/k<6/(U+1)<1/4,
 \qquad 0<\xi-\lfloor\xi\rfloor<1/4,
 \qquad \lfloor\xi\rfloor\equiv\binom{2R}{R}\pmod U.
\]
The interval makes $|Y_P-\lfloor\xi\rfloor|<2$.
Both integers are even, so they are equal. Since $Y_P=sN^2$
and $N^2\mid U$, we obtain
$N^2\mid\binom{2R}{R}$. The binomial criterion and the packing
lemma now give $\tau_2(S,T)=0$. Thus the complete Pell conclusion
was recovered using \eqref{eq:input-weaker-radix}, without an
unproved early inequality $B<q$.

\subsection{Decoding and positive necessity}

The subsequent first-block borrow argument applies unchanged.
It proves $d_0=0$ and decodes
\[
 \ell=\ell_0(B)+mq,\qquad e=e_0(B)-m,
 \qquad 0\le m<e_0(B).
\]
Since $g<q$, the first mask restricts its digits to the positive
positions of $\ell_0$. Its unit digit is zero. The inequality
$x<b<B$ then ensures that $\mathcal C=x+g$ has exactly the supplied
input as its unit digit, without a carry; its digit at position one
is $\delta$.

The sum of the coefficients of this digit polynomial is positive
and less than $1900b$. All degree and coefficient estimates in
Section~\ref{sec:unit-digit} still hold. The same high-digit window
bounds the digits of $(B-4)m$ by $Z\mathcal C(1)^2-1$; evaluating
its digit polynomial at four forces $m=0$. The canonical relaxed
mask then forces every ordinary homogeneous residual and the guard
to vanish, while its special target gives $\delta^2\in\{0,1\}$.
The guard and $x>0$ give $\delta=1$, so all original residuals vanish.

Conversely, choose genuine original coordinates, set $\delta=1$
and $u=x^2$, and choose a power of two $b$ greater than the input
and every encoded coordinate. The digit $\delta$ at position one
ensures $g\ge B>0$. The canonical code has the same degree and
coefficient bounds as before, giving positive $\alpha$, $\sigma$,
and the packed quotient. All masks hold. In this necessity direction
the intended code even satisfies $\mathcal C>B$, so the preceding
Pell necessity constructions retain their stronger original bounds.

\begin{theorem}\label{thm:best109}
For each r.e.\ set, the remapped effective admissible index gives
a membership-equivalent system in 33 positive unknowns and 21 equations
with a straight-line certificate of \textbf{109 arithmetic instructions}:
60 multiplications and 49 additions. Generating the fixed numerals
from $1$ adds seven instructions, giving 116 operations.
\end{theorem}
\begin{proof}
The preceding arguments establish equivalence with all domains
preserved. Replacing $xB+g$ by $x+g$ removes exactly one multiplication
from the 110-operation schedule. No other arithmetic instruction,
unknown, or equality is added. The complete primitive list and all
21 source residual identities, including the triangular corrections,
are verified by \file{round9\_1980\_certificate.py}.
\end{proof}

% Two independent successors to 109; together they give 107 instructions.
\section{A positive coefficient and a triangular Pell equation}
\label{sec:final-two-savings}

Two independent changes remove another addition and another multiplication.
The first uses positivity of an already computed coefficient to recover
a bound that previously required an explicit addition. The second
rewrites the first Pell equation using its odd root and proves that
its index modulus can be halved. Each separately gives 108 instructions;
together they give 107.

\subsection{Removing the square from the explicit bound}

Keep the fixed input-unit encoding of Section~\ref{sec:input-unit},
with $\mathcal C=x+g$ and $Y=\ell+eq$. Replace
\[
 Y+\mathcal C^2+\alpha=q^2
\]
by the equations
\begin{equation}\label{eq:positive-coefficient-bound}
 Y+\alpha=q^2,\qquad \Omega=Z\lambda-2e,
 \qquad \Omega>0.
\end{equation}
Here $\Omega$ is a new positive system unknown. Its defining value
is the negative of an existing register. Reverse that subtraction and
compute
\[
 P_1=(Z\lambda-2e)\mathcal C^2,
 \qquad P_2=B\lambda(1+q),\qquad S_3=P_2-P_1.
\]
The value of $S_3$ is unchanged. Reversing these signs has equal
arithmetic cost, and $\Omega$ is connected to the computed coefficient
by a free equality test. Retain the separate positive unknown
$\sigma=S_3$.

\begin{proposition}\label{prop:positive-coefficient108}
The replacements \eqref{eq:positive-coefficient-bound} give a
membership-equivalent system with 108 arithmetic instructions:
60 multiplications and 48 additions. It has 34 positive unknowns
and 22 equations.
\end{proposition}
\begin{proof}
Before decoding, the new bound gives $\ell<Y<q^2$ and $e<q$.
The geometric equation still gives $B\le q^2$ and $q>4b\ge8$.
Consequently
\[
 P_2=B\lambda(1+q)<2q^2(1+q)<3q^3.
\]
Since $\Omega$ and $\sigma$ are positive integers,
\[
 0<S_3=P_2-\Omega\mathcal C^2<P_2,
 \qquad \mathcal C^2<P_2/\Omega\le P_2<3q^3<q^4.
\]
Thus $g<\mathcal C<q^2$. This weaker bound is enough for the first
packing block, whose width is $q^2$. Together with $Y<q^2$ and
$S_3<q^4$, it gives $0<S<n=q^8$.

The mask bounds use only $\ell<q^2$, $B\le q^2$, and
$b<(B-Z)\lambda<q^2-1$. They are unchanged. Explicitly, if
$d_0=\lfloor(b-1)\ell/q^2\rfloor$ and
$\varepsilon_0=(b-1)\ell-d_0q^2$, then
\[
 T=(q^2-1-\varepsilon_0)
       +q^2\bigl((B-Z)\lambda-d_0\bigr)+q^4(B-4)\ell
\]
has normalized blocks in $[0,q^2)$, $[0,q^2)$, and $[0,q^4)$.
Hence $0\le T<n$ and $r\ge n^2-1\ge n$.
The relaxed Pell bootstrap proved in Section~\ref{sec:input-unit}
now applies: it uses $3L\le B\le n\le r$, $q,b\le n$, and the
unchanged Pell equations, without a further bound on $g$ or
$\mathcal C$. It gives $q=B^L$, the required powers of two,
and $\tau_2(S,T)=0$.

We recover the stronger bound before using it in digit decoding.
Now $q=B^L>B$, so $\lambda>q$. Since $e<q$ and $Z\ge4$,
\[
 2e<2q<2\lambda\le Z\lambda/2,
 \qquad \Omega>Z\lambda/2.
\]
It follows that
\[
 \mathcal C^2<\frac{P_2}{\Omega}
       <\frac{2B(1+q)}Z
       \le\frac{B(1+q)}2
       <\frac{q(q+1)}2<q^2.
\]
Thus $g<\mathcal C<q$ holds before the first-mask decoding and
high-mask quotient argument require it. Those arguments now apply
unchanged, proving canonical code recovery and membership.

Conversely, take the positive witnesses constructed for the
109-operation system. They have $e<q$, $\lambda>q$, and $S_3>0$.
Set $\Omega=Z\lambda-2e>0$ and replace the slack by
$\alpha=q^2-Y>0$. All other witnesses remain valid.

Only the addition of $\mathcal C^2$ to $Y$ disappears.
The coefficient subtraction, its product with $\mathcal C^2$, and
the subtraction defining $S_3$ have the same costs as before.
The positive unknown $\Omega$ introduces one free equality and no
arithmetic instruction beyond these existing registers.
The complete check is \file{round10\_1980\_certificate.py}.
\end{proof}

\subsection{Factoring the odd Pell root}

For this independent change to the 109-operation system, write
\[
\begin{gathered}
 N=n,\quad R=r,\quad U=wN^2,\quad Y_P=sN^2,
 \quad M=RY_P,\\
 a=M(U+1),\quad Q=UM^2,
 \quad P=2Q+1,\quad J=2R+1.
\end{gathered}
\]
The quantity $Q$ is already computed. The old first Pell norm
and index equations were
\[
 \bigl((2Q+1)^2-1\bigr)k^2+1=\tau_{\rm old}^2,
 \qquad k=R+1+2Qh_{\rm old}.
\]
Replace them by
\begin{equation}\label{eq:triangular-pell}
 \tau(\tau+1)=Q(Q+1)k^2,
 \qquad k=R+1+hQ,
 \qquad \tau,h>0.
\end{equation}
The mathematical Pell parameter $P=2Q+1$ need not be computed
by the new certificate.

\begin{proposition}\label{prop:triangular-pell108}
Equations \eqref{eq:triangular-pell} give a positive-domain
equivalent system with 108 arithmetic instructions:
59 multiplications and 49 additions. This independent variant
retains 33 positive unknowns and 21 equations.
\end{proposition}
\begin{proof}
For necessity, $P$ is odd, and the recurrence for $\chi_P(t)$
shows that every such coordinate is odd. The old positive norm
has $\tau_{\rm old}>1$. Thus
\[
 \tau=(\tau_{\rm old}-1)/2\ge1,
 \qquad h=2h_{\rm old}>0
\]
are integers satisfying \eqref{eq:triangular-pell}.
The factored norm is simply the old norm identity divided by four.
The certificate verifies the new polynomial equation directly;
this witness map does not introduce an uncharged division instruction.

For sufficiency, the initial coding bounds give $R\ge N\ge8$,
so $U,Y_P\ge64$, $M\ge512$, and $Q>R+1$.
Put $\widehat\tau=2\tau+1$. The new norm implies
\[
 \widehat\tau^2-(P^2-1)k^2=1,
 \qquad \widehat\tau>1.
\]
The new congruence gives $k\ge R+2$. The unchanged positive
interval gives $c>Y_Pk\ge64(R+2)>J$. The signed Pell block
therefore gives $c=\psi_a(J)$ and $d=\chi_a(J)$.

There is a positive index $t$ with
$k=\psi_P(t)$ and $\widehat\tau=\chi_P(t)$.
The elementary congruence $\psi_P(t)\equiv t\pmod{P-1}$
also holds modulo $Q$. Since $0<R+1<Q$, the new index equation
implies
\[
 t=R+1+vQ,\qquad v\ge0.
\]
If $v\ge1$, the elementary Pell growth estimates give
\[
 \frac ck\le
 \frac{(2a)^{2R}}{(2P-1)^{R+Q}}
 \le\left(\frac{2a}{2P-1}\right)^{2R}<\frac12.
\]
Here $Q\ge R$, and
$2a/(2P-1)=2M(U+1)/(4UM^2+1)<1/2$ follows from
$U\ge64$ and $M\ge512$. This contradicts $c/k>Y_P\ge64$.
Therefore $v=0$ and $k=\psi_P(R+1)$.

Applying the same Pell congruence now modulo $2Q$ gives
$k\equiv R+1\pmod{2Q}$. Since $k-R-1=hQ$, the new positive
integer $h$ is even. Consequently
\[
 h_{\rm old}=h/2>0,\qquad
 \tau_{\rm old}=2\tau+1>0
\]
restore both old equations, preserving every other witness.
This argument precedes the exponent relations and binomial rounding;
it does not assume either of them to prove the required evenness.

After $Q$ is available, the old norm block used seven operations:
form $2Q$, add two, multiply these two values, square $k$,
multiply, add one, and square $\tau_{\rm old}$.
The new block uses six: form $Q+1$ and $Q(Q+1)$, square $k$,
multiply, form $\tau+1$, and multiply $\tau(\tau+1)$.
The index product $hQ$ has the same cost as the former $2Qh_{\rm old}$,
so $2Q$ is no longer needed anywhere in the schedule.
The exact saving is one multiplication, as checked by
\file{round11\_1980\_pell\_certificate.py}.
\end{proof}

\subsection{Combining the two savings}

\begin{theorem}\label{thm:best107}
For each r.e.\ set, the fixed admissible input-unit encoding gives
a membership-equivalent system in 34 positive unknowns and 22 equations
whose straight-line certificate has \textbf{107 arithmetic instructions}:
59 multiplications and 48 additions. Equality tests and supplied
parameters are free. Generating the fixed numerals from $1$ adds
seven instructions, giving 114 operations.
\end{theorem}
\begin{proof}
Apply both replacements to the 109-operation system. Their
necessity maps affect disjoint witnesses: the coefficient-bound
change replaces $\alpha$ and supplies $\Omega$, while the Pell
change replaces $\tau$ and $h$. Thus both maps preserve the
other change's equations and positivity.

For soundness, first use $\Omega,\sigma>0$ to establish the
preliminary bounds $0<S<n$, $0\le T<n$, and $R\ge N\ge8$,
as in Proposition~\ref{prop:positive-coefficient108}.
The index proof in Proposition~\ref{prop:triangular-pell108}
then restores the old positive Pell witnesses without any exponent
or decoding assumption. The positive-coefficient proof can now
continue with those witnesses and recover membership. This order
establishes the equivalence without a circular appeal between the
two changes.

The combined schedule removes one addition and one multiplication
from 109 instructions. The coefficient change supplies the sole
extra positive unknown and free equality; the Pell change alters
neither count. The complete primitive list and all residual
identities are checked by \file{round12\_1980\_\allowbreak certificate.py}.
The numeral set and exponent $L=5^{16}$ are unchanged.
\end{proof}

\section{A smaller fixed support}\label{sec:modular-support}

We use a smaller support for the subsequent constructions. Replace the
support positions and exponent of Section~\ref{sec:input-unit}, keeping
its arithmetic structure. The resulting certificate still has 107
instructions. Fixed numerals are free in the complexity measure used
for the reductions below.

\subsection{Pair recovery over a finite field}

For $0\le i\le59$, let $r_i$ be the least nonnegative residue of $i^2$
modulo~61, and put
\[
 a_i=122i+r_i,\qquad v_i=1+3a_i.
\]
The sums $a_i+a_j$ distinguish unordered pairs, including repeated
indices. Indeed, $0\le r_i+r_j\le120<122$, so division by~122 recovers
the integer sum $s=i+j$ and the residue sum $r_i+r_j$. Over the field
with 61 elements these determine
\[
 i+j=s,\qquad ij=(s^2-r_i-r_j)/2.
\]
Thus the monic polynomial with roots $i,j$ is determined. Since the
indices $0,\ldots,59$ are distinct modulo~61, the integer pair is
determined as well. The division by~2 occurs in this proof of the
fixed support construction; it is not a certificate instruction.

Every $v_i$ is 1 modulo~3. Hence zero, a single positive weight, and
a pair sum are distinguished modulo~3. Within each class the weights
are distinct by the preceding argument. There are consequently
$1+60+1830=1891$ distinct homogeneous quadratic monomial weights in
the 61 coordinates. As before, the input has weight zero, $\delta$
has weight $v_0=1$, the original witnesses have weights $v_1,\ldots,v_{58}$,
and $u$ has weight $v_{59}$.

With $m=1832$ target rows, use the same formulas for their spacing:
\[
\begin{gathered}
 v_*=21607,\qquad D_*=4v_*+1=86429,\\
 T_*=(m+1)D_*+2v_*=158467571,\qquad t_j=T_*+jD_*,\\
 t_*=316719070,\qquad K=t_*+1=316719071,\qquad L=2^{32}.
\end{gathered}
\]
The exact margins are
\[
\begin{gathered}
 D_*>2v_*,\qquad T_*-2v_*>v_*,\qquad 2T_*-2v_*>t_*,\\
 3K+2=950157215<2^{30}<L.
\end{gathered}
\]
Different row intervals $[t_j-2v_*,t_j]$ are disjoint. A product of
a coefficient term and an ordinary quadratic monomial lies within
$2v_*$ of its row target, so no other row contributes there. Coefficient
terms lie above every low variable position, and products involving a
dummy row coordinate lie above every target. These are all the support
facts used in the preceding decoding proof.

Retain twice each ordinary homogeneous residual and the special
coefficient $\delta^2$. In particular, its coefficient 1 is placed at
$t_{m-1}-2$. Recompute $D,\ell_0,e_0,V$ with the new positions, choose
the power of two $z\ge2$ above the absolute coefficients of $D$, put
$Z=2z$, and choose a power of two
\[
 H>\max\{2Z^{2L+1},\ 4^{t_*+3}Z\,1900^2,\ 3L,\ 16\}.
\]
The number of coordinates remains 1893, so their coefficient sum remains
less than $1900b$. All bounds in the input-unit and positive-coefficient
proofs are therefore preserved. In particular, $e,\mathcal C<B^K$ and
$L>3K+2$ still give $2e\mathcal C^2<q$ after exponent recovery, and
evaluation at~4 still rules out the high-mask quotient ambiguity.
The fixed indices $Z,V,H$ are computed from the represented set,
independently of the queried input.

\subsection{The even exponent and the preserved certificate}

No preceding Pell argument requires $L$ to be odd. The signed congruence
uses the odd index $J=2r+1$. The other index argument uses only
$0<L<J<a$ and the congruence $\psi_a(t)\equiv t\pmod{a-1}$, valid for
every $t$. These hypotheses still hold because $3L<B\le n\le r$.
Necessity gives the positive integer
$\Delta=(\psi_a(L)-L)/(a-1)$ since $L>1$, and the positive gap follows
from $J>L$. The exponent comparison
$B^{3L}\le B^B\le n^n\le U^r<a$ is unchanged. Thus $q=B^L$ is
recovered, and every decoding and necessity argument still applies.

\begin{theorem}\label{thm:modular107}
The modular support gives a membership-equivalent system in 34 positive
unknowns and 22 equations with \textbf{107 arithmetic instructions}:
59 multiplications and 48 additions.
\end{theorem}

\begin{proof}
The support and exponent arguments establish the same encoding theorem.
The explicit schedule moves the existing register $a_-=a-1$ before E14 and
replace its computations of $4a$ and $4a-5$ by
\[
 a_4=4a_-,\qquad a_{4-5}=a_4-1.
\]
The latter is identically $4a-5$, so all 107 arithmetic instructions
are preserved. The complete schedule is serialized and checked by
\file{round13\_1980\_certificate.py}. It checks every primitive and
source residual identity and all 1891 monomial weights and support
margins. Its 59 multiplications and 48 additions have exactly the
same costs as in Appendix~\ref{app:107}.
\end{proof}

% A shifted main Pell parameter shares the existing B-4 mask register.
\section{A base-four congruence and the main Pell parameter}
\label{sec:main-base-four}

The third digit mask already computes $B-4$. This permits a useful joint
change: replace the base-two exponent test by a base-four test, and store
the main Pell parameter minus four. The two expressions $A-4$ and $A-B$
then share available values, eliminating one subtraction. The altered
Pell approximation lies strictly above its binomial expression; this
also makes the final rounding argument independent of parity.

Use the 107-operation system of Theorem~\ref{thm:best107}, with the
fixed modular-Sidon encoding and $L=2^{32}$ from
Section~\ref{sec:modular-support}. In this section $Y_P$ denotes
the Pell quantity $s n^2$, to distinguish it from the packed code.
Write
\[
\begin{gathered}
 N=n,\quad R=r,\quad U=wN^2,\quad Y_P=sN^2,\quad M=RY_P,\\
 a_0=M(U+1),\quad A=a_0+4,\quad Q=UM^2,\quad P=2Q+1,\\
 J=2R+1,\qquad W=bw,\qquad \xi=\frac{(U+1)^{2R}}{U^R}.
\end{gathered}
\]
The supplied positive unknown named $a$ in the certificate now means
$a_0$. Equation E12 still defines $a_0=M(U+1)$. Every main Pell
equation uses the mathematical parameter $A=a_0+4$, including
$G=1+(A+1)(f^2-1)$ in the signed auxiliary block and the modulus
$A-1$ in E20. The first Pell norm and index equations retain the forms
\begin{equation}\label{eq:base-four-first-pell}
 \tau(\tau+1)=Q(Q+1)k^2,\qquad k=R+1+hQ.
\end{equation}
The positive interval remains
\begin{equation}\label{eq:base-four-interval}
 c=kY_P+\eta,\qquad k=\eta+\zeta,
 \qquad\eta,\zeta>0,
 \qquad Y_P<\frac ck<Y_P+1.
\end{equation}
Replace E14 by
\begin{equation}\label{eq:base-four-congruence}
 d=W+c(A-4)+\gamma(8A-17).
\end{equation}
This is the source's $\chi$ exponent congruence for base four.
The second exponent congruence retains the form
\begin{equation}\label{eq:base-four-second-congruence}
 \mu=q+\kappa(A-B)+\rho(2AB-B^2-1).
\end{equation}

\subsection{Exact indices and the shifted approximation}

The pre-Pell part of Proposition~\ref{prop:positive-coefficient108}
uses only the encoding equations and gives
\[
 N=q^8,\qquad q>4b\ge8,\qquad
 3L\le B\le N\le R,\qquad b,q\le N.
\]
In particular $N,R\ge64$, $U,Y_P\ge4096$, $M\ge262144$,
$Q>R+1$, and $A>J>1$. These estimates precede both exponent
relations; the stronger $N\ge64$ will be needed for base four.

Since $h>0$, equation~\eqref{eq:base-four-first-pell} gives
$k\ge R+2$. The interval gives $c>Y_P(R+2)>J$.
The generic signed Pell argument of Proposition~\ref{lem:signed-pell}
therefore applies with parameter $A$ and gives
\begin{equation}\label{eq:base-four-main-index}
 c=\psi_A(J),\qquad d=\chi_A(J).
\end{equation}
The triangular norm gives the ordinary Pell equation for
$\widehat\tau=2\tau+1$ and $P=2Q+1$, so
$k=\psi_P(t)$ for some positive index $t$. The congruence
$\psi_P(t)\equiv t\pmod{2Q}$, reduced modulo $Q$, implies
\[
 t=R+1+vQ,\qquad v\ge0.
\]
If $v\ge1$, the elementary Pell growth bounds give
\[
 \frac ck\le
 \frac{(2A)^{2R}}{(2P-1)^{R+Q}}
 \le\left(\frac{2A}{2P-1}\right)^{2R}<\frac12.
\]
Here $Q\ge R$, and $2A/(2P-1)<1/2$ follows, for example, from
\[
 a_0<2UM,\qquad
 4A<8UM+16<4UM^2+1=2P-1.
\]
This contradicts~\eqref{eq:base-four-interval}. Consequently
\begin{equation}\label{eq:base-four-first-index}
 k=\psi_P(R+1),\qquad \widehat\tau=\chi_P(R+1).
\end{equation}
No exponent or rounding conclusion has been used to recover either
Pell index.

At these indices, the lower growth estimate gives
\begin{align}
 \frac ck
 &\ge\frac{(2A-1)^{2R}}{(2P)^R}\notag\\
 &=\xi\left(1+\frac7{2a_0}\right)^{2R}
          \left(1+\frac1{2Q}\right)^{-R}
 >\xi.\label{eq:base-four-lower}
\end{align}
Indeed $14Q>a_0$, so
$(1+7/(2a_0))^2>1+7/a_0>1+1/(2Q)$.
The upper growth estimate gives
\begin{align}
 \frac ck
 &\le\frac{(2A)^{2R}}{(2P-1)^R}
 <\xi\left(1+\frac4{a_0}\right)^{2R}\notag\\
 &<\xi\left(1+\frac{16R}{a_0}\right)
 =\xi\left(1+\frac{16}{Y_P(U+1)}\right).
 \label{eq:base-four-upper}
\end{align}
For the last step, binomial coefficients give
$(1+t)^m\le\sum_{j\ge0}(mt)^j=(1-mt)^{-1}<1+2mt$
when $0<mt<1/2$. Here $mt=8R/a_0=8/[Y_P(U+1)]<1/2$.

The lower bound and the interval imply $\xi<Y_P+1$.
Since $\xi>U^R$, integrality gives $Y_P\ge U^R$ and hence
\begin{equation}\label{eq:base-four-size}
 A>a_0=RY_P(U+1)>RU^R.
\end{equation}
This size bound is established before invoking an exponent lemma.
Also $\xi<Y_P+1<2Y_P$, so the upper bound yields
\begin{equation}\label{eq:base-four-error}
 0<\frac ck-\xi<\frac{32}{U+1}<\frac12.
\end{equation}

\subsection{The exponent relations and rounding}

Since $W=bw\le U$ and $N,R\ge64$, we have
\[
 4^{3J}=64\cdot64^{2R}
       \le64N^{2R}\le RU^R<A,
 \qquad W^3\le U^R<A.
\]
The $\chi$ version of source Lemma~2.22, explicitly stated after
that lemma, applies to~\eqref{eq:base-four-congruence} and
\eqref{eq:base-four-main-index}. It proves
\begin{equation}\label{eq:base-four-power}
 bw=4^J.
\end{equation}
Thus $b$ is a power of two, and
\[
 U=N^2w\ge Nbw=N4^J>4\cdot2^{2R}.
\]
E13, E19 and E20 give $0<\kappa<c$, its Pell norm, and
$\kappa\equiv L\pmod{A-1}$. Since $0<L<J<A$, source
Lemma~2.23 gives $\kappa=\psi_A(L)$ and $\mu=\chi_A(L)$.
Moreover
\[
 B^{3L}\le B^B\le N^N\le U^R<A,
 \qquad q^3\le N^N\le U^R<A.
\]
The same exponent lemma, applied now to
\eqref{eq:base-four-second-congruence}, proves $q=B^L$.

Let $F=\lfloor\xi\rfloor$. Source Lemma~2.24, or its direct
binomial expansion, gives
\begin{equation}\label{eq:base-four-fraction}
 0<\xi-F<\frac{2^{2R}}U<\frac14,
 \qquad F\equiv\binom{2R}{R}\pmod U.
\end{equation}
Combining the interval with~\eqref{eq:base-four-lower} and
\eqref{eq:base-four-error} gives
\[
 F<\xi<\frac ck<Y_P+1,
 \qquad Y_P<\frac ck<\xi+\frac12<F+\frac34.
\]
Both $F$ and $Y_P$ are integers. The first chain gives $F\le Y_P$;
the second gives $Y_P\le F$. Thus
\begin{equation}\label{eq:base-four-rounding}
 Y_P=F,\qquad N^2\mid\binom{2R}{R}.
\end{equation}
This argument requires no parity assertion. The power-of-two facts
for $b$ and $q$, the relation $q=B^L$, and the divisibility in
\eqref{eq:base-four-rounding} are precisely the conclusions needed
by the unchanged no-carry and modular-Sidon decoding proof.
They therefore imply membership in the represented r.e.\ set.

\subsection{Constructing positive witnesses}

For the converse, start from a solution of the represented system
and take its canonical modular-Sidon coding witnesses. Choose the
usual sufficiently large power of two $b$, put $B=Hb^2$,
$q=B^L$, $N=q^8$, and construct the unchanged packed $S,T,R$.
Their no-carry property gives $N^2\mid\binom{2R}{R}$.
These choices precede the Pell witnesses and supply $b\le R$
and $R\ge N\ge64$.

Put $J=2R+1$ and choose
\[
 w=4^J/b,\quad U=N^2w,\quad
 Y_P=\left\lfloor\frac{(U+1)^{2R}}{U^R}\right\rfloor,
 \quad s=Y_P/N^2.
\]
Writing $b=2^t$, the inequality $t\le b\le R<2J$ shows that
$w$ is a positive integer. The binomial expansion and central-binomial
divisibility make $s$ a positive integer. Define
$M=RY_P$, $a_0=M(U+1)$, $A=a_0+4$, $Q=UM^2$, $P=2Q+1$, and
\[
 c=\psi_A(J),\qquad d=\chi_A(J),\qquad k=\psi_P(R+1).
\]
The ratio estimates above hold at these chosen indices.
Here $\xi<Y_P+1$ is immediate from the definition of $Y_P$, so
\[
 Y_P<\xi<\frac ck<\xi+\frac12<Y_P+\frac34.
\]
Consequently
\[
 \eta=c-kY_P>0,\quad \zeta=k-\eta>0,\quad
 \tau=\frac{\chi_P(R+1)-1}{2},\quad
 h=\frac{k-R-1}{Q}
\]
are positive integers. Positivity and integrality of $\tau$ follow
from the oddness of the $\chi$ sequence for odd $P$; those of $h$
follow from the Pell congruence modulo $2Q$ and
$\psi_P(R+1)>R+1$. Thus the triangular norm, index equation,
positive interval, E12, and E15 all hold.

Choose positive $f,i$ with
$f^2-(A^2-1)(ic^2)^2=1$ using the same generic Pell divisibility
construction as in Proposition~\ref{lem:signed-pell}. Put
\[
 G=1+(A+1)(f^2-1),\quad I=\chi_G(J),\quad H_J=\psi_G(J),
 \qquad o=(I+d)/f,\quad j=(H_J-J)/c.
\]
Because $G\equiv-A\pmod f$, $G\equiv1\pmod c$, and $J$ is odd,
$o$ and $j$ are integers. They are positive since $I+d>0$ and
$\psi_G(J)>J$. These values satisfy the signed E17 with $of-d=I$.
Next choose
\[
 \kappa=\psi_A(L),\quad \mu=\chi_A(L),\quad
 \Delta=\frac{\kappa-L}{A-1},\quad \phi=c-\kappa.
\]
Since $J>L\ge2$, these are positive integers with the required
Pell norm, congruence and gap.

The exponent lemma supplies integer quotients $\gamma,\rho$ in
\eqref{eq:base-four-congruence} and
\eqref{eq:base-four-second-congruence}. To check their positivity,
use $\chi_A(t)=A\psi_A(t)-\psi_A(t-1)$ to obtain
\[
 d-(A-4)c>3c>W,
 \qquad \mu-(A-B)\kappa>(B-1)\kappa>q.
\]
For the last inequalities, $A>W^3,q^3$ and $J,L\ge2$ give
$c,\kappa\ge A$. The moduli are positive because $A>B\ge2$
and $A>4$. Hence $\gamma,\rho>0$.
All coding witnesses, including $\Omega$ and $\sigma$, remain
the canonical ones. This constructs every required positive witness.
In particular the construction re-chooses $w,s,a_0$ and their
dependent Pell coordinates; it does not preserve the predecessor's
base-two value of $w$.

\subsection{The complete arithmetic count}

\begin{theorem}\label{thm:best106}
For each r.e.\ set, the fixed modular-Sidon encoding gives a
membership-equivalent system in 34 positive unknowns and 22 equations
with a straight-line certificate of \textbf{106 arithmetic instructions}:
59 multiplications and 47 additions. Equality tests and supplied
parameters are free, as are fixed numerals.
\end{theorem}
\begin{proof}
The preceding arguments prove both directions over positive integers.
The arithmetic expressions are evaluated as
\[
\begin{gathered}
 A-4=a_0,\quad A-1=a_0+3,\quad A+1=(A-1)+2,\\
 A^2-1=(A-1)(A+1),\quad
 8A-17=8a_0+15,\quad A-B=a_0-(B-4).
\end{gathered}
\]
The previous subtraction computing the base-two coefficient $a-2$
disappears. The other expressions cost the same numbers of operations
as their predecessors; in particular $B-4$ is already computed by
the third mask, and the defining equation for $a_0$ retains its
existing factorization. No register for $A$ itself is required.

The checker \file{round14\_1980\_base\_four\_certificate.py}
verifies all source residuals, both triangular corrections, every
shifted-register identity, and the full primitive instruction list.
It counts 59 multiplications and 47 additions. The positive input
and equality counts are unchanged. Its generated schedule records
the exponent $L=2^{32}$.
\end{proof}

% One addition saved by forcing the first Pell index to be a multiple.
\section{Forcing a multiple of the Pell index}
\label{sec:index-multiple}

The first Pell index can be isolated with a single product in place of
an affine congruence. Arrange that its parameter is congruent to one
modulo the intended index. Then require its denominator to be a
positive multiple of that index. Pell growth excludes every multiple
after the first.

Retain the fixed encoding and the shifted main parameter from
Section~\ref{sec:main-base-four}, but change the intermediate scale to
\[
\begin{gathered}
 N=n,\quad R=r,\quad U=wN^2,\quad Y_P=sN^2,\quad D=R+1,\\
 M=DY_P,\quad a_0=M(U+1),\quad A=a_0+4,
 \quad Q=UM^2,\quad P=2Q+1,\quad J=2R+1.
\end{gathered}
\]
The supplied unknown $a$ still denotes $a_0$. Replace the old equation
$k=R+1+hQ$ by
\begin{equation}\label{eq:index-multiple}
 k=h(R+1)=hD,\qquad h>0.
\end{equation}
Retain the triangular norm and positive interval
\begin{equation}\label{eq:index-multiple-norm}
 \tau(\tau+1)=Q(Q+1)k^2,\qquad
 c=kY_P+\eta,\qquad k=\eta+\zeta,
 \qquad \tau,\eta,\zeta>0.
\end{equation}
The scale $M$ is an intermediate value, so the changed scale affects
the expanded source polynomials E9 and E12; equation
\eqref{eq:index-multiple} changes E11. All other source polynomials
in the supplied variables are unchanged.

\subsection{Isolating the first positive multiple}

The unchanged pre-Pell coding proof supplies
\[
 N=q^8,\quad q>4b\ge8,\quad
 3L\le B\le N\le R,\quad b,q\le N.
\]
Thus $N,R\ge64$, $U,Y_P\ge4096$, and $M\ge262144$.
Equation~\eqref{eq:index-multiple} gives $k\ge R+1$, and the interval
in~\eqref{eq:index-multiple-norm} gives
\[
 c>Y_Pk\ge Y_P(R+1)>2R+1=J.
\]
The generic signed Pell argument therefore gives
$c=\psi_A(J)$ and $d=\chi_A(J)$.

Put $\widehat\tau=2\tau+1$. The norm in
\eqref{eq:index-multiple-norm} is precisely
\[
 \widehat\tau^2-(P^2-1)k^2=1,\qquad \widehat\tau>1.
\]
Hence $k=\psi_P(t)$ for a positive integer $t$.
The new scale makes
\[
 P-1=2Q=2U(R+1)^2Y_P^2
\]
divisible by $D$. The elementary Pell congruence
$\psi_P(t)\equiv t\pmod{P-1}$ therefore gives
$k\equiv t\pmod D$. Equation~\eqref{eq:index-multiple} yields
\begin{equation}\label{eq:index-positive-multiple}
 t=mD=m(R+1),\qquad m\ge1.
\end{equation}

If $m\ge2$, then $t-1\ge2R+1$, and Pell growth gives
\[
 \frac ck\le\frac{(2A)^{2R}}{(2P-1)^{2R+1}}
 <\left(\frac{2A}{2P-1}\right)^{2R}<\frac12.
\]
The last inequality follows from $2A/(2P-1)<1/2$; with the shifted
parameter this is equivalent to
\[
 4UM(M-1)-4M-15=4M\bigl(U(M-1)-1\bigr)-15>0.
\]
The initial bounds already imply it. This contradicts $c/k>Y_P$.
Consequently
\begin{equation}\label{eq:index-multiple-exact}
 k=\psi_P(R+1),\qquad\widehat\tau=\chi_P(R+1).
\end{equation}
The conclusion uses no exponent relation or binomial rounding.
No parity condition on the positive quotient $h$ is required.

\subsection{Preserving the approximation and the encoding}

At the recovered indices, the principal ratio still cancels the scale:
\[
 \frac{(2M(U+1))^{2R}}{(4UM^2)^R}
 =\xi:=\frac{(U+1)^{2R}}{U^R}.
\]
The lower estimate of Section~\ref{sec:main-base-four} uses only
$14Q>a_0$, which remains true. Its upper estimate uses
\[
 \frac{8R}{a_0}
 =\frac{8R}{(R+1)Y_P(U+1)}
 <\frac8{Y_P(U+1)}<\frac12.
\]
Thus the same growth and binomial estimates give
\begin{equation}\label{eq:index-multiple-ratio}
 \xi<\frac ck
 <\xi\left(1+\frac{16R}{a_0}\right)
 <\xi\left(1+\frac{16}{Y_P(U+1)}\right).
\end{equation}
In particular $\xi<Y_P+1$, so $Y_P\ge U^R$ and
\[
 A>a_0=(R+1)Y_P(U+1)>RU^R,
 \qquad 0<\frac ck-\xi<\frac{32}{U+1}<\frac12.
\]
These are exactly the size and error bounds needed after index
recovery in Section~\ref{sec:main-base-four}; they have been
established here for the changed scale.

Explicitly, $N,R\ge64$ and $bw\le U$ give
\[
 4^{3J}\le64N^{2R}\le RU^R<A,
 \qquad (bw)^3\le U^R<A.
\]
The unchanged base-four congruence yields $bw=4^J$, so $b$ is a
power of two and $U\ge N4^J>4\cdot2^{2R}$.
The unchanged second Pell norm, gap and index congruence give
$\kappa=\psi_A(L)$ and $\mu=\chi_A(L)$.
The inequalities
\[
 B^{3L}\le B^B\le N^N\le U^R<A,
 \qquad q^3\le N^N\le U^R<A
\]
then give $q=B^L$ through the second exponent congruence.
For $F=\lfloor\xi\rfloor$, the binomial expansion gives
\[
 0<\xi-F<\frac14,\qquad
 F\equiv\binom{2R}{R}\pmod U.
\]
Now $F<\xi<c/k<Y_P+1$ and
$Y_P<c/k<\xi+1/2<F+3/4$ imply $Y_P=F$ by integrality.
Thus $N^2\mid\binom{2R}{R}$. The unchanged no-carry and
modular-Sidon decoding proves membership in the represented system.

\subsection{Positive witnesses and the saved addition}

Conversely, start with the canonical coding witnesses for a solution
of the represented system. They determine $b,B,q,N,R$ before any Pell
coordinates are chosen. As in Section~\ref{sec:main-base-four}, put
\[
 J=2R+1,\qquad w=4^J/b,\qquad U=N^2w,
 \qquad Y_P=\lfloor\xi\rfloor,\qquad s=Y_P/N^2.
\]
The positive integrality of $w$ and $s$ depends only on these coding
values and the binomial expansion, so it is unaffected by the scale
change. Now form $M=(R+1)Y_P$, $a_0=M(U+1)$, $A=a_0+4$,
$Q=UM^2$, $P=2Q+1$, and choose
\[
 c=\psi_A(J),\qquad d=\chi_A(J),\qquad k=\psi_P(R+1).
\]
Estimate~\eqref{eq:index-multiple-ratio} at these chosen indices
gives $Y_P<c/k<Y_P+3/4$, so
$\eta=c-kY_P$ and $\zeta=k-\eta$ are positive integers.
Choose
\[
 \tau=\frac{\chi_P(R+1)-1}{2},\qquad h=\frac{k}{R+1}.
\]
The oddness of $P$ makes $\tau$ positive integral. Since
$R+1\mid P-1$ and $\psi_P(R+1)\equiv R+1\pmod{P-1}$,
$h$ is a positive integer and satisfies the new product equation.
It is a supplied positive unknown; no division operation is omitted
from the certificate.

Choose the remaining dependent Pell witnesses using the generic
constructions at this new $A$ from Section~\ref{sec:main-base-four}.
In particular, after choosing positive $f,i$ for E16, take
\[
 G=1+(A+1)(f^2-1),\qquad
 o=\frac{\chi_G(J)+d}{f},\qquad
 j=\frac{\psi_G(J)-J}{c},
\]
and take $\kappa=\psi_A(L)$, $\mu=\chi_A(L)$,
$\Delta=(\kappa-L)/(A-1)$ and $\phi=c-\kappa$.
The same generic congruences and strict Pell growth make them positive
integers. The exponent quotients are positive because
\[
 d-(A-4)c>3c>4^J,\qquad
 \mu-(A-B)\kappa>(B-1)\kappa>B^L.
\]
Their moduli are positive, and the required inequalities
$A>(4^J)^3,(B^L)^3$ were preserved above. Thus every witness is
constructed with its required positive domain. All coding witnesses
remain valid, while $M$ and its dependent Pell coordinates are
recomputed; the old and new values of $h$ are not identified.

\begin{theorem}\label{thm:best105}
For each r.e.\ set, the fixed modular-Sidon encoding gives a
membership-equivalent system in 34 positive unknowns and 22 equations
with a straight-line certificate of \textbf{105 arithmetic instructions}:
59 multiplications and 46 additions. Fixed numerals, supplied parameters
and equality tests are free.
\end{theorem}
\begin{proof}
The preceding arguments prove soundness and construct all positive
witnesses for necessity. For the count, the existing instruction
$r_1=r+1$ is moved before the computation of $M$. Replacing
$M=rY_P$ by $M=r_1Y_P$ retains one multiplication. Replacing $hQ$
by $hr_1$ also retains one multiplication, but the old addition
$r_1+hQ$ disappears: the free equality now compares $k$ directly
with $hr_1$.

The other instructions retain their counts. Hence the 106-operation
certificate loses one addition and has 59 multiplications and 46
additions. The positive unknown and equation counts stay unchanged.
The checker \file{round16\_1980\_index\_multiple\_certificate.py}
verifies the three changed source residuals, every unchanged residual,
both triangular corrections, all modified register identities, and
the full serialized primitive certificate. Its fixed exponent is
$L=2^{32}$.
\end{proof}

% The interval scale itself suffices as the main Pell scale.
\section{Reusing the interval scale}
\label{sec:unit-scale}

There is another way to remove an operation from the base-four
construction. The scale $M=RY_P$ can be replaced by $M=Y_P$.
This saves its multiplication by $R$. The larger approximation error
is still small enough, provided it is bounded only after the
base-four exponent relation has been proved. A previously unused
upper bound on the packed index supplies the earlier estimates.

Start from Section~\ref{sec:main-base-four}, with the same fixed
modular-Sidon encoding and $L=2^{32}$. Set
\[
\begin{gathered}
 N=n,\quad R=r,\quad U=wN^2,\quad Y_P=sN^2,\quad M=Y_P,\\
 a_0=Y_P(U+1),\quad A=a_0+4,\quad
 Q=UY_P^2,\quad P=2Q+1,\quad J=2R+1.
\end{gathered}
\]
The supplied positive unknown $a$ continues to mean $a_0$.
Retain the first norm, affine index equation and positive interval:
\begin{equation}\label{eq:unit-scale-system}
\begin{gathered}
 \tau(\tau+1)=Q(Q+1)k^2,\qquad k=R+1+hQ,\\
 c=kY_P+\eta,\qquad k=\eta+\zeta,
 \qquad\tau,h,\eta,\zeta>0.
\end{gathered}
\end{equation}
Only the expanded source polynomials E9, E11 and E12 change, through
the definitions of $Q$ and $a_0$. The main parameter $A=a_0+4$,
signed auxiliary block and both exponent congruences retain their
forms in the supplied coordinates.

\subsection{A bound from the packing and the exact indices}

The pre-Pell coding proof gives
\[
 N=q^8,\quad q>4b\ge8,\quad 3L\le B\le N\le R,
 \quad b,q\le N,
 \quad0<S<N,\quad0\le T<N.
\]
The packed-index formula and integrality of the block values give
the useful additional bound
\begin{align}
 R&=S(N^2-N)+(T+1)(N^2-1)\notag\\
  &\le(N-1)(N^2-N)+N(N^2-1)
   =2N^3-2N^2<2N^3.\label{eq:unit-scale-rbound}
\end{align}
Thus $N,R\ge64$ and $U,Y_P\ge N^2\ge4096$. In particular
\begin{equation}\label{eq:unit-scale-initial}
 a_0>N^4>4N^3+1>J,\qquad
 Q\ge N^6>2N^3+1>R+1.
\end{equation}
These conclusions use only the coding bounds and positivity, before
either exponent relation is known.

The index equation gives $k\ge R+2$, and the interval gives
$c>Y_Pk>J$. The signed Pell block therefore implies
$c=\psi_A(J)$ and $d=\chi_A(J)$.
The first norm gives the ordinary Pell equation with root
$\widehat\tau=2\tau+1$ and parameter $P=2Q+1$, so
$k=\psi_P(t)$ for a positive index $t$.
The Pell congruence modulo $P-1=2Q$, reduced modulo $Q$, and
\eqref{eq:unit-scale-initial} give
\[
 t=R+1+vQ,\qquad v\ge0.
\]
If $v\ge1$, then Pell growth yields
\[
 \frac ck\le\frac{(2A)^{2R}}{(2P-1)^{R+Q}}
 \le\left(\frac{2A}{2P-1}\right)^{2R}<\frac12.
\]
Here $Q\ge R$, and the last inequality follows from
\[
 (2P-1)-4A=4Y_P\bigl(U(Y_P-1)-1\bigr)-15>0.
\]
This contradicts $c/k>Y_P$. Consequently
\begin{equation}\label{eq:unit-scale-indices}
 c=\psi_A(J),\quad d=\chi_A(J),\quad
 k=\psi_P(R+1),\quad\widehat\tau=\chi_P(R+1).
\end{equation}

\subsection{The larger error is bounded after exponent recovery}

Put $\xi=(U+1)^{2R}/U^R$. At the exact indices, the same cancellation
of $M$ in the principal ratio gives
\begin{equation}\label{eq:unit-scale-lower}
 \frac ck\ge
 \xi\left(1+\frac7{2a_0}\right)^{2R}
     \left(1+\frac1{2Q}\right)^{-R}>\xi,
\end{equation}
because $14Q>a_0$. The upper growth bound gives
\[
 \frac ck<\xi\left(1+\frac4{a_0}\right)^{2R}.
\]
The packed-index bound replaces the old factor $R$ in the scale:
\[
 \frac{8R}{a_0}<\frac{16N^3}{N^4}
 =\frac{16}{N}\le\frac14<\frac12.
\]
The elementary estimate $(1+t)^m\le(1-mt)^{-1}<1+2mt$
therefore gives
\begin{equation}\label{eq:unit-scale-upper}
 \frac ck<\xi\left(1+\frac{16R}{a_0}\right).
\end{equation}
The interval and the strict lower bound give $\xi<Y_P+1$.
Since $Y_P$ is integral and $\xi>U^R$, this implies
\begin{equation}\label{eq:unit-scale-size}
 Y_P\ge U^R,\qquad A>a_0=Y_P(U+1)>U^{R+1}.
\end{equation}
Also $\xi<Y_P+1<2Y_P$, so
\begin{equation}\label{eq:unit-scale-error-unbounded}
 0<\frac ck-\xi<\frac{32R}{U+1}.
\end{equation}
At this point the right-hand side is not assumed to be less than one.

Write $W=bw$. Since $W\le U$, $N\ge64$ and $U\ge4096$,
\eqref{eq:unit-scale-size} gives
\[
 4^{3J}\le64N^{2R}\le64U^R<U^{R+1}<A,
 \qquad W^3\le U^3\le U^R<A.
\]
The base-four exponent congruence now proves $W=4^J$ by source
Lemma~2.22. Hence $b$ is a power of two and
\[
 U=N^2w\ge Nbw=N4^J\ge64\cdot4^J>64R.
\]
Only now does~\eqref{eq:unit-scale-error-unbounded} yield
\begin{equation}\label{eq:unit-scale-small-error}
 0<\frac ck-\xi<\frac12.
\end{equation}

The gap, second Pell norm and index congruence give
$\kappa=\psi_A(L)$ and $\mu=\chi_A(L)$, since
$0<L<J<A$. Furthermore
\[
 B^{3L}\le B^B\le N^N\le U^R<A,
 \qquad q^3\le N^N\le U^R<A.
\]
The second exponent congruence therefore gives $q=B^L$.
Since also $U\ge N4^J>4\cdot2^{2R}$, the binomial expansion has
integer part $F=\lfloor\xi\rfloor$ with
\[
 0<\xi-F<\frac14,\qquad
 F\equiv\binom{2R}{R}\pmod U.
\]
The two inequalities
$F<\xi<c/k<Y_P+1$ and
$Y_P<c/k<\xi+1/2<F+3/4$
force $Y_P=F$. Thus $N^2\mid\binom{2R}{R}$, and the unchanged
no-carry and modular-Sidon decoding proves membership.

\subsection{Positive witnesses and the independent count}

For necessity, take the canonical coding witnesses, which determine
$b,B,q,N,R$ and satisfy~\eqref{eq:unit-scale-rbound}. Choose
\[
 J=2R+1,\quad w=4^J/b,\quad U=N^2w,\quad
 Y_P=\lfloor\xi\rfloor,\quad s=Y_P/N^2,
\]
and then form $a_0=Y_P(U+1)$, $A=a_0+4$, $Q=UY_P^2$,
$P=2Q+1$, $c=\psi_A(J)$, $d=\chi_A(J)$ and
$k=\psi_P(R+1)$.
The positive integrality of $w,s$ follows from the same power-of-two
and binomial arguments as in Section~\ref{sec:main-base-four}.
The bounds~\eqref{eq:unit-scale-lower}--\eqref{eq:unit-scale-upper}
apply at these indices, and $U\ge N4^J$ holds directly from the
chosen $w$. Hence
\[
 Y_P<\xi<\frac ck<\xi+\frac12<Y_P+\frac34.
\]
Choose the positive integers
\[
 \eta=c-kY_P,\qquad \zeta=k-\eta,\qquad
 \tau=\frac{\chi_P(R+1)-1}{2},\qquad
 h=\frac{k-R-1}{Q}.
\]
The oddness of $P$ gives the integrality of $\tau$, and the Pell
congruence modulo $2Q$ gives that of $h$. Strict Pell growth gives
their positivity.

All remaining dependent Pell witnesses are obtained by the generic
constructions of Section~\ref{sec:main-base-four} at this new $A$.
Explicitly, choose positive $f,i$ for E16 and put
\[
 G=1+(A+1)(f^2-1),\qquad
 o=\frac{\chi_G(J)+d}{f},\qquad
 j=\frac{\psi_G(J)-J}{c}.
\]
The congruences $G\equiv-A\pmod f$, $G\equiv1\pmod c$ and oddness
of $J$ make $o,j$ positive integers. Choose
$\kappa=\psi_A(L)$, $\mu=\chi_A(L)$,
$\Delta=(\kappa-L)/(A-1)$ and $\phi=c-\kappa$; these are
positive integers because $2\le L<J$.
The exponent quotients are positive since
\[
 d-(A-4)c>3c>4^J,\qquad
 \mu-(A-B)\kappa>(B-1)\kappa>B^L,
\]
using $A>(4^J)^3,(B^L)^3$ and $J,L\ge2$.
Their moduli are positive. All coding witnesses remain unchanged.
This constructs every required positive coordinate at the new scale.

\begin{theorem}\label{thm:unit-scale105}
For each r.e.\ set, reusing $Y_P$ as the Pell scale gives a
membership-equivalent system in 34 positive unknowns and 22 equations
with a straight-line certificate of \textbf{105 arithmetic instructions}:
58 multiplications and 47 additions. Fixed numerals, supplied parameters
and equality tests are free.
\end{theorem}
\begin{proof}
The preceding arguments establish soundness and positive necessity.
The predecessor already computes $Y_P=sN^2$. Delete the multiplication
$M=RY_P$ and reuse $Y_P$ wherever $M$ was used. The three instructions
\[
 UY_P=U\cdot Y_P,\qquad a_0=UY_P+Y_P,
 \qquad Q=(UY_P)Y_P
\]
retain their previous costs. The affine index equation retains its
multiplication and addition. Thus exactly one multiplication disappears
from the 106-operation certificate. The complete residual and primitive
checks are supplied by
\file{round18\_1980\_unit\_scale\_certificate.py}.
\end{proof}

This is independent of Theorem~\ref{thm:best105}: its saving is a
multiplication rather than an addition. The index-multiple argument
there cannot be applied automatically here, since $P-1=2UY_P^2$ is
not known to be divisible by $R+1$.

\subsection{Sharing the interval product with the norm}

The smaller scale permits another multiplication to be removed.
Let
\[
 D=UY_P,\qquad Q=DY_P=UY_P^2,\qquad H_P=Y_Pk.
\]
Both $D$ and $H_P$ are already computed: $D$ is needed for
$a_0=D+Y_P$, and $H_P$ is needed for $c=H_P+\eta$.
The polynomial identity
\begin{equation}\label{eq:shared-pell-norm}
 Q(Q+1)k^2
 =(UY_P^2)(UY_P^2+1)k^2
 =\bigl((UY_P)^2+U\bigr)(Y_Pk)^2
 =(D^2+U)H_P^2
\end{equation}
allows the norm to share these two products.
To remove the remaining use of the computed value $Q$, replace
the index equation in~\eqref{eq:unit-scale-system} by
\begin{equation}\label{eq:shared-pell-index}
 k=R+1+hD,\qquad h>0.
\end{equation}
The norm and every other source polynomial are unchanged.

We first prove that the smaller modulus still gives the exact index.
Before exponent decoding, the bounds above give
\[
 D=UY_P\ge N^4>2N^3+1>R+1,\qquad A>J.
\]
The new equation gives $k\ge R+2$, so the interval gives $c>J$.
The signed Pell block gives $c=\psi_A(J)$ and $d=\chi_A(J)$.
The unchanged triangular norm gives $k=\psi_P(t)$ for some
positive index $t$, where $P=2Q+1$. Since
$P-1=2Q=2DY_P$ is divisible by $D$, the Pell congruence and
\eqref{eq:shared-pell-index} imply
\[
 t=R+1+vD,\qquad v\ge0.
\]
If $v\ge1$, then $D>R$ and the same growth estimate gives
\[
 \frac ck\le\frac{(2A)^{2R}}{(2P-1)^{R+D}}
 \le\left(\frac{2A}{2P-1}\right)^{2R}<\frac12,
\]
contradicting $c/k>Y_P$. Therefore $t=R+1$.
This recovery uses only the preliminary coding bounds, the interval
and the two Pell blocks.

There are direct maps between the positive witnesses of this system
and those of Theorem~\ref{thm:unit-scale105}. From an old solution,
put $h=Y_Ph_{\rm old}$ and preserve every other witness. Then
$hD=h_{\rm old}Q$, so the new index equation holds.
Conversely, the recovered exact index and the stronger Pell congruence
modulo $2Q$ show that
\[
 h_{\rm old}=\frac{k-R-1}{Q}
\]
is an integer. It is positive because $\psi_P(R+1)>R+1$.
It restores the old index equation, preserving every other witness.
Equivalently, the new $h$ is a positive multiple of $Y_P$ and the
reverse map is $h_{\rm old}=h/Y_P$.
This proves equivalence over the positive domains before invoking
the unit-scale approximation, exponent or rounding conclusions.
Those conclusions and the complete decoding proof therefore transfer.

\begin{theorem}\label{thm:best104}
For each r.e.\ set, the fixed modular-Sidon encoding gives a
membership-equivalent system in 34 positive unknowns and 22 equations
whose straight-line certificate has \textbf{104 arithmetic instructions}:
57 multiplications and 47 additions. Fixed numerals, supplied parameters
and equality tests are free.
\end{theorem}
\begin{proof}
Positive-domain equivalence follows from the preceding witness maps
and Theorem~\ref{thm:unit-scale105}. For the count, move the existing
computation of $H_P=Y_Pk$ before the norm. The former five operations
\[
 Q=DY_P,\quad Q_+=Q+1,\quad F_Q=QQ_+,
 \quad k_2=k^2,\quad L_9=F_Qk_2
\]
are replaced by four:
\[
 D_2=D^2,\quad F_D=D_2+U,\quad H_2=H_P^2,
 \quad L_9=F_DH_2.
\]
Identity~\eqref{eq:shared-pell-norm} proves that they give the same
norm residual. The computation of $H_P$ is moved rather than copied,
and the root-side instructions $\tau+1$ and $\tau(\tau+1)$ are
unchanged. Both displayed lists contain one addition; the second
contains three multiplications rather than four.

The product $h_{\rm old}Q$ is replaced by $hD$ at the same cost.
Thus the calculated register $Q$ is no longer required anywhere in
the certificate. Its mathematical definition in the proof introduces
no additional supplied witness or uncharged computed register.
The total is 57 multiplications and 47 additions. The checker
\file{round19\_1980\_shared\_ratio\_product\_certificate.py}
verifies the identity, the forward witness-map polynomial, all source
residuals and primitive instructions, and the absence of the discarded
registers. It checks explicitly that only E11 changes as a source
polynomial.
\end{proof}

% One offset suffices when the target coefficients absorb a possible borrow.
\section{Allowing a borrow in the coefficient tests}\label{sec:borrow-mask}

The low part of the coefficient-testing offset can be removed. The
resulting digits may receive a borrow from the preceding position, but
the residuals can be encoded so that both possible digits give the same
answer. Starting from Theorem~\ref{thm:best104}, replace
\[
 S_3=B\lambda(1+q)-(Z\lambda-2e)\mathcal C^2
\]
by
\begin{equation}\label{eq:borrow-third-block}
 S_3=B\lambda q-(Z\lambda-2e)\mathcal C^2,
 \qquad \sigma=S_3>0.
\end{equation}
Here $\mathcal C=x+g$ is the digit code, distinguished from the main
Pell denominator $c$. The supplied positive witness $\sigma$ remains
part of the system. The change removes the instruction computing
$q+1$. It requires a new fixed coefficient index, constructed below.

\subsection{Coefficient targets and the admissible index}

Keep the modular support of Section~\ref{sec:modular-support}, including
the input at weight zero, $\delta$ at $v_0=1$, the original 58
coordinates at $v_1,\ldots,v_{58}$ and the guard coordinate $u$ at
$v_{59}$. In particular,
\[
\begin{gathered}
 v_i=1+3\bigl(122i+(i^2\bmod61)\bigr),\qquad 0\le i\le59,\\
 v_*=21607,\qquad t_j=158467571+86429j\quad(0\le j<1832),\\
 t_*=316719070,\qquad K=t_*+1=316719071,\qquad
 L=2^{32}>3K+2.
\end{gathered}
\]
The 1830 original homogeneous quadratic residuals come from the
integer row basis and padding of Section~\ref{sec:quadratic-masks}.
Adjoin the guard residual $u\delta-x^2$ as the 1831st ordinary row.
For each ordinary row replace its coefficient target by
\begin{equation}\label{eq:borrow-targets}
 G_j=4F_j^h+\delta^2\qquad(0\le j<1831),
 \qquad G_{1831}=\delta^2.
\end{equation}
Thus $G_{1830}=4(u\delta-x^2)+\delta^2$, and the last row remains
the special square test.

For a homogeneous quadratic monomial $I$, let $c_I\in\{1,2\}$ be its
coefficient in $\mathcal C^2$: it is 1 for a square and 2 for a
product of distinct coordinates. If $a_{j,I}$ is its coefficient in
$G_j$, define
\[
 \mathcal D(B)=\sum_{j,I}\frac{a_{j,I}}{c_I}
                   B^{t_j-w(I)}.
\]
Every displayed coefficient is integral. The coefficients from
$4F_j^h$ are divisible by both possible denominators, and the added
$\delta^2$ is a square with denominator 1. In the special row the
coefficient 1 is placed at $t_{1831}-2$.

The support separation proved in Section~\ref{sec:modular-support}
gives the exact identity
\begin{equation}\label{eq:borrow-coefficient-identity}
 [B^{t_j}]\,\mathcal D(B)\mathcal C(B)^2=G_j.
\end{equation}
Indeed, distinct monomials have distinct weights, different row
intervals are disjoint, and a product involving a dummy row coordinate
lies above every row target. The least position of $\mathcal D$ is
above $v_*$, so all low coordinate tests remain automatic.

Choose a power of two $z\ge2$ strictly above the absolute values of
the new coefficients of $\mathcal D$, and put $Z=2z$. Define
\begin{align*}
 \ell_0(B)&=\sum_{i=0}^{59}B^{v_i}+\sum_{j=0}^{1831}B^{t_j},\\
 e_0(B)&=z\sum_{j=0}^{K-1}B^j+\mathcal D(B),\\
 V&=\ell_0(Z)+e_0(Z)Z^L.
\end{align*}
All digits of $e_0$ below $K$ lie in $[1,Z-1]$; higher digits vanish.
Choose a power of two
\begin{equation}\label{eq:borrow-index-bound}
 H>\max\{2Z^{2L+1},\ 4^{t_*+3}Z\,1900^2,\ 3L,\ 16\}.
\end{equation}
The positions, coordinate count and size formula for $H$ are unchanged.
The indices $Z,V,H$ are rebuilt for the new coefficient polynomial.
They depend effectively on the represented set and are independent of
the queried positive input $x$.

\subsection{Bounds before the coefficients are decoded}

Retain the coding equations and positive conditions
\begin{align*}
 B&=Hb^2,& b&=x+\beta,& \lambda(B-1)&=q^2-1,& \theta&=B-Z,\\
 S_2&=\ell+eq=V+t\theta,& S_2+\alpha&=q^2,&
 \Omega&=Z\lambda-2e>0,& \mathcal C&=x+g.
\end{align*}
The packing remains
\begin{align*}
 n&=q^8,& S&=g+q^2(S_2+q^2S_3),\\
 T+1&=q^2(1+\theta\lambda)-(b-1)\ell+(B-4)\ell q^4,\\
 r&=S(n^2-n)+(T+1)(n^2-1).
\end{align*}
The Pell interval scale $Y_P=sn^2$ is distinct from the packed code
$S_2$.

All preliminary upper bounds from Sections~\ref{sec:final-two-savings}
and~\ref{sec:unit-scale} still apply, since the new positive offset
$B\lambda q$ is smaller than $B\lambda(1+q)$. More explicitly,
$S_2<q^2$ gives $e<q$ and $\ell<q^2$, while the geometric equation
gives $B\le q^2$. Positivity of $\Omega$ and $\sigma$ gives
\[
 \mathcal C^2<B\lambda q<3q^3<q^4,
 \qquad 0<S_3<q^4.
\]
Thus $g<q^2$ and the same normalized masks give
\[
 0<S<n,\qquad0\le T<n,\qquad
 n\le r\le2n^3-2n^2<2n^3.
\]
The Pell proof of Theorem~\ref{thm:best104} uses these bounds and the
positive interval, independently of the coefficient values in the
encoding. It therefore recovers $b$ as a power of two, $q=B^L$, and
$n^2\mid\binom{2r}{r}$.

After $q=B^L$ is known, $\lambda>q$ and $e<q$ imply
$\Omega>Z\lambda/2$. Equation~\eqref{eq:borrow-third-block} then gives
\begin{equation}\label{eq:borrow-code-bound}
 \mathcal C^2<\frac{2Bq}{Z}\le\frac{Bq}{2}<q^2.
\end{equation}
Hence $g<\mathcal C<q$ before either the input code or the quotient
of the packed polynomial is decoded.

The first two masks and the packed congruence have not changed.
The borrow-removal and base-$Z$ uniqueness arguments of
Sections~\ref{sec:reversed-packing} and~\ref{sec:input-unit} require
only the bounds just proved, the support of $\ell_0$, the digit bounds
for $e_0$, and~\eqref{eq:borrow-index-bound}. They therefore give
\begin{equation}\label{eq:borrow-alias}
 S_2=\ell_0(B)+e_0(B)q,\qquad
 \ell=\ell_0(B)+mq,\qquad e=e_0(B)-m,\qquad 0\le m<e_0(B).
\end{equation}
The first mask, using $g<q$, decodes the input digit exactly as $x$.
All other coordinates are nonnegative integers less than $b$; dummy
coordinates can occur only at the row positions. Consequently
\begin{equation}\label{eq:borrow-polynomial-bounds}
 \mathcal C<B^K,\qquad
 A_C:=\mathcal C(1)^2<(1900b)^2,\qquad
 0<e<B^K,\qquad 2e\mathcal C^2<q.
\end{equation}
The last estimate uses $L>3K+2$ and is valid before $m$ is excluded.

\subsection{The high digits still exclude the quotient ambiguity}

\begin{lemma}\label{lem:borrow-high-window}
Under the preceding hypotheses, the third no-carry condition forces
$m=0$ in~\eqref{eq:borrow-alias}.
\end{lemma}
\begin{proof}
Put $J_C=ZA_C<B/4$ and
\[
 W_C=Z\lambda\mathcal C^2=qW_1+W_0,
 \qquad0\le W_0<q.
\]
The coefficients of the polynomial $Z\lambda\mathcal C^2$ are
nonnegative and at most $J_C$. They are already its base-$B$ digits,
so no internal carries occur in this product. Since
$\deg\mathcal C^2\le2K-2$, these coefficients equal $J_C$ throughout
the interval $[2K-2,2L-1]$. In particular they equal $J_C$ at every
position $L,\ldots,L+K$, because $L>3K+2$.

By~\eqref{eq:borrow-polynomial-bounds}, exact division by $q$ gives
\[
 \left\lfloor\frac{S_3}{q}\right\rfloor
  =B\lambda-W_1+\varepsilon,
 \qquad
 \varepsilon=\left\lfloor
   \frac{2e\mathcal C^2-W_0}{q}\right\rfloor\in\{-1,0\}.
\]
The unit digit of $B\lambda$ is zero, and its digits at positions
$1,\ldots,2L$ are one. The first $K+1$ digits of $W_1$ are all
$J_C$. At the unit position the subtraction therefore gives digit
$B-J_C-1$ or $B-J_C$ and an outgoing borrow of one. At each following
position through $K$, the incoming borrow changes $1-J_C$ into
$-J_C$, giving digit $B-J_C$ and again an outgoing borrow of one.
Thus the digits of $S_3$ at positions $L,\ldots,L+K$ are all at least
$B-J_C-1$.

The product $(B-4)\ell_0(B)$ has digits $B-4$ at its support
positions and no internal carries, and it is less than $q$.
Accordingly the part of the third mask at and above position $L$ is
\[
 X=(B-4)m<B^{K+1}.
\]
If a digit $d$ has no binary overlap with a digit at least
$B-J_C-1$, then
$d\le B-1-(B-J_C-1)=J_C$.
Hence every digit of $X$ is at most $J_C$.

Let $F$ be the digit polynomial of $X$, so $F(B)=X$, its coefficients
are nonnegative, and $\deg F\le K$. Since $B-4\mid F(B)$, reduction
modulo $B-4$ gives $B-4\mid F(4)$. On the other hand,
\[
 0\le F(4)\le J_C\frac{4^{K+1}-1}{3}
       <\frac{B}{12}<B-4.
\]
Here~\eqref{eq:borrow-index-bound} and
$J_C<Z(1900b)^2$ give
$4^{K+2}J_C<B$, because $K=t_*+1$.
It follows that $F(4)=0$. Nonnegativity of its coefficients then
gives $F=0$, so $X=0$ and $m=0$.
\end{proof}

\subsection{Low carries determine the intended equations}

We now have $\ell=\ell_0(B)$ and $e=e_0(B)$.
At positions below $K$, the offset $qB\lambda$ makes no contribution,
and the negative tail of $e_0-z\lambda$ begins too high to contribute.
The raw coefficients there are precisely those of
$2\mathcal D(B)\mathcal C(B)^2$. Each has absolute value less than
$2zA_C=J_C<B/4$.

Normalize these signed coefficients starting at the unit position.
The initial carry is zero. If an incoming carry is in $\{-1,0\}$,
adding a raw coefficient of absolute value less than $B/4$ gives a
number strictly between $-B$ and $B$. Its Euclidean quotient by $B$
is therefore again -1 or 0. Induction proves that these are the only
incoming carries at every target position.

At a target, let $v$ be the raw coefficient and let
$\varepsilon\in\{-1,0\}$ be the incoming carry. The mask digit
$B-4$ permits exactly the digits $0,1,2,3$. Since $|v|<B/4$, a
negative $v+\varepsilon$ would normalize to a digit greater than 3,
and a nonnegative one cannot wrap around $B$. Thus the target test is
exactly
\begin{equation}\label{eq:borrow-local-test}
 0\le v+\varepsilon\le3.
\end{equation}
At the special row, $v=2\delta^2$. The test implies
$2\delta^2\le4$, so the nonnegative integer $\delta$ is 0 or 1.
At the guard, the raw value is
\[
 v=8(u\delta-x^2)+2\delta^2.
\]
If $\delta=0$, this is $-8x^2<0$, and a carry of -1 or 0 cannot
make it satisfy~\eqref{eq:borrow-local-test}. Hence $\delta=1$.

Every ordinary row now has $v=8F_j^h+2$. Depending on the carry,
\eqref{eq:borrow-local-test} says
\[
 0\le8F_j^h+1\le3
 \quad\hbox{or}\quad
 0\le8F_j^h+2\le3.
\]
Both imply $F_j^h=0$ by integrality. Since $\delta=1$, these are
exactly the original equations. At each low variable position the
raw coefficient and incoming carry are zero, since the least position
of $\mathcal D$ lies above all such positions. Those tests impose no
additional condition. This proves soundness of the new encoding.

\subsection{Positive witnesses and the resulting certificate}

Conversely, start with a solution of the represented system. Set
$\delta=1$, $u=x^2$, and every dummy coordinate to zero. Choose a
sufficiently large power of two $b$, and take the canonical
$B,q,\lambda,\ell_0,e_0,\mathcal C$ defined above. The positive digit
at weight one makes $g>0$. All previous fixed-index, packed-congruence,
bound and $\Omega$ witnesses remain positive, because the support and
digit bounds have the same form.

The changed witness $\sigma=S_3$ is positive as well. Indeed,
$\mathcal C<B^K$, $Z<B$ and $L>2K$ give
\[
 Z\mathcal C^2<B^{2K+1}<qB,
 \qquad
 S_3=\lambda(qB-Z\mathcal C^2)+2e\mathcal C^2>0.
\]
At every ordinary row and the special row the raw coefficient is
exactly 2. Its actual digit is consequently 1 or 2, both allowed by
the mask. The low coordinate tests have digit zero, and $\ell_0$ has
no other support. The first two masks work as before. Thus the
packed integers have no binary overlap, yielding
$n^2\mid\binom{2r}{r}$ for their newly calculated $r$.

All dependent positive Pell witnesses are chosen by the constructions
of Section~\ref{sec:unit-scale}, using this $r$. Those constructions
require only the packed ranges, the binomial divisibility, and the
fixed exponent bounds, which have just been established. This proves
necessity without identifying the old and new values of $r$.

\begin{theorem}\label{thm:best103}
For each r.e.\ set, the modified coefficient targets
\eqref{eq:borrow-targets} and admissible index above give a
membership-equivalent system in 34 positive unknowns and 22 equations
with a straight-line certificate of \textbf{103 arithmetic instructions}:
57 multiplications and 46 additions. Fixed numerals, supplied parameters
and equality tests are free.
\end{theorem}
\begin{proof}
The preceding arguments prove soundness and positive necessity.
The product $\mathit{Lam}=B\lambda$ is already required by the
geometric equation. Replace the instruction
$P_2=\mathit{Lam}(q+1)$ by $P_2=\mathit{Lam}\,q$ and delete the
single addition computing $q+1$. All other instruction counts remain
unchanged. The coefficient adjustment changes only the supplied fixed
index; it introduces no system unknown or arithmetic operation.
The exact primitive and source-residual identities, including the two
inherited triangular corrections, are checked by
\file{round20\_1980\_borrow\_mask\_certificate.py}.
The full high-window argument is Lemma~\ref{lem:borrow-high-window};
it needs no stronger bound on $H$ than the preceding construction.
\end{proof}

% The binomial scale itself is the target of the base-four congruence.
\section{A common witness for the binomial scale and exponent}
\label{sec:common-witness}

The first exponent test can use the binomial scale $U$ itself as its
target. That value is already computed, so the separate product $bw$
disappears. This change preserves the unit-scale norm and its index
modulus. Combined with the coefficient encoding of
Section~\ref{sec:borrow-mask}, it gives 102 operations.

Use the fixed coefficient index of Theorem~\ref{thm:best103}, including
its ordinary targets $4F_i^h+\delta^2$, special target $\delta^2$,
and exponent $L=2^{32}$. Retain the notation
\[
\begin{gathered}
 N=n,\quad R=r,\quad U=wN^2,\quad Y_P=sN^2,\\
 a_0=Y_P(U+1),\quad A=a_0+4,\quad
 D_P=UY_P,\quad Q=UY_P^2,\quad P=2Q+1,\quad J=2R+1.
\end{gathered}
\]
The positive supplied unknown $a$ still means $a_0$.
Keep the shared norm, first index equation and interval
\begin{equation}\label{eq:common-witness-pell}
\begin{gathered}
 \tau(\tau+1)=(D_P^2+U)(Y_Pk)^2,\qquad
 k=R+1+hD_P,\\
 c=kY_P+\eta,\qquad k=\eta+\zeta,
 \qquad\tau,h,\eta,\zeta>0.
\end{gathered}
\end{equation}
Change only E14 among the supplied-coordinate Pell polynomials,
replacing its target $bw$ by $U$:
\begin{equation}\label{eq:common-witness-exponent}
 d=U+c(A-4)+\gamma(8A-17).
\end{equation}
No parameter or witness is added. The defining identity $U=wN^2$
and the index modulus $D_P=UY_P$ are unchanged.

\subsection{Exact indices before the exponent equations}

The smaller positive offset in Section~\ref{sec:borrow-mask} retains
all preliminary coding bounds:
\[
\begin{gathered}
 N=q^8,\qquad q>4b\ge8,\qquad 3L\le B\le N\le R,\\
 0<S<N,\quad0\le T<N,\qquad
 R\le2N^3-2N^2<2N^3.
\end{gathered}
\]
These are established without using any Pell equation. Hence
$N,R\ge64$, $U,Y_P\ge N^2$, $a_0>N^4>J$, and
$D_P\ge N^4>R+1$.

The first index equation gives $k\ge R+2$, and the interval gives
$c>Y_Pk>J$. The signed block therefore gives
$c=\psi_A(J)$ and $d=\chi_A(J)$.
The norm in~\eqref{eq:common-witness-pell} is the triangular form
of the ordinary Pell equation with parameter $P=2Q+1$, so
$k=\psi_P(t)$ for some positive index $t$.
Since $P-1=2D_PY_P$ is divisible by $D_P$, the Pell congruence
and $D_P>R+1$ imply
\[
 t=R+1+vD_P,\qquad v\ge0.
\]
If $v\ge1$, then $D_P>R$ and Pell growth gives
\[
 \frac ck\le\frac{(2A)^{2R}}{(2P-1)^{R+D_P}}
 \le\left(\frac{2A}{2P-1}\right)^{2R}<\frac12.
\]
The last inequality follows from
$(2P-1)-4A=4Y_P[U(Y_P-1)-1]-15>0$.
This contradicts $c/k>Y_P$. Thus the exact indices are
\begin{equation}\label{eq:common-witness-indices}
 c=\psi_A(J),\quad d=\chi_A(J),\quad
 k=\psi_P(R+1),\quad 2\tau+1=\chi_P(R+1).
\end{equation}
The change to E14 has not entered this argument.

\subsection{Recovering the radices in the required order}

Put $\xi=(U+1)^{2R}/U^R$. The unit-scale lower estimate at
\eqref{eq:common-witness-indices} gives
\[
 \frac ck\ge
 \xi\left(1+\frac7{2a_0}\right)^{2R}
     \left(1+\frac1{2Q}\right)^{-R}>\xi,
\]
using $14Q>a_0$. The interval gives $\xi<Y_P+1$, so
\begin{equation}\label{eq:common-witness-size}
 Y_P\ge U^R,\qquad A>a_0=Y_P(U+1)>U^{R+1}.
\end{equation}
As in Section~\ref{sec:unit-scale}, the independent bound
$R<2N^3$ gives $8R/a_0<16/N<1/2$. The upper growth estimate and
$(1+t)^m\le(1-mt)^{-1}<1+2mt$ therefore yield
\begin{equation}\label{eq:common-witness-error}
 0<\frac ck-\xi<\frac{32R}{U+1}.
\end{equation}
This error need not be bounded by an absolute constant yet.

The changed exponent target satisfies
\[
 4^{3J}=64\cdot4096^R
 \le64N^{2R}\le64U^R<U^{R+1}<A,
 \qquad U^3\le U^R<A.
\]
The source $\chi$ exponent criterion applied to
\eqref{eq:common-witness-exponent} now gives
\begin{equation}\label{eq:common-witness-power}
 U=4^J.
\end{equation}
In particular $U>64R$, so~\eqref{eq:common-witness-error} is
strictly less than $1/2$.

The new first exponent conclusion determines the radices in a
different order. Since $N^2$ divides the power of two $U=wN^2$,
$N$ is a power of two. Since $N=q^8$, so is $q$.
No conclusion about $b$ is needed yet.

The unchanged second norm, gap and index congruence identify
$\kappa=\psi_A(L)$ and $\mu=\chi_A(L)$, since $0<L<J<A$.
Their exponent test has the size bounds
\[
 B^{3L}\le B^B\le N^N\le U^R<A,
 \qquad q^3\le N\le R<U^R<A.
\]
It therefore gives $q=B^L$. Now $B$ is a power of two, because
$q$ is a power of two and $L>0$. Finally $B=Hb^2$ implies that
$b$ is a power of two. Thus all binary radix facts are known before
the first-mask decoding uses them.

The binomial expansion has integer part $F=\lfloor\xi\rfloor$ with
\[
 0<\xi-F<\frac{2^{2R}}U<\frac14,
 \qquad F\equiv\binom{2R}{R}\pmod U.
\]
The interval and the ratio bounds give
\[
 F<\xi<\frac ck<Y_P+1,
 \qquad Y_P<\frac ck<\xi+\frac12<F+\frac34.
\]
Integrality forces $Y_P=F$. The unchanged definitions
$U=wN^2$ and $Y_P=sN^2$ now imply
\begin{equation}\label{eq:common-witness-binomial}
 N^2\mid\binom{2R}{R}.
\end{equation}

\subsection{Composition with the new coefficient encoding}

The conclusions just obtained are exactly those needed by
Section~\ref{sec:borrow-mask}: $b,q,N$ are powers of two,
$q=B^L$, and~\eqref{eq:common-witness-binomial} supplies binary
nonoverlap of the packed blocks. Positivity of the new offset then
gives
\[
 \Omega>Z\lambda/2,\qquad
 \mathcal C^2<2Bq/Z<q^2,
 \qquad\mathcal C=x+g.
\]
The first two masks recover the possible quotient ambiguity
$\ell=\ell_0+mq$, $e=e_0-m$ and the bounds required by
Lemma~\ref{lem:borrow-high-window}. That lemma forces $m=0$.
The low coefficients have incoming carries in $\{-1,0\}$; the
special row and guard force $\delta=1$, and every ordinary row
then forces $F_i^h=0$. These decoding arguments do not use E14
internally. They apply unchanged once its replacement has supplied
the required exponent and binomial conclusions, proving sufficiency.

For necessity, first use the new canonical coding construction of
Section~\ref{sec:borrow-mask}. It has $\delta=1$, $u=x^2$, zero
dummy coordinates and a sufficiently large power of two $b$.
Its changed $\sigma$ is positive because
\[
 S_3=\lambda(Bq-Z\mathcal C^2)+2e\mathcal C^2>0.
\]
The ordinary and special target digits are one or two, so the masks
pass. Form its resulting $S,T,R$. This supplies the packed bounds
and~\eqref{eq:common-witness-binomial} before choosing Pell witnesses.

Now set
\[
\begin{gathered}
 J=2R+1,\quad U=4^J,\quad w=U/N^2,\quad
 Y_P=\lfloor\xi\rfloor,\quad s=Y_P/N^2,\\
 a_0=Y_P(U+1),\quad A=a_0+4,\quad Q=UY_P^2,
 \quad P=2Q+1,\\
 c=\psi_A(J),\quad d=\chi_A(J),\quad k=\psi_P(R+1).
\end{gathered}
\]
The quotient $w$ is positive integral: canonical $N$ is a power of
two and $N^2\le R^2\le4^R<4^J=U$. The binomial expansion and
divisibility make $s$ positive integral. The same estimates at these
indices give $Y_P<c/k<Y_P+3/4$, so choose the positive integers
\[
 \eta=c-kY_P,\quad \zeta=k-\eta,\quad
 \tau=\frac{\chi_P(R+1)-1}{2},\quad
 h=\frac{k-R-1}{UY_P}.
\]
Odd $P$ gives the positive integrality of $\tau$. The congruence
modulo $P-1=2UY_P^2$ and strict Pell growth give that of $h$.
This construction needs no parity condition on $R$.

Choose the remaining Pell witnesses at the new $A$ by the generic
constructions already established in Section~\ref{sec:main-base-four}.
After choosing positive $f,i$ for E16, put
\[
 G=1+(A+1)(f^2-1),\qquad
 o=\frac{\chi_G(J)+d}{f},\qquad
 j=\frac{\psi_G(J)-J}{c},
\]
and take $\kappa=\psi_A(L)$, $\mu=\chi_A(L)$,
$\Delta=(\kappa-L)/(A-1)$ and $\phi=c-\kappa$.
Their congruences and strict growth make every required witness
positive integral. The exponent quotients are positive because
\[
 d-(A-4)c>3c>U,\qquad
 \mu-(A-B)\kappa>(B-1)\kappa>q,
\]
using $A>U^3,q^3$ and $J,L\ge2$; their moduli are positive.
Thus the new E14 and every other equation hold. The coding witnesses
are chosen first, and all dependent Pell coordinates are then chosen
for their new packed $R$.

\begin{theorem}\label{thm:best102}
For each r.e.\ set, the fixed coefficient encoding of
Section~\ref{sec:borrow-mask} gives a membership-equivalent system
in 34 positive unknowns and 22 equations whose straight-line
certificate has \textbf{102 arithmetic instructions}:
56 multiplications and 46 additions. Fixed numerals, supplied
parameters and equality tests are free.
\end{theorem}
\begin{proof}
The preceding arguments establish the complete positive-domain
equivalence and the order of the two constructions.
Starting from the 104-operation schedule, the offset change deletes
the addition $q+1$ and uses $(B\lambda)q$ in its place.
The common-witness change deletes the product $bw$ and uses the
existing register $U=wN^2$ directly in E14. These edits concern
disjoint register definitions, so they remove one addition and one
multiplication without introducing an instruction, unknown or equation.

The resulting count is 56 multiplications and 46 additions.
The rebuilt coefficient index affects only supplied constants.
The checker \file{round22\_1980\_composed\_certificate.py}
verifies the complete primitive list, all source residuals and both
triangular corrections. It also checks that the two schedule edits
commute and that both discarded registers are absent.
\end{proof}

% A coefficient alphabet independent of the represented set.
\section{Fixing the coefficient digit base at four}\label{sec:fixed-four}

The coefficient digit base can be fixed at $Z=4$ by expanding the
represented equations into elementary arithmetic gates. Their number
affects the fixed index, but not the number of unknowns in the final
system: the additional coordinates are digits of the existing code.
The exponent $L$ becomes a fixed index parameter in place of $Z$.
This identifies the already supplied $\theta=B-4$ with a separately
computed register of Theorem~\ref{thm:best102}, saving one addition.

\subsection{Primitive rows with a fixed coefficient alphabet}

Start with a finite polynomial system representing the desired set
over nonnegative witnesses and a positive input $x$. Split each equation
into an equality of polynomials with nonnegative coefficients, and
evaluate those polynomials by finite circuits of additions,
multiplications, copies, and constants zero and one. Every gate has a
fresh nonnegative output coordinate. Arbitrary integer coefficients are
implemented by finite arithmetic from one. Thus this construction
places no bound on the coefficients of the original system.

Replace repeated operands of an addition or multiplication by distinct
copied coordinates. With a separate homogenizing coordinate $\delta$,
the quadratic residuals have the following forms:
\[
\begin{array}{c|l}
\text{gate or constraint}&\text{homogeneous residual }F\\ \hline
\text{multiplication}&XY-W\delta\\
\text{addition}&(X+Y-W)\delta\\
\text{copy or equality}&(X-Y)\delta\\
\text{zero}&X\delta\\
\text{unit}&X\delta-\delta^2.
\end{array}
\]
Here $X,Y$ may include the input $x$, but are distinct from $\delta$;
outputs are fresh, and trivial identical equalities are omitted.
At $\delta=1$, these rows express the original circuit exactly.
Every solution of the original system extends by its nonnegative gate
values. Adjoin a fresh coordinate $u$ and the guard
$F_{\mathrm g}=u\delta-x^2$.

For every ordinary row, including the guard, use the target
\begin{equation}\label{eq:four-targets}
 G_j=2F_j+\delta^2,
 \qquad G_{\mathrm{special}}=\delta^2.
\end{equation}
For a homogeneous quadratic monomial $I$, let $c_I=1$ for a square
and $c_I=2$ for a product of distinct coordinates. Every coefficient
of every $G_j$, divided by $c_I$, belongs to
\begin{equation}\label{eq:four-coefficient-alphabet}
 \{-2,-1,0,1\}.
\end{equation}
Indeed, ordinary cross terms give $\pm1$, and the added $\delta^2$
gives $1$. The unit row gives $2X\delta-\delta^2$, with divided
coefficients $1,-1$. The guard gives
$2u\delta-2x^2+\delta^2$, with divided coefficients $1,-2,1$.
The special target has coefficient $1$. Distinct copied operands
prevent a repeated positive square or coincident cross terms from
enlarging this alphabet. The displayed orientation of the unit row
is part of the construction.

\subsection{Support and admissible indices of arbitrary finite size}

Let $m$ be the number of coordinates other than the input, before
adding dummy row coordinates. It includes $\delta$, all circuit
coordinates, and $u$. Give $x$ weight zero, and give the remaining
coordinates weights
\[
 v_i=3^i\quad(0\le i<m),\qquad v_0=1\text{ for }\delta,
 \qquad M=3^{m-1}.
\]
All quadratic monomial weights are distinct. The weights are zero,
a single $3^i$, or a sum of two such powers; their base-three digits
recover the unordered pair, including repeated coordinates and the
weight-zero input.

Let $s$ be the total number of targets, including the guard and the
special target. Define
\begin{align*}
 d_0&=4M+1,& t_0&=(s+1)d_0+2M,\\
 t_j&=t_0+jd_0\quad(0\le j<s),&
 t_*&=2sd_0+2M,& K&=t_*+1.
\end{align*}
Choose a power of two $L>3K+2$. The coefficient intervals
$[t_j-2M,t_j]$ are disjoint and lie above $M$. Moreover
\[
 2t_0-2M>t_*.
\]
Include a dummy coordinate at each row position $t_j$. This last
inequality makes every contribution involving a dummy lie above all
targets, so the dummies cannot alter a tested residual.

Write $a_{j,I}$ for the coefficient of monomial $I$ in $G_j$, and set
\begin{align*}
 \mathcal D(B)&=\sum_{j,I}\frac{a_{j,I}}{c_I}B^{t_j-w(I)},\\
 \ell_0(B)&=\sum_{i=0}^{m-1}B^{v_i}+\sum_{j=0}^{s-1}B^{t_j},\\
 e_0(B)&=2\sum_{j=0}^{K-1}B^j+\mathcal D(B),&
 V&=\ell_0(4)+e_0(4)4^L.
\end{align*}
Monomial uniqueness, disjoint row intervals, and the dummy inequality
give
\begin{equation}\label{eq:four-coefficient-identity}
 [B^{t_j}]\,\mathcal D(B)\mathcal C(B)^2=G_j.
\end{equation}
They also show that the coefficients of $\mathcal D$ lie in
\eqref{eq:four-coefficient-alphabet}. Consequently all digits of $e_0$
belong to $\{0,1,2,3\}$. Zero digits are permitted. Since $\mathcal D$
has no constant term, the unit digit of $e_0$ is exactly $2$; in
particular $e_0$ and $V$ are positive.

Put $c_*=m+s+1$, the number of code coordinates including the input.
Choose a power of two
\begin{equation}\label{eq:four-index-bound}
 H>\max\{2\,4^{2L+1},\ 4^{t_*+3}\,4c_*^2,\ 3L,\ 16\}.
\end{equation}
An admissible index is a triple $(V,H,L)$ obtained by this construction
from a finite representing system. All choices, including the circuit
size, are made independently of the queried input $x$. The supplied
positive parameters are exactly $(x,V,H,L)$. The variable number of
decoded coordinates introduces no new unknown in the final system.

\subsection{The fixed-four system and its preliminary bounds}

In the system of Theorem~\ref{thm:best102}, set $Z=4$ and replace the
former fixed exponent by the index parameter $L$. With the same charged
abbreviations and positive witnesses, the coding equations become
\begin{align*}
 B&=Hb^2,& b&=x+\beta,& \lambda(B-1)&=q^2-1,& \theta+4&=B,\\
 S_2&=\ell+eq=V+t\theta,& S_2+\alpha&=q^2,&
 \Omega&=4\lambda-2e>0,&\mathcal C&=x+g,\\
 S_3&=B\lambda q-\Omega\mathcal C^2,& \sigma&=S_3>0,&n&=q^8.
\end{align*}
The packing is
\begin{align*}
 S&=g+q^2(S_2+q^2S_3),\\
 T+1&=q^2(1+\theta\lambda)-(b-1)\ell+\theta\ell q^4,\\
 r&=S(n^2-n)+(T+1)(n^2-1).
\end{align*}
All common-witness Pell equations of Section~\ref{sec:common-witness}
are retained. In the second exponential congruence the expression
$a-(B-4)$ is now $a-\theta$.

Before decoding, $S_2<q^2$ gives $e<q$ and $\ell<q^2$.
The geometric equation and $H>16$ give
$B\le q^2$, $q>4b$, and $3L\le B\le n$.
Since $\Omega\ge1$ and $S_3>0$,
\[
 \mathcal C^2<B\lambda q<3q^3<q^4.
\]
The first block borrows fewer than $b$ units from the middle block;
the latter remains positive because $\theta=B-4>b$.
The width-$(2,2,4)$ argument of Section~\ref{sec:borrow-mask} therefore
gives
\[
 0<S<n,\qquad 0<T+1\le n,\qquad n\le r\le2n^3-2n^2.
\]
These are the hypotheses of the common-witness Pell argument. Its
exponent comparisons use $L\ge2$ and $3L\le B\le n\le r$, which
hold for every admissible index here. It follows that $b,B,q$ are
powers of two, $q=B^L$, and the required central-binomial divisibility
holds. No bound on the number of compiled rows is used in that argument.

\subsection{Canonical codes, zero digits, and the target tests}

Now $\lambda>q$, so $\Omega>2\lambda$. Positivity of $S_3$ gives
$\mathcal C^2<Bq/2<q^2$, hence $g<\mathcal C<q$.
Let $d=\lfloor(b-1)\ell/q^2\rfloor<b$. Subtracting $d$ from
$(B-4)\lambda$ changes only its unit digit. The middle binary mask
therefore bounds all nonunit digits of $S_2$ by three and its unit
digit by $3+d$. Thus
\[
 \ell\le S_2<\frac{4q^2}{B-1}+b.
\]
Using $B=Hb^2$ and $H>16$ shows $(b-1)\ell<q^2$, so $d=0$.
The canonical evaluation lemma applies to
$\ell_0(B)+e_0(B)B^L$, whose digits are below four and whose degree
is below $2L$. The first term of \eqref{eq:four-index-bound} supplies
the required evaluation bound. It yields
\[
 S_2=\ell_0(B)+e_0(B)q,\qquad
 \ell=\ell_0(B)+a_{\mathrm{al}}q,\qquad e=e_0(B)-a_{\mathrm{al}},
 \qquad0\le a_{\mathrm{al}}<e_0(B).
\]
The quantity $a_{\mathrm{al}}$ is the temporary quotient alias.

The first mask, with $g<q$, decodes the low positions specified by
$\ell_0$. Every coordinate is below $b$, and the unit digit of
$\mathcal C$ is $x$. Put
\[
 A_{\mathcal C}=\mathcal C(1)^2,
 \qquad J_{\mathcal C}=4A_{\mathcal C}.
\]
There are $c_*$ coordinates, so
$\mathcal C(1)<c_*b$ and $J_{\mathcal C}<B/4$.
Also $e,\mathcal C<B^K$ and $L>3K+2$ give
$2e\mathcal C^2<q$.

Lemma~\ref{lem:borrow-high-window} applies with $Z=4$.
Its coefficient plateau is $J_{\mathcal C}$: the digit of $S_3$
at $L$ is $B-J_{\mathcal C}-1$ or $B-J_{\mathcal C}$, and the
digits at $L+1,\ldots,L+K$ are $B-J_{\mathcal C}$.
The high third-mask part is $X=(B-4)a_{\mathrm{al}}$, of degree at most $K$.
No overlap bounds every digit of $X$ by $J_{\mathcal C}$.
For its digit polynomial $F_X$, evaluation at four gives
\[
 B-4\mid F_X(4),\qquad
 0\le F_X(4)\le
 J_{\mathcal C}\frac{4^{K+1}-1}{3}<\frac B{12}<B-4.
\]
The second bound on $H$ proves the strict inequality. Hence $X=0$
and $a_{\mathrm{al}}=0$: both codes are canonical. Zero digits of $e_0$ have caused
no change, since the evaluation lemma requires digits in $[0,4)$ and
the high-window argument requires $e_0>0$, not positivity of each digit.

Below $K$, the raw coefficients of $S_3$ are those of
$2\mathcal D\mathcal C^2$. The non-strict coefficient bound
$|\mathcal D_j|\le2$ is sufficient because
\[
 \bigl|[B^j]\,2\mathcal D\mathcal C^2\bigr|
 \le4\mathcal C(1)^2=J_{\mathcal C}<B/4.
\]
All incoming carries in this range are zero or minus one.
The mask $B-4$ permits just the digits $0,1,2,3$; a negative raw
value plus its incoming carry normalizes to a digit greater than three.
The special target has raw coefficient $2\delta^2$, forcing
$\delta\in\{0,1\}$. If $\delta=0$, the guard has raw coefficient
$-4x^2$ and fails. Thus $\delta=1$.
Each ordinary target now has raw coefficient $4F_j+2$. With either
possible incoming carry, it is permitted exactly when $F_j=0$.
The original system therefore holds at the actual queried input $x$.

Conversely, extend an original solution by its gate values, take
$\delta=1$ and $u=x^2$, and choose a power-of-two $b$ larger than all
coordinates. Use the canonical codes. The low coordinate tests vanish
by support separation, and every ordinary or special target has raw
coefficient $2$, hence normalized digit $1$ or $2$.
The inequalities $e_0,\ell_0<q$ give $S_2<q^2$ and
$\Omega>2\lambda$. Since $2K<L$, one has
$4\mathcal C^2<Bq$, which gives $S_3>0$.
All required positive slack witnesses therefore exist.
The positive Pell witnesses from Section~\ref{sec:common-witness}
complete the final system. This proves both directions.

\begin{theorem}\label{thm:best101}
For every recursively enumerable subset of the positive integers,
there is an admissible fixed index $(V,H,L)$ such that its membership
relation is represented by a system with 34 positive unknowns and
22 equations admitting a straight-line certificate of 101 operations:
56 multiplications and 45 additions. The supplied parameters are
$(x,V,H,L)$, with three fixed index parameters. Fixed numerals are free.
\end{theorem}
\begin{proof}
The preceding construction proves the representation and both
positive-domain implications. In the 102-operation certificate,
replace $Z$ by $4$, replace the old exponent numeral by $L$, and
replace all uses of the separately computed $B-4$ by $\theta$.
The already present equation $\theta+4=B$ justifies that substitution.
It removes one subtraction, with no added arithmetic, unknown, or
equation. The exact primitive and polynomial-residual check is recorded
in \texttt{round23\_1980\_fixed\_four\_certificate.py} and its JSON
receipt. Thus the count is $56+45=101$; no optimality is asserted.
\end{proof}

% A doubled-index auxiliary Pell norm and a Pell-valued fixed index.
\section{A doubled-index congruence in the Pell block}
\label{sec:doubled-index}

Two changes allow the main norm coefficient to reuse the first
exponential modulus. The auxiliary signed norm can compare doubled
Pell indices, and the final index equation can use a fixed Pell value
in place of the index itself. Together they remove both affine
parameters $A-1$ and $A+1$ from the arithmetic certificate, reducing
Theorem~\ref{thm:best101} by one addition.

Retain the fixed-four encoding of Section~\ref{sec:fixed-four} and the
common exponential witness of Section~\ref{sec:common-witness}. Write
\[
 A=a+4,\qquad D=A^2-1,\qquad J=2r+1,\qquad H_0=J+jc,
\]
where the supplied positive unknown $a$ still satisfies
\[
 a=Y_P(U+1),\qquad U=wn^2,\qquad Y_P=sn^2.
\]
The first exponential equation and its modulus remain
\begin{equation}\label{eq:doubled-first-exponent}
 M=8a+15=8A-17,\qquad d=U+ac+\gamma M.
\end{equation}

Replace the signed auxiliary norm by
\begin{equation}\label{eq:doubled-auxiliary}
 \begin{gathered}
 v=of-d^2,\qquad K=D(f^2-1),\qquad G=2K-1,\\
 v(v+1)=K(K-1)H_0^2.
 \end{gathered}
\end{equation}
The quantities $v,K,G$ are abbreviations; $v$ is a computed integer
and need not be positive. No unknown is added. If $L$ is the underlying
power-of-two exponent used to construct the coefficient index, replace
the supplied component $L$ by
\[
 T_L=\psi_4(L)
\]
and replace the final index equation by
\begin{equation}\label{eq:doubled-fixed-index}
 \kappa=T_L+\Delta a.
\end{equation}
Thus the four supplied inputs are $(x,V,H,T_L)$. The admissible index
is constructed by choosing the same finite circuit layout and exponent
$L$ as before, constructing $V,H$ from them, and supplying $\psi_4(L)$.
These choices are independent of $x$. Every other source equation is
unchanged.

\subsection{Recovering the main index from doubled coordinates}

\begin{lemma}\label{lem:doubled-main-index}
Let $A>1$, let $J>1$ be odd, and suppose $c>J$. Over positive integer
unknowns, the two norms
\[
 d^2-(A^2-1)c^2=1,
 \qquad f^2-(A^2-1)(ic^2)^2=1
\]
together with \eqref{eq:doubled-auxiliary} imply
$c=\psi_A(J)$ and $d=\chi_A(J)$. Conversely, when those two identities
hold, positive $f,i,j,o$ satisfying the auxiliary equations exist.
\end{lemma}
\begin{proof}
Pell classification gives positive indices $p,m$ such that
\[
 d=\chi_A(p),\quad c=\psi_A(p),\qquad
 f=\chi_A(m),\quad ic^2=\psi_A(m).
\]
The divisibility fact used in Proposition~\ref{prop:signed-pell-new}
gives $c\mid m$ from $c^2\mid\psi_A(m)$. The case $p=1$ is excluded
by $c>J>1$. For $p=2$ one has $\psi_A(2)=2A\ge4$; for $p\ge3$,
\[
 \psi_A(p)\ge(2A-1)^{p-1}\ge3^{p-1}\ge2p.
\]
Consequently
\begin{equation}\label{eq:doubled-comparison-bound}
 0<2p\le c\le m.
\end{equation}

Here $K>1$ and $G>1$. The new norm is exactly
\[
 (2v+1)^2-(G^2-1)H_0^2=1.
\]
Its right-hand product $K(K-1)H_0^2$ is positive, so $v\ge1$ or
$v\le-2$. Taking $I=|2v+1|$, another application of Pell
classification gives
\[
 I=\chi_G(t),\qquad H_0=\psi_G(t)
\]
for a positive index $t$. Reduction modulo $f$ gives
\[
 \begin{aligned}
 G&=2D(f^2-1)-1\equiv-\chi_A(2),\\
 2v+1&=2of-2d^2+1\equiv-\chi_A(2p).
 \end{aligned}
\]
The identities
$\chi_{-z}(t)=(-1)^t\chi_z(t)$ and
$\chi_{\chi_A(2)}(t)=\chi_A(2t)$ therefore imply
\[
 \chi_A(2t)\equiv\pm\chi_A(2p)\pmod{\chi_A(m)}.
\]
By Lemma~\ref{lem:signed-chi-new}, whose comparison hypothesis is
\eqref{eq:doubled-comparison-bound},
\[
 2t\equiv\pm2p\pmod{2m},\qquad
 t\equiv\pm p\pmod m,\qquad t\equiv\pm p\pmod c.
\]
The middle implication divides an integer equality by two; it does
not assume that two is invertible modulo $m$.

Since $c\mid f^2-1$, one has $G\equiv-1\pmod c$. The recurrence for
$\psi$ then gives
\[
 H_0\equiv\psi_{-1}(t)=(-1)^{t-1}t\pmod c.
\]
Together with $H_0=J+jc$, this proves $J\equiv\pm p\pmod c$.
Both $J$ and $p$ lie in $(0,c)$, so $J=p$ or $J+p=c$. The latter
contradicts $\psi_A(p)\equiv p\pmod2$ and the oddness of $J$.
Thus $p=J$. No oddness condition on $f$ was used in this direction.

For necessity, start with $c=\psi_A(J)$ and $d=\chi_A(J)$. Choose
\[
 m=2cJ,\qquad f=\chi_A(m),\qquad i=\frac{\psi_A(m)}{c^2}.
\]
The quotient is a positive integer: in
$(d+c\sqrt D)^{2c}$ the term with one square root contains
$(2c)c=2c^2$, and every remaining square-root term contains at least
$c^3$. Moreover $m$ is even, so
$f=2\chi_A(m/2)^2-1$ is odd. Put
\[
 K=D(f^2-1),\quad G=2K-1,\quad
 I=\chi_G(J),\quad H_0=\psi_G(J).
\]
The parameter $G$ is odd and exceeds one, so $I$ is odd and greater
than one. Because $J$ is odd,
\[
 I\equiv-\chi_A(2J)=1-2d^2\pmod f.
\]
Thus $I+2d^2-1$ is divisible by both $2$ and the odd number $f$.
The choices
\[
 v=\frac{I-1}{2}>0,\qquad
 o=\frac{I+2d^2-1}{2f}>0
\]
are integers and satisfy $v=of-d^2$. Also $G\equiv-1\pmod c$ and
$J$ is odd, whence $H_0\equiv J\pmod c$. Strict Pell growth makes
$j=(H_0-J)/c$ a positive integer. The Pell norm for $I,H_0$ gives
\eqref{eq:doubled-auxiliary}.
\end{proof}

The required hypotheses $A>1$ and $1<J<c$ are established by the
unchanged preliminary packing and interval arguments. Thus this lemma
is available at the same point as the preceding signed-block lemma,
before either exponential relation or binary decoding is used.

\subsection{Encoding the fixed index by a Pell value}

\begin{lemma}\label{lem:doubled-fixed-index}
In the new system, \eqref{eq:doubled-fixed-index} and the unchanged
second norm and gap force
$\kappa=\psi_A(L)$ and $\mu=\chi_A(L)$.
Conversely, those coordinates give a positive integer $\Delta$
satisfying \eqref{eq:doubled-fixed-index}.
\end{lemma}
\begin{proof}
After Lemma~\ref{lem:doubled-main-index}, the common-witness argument
is unchanged through the first exponential equation. The ratio lower
bound first gives $Y_P\ge U^r$ and
$a=Y_P(U+1)>U^{r+1}$. Equation~\eqref{eq:doubled-first-exponent}
then gives $U=4^J$. Since $r+1>3$, the modulus in the new index
equation satisfies the stronger bound
\begin{equation}\label{eq:doubled-modulus-growth}
 a>U^{r+1}>4^{3J}.
\end{equation}
This is a bound for $a=A-4$, not merely for $A$.

The unchanged second norm and $c=\kappa+\phi$ give
\[
 \kappa=\psi_A(t_0),\qquad \mu=\chi_A(t_0),\qquad 0<t_0<J.
\]
Since $A\equiv4\pmod a$, the integer polynomial recurrence for
$\psi$ and \eqref{eq:doubled-fixed-index} yield
\[
 \psi_4(t_0)\equiv\psi_4(L)\pmod a.
\]
The underlying encoding still satisfies $3L\le B\le n\le r$,
so $L<J$. For $1\le s<J$,
\[
 0<\psi_4(s)\le8^{s-1}<4^{3J}<a.
\]
The two residues therefore agree as integers. Strict growth of
$\psi_4$ gives $t_0=L$.

Conversely, for the intended coordinates the quotient
\[
 \Delta=\frac{\psi_A(L)-\psi_4(L)}a
\]
is integral by the same recurrence and strictly positive because
$A>4$ and $L\ge2$. Strict growth in the base follows, for example,
from the positive binomial expansion of the Pell solution.
\end{proof}

Once this index is recovered, the second exponent test proves
$q=B^L$ with the same bounds as before. The common-witness proof has
already supplied the power-of-two radix and the requisite binomial
divisibility. All fixed-four coefficient decoding is consequently
unchanged. In necessity, the new fixed component $T_L$ is chosen with
$V,H$ before querying $x$; the remaining Pell witnesses are chosen by
the common-witness construction and the two lemmas above.

\subsection{The complete count}

The earlier certificate computed the norm coefficient by
\[
 a_-=a+3,\qquad a_+=a_-+2,\qquad D=a_-a_+.
\]
Neither affine value is required after the two changes. Since
$M=8a+15$ is already computed for the first exponent test, use instead
\[
 a_2=a^2,\qquad D=a_2+M.
\]
This replaces three instructions by two. The auxiliary block has the
same length as before: use the existing $d^2$ in $v=of-d^2$;
replace the former root square and final addition of one by $v+1$
and $v(v+1)$; and use the three coefficient instructions
$K=D\,AE$, $K_-=K-1$, and $KK_-$, where the existing E16 register
is $AE=D(ic^2)^2$. The final index equation still takes one product
and one addition.

For exact comparison with the source equations, let
\[
 F_{16}=f^2-AE-1,\qquad
 K_s=D(f^2-1),\qquad K_c=D\,AE.
\]
The changed norm has the triangular identity
\[
 F_{17}^{\mathrm{calc}}
 =F_{17}^{\mathrm{source}}
  +D F_{16}(K_s+K_c-1)H_0^2.
\]
The executable certificate checks this identity, the inherited E7 and
E18 corrections from $\theta=B-4$, every other residual, and every
primitive instruction. The fixed-index and positive-domain conclusions
are supplied by the proofs above, independently of these symbolic
checks.

\begin{theorem}\label{thm:best100-pell}
For every represented set, an admissible fixed index $(V,H,T_L)$
can be chosen such that the system described above represents it over
34 positive integer unknowns with 22 equality tests. Its straight-line
certificate has
\[
 \boxed{100=56\text{ multiplications}+44\text{ additions}}.
\]
Fixed numerals and supplied index parameters are free. Relative to
Theorem~\ref{thm:best101}, the saving is one addition.
\end{theorem}

The full source equations, primitive schedule, domains and residual
identities are recorded by
\texttt{round24\_1980\_doubled\_pell\_certificate.py} and its JSON receipt.

% An independent 100-operation encoding, with all Pell equations retained.
\section{Halving the coefficient and sharing the geometric factor}
\label{sec:half-center}

The fixed-four construction admits another saving, independent of the
Pell change in the preceding section. Starting from
Theorem~\ref{thm:best101}, replace its positive coefficient and third
block by
\begin{equation}\label{eq:half-center-block}
 \Omega=2\lambda-e>0,\qquad
 S_3=\theta\lambda q-\Omega\mathcal C^2=\sigma>0,
 \qquad \mathcal C=x+g,\quad\theta=B-4.
\end{equation}
The factor $\theta\lambda$ can now be shared with the geometric
equation. To accommodate the smaller coefficient, rebuild the fixed
index by pairing ordinary equation tests and placing the special
target first. No packed bound or Pell equation is removed.

\subsection{Paired rows and a first special target}

Use the finite primitive circuit construction of
Section~\ref{sec:fixed-four}, with distinct copied operands and fresh
outputs, all separate from the homogenizing coordinate $\delta$.
For every multiplication, addition, copy, final equality, or zero
residual
\[
 XY-W\delta,\qquad (X+Y-W)\delta,\qquad
 (X-Y)\delta,\qquad X\delta,
\]
include both $F$ and $-F$ among the ordinary rows. Omit trivial
identical equalities. Replace each unit gate $X=1$ by two unpaired
rows
\begin{equation}\label{eq:half-center-unit}
 F_{X,1}=X\delta-\delta^2,\qquad
 F_{X,2}=X\delta-X^2.
\end{equation}
Include the unpaired guard $F_{\mathrm g}=u\delta-x^2$, with $u$
fresh. Every ordinary row has target
\begin{equation}\label{eq:half-center-target}
 G_j=2F_j+\delta^2.
\end{equation}
The special target $G_0=\delta^2$ is placed \emph{first}, before all
ordinary targets.

After division by the monomial coefficient in $\mathcal C(T)^2$,
the coefficient alphabet is still $\{-2,-1,0,1\}$. Paired primitive
rows have cross-term quotients $\pm1$ and the added square
$\delta^2$ has quotient $1$. The two unit targets are
\[
 2X\delta-\delta^2,\qquad
 2X\delta-2X^2+\delta^2,
\]
and the guard has the latter coefficient pattern. The special
target has its single coefficient $1$. Thus the pairing does not
enlarge the digit alphabet. At a genuine circuit solution with
$\delta=1$, every ordinary $F_j$ is zero, including both unit rows;
choose $u=x^2$ for the guard.

\subsection{The rebuilt fixed index and the absence of an initial carry}

Apply the dynamic support construction of Section~\ref{sec:fixed-four}
to this finite list. With $m$ positive-weight coordinates and $s$
targets, write
\begin{align*}
 v_i&=3^i\quad(0\le i<m),& M&=3^{m-1},& d_0&=4M+1,\\
 t_0&=(s+1)d_0+2M,& t_j&=t_0+jd_0,& t_*&=2sd_0+2M,\\
 K&=t_*+1,& L&>3K+2,
\end{align*}
where $L$ is a power of two, $x$ has weight zero, and $\delta$ has
weight $v_0=1$. The special target is at $t_0$. Include the usual
dummy coordinate at each $t_j$. Monomial uniqueness, separation of
the intervals $[t_j-2M,t_j]$, and $2t_0-2M>t_*$ give
\begin{equation}\label{eq:half-center-coefficient}
 [T^{t_j}]\,\mathcal D(T)\mathcal C(T)^2=G_j.
\end{equation}
In particular, no term involving a dummy contributes to any target.

Define the fixed polynomials and parameters by
\begin{align*}
 \ell_0(T)&=\sum_{i=0}^{m-1}T^{v_i}+\sum_{j=0}^{s-1}T^{t_j},\\
 e_0(T)&=2\sum_{j=0}^{K-1}T^j+\mathcal D(T),&
 V&=\ell_0(4)+e_0(4)4^L.
\end{align*}
The coefficients of $e_0$ are in $\{0,1,2,3\}$, with unit digit
$2$. Put $c_*=m+s+1$, counting all code coordinates including the
input, and choose a power of two
\begin{equation}\label{eq:half-center-index-bound}
 H>\max\{2\,4^{2L+1},\ 4^{t_*+3}\,4c_*^2,\ 3L,\ 16\}.
\end{equation}
The resulting $(V,H,L)$ is an admissible index. Its construction is
independent of the queried input and introduces no additional
unknown in the final system.

The first target has a stronger support property than the others.
Its sole coefficient in $\mathcal D$ is $+1$ at $t_0-2$.
Every later row has its entire coefficient support above $t_0$,
because
\[
 t_1-2M=t_0+d_0-2M>t_0.
\]
Consequently every raw coefficient of
$\mathcal D(T)\mathcal C(T)^2$ through position $t_0$ is
nonnegative. After coordinate decoding these coefficients are below
$B$, as proved below. Hence the incoming carry at the special target
is exactly zero. This argument assumes neither $\delta=1$ nor any
ordinary equation.

\subsection{Bounds before the Pell argument}

Retain $B=Hb^2$, $b=x+\beta$, $n=q^8$, and
\begin{equation}\label{eq:half-center-coding}
 \lambda(B-1)=q^2-1,\quad \theta+4=B,\quad
 S_2=\ell+eq=V+t\theta,\quad S_2+\alpha=q^2.
\end{equation}
All supplied unknowns, including $\Omega,\sigma,\alpha$, remain
positive integers. Keep the packing
\begin{align}
 S&=g+q^2(S_2+q^2S_3),\notag\\
 T+1&=q^2(1+\theta\lambda)-(b-1)\ell+\theta\ell q^4,\notag\\
 r&=S(n^2-n)+(T+1)(n^2-1).\label{eq:half-center-packing}
\end{align}

The packed bound gives $e<q$ and $\ell<q^2$. Geometry and
\eqref{eq:half-center-index-bound} give
\[
 B\le q^2,\qquad q>4b\ge8,\qquad 3L\le B\le n.
\]
Since $\Omega\ge1$ and $S_3>0$, the new offset gives
\begin{equation}\label{eq:half-center-preliminary-C}
 \mathcal C^2<\theta\lambda q<B\lambda q<3q^3<q^4.
\end{equation}
Thus $g<\mathcal C<q^2$ and $0<S_3<q^4$. The mask $T+1$ is
unchanged; its first block borrows fewer than $b$ units from the
middle block, which remains positive since $\theta=B-4>b$.
The same width-$(2,2,4)$ bounds therefore give
\begin{equation}\label{eq:half-center-packing-bounds}
 0<S<n,\qquad 0<T+1\le n,\qquad
 n\le r\le2n^3-2n^2<2n^3.
\end{equation}

These bounds have been established before any Pell equation is
used. In particular $n,r\ge64$, $U=w n^2\ge n^2$ and
$Y_P=s n^2\ge n^2$. The main parameter satisfies
$a=Y_P(U+1)>n^4>2r+1$, and the first index modulus satisfies
$UY_P\ge n^4>r+1$. Moreover $L\ge2$ and
$3L\le B\le n\le r$. Thus every preliminary hypothesis of the
common-witness Pell proof in Section~\ref{sec:common-witness} is
available. All its equations are retained. That proof yields powers
of two $b,B,q$, the exact equality $q=B^L$, and
\begin{equation}\label{eq:half-center-binomial}
 n^2\mid\binom{2r}{r}.
\end{equation}
No internal target equation has entered this argument.

Now $\lambda>q$ and $e<q$, so $\Omega=2\lambda-e>\lambda$.
Positivity in \eqref{eq:half-center-block} improves the code bound to
\begin{equation}\label{eq:half-center-post-C}
 \mathcal C^2<\theta q<Bq<q^2,
 \qquad g<\mathcal C<q.
\end{equation}
The original factor one-half in the bound for $\mathcal C^2$ is no
longer asserted; the strict inequality $\mathcal C<q$ is what the
subsequent decoding needs.

\subsection{Canonical decoding and the revised high window}

By \eqref{eq:half-center-packing-bounds} and
\eqref{eq:half-center-binomial}, the packed binary no-carry tests
hold. Let $d=\lfloor(b-1)\ell/q^2\rfloor<b$. Subtracting $d$
from $(B-4)\lambda$ changes only its unit digit. The middle mask
therefore bounds all nonunit digits of $S_2$ by three and its unit
digit by $3+d$. It follows that
\[
 \ell\le S_2<\frac{4q^2}{B-1}+b.
\]
Together with $B=Hb^2$ and $H>16$, this forces
$(b-1)\ell<q^2$, hence $d=0$. The base-four evaluation lemma now
applies to the fixed polynomial defining $V$: its degree is below
$2L$, all digits are below four, and the first term of
\eqref{eq:half-center-index-bound} supplies the evaluation bound.
Consequently
\begin{equation}\label{eq:half-center-alias}
 \begin{gathered}
 S_2=\ell_0(B)+e_0(B)q,\\
 \ell=\ell_0(B)+a_{\mathrm{al}}q,\qquad
 e=e_0(B)-a_{\mathrm{al}},\qquad
 0\le a_{\mathrm{al}}<e_0(B).
 \end{gathered}
\end{equation}

The first mask and $g<q$ decode the low positions specified by
$\ell_0$. The unit digit of $g$ is zero; every coordinate is a
nonnegative integer below $b$, and the unit digit of $\mathcal C$
is $x$. Dummy digits occur only at row positions. Write
\[
 A_{\mathcal C}=\mathcal C(1)^2,
 \qquad J_{\mathcal C}=2A_{\mathcal C}.
\]
Then $\mathcal C(1)<c_*b$, and
\eqref{eq:half-center-index-bound} gives $J_{\mathcal C}<B/4$.
Also $e,\mathcal C<B^K$ and $L>3K+2$ give $e\mathcal C^2<q$.
Since $g>0$ and $x>0$, the decoded coefficient sum is at least two,
so
\begin{equation}\label{eq:half-center-coefficient-sum}
 A_{\mathcal C}\ge4.
\end{equation}
This holds before identifying $\delta$, even if the nonzero digit
of $g$ occupies a dummy position.

Write
\[
 2\lambda\mathcal C^2=qV_1+V_0,\qquad0\le V_0<q.
\]
Its polynomial coefficients are at most $J_{\mathcal C}<B$, so
they are its base-$B$ digits. The stable convolution interval
$L,\ldots,L+K$ has all digits equal to $J_{\mathcal C}$, because
$L>3K+2$. Exact division of \eqref{eq:half-center-block} by $q$
therefore gives
\[
 \left\lfloor\frac{S_3}{q}\right\rfloor
 =\theta\lambda-V_1+\epsilon,\qquad
 \epsilon=\left\lfloor
 \frac{e\mathcal C^2-V_0}{q}\right\rfloor\in\{-1,0\}.
\]
Every digit of $\theta\lambda$ through position $2L-1$ is $B-4$.
Subtracting $V_1$ creates no borrow in the window: its first digit
$B-4-J_{\mathcal C}+\epsilon$ is positive, as are all subsequent
digits. Thus
\begin{equation}\label{eq:half-center-high-digits}
 (S_3)_L=B-4-J_{\mathcal C}+\epsilon,\qquad
 (S_3)_{L+j}=B-4-J_{\mathcal C}\quad(1\le j\le K).
\end{equation}

The high part of the third mask is
$X=(B-4)a_{\mathrm{al}}<B^{K+1}$. Binary nonoverlap bounds each
digit of $X$ by
\[
 J_{\mathcal C}+4=2A_{\mathcal C}+4\le4A_{\mathcal C},
\]
using \eqref{eq:half-center-coefficient-sum}. If $F_X$ is its digit
polynomial, then $B-4\mid F_X(4)$, whereas
\[
 0\le F_X(4)\le
 4A_{\mathcal C}\frac{4^{K+1}-1}{3}<\frac B{12}<B-4.
\]
The strict bound follows from
\eqref{eq:half-center-index-bound}, without enlarging $H$.
Therefore $F_X(4)=0$, $X=0$, and $a_{\mathrm{al}}=0$.
Both supplied codes are canonical.

\subsection{The low target tests and positive necessity}

With the canonical codes, the raw coefficients of $S_3$ below $K$
are those of $\mathcal D(T)\mathcal C(T)^2$. The positive offset
starts at $L>K$, and the missing baseline tail starts at $K$.
Every raw coefficient has absolute value at most
$2A_{\mathcal C}<B/4$, so every incoming carry is zero or minus one.
At an ordinary target put $v=2F+\delta^2$. The mask $B-4$ permits
exactly the digits $0,1,2,3$ and is equivalent to
\begin{equation}\label{eq:half-center-allowed}
 0\le v+\epsilon\le3,\qquad\epsilon\in\{-1,0\}.
\end{equation}
A negative value would normalize to a digit greater than three.

At the first special target there is no incoming carry, by its
support property. The target is $\delta^2$, so
$0\le\delta^2\le3$, forcing $\delta\in\{0,1\}$. If $\delta=0$,
the guard has $v=-2x^2<0$ and cannot pass. Thus $\delta=1$.
For either possible carry, \eqref{eq:half-center-allowed} now says
\[
 0\le2F+1+\epsilon\le3
 \quad\Longleftrightarrow\quad F\in\{0,1\}.
\]
Paired rows $F$ and $-F$ force $F=0$. The unit rows give
\[
 X-1\in\{0,1\},\qquad X-X^2\in\{0,1\}.
\]
The first restricts $X$ to $\{1,2\}$, and the second restricts it
to $\{0,1\}$, so $X=1$. The guard need not vanish in sufficiency:
it excludes $\delta=0$, and its coordinate $u$ is absent from the
original circuit. Every original gate and equality has now been
recovered exactly at the input $x$. The low coordinate tests are
automatic because $\mathcal D$ starts above their positions.

Conversely, take a solution of the represented system, extend it by
its circuit values, and choose $\delta=1$, $u=x^2$, and zero dummy
coordinates. Choose a power-of-two $b$ larger than every coordinate
and use the canonical codes and $q=B^L$. Every ordinary $F$ is
zero. The special target has raw coefficient one and carry zero;
each ordinary target has raw coefficient one and actual digit zero
or one. These digits are allowed, and the low variable tests vanish.
The first two masks hold as in Section~\ref{sec:fixed-four}.

All positive slack witnesses exist. In particular, the unit digit
of $e_0$ is two, the positive $\delta$ digit makes $g>0$, and
$S_2<q^2$. The packed-congruence quotient is positive because
$P(B)>P(4)$ for the nonconstant polynomial
$P(T)=\ell_0(T)+e_0(T)T^L$ with nonnegative coefficients.
Also $\Omega=2\lambda-e>\lambda>0$. Since
$\mathcal C<B^K$ and $L>3K+2$,
\[
 2\mathcal C^2<q\le\theta q,
 \qquad
 S_3=\lambda(\theta q-2\mathcal C^2)+e\mathcal C^2>0.
\]
The three masks give \eqref{eq:half-center-binomial}, and their
packing satisfies \eqref{eq:half-center-packing-bounds}. The complete
positive Pell witness construction of Section~\ref{sec:common-witness}
therefore applies to the newly calculated $r$. This proves necessity
for the rebuilt fixed index.

\subsection{Seven shared instructions in place of eight}

The predecessor uses eight instructions for the relevant quantities:
\[
 \begin{gathered}
 L_2=\lambda+q^2,\quad \Lambda=B\lambda,\quad R_2=\Lambda+1,\\
 z_\lambda=4\lambda,\quad e_2=2e,\quad
 t_2=z_\lambda-e_2,\\
 T_{\mathrm{coef}}=L_2-z_\lambda,\qquad P_2=\Lambda q.
 \end{gathered}
\]
Use instead the seven instructions
\begin{equation}\label{eq:half-center-seven}
 \begin{gathered}
 U_\lambda=\theta\lambda,\qquad \lambda_2=2\lambda,\qquad
 T_{\mathrm{coef}}=U_\lambda+1,\\
 \lambda_3=\lambda_2+\lambda,\qquad
 R_2=T_{\mathrm{coef}}+\lambda_3,\\
 t_2=\lambda_2-e,\qquad P_2=U_\lambda q.
 \end{gathered}
\end{equation}
Test $q^2=R_2$ and $t_2=\Omega$. In conjunction with
$\theta+4=B$, the first test is exactly
$\lambda(B-1)=q^2-1$. All other arithmetic instructions remain,
including the packed upper bound and $S_3=P_2-t_2\mathcal C^2$.

\begin{theorem}\label{thm:best100-encoding}
Every recursively enumerable subset of the positive integers has an
admissible fixed index $(V,H,L)$ for a system with 34 positive unknowns
and 22 equations admitting a straight-line certificate of 100
operations: 55 multiplications and 45 additions. The supplied
parameters are $(x,V,H,L)$; fixed numerals are free.
\end{theorem}
\begin{proof}
The construction above proves both positive-domain implications.
The replacement \eqref{eq:half-center-seven} has three multiplications
and four additions, instead of four multiplications and four
additions, and every other instruction is retained. Hence the
101-operation certificate loses one multiplication.

The exact verification is
\texttt{round25\_1980\_half\_center\_certificate.py} and its JSON
receipt. It checks the full 100-instruction primitive list and all
22 polynomial residuals. If $F_3$ is the original geometric residual
and $F_4=\theta+4-B$, the new geometric test has residual
$F_3-\lambda F_4$. The new packing and positive-$S_3$ source equations
use $\theta$ explicitly; they need no geometric correction.
The existing signed Pell and second-exponent corrections are also
verified. Thus the count concerns an explicit certificate for the
proved system; no optimality is claimed.
\end{proof}

% Compose the independently proved coding and Pell reductions.
\section{A certificate of 99 operations}\label{sec:composed-99}

The reductions in Sections~\ref{sec:half-center} and
\ref{sec:doubled-index} can be applied together. Their arithmetic
changes occupy disjoint parts of the certificate. The composition
also preserves the order of the mathematical proof: the new packing
first supplies the bounds for the new Pell block; that block then
recovers the exponent and binary radix needed by the new digit tests.

\begin{theorem}\label{thm:best99}
For every recursively enumerable subset of the positive integers,
there is an admissible fixed index $(V,H,T_L)$ such that its membership
relation is represented by a system with 34 positive unknowns and
22 equations admitting a straight-line certificate of 99 operations:
55 multiplications and 44 additions. Here $T_L=\psi_4(L)$ for the
underlying exponent of the index construction. The supplied parameters
are $(x,V,H,T_L)$, with three fixed index parameters, and fixed numerals
are free.
\end{theorem}

\begin{proof}
Construct the coefficient index using the entire paired-row
construction of Section~\ref{sec:half-center}. In particular, the
special target $\delta^2$ comes first; multiplication, addition,
copy, equality, and zero rows are paired with their negatives; each
unit gate uses the two unpaired rows
$X\delta-\delta^2$ and $X\delta-X^2$; and the guard is
$u\delta-x^2$. Every ordinary target is $2F+\delta^2$.

If there are $m$ positive-weight coordinates and $s$ targets, use
\[
\begin{gathered}
 v_i=3^i\quad(0\le i<m),\qquad M=3^{m-1},\qquad d_0=4M+1,\\
 t_j=(s+1)d_0+2M+jd_0\quad(0\le j<s),\\
 t_*=2sd_0+2M,\qquad K_0=t_*+1,\qquad c_*=m+s+1.
\end{gathered}
\]
Choose a power of two $L>3K_0+2$, form the corresponding
$\ell_0,e_0$, and set
\[
 V=\ell_0(4)+e_0(4)4^L,\qquad
 H>\max\{2\,4^{2L+1},\ 4^{t_*+3}\,4c_*^2,\ 3L,\ 16\},
\]
with $H$ a power of two. Finally supply $T_L=\psi_4(L)$ instead
of $L$. Thus $V,H$ come from the new paired layout throughout;
they are not reused from an earlier unpaired index. These choices
are independent of the queried input $x$.

Keep the coding equations and packing
\eqref{eq:half-center-block}--\eqref{eq:half-center-packing}, with
\[
 \Omega=2\lambda-e>0,\qquad
 \sigma=\theta\lambda q-\Omega\mathcal C^2>0,
 \qquad \theta=B-4,
\]
and apply both Pell changes
\eqref{eq:doubled-auxiliary} and \eqref{eq:doubled-fixed-index}.
All other common-witness Pell equations remain unchanged.

Before either code is decoded, the positive bound $S_2+\alpha=q^2$
gives $e<q$ and $\ell<q^2$. The geometric equation gives
$B\le q^2$, and the index bound gives $q>4b\ge8$ and
$3L\le B\le n$. Positivity of $\Omega$ and $\sigma$ gives
\[
 \mathcal C^2<\theta\lambda q<3q^3<q^4.
\]
The preliminary borrow argument of Section~\ref{sec:half-center}
therefore proves
\begin{equation}\label{eq:composed99-preliminary}
 0<S<n,\qquad 0<T+1\le n,\qquad
 n\le r\le2n^3-2n^2<2n^3.
\end{equation}
No coefficient target has been tested at this stage.

Put $U=wn^2$, $Y_P=sn^2$, $a=Y_P(U+1)$, $A=a+4$, and $J=2r+1$.
Then $U,Y_P\ge n^2$, $UY_P\ge n^4>r+1$, and the first-index
congruence and interval give $c>J$. Also $A>n^4>J$.
Thus every hypothesis of Lemma~\ref{lem:doubled-main-index} holds:
$A>1$, $1<J<c$, and $J$ is odd. That lemma proves
\[
 c=\psi_A(J),\qquad d=\chi_A(J).
\]
In particular, its signed comparison does not depend on either
coefficient decoding or an already known value of $q$.

The unchanged first Pell index and lower ratio argument now give,
with $\xi=(U+1)^{2r}/U^r$,
\[
 k=\psi_{2UY_P^2+1}(r+1),\qquad
 \xi<c/k<Y_P+1,\qquad
 Y_P\ge U^r,\qquad a>U^{r+1}.
\]
The first exponential congruence yields $U=4^J$. Therefore $n$ and
$q$ are powers of two, and $a>4^{3J}$.
The positive gap and second norm give
$\kappa=\psi_A(t)$ for some $0<t<J$. Admissibility gives
$L<J$. Lemma~\ref{lem:doubled-fixed-index} applies to
$\kappa=T_L+\Delta a$: reduction modulo $a$ gives
$\psi_4(t)\equiv\psi_4(L)$, while both values lie strictly between
zero and $a$. It follows that $t=L$.

The remaining exponent comparison uses
\[
 B^{3L}\le B^B\le n^n\le U^r<A,
 \qquad q^3\le n^n\le U^r<A,
\]
so it yields $q=B^L$, and hence powers of two $B$ and $b$.
The unchanged upper ratio and binomial arguments give
\begin{equation}\label{eq:composed99-binomial}
 n^2\mid\binom{2r}{r}.
\end{equation}
These conclusions precede all internal circuit tests.

Now $\lambda>q$ and $e<q$, so $\Omega>\lambda$ and
$\mathcal C^2<\theta q<Bq<q^2$. Thus $g<\mathcal C<q$.
All post-Pell hypotheses of Section~\ref{sec:half-center} have been
restored. Its first two masks remove the preliminary borrow and
identify the packed coefficient polynomial, up to the quotient alias.
The first mask bounds every coordinate by $b$. Positivity of $x,g$
gives $\mathcal C(1)\ge2$ before $\delta$ is identified. Accordingly,
with $A_{\mathcal C}=\mathcal C(1)^2$, the modified high-window bound
$2A_{\mathcal C}+4\le4A_{\mathcal C}$ and the displayed bound on $H$
remove the alias. The two codes are canonical.

The first special target has incoming carry zero and forces
$\delta\in\{0,1\}$. The guard excludes zero. At $\delta=1$, each
ordinary row tests $F\in\{0,1\}$; its paired negative forces $F=0$,
and the two unit rows force the unit coordinate to equal one.
This recovers the represented system at the actual input $x$.

For the converse, take a solution of that system, its circuit values,
$\delta=1$, $u=x^2$, and zero dummy digits. Choose a power-of-two
$b$ larger than every coordinate and use the canonical codes.
Then $\beta,\theta,\lambda,\alpha$ and the packed-congruence quotient
are positive integers. The half-center necessity proof gives positive
$\Omega,\sigma$ and all three no-carry conditions. Define $r$ by
the packing; \eqref{eq:composed99-preliminary} and
\eqref{eq:composed99-binomial} follow.

Choose $J=2r+1$, $U=4^J$, and $w=U/n^2$. This quotient is a positive
integer since $n$ is a power of two and $n^2\le r^2<4^J$.
Take $Y_P=\lfloor\xi\rfloor$ and $s=Y_P/n^2$, integral by the
binomial congruence, and define $a=Y_P(U+1)$.
The common-witness construction supplies positive
$k,c,d,\tau,h,\eta,\zeta,\gamma,\rho$.
Choose
\[
 \kappa=\psi_A(L),\quad \mu=\chi_A(L),\quad
 \phi=c-\kappa,\quad
 \Delta=\frac{\psi_A(L)-\psi_4(L)}{a}.
\]
Here $L<J$ makes $\phi>0$, and $A=a+4>4$, $L\ge2$ make
$\Delta$ a positive integer.

Finally choose the auxiliary witnesses exactly as in
Lemma~\ref{lem:doubled-main-index}: set
\[
\begin{gathered}
 m_0=2cJ,\qquad f=\chi_A(m_0),\qquad
 i=\psi_A(m_0)/c^2,\\
 K=(A^2-1)(f^2-1),\qquad G=2K-1,\qquad
 I=\chi_G(J),\qquad H_0=\psi_G(J),\\
 o=\frac{I+2d^2-1}{2f},\qquad j=\frac{H_0-J}{c}.
\end{gathered}
\]
The proved divisibility, oddness of $f$, and Pell congruences make
all four auxiliary unknowns positive integers and establish the
doubled norm. This last choice changes no coding, interval, gap,
or exponential witness. Thus every required positive unknown exists.

For the arithmetic count, the half-center substitution saves one
multiplication from Theorem~\ref{thm:best101}; the doubled-index
substitution saves one addition. They change disjoint source rows and
disjoint instruction groups. The exact checker
\texttt{round26\_1980\_composed\_certificate.py} verifies that both
orders give the same complete schedule and source system, together
with all primitive instructions, 22 equality tests, and three
triangular residual corrections. The resulting count is
$(56-1)+(45-1)=55+44=99$.
\end{proof}

% A unit baseline and a direct guard preserve the fixed-four alphabet.
\section{A coefficient baseline of one}\label{sec:unit-center}

The coefficient baseline in Theorem~\ref{thm:best99} can be reduced
from two to one. Copies remove the negative square coefficient in
the unit test, while a direct guard excludes a zero homogenizing
coordinate before any copy relation is used. The resulting
coefficient is
\begin{equation}\label{eq:unit-center-block}
 \Omega=\lambda-e>0,\qquad
 S_3=\theta\lambda q-\Omega\mathcal C^2=\sigma>0,
 \qquad\mathcal C=x+g,\quad\theta=B-4.
\end{equation}
The offset, geometric equation, packed upper bound and all Pell
equations are retained. Computing $3\lambda$ directly now saves
one addition, giving 98 operations.

\subsection{Copied unit tests and a direct guard}

Keep the paired primitive rows of Section~\ref{sec:half-center}:
each multiplication, addition, copy, final equality and zero
residual $F$ occurs together with $-F$, with targets
$2F+\delta^2$. Operands are distinct copied coordinates, outputs
are fresh, and circuit coordinates are distinct from $\delta$.

For each unit gate $X=1$, introduce a fresh nonnegative coordinate
$X'$ and the paired copy rows $(X-X')\delta$ and $(X'-X)\delta$.
Replace its two unpaired unit rows by
\begin{equation}\label{eq:unit-center-unit}
 F_{X,1}=X\delta-\delta^2,\qquad
 F_{X,2}=X\delta-XX'.
\end{equation}
All four rows still use targets $2F+\delta^2$. The negative product
$XX'$ has distinct factors, so its divided target coefficient is
$-1$, in place of the $-2$ of the old square $X^2$.

Give the fresh guard coordinate $u$ the direct target
\begin{equation}\label{eq:unit-center-guard}
 G_{\mathrm g}=2u\delta-x^2+\delta^2.
\end{equation}
This target is not required to be $2F+\delta^2$ for an integer
polynomial $F$. It will exclude $\delta=0$ directly. Replacing
$x^2$ by a copied product would not do this: its homogenized copy
equation would impose no condition when $\delta=0$.
Keep the special target $G_0=\delta^2$ first.

Every target coefficient, divided by its monomial coefficient in
$\mathcal C(T)^2$, now belongs to
\begin{equation}\label{eq:unit-center-alphabet}
 \{-1,0,1\}.
\end{equation}
Paired primitive rows have cross-term quotients $\pm1$ and added
$\delta^2$ coefficient $1$. The first unit target has coefficients
$1,-1$ after division; the second has cross-term quotients $1,-1$
and square coefficient $1$. The direct guard has cross-term quotient
$1$, coefficient $-1$ on $x^2$, and coefficient $1$ on $\delta^2$.
Thus all these divisions are integral.

\subsection{The explicit admissible index}

Apply the support construction to the entire new finite list,
including all unit-copy coordinates and paired copy rows. Let $m$
be the number of positive-weight coordinates and $s$ the number of
targets. The input $x$ has weight zero and $\delta$ has weight one.
Set
\begin{align*}
 v_i&=3^i\quad(0\le i<m),& M&=3^{m-1},& d_0&=4M+1,\\
 t_j&=(s+1)d_0+2M+jd_0\quad(0\le j<s),\\
 t_*&=2sd_0+2M,& K_0&=t_*+1,&c_*&=m+s+1.
\end{align*}
The special target occupies $t_0$. Include a dummy coordinate at
every target position and choose a power of two $L>3K_0+2$.
The complementary coefficient polynomial $\mathcal D$ has
coefficients in \eqref{eq:unit-center-alphabet} and satisfies
\[
 [T^{t_j}]\,\mathcal D(T)\mathcal C(T)^2=G_j.
\]
The base-three monomial uniqueness and disjoint row intervals prove
this identity exactly; $2t_0-2M>t_*$ excludes every dummy contribution.
The support of $\mathcal D$ lies above every low coordinate and
below $K_0$. Through $t_0$ its only term is $+T^{t_0-2}$, since
every later row starts above $t_0$.

Define
\begin{align}
 \ell_0(T)&=\sum_{i=0}^{m-1}T^{v_i}+\sum_{j=0}^{s-1}T^{t_j},\notag\\
 e_0(T)&=\sum_{j=0}^{K_0-1}T^j+\mathcal D(T),&
 V&=\ell_0(4)+e_0(4)4^L.\label{eq:unit-center-index}
\end{align}
The digits of $e_0$ are in $\{0,1,2\}$ and its unit digit is one.
In particular, $e_0>0$, while the packed digit polynomial has all
digits below four and degree below $2L$. Choose a power of two
\begin{equation}\label{eq:unit-center-H}
 H>\max\{2\,4^{2L+1},\ 4^{t_*+3}\,4c_*^2,\ 3L,\ 16\}.
\end{equation}
Finally supply $T_L=\psi_4(L)$. The admissible index is
$(V,H,T_L)$; $L$ is the underlying exponent defining it, not an
additional system input. Both $V$ and $H$ are rebuilt from the new
layout and baseline. All choices depend only on the represented
system, independently of the queried positive input $x$.

\subsection{Establish the Pell hypotheses before decoding}

Besides \eqref{eq:unit-center-block}, retain
\begin{align*}
 B&=Hb^2,& b&=x+\beta,&n&=q^8,\\
 \lambda(B-1)&=q^2-1,&\theta+4&=B,\\
 S_2&=\ell+eq=V+t\theta,&S_2+\alpha&=q^2,\\
 S&=g+q^2(S_2+q^2S_3),\\
 T+1&=q^2(1+\theta\lambda)-(b-1)\ell+\theta\ell q^4,\\
 r&=S(n^2-n)+(T+1)(n^2-1).
\end{align*}
Every supplied unknown remains positive. The packed bound gives
$e<q$ and $\ell<q^2$, and geometry gives $B\le q^2$.
The bound on $H$ gives $q>4b\ge8$ and $3L\le B\le n$.
Since $\Omega\ge1$, positivity of the new block gives directly
\[
 \mathcal C^2<\theta\lambda q<B\lambda q<3q^3<q^4.
\]
Thus $g<\mathcal C<q^2$ and $0<S_3<q^4$. The first-block borrow
is fewer than $b$, while $\theta\lambda>b$, so the unchanged
width-$(2,2,4)$ estimates give
\begin{equation}\label{eq:unit-center-ranges}
 0<S<n,\quad0<T+1\le n,\quad
 n\le r\le2n^3-2n^2<2n^3.
\end{equation}
These are direct bounds for the new coefficient; no comparison
with an old value of $S_3$ is assumed.

Write $U=wn^2$, $Y_P=sn^2$, $a=Y_P(U+1)$, $A=a+4$ and $J=2r+1$.
Then $n,r\ge64$, $U,Y_P\ge n^2$, $UY_P\ge n^4>r+1$, and
$a>n^4>J$. The first index congruence and positive interval give
$c>J$. Thus $A>1$, $1<J<c$ and odd $J$ are available, as are
$L\ge2$ and $3L\le B\le n\le r$.

All Pell source equations of Theorem~\ref{thm:best99} are unchanged.
Their proof therefore applies in its original order. First the
doubled-index auxiliary norm identifies $c=\psi_A(J)$,
$d=\chi_A(J)$. The first norm and index give
$k=\psi_{2UY_P^2+1}(r+1)$. The lower ratio estimate gives
$Y_P\ge U^r$ and $a>U^{r+1}$, before the first exponent equation
yields $U=4^J$. In particular $a>4^{3J}$.
The positive gap and the equation $\kappa=T_L+\Delta a$ now
identify the second Pell index as $L$: its two residues
$\psi_4(t)$ and $\psi_4(L)$ lie strictly between zero and $a$.
The second exponent comparison uses
\[
 B^{3L}\le B^B\le n^n\le U^r<A,
 \qquad q^3\le n^n\le U^r<A,
\]
and yields $q=B^L$. The upper ratio and binomial rounding arguments
then give
\begin{equation}\label{eq:unit-center-decoded-powers}
 b,B,q\text{ are powers of two},\qquad
 q=B^L,\qquad n^2\mid\binom{2r}{r}.
\end{equation}
Every internal coefficient test still awaits decoding.

The weaker coefficient requires a revised post-Pell estimate. Since
$L\ge2$ and $B>16$,
\[
 \lambda\ge B^{2L-1}\ge Bq>2q,
 \qquad q=B^L\ge B^2>2B.
\]
Therefore $e<q$ implies $\Omega=\lambda-e>\lambda/2$, and
\eqref{eq:unit-center-block} yields
\begin{equation}\label{eq:unit-center-C-bound}
 \mathcal C^2<2\theta q<2Bq<q^2,
 \qquad g<\mathcal C<q.
\end{equation}

\subsection{Canonical codes and the smaller coefficient plateau}

The first two masks, congruence and bound are unchanged. Their
borrow-removal proof gives $d=\lfloor(b-1)\ell/q^2\rfloor=0$:
the middle mask first implies
$\ell\le S_2<4q^2/(B-1)+b$, which forces $(b-1)\ell<q^2$.
The base-four evaluation lemma applies to
\eqref{eq:unit-center-index}; it permits digits in $\{0,1,2\}$
and does not require the digit three. Consequently
\[
 S_2=\ell_0(B)+e_0(B)q,\qquad
 \ell=\ell_0(B)+a_{\mathrm{al}}q,\qquad
 e=e_0(B)-a_{\mathrm{al}},\qquad
 0\le a_{\mathrm{al}}<e_0(B).
\]
Using \eqref{eq:unit-center-C-bound}, the first mask decodes every
coordinate below $b$, with input $x$ at the unit position and dummy
digits only at row positions. Thus
\[
 e,\mathcal C<B^{K_0},\qquad e\mathcal C^2<q,
 \qquad A_{\mathcal C}=\mathcal C(1)^2<(c_*b)^2.
\]
Since $x,g>0$, one has $\mathcal C(1)\ge2$ and
$A_{\mathcal C}\ge4$ before recovering $\delta$.

Now the polynomial $\lambda\mathcal C^2$ has coefficients at most
$J_{\mathcal C}=A_{\mathcal C}<B/4$, and its digits from $L$
through $L+K_0$ all equal $J_{\mathcal C}$. Write
\[
 \lambda\mathcal C^2=qV_1+V_0,\qquad0\le V_0<q.
\]
Exact division gives
\[
 \left\lfloor\frac{S_3}{q}\right\rfloor
 =\theta\lambda-V_1+\epsilon,\qquad
 \epsilon=\left\lfloor
 \frac{e\mathcal C^2-V_0}{q}\right\rfloor\in\{-1,0\}.
\]
There is no borrow, since the digits of $\theta\lambda$ are $B-4$
throughout this interval. The first digit of $S_3$ in the window is
$B-4-J_{\mathcal C}+\epsilon$, and every subsequent one is
$B-4-J_{\mathcal C}$. Binary nonoverlap therefore bounds every
digit of $X=(B-4)a_{\mathrm{al}}$ by
\[
 J_{\mathcal C}+4=A_{\mathcal C}+4\le4A_{\mathcal C}.
\]
For its digit polynomial $F_X$,
\[
 B-4\mid F_X(4),\qquad
 0\le F_X(4)\le
 4A_{\mathcal C}\frac{4^{K_0+1}-1}{3}<\frac B{12}<B-4.
\]
The unchanged bound \eqref{eq:unit-center-H} supplies the strict
inequality. Hence $X=0$ and $a_{\mathrm{al}}=0$, so both codes
are canonical before any circuit row is tested.

\subsection{Recover the special coordinate before the copies}

Below $K_0$, the raw coefficients of $S_3$ are now those of
$\mathcal D(T)\mathcal C(T)^2$. Each has absolute value at most
$A_{\mathcal C}<B/4$, so incoming carries are zero or minus one.
At a target raw value $v$, the mask $B-4$ is equivalent to
\begin{equation}\label{eq:unit-center-target-test}
 0\le v+\epsilon\le3,\qquad\epsilon\in\{-1,0\}.
\end{equation}
The first special target has incoming carry zero: its single
positive coefficient precedes all later row intervals, making
every raw coefficient through that target nonnegative and below
$B$. Thus $0\le\delta^2\le3$, so $\delta\in\{0,1\}$.

At $\delta=0$, the direct guard \eqref{eq:unit-center-guard} has
raw value $-x^2<0$, which fails
\eqref{eq:unit-center-target-test}. Hence $\delta=1$ without any
copy assumption. Ordinary targets now test $F\in\{0,1\}$, and
paired rows force $F=0$. In particular, the paired unit copies give
$X'=X$. The two rows \eqref{eq:unit-center-unit} then impose
\[
 X-1\in\{0,1\},\qquad X-X^2\in\{0,1\},
\]
forcing $X=1$. Every original circuit gate and equality is recovered
at the actual input. The guard is auxiliary; it is not asserted to
be an original equation or to have zero target value.

\subsection{Positive necessity and the exact saving}

Given a solution of the represented system, take its circuit values,
$\delta=1$, all unit copies $X'=X=1$, zero dummy coordinates, and
\[
 u=\left\lceil\frac{x^2}{2}\right\rceil.
\]
The direct guard then has raw value
$2u-x^2+1=1+(x^2\bmod2)$, which is one or two and survives either
incoming carry. All ordinary targets have raw value one, and the
first special target has value one with carry zero. The low
coordinate tests vanish by support separation.

Choose a power-of-two $b$ larger than every coordinate and use the
canonical codes, $B=Hb^2$, $q=B^L$ and $n=q^8$. Their positive
slack witnesses exist as before. In particular, $e_0$ has unit
digit one, the $\delta$ digit makes $g>0$, and the nonconstant
polynomial $\ell_0(T)+e_0(T)T^L$ has nonnegative coefficients, so
its value at $B$ exceeds its value at four and the congruence
quotient $t$ is positive. Also $\Omega=\lambda-e>0$ and
\[
 \mathcal C^2<q\le\theta q,
 \qquad
 S_3=\lambda(\theta q-\mathcal C^2)+e\mathcal C^2>0.
\]
All three masks hold. Their newly calculated $r$ satisfies
\eqref{eq:unit-center-ranges} and the binomial divisibility in
\eqref{eq:unit-center-decoded-powers}.

The positive Pell construction of Section~\ref{sec:composed-99}
now applies with these new coding values. Choose $U=4^{2r+1}$,
$w=U/n^2$, $Y_P=\lfloor(U+1)^{2r}/U^r\rfloor$, and $s=Y_P/n^2$.
Its binomial congruence makes the latter quotient integral.
Choose the main and first Pell coordinates and interval witnesses,
then the second coordinates at $L$, the gap and positive exponent
quotients. In particular
$\Delta=(\psi_{a+4}(L)-\psi_4(L))/a$ is a positive integer.
Finally choose the doubled-index auxiliary witnesses at the new
main parameter. This changes no coding value. Every required
positive unknown is therefore supplied in the same proved order.

The former seven-instruction block computed $2\lambda$, then
$3\lambda=2\lambda+\lambda$. Replace it by the six instructions
\begin{equation}\label{eq:unit-center-six}
 \begin{gathered}
 U_\lambda=\theta\lambda,\qquad
 T_{\mathrm{coef}}=U_\lambda+1,\qquad \lambda_3=3\lambda,\\
 R_2=T_{\mathrm{coef}}+\lambda_3,\qquad
 t_2=\lambda-e,\qquad P_2=U_\lambda q.
 \end{gathered}
\end{equation}
The tests $q^2=R_2$ and $\theta+4=B$ still give exactly
$\lambda(B-1)=q^2-1$. The new block has three multiplications and
three additions; the old block had three multiplications and four
additions. Every other arithmetic instruction is retained.

\begin{theorem}\label{thm:best98}
Every recursively enumerable subset of the positive integers has an
admissible fixed index $(V,H,T_L)$ for a system with 34 positive
unknowns and 22 equations admitting a straight-line certificate of
98 operations: 55 multiplications and 43 additions. Here
$T_L=\psi_4(L)$ for the underlying exponent of the rebuilt index.
The supplied parameters are $(x,V,H,T_L)$, and fixed numerals are free.
\end{theorem}
\begin{proof}
The preceding construction proves the representation in both
directions and gives every required positive witness. Formula
\eqref{eq:unit-center-six} saves one addition from
Theorem~\ref{thm:best99}, giving $55+43=98$.
The exact checker
\texttt{round27\_1980\_unit\_center\_certificate.py} verifies all
98 primitive instructions, all 22 source residuals, and the same
three residual corrections as its predecessor. Only the packed
$r$ equation and the positive-$\sigma$ and positive-$\Omega$
equations change; geometry and every Pell source equation are
unchanged. Its finite regressions are distinct from the universal
proof above. No optimality is asserted.
\end{proof}

% A relaxed auxiliary norm whose integrality follows from the size bounds.
\section{Recovering the integral auxiliary Pell sequence}
\label{sec:relaxed-auxiliary}

The auxiliary norm can be weakened while preserving the index information
used by the doubled signed block. Starting with Theorem~\ref{thm:best99},
replace only E16 by
\begin{equation}\label{eq:relaxed-auxiliary-norm}
 (ic^2)^2=D(f^2-1),\qquad D=A^2-1,\qquad A=a+4.
\end{equation}
The previous norm was $f^2-1=D(ic^2)^2$. The two individual equations
are not equivalent: the new equation can admit additional rational
Pell solutions. The existing size bounds exclude the unwanted
solutions and recover the divisibility needed by the signed-index
argument. At the arithmetic level, the square $(ic^2)^2$ is now
exactly the coefficient $K=D(f^2-1)$ used by that argument, saving
one multiplication.

\subsection{A preliminary bound on the main index}

Keep the definitions
\[
 U=wn^2,\quad Y_P=sn^2,\quad a=Y_P(U+1),\quad
 A=a+4,\quad P=2UY_P^2+1,\quad J=2r+1.
\]
Before any auxiliary norm is used, the coding bootstrap supplies
\[
 n\ge64,\qquad n\le r<2n^3,\qquad U,Y_P\ge n^2,
 \qquad UY_P>r+1.
\]
The unchanged main norm and first triangular norm give positive
indices $p,t$ with
\[
 d=\chi_A(p),\quad c=\psi_A(p),\qquad
 2\tau+1=\chi_P(t),\quad k=\psi_P(t).
\]
Since $P\equiv1\pmod{UY_P}$ and $k=r+1+hUY_P$, the congruence
for $\psi_P$ gives $t\equiv r+1\pmod{UY_P}$. Thus $t\ge r+1$.
Also
\[
 P-A=UY_P(2Y_P-1)-Y_P-3>0.
\]
If $p\le r+1$, Pell growth would give
\[
 c\le(2A)^{p-1}\le(2A)^r
  <(2P-1)^r\le k,
\]
contrary to $c=kY_P+\eta>k$. Consequently
\begin{equation}\label{eq:relaxed-main-index-bound}
 p\ge r+2\ge66,\qquad c>A^6>A D^2.
\end{equation}
For the second assertion, use
$\psi_A(p)\ge(2A-1)^{p-1}$ and $D<A^2$.
This argument uses neither the auxiliary norm, the signed index
equation, an exponential identity, nor internal coefficient decoding.

\subsection{The rank argument}

\begin{lemma}\label{lem:relaxed-pell-integrality}
Let $A>1$, $D=A^2-1$, and let
$c=\psi_A(p)$, $d=\chi_A(p)$ for a positive index $p$.
Suppose $c>A D^2$ and positive integers $f,i$ satisfy
\eqref{eq:relaxed-auxiliary-norm}. Then there is a positive index $m$
such that
\[
 f=\chi_A(m),\qquad p\mid m,\qquad c\mid m,
 \qquad ic^2=D\psi_A(m).
\]
\end{lemma}
\begin{proof}
Write $D=d_0s_0^2$ with $d_0>1$ squarefree. We first justify the
use of a fundamental unit within the integer-coefficient ring
$\mathbb Z[\sqrt{d_0}]$. The known unit $A+s_0\sqrt{d_0}$ ensures
that positive solutions to $x^2-d_0y^2=1$ exist. Choose the least
positive unit
\[
 \varepsilon=F+y_0\sqrt{d_0}>1,
 \qquad F,y_0\in\mathbb Z_{>0},\qquad
 F^2-d_0y_0^2=1.
\]
Such a choice exists by minimizing its first coordinate. Every other
positive integer-coefficient unit $\beta>1$ is a power of
$\varepsilon$. Indeed, choose $j\ge0$ with
$\varepsilon^j\le\beta<\varepsilon^{j+1}$. The quotient
$\beta\varepsilon^{-j}$ has integer coefficients and norm one,
because $\varepsilon^{-1}=F-y_0\sqrt{d_0}$. If it belonged to
$(1,\varepsilon)$, both coordinates would be positive integers,
contradicting minimality. It is therefore one. No assertion about
the full ring of algebraic integers is needed.

There is accordingly an $e\ge1$ such that
\[
 A=\chi_F(e),\qquad s_0=y_0S,\qquad
 S=\psi_F(e),\qquad D=(F^2-1)S^2.
\]
Composition of the Pell solutions gives
\begin{equation}\label{eq:relaxed-composed-coordinate}
 \psi_F(ep)=Sc.
\end{equation}
Put $R_0=ic^2$. The relaxed equation gives
$R_0^2=d_0s_0^2(f^2-1)$. Thus $s_0\mid R_0$, and squarefreeness
implies $d_0\mid R_0/s_0$. The positive integer
$y=R_0/(d_0s_0)$ satisfies $f^2-d_0y^2=1$. Hence for some
$n_0>0$,
\begin{equation}\label{eq:relaxed-field-coordinate}
 f=\chi_F(n_0),\qquad
 R_0=(F^2-1)S\psi_F(n_0).
\end{equation}

We use the strong divisibility identity
\begin{equation}\label{eq:relaxed-strong-divisibility}
 \gcd(\psi_F(u),\psi_F(v))=\psi_F(\gcd(u,v)).
\end{equation}
For completeness, the Pell norm gives
$\gcd(\chi_F(j),\psi_F(j))=1$. The addition formula therefore
reduces $\gcd(\psi_F(u),\psi_F(u+v))$ to
$\gcd(\psi_F(u),\psi_F(v))$, and the Euclidean algorithm proves
\eqref{eq:relaxed-strong-divisibility}.

Let $T=(F^2-1)S$, $g=\gcd(n_0,ep)$, and
\[
 h=\gcd(\psi_F(n_0),\psi_F(ep))=\psi_F(g).
\]
Because $c^2\mid R_0=T\psi_F(n_0)$,
\[
 \begin{aligned}
 \gcd(c^2,T\psi_F(ep))
   &=\gcd(c^2,Dc)\\
   &=c\gcd(c,D)\ge c
 \end{aligned}
\]
divides $T h$. Since $T=D/S\le D$,
\begin{equation}\label{eq:relaxed-gcd-lower-bound}
 h\ge c/T\ge c/D.
\end{equation}
If $g<ep$, then $2g\le ep$, since $g$ is a proper divisor. The
duplication identity gives
\[
 2h^2<2\chi_F(g)\psi_F(g)=\psi_F(2g)
       \le\psi_F(ep)=Sc.
\]
Here $S<A$, by $D=(F^2-1)S^2$ and $F\ge2$. Thus $h^2<Ac/2$,
whereas \eqref{eq:relaxed-gcd-lower-bound} and $c>AD^2$ give
$h^2\ge c^2/D^2>Ac$, a contradiction. Therefore $ep\mid n_0$.
Writing $n_0=em$ proves $p\mid m$ and, by composition,
\begin{equation}\label{eq:relaxed-integrality-restored}
 f=\chi_A(m),\qquad R_0=D\psi_A(m).
\end{equation}

To obtain $c\mid m$, write $m=pk$. Expansion of
$(d+c\sqrt D)^k$, followed by division of the square-root
coefficient by $c$, gives
\[
 \frac{\psi_A(pk)}c\equiv k d^{k-1}\pmod c.
\]
As $c^2\mid D\psi_A(m)$ and $\gcd(d,c)=1$, this implies
$c\mid Dk$. The main Pell expansion modulo $D$ also gives
\[
 c=\psi_A(p)\equiv pA^{p-1}\pmod D,
 \qquad \gcd(c,D)=\gcd(p,D),
\]
using $\gcd(A,D)=1$. Put $g_0=\gcd(c,D)$. Then $c/g_0\mid k$
and $g_0\mid p$, so $c\mid pk=m$.
\end{proof}

The lemma recovers precisely the integral sequence and auxiliary
index divisibility needed below. It does not assert
$c^2\mid\psi_A(m)$, which is stronger than necessary and is not
an automatic consequence of the relaxed norm.

\subsection{The signed index and all positive witnesses}

By \eqref{eq:relaxed-main-index-bound},
Lemma~\ref{lem:relaxed-pell-integrality} applies before the signed
auxiliary norm. Its source polynomial remains
\[
 v=of-d^2,\qquad K=D(f^2-1)=(ic^2)^2,\qquad
 v(v+1)=K(K-1)(J+jc)^2.
\]
For $G=2K-1$, reduction modulo $f$ still gives
\[
 G\equiv-\chi_A(2),\qquad
 2v+1\equiv-\chi_A(2p)\pmod f.
\]
The conclusions $f=\chi_A(m)$ and $c\mid m$ give the comparison
bound $0<2p\le c\le m$. Also $K=(ic^2)^2$ gives
$G\equiv-1\pmod c$ directly. The signed chi step-down and parity
argument in the proof of Lemma~\ref{lem:doubled-main-index} therefore
apply unchanged: they yield $J\equiv\pm p\pmod c$, exclude
$J+p=c$, and conclude $p=J$. The old stronger divisibility of the
second Pell coordinate is not used at this step.

The exact first index, ratio estimates, common exponential witness,
fixed-index recovery through $T_L=\psi_4(L)$, and coefficient
decoding now proceed in the order proved for
Theorem~\ref{thm:best99}. All their hypotheses have been recovered
before use. This proves sufficiency of the modified system.

For necessity, start with that theorem's positive witness construction.
Its auxiliary index is the even number $m=2cJ$, and
\[
 f=\chi_A(m),\qquad i_{\mathrm{old}}=\psi_A(m)/c^2>0.
\]
Keep $f$ and set $i_{\mathrm{new}}=D i_{\mathrm{old}}>0$. Then
\[
 (i_{\mathrm{new}}c^2)^2
  =D^2\psi_A(m)^2=D(f^2-1).
\]
The source value $K=D(f^2-1)$ has not changed. Thus the same odd
$f$, the same $G$, and the positive $o,j$ constructed in
Lemma~\ref{lem:doubled-main-index} still satisfy the signed norm.
All other witnesses and fixed index components are preserved.

\subsection{One multiplication is removed}

Previously E16 and the shared coefficient $K$ required six
instructions:
\[
 \begin{gathered}
 z=ic^2,\quad z_2=z^2,\quad AE=Dz_2,\\
 f_2=f^2,\quad R_{16}=AE+1,\quad K=D\,AE.
 \end{gathered}
\]
The relaxed equation needs five:
\[
 z=ic^2,\quad K=z^2,\quad f_2=f^2,\quad
 f_-=f_2-1,\quad R_{16}=D f_-.
\]
Test $K=R_{16}$ and use the already computed square $K$ in the
signed norm. All other instructions remain unchanged.

The exact triangular correction also becomes simpler. If
\[
 F_{16}=(ic^2)^2-D(f^2-1),\quad
 K_s=D(f^2-1),\quad K_c=(ic^2)^2,
\]
then
\[
 F_{17}^{\mathrm{calc}}
 =F_{17}^{\mathrm{source}}
  -F_{16}(K_s+K_c-1)(J+jc)^2.
\]
The complete verifier checks this correction, the retained geometric
and second-exponent corrections, all source residuals and primitive
instructions, and the input domains. The Pell integrality and index
claims are proved above, independently of the symbolic checks.

\begin{theorem}\label{thm:best98-pell}
Replacing E16 by \eqref{eq:relaxed-auxiliary-norm} in the system of
Theorem~\ref{thm:best99} preserves its existential representation over
34 positive integer unknowns and 22 equality tests. Its straight-line
certificate has
\[
 \boxed{98=54\text{ multiplications}+44\text{ additions}}.
\]
Fixed numerals and supplied index parameters are free. The saving
relative to Theorem~\ref{thm:best99} is one multiplication.
\end{theorem}

The full schedule, source polynomials and exact residual identities are
recorded in the standalone checker
\begin{center}\small
\texttt{round28\_1980\_relaxed\_auxiliary\_certificate.py}
\end{center}
and its JSON receipt.

% Compose the unit coefficient baseline and relaxed auxiliary norm.
\section{A certificate of 97 operations}\label{sec:composed-97}

The unit-centered encoding of Section~\ref{sec:unit-center} and the
relaxed auxiliary norm of Section~\ref{sec:relaxed-auxiliary} compose.
The arithmetic substitutions are disjoint, but their mathematical
compatibility also requires an ordered proof: the new coding bounds
must force a sufficiently large main Pell index before the relaxed
norm is used, and the Pell conclusions must precede digit decoding.

\begin{theorem}\label{thm:best97}
Every recursively enumerable subset of the positive integers has an
admissible fixed index $(V,H,T_L)$ for a system with 34 positive
unknowns and 22 equations admitting a straight-line certificate of
97 operations: 54 multiplications and 43 additions. Here
$T_L=\psi_4(L)$ for the underlying exponent of the index construction.
The supplied parameters are $(x,V,H,T_L)$, with three fixed index
parameters besides the positive input $x$. Fixed numerals and
equality tests are free.
\end{theorem}

\begin{proof}
Use the entire construction of Section~\ref{sec:unit-center} for
the represented system. Compile it into nonnegative primitive
circuit coordinates with distinct operands and fresh outputs.
Pair multiplication, addition, copy, zero and final equality rows
$F$ with their negatives, using targets $2F+\delta^2$. For each
unit coordinate $X$, introduce $X'$, include the paired copy rows
$(X-X')\delta$ and $(X'-X)\delta$, and the two unpaired unit rows
\[
 X\delta-\delta^2,\qquad X\delta-XX'.
\]
Place the special target $\delta^2$ first, and include the direct
guard $2u\delta-x^2+\delta^2$. The actual input square is used in
the guard, independently of any copy. Every divided target
coefficient belongs to $\{-1,0,1\}$.

With $m$ positive-weight coordinates and $s$ targets, use
\[
\begin{gathered}
 v_i=3^i\quad(0\le i<m),\qquad v_0=1\text{ for }\delta,
 \qquad M=3^{m-1},\qquad d_0=4M+1,\\
 t_j=(s+1)d_0+2M+jd_0\quad(0\le j<s),\\
 t_*=2sd_0+2M,\qquad K_0=t_*+1,\qquad c_*=m+s+1.
\end{gathered}
\]
The input has weight zero. Include a dummy coordinate at each $t_j$.
The complementary polynomial $\mathcal D$ satisfies
$[T^{t_j}]\mathcal D(T)\mathcal C(T)^2=G_j$ exactly, by the proved
monomial and row separation. Form
\[
 \ell_0(T)=\sum_iT^{v_i}+\sum_jT^{t_j},\qquad
 e_0(T)=\sum_{j=0}^{K_0-1}T^j+\mathcal D(T).
\]
Choose a power of two $L>3K_0+2$ and a power of two
\begin{equation}\label{eq:composed97-index}
 \begin{gathered}
 H>\max\{2\,4^{2L+1},\ 4^{t_*+3}\,4c_*^2,\ 3L,\ 16\},\\
 V=\ell_0(4)+e_0(4)4^L,
 \qquad T_L=\psi_4(L).
 \end{gathered}
\end{equation}
These three index values use the complete new coordinate and target
list. They are fixed independently of $x$; the underlying $L$ is
not an additional supplied parameter. The digits of $e_0$ are in
$\{0,1,2\}$, with unit digit one. The special target has no incoming
carry by the support proof in Section~\ref{sec:unit-center}.

Retain the coding equations and packing
\begin{align*}
 B&=Hb^2,&\mathcal C&=x+g,&b&=x+\beta,\\
 \lambda(B-1)&=q^2-1,&\theta+4&=B,&n&=q^8,\\
 S_2&=\ell+eq=V+t\theta,&S_2+\alpha&=q^2,\\
 \Omega&=\lambda-e>0,&
 S_3&=\theta\lambda q-\Omega\mathcal C^2=\sigma>0,\\
 S&=g+q^2(S_2+q^2S_3),\\
 T+1&=q^2(1+\theta\lambda)-(b-1)\ell+\theta\ell q^4,\\
 r&=S(n^2-n)+(T+1)(n^2-1).
\end{align*}
In the Pell block write $A=a+4$ and $D=A^2-1$. Replace only its
auxiliary norm by
\begin{equation}\label{eq:composed97-auxiliary}
 (ic^2)^2=D(f^2-1).
\end{equation}
The doubled signed norm still uses
\begin{equation}\label{eq:composed97-signed}
 \begin{gathered}
 v=of-d^2,\qquad K_{\mathrm{aux}}=D(f^2-1),\qquad J=2r+1,\\
 v(v+1)=K_{\mathrm{aux}}(K_{\mathrm{aux}}-1)(J+jc)^2.
 \end{gathered}
\end{equation}
The fixed index equation remains $\kappa=T_L+\Delta a$, and every
other Pell source equation is unchanged. The integer register $v$
may be signed; every supplied unknown, including $o$, is positive.

Before using any auxiliary Pell equation or internal target,
the positive packed bound and geometry give
\[
 e<q,\quad\ell<q^2,\quad B\le q^2,\quad
 q>4b\ge8,\quad3L\le B\le n.
\]
Also $\Omega\ge1$ implies
$\mathcal C^2<\theta\lambda q<3q^3<q^4$ and $0<S_3<q^4$.
The unit-center preliminary borrow argument consequently proves
\begin{equation}\label{eq:composed97-ranges}
 0<S<n,\qquad0<T+1\le n,\qquad
 n\le r\le2n^3-2n^2<2n^3.
\end{equation}
Put $U=wn^2$, $Y_P=sn^2$, $a=Y_P(U+1)$ and $P=2UY_P^2+1$.
Then $n\ge64$, $U,Y_P\ge n^2$, $UY_P>r+1$ and $a>n^4>J$.
The positive interval and first-index equation give $c>J$.

The unchanged main norm and triangular first norm give positive
indices $p,t$ such that
\[
 c=\psi_A(p),\quad d=\chi_A(p),\qquad
 k=\psi_P(t),\quad2\tau+1=\chi_P(t).
\]
Since $P\equiv1\pmod{UY_P}$ and $k=r+1+hUY_P$,
$t\equiv r+1\pmod{UY_P}$. The inequality $0<r+1<UY_P$ therefore
gives $t\ge r+1$, without knowing the main index.
Moreover
\[
 P-A=UY_P(2Y_P-1)-Y_P-3>0.
\]
If $p\le r+1$, then
\[
 c\le(2A)^{p-1}\le(2A)^r
   <(2P-1)^r\le\psi_P(t)=k,
\]
contradicting $c=kY_P+\eta>k$. Thus
\begin{equation}\label{eq:composed97-large-main-index}
 p\ge r+2\ge66,
 \qquad c=\psi_A(p)>A^6>A D^2.
\end{equation}
The growth bound $\psi_A(p)\ge(2A-1)^{p-1}$ proves the second
inequality. Neither the relaxed norm, the signed norm nor an
exponential equation has been used here.

Lemma~\ref{lem:relaxed-pell-integrality} now applies to
\eqref{eq:composed97-auxiliary} and
\eqref{eq:composed97-large-main-index}. Its integer-coefficient
fundamental-unit and strong-divisibility argument supplies an index
$m_0$ with
\begin{equation}\label{eq:composed97-recovered-auxiliary}
 f=\chi_A(m_0),\qquad p\mid m_0,\qquad c\mid m_0,
 \qquad ic^2=D\psi_A(m_0).
\end{equation}
These are exactly the needed conclusions; the stronger old
divisibility $c^2\mid\psi_A(m_0)$ is not asserted.

Equation \eqref{eq:composed97-auxiliary} identifies
$K_{\mathrm{aux}}=(ic^2)^2$, so the signed auxiliary parameter
$G=2K_{\mathrm{aux}}-1$ is $-1$ modulo $c$. In conjunction with
\eqref{eq:composed97-recovered-auxiliary}, the doubled signed-index
argument has all its hypotheses, including
$0<2p\le c\le m_0$. It gives $J\equiv\pm p\pmod c$.
Both indices lie strictly between zero and $c$. The possibility
$J+p=c$ contradicts odd $J$ and $\psi_A(p)\equiv p\pmod2$.
Hence $p=J$, and
\[
 c=\psi_A(J),\qquad d=\chi_A(J).
\]

The subsequent common-witness proof proceeds in its established
order. It recovers the exact first index, proves $Y_P\ge U^r$ and
$a>U^{r+1}$, and then obtains $U=4^J$ from the first exponent
equation. The fixed value $T_L=\psi_4(L)$, reduced modulo $a$,
determines the second Pell index as $L$: both possible residues lie
in $(0,a)$ because $a>4^{3J}$. Its exponent comparison gives
$q=B^L$. The necessary bounds follow from
\eqref{eq:composed97-ranges} and $3L\le B\le n\le r$.
Finally the unchanged upper ratio and binomial fractional-part
estimates give
\begin{equation}\label{eq:composed97-binomial}
 b,B,q\text{ are powers of two},\qquad
 q=B^L,\qquad n^2\mid\binom{2r}{r}.
\end{equation}
These conclusions precede coefficient decoding.

Now $\lambda\ge B^{2L-1}\ge Bq>2q$ and $q\ge B^2>2B$.
Thus $\Omega=\lambda-e>\lambda/2$ and
\[
 \mathcal C^2<2\theta q<2Bq<q^2,
 \qquad g<\mathcal C<q.
\]
The first two masks remove their preliminary borrow and recover the
combined coefficient polynomial up to its quotient alias. They
also bound each decoded coordinate by $b$. Before identifying
$\delta$, positivity of $x,g$ gives $\mathcal C(1)\ge2$.
With $A_{\mathcal C}=\mathcal C(1)^2$, the high plateau of
$\lambda\mathcal C^2$ has height $A_{\mathcal C}$. The high-mask
digits of $(B-4)a_{\mathrm{al}}$ are consequently bounded by
$A_{\mathcal C}+4\le4A_{\mathcal C}$. The bound on $H$ in
\eqref{eq:composed97-index} and evaluation at four force the alias
to vanish, by the complete argument in Section~\ref{sec:unit-center}.
Both codes are canonical.

The first special target has incoming carry zero and forces
$\delta\in\{0,1\}$. The direct guard excludes zero before any
copy equation is used. At $\delta=1$, paired rows force $F=0$;
unit copies agree and the two unit rows force $X=1$. All original
circuit equations therefore hold at the actual input $x$. This
proves sufficiency of the composition.

For necessity, take a solution of the represented system, its
circuit values, $\delta=1$, unit copies $X'=X=1$, and zero dummy
digits. Put $u=\lceil x^2/2\rceil$, so the guard target is one or
two. Choose a power-of-two $b$ larger than every coordinate and
use the canonical codes, $B=Hb^2$, $q=B^L$ and $n=q^8$.
Section~\ref{sec:unit-center} supplies all positive coding witnesses,
the three masks, \eqref{eq:composed97-ranges} and
\eqref{eq:composed97-binomial} for the newly calculated $r$.

Construct the full doubled-index Pell witnesses as in
Section~\ref{sec:composed-99}. Choose $U=4^J$,
$Y_P=\lfloor(U+1)^{2r}/U^r\rfloor$, $a=Y_P(U+1)$, the main and
first Pell coordinates, positive interval witnesses, the second
coordinates at $L$, and the positive exponent quotients. All these
choices use the new fixed value $T_L=\psi_4(L)$. In particular
$w=U/n^2$, $s=Y_P/n^2$ and
$\Delta=(\psi_{a+4}(L)-\psi_4(L))/a$ are positive integers.
Its auxiliary construction chooses
\[
 m_0=2cJ,\qquad f=\chi_A(m_0),\qquad
 i_{\mathrm{old}}=\frac{\psi_A(m_0)}{c^2}>0,
\]
and supplies positive $o,j$ for the doubled signed norm. Replace
only the auxiliary witness
\[
 i=D i_{\mathrm{old}}>0.
\]
Then
\[
 (ic^2)^2=D^2\psi_A(m_0)^2=D(f^2-1),
\]
while the mathematical value $K_{\mathrm{aux}}=D(f^2-1)$ is
unchanged. Consequently the same $f,o,j$ still satisfy the signed
equation. The witness $i$ occurs nowhere else. Every one of the
34 unknowns has therefore received a positive integer value.

For the count, the unit-centered encoding saves one addition from
the 99-operation certificate, and the relaxed auxiliary norm saves
one multiplication from that same certificate. The full checker
\texttt{round29\_1980\_composed\_certificate.py} verifies that both
orders give the same 97 primitive instructions and the same 22
source residuals. Relative to 99, only the packed equation,
auxiliary norm, and positive-$\sigma$ and positive-$\Omega$
equations change. Hence the final count is $54+43=97$.

The three residual corrections are acyclic: the geometric and
second-exponent corrections use the independently exact equation
$\theta+4=B$, while the signed norm correction is
\[
 -F_{16}(K_{\mathrm{source}}+K_{\mathrm{calculated}}-1)(J+jc)^2,
 \qquad F_{16}=(ic^2)^2-D(f^2-1),
\]
whose auxiliary equality is checked exactly. Each subtraction is
one reversed addition equation with a free integer intermediate.
Thus 97 is the full addition/multiplication certificate length for
the proved positive-domain system. No optimality is claimed.
\end{proof}

% Linear radix scaling, paired zero tests, and two-digit reset padding.
\section{Linear radix scaling and a certificate of 96 operations}
\label{sec:linear-radix}

The quadratic radix scale can be replaced by a linear one when the
coefficient tests are allowed two base-$B$ digits. An explicit bound
on the indicator code removes its quotient alias. Positive padding
then resets the incoming carry at every ordinary target, and a
second unit test excludes a scale of order $\sqrt B$.

Starting from Section~\ref{sec:composed-97}, replace the radix and
third packed block by
\begin{equation}\label{eq:linear-radix-block}
 B=Hb,\qquad \mathcal C=x+g,\qquad
 \Omega=\lambda-e>0,\qquad
 \sigma=\Omega(q-\mathcal C^2)>0,
\end{equation}
and add the positive gap equation
\begin{equation}\label{eq:linear-radix-gap}
 \ell+\alpha_2=q.
\end{equation}
Every Pell equation is retained, with $B=Hb$ wherever the radix
occurs. The coefficient index is rebuilt below.

\begin{theorem}\label{thm:best96}
Every recursively enumerable subset of the positive integers has
an admissible fixed index $(V,H,T_L)$ for a system with 35 positive
unknowns and 23 equations admitting a straight-line certificate of
96 operations: 52 multiplications and 44 additions. Here
$T_L=\psi_4(L)$ for the underlying exponent in the index
construction. The supplied parameters are $(x,V,H,T_L)$, with
three fixed index parameters besides the positive input $x$.
Fixed numerals and equality tests are free.
\end{theorem}

The following subsections prove the theorem. The complete arithmetic
table is in Appendix~\ref{app:96}. The exact source and primitive
checker is located in \file{../verification/}:
\begin{center}
\small\file{round30\_1980\_linear\_radix\_certificate.py}.
\end{center}

\subsection{The homogeneous circuit and its normalization rows}

Compile the represented finite Diophantine system into an
arithmetic circuit over nonnegative integers, using the actual
positive input $x$, addition, multiplication, zero and one.
Introduce a homogeneous coordinate $\delta$. Repeated operands
are replaced by fresh copies, and every gate output is fresh.
All circuit coordinates other than $x$ have positive weights in
the code.

For a multiplication, addition or zero equation use, respectively,
\begin{equation}\label{eq:linear-radix-gates}
 2XY-2W\delta=0,\qquad
 2(X+Y-W)\delta=0,\qquad
 2X\delta=0.
\end{equation}
All products in these expressions have distinct factors. A copy
$X'=X$ can instead be imposed by
\begin{equation}\label{eq:linear-radix-copy}
 X'^2-X^2=0.
\end{equation}
Nonnegativity makes this an exact copy equation. It also supplies
arbitrarily many distinct copies of $\delta$ without introducing
a diagonal coefficient of magnitude two. A circuit occurrence of
the constant one is represented by $\delta$ or a copy of $\delta$.
Whenever using $\delta$ directly would create $\delta\delta$ in
\eqref{eq:linear-radix-gates}, use a fresh square-difference copy
instead. Likewise give repeated gate operands distinct copies.
Thus every product in \eqref{eq:linear-radix-gates} has distinct
factors, including when a gate operand is the constant one. Final
equalities can use either distinct-coordinate products with
$\delta$ or square differences. Include both $F$ and $-F$ for
every zero equation $F=0$ in this construction.

For normalization, introduce distinct coordinates
$X,Y,\delta',u$ and include the following equations and their
negatives:
\begin{equation}\label{eq:linear-radix-normalization}
 \begin{gathered}
 X^2-x^2=0,\qquad Y^2-x^2=0,\qquad
 \delta'^2-\delta^2=0,\\
 2xX-2\delta\delta'-2u\delta=0.
 \end{gathered}
\end{equation}
Once their zero tests have been decoded, these equations give
\begin{equation}\label{eq:linear-radix-guard}
 X=Y=x,\qquad \delta'=\delta,\qquad
 x^2=\delta^2+u\delta.
\end{equation}
Thus $\delta>0$ because $x>0$, and $\delta\le x$.
Conversely, $\delta=\delta'=1$, $X=Y=x$ and $u=x^2-1$ are
nonnegative witnesses for every positive $x$, including $x=1$.

Finally include two unpaired targets with raw polynomial
$\delta^2$. The second will receive extra positive padding.
These two tests force $\delta=1$ below; the argument does not
presume an incoming carry at an unprotected special target.

Divide every square coefficient of a target polynomial by one
and every distinct-coordinate product coefficient by two. All
results lie in $\{-1,0,1\}$. This includes
\eqref{eq:linear-radix-gates}, \eqref{eq:linear-radix-copy},
both signs of \eqref{eq:linear-radix-normalization}, and the
two unit targets. The number of rows and encoded coordinates
may depend on the represented system; they are digits of the
integer codes, not additional unknowns of the final system.

\subsection{Two-digit support and reset positions}

Let $m$ count all non-input circuit and normalization coordinates.
Give them weights
\[
 v_i=4\,3^i\quad(0\le i<m),\qquad M=4\,3^{m-1}.
\]
The input $x$ has weight zero. Quadratic monomial weights determine
their unordered coordinate pair, including squares and products
with $x$: divide by four and read the base-three digits. All these
weights are multiples of four and are at most $2M$.

Let $s$ be the number of targets, counting both signs of each
zero equation and the two final unit tests. Put the unit tests
last, and distinguish the last target $t_*$ as the additional-carry
test. Define
\begin{equation}\label{eq:linear-radix-support}
 \begin{gathered}
 d_0=4M+4,\qquad
 t_j=(s+1)d_0+2M+jd_0\quad(0\le j<s),\\
 t_*=t_{s-1}=2sd_0+2M,\qquad K=t_*+2.
 \end{gathered}
\end{equation}
All $t_j$ are multiples of four. For a target $G_j$ and a monomial
of weight $w$ with divided coefficient $a$, put $aT^{t_j-w}$ in
$\mathcal D_{\rm main}(T)$. The intervals
$[t_j-2M,t_j]$ are disjoint, so these instructions do not collide.
Add the reset polynomial and the last-target padding:
\begin{equation}\label{eq:linear-radix-padding}
 \begin{aligned}
 \mathcal D_{\rm reset}(T)&=\sum_jT^{t_j-2},\\
 \mathcal D_{\rm extra}(T)&=T^{t_*-1}
                  +T^{t_*-1-v_X}+T^{t_*-1-v_Y},\\
 \mathcal D(T)&=\mathcal D_{\rm main}(T)
                 +\mathcal D_{\rm reset}(T)
                 +\mathcal D_{\rm extra}(T).
 \end{aligned}
\end{equation}
The three parts have exponents congruent to $0$, $2$ and $3$
modulo four, respectively. The extra exponents are distinct.
All $\mathcal D$ coefficients therefore still lie in
$\{-1,0,1\}$. Every exponent is positive, greater than $M$,
and less than $K$.

Include a dummy coordinate at each $t_j$ and $t_j+1$, and set
\begin{equation}\label{eq:linear-radix-codes}
 \begin{aligned}
 \ell_0(T)&=\sum_iT^{v_i}+\sum_j(T^{t_j}+T^{t_j+1}),\\
 e_0(T)&=\sum_{h=0}^{K-1}T^h+\mathcal D(T).
 \end{aligned}
\end{equation}
The digits of $\ell_0$ are zero or one, and those of $e_0$ are
zero, one or two. The unit digit of $e_0$ is one. Let $c_*$ count
the input, all $m$ true coordinates and the $2s$ dummies, and put
\[
 D_1=\sum_h\bigl|[T^h]\mathcal D(T)\bigr|.
\]
Choose a power of two $L>3K+2$ and a power of two $H$ such that
\begin{equation}\label{eq:linear-radix-index}
 \begin{gathered}
 H>\max\{2\,4^{2L+1},\ 64\max(1,D_1)c_*^2,\ 3L,\ 64\},\\
 V=\ell_0(4)+e_0(4)4^L,\qquad T_L=\psi_4(L).
 \end{gathered}
\end{equation}
These choices depend only on the represented system and precede
the queried input $x$.

No arbitrary decoded dummy value contributes to
$\mathcal D(T)\mathcal C(T)^2$ through position $t_*+1$.
Indeed, the least $\mathcal D$ exponent is at least $t_0-2M$,
the least dummy weight is $t_0$, and
\begin{equation}\label{eq:linear-radix-dummies}
 2t_0-2M>t_*+1.
\end{equation}
This matters because the dummies at $t_j+1$ are not in the zero
residue class modulo four. Up to the last tested digit, however,
they have no effect on the product.

\subsection{Preliminary bounds and the Pell conclusions}

Retain the coding equations and packing
\begin{equation}\label{eq:linear-radix-packing}
 \begin{aligned}
 \mathcal C&=x+g,& b&=x+\beta,& B&=Hb,&\theta&=B-4,\\
 \lambda(B-1)&=q^2-1,& n&=q^8,\\
 S_2&=\ell+eq=V+t\theta,&S_2+\alpha&=q^2,\\
 S&=g+q^2(S_2+q^2\sigma),\\
 T+1&=q^2(1+\theta\lambda)-(b-1)\ell+\theta\ell q^4,\\
 r&=S(n^2-n)+(T+1)(n^2-1).
 \end{aligned}
\end{equation}
All supplied unknowns, including $\alpha$, $\alpha_2$,
$\Omega$ and $\sigma$, are positive integers. Before any circuit
row or Pell conclusion is used, positivity gives
\begin{equation}\label{eq:linear-radix-preliminary}
 \ell<q,\quad e<q,\quad S_2<q^2,\quad
 \mathcal C^2<q,\quad g<q,\quad
 0<\sigma<\lambda q<q^3.
\end{equation}
The last inequality follows from
$\lambda=(q^2-1)/(B-1)<q^2$. Geometry gives $B\le q^2$.
Also $b\ge2$ and \eqref{eq:linear-radix-index} imply
$q\ge6$, $b<B<n$ and $3L\le B\le n$. The former estimate
$q>4b$ is not used or asserted for the linear radix. Moreover,
\[
 \theta\lambda\ge B-4=Hb-4>b.
\]
The integer bounds in \eqref{eq:linear-radix-preliminary} give
$0<S<q^7<n$. For the other packed number,
\begin{equation}\label{eq:linear-radix-packed-range}
 \begin{aligned}
 T+1&>q^2\theta\lambda-b\ell>bq^2-bq>0,\\
 T+1&<q^4(1+\theta\ell)<Bq^5\le q^7<n.
 \end{aligned}
\end{equation}
The upper bound uses $1+\theta\lambda<q^2$ and
$1+\theta\ell\le1+(B-4)(q-1)<Bq$. Consequently
\begin{equation}\label{eq:linear-radix-r-range}
 n\le n^2-1\le r<2n^3.
\end{equation}

These are all the packing-size inputs to the main Pell and
binomial arguments. Write
\[
 \begin{gathered}
 U=wn^2,\qquad Y_{\rm p}=s_{\rm p}n^2,
 \qquad a=Y_{\rm p}(U+1),\qquad A=a+4,\\
 D_{\rm p}=A^2-1,\qquad P=2UY_{\rm p}^2+1,
 \qquad J=2r+1.
 \end{gathered}
\]
The notation $Y_{\rm p},s_{\rm p}$ distinguishes the Pell values
from the encoded input copy $Y$. We have
$U,Y_{\rm p}\ge n^2$, $UY_{\rm p}\ge n^4>r+1$ and
$a>n^4>J$. The first index equation and positive interval give
$c>J$. The norm classifications and $P>A$ imply that the preliminary
main index is at least $r+2$, hence $c>AD_{\rm p}^2$, exactly as
in Section~\ref{sec:composed-97}.

The retained relaxed auxiliary norm
\[
 (ic^2)^2=D_{\rm p}(f^2-1)
\]
therefore has the same fundamental-unit and strong-divisibility
hypotheses. It recovers the required auxiliary multiple of $c$.
The doubled signed-index argument identifies
$c=\psi_A(J)$ and $d=\chi_A(J)$. The first-index and ratio
argument gives $Y_{\rm p}\ge U^r$ and $a>U^{r+1}$ before either
exponent relation is used. The common-witness exponent equation
then gives $U=4^J$. The retained fixed congruence
$\kappa=T_L+\Delta a$ determines the second index as $L$, and
the second exponent equation gives $q=B^L$.

Every exponent-size hypothesis follows from
\eqref{eq:linear-radix-r-range} and $3L\le B\le n\le r$ as
in the preceding proof. None requires $B=Hb^2$. Since $n^2$ divides
$U=4^J$, $q$ is a power of two; $q=B^L$ makes $B$ a power of two;
and the fixed power of two $H$ divides $B=Hb$, so $b$ is a power
of two. The unchanged upper ratio and binomial fractional-part
estimates give
\begin{equation}\label{eq:linear-radix-binomial}
 n^2\mid\binom{2r}{r}.
\end{equation}
This applies the full Pell argument before any coefficient
equation or unit value is assumed. Its necessity construction
will be used only after the new coding witnesses have been chosen.

\subsection{Canonical codes without a high-mask estimate}

After the Pell conclusions, $b<B<q$ and $\ell<q$. Therefore
\begin{equation}\label{eq:linear-radix-no-borrow}
 q^2-1-(b-1)\ell>0.
\end{equation}
The packed number $T$ now has the exact three-block expansion
\[
 T=\bigl(q^2-1-(b-1)\ell\bigr)
                +q^2\theta\lambda+q^4\theta\ell.
\]
The blocks have widths $(2,2,4)$ in radix $q$. The binary
no-carry packing lemma, using
\eqref{eq:linear-radix-packing}, \eqref{eq:linear-radix-r-range}
and \eqref{eq:linear-radix-binomial}, gives
\begin{equation}\label{eq:linear-radix-masks}
 \begin{gathered}
 g\mathbin{\&}\bigl(q^2-1-(b-1)\ell\bigr)=0,\\
 S_2\mathbin{\&}\theta\lambda=0,\qquad
 \sigma\mathbin{\&}\theta\ell=0.
 \end{gathered}
\end{equation}
Here $\mathbin{\&}$ denotes bitwise intersection only in the
proof of the polynomial system's meaning.

Since $q=B^L$ and $\theta=B-4$, the middle mask restricts the
base-$B$ digits of $S_2$ to $\{0,1,2,3\}$ in fewer than $2L$
positions. If $F(T)$ is its digit polynomial, then $F(B)=S_2$,
and the congruence in \eqref{eq:linear-radix-packing} gives
$F(4)\equiv V\pmod{B-4}$. Both $F(4)$ and $V$ are less than
$4^{2L}$, while \eqref{eq:linear-radix-index} makes $B-4$
larger than their possible difference. Hence $F(4)=V$.
Uniqueness of the base-four expansion gives
\[
 S_2=\ell_0(B)+e_0(B)q.
\]
Both $\ell,e$ and $\ell_0(B),e_0(B)$ are less than $q$.
Uniqueness of quotient and remainder on division by $q$ yields
\begin{equation}\label{eq:linear-radix-canonical}
 \ell=\ell_0(B),\qquad e=e_0(B).
\end{equation}
No high-mask estimate proportional to the square of the digit
bound is required.

The first mask, together with $b$ and $B$ being powers of two,
now bounds every coordinate digit of $g$ by $b-1$ and permits
nonzero digits only at the indicator positions. Its unit digit
is zero, so the unit coordinate of $\mathcal C$ is the queried
$x$. Write
\[
 \mathcal C(T)=x+\sum_i z_iT^{v_i}
              +\sum_j(d_jT^{t_j}+d'_jT^{t_j+1}).
\]
All decoded coordinates are nonnegative and below $b$. Put
$A_{\mathcal C}=\mathcal C(1)^2$. Then
\begin{equation}\label{eq:linear-radix-coefficient-bound}
 A_{\mathcal C}<(c_*b)^2<B^2/64.
\end{equation}
Also $e,\mathcal C<B^K$, and $L>3K+2$ gives
$e\mathcal C^2<q$. Through position $K-1$, the raw signed
coefficients of $\sigma$ are those of
$\mathcal D(T)\mathcal C(T)^2$: the $q$ terms in
\eqref{eq:linear-radix-block} start at $L$, and the difference
between the finite baseline and $\lambda$ starts at $K$.
Each raw coefficient has absolute value at most $A_{\mathcal C}$,
because the coefficients of $\mathcal D$ have absolute value
at most one.

\subsection{Reset padding and exact paired zero tests}

Consider the signed base-$B$ carry recursion
\begin{equation}\label{eq:linear-radix-carry}
 \begin{gathered}
 h_0=0,\qquad a_j=[T^j]\mathcal D(T)\mathcal C(T)^2,\\
 h_{j+1}=\left\lfloor\frac{a_j+h_j}{B}\right\rfloor,
 \qquad \mathrm{digit}_j=a_j+h_j-Bh_{j+1}.
 \end{gathered}
\end{equation}
It computes the actual low digits of the positive integer
$\sigma$. Bound \eqref{eq:linear-radix-coefficient-bound}
gives $|h_j|<B/8$ by induction. Indeed,
\[
 \left|\left\lfloor\frac{a_j+h_j}{B}\right\rfloor\right|
 \le \frac B{64}+\frac18+1<\frac B8
\]
when $|a_j|<B^2/64$, $|h_j|<B/8$ and $B>64$.

By \eqref{eq:linear-radix-dummies}, dummies do not affect the
tested range. All true non-input weights are zero modulo four.
The raw coefficients can therefore be nonzero only in residue
classes $0$, $2$ and $3$ modulo four; class $1$ is empty. The
carry into $t_j-2$, following an empty class-$1$ position, is
consequently either $-1$ or zero.

At $t_j-2$, the raw coefficient is nonnegative and at least
$x^2$: the row's own reset contributes $x^2$, and every other
reset contribution is nonnegative. It is less than $B^2/64$.
The carry into $t_j-1$ is therefore a nonnegative integer
$\alpha_j<B/64$.

At every target except $t_*$, the raw coefficient at $t_j-1$
is zero. Only $\mathcal D_{\rm extra}$ contributes in this
residue class, and its contributions are within distance $2M$
of $t_*-1$; different targets are separated by $4M+4$.
Dividing $\alpha_j<B$ at this empty position leaves exactly
zero incoming carry at $t_j$.

There is no raw coefficient at $t_j+1$. Thus the two tested
digits at $t_j,t_j+1$ form the residue modulo $B^2$ of the raw
target polynomial $G_j$. The third mask permits both digits
to lie only in $\{0,1,2,3\}$.

For a paired zero equation, $|G_j|<B^2/64$. A negative nonzero
value has a residue modulo $B^2$ whose high digit exceeds three,
so it fails the mask. Both $G_j$ and $-G_j$ have their own
reset and zero incoming carry. They can both pass exactly when
$G_j=0$. This proves every circuit and normalization equation
exactly, before $\delta$ is known to be one. In particular
\eqref{eq:linear-radix-guard} holds.

The first unpaired $\delta^2$ target also receives zero incoming
carry. The sole exception to this rule is the deliberately
modified last target.

\subsection{The last carry fixes the unit}

The first unit test has raw value $\delta^2$ and zero incoming
carry. Its two allowed digits give
\begin{equation}\label{eq:linear-radix-unit-digits}
 \delta^2=kB+s,\qquad k,s\in\{0,1,2,3\}.
\end{equation}
The guard gives $0<\delta\le x$, and the coordinate bound gives
$\delta<b=B/H<B/4$. Since $B$ is a sufficiently large power
of two, $s=2$ or $3$ is impossible modulo four. If $s=1$, the
square roots of one modulo $B$ are congruent to
$1,B/2-1,B/2+1,B-1$. The bound $\delta<B/4$ forces
$\delta=1$. The only remaining case for $\delta>1$ is $s=0$,
and therefore
\begin{equation}\label{eq:linear-radix-large-unit}
 \delta^2\ge B,\qquad x^2\ge B.
\end{equation}

At $t_*-1$, monomial uniqueness and the input-copy equations
make the raw extra-padding coefficient exactly
\[
 x^2+2xX+2xY=5x^2.
\]
The reset at $t_*-2$ remains present and supplies a nonnegative
carry $\alpha_*<B/64$. Thus the incoming carry at the last
target is
\begin{equation}\label{eq:linear-radix-unit-carry}
 k_{\rm pad}=\left\lfloor\frac{5x^2+\alpha_*}{B}\right\rfloor
 \ge0.
\end{equation}
The global coefficient bound also gives $k_{\rm pad}<B/8<B$.
If \eqref{eq:linear-radix-large-unit} held, then
$k_{\rm pad}\ge5$. Since $\delta^2$ is a multiple of $B$,
the low digit at the last target would be $k_{\rm pad}$,
between five and $B-1$, contradicting its mask. Therefore
$\delta=1$.

The circuit equations now recover the represented system at the
actual positive input $x$. This proves sufficiency.

\subsection{Necessity and all positive supplied witnesses}

Given a solution of the represented system, take its circuit
values, $\delta=\delta'=1$, $X=Y=x$, $u=x^2-1$, all fresh
copies equal, and every dummy zero. Each paired row has value
zero, and both unpaired unit targets have value one.

The fixed index was already chosen in
\eqref{eq:linear-radix-support}--\eqref{eq:linear-radix-index}.
Now choose a power of two $b$ arbitrarily large relative to the
finitely many coordinate values. Besides exceeding every
coordinate, choose it so that every raw coefficient of
$\mathcal D(T)\mathcal C(T)^2$ has absolute value less than
$B/4$ for $B=Hb$, and $5x^2<B/4$. This is possible because
the coordinates, $\mathcal D$ and their raw coefficients are
fixed before $b$ is chosen, whereas $Hb$ grows without bound.

Take the canonical codes, $q=B^L$ and $n=q^8$. The support
margin gives $\mathcal C^2<q$, and
$\lambda=(q^2-1)/(B-1)>e$. Hence $\Omega$ and $\sigma$ are
positive. The values
\[
 \begin{gathered}
 \alpha_2=q-\ell,\qquad
 \alpha=q^2-(\ell+eq),\qquad \beta=b-x,\\
 t=\frac{\ell+eq-V}{B-4}
 \end{gathered}
\]
are positive integers. Positivity of $t$ follows by evaluating
the nonconstant nonnegative polynomial
$\ell_0(T)+e_0(T)T^L$ at $B>4$. The coordinate $\delta=1$
ensures $g>0$.

The first two masks hold by the coordinate and digit bounds.
For the third, all variable positions and their preceding
positions lie below the least $\mathcal D$ exponent, so their
raw coefficients and carries are zero. Every paired target has
zero value and receives zero incoming carry from its reset.
The first unit target has value one. At the last target, the
chosen small-coefficient bounds give $\alpha_*=0$ and
$k_{\rm pad}=0$, so its tested digits are also one and zero.
Thus all masks hold. The packing lemma gives
\eqref{eq:linear-radix-binomial} for the newly determined $r$,
with the required bounds already established.

Construct the positive Pell witnesses at this new $r$ as in
Section~\ref{sec:composed-97}: choose $U=4^J$,
\[
 Y_{\rm p}=\left\lfloor\frac{(U+1)^{2r}}{U^r}\right\rfloor,
 \qquad w=\frac U{n^2},\qquad
 s_{\rm p}=\frac{Y_{\rm p}}{n^2},\qquad a=Y_{\rm p}(U+1).
\]
Take the main and first Pell coordinates at $J$ and $r+1$,
their positive interval and first-index quotients, and the
second coordinates at $L$. The second index uses the newly
fixed $T_L=\psi_4(L)$. The exponent and rounding bounds give
every required positive quotient as in that proof.

Choose the doubled-index auxiliary exponent last as
$m_{\rm aux}=2cJ$. Set $f=\chi_A(m_{\rm aux})$ and
$i_{\rm old}=\psi_A(m_{\rm aux})/c^2>0$; the doubled-index
construction supplies positive $o,j$. Use
$i=D_{\rm p}i_{\rm old}$. This gives the relaxed norm
$(ic^2)^2=D_{\rm p}(f^2-1)$ and preserves the signed auxiliary
value. No coding witness changes. Together with the new positive
$\alpha_2$, all 35 supplied unknowns are positive integers.
This proves necessity and completes the mathematical proof of
Theorem~\ref{thm:best96}.

\subsection{The complete arithmetic count}

Relative to the 97-operation certificate, replacing $B=Hb^2$
by $Hb$ removes one multiplication. Replacing
\[
 \sigma=\theta\lambda q-(\lambda-e)\mathcal C^2
 \quad\hbox{by}\quad
 \sigma=(\lambda-e)(q-\mathcal C^2)
\]
replaces two multiplications and one subtraction by one
multiplication and one subtraction, removing another
multiplication. Equation~\eqref{eq:linear-radix-gap} costs
one addition. All normalization, padding and two-digit support
changes belong to the fixed index rather than to the arithmetic
of the final system. Therefore
\[
 97-1-1+1=96,\qquad
 (54\text{ multiplications},43\text{ additions})
 \longrightarrow(52\text{ multiplications},44\text{ additions}).
\]

The standalone checker verifies all 96 primitive instructions,
the complete 23 source residuals and their acyclic corrections.
The geometric and second-exponent corrections depend only on
the exact equality $\theta+4=B$. The signed auxiliary
correction depends only on the exact relaxed auxiliary norm.
The sparse coefficient/carry regression supplies additional
finite evidence; it does not replace the support, carry, index
or positive-witness arguments above. The result is an upper
bound, not an optimality claim.

% Affine radix, three-digit targets, and two independent unit carries.
% Prepared for integration after the round32 certificate and proof audits.
\section{Affine radix scaling and a certificate of 95 operations}
\label{sec:affine-radix}

An affine radix saves one further operation. Its first mask forbids
one fixed bit in each coordinate digit. Three coordinates with that
bit clear represent an arbitrary nonnegative circuit value. The
larger possible digits require three-digit target windows; reset
positions make paired zero equations exact, and two different
padding carries exclude every nonunit homogeneous scaling.

Start with the system of Section~\ref{sec:linear-radix}, replacing
its two positive gaps by the combined positive gap
\begin{equation}\label{eq:affine-radix-gap}
 \ell+e+\alpha=q.
\end{equation}
This preserves the two additions used for the gaps and removes
one positive unknown and one equation. Indeed, it implies
$\ell,e<q$ and
\[
 S_2=\ell+eq\le(e+1)(q-1)\le(q-1)^2<q^2.
\]
Conversely, the canonical codes below satisfy $\ell+e<q$.
Now replace the radix and the first packed mask by
\begin{equation}\label{eq:affine-radix-change}
 H=H_0+1,\qquad B=H+b,\qquad q^2-1-b\ell,
\end{equation}
where $H_0$ is a sufficiently large fixed power of two. The
supplied fixed index remains $(V,H,T_L)$; $H_0=H-1$ is not an
additional parameter. Every Pell equation is retained.

\begin{theorem}\label{thm:best95}
Every recursively enumerable subset of the positive integers has
an admissible fixed index $(V,H,T_L)$ for a system with 34 positive
unknowns and 22 equations admitting a straight-line certificate of
95 operations: 51 multiplications and 44 additions. Here
$T_L=\psi_4(L)$ for the underlying exponent in the index
construction. The supplied parameters are $(x,V,H,T_L)$, with
three fixed index parameters besides the positive input $x$.
Fixed numerals and equality tests are free.
\end{theorem}

The following proof includes the rebuilt coefficient index and
all positive witnesses. Appendix~\ref{app:95} gives the complete
arithmetic table. The exact source and primitive checker is in
\file{../verification/}:
\begin{center}
\small\file{round32\_1980\_affine\_radix\_certificate.py}.
\end{center}

\subsection{Logical values and physical coordinates}

Compile the represented finite Diophantine system into a circuit
over nonnegative integers with actual positive input $x$,
addition, multiplication, zero and one. Signed integer witnesses
can first be represented by differences of nonnegative witnesses.
Introduce a homogeneous coordinate $\delta$, retained as one
physical coordinate. Every other logical non-input coordinate is
the sum of three distinct nonnegative physical coordinates.
Different logical coordinates have disjoint groups. The input $x$
is unsplit and has weight zero.

Every nonnegative integer $z$ has such a representation with the
$H_0$ bit clear in each entry. If that bit of $z$ is zero, take
$(z,0,0)$. If it is one, take
\begin{equation}\label{eq:affine-radix-split}
 (z-H_0,\ H_0/2,\ H_0/2).
\end{equation}
All entries are nonnegative and have the required bit clear.
They become base-$B$ digits once $B$ exceeds these finitely many
values. Their choice depends on $z,H_0$, not on how much larger
$B$ is subsequently chosen. In the reverse direction, the sum
of three decoded physical digits is a nonnegative logical value;
no further bound on that sum is needed.

Use fresh outputs and fresh square-difference copies whenever a
product would repeat a logical factor. For multiplication,
addition and zero equations use respectively
\begin{equation}\label{eq:affine-radix-gates}
 2XY-2W\delta=0,\qquad
 2(X+Y-W)\delta=0,\qquad 2X\delta=0.
\end{equation}
A copy $X'=X$ can be imposed by
\begin{equation}\label{eq:affine-radix-copy}
 X'^2-X^2=0.
\end{equation}
Nonnegativity makes this exact. Constant-one occurrences use
$\delta$ or a fresh logical copy of $\delta$. Whenever a use
in \eqref{eq:affine-radix-gates} would create $\delta\delta$,
use a fresh copy instead. Repeated operands are copied too.
Every displayed product thus has distinct logical factors and
expands into distinct physical factors. Final equalities may
use square differences or distinct-factor products with
$\delta$. Include both $F$ and $-F$ for every zero equation
$F=0$.

Introduce fresh logical coordinates $X,Y,Z,\delta',u$ with
disjoint physical groups and include both signs of
\begin{equation}\label{eq:affine-radix-normalization}
 \begin{gathered}
 X^2-x^2=0,\quad Y^2-x^2=0,\quad Z^2-x^2=0,\quad
 \delta'^2-\delta^2=0,\\
 2xX-2\delta\delta'-2u\delta=0.
 \end{gathered}
\end{equation}
Their exact zero tests will give
\begin{equation}\label{eq:affine-radix-guard}
 X=Y=Z=x,\qquad \delta'=\delta,\qquad
 x^2=\delta^2+u\delta.
\end{equation}
Consequently $\delta>0$ and $\delta\le x$. Conversely, every
positive input has the nonnegative normalization witnesses
\begin{equation}\label{eq:affine-radix-unit-witness}
 \delta=\delta'=1,\quad X=Y=Z=x,\quad u=x^2-1.
\end{equation}

Add three unpaired targets whose raw polynomial is $\delta^2$,
and put them last. The first is ordinary; the second and third
will have different padding at their preceding positions.

Expand each logical row in the physical coordinates and $x$.
Divide square coefficients by one and distinct-coordinate product
coefficients by two. Every divided coefficient lies in
$\{-1,0,1\}$. A square of a three-coordinate sum has diagonal
coefficients one and cross coefficients two. Disjoint groups and
fresh copies prevent two logical monomials from contributing to
the same physical monomial in
\eqref{eq:affine-radix-gates}--\eqref{eq:affine-radix-normalization}.
The three unit targets have divided coefficient one. These rows
and physical coordinates affect the fixed index, not the number
of unknowns of the final system.

\subsection{Support, reset positions, and the fixed index}

Let $m$ count all true non-input physical coordinates, including
the unsplit $\delta$. Assign weights
\[
 v_i=6\,3^i\quad(0\le i<m),\qquad M=6\,3^{m-1}.
\]
The input has weight zero. Quadratic monomial weights are
multiples of six, at most $2M$, and determine their unordered
coordinate pair, including squares and products with $x$.
Divide by six and read the base-three digits to see this.

Let $s$ count all targets, including both signs of every zero
equation and the three final unit targets. Define
\begin{equation}\label{eq:affine-radix-support}
 \begin{gathered}
 d_0=4M+6,\qquad
 t_j=(s+1)d_0+2M+jd_0\quad(0\le j<s),\\
 t_{\max}=t_{s-1}=2sd_0+2M,\qquad K=t_{\max}+3.
 \end{gathered}
\end{equation}
All targets are multiples of six. Write
$\tau_5=t_{s-2}$ and $\tau_7=t_{s-1}$ for the two modified
unit targets. For a target $G_j$, place divided coefficient $a$
of monomial weight $w$ at exponent $t_j-w$ in
$\mathcal D_{\rm main}(T)$. The intervals $[t_j-2M,t_j]$
are disjoint. Multiplication by the physical coordinate
polynomial squared recovers exactly $G_j$ at $t_j$.

Add the resets
\begin{equation}\label{eq:affine-radix-resets}
 \mathcal D_{\rm reset}(T)=\sum_j T^{t_j-3}.
\end{equation}
Let $\mathcal P(X)$ be the three physical indices belonging to
logical coordinate $X$, and set
\begin{equation}\label{eq:affine-radix-padding}
 \begin{aligned}
 \mathcal D_5(T)&=T^{\tau_5-1}
       +\sum_{p\in\mathcal P(X)\cup\mathcal P(Y)}
                                      T^{\tau_5-1-v_p},\\
 \mathcal D_7(T)&=T^{\tau_7-1}
       +\sum_{p\in\mathcal P(X)\cup\mathcal P(Y)\cup\mathcal P(Z)}
                                      T^{\tau_7-1-v_p},\\
 \mathcal D(T)&=\mathcal D_{\rm main}(T)
       +\mathcal D_{\rm reset}(T)+\mathcal D_5(T)+\mathcal D_7(T).
 \end{aligned}
\end{equation}
The three types have exponent residues zero, three, and five
modulo six. Each padding polynomial has distinct exponents, and
their intervals are disjoint because $\tau_7-\tau_5=d_0>M$.
All $\mathcal D$ coefficients remain in $\{-1,0,1\}$.
Every exponent is positive, greater than $M$, and below $K$.

Allow dummy physical digits at all three tested positions
$t_j,t_j+1,t_j+2$, and define
\begin{equation}\label{eq:affine-radix-codes}
 \begin{aligned}
 \ell_0(T)&=\sum_iT^{v_i}
              +\sum_j(T^{t_j}+T^{t_j+1}+T^{t_j+2}),\\
 e_0(T)&=\sum_{h=0}^{K-1}T^h+\mathcal D(T).
 \end{aligned}
\end{equation}
The digits of $\ell_0$ are zero or one, those of $e_0$ are zero,
one or two, and the unit digit of $e_0$ is one. Put
\[
 c_*=1+m+3s,\qquad D_1=\sum_h|[T^h]\mathcal D(T)|.
\]
Choose a power of two $L>3K+2$. After the complete layout is
fixed, choose a power of two $H_0$ satisfying
\begin{equation}\label{eq:affine-radix-index}
 \begin{gathered}
 H_0>\max\{2\,4^{2L+1},\ 128\max(1,D_1)c_*^2,\ 3L,\ 1024\},\\
 H=H_0+1,\qquad V=\ell_0(4)+e_0(4)4^L,\qquad T_L=\psi_4(L).
 \end{gathered}
\end{equation}
These choices precede the queried input. The fixed integer $H$
need not be a power of two.

No decoded dummy contributes to
$\mathcal D(T)\mathcal C(T)^2$ through $t_{\max}+2$.
The least $\mathcal D$ exponent is at least $t_0-2M$, the
least dummy weight is $t_0$, and
\begin{equation}\label{eq:affine-radix-dummies}
 2t_0-2M>t_{\max}+2.
\end{equation}
Thus the residue analysis below involves only true coordinates
of weight zero modulo six. This exclusion is needed because
some dummy weights are not zero modulo six.

The spacing also makes the relevant padding coefficients exact.
At $t_j-3$ the own reset contributes $x^2$. Another reset would
require a monomial weight of absolute value at least
$d_0>2M$, so cannot contribute. At $\tau_5-1$, the extra terms
give exactly $x^2+2xX+2xY$; at $\tau_7-1$, they give exactly
$x^2+2xX+2xY+2xZ$. Other targets' preceding positions receive
no extra term. Monomial uniqueness identifies these terms, the
different bands are separated by more than $4M$, and the other
parts of $\mathcal D$ have different residues modulo six.

\subsection{Preliminary bounds and the Pell conclusions}

The source definitions and coding equations are
\begin{equation}\label{eq:affine-radix-packing}
 \begin{gathered}
 \mathcal C=x+g,\quad b=x+\beta,\quad B=H+b,\quad\theta=B-4,\\
 \lambda(B-1)=q^2-1,\quad
 S_2=\ell+eq=V+t\theta,\quad\ell+e+\alpha=q,\\
 \Omega=\lambda-e>0,\quad
 \sigma=\Omega(q-\mathcal C^2)>0,\quad n=q^8,\\
 S=g+q^2(S_2+q^2\sigma),\\
 T+1=q^2(1+\theta\lambda)-b\ell+\theta\ell q^4,\\
 r=S(n^2-n)+(T+1)(n^2-1).
 \end{gathered}
\end{equation}
All supplied unknowns, including $\alpha,\beta,\Omega,\sigma$,
are positive integers. Before using any Pell conclusion or
decoded coefficient row, positivity gives
\begin{equation}\label{eq:affine-radix-preliminary}
 \ell,e<q,\quad S_2<q^2,\quad\mathcal C^2<q,\quad g<q,
 \quad0<\sigma<\lambda q<q^3.
\end{equation}
The combined gap even gives $S_2\le(q-1)^2$. Geometry gives
$B\le q^2$. Since $b\ge2$ and $H>1024$, we have $q\ge6$,
$b<B<n$, $3L\le B\le n$, and
\[
 \theta\lambda\ge B-4=H+b-4>b.
\]
Integer gaps give
\[
 0<S\le q-1+q^2(q^2-1)+q^4(q^3-1)<q^7<n.
\]
For the other packed number,
\begin{equation}\label{eq:affine-radix-packed-range}
 \begin{aligned}
 T+1&>q^2\theta\lambda-b\ell>bq^2-bq>0,\\
 T+1&<q^4(1+\theta\ell)<Bq^5\le q^7<n.
 \end{aligned}
\end{equation}
The upper bound uses $1+\theta\lambda<q^2$ and
$1+\theta\ell\le1+(B-4)(q-1)<Bq$. Consequently
\begin{equation}\label{eq:affine-radix-r-range}
 n\le n^2-1\le r<2n^3.
\end{equation}

These are the packing and growth hypotheses of the retained Pell
block. To specify the proof order, write
\[
 \begin{gathered}
 U=wn^2,\quad Y_{\rm p}=s_{\rm p}n^2,\quad
 a=Y_{\rm p}(U+1),\quad A=a+4,\\
 D_{\rm p}=A^2-1,\quad P=2UY_{\rm p}^2+1,\quad J=2r+1.
 \end{gathered}
\]
The subscripts distinguish Pell values from encoded logical
coordinates. We have $U,Y_{\rm p}\ge n^2$,
$UY_{\rm p}\ge n^4>r+1$, and $a>n^4>J$. The positive
interval and first index equation imply $c>J$. Classification
of the main and first norms, together with $P>A$, gives
preliminary main index at least $r+2$. Thus $c>AD_{\rm p}^2$,
as in Sections~\ref{sec:composed-97} and
\ref{sec:relaxed-auxiliary}.

The retained relaxed norm
\[
 (ic^2)^2=D_{\rm p}(f^2-1)
\]
has the required fundamental-unit and divisibility hypotheses.
The doubled signed-index argument gives $c=\psi_A(J)$ and
$d=\chi_A(J)$. The first-index and ratio argument gives
$Y_{\rm p}\ge U^r$ and $a>U^{r+1}$ before either exponent
relation is used. The common-witness exponent equation gives
$U=4^J$. The fixed congruence $\kappa=T_L+\Delta a$
determines the second index as $L$, and its exponent equation
gives $q=B^L$. All exponent-size inequalities follow from
\eqref{eq:affine-radix-r-range} and $3L\le B\le n\le r$
as in those proofs.

Since $n^2\mid U=4^J$, $q$ is a power of two; $q=B^L$ then
makes $B$ a power of two. From $B>H_0+1$ we get $B\ge2H_0$.
Neither $b$ nor $H$ is asserted to be a power of two, and no
Pell step requires it. The retained upper ratio and binomial
fractional-part argument gives
\begin{equation}\label{eq:affine-radix-binomial}
 n^2\mid\binom{2r}{r}.
\end{equation}
All these conclusions precede canonical decoding, internal
equations, and the proof that $\delta=1$.

\subsection{Canonical codes and the forbidden bit}

After the Pell conclusions, $b<B<q$ and $\ell<q$, so
\begin{equation}\label{eq:affine-radix-no-borrow}
 q^2-1-b\ell>0.
\end{equation}
The exact block expansion is
\[
 T=(q^2-1-b\ell)+q^2\theta\lambda+q^4\theta\ell,
\]
with widths $(2,2,4)$ in radix $q$. The binary no-carry
packing lemma, using \eqref{eq:affine-radix-packing},
\eqref{eq:affine-radix-r-range}, and
\eqref{eq:affine-radix-binomial}, gives
\begin{equation}\label{eq:affine-radix-masks}
 g\mathbin{\&}(q^2-1-b\ell)=0,\qquad
 S_2\mathbin{\&}\theta\lambda=0,\qquad
 \sigma\mathbin{\&}\theta\ell=0.
\end{equation}
The ampersand is bitwise intersection, used only to interpret
the polynomial system.

Since $q=B^L$ and $\theta=B-4$, the middle mask bounds each
base-$B$ digit of $S_2$ by three. Its digit polynomial $F(T)$
has degree less than $2L$ and $F(B)=S_2$. The source congruence
gives $F(4)\equiv V\pmod{B-4}$. Both $F(4),V<4^{2L}$,
whereas \eqref{eq:affine-radix-index} makes $B-4$ larger than
their possible difference. Hence $F(4)=V$. Uniqueness of
base-four expansion yields $S_2=\ell_0(B)+e_0(B)q$.
Both $\ell,e<q$, and both fixed codes are below $B^K<q$;
quotient and remainder uniqueness gives
\begin{equation}\label{eq:affine-radix-canonical}
 \ell=\ell_0(B),\qquad e=e_0(B).
\end{equation}
This removes the quotient alias without a quadratic digit bound.

The integer $q^2-1$ has every base-$B$ digit equal to $B-1$.
Subtracting $b\ell_0(B)$ causes no borrow. At indicator
positions the resulting digit is
\[
 B-1-b=H-1=H_0,
\]
and elsewhere it is $B-1$. Since $B,H_0$ are powers of two
with $H_0<B$, the first mask says exactly that each physical
digit of $g$ has its $H_0$ bit clear and that digits outside
the indicator support vanish. Every digit is nonnegative and
below $B$. The unit position is outside the support, so $g$
has unit digit zero. Since $x<b<B$, the unit coordinate of
$\mathcal C=x+g$ is exactly the queried input. Write
\[
 \mathcal C(T)=x+\sum_i z_iT^{v_i}
       +\sum_j(d_jT^{t_j}+d'_jT^{t_j+1}+d''_jT^{t_j+2}).
\]
Every physical digit, including the unsplit $\delta$, is below
$B$. The input is below $B$ too. Therefore
\begin{equation}\label{eq:affine-radix-coefficient-bound}
 A_{\mathcal C}=\mathcal C(1)^2<(c_*B)^2<B^3/64.
\end{equation}
The last inequality follows from the fixed bound and $B\ge2H_0$.
Every raw coefficient of $\mathcal D(T)\mathcal C(T)^2$ has
absolute value at most $A_{\mathcal C}$: distinct exponents
of $\mathcal D$ select distinct nonnegative coefficients of
$\mathcal C(T)^2$, and each multiplier has absolute value at
most one.

Also $e,\mathcal C<B^K$, hence $e\mathcal C^2<q$ because
$L>3K+2$. Through position $K-1$, the raw signed coefficients
of $\sigma$ are those of $\mathcal D(T)\mathcal C(T)^2$.
Indeed,
\[
 \sigma=(\lambda-e)q+(e-\lambda)\mathcal C^2.
\]
The first summand starts at $L$, and the coefficient polynomial
of $e-\lambda$ agrees with $\mathcal D(T)$ below $K$.
Terms starting at $K$ or $L$ cannot affect lower digits.
This identifies the actual low digits even when some raw
coefficients are negative.

\subsection{Reset padding and exact three-digit zero tests}

Use the signed carry recursion
\begin{equation}\label{eq:affine-radix-carry}
 \begin{gathered}
 h_0=0,\qquad a_j=[T^j]\mathcal D(T)\mathcal C(T)^2,\\
 h_{j+1}=\left\lfloor\frac{a_j+h_j}{B}\right\rfloor,
 \qquad\mathrm{digit}_j=a_j+h_j-Bh_{j+1}.
 \end{gathered}
\end{equation}
It computes the low base-$B$ digits of the positive integer
$\sigma$. The coefficient bound gives $|h_j|<B^2/8$ by
induction throughout the tested range: the induction step has
upper bound $B^2/64+B/8+1<B^2/8$ for $B>64$.

By \eqref{eq:affine-radix-dummies}, dummies do not contribute.
Only residue classes zero, three and five modulo six can have
nonzero raw coefficients; classes one, two and four are empty.
Two consecutive empty positions of classes one and two reduce
a carry of magnitude below $B^2/8$ to $-1$ or zero, since their
combined operation is division by $B^2$ followed by the floor.
Thus the incoming carry at every reset $t_j-3$ is in
$\{-1,0\}$.

Its raw coefficient is exactly $x^2$. Since $1\le x<B$, the
sum with the incoming carry is nonnegative and below $B^2$.
The carry to $t_j-2$ is a nonnegative integer below $B$.
That position is empty, so the carry to $t_j-1$ is exactly zero.
Except at $\tau_5,\tau_7$, the raw coefficient at $t_j-1$
is also zero. Every ordinary target therefore receives zero
incoming carry. Its next two raw coefficients are zero as well,
so its three tested digits form the residue of $G_j$ modulo
$B^3$. The third mask requires each digit to be in
$\{0,1,2,3\}$.

For a paired zero row, $|G_j|<B^3/64$. A negative nonzero
value has residue modulo $B^3$ greater than $63B^3/64$;
its top digit exceeds three, so it fails the mask. Both signs
have separate resets and zero incoming carry. They can both
pass only at zero, which passes both. Every circuit equation
and every normalization equation is now exact, without assuming
a unit value. In particular \eqref{eq:affine-radix-guard} holds.

The first unpaired unit target also receives zero carry. At
$\tau_5-1$ and $\tau_7-1$, the exact padding coefficients
and the input-copy equations give $5x^2$ and $7x^2$, with
zero incoming carry. Thus the target carries are exactly
\begin{equation}\label{eq:affine-radix-two-carries}
 h_5=\left\lfloor\frac{5x^2}{B}\right\rfloor,
 \qquad h_7=\left\lfloor\frac{7x^2}{B}\right\rfloor.
\end{equation}
Both targets still have raw value $\delta^2$ and two empty
following raw positions. Their tested digits are the residues
of $\delta^2+h_5$ and $\delta^2+h_7$ modulo $B^3$.

\subsection{Two different padding carries fix the unit}

The guard gives $0<\delta\le x<B$; the unsplit physical
digit bound also gives $\delta<B$ directly. Hence
$\delta^2<B^2$. The first unit test implies
\begin{equation}\label{eq:affine-radix-unit-digits}
 \delta^2=kB+s,\qquad k,s\in\{0,1,2,3\}.
\end{equation}
Since $4\mid B$, a square cannot have $s=2$ or $3$. If $s=1$,
then $\delta^2\le3B+3<B^2/16$ for $B\ge64$, so
$\delta<B/4$. The roots of one modulo a power of two $B\ge8$
are congruent to $1,B/2-1,B/2+1,B-1$; only the first lies in
$(0,B/4)$. Thus $\delta=1$.

For $\delta>1$ the only remaining case is
\begin{equation}\label{eq:affine-radix-large-unit}
 \delta^2=kB,\qquad1\le k\le3,\qquad x^2\ge kB.
\end{equation}
For $c=5,7$, write
\[
 h_c=m_cB+r_c,\qquad0\le r_c<B.
\]
Since $x<B$, we have $h_c<cB$. The tested integer is
$(k+m_c)B+r_c<(3+c)B<B^2$. There is no reduction modulo
$B^3$ and no overflow into its third digit. Acceptance forces
\begin{equation}\label{eq:affine-radix-carry-digits}
 0\le r_5,r_7\le3,\qquad
 0\le m_5,m_7\le3-k\le2.
\end{equation}
The two floors are related. For $z=x^2/B$, using braces for
fractional part,
\begin{equation}\label{eq:affine-radix-floor-relation}
 7h_5-5h_7=5\{7z\}-7\{5z\},\qquad
 -7<7h_5-5h_7<5.
\end{equation}
Thus its absolute value is less than seven. Substituting the
quotients and remainders and using $|7r_5-5r_7|\le21$ gives
\[
 |B(7m_5-5m_7)|<28.
\]
Since $B>32$, we obtain $7m_5=5m_7$. Its nonnegative
solutions are $(m_5,m_7)=(5h,7h)$. Bound
\eqref{eq:affine-radix-carry-digits} forces both quotients to
be zero. Then $h_5=r_5\le3$, whereas
\eqref{eq:affine-radix-large-unit} gives $h_5\ge5k\ge5$,
a contradiction. Therefore $\delta=1$.

The paired circuit equations now recover the represented system
at its actual positive input. This proves sufficiency.

\subsection{Necessity and all positive supplied witnesses}

Given a solution of the represented system, assign its logical
circuit values and fresh copies, and use
\eqref{eq:affine-radix-unit-witness} for normalization.
Split each non-input logical value other than $\delta$ into
three physical coordinates as above; retain $\delta=1$
unsplit. All physical values have their $H_0$ bit clear,
including $\delta$ since $H_0>1$. Set every dummy to zero.
Each paired target has value zero and the three unit targets
have value one.

The fixed index was chosen before this input. With the physical
assignment fixed, choose $B$ an arbitrarily large power of two.
Require $B>H+x$ and $B$ larger than every physical coordinate.
Also require every raw coefficient of
$\mathcal D(T)\mathcal C(T)^2$ to have absolute value below
$B/4$, and require $7x^2<B/4$. These conditions are possible
because all those finitely many values are fixed while $B$
grows without bound. Set
\[
 \begin{gathered}
 b=B-H,\quad\beta=b-x>0,\quad
 \ell=\ell_0(B),\quad e=e_0(B),\\
 q=B^L,\qquad n=q^8.
 \end{gathered}
\]
The sum of the two code polynomials has digits at most three,
so $\ell+e<B^K<q$. The support margin gives
$\mathcal C^2<q$ and $\lambda=(q^2-1)/(B-1)>e$.
Consequently
\[
 \begin{gathered}
 \alpha=q-\ell-e,\quad\Omega=\lambda-e,\quad
 \sigma=\Omega(q-\mathcal C^2),\\
 t=\frac{\ell+eq-V}{B-4}
 \end{gathered}
\]
are positive integers. For $t$, evaluate the nonconstant
nonnegative polynomial $\ell_0(T)+e_0(T)T^L$ at $B>4$;
its difference from its value at four is positive and divisible
by $B-4$. The encoded coordinate $\delta=1$ ensures $g>0$.

The first mask holds because the physical digits have the
specified bit clear; the second holds because the combined
code digits are at most three. For the third mask, each true
variable position lies below the least $\mathcal D$ exponent
and has zero raw coefficient and zero carry. Every paired
target receives zero carry from its reset and has value zero.
The first unit target has value one. Since $7x^2<B/4$,
both special carries vanish, so those targets also have digits
$(1,0,0)$. All masks hold. The packing lemma supplies
\eqref{eq:affine-radix-binomial} for the newly determined $r$;
the preliminary bounds give its required size estimates.

Construct every positive Pell witness for this $r$ as in
Section~\ref{sec:composed-97}. Choose $U=4^J$ and
\[
 Y_{\rm p}=\left\lfloor\frac{(U+1)^{2r}}{U^r}\right\rfloor,
 \qquad w=\frac U{n^2},\qquad
 s_{\rm p}=\frac{Y_{\rm p}}{n^2},\qquad a=Y_{\rm p}(U+1).
\]
Take the main and first Pell coordinates at $J$ and $r+1$,
their positive interval and first-index quotients, and the
second Pell coordinates at $L$. That index uses the newly
fixed $T_L=\psi_4(L)$. The retained exponent and rounding
bounds give the positive exponential quotients. They require
$B,q$ to be powers of two, as they are, but place no such
requirement on $b$ or $H$.

Choose the doubled-index auxiliary exponent last as
$m_{\rm aux}=2cJ$. Set $f=\chi_A(m_{\rm aux})$,
$i_{\rm old}=\psi_A(m_{\rm aux})/c^2>0$ and
$i=D_{\rm p}i_{\rm old}$. The doubled-index construction
supplies positive $o,j$. The relaxed auxiliary norm holds
and its signed auxiliary value is unchanged. Together with
the new positive $\beta$ and combined-bound $\alpha$, all
34 supplied unknowns are positive integers. This proves
necessity and completes the proof of Theorem~\ref{thm:best95}.

\subsection{The complete arithmetic count}

The combined-bound 96-operation schedule uses a multiplication
for $B=Hb$ and a subtraction for $b-1$. The new definition
$B=H+b$ uses one addition instead of one multiplication.
Replacing $(b-1)\ell$ by $b\ell$ preserves that multiplication
and removes the subtraction. Therefore
\[
 (52\text{ multiplications},44\text{ additions})
 \longrightarrow
 (51\text{ multiplications},44\text{ additions}),
 \qquad 51+44=95.
\]
There are 34 positive unknowns and 22 source equations. The
logical coordinates, physical triples, three-digit windows,
resets, and both padding patterns change only the fixed integer
index. No numeral construction is charged. Each subtraction
can be realized by one reversed addition equation with a free
integer intermediate.

The complete arithmetic checker compares the changed radix
and first-mask expressions with all source residuals, including
the retained acyclic corrections for geometry, the second
exponent, and the relaxed signed norm. The separate sparse
encoding regression gives finite evidence for coefficient
isolation and carries. It does not replace the universal
support, unit, Pell, or positive-witness arguments above.
The result is an upper bound, with no optimality assertion.

% Translate the fixed affine-radix parameter to remove one addition.
\section{Translating the fixed index: 94 operations}
\label{sec:shifted-affine}

The affine-radix construction permits one more saving by translating
a fixed index parameter. This changes neither its coefficient
encoding nor any positive witness. Let $H_{\rm old}=H_0+1$ be
the fixed parameter of Section~\ref{sec:affine-radix}, and supply
instead
\begin{equation}\label{eq:shifted-affine-index}
 H=H_{\rm old}-4=H_0-3.
\end{equation}
The actual radix is still
\begin{equation}\label{eq:shifted-affine-radix}
 B=H+b+4=H_{\rm old}+b,
\end{equation}
so the fixed-bit mask and every part of the encoding are unchanged.
In particular, $B-1-b=H_0$ as before. The positive source equation
for $\theta$ becomes
\begin{equation}\label{eq:shifted-affine-theta}
 \theta=H+b.
\end{equation}
All other source equations are obtained by replacing the former
fixed parameter $H_{\rm old}$ by $H+4$. Every occurrence of the
mathematical radix therefore has the value in
\eqref{eq:shifted-affine-radix}.

\begin{theorem}\label{thm:best94}
Every recursively enumerable subset of the positive integers has
an admissible fixed index $(V,H,T_L)$ for a system with 34 positive
unknowns and 22 equations admitting a straight-line certificate of
94 operations: 51 multiplications and 43 additions. Here
$H=H_0-3$, with $H_0$ chosen by
\eqref{eq:affine-radix-index}, and $T_L=\psi_4(L)$.
The supplied parameters are $(x,V,H,T_L)$; fixed numerals and
equality tests are free.
\end{theorem}

\begin{proof}
Fix an index $(V,H_{\rm old},T_L)$ furnished by
Theorem~\ref{thm:best95}. The translated index is
$(V,H_{\rm old}-4,T_L)$, with a positive second component because
$H_0>1024$. For every positive input $x$ and every tuple $\mathbf z$
of the 34 positive supplied unknowns, define the new source
residuals by
\[
 F_i^{94}(x,V,H,T_L;\mathbf z)
   =F_i^{95}(x,V,H+4,T_L;\mathbf z),\qquad 1\le i\le22.
\]
This is an exact polynomial substitution. The old radix equation
$\theta+4-(H_{\rm old}+b)=0$ simplifies to
\eqref{eq:shifted-affine-theta}; the remaining equations are
exactly their old values at the translated fixed index. Thus the
identity map on $\mathbf z$ is a bijection between the positive
solutions of the two source systems, in both directions.

For sufficiency, a solution of the translated system is a solution
of the 95-operation system at $H_{\rm old}=H+4$. Its complete
Pell, canonical-code, carry and unit arguments recover the
represented set at the actual input $x$. For necessity, use every
positive witness of that construction unchanged. In particular,
choosing a sufficiently large power of two $B$ gives
\[
 b=B-H-4=B-H_{\rm old},\qquad\beta=b-x>0,
\]
and preserves the same $\alpha$, codes, packing, and positive
Pell witnesses. The value $H_0=H+3$ is the same fixed power of
two, so its forbidden bit and all support bounds are identical.
This proves the full source-system equivalence without imposing
any new restriction on the witnesses.

It remains to realize the translated source equations with one
fewer primitive operation. In the 95-operation schedule the
two additions
\[
 B_{\rm reg}=H_{\rm old}+b,\qquad L_3=\theta+4
\]
are followed by the free equality test $L_3=B_{\rm reg}$.
The register $B_{\rm reg}$ has no other use in that schedule.
All other radix-dependent expressions already use the existing
$\theta$ registers, together with the exact source relation
$\theta+4=B$.

In the translated schedule retain just
\[
 \theta_{\rm sum}=H+b
\]
and replace the equality test by $\theta=\theta_{\rm sum}$.
Delete the instruction for $L_3$. Renaming the retained register
has no arithmetic cost. The new register is the value
of $\theta$; it is not the mathematical radix, which remains
$B=H+b+4$ and is never constructed as a separate register.
No instruction depending on it needs correction,
since its sole use is this equality. The other primitive
instructions and all other equality tests are retained, giving
\[
 51\text{ multiplications}+43\text{ additions}=94
\]
with 34 positive unknowns and 22 source equations. Translating
the fixed integer $H_{\rm old}$ has no numeral-construction cost.
\end{proof}

Appendix~\ref{app:94} gives the complete 94-instruction table.

The standalone checker is in \file{../verification/}:
\begin{center}
\small\file{round33\_1980\_shifted\_affine\_certificate.py}.
\end{center}
It checks the entire primitive schedule and all 22 translated
source residuals, including the retained acyclic corrections.
The unchanged coefficient construction retains the support,
carry, unit, and positive-witness proof of
Section~\ref{sec:affine-radix}. The number 94 is a proved upper
bound, with no assertion of optimality.

% Distributivity shortens the unchanged shifted-affine certificate.
\section{Factoring the packed mask: 93 operations}
\label{sec:factored-mask}

The system and fixed indices of Section~\ref{sec:shifted-affine}
admit a shorter arithmetic certificate by factoring the indicator
code out of two terms of the packed mask. No source equation,
encoding, admissibility condition, or positive unknown changes.

\begin{theorem}\label{thm:best93}
The universal system of Theorem~\ref{thm:best94}, with exactly
the same admissible fixed indices and positive witnesses, admits
a straight-line certificate of 93 operations: 50 multiplications
and 43 additions. There remain 34 positive unknowns and 22
equations. Its supplied parameters are $(x,V,H,T_L)$, with
three fixed index parameters besides the positive input $x$.
The literal numerals are $\{1,3,8,15\}$; fixed numerals and
equality tests are free.
\end{theorem}

\begin{proof}
The existing register $T_{\rm coef}=1+\theta\lambda$ gives
the packed mask in the form
\[
 T+1=q^2T_{\rm coef}-\ell b+\theta\ell q^4.
\]
The 94-operation schedule evaluates this expression using four
multiplications and two additions or subtractions:
\[
 \begin{aligned}
 T_{q^4}&=T_{\rm coef}q^2,&
 L_b&=\ell b,& T_2&=T_{q^4}-L_b,\\
 L_{q^4}&=\ell q^4,&
 M&=\theta L_{q^4},& T+1&=T_2+M.
 \end{aligned}
\]
Replace these six instructions by the following five:
\begin{equation}\label{eq:factored-mask-five}
 \begin{aligned}
 T_{q^4}&=T_{\rm coef}q^2,\qquad
 Q_\theta=\theta q^4,\qquad G=Q_\theta-b,\\
 I&=\ell G,\qquad T+1=T_{q^4}+I.
 \end{aligned}
\end{equation}
These use three multiplications and two additions or subtractions.
Distributivity gives the unconditional integer identity
\begin{equation}\label{eq:factored-mask-identity}
 q^2T_{\rm coef}+\ell(\theta q^4-b)
   =q^2T_{\rm coef}-\ell b+\theta\ell q^4.
\end{equation}
Thus the retained packed output, named \texttt{T} in the
certificate, has its old value for every integer assignment,
before any source equation is assumed. The deleted temporary
registers have no consumers outside this block and occur in
none of the free equality tests.

Every other instruction and equality test is retained. The
source residuals are therefore exactly those of the 94-operation
system. The identity map on the input parameters and all 34
positive supplied unknowns is the positive-witness map in both
directions. Each schedule determines its own temporary registers
from that same tuple. Consequently the full universality,
canonical-code, carry, unit, and Pell proof transfers without
any changed hypothesis.

The new temporary registers introduce no supplied unknown or
domain condition. They are in fact positive on the existing
domain: $\theta=H+b>b$, $q\ge1$, and hence
$\theta q^4-b>0$. Each subtraction is certified by its
reversed addition, as before. The complete count is
\[
 94-1=93=50\text{ multiplications}+43\text{ additions}.
\]
No numeral is introduced or changed.
\end{proof}

Appendix~\ref{app:93} gives all 93 instructions and the 22
unchanged free equality tests. The complete checker is in
\file{../verification/}:
\begin{center}
\small\file{round34\_1980\_factored\_mask\_certificate.py}.
\end{center}
It verifies the exact six-to-five replacement, the absence of
outside uses for deleted temporaries, equality of every retained
register, all primitive checks, and all source residuals with
their unchanged acyclic corrections. The shorter schedule is
a certificate for the same proved system; 93 is an upper bound,
with no optimality assertion.

% Full proof of the nine-operation replacement for the signed-index block.
\section{A smaller auxiliary Pell parameter: 92 operations}
\label{sec:half-parameter-pell}

The square $K=(ic^2)^2$ already present in the relaxed auxiliary norm
allows a smaller Pell parameter. Replacing the ten-operation signed
block by a nine-operation block saves one multiplication. The new
block uses one additional positive unknown and one additional equation.
Its positive witnesses are constructed afresh; the coefficient encoding
and fixed index do not change.

\begin{theorem}\label{thm:best92}
Every recursively enumerable subset of the positive integers has an
admissible fixed index $(V,H,T_L)$ for a system with 35 positive
unknowns and 23 equations admitting a straight-line certificate of
92 operations: 49 multiplications and 43 additions. The fixed index
and coefficient encoding are those of Theorem~\ref{thm:best93}.
The supplied inputs are $(x,V,H,T_L)$, with three fixed components
besides the positive input $x$. Fixed numerals and equality tests
are free.
\end{theorem}

\subsection{The new equations}

Retain the definitions
\begin{equation}\label{eq:half-pell-parameters}
 J=2r+1,\qquad A=a+4,\qquad D=A^2-1,\qquad
 R=ic^2,\qquad K=R^2.
\end{equation}
In particular, the relaxed auxiliary equation is retained:
\begin{equation}\label{eq:half-pell-relaxed}
 R^2=D(f^2-1).
\end{equation}
The former signed equation was
\[
 v(v+1)=K(K-1)(J+jc)^2,\qquad v=of-d^2.
\]
Replace it by
\begin{equation}\label{eq:half-pell-linear}
 u=J+jc,\qquad u=c+of,
\end{equation}
and
\begin{equation}\label{eq:half-pell-norm}
 K(u^2-y_{\rm aux}^2)=1-y_{\rm aux}^2.
\end{equation}
The quantity $u$ is an arithmetic register, not another supplied
unknown. The old variables $o,j$ and the new variable $y_{\rm aux}$
are positive integers. The last equation is equivalently
\begin{equation}\label{eq:half-pell-standard-norm}
 (Ru)^2-(R^2-1)y_{\rm aux}^2=1.
\end{equation}
Its ordinary Pell parameter is $R$, whereas the preceding signed
construction used $2R^2-1$. The certificate does not have to
construct either $Ru$ or that larger parameter.

\subsection{Independent bounds before the new block}

The preliminary proof of Section~\ref{sec:affine-radix}, retained
by the subsequent arithmetic improvements, gives
\begin{equation}\label{eq:half-pell-preliminary}
 \begin{gathered}
 n=q^8,\qquad n\ge64,\qquad n\le r<2n^3,\qquad 3L\le B\le n,\\
 U=wn^2,\qquad Y_P=s_Pn^2,\qquad a=Y_P(U+1),\qquad
 P=2UY_P^2+1,\\
 U,Y_P\ge n^2,\qquad UY_P>r+1,\qquad a>J,\qquad c>J>1.
 \end{gathered}
\end{equation}
The subscripts distinguish these Pell quantities from the encoded
logical coordinates. These bounds use positivity, the combined gap,
geometry, packing, and the positive interval. They precede coefficient
decoding and any use of either auxiliary signed norm.

Classification of the unchanged main and first norms gives positive
indices $p,t$ with
\[
 c=\psi_A(p),\qquad d=\chi_A(p),\qquad k=\psi_P(t).
\]
The first-index equation implies $t\ge r+1$. Since $P>A$ and
$c>Y_Pk>k$, the growth comparison in
Section~\ref{sec:relaxed-auxiliary} forces
\begin{equation}\label{eq:half-pell-large-index}
 p\ge r+2\ge66,\qquad c>AD^2.
\end{equation}
The retained relaxed-norm argument therefore applies to
\eqref{eq:half-pell-relaxed}. Its integer-coefficient fundamental-unit
and divisibility proof supplies a positive index $m$ such that
\begin{equation}\label{eq:half-pell-auxiliary-index}
 f=\chi_A(m),\qquad p\mid m,\qquad c\mid m,
 \qquad R=D\psi_A(m).
\end{equation}
That proof uses neither the former nor the new signed-index equation.
Elementary Pell growth now gives
\begin{equation}\label{eq:half-pell-comparison}
 0<2p\le c\le m,\qquad R=ic^2\ge c^2>A>1.
\end{equation}
In particular, $R$ is an integer Pell parameter greater than one.
No exact value of $p$, no parity of $r$, no conclusion $q=B^L$,
and no decoded circuit equation has been assumed here.

\subsection{The odd-index polynomial identity}

For the standard Pell sequences, write
\[
 \begin{gathered}
 \chi_X(0)=1,\quad\chi_X(1)=X,\qquad
 \psi_X(0)=0,\quad\psi_X(1)=1,\\
 z(n+2)=2Xz(n+1)-z(n).
 \end{gathered}
\]
For every $h\ge0$ there is an integer polynomial $Q_h(T)$ with
\begin{equation}\label{eq:half-pell-Q-definition}
 \chi_X(2h+1)=XQ_h(X^2).
\end{equation}
Its initial values and doubled-step recurrence are
\begin{equation}\label{eq:half-pell-Q-recurrence}
 \begin{gathered}
 Q_0(T)=1,\qquad Q_1(T)=4T-3,\\
 Q_{h+2}(T)=(4T-2)Q_{h+1}(T)-Q_h(T).
 \end{gathered}
\end{equation}
The doubled recurrence follows by applying the ordinary recurrence
twice, or from $\chi_X(2)=2X^2-1$ and the addition identities.
Thus these formulas prove divisibility by $X$ as an integer-polynomial
identity; they do not require division modulo an integer.

The odd subsequence $\psi_A(2h+1)$ has recurrence coefficient
$4A^2-2$ and initial values $1,4A^2-1$. Consequently
$(-1)^h\psi_A(2h+1)$ has coefficient $2-4A^2$ and initial
values $1,1-4A^2$. Substituting $T=1-A^2$ in
\eqref{eq:half-pell-Q-recurrence} gives exactly these data.
Induction proves
\begin{equation}\label{eq:half-pell-Q-complement}
 Q_h(1-A^2)=(-1)^h\psi_A(2h+1).
\end{equation}
Similarly, substitution of $T=0$ gives
\begin{equation}\label{eq:half-pell-Q-zero}
 Q_h(0)=(-1)^h(2h+1).
\end{equation}
Both identities hold for every nonnegative $h$. Their two uses below
have different moduli.

\subsection{Sufficiency: the exact main index}

Suppose the new source equations have positive integer witnesses.
By \eqref{eq:half-pell-standard-norm}, positivity, and $R>1$,
ordinary Pell classification gives a positive index $s$ such that
\begin{equation}\label{eq:half-pell-classification}
 Ru=\chi_R(s),\qquad y_{\rm aux}=\psi_R(s).
\end{equation}
The index $s$ is odd. Indeed, the chi recurrence at $R=0$ gives
$\chi_R(2h)\equiv(-1)^h\pmod R$. An even $s$ would make
$R>1$ divide one, since $R$ divides the first coordinate in
\eqref{eq:half-pell-classification}. Write $s=2h+1$.
Equation~\eqref{eq:half-pell-Q-definition} gives the exact equality
\begin{equation}\label{eq:half-pell-u-polynomial}
 u=Q_h(R^2).
\end{equation}

Reduction of \eqref{eq:half-pell-relaxed} modulo $f$ gives
\[
 R^2=D(f^2-1)\equiv1-A^2\pmod f.
\]
Equations \eqref{eq:half-pell-Q-complement},
\eqref{eq:half-pell-u-polynomial}, and the second equality in
\eqref{eq:half-pell-linear} therefore imply
\[
 (-1)^h\psi_A(s)\equiv u\equiv c=\psi_A(p)\pmod f.
\]
Squaring and using the doubling identity
\[
 \chi_A(2z)=1+2D\psi_A(z)^2
\]
gives
\begin{equation}\label{eq:half-pell-chi-comparison}
 \chi_A(2s)\equiv\chi_A(2p)\pmod{\chi_A(m)}.
\end{equation}

Apply Lemma~\ref{lem:signed-chi-new}: for $A>1$ and
$0<k\le m$, the congruence
$\chi_A(n)\equiv\pm\chi_A(k)\pmod{\chi_A(m)}$ implies
$n\equiv\pm k\pmod{2m}$. The positive-sign case also follows
from the ordinary chi step-down statement with its stronger
modulus $4m$. Here its comparison-index hypothesis is exactly
$0<2p\le m$ from \eqref{eq:half-pell-comparison}. Hence
\begin{equation}\label{eq:half-pell-stepdown}
 2s\equiv\pm2p\pmod{2m},\qquad
 s\equiv\pm p\pmod m,\qquad s\equiv\pm p\pmod c.
\end{equation}
The middle implication divides an integer equality
$2s=\pm2p+2mz$ by two. It does not assume that two is invertible
in an even residue ring.

Independently, $c\mid R$ because $R=ic^2$. Equations
\eqref{eq:half-pell-Q-zero} and
\eqref{eq:half-pell-u-polynomial} give
\[
 u\equiv(-1)^h s\pmod c.
\]
The first equality in \eqref{eq:half-pell-linear}, together with
\eqref{eq:half-pell-stepdown}, now yields
\begin{equation}\label{eq:half-pell-main-congruence}
 J\equiv\pm p\pmod c.
\end{equation}
Both $J$ and $p$ lie strictly between zero and $c$, by
\eqref{eq:half-pell-preliminary} and Pell growth. Thus either
$J=p$ or $J+p=c$. The latter is impossible: the psi recurrence
gives $\psi_A(p)\equiv p\pmod2$ for every integer $A$, whereas
$J=2r+1$ is odd. Therefore $J+p$ has parity opposite to
$c=\psi_A(p)$. We have proved
\begin{equation}\label{eq:half-pell-main-index}
 p=J,\qquad c=\psi_A(J),\qquad d=\chi_A(J).
\end{equation}
This deduction did not assume that $r$ is even. In particular, it
does not use the subsequent binary decoding circularly.

\subsection{The remaining sufficiency proof}

Once \eqref{eq:half-pell-main-index} is available, the retained
first-index and ratio arguments recover the first index $r+1$,
prove $Y_P\ge U^r$ and $a>U^{r+1}$, and use the common-witness
exponent equation to obtain $U=4^J$. The unchanged fixed value
$T_L=\psi_4(L)$, gap, and second Pell equations recover the second
index $L$ and $q=B^L$. The same bounds give
$n^2\mid\binom{2r}{r}$.

These are exactly the hypotheses and conclusions of
Section~\ref{sec:affine-radix}. The replaced auxiliary variables
$o,j$ occur nowhere in these later equations. Since $n^2\mid U$,
the integers $q$ and then $B$ are powers of two. Neither $b$ nor
$H$ is required to be a power of two. The unchanged masks recover
the canonical codes, the three-way physical coordinates, all paired
rows, and the unit $\delta=1$. They therefore recover the represented
system at its actual positive input $x$.

The supplied fixed index remains
\[
 H=H_0-3,\qquad V=\ell_0(4)+e_0(4)4^L,\qquad
 T_L=\psi_4(L),\qquad B=H+b+4.
\]
The new block adds no index parameter, compiler convention, or
support requirement.

\subsection{Necessity: parity and positive auxiliary witnesses}

Start with a solution of the represented Diophantine system. Choose
its canonical physical encoding and all positive coding and Pell
witnesses as in Section~\ref{sec:affine-radix}, translating the
fixed index and retaining the factored mask as in
Sections~\ref{sec:shifted-affine} and \ref{sec:factored-mask}.
All those choices remain available.

The canonical indicator has zero unit digit, so $B\mid\ell$.
The integer $g$ also has zero unit digit, so $B\mid g$.
Since $B$ and $q=B^L$ are even, both packed integers
\[
 \begin{aligned}
 S&=g+q^2(\ell+eq+q^2\sigma),\\
 T+1&=q^2(1+\theta\lambda)+\ell(\theta q^4-b)
 \end{aligned}
\]
are even. Hence $r=S(n^2-n)+(T+1)(n^2-1)$ is even, and
\begin{equation}\label{eq:half-pell-necessity-parity}
 J=2r+1\equiv1\pmod4.
\end{equation}

Keep the canonical relaxed auxiliary witnesses
\begin{equation}\label{eq:half-pell-necessity-auxiliary}
 \begin{gathered}
 m=2cJ,\qquad f=\chi_A(m),\qquad
 i_{\rm old}=\frac{\psi_A(m)}{c^2}>0,\\
 i=Di_{\rm old},\qquad R=ic^2=D\psi_A(m),\qquad K=R^2.
 \end{gathered}
\end{equation}
The quotient is integral: expanding
$(d+c\sqrt D)^{2c}$, its term with one square root contains
$2c^2$, while every later square-root term contains at least $c^3$.
Positivity follows from $m>0$. The Pell norm gives
\eqref{eq:half-pell-relaxed}. Thus the same $i,f$ can be retained.
Keep every other canonical witness except $o,j$, and define
\begin{equation}\label{eq:half-pell-new-witnesses}
 \begin{gathered}
 y_{\rm aux}=\psi_R(J),\qquad
 u=\frac{\chi_R(J)}R=Q_{(J-1)/2}(R^2),\\
 o=\frac{u-c}{f},\qquad j=\frac{u-J}{c}.
 \end{gathered}
\end{equation}
The division defining $u$ is exact by
\eqref{eq:half-pell-Q-definition}. Equation
\eqref{eq:half-pell-necessity-parity} makes $(J-1)/2$ even, so
\eqref{eq:half-pell-Q-complement} and
\eqref{eq:half-pell-Q-zero} give
\[
 u\equiv\psi_A(J)=c\pmod f,\qquad u\equiv J\pmod c.
\]
Both quotients in \eqref{eq:half-pell-new-witnesses} are therefore
integers.

They are strictly positive. For $R>1$, the positive chi sequence
is strictly increasing, and its recurrence gives
\[
 \chi_R(k+1)>(2R-1)\chi_R(k)\qquad(k\ge1).
\]
Since $\chi_R(1)=R$ and $J\ge3$, we obtain
\begin{equation}\label{eq:half-pell-positive-quotients}
 \frac{\chi_R(J)}R>(2R-1)^{J-1}
   >(2A)^{J-1}>\psi_A(J)=c>J.
\end{equation}
The middle inequality uses the integer bound $R\ge c^2>A$,
which implies $2R-1>2A$; the next is the usual strict psi upper
bound at index $J\ge3$. Thus $o,j>0$, and $y_{\rm aux}>0$ too.

The Pell equation for $(\chi_R(J),\psi_R(J))$ is precisely
\eqref{eq:half-pell-standard-norm}, hence
\eqref{eq:half-pell-norm}. The definitions give both equalities
in \eqref{eq:half-pell-linear}. Every other source equation is
retained and uses none of the changed auxiliary values. All 35
supplied unknowns are therefore positive integers. This proves
necessity without asserting an identity map for the auxiliary
coordinates, and shows exactly why $J\equiv1\pmod4$ was needed
only in this direction.

\subsection{The complete arithmetic count}

The former signed block used ten operations:
\[
 \begin{aligned}
 O&=of,& v&=O-d^2,& v_1&=v+1,& L&=vv_1,\\
 K_1&=K-1,& K_2&=KK_1,\\
 J_c&=jc,& u&=J+J_c,& U_2&=u^2,& P&=K_2U_2.
 \end{aligned}
\]
The quantities $d^2,K,c,J$ already exist outside this block.
Replace it by nine operations:
\begin{equation}\label{eq:half-pell-nine-operations}
 \begin{aligned}
 J_c&=jc,& u&=J+J_c,& O&=of,\\
 u_{\rm rhs}&=c+O,& U_2&=u^2,& Y_2&=y_{\rm aux}^2,\\
 G&=U_2-Y_2,& L&=KG,& P&=1-Y_2.
 \end{aligned}
\end{equation}
The free tests are $u=u_{\rm rhs}$ and $L=P$. The old block
has six multiplications and four additions or subtractions; the new
block has five multiplications and four additions or subtractions.
Every other primitive remains unchanged. The complete count is
\[
 93-1=92=49\text{ multiplications}+43\text{ additions}.
\]
The variable $d$ and its squared register remain in the main norm;
their absence from the new auxiliary block gives no further saving.

The new intermediates are ordinary integer registers. In fact $G$
and $P$ in \eqref{eq:half-pell-nine-operations} are negative on
solutions. Since $u=J+jc>1$, Pell index one is excluded in
\eqref{eq:half-pell-classification}; thus $y_{\rm aux}>1$ and
\[
 u^2-y_{\rm aux}^2=\frac{1-y_{\rm aux}^2}{K}<0.
\]
These subtractions are represented by reversed addition equations.
No positivity condition is imposed on the intermediate gaps;
all 35 supplied unknowns, including $y_{\rm aux}$, are positive.

The checker writes the source norm with $K_{\rm source}=D(f^2-1)$
and uses the already computed $K=(ic^2)^2$ in the primitive product.
The difference of its residuals is exactly
\[
 (K-K_{\rm source})(u^2-y_{\rm aux}^2).
\]
Its first factor is the independently exact retained equation
\eqref{eq:half-pell-relaxed}, so this is an acyclic correction.
It replaces the preceding signed norm's quadratic-factor correction.
The new linear congruence check is exact, and the geometric and
second-exponent corrections remain as before.

Appendix~\ref{app:92} gives the complete table of 92 instructions
and all 23 equality tests. The checker is in
\file{../verification/}:
\begin{center}
\small\file{round35\_1980\_half\_parameter\_pell\_certificate.py}.
\end{center}
Its JSON receipt records the full primitive schedule, domains, source
residuals, and remaining corrections. Those arithmetic checks and the
positive-witness proof above establish Theorem~\ref{thm:best92}.
The number 92 is an upper bound; no optimality assertion or new
compiled Lean theorem is claimed.

% Full helper-compiler proof of the 91-operation product-bound system.
\section{A positive product bound and two copied inputs: 91 operations}
\label{sec:product-bound}

The coefficient polynomial can be encoded as the difference of the two
already supplied codes. This removes the subtraction $q-C^2$ from the
product used in the third binary mask. Every negative coefficient then
creates an additional tested position. Two encoded copies of the input,
and positive complementary terms, make those extra tests exact without
adding arithmetic operations to the final system.

\begin{theorem}\label{thm:best91}
Every recursively enumerable subset of the positive integers has an
admissible fixed index $(V,H,T_L)$ for a system with 34 positive unknowns
and 22 equations admitting a straight-line certificate of 91 operations:
49 multiplications and 42 additions. The supplied inputs are
$(x,V,H,T_L)$, with three fixed components besides the positive queried
input $x$. Fixed numerals and equality tests are free.
\end{theorem}

\subsection{The changed source and its preliminary bounds}

Start from the system of Theorem~\ref{thm:best92}, retaining its
half-parameter Pell block and factored packed mask. Replace the coding
bound and product by
\begin{equation}\label{eq:product-bound-source}
 \ell+\sigma+\alpha=q,\qquad
 \sigma=(e-\ell)C^2,\qquad C=x+g.
\end{equation}
All supplied unknowns in these equations are positive integers. The
previous equations were
\[
 \ell+e+\alpha=q,\qquad
 \Omega=\lambda-e>0,\qquad
 \sigma=\Omega(q-C^2).
\]
The new system does not supply $\Omega$. Indeed, $C\ge2$ and
$\sigma>0$ imply $e-\ell>0$. If desired, a supplied positive
$\Omega=e-\ell$ and its free equality can be restored uniquely.
Their elimination therefore preserves the positive solution set.

Write $\Omega=e-\ell$ only as notation for this computed positive
integer. Equations~\eqref{eq:product-bound-source} give, before any
Pell or coefficient decoding,
\begin{equation}\label{eq:product-bound-ranges}
 \ell<q,\qquad \sigma<q,\qquad
 C^2\le\sigma<q,\qquad
 e=\ell+\Omega\le\ell+\sigma<q.
\end{equation}
The remaining definitions are
\begin{equation}\label{eq:product-bound-packing}
\begin{aligned}
 &H=H_0-3,\quad B=H+b+4=H_0+1+b,\quad
   \theta=B-4,\quad b=x+\beta,\\
 &\lambda(B-1)=q^2-1,\quad
   S_2=\ell+eq=V+t\theta,\quad n=q^8,\\
 &S=g+q^2(S_2+q^2\sigma),\\
 &T_+=q^2(1+\theta\lambda)+\ell(\theta q^4-b),\\
 &r=S(n^2-n)+T_+(n^2-1).
\end{aligned}
\end{equation}
Here $H_0$ is a sufficiently large fixed power of two, specified below.
In particular $S_2<q^2$ and $g<C<q$. Geometry gives $B\le q^2$;
the fixed threshold gives $3L\le B\le n$ and $\theta\lambda>b$.
The estimates used in Section~\ref{sec:affine-radix} consequently give
\begin{equation}\label{eq:product-bound-pre-pell}
 0<S,T_+<q^7<n,\qquad n^2-1\le r<2n^3.
\end{equation}
For example,
\[
 T_+>q^2\theta\lambda-b\ell>bq^2-bq>0,
 \qquad
 T_+<q^4(1+\theta\ell)<Bq^5\le q^7.
\]
The upper estimate uses $1+\theta\lambda<q^2$ and
$1+\theta\ell<Bq$. The bound on $S$ is stronger than its predecessor
because $\sigma<q$.

These are the same independent hypotheses used by the half-parameter
Pell proof in Section~\ref{sec:half-parameter-pell}. That proof now
applies in its stated order, before any coefficient test. It yields
the required central-binomial divisibility, $U=4^{2r+1}$, $q=B^L$,
and that $B,q,n$ are powers of two. Its sufficiency proof does not
assume even $r$. No discarded formula for $\sigma$ enters that Pell
argument.

\subsection{Nonnegative coordinates and two helpers}

Compile the represented finite Diophantine system into a circuit over
nonnegative integers, using the fresh-copy conventions of
Section~\ref{sec:affine-radix}. The actual positive input $x$ has
weight zero and is not split. Keep the homogeneous unit $\delta$ as
one physical coordinate. Every other logical value is a sum of three
distinct nonnegative physical coordinates. Different logical values
have disjoint physical groups.

For a nonnegative value $z$, the fixed forbidden-bit decomposition is
\begin{equation}\label{eq:product-bound-split}
 (z,0,0)\quad\hbox{if its $H_0$ bit is zero},\qquad
 (z-H_0,H_0/2,H_0/2)\quad\hbox{otherwise}.
\end{equation}
All pieces have that bit zero and are nonnegative. They can be fixed
before choosing an arbitrarily large power-of-two radix $B$.

Introduce two reserved logical helpers $V_0,V_1$, each the sum of its
own three physical coordinates. Every ordinary circuit row uses
$V_0$ in place of the input and mentions neither $x$ nor $V_1$.
Every monomial of such a row therefore has two non-input physical
factors. Fresh copies ensure that each nonzero physical square
coefficient is $\pm1$ and each nonzero physical cross coefficient is
$\pm2$. Dividing by the multiplicity in $C(T)^2$ gives coefficients
in $\{-1,1\}$.

The first four main rows are helper rows. For each $i=0,1$, use a
seed with base polynomial $-x^2$ and a reverse row with base
polynomial $x^2-V_i^2$. The seed will be augmented so that its actual
tested polynomial is $V_i^2-x^2$. The reverse row keeps its base
polynomial. This is the only intentional change to a main row caused
by the augmentation.

Include both signs of every ordinary zero row. In addition to the
circuit rows, include both signs of
\begin{equation}\label{eq:product-bound-normalization}
\begin{gathered}
 X^2-V_0^2,\qquad Y^2-V_0^2,\qquad Z^2-V_0^2,
 \qquad \delta'^2-\delta^2,\\
 2V_0X-2\delta\delta'-2u\delta.
\end{gathered}
\end{equation}
All these logical groups are distinct. Finish with three unpaired
$\delta^2$ rows. The last two receive the fivefold and sevenfold
padding described below. They have no negative coefficients.

\subsection{Weights and the positive compensation terms}

Let $m$ count the true non-input physical coordinates, including
$\delta$. Give them distinct weights
\begin{equation}\label{eq:product-bound-weights}
 v_i=8\cdot7^i\quad(0\le i<m),\qquad M=8\cdot7^{m-1}.
\end{equation}
After division by eight, a sum of at most six such weights has
base-seven digits below seven. Equality of any two such sums is
therefore equality of their multisets. In particular, a sum with
formal degree four cannot equal a sum with formal degree at most two.
The actual input has weight zero; its seed exception is treated
separately.

Let $s$ be the number of main rows, including the four helper rows
and the three final unit rows. Put
\begin{equation}\label{eq:product-bound-layout}
\begin{gathered}
 d_0=12M+8,\qquad
 t_j=(s+2)d_0+8M+jd_0\quad(0\le j<s),\\
 t_*=t_{s-1}=(2s+1)d_0+8M,\qquad K=t_*+3.
\end{gathered}
\end{equation}
For every divided base coefficient $a$ at a physical monomial of
weight $w$ in row $j$, put $aT^{t_j-w}$ in a complementary polynomial.

For each negative divided coefficient at $r=t_j-w$, choose a helper:
$V_1$ for an ordinary row, its own $V_i$ for seed $i$, and the other
$V_{1-i}$ for reverse row $i$. If its physical group is $P$, add
the six positive terms
\begin{equation}\label{eq:product-bound-augmentation}
 \sum_{h\le h',\ h,h'\in P}T^{r-v_h-v_{h'}}.
\end{equation}
The cross multiplicities in $C(T)^2$ recover the square of the sum
of the three helper coordinates. Mark every main target and every
negative base exponent as a tested start. Denote their set by
$\mathcal R$. In a seed, the negative exponent and main target are
the same start.

There are four relevant coefficient identities. First, an ordinary
or reverse augmentation has complementary weight
$w_{\rm negative}+w_{\rm helper}$ of degree four. It cannot affect
the main target, because a term of $C^2$ has degree at most two.
Second, at a negative position $r=t_j-w_{\rm negative}$, the negative
base term contributes $-x^2$. Another non-input base monomial could
contribute only if
\[
 w_{\rm base}=w_{\rm negative}+w_C.
\]
Degree and multiset uniqueness force $w_C=0$ and the identical base
monomial. The positive input square in a reverse row has weight zero
and cannot satisfy this equality.

Third, an augmentation contributes at that position only if
\begin{equation}\label{eq:product-bound-multiset}
 w_{{\rm negative},1}+w_{\rm helper}=w_{{\rm negative},2}+w_C.
\end{equation}
Every factor of a negative monomial in this row is disjoint from the
helper group. Both helper factors in~\eqref{eq:product-bound-multiset}
must consequently belong to the monomial on the right. Uniqueness
forces the identical negative pair and the identical helper pair.
Thus the extra tested polynomial is exactly
\begin{equation}\label{eq:product-bound-extra-target}
 V_{\rm helper}^2-x^2.
\end{equation}
Fourth, a seed has only $-x^2$ at its main target and its degree-two
helper complements. Its actual target is exactly $V_i^2-x^2$.

These arguments also exclude coefficient accumulation. Within a
non-seed augmentation, the partition into a negative pair and a
helper pair is unique by their disjoint groups. Seed helper pairs
have distinct weights. Base and augmented complementary exponents
have different degrees. Different row bands are separated by
$d_0>12M$. Thus all coefficients constructed so far remain $\pm1$.

\subsection{Resets, unit padding, and the two code polynomials}

Add a positive reset at every distinct tested start:
\begin{equation}\label{eq:product-bound-resets}
 D_{\rm reset}(T)=\sum_{r\in\mathcal R}T^{r-4}.
\end{equation}
Write $t_5=t_{s-2}$ and $t_7=t_{s-1}$ for the last two unit targets.
If $P(X)$ denotes the three physical indices of $X$, add
\begin{equation}\label{eq:product-bound-pads}
\begin{aligned}
 D_5(T)&=T^{t_5-1}
   +\sum_{h\in P(X)\cup P(Y)}T^{t_5-1-v_h},\\
 D_7(T)&=T^{t_7-1}
   +\sum_{h\in P(X)\cup P(Y)\cup P(Z)}T^{t_7-1-v_h}.
\end{aligned}
\end{equation}
Let $D(T)$ be the full complementary polynomial. Its base and
augmentation exponents are zero modulo eight, its resets are four
modulo eight, and its padding exponents are seven modulo eight.
These parts do not collide. Resets have distinct exponents, and the
padding terms are distinct within their disjoint final bands.
Consequently every nonzero coefficient of $D$ is still $\pm1$.

Define
\begin{equation}\label{eq:product-bound-codes}
 \ell_0(T)=\sum_{i=0}^{m-1}T^{v_i}
        +\sum_{r\in\mathcal R}(T^r+T^{r+1}+T^{r+2}),
 \qquad e_0(T)=\ell_0(T)+D(T).
\end{equation}
All indicator positions are distinct. Every negative coefficient of
$D$ occurs at a tested start and is canceled in $e_0$ by an indicator
digit one. Hence the digits of $\ell_0$ are zero or one and those of
$e_0$ are zero, one, or two. Both polynomials have zero constant
coefficient. The highest nonzero $D$ coefficient is the positive
padding term at $t_*-1$. Its leading $+1$ and remaining digits
$\pm1$ imply $D(B)>0$ for every integer $B\ge2$.

The indicator also permits dummy coordinate digits at the tested
positions. None can affect the low coefficient tests: the least
$D$ exponent is at least $t_0-4M$, the least dummy weight is at
least $t_0-2M$, and
\begin{equation}\label{eq:product-bound-dummies}
 2t_0-6M>t_*+2.
\end{equation}
Every true-coordinate weight lies below the least $D$ exponent.
Its raw coefficient and incoming carry are therefore zero.

Let $c_*$ count the input, all true physical coordinates, and all
permitted dummies, and put $D_1=\sum_h|[T^h]D(T)|$. Choose a power
of two $L>3K+2$ and then a power of two $H_0$ such that
\begin{equation}\label{eq:product-bound-fixed-threshold}
 H_0>\max\bigl(
   2\cdot4^{2L+1},\ 128\max(1,D_1)c_*^2,\ 3L,\ 1024\bigr).
\end{equation}
The supplied fixed index is
\begin{equation}\label{eq:product-bound-index}
 H=H_0-3,\qquad
 V=\ell_0(4)+4^Le_0(4),\qquad T_L=\psi_4(L).
\end{equation}
These choices precede the queried input and its witnesses.

\subsection{Canonical decoding after the Pell conclusions}

The bounds~\eqref{eq:product-bound-ranges} and the conclusion $q=B^L$
give $b\ell<Bq<q^2$. Thus
\begin{equation}\label{eq:product-bound-mask-blocks}
 T_+-1=(q^2-1-b\ell)+q^2\theta\lambda+q^4\theta\ell
\end{equation}
is the genuine block decomposition: its first two blocks are below
$q^2$, and the last is below $q^4$. The packed binomial condition
therefore gives the three independent binary masks for $g,S_2,\sigma$.

The middle mask is $\theta\lambda=(B-4)\sum_{j<2L}B^j$.
It requires every base-$B$ digit of $S_2$ to be at most three.
Write its digit polynomial as $F(T)$, with degree below $2L$.
Since $S_2\equiv V\pmod{B-4}$, one has
\[
 F(4)\equiv V\pmod{B-4}.
\]
Both sides are bounded by $4^{2L}-1$, while
\eqref{eq:product-bound-fixed-threshold} makes $B-4$ larger than
their possible difference. Hence $F(4)=V$. Uniqueness of the
base-four digit expansion gives
\[
 F(T)=\ell_0(T)+T^Le_0(T).
\]
Because $\ell,e<q$, the low and high $q$-blocks are their actual
values. Consequently $\ell=\ell_0(B)$ and $e=e_0(B)$.

The first mask now has digit $H_0=B-1-b$ at an indicator position
and digit $B-1$ elsewhere. It forces all non-indicator digits of
$g$ to vanish and clears the $H_0$ bit of each permitted digit.
The indicator has zero unit digit, so $g$ has zero unit digit.
Since $x<B$, adding $x$ introduces no carry, and $C=x+g$ is exactly
the physical coordinate polynomial evaluated at $B$.

Finally $\theta\ell_0(B)$ has digit $B-4$ at every indicator
position and zero elsewhere. The third mask says that each
corresponding digit of
\[
 \sigma=(e-\ell)C^2=D(B)C(B)^2
\]
is at most three. This includes every extra negative-coefficient
position and all three digits of each tested window.

\subsection{Signed carries and exact zero equations}

Every decoded coordinate digit is less than $B$, and $B\ge2H_0$.
Thus every raw coefficient of $D(T)C(T)^2$ has absolute value at most
\begin{equation}\label{eq:product-bound-cubic}
 D_1C(1)^2<D_1c_*^2B^2<B^3/256.
\end{equation}
The signed carry recurrence keeps every incoming carry in absolute
value below $B^2/128$. By~\eqref{eq:product-bound-dummies}, the only
raw residue classes through $t_*+2$ are zero, four, and seven
modulo eight.

At a tested start $r$, the positions $r-6,r-5$ are empty. Two
successive divisions by $B$ reduce the carry entering $r-4$ to
$-1$ or zero. Its raw reset polynomial is nonnegative and contains
$x^2$: the reset belonging to $r$ contributes $x^2$, and all other
reset contributions are positive. The other parts of $D$ have
different residue classes. Since $x\ge1$, the outgoing carry is
nonnegative and below $B^2$. The two empty positions $r-3,r-2$
therefore reduce it to zero.

At every start other than $t_5,t_7$, the position $r-1$ is also
empty. The following positions $r+1,r+2$ are always empty. Such a
test therefore sees its exact raw value modulo $B^3$. A negative
nonzero value satisfying~\eqref{eq:product-bound-cubic} has a top
digit greater than three and fails. The seed and reverse tests
thus give $V_i^2\ge x^2$ and $V_i^2\le x^2$, respectively. Since
these values are nonnegative,
\begin{equation}\label{eq:product-bound-helpers-fixed}
 V_0=V_1=x.
\end{equation}
All extra target values~\eqref{eq:product-bound-extra-target} now
vanish. Both signs of every ordinary row must pass, so every
ordinary zero equation is exact. In particular,
\begin{equation}\label{eq:product-bound-unit-guard}
 X=Y=Z=x,\qquad \delta'=\delta,\qquad
 x^2=\delta^2+u\delta,
 \qquad 0<\delta\le x<B.
\end{equation}
This establishes the input and guard constraints before using any
unit test.

At the two remaining preceding positions, the exact raw values
from~\eqref{eq:product-bound-pads} are
$x^2+2xX+2xY$ and $x^2+2xX+2xY+2xZ$. Their incoming carries
are zero. By~\eqref{eq:product-bound-unit-guard}, the three unit
tests are therefore the residues modulo $B^3$ of
\begin{equation}\label{eq:product-bound-three-units}
 \delta^2,\qquad
 \delta^2+\lfloor5x^2/B\rfloor,\qquad
 \delta^2+\lfloor7x^2/B\rfloor.
\end{equation}
There is no padding at an extra negative start: the two padded
rows contain no negative term, and other row bands are disjoint.

\subsection{The unit is one}

For completeness, the three-unit argument can be applied directly
to~\eqref{eq:product-bound-three-units}. Since $\delta<B$, the first
test gives
\[
 \delta^2=kB+s,\qquad k,s\in\{0,1,2,3\}.
\]
A square modulo four cannot have $s=2$ or $s=3$. If $s=1$, then
$\delta^2\le3B+3<B^2/16$, so $0<\delta<B/4$. The four square
roots of one modulo a power of two $B\ge8$ are represented by
$1,B/2-1,B/2+1,B-1$. Only the first lies in this range, giving
$\delta=1$.

Otherwise a nonunit solution would have $\delta^2=kB$ with
$1\le k\le3$ and $x^2\ge kB$. Write
\[
 h_c=\lfloor cx^2/B\rfloor=m_cB+r_c
 \quad(c=5,7),\qquad 0\le r_c<B.
\]
Because $x<B$, one has $h_c<cB$. The padded tests then force
$r_5,r_7\le3$ and $m_5,m_7\le3-k\le2$. On the other hand,
\[
 -7<7h_5-5h_7<5,\qquad |7r_5-5r_7|\le21.
\]
It follows that $|B(7m_5-5m_7)|<28$. Since $B>32$,
$7m_5=5m_7$, and their bounds imply $m_5=m_7=0$. This contradicts
$h_5\ge5k\ge5$ and $h_5=r_5\le3$. Thus $\delta=1$ in every
solution. The decoded circuit consequently represents the original
system at its actual input $x$, proving sufficiency.

\subsection{Necessity, including the auxiliary parity condition}

Given a represented solution, assign both helpers the value $x$,
set $\delta=\delta'=1$, $X=Y=Z=x$, and $u=x^2-1$, and supply all
circuit values and fresh copies. Split the non-input logical values
by~\eqref{eq:product-bound-split}; keep $\delta$ unsplit and set
all dummies to zero. Every ordinary, helper, and extra target has
raw value zero, and all three unit targets have raw value one.

After these finitely many physical values are fixed, choose $B$ to
be an arbitrarily large power of two. Require $B>H_0+1+x$, every
raw coefficient in absolute value below $B/4$, and $7x^2<B/4$.
Set
\begin{equation}\label{eq:product-bound-necessary-codes}
\begin{aligned}
 &b=B-H_0-1,\quad \beta=b-x,\quad g=C(B)-x,\\
 &\ell=\ell_0(B),\quad e=e_0(B),\quad q=B^L,\quad n=q^8,\\
 &\sigma=D(B)C(B)^2,\qquad \alpha=q-\ell-\sigma.
\end{aligned}
\end{equation}
The leading coefficient of $D$ makes $\sigma>0$. The encoded unit
makes $g>0$. Since $\deg(D C^2)<3K<L$, choosing $B$ still larger
makes $\alpha>0$. All forbidden-bit tests pass, and all middle-code
digits are at most two. Resets clear every signed carry. The two
padded unit carries are zero because $7x^2<B/4$. All third-mask
windows are therefore $(0,0,0)$ or $(1,0,0)$.

Geometry supplies a positive integer $\lambda$, and
\[
 t=\frac{\ell+eq-V}{B-4}>0
\]
is an integer: the numerator is the difference at $B$ and at four
of the nonconstant polynomial $\ell_0(T)+T^Le_0(T)$, which has
nonnegative coefficients. The three masks give the required binomial
divisibility for the newly determined $r$, whose growth bounds were
proved in~\eqref{eq:product-bound-pre-pell}.

Construct all non-auxiliary Pell witnesses as in
Section~\ref{sec:affine-radix}, with $U=4^{2r+1}$, the binomial
ratio value $Y_P$, the main and first Pell coordinates, and second
index $L$ with fixed $T_L=\psi_4(L)$. The half-parameter auxiliary
construction has one additional necessity hypothesis. Here $B$
divides both $g$ and $\ell$, since their unit digits vanish. Thus
both $S$ and $T_+$ in~\eqref{eq:product-bound-packing} are even,
and so is $r$. Consequently $J=2r+1$ is one modulo four.

Write $A=a+4$ and $D_P=A^2-1$, distinguishing the Pell coefficient
from the encoding polynomial $D(T)$. The positive construction in
Section~\ref{sec:half-parameter-pell} now applies with
\begin{equation}\label{eq:product-bound-auxiliary-witnesses}
\begin{gathered}
 m=2cJ,\qquad f=\chi_A(m),\qquad
 i=D_P\psi_A(m)/c^2,\qquad R=ic^2,\\
 y_{\rm aux}=\psi_R(J),\qquad
 u_{\rm aux}=\chi_R(J)/R,\\
 o=(u_{\rm aux}-c)/f,\qquad
 j=(u_{\rm aux}-J)/c.
\end{gathered}
\end{equation}
That proof establishes integrality and strict positivity using
$J\equiv1\pmod4$, and gives exactly the retained auxiliary norm
and congruence. All changed positive coding witnesses were supplied
in~\eqref{eq:product-bound-necessary-codes}. The difference $e-\ell$
is positive but is a computed register, not another supplied unknown.
This proves necessity for all 34 positive unknowns.

\subsection{Complete arithmetic certificate}

Relative to the 92-operation schedule, replace the three instructions
\[
 t_2=\lambda-e,\qquad b_\Sigma=\ell+e,\qquad
 S_3=t_2(q-C^2)
\]
by $t_2=e-\ell$, $b_\Sigma=\ell+\sigma$, and $S_3=t_2C^2$,
and delete the subtraction $q-C^2$. The supplied $\Omega$ and its
free equality are also deleted. Exactly one arithmetic operation
disappears, giving 49 multiplications and 42 additions.

The complete checker states all 22 source residuals afresh and
compares them with every certificate equality. If
$P_{\rm old}=(\lambda-e)(q-C^2)$ and $P_{\rm new}=(e-\ell)C^2$,
the aligned predecessor residuals change only by
\begin{equation}\label{eq:product-bound-source-differences}
 \sigma-e,\qquad
 -q^4(n^2-n)(P_{\rm new}-P_{\rm old}),\qquad
 P_{\rm new}-P_{\rm old},
\end{equation}
in the bound, packed-$r$, and product equations respectively.
The old positive-factor equation is omitted; the auxiliary congruence
is unchanged. To state the inherited residual corrections precisely,
let $F_j$ denote source residual $j$ in the checker's zero-based order,
and put $u_*=2r+1+jc$. The nonzero corrections are
\[
\begin{aligned}
 \mathcal E_2&=-\lambda F_3,\\
 \mathcal E_{16}&=F_{15}(u_*^2-y_{\rm aux}^2),\\
 \mathcal E_{17}&=F_3\bigl(\kappa+
       \rho(F_3-2(A-B))\bigr).
\end{aligned}
\]
Residuals $F_3$ and $F_{15}$ are independently exact geometry and
relaxed-norm equations. These corrections are therefore triangular
and vanish without circular reasoning. The generated schedule gives
all 91 primitive instructions and 22 free equality tests.

The general compiler and positive-domain proof received two
independent reviews. A supplementary sparse regression checks exact
symbolic target and padding identities in a sample with 35 true
coordinates, 25 main rows, 128 tested starts, 418 indicator digits,
and 976 complementary coefficients. Six canonical assignments and
four wide signed-carry assignments pass. These finite checks support
the general proof and do not assert exhaustive coverage of circuits
or integer witnesses. The construction proves an upper bound; it
does not claim that 91 is optimal.

% Complete binary helper compiler and matching base-two Pell proof.
\section{Binary coefficient codes and a matching Pell shift: 90 operations}
\label{sec:binary-product}

Changing the helper reverse rows makes every positive complementary
coefficient lie outside the indicator. Both fixed codes then have binary
digits. The geometric identity uses $\lambda$ directly in place of
$3\lambda$, saving one multiplication. A simultaneous change of the
main Pell base from $a+4$ to $a+2$ preserves all shared expressions.

\begin{theorem}\label{thm:best90}
Every recursively enumerable subset of the positive integers has an
admissible fixed index $(V,H,T_L)$ for a system with 34 positive unknowns
and 22 equations admitting a straight-line certificate of 90 operations:
48 multiplications and 42 additions. The supplied inputs are
$(x,V,H,T_L)$, with three fixed components besides the positive queried
input $x$. Fixed numerals and equality tests are free.
\end{theorem}

\subsection{The binary helper compiler}

Retain the physical representation of Section~\ref{sec:product-bound}.
The actual positive input $x$ has weight zero and is not split; the
homogeneous unit $\delta$ is one physical coordinate. Every other
logical value is a sum of three distinct nonnegative physical
coordinates, with disjoint groups for distinct logical values. The
fixed forbidden-bit decomposition of a value $z$ is $(z,0,0)$ if its
$H_0$ bit is zero, and $(z-H_0,H_0/2,H_0/2)$ otherwise.

Use the same two reserved helper triples $V_0,V_1$. Every ordinary
circuit row uses $V_0$ as input and mentions neither $x$ nor $V_1$.
Fresh square-difference copies handle repeated operands and repeated
unit factors. Expanded physical square coefficients are $\pm1$ and
cross coefficients are $\pm2$, so their divided coefficients are
$\pm1$.

For each $i=0,1$, retain the seed base row $-x^2$, augmented by its
own helper so that its tested value is $V_i^2-x^2$. Replace the
reverse row by
\begin{equation}\label{eq:binary-product-reverse}
 2x\delta-2V_i\delta.
\end{equation}
Its positive monomial has one non-input physical factor, while all
its negative monomials have two. Its augmenting helper is the other
$V_{1-i}$. The implication $V_i=x$ will be proved only after
$\delta>0$ has been recovered.

Include both signs of each ordinary zero equation, including
\begin{equation}\label{eq:binary-product-normalization}
\begin{gathered}
 X^2-V_0^2,\quad Y^2-V_0^2,\quad Z^2-V_0^2,\quad
 \delta'^2-\delta^2,\\
 2V_0X-2\delta\delta'-2u\delta.
\end{gathered}
\end{equation}
The logical groups are distinct. Keeping the extra copies $Y,Z$ is
convenient and has no effect on the final operation count. Finish with
two unpaired $\delta^2$ rows. Only the last is padded, by a preceding
raw coefficient $2xX$.

For $m$ true non-input physical coordinates, use
\[
 v_i=8\cdot7^i\quad(0\le i<m),\qquad M=8\cdot7^{m-1}.
\]
Sums of at most six such weights have unique multisets. If $s$ counts
the main rows, including the four helper rows and two unit rows, put
\begin{equation}\label{eq:binary-product-layout}
\begin{gathered}
 d_0=12M+8,\qquad
 t_j=(s+2)d_0+8M+jd_0\quad(0\le j<s),\\
 t_*=t_{s-1}=(2s+1)d_0+8M,\qquad K=t_*+3.
\end{gathered}
\end{equation}
Place each divided base coefficient $a$ of monomial weight $w$ at
$t_j-w$. For a negative coefficient at $r=t_j-w$, add the six
positive helper complements
\begin{equation}\label{eq:binary-product-complements}
 \sum_{h\le h',\ h,h'\in P}T^{r-v_h-v_{h'}},
\end{equation}
where $P$ is the own helper group for a seed, the other group for a
reverse row, and the $V_1$ group for an ordinary row. Let $\mathcal R$
contain every main target and every negative base exponent. The seed
target coincides with its negative exponent. Add one reset $T^{r-4}$
for each distinct $r\in\mathcal R$, and add the final padding
\begin{equation}\label{eq:binary-product-padding}
 D_{\rm pad}(T)=\sum_{h\in P(X)}T^{t_*-1-v_h}.
\end{equation}
Call the resulting polynomial $D(T)$.

The coefficient-isolation proof from Section~\ref{sec:product-bound}
has one new degree case. Helper complements in non-seed rows have
degree four, so cannot affect a main target after multiplication by
$C(T)^2$. At a negative position, another base monomial could
contribute only if
$w_{\rm base}=w_{\rm negative}+w_C$. Ordinary base monomials have
degree two; the new positive reverse monomial has degree one and
cannot satisfy this equality with a degree-two negative monomial.
At an extra negative-position target, an augmentation contributes
only if
\[
 w_{{\rm negative},1}+w_{\rm helper}
   =w_{{\rm negative},2}+w_C.
\]
The negative factors and helper group are disjoint. Multiset uniqueness
forces the same negative pair and helper pair, giving exactly
$V_{\rm helper}^2-x^2$. A seed separately gives $V_i^2-x^2$ at its
coincident target. Reverse main targets remain
\eqref{eq:binary-product-reverse}.

The same uniqueness and row spacing prevent coefficient accumulation.
Base and augmentation terms have residue zero modulo eight; resets
have residue four and padding has residue seven. All nonzero $D$
coefficients therefore remain $\pm1$. Define
\begin{equation}\label{eq:binary-product-codes}
 \ell_0(T)=\sum_{i=0}^{m-1}T^{v_i}
     +\sum_{r\in\mathcal R}(T^r+T^{r+1}+T^{r+2}),
 \qquad e_0(T)=\ell_0(T)+D(T).
\end{equation}

Every negative $D$ coefficient is at an indicator start. Every
positive coefficient is now outside the indicator. Ordinary positive
base complements are distinct degree-two monomials; reverse positive
complements have degree one; non-seed helper complements have degree
four; and seed helper complements differ from their sole main start.
Reset and padding residues are outside the three tested residues.
Other row bands and the low true-coordinate region are disjoint.
In particular, the previous positive $x^2$ term at the reverse main
target has disappeared. Thus both $\ell_0$ and $e_0$ have only
digits zero and one.

Both code polynomials have zero constant coefficient and support below
$K$. The leading $D$ term is the last reset $T^{t_*-4}$ with coefficient
$+1$; all other terms in the final positive-only row have smaller degree.
Consequently $D(B)>0$ for every integer $B\ge2$.

Permit dummy coordinate digits at the indicator windows. Their least
weight is at least $t_0-2M$, while the least $D$ exponent is at least
$t_0-4M$. Their sum exceeds $t_*+2$, so no dummy affects any tested
coefficient or reset. The least $D$ exponent also exceeds $M$, so
every true-coordinate indicator position has raw coefficient and
incoming carry zero.

\subsection{Fixed constants and bounds before Pell decoding}

Let $c_*$ count the input, true coordinates, and permitted dummies,
and let $D_1=\sum_h|[T^h]D(T)|$. Choose a power of two $L>3K+2$,
and then a power of two $H_0$ such that
\begin{equation}\label{eq:binary-product-threshold}
 H_0>\max\bigl(2\cdot2^{2L+1},
       128\max(1,D_1)c_*^2,\ 3L,\ 1024\bigr).
\end{equation}
The fixed index is
\begin{equation}\label{eq:binary-product-index}
 H=H_0-1,\qquad V=\ell_0(2)+2^Le_0(2),\qquad
 T_L=\psi_2(L).
\end{equation}
It is chosen before the queried input.

The source retains the positive product bound and uses the new alphabet:
\begin{equation}\label{eq:binary-product-source}
\begin{gathered}
 C=x+g,\qquad \sigma=(e-\ell)C^2,\qquad
 \ell+\sigma+\alpha=q,\\
 b=x+\beta,\qquad \theta=H+b=B-2,\qquad
 B=H_0+1+b,\qquad \lambda(B-1)=q^2-1.
\end{gathered}
\end{equation}
There is no supplied $\Omega$. Positivity gives $C\ge2$, $e-\ell\ge1$,
and therefore
\begin{equation}\label{eq:binary-product-ranges}
 \ell<q,\qquad C^2\le\sigma<q,\qquad
 e\le\ell+\sigma<q,\qquad g<C<q.
\end{equation}
Set
\begin{equation}\label{eq:binary-product-packing}
\begin{aligned}
 &S_2=\ell+eq=V+t\theta,\qquad n=q^8,\\
 &S=g+q^2(S_2+q^2\sigma),\\
 &T_+=q^2(1+\theta\lambda)+\ell(\theta q^4-b),\\
 &r=S(n^2-n)+T_+(n^2-1).
\end{aligned}
\end{equation}
Geometry gives $B\le q^2$, $\theta\lambda>b$, and
$1+\theta\lambda=q^2-\lambda<q^2$. Hence
\[
 T_+>bq^2-bq>0,\qquad
 T_+<q^4(1+\theta\ell)<Bq^5\le q^7,
 \qquad 0<S<q^7<n.
\]
Together with the fixed threshold, these imply
\begin{equation}\label{eq:binary-product-preliminary}
 n\ge64,\qquad n\le n^2-1\le r<2n^3,\qquad
 2\le L,\qquad 3L\le B\le n.
\end{equation}
No exponent or coefficient decoding has been used.

\subsection{The main Pell base is now \texorpdfstring{$a+2$}{a+2}}

Distinguish the Pell discriminant from the coefficient polynomial:
\begin{equation}\label{eq:binary-product-pell-definitions}
\begin{gathered}
 U=wn^2,\quad Y_P=s_Pn^2,\quad E=UY_P,\quad Q=UY_P^2,\quad
 P=2Q+1,\\
 a=Y_P(U+1),\qquad A=a+2,\qquad D_P=A^2-1,\qquad J=2r+1.
\end{gathered}
\end{equation}
Retain the first norm, positive interval, and index congruence:
\begin{equation}\label{eq:binary-product-first-pell}
\begin{gathered}
 \tau(\tau+1)=(E^2+U)(Y_Pk)^2,\\
 c=Y_Pk+\eta,\qquad k=\eta+\zeta,\qquad
 k=r+1+hE.
\end{gathered}
\end{equation}
The main norm, relaxed auxiliary norm, and half-parameter block are
\begin{equation}\label{eq:binary-product-auxiliary}
\begin{gathered}
 d^2=1+D_Pc^2,\qquad R=ic^2,\qquad
 K_P=R^2=D_P(f^2-1),\\
 u_*=J+jc=c+of,\qquad
 K_P(u_*^2-y_{\rm aux}^2)=1-y_{\rm aux}^2.
\end{gathered}
\end{equation}
Replace the first exponent congruence by
\begin{equation}\label{eq:binary-product-first-exponent}
 d=U+ac+\gamma(4a+3).
\end{equation}
The second norm, positive gap, fixed index, and exponent congruence are
\begin{equation}\label{eq:binary-product-second-pell}
\begin{gathered}
 c=\kappa+\phi,\qquad \mu^2=1+D_P\kappa^2,\qquad
 \kappa=T_L+\Delta a,\\
 \mu=q+\kappa(A-B)+\rho\bigl(D_P-(A-B)^2\bigr).
\end{gathered}
\end{equation}
Because $A=a+2$ and $B=\theta+2$, the shared difference is still
$A-B=a-\theta$.

\subsection{Exact indices before either exponential}

The preliminary bounds give
\[
 U,Y_P\ge n^2\ge4096,\qquad
 E\ge n^4>r+1,\qquad a>n^4>J.
\]
Classification of the first norm yields $k=\psi_P(t)$ at a positive
index. Its congruence modulo $E$ gives
$t=r+1+vE$ with $v\ge0$. Thus $t\ge r+1$. Since $P>A$ and
$c>Y_Pk>k$, comparison with the main norm
$c=\psi_A(p)$, $d=\chi_A(p)$ gives
\[
 p\ge t+1\ge r+2\ge66,\qquad c>AD_P^2,\qquad c>J.
\]
These are the generic hypotheses of the relaxed auxiliary and
half-parameter proofs above. They use only $A>1$, not its former
shift by four or evenness of $r$. Applying them to
\eqref{eq:binary-product-auxiliary} recovers
\[
 p=J,\qquad c=\psi_A(J),\qquad d=\chi_A(J).
\]

If $v\ge1$, ordinary Pell growth and $E>r$ give
\[
 \frac ck
 \le\frac{(2A)^{2r}}{(2P-1)^{r+E}}
 \le\left(\frac{2A}{2P-1}\right)^{2r}<\frac12,
\]
because the changed exact difference is
\begin{equation}\label{eq:binary-product-index-exclusion}
 (2P-1)-4A=4Y_P\bigl(U(Y_P-1)-1\bigr)-7>0.
\end{equation}
This contradicts $c/k>Y_P$. Consequently $k=\psi_P(r+1)$, still
without either exponential identity or a decoded circuit equation.

\subsection{Ratio growth, exponent recovery, and exact rounding}

Put $\xi=(U+1)^{2r}/U^r$. Standard lower and upper Pell bounds give
\begin{equation}\label{eq:binary-product-ratio}
\begin{aligned}
 \frac ck
 &\ge \xi\,
   \frac{(1+3/(2a))^{2r}}{(1+1/(2Q))^r}>\xi,\\
 \frac ck&<\xi(1+2/a)^{2r}.
\end{aligned}
\end{equation}
The strict lower bound uses
$(1+3/(2a))^2>1+3/a>1+1/(2Q)$, since $6Q>a$.
Also
$4r/a<8/n\le1/8<1/2$. The elementary estimate
$(1+t)^m\le(1-mt)^{-1}<1+2mt$ therefore gives
$c/k<\xi(1+8r/a)$.

The positive interval and the lower ratio bound imply $\xi<Y_P+1$.
Since $\xi>U^r$ and $Y_P$ is an integer,
\begin{equation}\label{eq:binary-product-ratio-growth}
 Y_P\ge U^r,\qquad a=Y_P(U+1)>U^{r+1}.
\end{equation}
Using $\xi<Y_P+1<2Y_P$ in the upper estimate gives
\begin{equation}\label{eq:binary-product-ratio-error}
 0<c/k-\xi<16r/(U+1).
\end{equation}
The right side is not yet asserted small.

The chi-congruence exponent criterion used above applies to
\eqref{eq:binary-product-first-exponent} with base two, since
$A-2=a$ and $4A-5=4a+3$. Its growth assumptions hold before its use:
\[
 2^{3J}=8\cdot64^r<U^{r+1}<a<A,\qquad
 U^3\le U^r<a<A.
\]
It follows that $U=2^J$. Since $U=wn^2$ and $n=q^8$, both $n$ and
$q$ are powers of two.

For the second index, the norm and gap in
\eqref{eq:binary-product-second-pell} give
$\kappa=\psi_A(t_0)$, $\mu=\chi_A(t_0)$ with $0<t_0<J$. Moreover,
$0<L<J$ follows from $3L\le B\le n\le r$. Reduction of the fixed
index congruence modulo $a$, using $A\equiv2\pmod a$, gives
$\psi_2(t_0)\equiv\psi_2(L)\pmod a$. Both quantities are strictly
smaller than $a$: for $0<t<J$,
\[
 0<\psi_2(t)\le4^{t-1}<2^{3J}<a.
\]
Thus they are equal as integers, and strict growth implies $t_0=L$.
The second exponent criterion has the required bounds
\[
 B^{3L}\le B^B\le n^n\le U^r<a<A,\qquad
 q^3\le n^n\le U^r<a<A.
\]
Its congruence in~\eqref{eq:binary-product-second-pell} therefore
gives $q=B^L$. As $q$ is a power of two, so is $B$.

Now $U=2^J=2\cdot4^r>32r$, so
\eqref{eq:binary-product-ratio-error} is less than $1/2$. Split the
binomial expansion as $\xi=F+\varepsilon$, where
\begin{equation}\label{eq:binary-product-binomial-split}
 F=\sum_{j=r}^{2r}\binom{2r}{j}U^{j-r},
 \qquad
 \varepsilon=\sum_{j=1}^{r}\frac{\binom{2r}{r-j}}{U^j}.
\end{equation}
Exact symmetry gives the sharper bound
\begin{equation}\label{eq:binary-product-tail}
 0<\varepsilon
 \le\frac{4^r-\binom{2r}{r}}{2U}<\frac14.
\end{equation}
Hence $F$ is the integer part of $\xi$ and
$F\equiv\binom{2r}{r}\pmod U$. The interval and ratio bounds yield
\[
 F<\xi<c/k<Y_P+1,\qquad
 Y_P<c/k<\xi+\tfrac12<F+\tfrac34.
\]
Integrality forces $Y_P=F$. Since $n^2$ divides $Y_P$ and $U$, it
divides $\binom{2r}{r}$. This establishes every hypothesis for the
packed binary masks without relying on a decoded coefficient row.

\subsection{Canonical codes, carry clearing, and the unit}

From $b<B<q$ and $\ell<q$, one has $b\ell<q^2$. Thus the exact
normalized mask is
\[
 T_+-1=(q^2-1-b\ell)+q^2\theta\lambda+q^4\theta\ell,
\]
with block widths $q^2,q^2,q^4$. The three independent masks follow
from central-binomial divisibility. The middle mask is
$(B-2)\sum_{j<2L}B^j$, so all digits of $S_2$ are binary. If $F_0(T)$
is its digit polynomial, the fixed congruence gives
$F_0(2)\equiv V\pmod{B-2}$. Both nonnegative quantities are below
$2^{2L}$, and~\eqref{eq:binary-product-threshold} makes their
difference smaller than $B-2$. Binary uniqueness therefore gives
\[
 F_0(T)=\ell_0(T)+T^Le_0(T),\qquad
 \ell=\ell_0(B),\qquad e=e_0(B).
\]
There is no ambiguity between the two codes because $\ell,e<q$.

The first mask has digit $H_0=B-1-b$ at each indicator position and
$B-1$ elsewhere. It confines $g$ to permitted digits and clears
their fixed bit. Its unit digit is zero; $x<B$ then makes $C=x+g$
exactly the physical polynomial evaluated at $B$. The third mask
requires binary digits of $\sigma=D(B)C(B)^2$ at every indicator.

Since $B\ge2H_0$, all raw coefficients satisfy
\[
 |[T^h]D(T)C(T)^2|
 \le D_1C(1)^2<D_1c_*^2B^2<B^3/256.
\]
The signed carry recurrence keeps incoming carries below $B^2/128$
in absolute value. Dummy exclusion makes the low raw residues
$0,4,7$ modulo eight. At a start $r$, empty positions $r-6,r-5$
reduce the carry entering the reset at $r-4$ to $-1$ or zero.
Its nonnegative raw polynomial contains $x^2\ge1$, so its outgoing
carry is nonnegative and below $B^2$. The empty positions $r-3,r-2$
reduce it to zero. Every $r-1$ is empty except the last unit pad,
and every $r+1,r+2$ is empty.

Every ordinary, helper, and extra target therefore tests its exact
raw value modulo $B^3$. Any negative nonzero value in the proved
range has top digit greater than one and fails. Use these facts in
the following order. The seeds give $V_i\ge x>0$. Both signs of
every ordinary row force it to vanish independently of the helper
values. In particular~\eqref{eq:binary-product-normalization} gives
\[
 X=Y=Z=V_0,\qquad \delta'=\delta,\qquad
 V_0^2=\delta^2+u\delta,
\]
so $\delta>0$. The reverse target $2\delta(x-V_i)$ cannot be
negative, proving $V_i\le x$ and hence $V_0=V_1=x$. All extra
targets now vanish, and the guard gives $0<\delta\le x<B$.

The final preceding coefficient from
\eqref{eq:binary-product-padding} is exactly $2xX=2x^2$. Its
incoming carry is zero. The two unit tests are consequently
\[
 \delta^2,\qquad \delta^2+\lfloor2x^2/B\rfloor
\]
with two empty following positions. The first gives
$\delta^2=kB+s$, $k,s\in\{0,1\}$. If $s=1$, then
$\delta^2\le B+1<B^2/16$, and the four roots of one modulo a
power of two leave only $\delta=1$ in $(0,B/4)$. Otherwise a
nonunit solution has $\delta^2=B$ and $x^2\ge B$.
Writing $h=\lfloor2x^2/B\rfloor$ gives $2\le h<2B$. Binary digits
of $B+h<3B<B^2$ would force
$1+\lfloor h/B\rfloor\le1$ and then $h\le1$, a contradiction.
Thus $\delta=1$, and the exact homogeneous circuit recovers the
represented system at its actual input.

\subsection{Positive necessity at the new Pell base}

Given a represented solution, choose $V_0=V_1=x$,
$\delta=\delta'=1$, $X=Y=Z=x$, $u=x^2-1$, all circuit values and
fresh copies, and their fixed forbidden-bit splits. Set all dummy
digits to zero. Every ordinary, seed, reverse, and extra target is
zero, and both unit targets are one.

Choose an arbitrarily large power of two $B$ after these physical
values have been fixed. Require $B>H_0+1+x$, every raw coefficient
in absolute value below $B/4$, and $2x^2<B/4$. Set
\begin{equation}\label{eq:binary-product-necessary-coding}
\begin{aligned}
 &b=B-H_0-1,\quad \beta=b-x,\quad g=C(B)-x,\\
 &\ell=\ell_0(B),\quad e=e_0(B),\quad q=B^L,\quad n=q^8,\\
 &\sigma=D(B)C(B)^2,\qquad \alpha=q-\ell-\sigma.
\end{aligned}
\end{equation}
The positive leading coefficient of $D$ and the unit coordinate give
$\sigma,g>0$. The degree bound $\deg(DC^2)<3K<L$ allows $B$ large
enough that $\alpha>0$. The first-mask bits clear, the middle digits
are binary, and every third window is $(0,0,0)$ or $(1,0,0)$; the
last pad has zero carry. Geometry supplies positive $\lambda$, and
the congruence quotient $(\ell+eq-V)/(B-2)$ is positive integral by
evaluation of the nonconstant nonnegative polynomial
$\ell_0(T)+T^Le_0(T)$ at $B$ and at two.

The masks give central-binomial divisibility for the resulting $r$,
with bounds~\eqref{eq:binary-product-preliminary}. Also $B$ divides
both $g$ and $\ell$, so~\eqref{eq:binary-product-packing} makes
$S,T_+,r$ even. Put $J=2r+1$ and construct the new Pell witnesses:
\begin{equation}\label{eq:binary-product-necessary-pell}
\begin{gathered}
 U=2^J,\quad w=U/n^2,\quad
 Y_P=\lfloor(U+1)^{2r}/U^r\rfloor,\quad s_P=Y_P/n^2,\\
 a=Y_P(U+1),\quad A=a+2,\quad E=UY_P,\quad Q=UY_P^2,\quad
 P=2Q+1,\\
 c=\psi_A(J),\quad d=\chi_A(J),\quad k=\psi_P(r+1).
\end{gathered}
\end{equation}
The quotient $w$ is positive integral because $n^2$ and $U$ are
powers of two and $n^2\le r^2<2^{2r+1}$. Binomial divisibility and
\eqref{eq:binary-product-tail} make $s_P$ positive integral.
The same ratio bounds at these chosen indices give
$Y_P<\xi<c/k<Y_P+3/4$. Thus
\[
 \eta=c-kY_P>0,\qquad \zeta=k-\eta>0.
\]
Set $\tau=(\chi_P(r+1)-1)/2$ and
$h=(k-r-1)/E$. The odd parameter $P$ makes $\tau$ integral;
congruence modulo $P-1=2Q$ gives integrality of $h$, and strict
Pell growth gives their positivity.

For the relaxed and half-parameter auxiliaries choose
\[
 m=2cJ,\quad f=\chi_A(m),\quad
 i=D_P\psi_A(m)/c^2,\quad R=ic^2,
\]
then take
\[
 y_{\rm aux}=\psi_R(J),\qquad
 u_*=\chi_R(J)/R,\qquad
 o=(u_*-c)/f,\qquad j=(u_*-J)/c.
\]
The generic positive construction of
Section~\ref{sec:half-parameter-pell} applies because even $r$
gives $J\equiv1\pmod4$. It proves every required divisibility and
strict positivity at the new value of $A$.

Finally choose $\kappa=\psi_A(L)$, $\mu=\chi_A(L)$,
$\phi=c-\kappa$, and
$\Delta=(\psi_A(L)-\psi_2(L))/a$. The inequalities $0<L<J$,
$A>2$, and $L\ge2$ give $\phi,\Delta>0$; polynomial congruence
modulo $a$ makes $\Delta$ integral. Both exponent congruences have
integer quotients. Their positivity follows explicitly from
\[
 d-ac=2c-\psi_A(J-1)>c>U,
\]
and
\[
 \mu-(A-B)\kappa
   =B\kappa-\psi_A(L-1)>(B-1)\kappa>q.
\]
Here $c,\kappa\ge A$, and the established growth gives
$A>U^3,q^3$. Both moduli are positive. This supplies every positive
unknown at the changed base; it does not retain old base-four
witnesses without justification.

\subsection{Exact operation count and verification}

Geometry now reads
\[
 q^2=(1+\theta\lambda)+\lambda.
\]
The former instruction $3\lambda$ is deleted. The shared exponent
modulus and discriminant are now $4a+3$ and $a^2+(4a+3)$ in place
of $8a+15$ and $a^2+(8a+15)$, with the same count; $A-B=a-\theta$
still uses one subtraction. Thus exactly one multiplication disappears:
48 multiplications and 42 additions.

The complete checker states all 22 source residuals afresh. If
$E_0=4a+12$ is the old discriminant minus the new one, their only
changes, in zero-based source order, are
\[
\begin{array}{c|l}
2&2\lambda\\
13&\gamma E_0\\
14&c^2E_0\\
15&(f^2-1)E_0\\
16&-(f^2-1)(u_*^2-y_{\rm aux}^2)E_0\\
17&\rho E_0\\
18&-\kappa^2E_0 .
\end{array}
\]
The three inherited residual corrections remain triangular, using
independently exact geometry and relaxed-norm equations. Every
primitive, equality, source difference, and supplied input is checked.
The generated table lists all 90 instructions and 22 free equalities.

The full compiler proof and exact certificate passed independent
review. The sparse regression checks a sample with 35 true coordinates,
24 main rows, 121 tested starts, 397 indicator digits, and 912
complementary coefficients. It verifies binary code digits, every
extra target, the exact pad, dummy exclusion, six canonical assignments,
four wide signed-carry assignments, and 1,252 local unit cases.
A separate exact Pell regression checks the changed ratio, norms,
congruences, and positive witnesses through indices $r=64,66$.
These finite receipts supplement the general proofs and do not assert
exhaustive coverage or optimality.

% Narrow binary packing with a complete post-exponent range argument.
\section{Narrow binary packing: 89 operations}
\label{sec:narrow-binary-product}

The binary coefficient compiler of Section~\ref{sec:binary-product}
permits a narrower packing. A change in its lowest mask digit removes
one multiplication from the power chain. The proof must recover the
mask's size after exponent decoding and then exclude a possible carry
into the fixed-code field; neither property is assumed in advance.

\begin{theorem}\label{thm:best89}
Every recursively enumerable subset of the positive integers has the
same admissible fixed index $(V,H,T_L)$ as in Theorem~\ref{thm:best90}
for a system with 34 positive unknowns and 22 equations admitting a
straight-line certificate of 89 operations: 47 multiplications and
42 additions. The supplied inputs are $(x,V,H,T_L)$, with three fixed
components besides the positive queried input $x$. Fixed numerals and
equality tests are free.
\end{theorem}

\subsection{The changed source and preliminary bounds}

Retain the compiler, fixed index, product equation, positive bounds,
geometry, and Pell equations of Section~\ref{sec:binary-product}.
In particular,
\begin{equation}\label{eq:narrow-binary-retained}
\begin{gathered}
 C=x+g,\qquad \sigma=(e-\ell)C^2,\qquad
 \ell+\sigma+\alpha=q,\\
 b=x+\beta,\qquad \theta=H+b=B-2,\qquad
 \lambda(B-1)=q^2-1,\\
 S_2=\ell+eq=V+t\theta.
\end{gathered}
\end{equation}
All supplied unknowns are strictly positive. Replace only the packing
definitions and their index equation by
\begin{equation}\label{eq:narrow-binary-packing}
\begin{aligned}
 n&=q^4, & S&=g+q\sigma+q^2S_2,\\
 T_+&=q^2\theta\lambda+\ell(\theta q-b), &
 r&=S(n^2-n)+T_+(n^2-1).
\end{aligned}
\end{equation}
Thus the unchanged bound gives
\[
 0<g,\sigma,\ell,e<q,\qquad 0<S<q^4=n.
\]
Before any exponent is decoded, geometry gives $B\le q^2$ and
$\theta\lambda>b$. Therefore
\[
 T_+>bq^2-b\ell>0,\qquad
 T_+<q^4+\theta\ell q<2q^4=2n.
\]
Here $\theta\lambda=q^2-1-\lambda<q^2$ and
$\theta\ell q<Bq^2\le q^4$. The fixed compiler threshold then gives
\begin{equation}\label{eq:narrow-binary-preliminary}
 n\ge64,\qquad n\le n^2-1\le r<3n^3,\qquad
 2\le L,\qquad 3L\le B\le n.
\end{equation}
In particular, the narrower bound $T_+<n$ has not yet been used.

\subsection{Pell recovery with the wider preliminary range}

Use the unchanged Pell definitions and equations
\eqref{eq:binary-product-pell-definitions}--\eqref{eq:binary-product-second-pell}
at this new $n$ and $r$. Their proof extends to
\eqref{eq:narrow-binary-preliminary}. Indeed,
\[
 U,Y_P\ge n^2,\qquad
 E\ge n^4>3n^3\ge r+1,\qquad
 a>n^4>6n^3>2r+1.
\]
The first Pell index is at least $r+1$, the main index is at least
$r+2$, and the relaxed auxiliary and half-parameter arguments recover
the exact main index $J=2r+1$. The same Pell growth comparison then
recovers the first index $r+1$.

For $\xi=(U+1)^{2r}/U^r$, the lower ratio bound remains
$c/k>\xi$. The only changed estimate in the upper bound is
\[
 \frac{4r}{a}<\frac{12}{n}\le\frac3{16}<\frac12.
\]
Consequently the same elementary estimate gives
\[
 \frac ck<\xi\left(1+\frac{8r}{a}\right),\qquad
 Y_P\ge U^r,\qquad a>U^{r+1},\qquad
 0<\frac ck-\xi<\frac{16r}{U+1}.
\]
The first exponent criterion proves $U=2^{2r+1}$, independently of
mask decoding. The positive gap $c>\kappa$ and fixed index congruence
give the second Pell index $L$, since $0<L<J$ and both values
$\psi_2(L),\psi_2(t_0)$ are less than $a$. Its size hypotheses are
unchanged:
\[
 B^{3L}\le B^B\le n^n\le U^r<a,
 \qquad q^3\le n^n\le U^r<a.
\]
Thus the second exponent criterion proves $q=B^L$. Both $B$ and $q$
are powers of two. Exact binomial symmetry and the ratio estimate give
the same rounding conclusion as before,
\begin{equation}\label{eq:narrow-binary-pell-output}
 n^2\mid\binom{2r}{r}.
\end{equation}
This proof applies to arbitrary positive source solutions. It has used
neither canonical fixed codes nor nonintersection of packed fields.

\subsection{Recover the packing range after exponent decoding}

Now $q=B^L$ with $L\ge2$, so
\[
 \lambda=\frac{q^2-1}{B-1}\ge B+1>\theta.
\]
Using $\ell<q$ in the exact identity
\[
 T_+=q^4-q^2-\lambda q^2+\ell(\theta q-b)
\]
shows $T_+<q^4=n$. Together with $0<S<n$, the ordinary packed
binary lemma therefore applies to \eqref{eq:narrow-binary-pell-output}:
\begin{equation}\label{eq:narrow-binary-disjoint}
 S\mathbin{\&}(T_+-1)=0.
\end{equation}
There is no appeal to that lemma outside its proved range.

\subsection{A spill into the fixed-code mask is impossible}

Write
\[
 b\ell=k_0q+r_0,\qquad 0\le r_0<q,
 \qquad \theta\ell-1-k_0=s_0q+u_0,\qquad 0\le u_0<q.
\]
Since $b<B\le q$, we have $0\le k_0\le b-1$. Also
$\theta\ell-1-k_0\ge\theta-b=H>0$, and it is less than
$\theta q$. Hence $0\le s_0\le B-3$. The normalized mask is
\begin{equation}\label{eq:narrow-binary-normalized}
 T_+-1=(q-1-r_0)+qu_0+q^2(\theta\lambda+s_0).
\end{equation}
The last coefficient is below $q^2$, because
$\theta\lambda=q^2-1-\lambda$ and $s_0<\lambda$.
The corresponding fields of $S$ are $g,\sigma,S_2$, of widths
$q,q,q^2$. Their bounds justify separate extraction from
\eqref{eq:narrow-binary-disjoint}.

The base-$B$ digits of $\theta\lambda$ are all $B-2$. If $s_0=1$,
the unit mask digit becomes $B-1$, while the others remain $B-2$.
Every digit of $S_2$, and therefore of $\ell$, is then Boolean, so
\[
 \ell\le\frac{q-1}{B-1}<\frac q{B-1},\qquad \theta\ell<q.
\]
This contradicts $s_0=1$. If $s_0\ge2$, the first two mask digits
become $s_0-2,B-1$ and all later digits remain $B-2$. The unit
digit of $\ell$ is at most $B-1$, its next digit is zero, and its
higher digits are Boolean. Since $L\ge2$,
\[
 \ell\le B-1+\sum_{i=2}^{L-1}B^i
       =\frac{q-2B+1}{B-1}<\frac q{B-1}.
\]
Again $\theta\ell<q$, contradicting $s_0\ge2$. Thus $s_0=0$.

The unperturbed middle mask now proves that $S_2$ has a Boolean
digit polynomial $F(T)$ of degree below $2L$. Reduction of
$S_2=V+t(B-2)$ modulo $B-2$ gives $F(2)\equiv V\pmod{B-2}$.
The fixed threshold \eqref{eq:binary-product-threshold} puts both
integers below $B-2$. Binary uniqueness therefore recovers
\begin{equation}\label{eq:narrow-binary-canonical}
 \ell=\ell_0(B),\qquad e=e_0(B).
\end{equation}
The unchanged support bound $K+1<L$ implies $b\ell<q$ and
$\theta\ell<q$. Consequently $k_0=0$, and the other two masks are
exactly
\begin{equation}\label{eq:narrow-binary-individual-masks}
 g\mathbin{\&}(q-1-b\ell)=0,
 \qquad \sigma\mathbin{\&}(\theta\ell-1)=0.
\end{equation}

\subsection{The altered low mask preserves every circuit test}

The first mask in \eqref{eq:narrow-binary-individual-masks} is the
same coordinate mask as before, with its unused high block removed.
At an indicator position its digit is $H_0=B-1-b$, and elsewhere
it is $B-1$. Since $g<q$, the removed high block makes no difference.
It recovers the physical coordinate polynomial $C(B)=x+g$ exactly.

The second mask differs from $\theta\ell$ only at and below the
least indicator position $v_0=8$. Below $v_0$ its digits are $B-1$;
at $v_0$ the digit is $B-3$ instead of $B-2$. But the unchanged
compiler has $\min\operatorname{supp}D>M\ge v_0$. Independently of
the values of the coordinate digits,
\[
 \sigma=D(B)(x+g)^2\equiv0\pmod{B^{v_0+1}}.
\]
Every altered digit of $\sigma$ is therefore zero. The new test is
equivalent to the old product mask. All coefficient-isolation,
carry-reset, helper-normalization and unit arguments of
Section~\ref{sec:binary-product} apply unchanged. This proves the
soundness direction for the same fixed index and the same raw input.

\subsection{Positive converse and exact operation count}

Given a represented input, retain the canonical coding witnesses of
Section~\ref{sec:binary-product} at a sufficiently large power-of-two
radix. They satisfy all unchanged bounds and give the three masks
just proved, including the modified low product mask. Form the new
$S,T_+,n$ by \eqref{eq:narrow-binary-packing} and define its actual
new $r$. The packed lemma gives
$n^2\mid\binom{2r}{r}$ and $r\ge n^2-1$.

The canonical $g$ and $\ell$ have zero radix-unit digit. Thus $S$
and $T_+$ are even, and the new $r$ is even. Apply the positive
Pell construction of Section~\ref{sec:binary-product} afresh at this
$n,r$: set $U=2^{2r+1}$ and
$Y_P=\lfloor(U+1)^{2r}/U^r\rfloor$, use the main index $2r+1$,
the first index $r+1$, and the second index $L$. The same exact
congruences, ratio gaps and even-index half-parameter construction
make every quotient and every remaining supplied witness positive.
No old packed index or Pell witness is reused.

Both packings still cost the same number of operations. The new
power chain computes $q^2,q^4,n^2=q^8$ in three multiplications,
where the old one computed $q^2,q^4,q^8,n^2=q^{16}$ in four.
The baseline term $q^2$ formerly supplied by
$q^2(1+\theta\lambda)$ is absent, so the reordered mask incurs no
extra addition. This gives 47 multiplications and 42 additions,
proving Theorem~\ref{thm:best89}.

The complete source checker compares all 22 independently written
polynomial residuals with all 89 primitive instructions, including
the same acyclic norm and geometric corrections. The full table is
Appendix~\ref{app:89}. The companion files
\file{BINARY\_PRODUCT\_89\_PROOF.md} and
\file{BASE\_TWO\_PELL\_89\_PROOF.md} state the outer and Pell proofs
separately. Finite regression checks corroborate these arguments;
they do not replace either unbounded positive-domain direction.


\section{Theorem 5 in a precise form}\label{sec:thm5}

Let $T$ be an effectively axiomatizable theory: its language is countable
with a G\"odel numbering under which the sets of axioms and of instances of
the inference rules are recursive. Let $\mathrm{Thm}_T\subseteq\mathbb N$ be
the set of G\"odel numbers of theorems of $T$.

\begin{lemma}\label{lem:re}
$\mathrm{Thm}_T$ is recursively enumerable, and there is an admissible
coding triple $\langle z,u,y\rangle$ with
$\mathrm{Thm}_T\cap\mathbb Z_{>0}=W_{\langle z,u,y\rangle}$.
\end{lemma}

\begin{proof}
The first assertion is standard: enumerate all finite sequences of formulas,
check recursively whether each is a derivation, and output the G\"odel
number of its last formula. For the second, every r.e.\ set is Diophantine
\cite{DPR,Mat70}; by \cite[\S5]{J82} it has a representation with degree
$\delta=4$ in $\nu=58$ unknowns, and \cite[(4.1)]{J82} turns such a
representation into an admissible triple $\langle z,u,y\rangle$ with
$x\in\mathrm{Thm}_T\Leftrightarrow x\in W_{\langle z,u,y\rangle}$ for $x>0$.
\end{proof}

\begin{theorem}[Theorem 5 of \cite{J80}, certificate form]\label{thm:5}
Let $T$ be an effectively axiomatizable theory and let $\langle z,u,y\rangle$
be the triple of Lemma~\ref{lem:re}. For every proposition $P$ with G\"odel
number $x>0$ the following are equivalent:
\begin{enumerate}[nosep,label=\textup{(\alph*)}]
\item $P$ has a proof in $T$;
\item there are 32 positive integers satisfying the twenty equations of
 $\Sigma$ with parameters $x,z,u,y$;
\item there is a valid straight-line certificate of the form given in the
 Appendix~\ref{app:129}: 32 positive integers, the auxiliary integers, and 129 true
 statements of the form ``integer $+$ integer $=$ integer'' or
 ``integer $\cdot$ integer $=$ integer'' \textup{(}139 if the numerals are
 generated from $1$\textup{)}, together with 20 equality tests;
\item for the packed admissible index $(Z,V,H)$ constructed from
 $\langle z,u,y\rangle$ in Section~\ref{sec:optimization}, there is a valid
 certificate of Appendix~\ref{app:120}: 32 positive witnesses, auxiliary
 integers, 120 true addition or multiplication statements, and 20 equality
 tests.
\item for the admissible index $(Z,V,H)$ constructed from the quadratic
 representation of $\mathrm{Thm}_T$ in Section~\ref{sec:short-masks},
 there is a valid certificate of Appendix~\ref{app:114}: 33 positive
 witnesses, auxiliary integers, 114 true addition or multiplication
 statements \textup{(}123 with the fixed numerals generated from
 $1$\textup{)}, and 21 equality tests.
\item for the direct quadratic index of Section~\ref{sec:quadratic-masks},
 there is a valid certificate of Appendix~\ref{app:112}: 32 positive
 witnesses, auxiliary integers, 112 true addition or multiplication
 statements \textup{(}119 with the fixed numerals generated from
 $1$\textup{)}, and 20 equality tests.
\item for the reversed quadratic index of Section~\ref{sec:reversed-packing},
 there is a valid certificate of Appendix~\ref{app:111}: 33 positive
 witnesses, auxiliary integers, 111 true addition or multiplication
 statements \textup{(}118 with the fixed numerals generated from
 $1$\textup{)}, and 21 equality tests.
\item for the homogeneous index of Section~\ref{sec:unit-digit},
 there is a valid certificate of Appendix~\ref{app:110}: 33 positive
 witnesses, auxiliary integers, 110 true addition or multiplication
 statements \textup{(}117 with the fixed numerals generated from
 $1$\textup{)}, and 21 equality tests.
\item for the input-unit index of Section~\ref{sec:input-unit},
 there is a valid certificate of Appendix~\ref{app:109}: 33 positive
 witnesses, auxiliary integers, 109 true addition or multiplication
 statements \textup{(}116 with the fixed numerals generated from
 $1$\textup{)}, and 21 equality tests.
\item for the same fixed index, there is a valid certificate of
 Appendix~\ref{app:107}: 34 positive witnesses, auxiliary integers,
 107 true addition or multiplication statements \textup{(}114 with the
 fixed numerals generated from $1$\textup{)}, and 22 equality tests.
\item for the modular index of Section~\ref{sec:modular-support}, there
 is a valid certificate with 34 positive witnesses, auxiliary integers,
 107 true addition or multiplication statements, and 22 equality tests.
\item for the same modular index and the shifted main Pell parameter of
 Section~\ref{sec:main-base-four}, there is a valid certificate of
 Appendix~\ref{app:106}: 34 positive witnesses, auxiliary integers,
 106 true addition or multiplication statements, and 22 equality tests.
\item for the same modular index and the positive index-multiple
 condition of Section~\ref{sec:index-multiple}, there is a valid
 certificate of Appendix~\ref{app:105}: 34 positive witnesses,
 auxiliary integers, 105 true addition or multiplication statements,
 and 22 equality tests.
\item for the same modular index and the shared-product Pell system of
 Section~\ref{sec:unit-scale}, there is a valid certificate of
 Appendix~\ref{app:104}: 34 positive witnesses, auxiliary integers,
 104 true addition or multiplication statements, and 22 equality tests.
\item for the new coefficient index of Section~\ref{sec:borrow-mask},
 there is a valid certificate of Appendix~\ref{app:103}: 34 positive
 witnesses, auxiliary integers, 103 true addition or multiplication
 statements, and 22 equality tests.
\item for that coefficient index and the common exponent witness of
 Section~\ref{sec:common-witness}, there is a valid certificate of
 Appendix~\ref{app:102}: 34 positive witnesses, auxiliary integers,
 102 true addition or multiplication statements, and 22 equality tests.
\item for the fixed-four circuit index $(V,H,L)$ of
 Section~\ref{sec:fixed-four}, there is a valid certificate of
 Appendix~\ref{app:101}: 34 positive witnesses, auxiliary integers,
 101 true addition or multiplication statements, and 22 equality tests.
\item for the fixed index $(V,H,\psi_4(L))$ and doubled-index Pell
 system of Section~\ref{sec:doubled-index}, there is a valid certificate
 of Appendix~\ref{app:100-pell}: 34 positive witnesses, auxiliary
 integers, 100 true addition or multiplication statements, and 22 equality tests.
\item for the paired circuit index $(V,H,L)$ of
 Section~\ref{sec:half-center}, there is a valid certificate of
 Appendix~\ref{app:100-encoding}: 34 positive witnesses, auxiliary
 integers, 100 true addition or multiplication statements, and 22 equality tests.
\item for the combined index $(V,H,\psi_4(L))$ of
 Section~\ref{sec:composed-99}, there is a valid certificate of
 Appendix~\ref{app:99}: 34 positive witnesses, auxiliary integers,
 99 true addition or multiplication statements, and 22 equality tests.
\item for the unit-centered index $(V,H,\psi_4(L))$ of
 Section~\ref{sec:unit-center}, there is a valid certificate of
 Appendix~\ref{app:98-encoding}: 34 positive witnesses, auxiliary integers,
 98 true addition or multiplication statements, and 22 equality tests.
\item for the index of Section~\ref{sec:composed-99} with the relaxed
 auxiliary Pell norm of Section~\ref{sec:relaxed-auxiliary}, there is a
 valid certificate of Appendix~\ref{app:98-pell}: 34 positive witnesses,
 auxiliary integers, 98 true addition or multiplication statements,
 and 22 equality tests.
\item for the combined index $(V,H,\psi_4(L))$ of
 Section~\ref{sec:composed-97}, there is a valid certificate of
 Appendix~\ref{app:97}: 34 positive witnesses, auxiliary integers,
 97 true addition or multiplication statements, and 22 equality tests.
\item for the reset-padded index $(V,H,\psi_4(L))$ of
 Section~\ref{sec:linear-radix}, there is a valid certificate of
 Appendix~\ref{app:96}: 35 positive witnesses, auxiliary integers,
 96 true addition or multiplication statements, and 23 equality tests.
 Fixed numerals are free.
\item for the affine-radix index $(V,H,\psi_4(L))$ of
 Section~\ref{sec:affine-radix}, there is a valid certificate of
 Appendix~\ref{app:95}: 34 positive witnesses, auxiliary integers,
 95 true addition or multiplication statements, and 22 equality tests.
 Fixed numerals are free.
\item for the translated index $(V,H,\psi_4(L))$ of
 Section~\ref{sec:shifted-affine}, there is a valid certificate of
 Appendix~\ref{app:94}: 34 positive witnesses, auxiliary integers,
 94 true addition or multiplication statements, and 22 equality tests.
 Fixed numerals are free.
\item[(aa)] for the same translated index and the same system, there is
 a valid certificate of Appendix~\ref{app:93}: 34 positive witnesses,
 auxiliary integers, 93 true addition or multiplication statements,
 and 22 equality tests. Fixed numerals are free.
\item[(ab)] for the same translated index and the smaller auxiliary
 Pell parameter of Section~\ref{sec:half-parameter-pell}, there is
 a valid certificate of Appendix~\ref{app:92}: 35 positive witnesses,
 auxiliary integers, 92 true addition or multiplication statements,
 and 23 equality tests. Fixed numerals are free.
\item[(ac)] for the product-bound coefficient index of
 Section~\ref{sec:product-bound} and the retained smaller auxiliary
 Pell parameter, there is a valid certificate of Appendix~\ref{app:91}:
 34 positive witnesses, auxiliary integers, 91 true addition or
 multiplication statements, and 22 equality tests. Fixed numerals
 are free.
\item[(ad)] for the binary coefficient index $(V,H,\psi_2(L))$ of
 Section~\ref{sec:binary-product} and its matching Pell shift, there
 is a valid certificate of Appendix~\ref{app:90}: 34 positive witnesses,
 auxiliary integers, 90 true addition or multiplication statements,
 and 22 equality tests. Fixed numerals are free.
\item[(ae)] for the same binary coefficient index and the narrow
 packing of Section~\ref{sec:narrow-binary-product}, there is a valid
 certificate of Appendix~\ref{app:89}: 34 positive witnesses,
 auxiliary integers, 89 true addition or multiplication statements,
 and 22 equality tests. Fixed numerals are free.
\end{enumerate}
Moreover, the seventeen polynomial equations of Theorem~3 that the
certificate verifies contain 100 indicated additions, subtractions and
multiplications, which is the number reported in \cite{J80}.
\end{theorem}

\begin{proof}
(a)$\Leftrightarrow$(b) is Lemma~\ref{lem:re} and Theorem~\ref{thm:sigma};
(b)$\Leftrightarrow$(c) is Theorem~\ref{thm:129}, and
(a)$\Leftrightarrow$(d) is Theorem~\ref{thm:120}.
The quadratic-system construction of \cite[\S5]{J82} and
Corollary~\ref{thm:best114} give (a)$\Leftrightarrow$(e), and
Theorem~\ref{thm:best112} gives (a)$\Leftrightarrow$(f), and
Theorem~\ref{thm:best111} gives (a)$\Leftrightarrow$(g).
Theorems~\ref{thm:best110} and~\ref{thm:best109} give the equivalences
with (h) and (i), respectively. Theorem~\ref{thm:best107} gives the
equivalence with (j). Theorems~\ref{thm:modular107} and~\ref{thm:best106}
give the equivalences with (k) and (l), and
Theorem~\ref{thm:best105} gives the equivalence with (m), and
Theorem~\ref{thm:best104} gives the equivalence with (n).
Theorems~\ref{thm:best103} and~\ref{thm:best102} give the equivalences
with (o) and (p). Theorem~\ref{thm:best101} gives the equivalence with (q).
Theorems~\ref{thm:best100-pell}, \ref{thm:best100-encoding} and~\ref{thm:best99}
give the equivalences with (r), (s) and (t), respectively.
Theorems~\ref{thm:best98}, \ref{thm:best98-pell} and~\ref{thm:best97}
give the equivalences with (u), (v) and (w), respectively.
Theorem~\ref{thm:best96} gives the equivalence with (x).
Theorem~\ref{thm:best95} gives the equivalence with (y).
Theorem~\ref{thm:best94} gives the equivalence with (z).
Theorem~\ref{thm:best93} gives the equivalence with (aa).
Theorem~\ref{thm:best92} gives the equivalence with (ab).
Theorem~\ref{thm:best91} gives the equivalence with (ac).
Theorem~\ref{thm:best90} gives the equivalence with (ad).
Theorem~\ref{thm:best89} gives the equivalence with (ae).
The last sentence is
Proposition~\ref{prop:100}.
\end{proof}

\begin{remark}[what the ``other proof'' is]
Each certificate in (c)--(ae) is the ``proof consisting of additions and
multiplications of integers'' of Theorem~5. Its correctness is checked by
performing the listed operations and comparing results; no knowledge of $T$,
of logic, or of the meaning of the integers is needed, because the
soundness directions are guaranteed by Theorems~\ref{thm:sigma}
and~\ref{thm:120}, Corollary~\ref{thm:best114}, and
Theorems~\ref{thm:best112}, \ref{thm:best111}, \ref{thm:best110},
\ref{thm:best109}, \ref{thm:best107}, \ref{thm:modular107},
\ref{thm:best106}, \ref{thm:best105}, \ref{thm:best104},
\ref{thm:best103}, \ref{thm:best102}, \ref{thm:best101},
\ref{thm:best100-pell}, \ref{thm:best100-encoding}, \ref{thm:best99},
\ref{thm:best98}, \ref{thm:best98-pell}, \ref{thm:best97}, \ref{thm:best96},
\ref{thm:best95}, \ref{thm:best94}, \ref{thm:best93}, \ref{thm:best92},
\ref{thm:best91}, \ref{thm:best90}, and~\ref{thm:best89}
for the respective fixed
indices. Three things are not
claimed. The integers are astronomically large: the original construction
requires $b>y\ge2^{5^{59}}$, while the packed construction has the huge
fixed scale $H$, with still larger $q,n,r,a,\kappa,\mu$. The count bounds neither their lengths nor
the number of bit operations. Nothing is said about finding the certificate,
which is as hard as finding a proof in $T$. And the certificate does not
preserve the length or structure of the proof in $T$; a derivation is
recoverable from it only through the encoding of Theorem~\ref{thm:sigma},
as the editorial footnote to Theorem~5 already says.
\end{remark}

\section{How the count arose: typographical complexity}\label{sec:typo}

That an author of Jones's care should leave the cost of $q^{16}$ out of a
count of ``additions and multiplications'' is at first bewildering. The
papers themselves, read in order, make it intelligible. What Jones counted
was the typographical size of the system; the proof-theoretic reading was
attached to that number afterwards, and the convention travelled from one
paper to the next without being re-examined.

\medskip\noindent\textbf{The measure was defined as a count of what is
written.} The 1978 paper introduces $o$ as follows:
\begin{quote}\small
A natural measure of the size or complexity of a system of Diophantine
equations is the number of operations (additions and multiplications),
necessary to write the system. If we regard subtraction as merely a special
form of addition, then the system (1.3) will be seen to involve 243
operations. (Combining the 36 equations into one equation would increase
this number to 350.) \hfill\cite[p.~338]{J78}
\end{quote}
and, a few lines later,
\begin{quote}\small
As a measure of complexity, the number of operations in a Diophantine
equation is very closely associated with the number of symbols necessary to
write down the equation. \hfill\cite[p.~338]{J78}
\end{quote}
Operations ``necessary to write the system'' and ``the number of symbols
necessary to write down the equation'' are counts of signs on the page. An
exponent is a superscript, not an operation sign, and under this measure a
monomial $q^{16}$ is no longer than $q$ by more than one small numeral. That
is precisely Definition~\ref{def:indicated}, and it is why the rule
reproduces 87, 241 and 100.

\medskip\noindent\textbf{The convention was inherited from 1976.} The
primality paper states its bound in the same words and adds the only comment
on its derivation that any of the papers contains:
\begin{quote}\small
If $p$ is a composite number, then there is a proof that $p$ is composite
consisting of a single multiplication. Julia Robinson remarks in [14] that
prior to the solution of Hilbert's tenth problem in 1970, it was not known
that there was a similar proof establishing primality in a bounded number of
steps. Yet it follows from [4] [8] that this is so. And our construction
permits us to compute a bound on this number.\\
\textsc{Theorem 5.} If $p$ is a prime number, then there is a proof that $p$
is prime consisting of only 87 additions and multiplications.\\
The number is easily calculated from the equations of Theorem 2.12.
\hfill\cite[p.~451]{JSWW}
\end{quote}
``Easily calculated'' is a count, not a construction. The system of
Theorem~2.12 contains almost nothing but squares and a few cubes, so on that
system the sign count and an honest straight-line count happen to agree
(Proposition~\ref{prop:calibration} and \cite[\S1]{JSWW}, editorial verification note).
The habit that was harmless there was carried over to Theorem~3 of 1980,
where the powers $q^{16}$, $q^8$, $q^7$, $q^5$, $q^3$, $b^5$, $n^6$, $c^4$
and a fourth power of a trinomial make the sign count diverge from the
initial certificate count. Section~\ref{sec:optimization} reduces the
latter by changing both the calculations and the encoding.

\medskip\noindent\textbf{Powers were treated as notation, as numerals were.}
The 1982 paper restates the measure and immediately draws the
proof-theoretic conclusion:
\begin{quote}\small
Perhaps therefore a better measure of size and complexity of diophantine
equations is to consider the total number of operations, $o$ necessary to
evaluate the polynomial. By an arithmetical operation we shall mean an
addition or multiplication. We will consider subtraction as merely a signed
form of addition. From the equations of Theorem~3, modified to do away with
$q=b^{5^{60}}$, it can be seen that $o=100$. Thus 100 arithmetical
operations are universally sufficient to determine membership of a number
in any r.e.\ set. \hfill\cite[p.~553]{J82}
\end{quote}
``It can be seen'' (1980 has ``it can be shown'') again describes reading
the number off the page. In the same sentence the exponent $5^{60}$ and the
coefficients $2$, $4$, $5$ are numerals that cost nothing, and a power $x^k$
with a fixed exponent is a fixed-degree monomial, a piece of notation of the
same kind. For measuring the size of a polynomial this is a defensible
choice; for counting the steps of a calculator that only adds and multiplies
it is not; and the paper does not separate the two readings.

\medskip\noindent\textbf{The inconsistency was not a considered position.}
The parenthesis ``would increase this number to 350'' in the first quotation
charges, for each of the 36 equations, one subtraction for transposing and
one squaring, plus 35 additions to sum: $243+36+36+35=350$. So the squaring
of a transposed equation cost one operation while the $z^{18}$ inside
(1.3.06) cost nothing. That is the signature of counting the operation signs
that become visible when the equations are rewritten, not of applying a
definition in which exponentiation is free.

\medskip\noindent\textbf{The proof-theoretic gloss was added on top.} Both
1978 and 1980 pass from the size measure to a statement about proofs in one
step:
\begin{quote}\small
Hence from Theorem~3 we obtain an absolute epistemological upper bound on
the complexity of mathematical proofs. \hfill\cite[p.~338]{J78}
\end{quote}
\begin{quote}\small
These solutions faithfully reflect the logical complexity of the original
proof, for the entire deduction is effectively recoverable from the
solution. Hence we have \textsc{Theorem 5}. \hfill\cite[p.~862]{J80}
\end{quote}
The step from ``the theorems form an r.e.\ set, so provability is
solvability of a fixed system'' to ``a proof consisting of $o$ additions and
multiplications'' is sound only for the calculator count of
Definition~\ref{def:certificate}; the number that was actually computed is
the sign count. The corrected edition already softened ``faithfully
reflect'' to ``effectively recoverable, without preservation of length or
complexity''; the present article supplies the number that the
proof-theoretic statement needs: an initial 129-operation certificate,
improved to 89 by Sections~\ref{sec:optimization}--\ref{sec:narrow-binary-product}, alongside the number
that was printed.

\medskip\noindent\textbf{Nobody checked.} The count was announced without a
derivation, repeated verbatim in the full paper, and reported since as a
curiosity; earlier revisions of this edition also left it as ``the
author's reported count''. A number that is never derived in print is never
corrected. In sum: Jones counted the typographical complexity of his system,
said so in 1978, and did not notice that the corollary about proofs requires
the other count.

\section{Limits of this verification}

Proposition~\ref{prop:100} and Theorem~\ref{thm:129} are exact computations
on the transcribed systems and are reproduced by the companion program.
Theorem~\ref{thm:sigma} relies on Lemmas~2.22, 2.25 and 2.26 and the proof of
Theorem~3 of \cite{J82}, as corrected in this edition; the Lean
formalization of those results is in progress and is reported in
\file{Lean/STATUS.md}. The identification of Jones's convention in
Remark~\ref{rem:evidence} is an inference from the published numbers, not a
statement found in the papers: the two exact matches (87 and 100) and the
near match (241 against 243) are the evidence, and the unexplained two signs
in (1.3) and the treatment of the squaring in the ``350'' remark are the
residue. The improved results additionally use the recoding, Pell, and
interval proofs in Sections~\ref{sec:optimization}--\ref{sec:narrow-binary-product};
symbolic residual verification does not establish their positive-domain
or decoding arguments. Those
arguments are supplied as mathematical proofs and have not been
formalized in Lean. All operation bounds in this article are upper bounds;
no optimality is claimed.

\begin{thebibliography}{9}\small
\bibitem{DPR} M.~Davis, H.~Putnam and J.~Robinson, \textit{The decision
problem for exponential diophantine equations}, Ann.\ of Math.\ (2)
\textbf{74} (1961), 425--436.
\bibitem{J78} J.~P.~Jones, \textit{Three universal representations of
recursively enumerable sets}, J.\ Symbolic Logic \textbf{43} (1978),
335--351. Corrected edition: \file{Papers/1978/}.
\bibitem{J80} J.~P.~Jones, \textit{Undecidable diophantine equations},
Bull.\ Amer.\ Math.\ Soc.\ (N.S.) \textbf{3} (1980), 859--862. Corrected
edition: \file{Papers/1980/}.
\bibitem{J82} J.~P.~Jones, \textit{Universal diophantine equation},
J.\ Symbolic Logic \textbf{47} (1982), 549--571. Corrected edition:
\file{Papers/1982/}.
\bibitem{JSWW} J.~P.~Jones, D.~Sato, H.~Wada and D.~Wiens,
\textit{Diophantine representation of the set of prime numbers}, Amer.\
Math.\ Monthly \textbf{83} (1976), 449--464. Corrected edition:
\file{Papers/1976/}; the 87-operation certificate is
\file{Papers/verification/jones1976\_primality87.md}.
\bibitem{Mat70} Ju.~V.~\Matijasevic, \textit{Enumerable sets are
diophantine}, Dokl.\ Akad.\ Nauk SSSR \textbf{191} (1970), 279--282;
English transl.\ Soviet Math.\ Dokl.\ \textbf{11} (1970), 354--357.
\end{thebibliography}

\appendix
\renewcommand{\thesection}{\AlphAlph{\value{section}}}
\section{The original 129-instruction certificate}\label{app:129}

Names: the parameters $x,z,u,y$; the unknowns of Theorem~3 with
$al=\alpha$, $ga=\gamma$, $eta=\eta$, $th=\theta$, $la=\lambda$,
$tau=\tau$, $phi=\varphi$; the new unknowns $ka=\kappa$, $mu=\mu$,
$rho=\rho$, $Delta=\Delta$. Every other name is an auxiliary integer supplied
with the certificate. The column ``Eq.'' gives the equation of $\Sigma$
served by the instruction; ``pow'' marks the shared powers of $q$ and $b$.
The last column is the instruction in addition/multiplication form. The
twenty equality tests that follow the instructions are
\[
\begin{array}{@{}l@{}}
 L1=R1,\ L2=R2,\ L3=b5,\ l=R4,\ e=R5,\ n=q16,\ r=R7,\\
 p=R8,\ L9=R9,\ L10=k2,\ k=R11,\ a=R12,\ c=R13,\ d=R14,\\
 L15=R15,\ L16=R16,\ L17=R17,\ mu=R18,\ L19=R19,\ ka=R20.
\end{array}
\]

\small
% Generated from round4_1980_operation_count.py (SCHEDULE); do not edit by hand.
\begin{longtable}{@{}rlll@{}}
\toprule
No. & Eq. & Assignment & Addition/multiplication check\\
\midrule
\endfirsthead
\toprule
No. & Eq. & Assignment & Addition/multiplication check\\
\midrule
\endhead
\bottomrule
\endfoot
1 & pow & $\mathit{q2} = \mathit{q} \cdot \mathit{q}$ & $\mathit{q} \cdot \mathit{q} = \mathit{q2}$ \\
2 & pow & $\mathit{q3} = \mathit{q2} \cdot \mathit{q}$ & $\mathit{q2} \cdot \mathit{q} = \mathit{q3}$ \\
3 & pow & $\mathit{q4} = \mathit{q2} \cdot \mathit{q2}$ & $\mathit{q2} \cdot \mathit{q2} = \mathit{q4}$ \\
4 & pow & $\mathit{q8} = \mathit{q4} \cdot \mathit{q4}$ & $\mathit{q4} \cdot \mathit{q4} = \mathit{q8}$ \\
5 & pow & $\mathit{q16} = \mathit{q8} \cdot \mathit{q8}$ & $\mathit{q8} \cdot \mathit{q8} = \mathit{q16}$ \\
6 & pow & $\mathit{b2} = \mathit{b} \cdot \mathit{b}$ & $\mathit{b} \cdot \mathit{b} = \mathit{b2}$ \\
7 & pow & $\mathit{b4} = \mathit{b2} \cdot \mathit{b2}$ & $\mathit{b2} \cdot \mathit{b2} = \mathit{b4}$ \\
8 & pow & $\mathit{b5} = \mathit{b4} \cdot \mathit{b}$ & $\mathit{b4} \cdot \mathit{b} = \mathit{b5}$ \\
9 & E1 & $\mathit{g2} = \mathit{g} \cdot \mathit{g}$ & $\mathit{g} \cdot \mathit{g} = \mathit{g2}$ \\
10 & E1 & $\mathit{el} = \mathit{e} \cdot \mathit{l}$ & $\mathit{e} \cdot \mathit{l} = \mathit{el}$ \\
11 & E1 & $\mathit{elg2} = \mathit{el} \cdot \mathit{g2}$ & $\mathit{el} \cdot \mathit{g2} = \mathit{elg2}$ \\
12 & E1 & $\mathit{L1} = \mathit{elg2} + \mathit{al}$ & $\mathit{elg2} + \mathit{al} = \mathit{L1}$ \\
13 & E1 & $\mathit{xy} = \mathit{x} \cdot \mathit{y}$ & $\mathit{x} \cdot \mathit{y} = \mathit{xy}$ \\
14 & E1 & $\mathit{bxy} = \mathit{b} - \mathit{xy}$ & $\mathit{bxy} + \mathit{xy} = \mathit{b}$ \\
15 & E1 & $\mathit{R1} = \mathit{bxy} \cdot \mathit{q2}$ & $\mathit{bxy} \cdot \mathit{q2} = \mathit{R1}$ \\
16 & E2 & $\mathit{L2} = \mathit{la} + \mathit{q4}$ & $\mathit{la} + \mathit{q4} = \mathit{L2}$ \\
17 & E2 & $\mathit{Lam} = \mathit{la} \cdot \mathit{b5}$ & $\mathit{la} \cdot \mathit{b5} = \mathit{Lam}$ \\
18 & E2 & $\mathit{R2} = \mathit{Lam} + 1$ & $\mathit{Lam} + 1 = \mathit{R2}$ \\
19 & E3 & $\mathit{z2} = 2 \cdot \mathit{z}$ & $2 \cdot \mathit{z} = \mathit{z2}$ \\
20 & E3 & $\mathit{L3} = \mathit{th} + \mathit{z2}$ & $\mathit{th} + \mathit{z2} = \mathit{L3}$ \\
21 & E4 & $\mathit{tth} = \mathit{t} \cdot \mathit{th}$ & $\mathit{t} \cdot \mathit{th} = \mathit{tth}$ \\
22 & E4 & $\mathit{R4} = \mathit{u} + \mathit{tth}$ & $\mathit{u} + \mathit{tth} = \mathit{R4}$ \\
23 & E5 & $\mathit{mth} = \mathit{m} \cdot \mathit{th}$ & $\mathit{m} \cdot \mathit{th} = \mathit{mth}$ \\
24 & E5 & $\mathit{R5} = \mathit{y} + \mathit{mth}$ & $\mathit{y} + \mathit{mth} = \mathit{R5}$ \\
25 & E7 & $\mathit{lq2} = \mathit{l} \cdot \mathit{q2}$ & $\mathit{l} \cdot \mathit{q2} = \mathit{lq2}$ \\
26 & E7 & $\mathit{S2} = \mathit{e} + \mathit{lq2}$ & $\mathit{e} + \mathit{lq2} = \mathit{S2}$ \\
27 & E7 & $\mathit{zl} = \mathit{z} \cdot \mathit{la}$ & $\mathit{z} \cdot \mathit{la} = \mathit{zl}$ \\
28 & E7 & $\mathit{ezl} = \mathit{e} - \mathit{zl}$ & $\mathit{ezl} + \mathit{zl} = \mathit{e}$ \\
29 & E7 & $\mathit{t2} = 2 \cdot \mathit{ezl}$ & $2 \cdot \mathit{ezl} = \mathit{t2}$ \\
30 & E7 & $\mathit{xb5} = \mathit{x} \cdot \mathit{b5}$ & $\mathit{x} \cdot \mathit{b5} = \mathit{xb5}$ \\
31 & E7 & $\mathit{c0} = \mathit{xb5} + 1$ & $\mathit{xb5} + 1 = \mathit{c0}$ \\
32 & E7 & $\mathit{C} = \mathit{c0} + \mathit{g}$ & $\mathit{c0} + \mathit{g} = \mathit{C}$ \\
33 & E7 & $\mathit{C2} = \mathit{C} \cdot \mathit{C}$ & $\mathit{C} \cdot \mathit{C} = \mathit{C2}$ \\
34 & E7 & $\mathit{C4} = \mathit{C2} \cdot \mathit{C2}$ & $\mathit{C2} \cdot \mathit{C2} = \mathit{C4}$ \\
35 & E7 & $\mathit{P1} = \mathit{t2} \cdot \mathit{C4}$ & $\mathit{t2} \cdot \mathit{C4} = \mathit{P1}$ \\
36 & E7 & $\mathit{q4p1} = \mathit{q4} + 1$ & $\mathit{q4} + 1 = \mathit{q4p1}$ \\
37 & E7 & $\mathit{P2} = \mathit{Lam} \cdot \mathit{q4p1}$ & $\mathit{Lam} \cdot \mathit{q4p1} = \mathit{P2}$ \\
38 & E7 & $\mathit{S3} = \mathit{P1} + \mathit{P2}$ & $\mathit{P1} + \mathit{P2} = \mathit{S3}$ \\
39 & E7 & $\mathit{S3q4} = \mathit{S3} \cdot \mathit{q4}$ & $\mathit{S3} \cdot \mathit{q4} = \mathit{S3q4}$ \\
40 & E7 & $\mathit{Sin} = \mathit{S2} + \mathit{S3q4}$ & $\mathit{S2} + \mathit{S3q4} = \mathit{Sin}$ \\
41 & E7 & $\mathit{Sq3} = \mathit{Sin} \cdot \mathit{q3}$ & $\mathit{Sin} \cdot \mathit{q3} = \mathit{Sq3}$ \\
42 & E7 & $\mathit{S} = \mathit{g} + \mathit{Sq3}$ & $\mathit{g} + \mathit{Sq3} = \mathit{S}$ \\
43 & E7 & $\mathit{n2} = \mathit{n} \cdot \mathit{n}$ & $\mathit{n} \cdot \mathit{n} = \mathit{n2}$ \\
44 & E7 & $\mathit{n2n} = \mathit{n2} - \mathit{n}$ & $\mathit{n2n} + \mathit{n} = \mathit{n2}$ \\
45 & E7 & $\mathit{SA} = \mathit{S} \cdot \mathit{n2n}$ & $\mathit{S} \cdot \mathit{n2n} = \mathit{SA}$ \\
46 & E7 & $\mathit{bl} = \mathit{b} \cdot \mathit{l}$ & $\mathit{b} \cdot \mathit{l} = \mathit{bl}$ \\
47 & E7 & $\mathit{q3bl} = \mathit{q3} - \mathit{bl}$ & $\mathit{q3bl} + \mathit{bl} = \mathit{q3}$ \\
48 & E7 & $\mathit{T1} = \mathit{q3bl} + \mathit{l}$ & $\mathit{q3bl} + \mathit{l} = \mathit{T1}$ \\
49 & E7 & $\mathit{thl} = \mathit{th} \cdot \mathit{la}$ & $\mathit{th} \cdot \mathit{la} = \mathit{thl}$ \\
50 & E7 & $\mathit{thlq3} = \mathit{thl} \cdot \mathit{q3}$ & $\mathit{thl} \cdot \mathit{q3} = \mathit{thlq3}$ \\
51 & E7 & $\mathit{T2} = \mathit{T1} + \mathit{thlq3}$ & $\mathit{T1} + \mathit{thlq3} = \mathit{T2}$ \\
52 & E7 & $\mathit{b5m2} = \mathit{b5} - 2$ & $\mathit{b5m2} + 2 = \mathit{b5}$ \\
53 & E7 & $\mathit{b5m2q8} = \mathit{b5m2} \cdot \mathit{q8}$ & $\mathit{b5m2} \cdot \mathit{q8} = \mathit{b5m2q8}$ \\
54 & E7 & $\mathit{T} = \mathit{T2} + \mathit{b5m2q8}$ & $\mathit{T2} + \mathit{b5m2q8} = \mathit{T}$ \\
55 & E7 & $\mathit{n2m1} = \mathit{n2} - 1$ & $\mathit{n2m1} + 1 = \mathit{n2}$ \\
56 & E7 & $\mathit{TB} = \mathit{T} \cdot \mathit{n2m1}$ & $\mathit{T} \cdot \mathit{n2m1} = \mathit{TB}$ \\
57 & E7 & $\mathit{R7} = \mathit{SA} + \mathit{TB}$ & $\mathit{SA} + \mathit{TB} = \mathit{R7}$ \\
58 & E12 & $\mathit{wn2} = \mathit{w} \cdot \mathit{n2}$ & $\mathit{w} \cdot \mathit{n2} = \mathit{wn2}$ \\
59 & E12 & $\mathit{wn2p1} = \mathit{wn2} + 1$ & $\mathit{wn2} + 1 = \mathit{wn2p1}$ \\
60 & E12 & $\mathit{sn2} = \mathit{s} \cdot \mathit{n2}$ & $\mathit{s} \cdot \mathit{n2} = \mathit{sn2}$ \\
61 & E12 & $\mathit{rsn2} = \mathit{r} \cdot \mathit{sn2}$ & $\mathit{r} \cdot \mathit{sn2} = \mathit{rsn2}$ \\
62 & E12 & $\mathit{R12} = \mathit{wn2p1} \cdot \mathit{rsn2}$ & $\mathit{wn2p1} \cdot \mathit{rsn2} = \mathit{R12}$ \\
63 & E8 & $\mathit{rsn2sq} = \mathit{rsn2} \cdot \mathit{rsn2}$ & $\mathit{rsn2} \cdot \mathit{rsn2} = \mathit{rsn2sq}$ \\
64 & E8 & $\mathit{wsq} = \mathit{wn2} \cdot \mathit{rsn2sq}$ & $\mathit{wn2} \cdot \mathit{rsn2sq} = \mathit{wsq}$ \\
65 & E8 & $\mathit{R8} = 2 \cdot \mathit{wsq}$ & $2 \cdot \mathit{wsq} = \mathit{R8}$ \\
66 & E9 & $\mathit{p2} = \mathit{p} \cdot \mathit{p}$ & $\mathit{p} \cdot \mathit{p} = \mathit{p2}$ \\
67 & E9 & $\mathit{p2m1} = \mathit{p2} - 1$ & $\mathit{p2m1} + 1 = \mathit{p2}$ \\
68 & E9 & $\mathit{k2} = \mathit{k} \cdot \mathit{k}$ & $\mathit{k} \cdot \mathit{k} = \mathit{k2}$ \\
69 & E9 & $\mathit{P9} = \mathit{p2m1} \cdot \mathit{k2}$ & $\mathit{p2m1} \cdot \mathit{k2} = \mathit{P9}$ \\
70 & E9 & $\mathit{L9} = \mathit{P9} + 1$ & $\mathit{P9} + 1 = \mathit{L9}$ \\
71 & E9 & $\mathit{R9} = \mathit{tau} \cdot \mathit{tau}$ & $\mathit{tau} \cdot \mathit{tau} = \mathit{R9}$ \\
72 & E10 & $\mathit{ksn2} = \mathit{k} \cdot \mathit{sn2}$ & $\mathit{k} \cdot \mathit{sn2} = \mathit{ksn2}$ \\
73 & E10 & $\mathit{cm} = \mathit{c} - \mathit{ksn2}$ & $\mathit{cm} + \mathit{ksn2} = \mathit{c}$ \\
74 & E10 & $\mathit{cm2} = \mathit{cm} \cdot \mathit{cm}$ & $\mathit{cm} \cdot \mathit{cm} = \mathit{cm2}$ \\
75 & E10 & $\mathit{f4} = 4 \cdot \mathit{cm2}$ & $4 \cdot \mathit{cm2} = \mathit{f4}$ \\
76 & E10 & $\mathit{L10} = \mathit{f4} + \mathit{eta}$ & $\mathit{f4} + \mathit{eta} = \mathit{L10}$ \\
77 & E11 & $\mathit{r1} = \mathit{r} + 1$ & $\mathit{r} + 1 = \mathit{r1}$ \\
78 & E11 & $\mathit{pm1} = \mathit{p} - 1$ & $\mathit{pm1} + 1 = \mathit{p}$ \\
79 & E11 & $\mathit{hpm1} = \mathit{h} \cdot \mathit{pm1}$ & $\mathit{h} \cdot \mathit{pm1} = \mathit{hpm1}$ \\
80 & E11 & $\mathit{R11} = \mathit{r1} + \mathit{hpm1}$ & $\mathit{r1} + \mathit{hpm1} = \mathit{R11}$ \\
81 & E13 & $\mathit{tr1} = \mathit{r1} + \mathit{r}$ & $\mathit{r1} + \mathit{r} = \mathit{tr1}$ \\
82 & E13 & $\mathit{tk} = \mathit{tr1} + \mathit{ka}$ & $\mathit{tr1} + \mathit{ka} = \mathit{tk}$ \\
83 & E13 & $\mathit{R13} = \mathit{tk} + \mathit{phi}$ & $\mathit{tk} + \mathit{phi} = \mathit{R13}$ \\
84 & E14 & $\mathit{bw} = \mathit{b} \cdot \mathit{w}$ & $\mathit{b} \cdot \mathit{w} = \mathit{bw}$ \\
85 & E14 & $\mathit{am2} = \mathit{a} - 2$ & $\mathit{am2} + 2 = \mathit{a}$ \\
86 & E14 & $\mathit{cam2} = \mathit{c} \cdot \mathit{am2}$ & $\mathit{c} \cdot \mathit{am2} = \mathit{cam2}$ \\
87 & E14 & $\mathit{D1} = \mathit{bw} + \mathit{cam2}$ & $\mathit{bw} + \mathit{cam2} = \mathit{D1}$ \\
88 & E14 & $\mathit{a4} = 4 \cdot \mathit{a}$ & $4 \cdot \mathit{a} = \mathit{a4}$ \\
89 & E14 & $\mathit{a4m5} = \mathit{a4} - 5$ & $\mathit{a4m5} + 5 = \mathit{a4}$ \\
90 & E14 & $\mathit{gam} = \mathit{ga} \cdot \mathit{a4m5}$ & $\mathit{ga} \cdot \mathit{a4m5} = \mathit{gam}$ \\
91 & E14 & $\mathit{R14} = \mathit{D1} + \mathit{gam}$ & $\mathit{D1} + \mathit{gam} = \mathit{R14}$ \\
92 & E15 & $\mathit{a2} = \mathit{a} \cdot \mathit{a}$ & $\mathit{a} \cdot \mathit{a} = \mathit{a2}$ \\
93 & E15 & $\mathit{A} = \mathit{a2} - 1$ & $\mathit{A} + 1 = \mathit{a2}$ \\
94 & E15 & $\mathit{c2} = \mathit{c} \cdot \mathit{c}$ & $\mathit{c} \cdot \mathit{c} = \mathit{c2}$ \\
95 & E15 & $\mathit{Ac2} = \mathit{A} \cdot \mathit{c2}$ & $\mathit{A} \cdot \mathit{c2} = \mathit{Ac2}$ \\
96 & E15 & $\mathit{R15} = \mathit{Ac2} + 1$ & $\mathit{Ac2} + 1 = \mathit{R15}$ \\
97 & E15 & $\mathit{L15} = \mathit{d} \cdot \mathit{d}$ & $\mathit{d} \cdot \mathit{d} = \mathit{L15}$ \\
98 & E16 & $\mathit{ic2} = \mathit{i} \cdot \mathit{c2}$ & $\mathit{i} \cdot \mathit{c2} = \mathit{ic2}$ \\
99 & E16 & $\mathit{ic22} = \mathit{ic2} \cdot \mathit{ic2}$ & $\mathit{ic2} \cdot \mathit{ic2} = \mathit{ic22}$ \\
100 & E16 & $\mathit{AE} = \mathit{A} \cdot \mathit{ic22}$ & $\mathit{A} \cdot \mathit{ic22} = \mathit{AE}$ \\
101 & E16 & $\mathit{R16} = \mathit{AE} + 1$ & $\mathit{AE} + 1 = \mathit{R16}$ \\
102 & E16 & $\mathit{L16} = \mathit{f} \cdot \mathit{f}$ & $\mathit{f} \cdot \mathit{f} = \mathit{L16}$ \\
103 & E17 & $\mathit{of} = \mathit{o} \cdot \mathit{f}$ & $\mathit{o} \cdot \mathit{f} = \mathit{of}$ \\
104 & E17 & $\mathit{dof} = \mathit{d} + \mathit{of}$ & $\mathit{d} + \mathit{of} = \mathit{dof}$ \\
105 & E17 & $\mathit{L17} = \mathit{dof} \cdot \mathit{dof}$ & $\mathit{dof} \cdot \mathit{dof} = \mathit{L17}$ \\
106 & E17 & $\mathit{d2ma} = \mathit{L15} - \mathit{a}$ & $\mathit{d2ma} + \mathit{a} = \mathit{L15}$ \\
107 & E17 & $\mathit{f2d} = \mathit{L16} \cdot \mathit{d2ma}$ & $\mathit{L16} \cdot \mathit{d2ma} = \mathit{f2d}$ \\
108 & E17 & $\mathit{G} = \mathit{a} + \mathit{f2d}$ & $\mathit{a} + \mathit{f2d} = \mathit{G}$ \\
109 & E17 & $\mathit{G2} = \mathit{G} \cdot \mathit{G}$ & $\mathit{G} \cdot \mathit{G} = \mathit{G2}$ \\
110 & E17 & $\mathit{G2m1} = \mathit{G2} - 1$ & $\mathit{G2m1} + 1 = \mathit{G2}$ \\
111 & E17 & $\mathit{jc} = \mathit{j} \cdot \mathit{c}$ & $\mathit{j} \cdot \mathit{c} = \mathit{jc}$ \\
112 & E17 & $\mathit{H} = \mathit{tr1} + \mathit{jc}$ & $\mathit{tr1} + \mathit{jc} = \mathit{H}$ \\
113 & E17 & $\mathit{H2} = \mathit{H} \cdot \mathit{H}$ & $\mathit{H} \cdot \mathit{H} = \mathit{H2}$ \\
114 & E17 & $\mathit{P17} = \mathit{G2m1} \cdot \mathit{H2}$ & $\mathit{G2m1} \cdot \mathit{H2} = \mathit{P17}$ \\
115 & E17 & $\mathit{R17} = \mathit{P17} + 1$ & $\mathit{P17} + 1 = \mathit{R17}$ \\
116 & E18 & $\mathit{amb} = \mathit{a} - \mathit{b}$ & $\mathit{amb} + \mathit{b} = \mathit{a}$ \\
117 & E18 & $\mathit{kamb} = \mathit{ka} \cdot \mathit{amb}$ & $\mathit{ka} \cdot \mathit{amb} = \mathit{kamb}$ \\
118 & E18 & $\mathit{qk} = \mathit{q} + \mathit{kamb}$ & $\mathit{q} + \mathit{kamb} = \mathit{qk}$ \\
119 & E18 & $\mathit{amb2} = \mathit{amb} \cdot \mathit{amb}$ & $\mathit{amb} \cdot \mathit{amb} = \mathit{amb2}$ \\
120 & E18 & $\mathit{Mod} = \mathit{A} - \mathit{amb2}$ & $\mathit{Mod} + \mathit{amb2} = \mathit{A}$ \\
121 & E18 & $\mathit{rM} = \mathit{rho} \cdot \mathit{Mod}$ & $\mathit{rho} \cdot \mathit{Mod} = \mathit{rM}$ \\
122 & E18 & $\mathit{R18} = \mathit{qk} + \mathit{rM}$ & $\mathit{qk} + \mathit{rM} = \mathit{R18}$ \\
123 & E19 & $\mathit{ka2} = \mathit{ka} \cdot \mathit{ka}$ & $\mathit{ka} \cdot \mathit{ka} = \mathit{ka2}$ \\
124 & E19 & $\mathit{Aka2} = \mathit{A} \cdot \mathit{ka2}$ & $\mathit{A} \cdot \mathit{ka2} = \mathit{Aka2}$ \\
125 & E19 & $\mathit{R19} = \mathit{Aka2} + 1$ & $\mathit{Aka2} + 1 = \mathit{R19}$ \\
126 & E19 & $\mathit{L19} = \mathit{mu} \cdot \mathit{mu}$ & $\mathit{mu} \cdot \mathit{mu} = \mathit{L19}$ \\
127 & E20 & $\mathit{am1} = \mathit{a} - 1$ & $\mathit{am1} + 1 = \mathit{a}$ \\
128 & E20 & $\mathit{Dam1} = \mathit{Delta} \cdot \mathit{am1}$ & $\mathit{Delta} \cdot \mathit{am1} = \mathit{Dam1}$ \\
129 & E20 & $\mathit{R20} = \mathit{Dam1} + 5^{60}$ & $\mathit{Dam1} + 5^{60} = \mathit{R20}$ \\
\end{longtable}


\normalsize
\section{The 120-instruction certificate}\label{app:120}

The four inputs are $x,Z,V,H$, with $(Z,V,H)$ a fixed admissible packed
index. The positive witnesses are those of Appendix~\ref{app:129}, with
$m$ removed and $\beta$ added. Here $\mathit{beta}=\beta$;
the other Greek names use the same transliterations. Every additional
name denotes a supplied integer whose value is constrained by the listed
checks. In particular, $B=Hb^4$ and $C=1+xB+g$ are computed, while $c$
is the distinct Pell witness. The literal $5^{59}$ is denoted explicitly.
All assignments and all 20 final equality tests are generated from
\file{round4\_1980\_packed\_certificate.py}; reversing each subtraction
gives the addition in the last column. The E7 test is equivalent to its
source equation after the E2 and E3 tests, as proved in
Section~\ref{sec:optimization}.

\small
% Generated by round4_1980_packed_schedule_table.py; do not edit.
\begin{longtable}{@{}rll@{}}
\toprule
No. & Assignment & Addition/multiplication check\\
\midrule
\endfirsthead
\toprule
No. & Assignment & Addition/multiplication check\\
\midrule
\endhead
\bottomrule
\endfoot
1 & $\mathit{q2} = \mathit{q} \cdot \mathit{q}$ & $\mathit{q} \cdot \mathit{q} = \mathit{q2}$ \\
2 & $\mathit{q3} = \mathit{q2} \cdot \mathit{q}$ & $\mathit{q2} \cdot \mathit{q} = \mathit{q3}$ \\
3 & $\mathit{q4} = \mathit{q2} \cdot \mathit{q2}$ & $\mathit{q2} \cdot \mathit{q2} = \mathit{q4}$ \\
4 & $\mathit{q8} = \mathit{q4} \cdot \mathit{q4}$ & $\mathit{q4} \cdot \mathit{q4} = \mathit{q8}$ \\
5 & $\mathit{q16} = \mathit{q8} \cdot \mathit{q8}$ & $\mathit{q8} \cdot \mathit{q8} = \mathit{q16}$ \\
6 & $\mathit{b2} = \mathit{b} \cdot \mathit{b}$ & $\mathit{b} \cdot \mathit{b} = \mathit{b2}$ \\
7 & $\mathit{b4} = \mathit{b2} \cdot \mathit{b2}$ & $\mathit{b2} \cdot \mathit{b2} = \mathit{b4}$ \\
8 & $\mathit{B} = \mathit{H} \cdot \mathit{b4}$ & $\mathit{H} \cdot \mathit{b4} = \mathit{B}$ \\
9 & $\mathit{el} = \mathit{e} \cdot \mathit{l}$ & $\mathit{e} \cdot \mathit{l} = \mathit{el}$ \\
10 & $\mathit{bx} = \mathit{x} + \mathit{beta}$ & $\mathit{x} + \mathit{beta} = \mathit{bx}$ \\
11 & $\mathit{L2} = \mathit{la} + \mathit{q4}$ & $\mathit{la} + \mathit{q4} = \mathit{L2}$ \\
12 & $\mathit{Lam} = \mathit{la} \cdot \mathit{B}$ & $\mathit{la} \cdot \mathit{B} = \mathit{Lam}$ \\
13 & $\mathit{R2} = \mathit{Lam} + 1$ & $\mathit{Lam} + 1 = \mathit{R2}$ \\
14 & $\mathit{L3} = \mathit{th} + \mathit{Z}$ & $\mathit{th} + \mathit{Z} = \mathit{L3}$ \\
15 & $\mathit{lq2} = \mathit{l} \cdot \mathit{q2}$ & $\mathit{l} \cdot \mathit{q2} = \mathit{lq2}$ \\
16 & $\mathit{S2} = \mathit{e} + \mathit{lq2}$ & $\mathit{e} + \mathit{lq2} = \mathit{S2}$ \\
17 & $\mathit{pack\_tth} = \mathit{t} \cdot \mathit{th}$ & $\mathit{t} \cdot \mathit{th} = \mathit{pack\_tth}$ \\
18 & $\mathit{pack\_rhs} = \mathit{V} + \mathit{pack\_tth}$ & $\mathit{V} + \mathit{pack\_tth} = \mathit{pack\_rhs}$ \\
19 & $\mathit{zl} = \mathit{Z} \cdot \mathit{la}$ & $\mathit{Z} \cdot \mathit{la} = \mathit{zl}$ \\
20 & $\mathit{twice\_e} = 2 \cdot \mathit{e}$ & $2 \cdot \mathit{e} = \mathit{twice\_e}$ \\
21 & $\mathit{t2} = \mathit{twice\_e} - \mathit{zl}$ & $\mathit{t2} + \mathit{zl} = \mathit{twice\_e}$ \\
22 & $\mathit{xb5} = \mathit{x} \cdot \mathit{B}$ & $\mathit{x} \cdot \mathit{B} = \mathit{xb5}$ \\
23 & $\mathit{c0} = \mathit{xb5} + 1$ & $\mathit{xb5} + 1 = \mathit{c0}$ \\
24 & $\mathit{C} = \mathit{c0} + \mathit{g}$ & $\mathit{c0} + \mathit{g} = \mathit{C}$ \\
25 & $\mathit{C2} = \mathit{C} \cdot \mathit{C}$ & $\mathit{C} \cdot \mathit{C} = \mathit{C2}$ \\
26 & $\mathit{elC2} = \mathit{el} \cdot \mathit{C2}$ & $\mathit{el} \cdot \mathit{C2} = \mathit{elC2}$ \\
27 & $\mathit{L1} = \mathit{elC2} + \mathit{al}$ & $\mathit{elC2} + \mathit{al} = \mathit{L1}$ \\
28 & $\mathit{C4} = \mathit{C2} \cdot \mathit{C2}$ & $\mathit{C2} \cdot \mathit{C2} = \mathit{C4}$ \\
29 & $\mathit{P1} = \mathit{t2} \cdot \mathit{C4}$ & $\mathit{t2} \cdot \mathit{C4} = \mathit{P1}$ \\
30 & $\mathit{q4p1} = \mathit{q4} + 1$ & $\mathit{q4} + 1 = \mathit{q4p1}$ \\
31 & $\mathit{P2} = \mathit{Lam} \cdot \mathit{q4p1}$ & $\mathit{Lam} \cdot \mathit{q4p1} = \mathit{P2}$ \\
32 & $\mathit{S3} = \mathit{P1} + \mathit{P2}$ & $\mathit{P1} + \mathit{P2} = \mathit{S3}$ \\
33 & $\mathit{S3q4} = \mathit{S3} \cdot \mathit{q4}$ & $\mathit{S3} \cdot \mathit{q4} = \mathit{S3q4}$ \\
34 & $\mathit{Sin} = \mathit{S2} + \mathit{S3q4}$ & $\mathit{S2} + \mathit{S3q4} = \mathit{Sin}$ \\
35 & $\mathit{Sq3} = \mathit{Sin} \cdot \mathit{q3}$ & $\mathit{Sin} \cdot \mathit{q3} = \mathit{Sq3}$ \\
36 & $\mathit{S} = \mathit{g} + \mathit{Sq3}$ & $\mathit{g} + \mathit{Sq3} = \mathit{S}$ \\
37 & $\mathit{n2} = \mathit{n} \cdot \mathit{n}$ & $\mathit{n} \cdot \mathit{n} = \mathit{n2}$ \\
38 & $\mathit{n2n} = \mathit{n2} - \mathit{n}$ & $\mathit{n2n} + \mathit{n} = \mathit{n2}$ \\
39 & $\mathit{SA} = \mathit{S} \cdot \mathit{n2n}$ & $\mathit{S} \cdot \mathit{n2n} = \mathit{SA}$ \\
40 & $\mathit{Tcoef} = \mathit{L2} - \mathit{zl}$ & $\mathit{Tcoef} + \mathit{zl} = \mathit{L2}$ \\
41 & $\mathit{Tq3} = \mathit{Tcoef} \cdot \mathit{q3}$ & $\mathit{Tcoef} \cdot \mathit{q3} = \mathit{Tq3}$ \\
42 & $\mathit{bm1} = \mathit{b} - 1$ & $\mathit{bm1} + 1 = \mathit{b}$ \\
43 & $\mathit{lbm1} = \mathit{l} \cdot \mathit{bm1}$ & $\mathit{l} \cdot \mathit{bm1} = \mathit{lbm1}$ \\
44 & $\mathit{T2} = \mathit{Tq3} - \mathit{lbm1}$ & $\mathit{T2} + \mathit{lbm1} = \mathit{Tq3}$ \\
45 & $\mathit{b5m2} = \mathit{B} - 2$ & $\mathit{b5m2} + 2 = \mathit{B}$ \\
46 & $\mathit{b5m2q8} = \mathit{b5m2} \cdot \mathit{q8}$ & $\mathit{b5m2} \cdot \mathit{q8} = \mathit{b5m2q8}$ \\
47 & $\mathit{T} = \mathit{T2} + \mathit{b5m2q8}$ & $\mathit{T2} + \mathit{b5m2q8} = \mathit{T}$ \\
48 & $\mathit{n2m1} = \mathit{n2} - 1$ & $\mathit{n2m1} + 1 = \mathit{n2}$ \\
49 & $\mathit{TB} = \mathit{T} \cdot \mathit{n2m1}$ & $\mathit{T} \cdot \mathit{n2m1} = \mathit{TB}$ \\
50 & $\mathit{R7} = \mathit{SA} + \mathit{TB}$ & $\mathit{SA} + \mathit{TB} = \mathit{R7}$ \\
51 & $\mathit{wn2} = \mathit{w} \cdot \mathit{n2}$ & $\mathit{w} \cdot \mathit{n2} = \mathit{wn2}$ \\
52 & $\mathit{sn2} = \mathit{s} \cdot \mathit{n2}$ & $\mathit{s} \cdot \mathit{n2} = \mathit{sn2}$ \\
53 & $\mathit{rsn2} = \mathit{r} \cdot \mathit{sn2}$ & $\mathit{r} \cdot \mathit{sn2} = \mathit{rsn2}$ \\
54 & $\mathit{UM} = \mathit{wn2} \cdot \mathit{rsn2}$ & $\mathit{wn2} \cdot \mathit{rsn2} = \mathit{UM}$ \\
55 & $\mathit{R12} = \mathit{UM} + \mathit{rsn2}$ & $\mathit{UM} + \mathit{rsn2} = \mathit{R12}$ \\
56 & $\mathit{wsq} = \mathit{UM} \cdot \mathit{rsn2}$ & $\mathit{UM} \cdot \mathit{rsn2} = \mathit{wsq}$ \\
57 & $\mathit{R8} = 2 \cdot \mathit{wsq}$ & $2 \cdot \mathit{wsq} = \mathit{R8}$ \\
58 & $\mathit{p2} = \mathit{p} \cdot \mathit{p}$ & $\mathit{p} \cdot \mathit{p} = \mathit{p2}$ \\
59 & $\mathit{p2m1} = \mathit{p2} - 1$ & $\mathit{p2m1} + 1 = \mathit{p2}$ \\
60 & $\mathit{k2} = \mathit{k} \cdot \mathit{k}$ & $\mathit{k} \cdot \mathit{k} = \mathit{k2}$ \\
61 & $\mathit{P9} = \mathit{p2m1} \cdot \mathit{k2}$ & $\mathit{p2m1} \cdot \mathit{k2} = \mathit{P9}$ \\
62 & $\mathit{L9} = \mathit{P9} + 1$ & $\mathit{P9} + 1 = \mathit{L9}$ \\
63 & $\mathit{R9} = \mathit{tau} \cdot \mathit{tau}$ & $\mathit{tau} \cdot \mathit{tau} = \mathit{R9}$ \\
64 & $\mathit{ksn2} = \mathit{k} \cdot \mathit{sn2}$ & $\mathit{k} \cdot \mathit{sn2} = \mathit{ksn2}$ \\
65 & $\mathit{cm} = \mathit{c} - \mathit{ksn2}$ & $\mathit{cm} + \mathit{ksn2} = \mathit{c}$ \\
66 & $\mathit{cm2} = \mathit{cm} \cdot \mathit{cm}$ & $\mathit{cm} \cdot \mathit{cm} = \mathit{cm2}$ \\
67 & $\mathit{f4} = 4 \cdot \mathit{cm2}$ & $4 \cdot \mathit{cm2} = \mathit{f4}$ \\
68 & $\mathit{L10} = \mathit{f4} + \mathit{eta}$ & $\mathit{f4} + \mathit{eta} = \mathit{L10}$ \\
69 & $\mathit{r1} = \mathit{r} + 1$ & $\mathit{r} + 1 = \mathit{r1}$ \\
70 & $\mathit{pm1} = \mathit{p} - 1$ & $\mathit{pm1} + 1 = \mathit{p}$ \\
71 & $\mathit{hpm1} = \mathit{h} \cdot \mathit{pm1}$ & $\mathit{h} \cdot \mathit{pm1} = \mathit{hpm1}$ \\
72 & $\mathit{R11} = \mathit{r1} + \mathit{hpm1}$ & $\mathit{r1} + \mathit{hpm1} = \mathit{R11}$ \\
73 & $\mathit{tr1} = \mathit{r1} + \mathit{r}$ & $\mathit{r1} + \mathit{r} = \mathit{tr1}$ \\
74 & $\mathit{R13} = \mathit{ka} + \mathit{phi}$ & $\mathit{ka} + \mathit{phi} = \mathit{R13}$ \\
75 & $\mathit{bw} = \mathit{b} \cdot \mathit{w}$ & $\mathit{b} \cdot \mathit{w} = \mathit{bw}$ \\
76 & $\mathit{am2} = \mathit{a} - 2$ & $\mathit{am2} + 2 = \mathit{a}$ \\
77 & $\mathit{cam2} = \mathit{c} \cdot \mathit{am2}$ & $\mathit{c} \cdot \mathit{am2} = \mathit{cam2}$ \\
78 & $\mathit{D1} = \mathit{bw} + \mathit{cam2}$ & $\mathit{bw} + \mathit{cam2} = \mathit{D1}$ \\
79 & $\mathit{a4} = 4 \cdot \mathit{a}$ & $4 \cdot \mathit{a} = \mathit{a4}$ \\
80 & $\mathit{a4m5} = \mathit{a4} - 5$ & $\mathit{a4m5} + 5 = \mathit{a4}$ \\
81 & $\mathit{gam} = \mathit{ga} \cdot \mathit{a4m5}$ & $\mathit{ga} \cdot \mathit{a4m5} = \mathit{gam}$ \\
82 & $\mathit{R14} = \mathit{D1} + \mathit{gam}$ & $\mathit{D1} + \mathit{gam} = \mathit{R14}$ \\
83 & $\mathit{a2} = \mathit{a} \cdot \mathit{a}$ & $\mathit{a} \cdot \mathit{a} = \mathit{a2}$ \\
84 & $\mathit{A} = \mathit{a2} - 1$ & $\mathit{A} + 1 = \mathit{a2}$ \\
85 & $\mathit{c2} = \mathit{c} \cdot \mathit{c}$ & $\mathit{c} \cdot \mathit{c} = \mathit{c2}$ \\
86 & $\mathit{Ac2} = \mathit{A} \cdot \mathit{c2}$ & $\mathit{A} \cdot \mathit{c2} = \mathit{Ac2}$ \\
87 & $\mathit{R15} = \mathit{Ac2} + 1$ & $\mathit{Ac2} + 1 = \mathit{R15}$ \\
88 & $\mathit{L15} = \mathit{d} \cdot \mathit{d}$ & $\mathit{d} \cdot \mathit{d} = \mathit{L15}$ \\
89 & $\mathit{ic2} = \mathit{i} \cdot \mathit{c2}$ & $\mathit{i} \cdot \mathit{c2} = \mathit{ic2}$ \\
90 & $\mathit{ic22} = \mathit{ic2} \cdot \mathit{ic2}$ & $\mathit{ic2} \cdot \mathit{ic2} = \mathit{ic22}$ \\
91 & $\mathit{AE} = \mathit{A} \cdot \mathit{ic22}$ & $\mathit{A} \cdot \mathit{ic22} = \mathit{AE}$ \\
92 & $\mathit{R16} = \mathit{AE} + 1$ & $\mathit{AE} + 1 = \mathit{R16}$ \\
93 & $\mathit{L16} = \mathit{f} \cdot \mathit{f}$ & $\mathit{f} \cdot \mathit{f} = \mathit{L16}$ \\
94 & $\mathit{of} = \mathit{o} \cdot \mathit{f}$ & $\mathit{o} \cdot \mathit{f} = \mathit{of}$ \\
95 & $\mathit{dof} = \mathit{d} + \mathit{of}$ & $\mathit{d} + \mathit{of} = \mathit{dof}$ \\
96 & $\mathit{L17} = \mathit{dof} \cdot \mathit{dof}$ & $\mathit{dof} \cdot \mathit{dof} = \mathit{L17}$ \\
97 & $\mathit{d2ma} = \mathit{L15} - \mathit{a}$ & $\mathit{d2ma} + \mathit{a} = \mathit{L15}$ \\
98 & $\mathit{f2d} = \mathit{L16} \cdot \mathit{d2ma}$ & $\mathit{L16} \cdot \mathit{d2ma} = \mathit{f2d}$ \\
99 & $\mathit{G} = \mathit{a} + \mathit{f2d}$ & $\mathit{a} + \mathit{f2d} = \mathit{G}$ \\
100 & $\mathit{G2} = \mathit{G} \cdot \mathit{G}$ & $\mathit{G} \cdot \mathit{G} = \mathit{G2}$ \\
101 & $\mathit{G2m1} = \mathit{G2} - 1$ & $\mathit{G2m1} + 1 = \mathit{G2}$ \\
102 & $\mathit{jc} = \mathit{j} \cdot \mathit{c}$ & $\mathit{j} \cdot \mathit{c} = \mathit{jc}$ \\
103 & $\mathit{H17} = \mathit{tr1} + \mathit{jc}$ & $\mathit{tr1} + \mathit{jc} = \mathit{H17}$ \\
104 & $\mathit{H2} = \mathit{H17} \cdot \mathit{H17}$ & $\mathit{H17} \cdot \mathit{H17} = \mathit{H2}$ \\
105 & $\mathit{P17} = \mathit{G2m1} \cdot \mathit{H2}$ & $\mathit{G2m1} \cdot \mathit{H2} = \mathit{P17}$ \\
106 & $\mathit{R17} = \mathit{P17} + 1$ & $\mathit{P17} + 1 = \mathit{R17}$ \\
107 & $\mathit{amb} = \mathit{a} - \mathit{B}$ & $\mathit{amb} + \mathit{B} = \mathit{a}$ \\
108 & $\mathit{kamb} = \mathit{ka} \cdot \mathit{amb}$ & $\mathit{ka} \cdot \mathit{amb} = \mathit{kamb}$ \\
109 & $\mathit{qk} = \mathit{q} + \mathit{kamb}$ & $\mathit{q} + \mathit{kamb} = \mathit{qk}$ \\
110 & $\mathit{amb2} = \mathit{amb} \cdot \mathit{amb}$ & $\mathit{amb} \cdot \mathit{amb} = \mathit{amb2}$ \\
111 & $\mathit{Mod} = \mathit{A} - \mathit{amb2}$ & $\mathit{Mod} + \mathit{amb2} = \mathit{A}$ \\
112 & $\mathit{rM} = \mathit{rho} \cdot \mathit{Mod}$ & $\mathit{rho} \cdot \mathit{Mod} = \mathit{rM}$ \\
113 & $\mathit{R18} = \mathit{qk} + \mathit{rM}$ & $\mathit{qk} + \mathit{rM} = \mathit{R18}$ \\
114 & $\mathit{ka2} = \mathit{ka} \cdot \mathit{ka}$ & $\mathit{ka} \cdot \mathit{ka} = \mathit{ka2}$ \\
115 & $\mathit{Aka2} = \mathit{A} \cdot \mathit{ka2}$ & $\mathit{A} \cdot \mathit{ka2} = \mathit{Aka2}$ \\
116 & $\mathit{R19} = \mathit{Aka2} + 1$ & $\mathit{Aka2} + 1 = \mathit{R19}$ \\
117 & $\mathit{L19} = \mathit{mu} \cdot \mathit{mu}$ & $\mathit{mu} \cdot \mathit{mu} = \mathit{L19}$ \\
118 & $\mathit{am1} = \mathit{a} - 1$ & $\mathit{am1} + 1 = \mathit{a}$ \\
119 & $\mathit{Dam1} = \mathit{Delta} \cdot \mathit{am1}$ & $\mathit{Delta} \cdot \mathit{am1} = \mathit{Dam1}$ \\
120 & $\mathit{R20} = \mathit{Dam1} + 5^{59}$ & $\mathit{Dam1} + 5^{59} = \mathit{R20}$ \\
\end{longtable}

The twenty free equality tests are
\[\begin{array}{@{}l@{}}
\mathit{L1}=\mathit{q2},\quad \mathit{b}=\mathit{bx},\quad \mathit{L2}=\mathit{R2},\\
\mathit{L3}=\mathit{B},\quad \mathit{S2}=\mathit{pack\_rhs},\quad \mathit{n}=\mathit{q16},\\
\mathit{r}=\mathit{R7},\quad \mathit{p}=\mathit{R8},\quad \mathit{L9}=\mathit{R9},\\
\mathit{L10}=\mathit{k2},\quad \mathit{k}=\mathit{R11},\quad \mathit{a}=\mathit{R12},\\
\mathit{c}=\mathit{R13},\quad \mathit{d}=\mathit{R14},\quad \mathit{L15}=\mathit{R15},\\
\mathit{L16}=\mathit{R16},\quad \mathit{L17}=\mathit{R17},\quad \mathit{mu}=\mathit{R18},\\
\mathit{L19}=\mathit{R19},\quad \mathit{ka}=\mathit{R20}
\end{array}\]


\normalsize
\section{The 114-instruction certificate}\label{app:114}

The four inputs are $x,Z,V,H$, with $(Z,V,H)$ the fixed admissible index
of Section~\ref{sec:short-masks}. The 33 positive witnesses are
\[
\begin{gathered}
 a,b,c,d,e,f,g,h,i,j,k,\ell,p,q,r,s,t,w,\\
 \alpha,\beta,\gamma,\eta,\theta,\lambda,\tau,\varphi,
 \kappa,\mu,\rho,\Delta,n,o,\zeta.
\end{gathered}
\]
The names $\mathit{l},\mathit{al},\mathit{beta},\mathit{ga},
\mathit{eta},\mathit{th},\mathit{la},\mathit{tau},\mathit{phi},
\mathit{ka},\mathit{mu},\mathit{rho},\mathit{Delta},\mathit{zeta}$
denote the corresponding Greek letters (with $\mathit l=\ell$).
All other names denote auxiliary integers constrained by the listed
arithmetic checks. In particular, the signed value $of-d$ is allowed;
positivity is required only of the declared parameters and witnesses.
The fixed exponent literal is $L=5^{64}$.

The table and all 21 equality tests are generated from
\file{round5\_1980\_certificate.py}. Its exact polynomial checks include
the two triangular residual identities for (E7) and (E17). The table
contains 63 multiplication checks and 51 addition checks. The optional
nine-step numeral construction in Section~\ref{sec:short-masks} raises
this to 123 while leaving the index parameters supplied. The proofs of
Corollary~\ref{thm:best114} establish the positive-domain and decoding
equivalences needed to interpret these arithmetic checks.

\small
% Generated by round5_1980_schedule_table.py; do not edit.
\begin{longtable}{@{}rll@{}}
\toprule
No. & Assignment & Addition/multiplication check\\
\midrule
\endfirsthead
\toprule
No. & Assignment & Addition/multiplication check\\
\midrule
\endhead
\bottomrule
\endfoot
1 & $\mathit{q2} = \mathit{q} \cdot \mathit{q}$ & $\mathit{q} \cdot \mathit{q} = \mathit{q2}$ \\
2 & $\mathit{q4} = \mathit{q2} \cdot \mathit{q2}$ & $\mathit{q2} \cdot \mathit{q2} = \mathit{q4}$ \\
3 & $\mathit{q8} = \mathit{q4} \cdot \mathit{q4}$ & $\mathit{q4} \cdot \mathit{q4} = \mathit{q8}$ \\
4 & $\mathit{b2} = \mathit{b} \cdot \mathit{b}$ & $\mathit{b} \cdot \mathit{b} = \mathit{b2}$ \\
5 & $\mathit{b4} = \mathit{b2} \cdot \mathit{b2}$ & $\mathit{b2} \cdot \mathit{b2} = \mathit{b4}$ \\
6 & $\mathit{B} = \mathit{H} \cdot \mathit{b4}$ & $\mathit{H} \cdot \mathit{b4} = \mathit{B}$ \\
7 & $\mathit{bx} = \mathit{x} + \mathit{beta}$ & $\mathit{x} + \mathit{beta} = \mathit{bx}$ \\
8 & $\mathit{L2} = \mathit{la} + \mathit{q2}$ & $\mathit{la} + \mathit{q2} = \mathit{L2}$ \\
9 & $\mathit{Lam} = \mathit{la} \cdot \mathit{B}$ & $\mathit{la} \cdot \mathit{B} = \mathit{Lam}$ \\
10 & $\mathit{R2} = \mathit{Lam} + 1$ & $\mathit{Lam} + 1 = \mathit{R2}$ \\
11 & $\mathit{L3} = \mathit{th} + \mathit{Z}$ & $\mathit{th} + \mathit{Z} = \mathit{L3}$ \\
12 & $\mathit{lq2} = \mathit{l} \cdot \mathit{q}$ & $\mathit{l} \cdot \mathit{q} = \mathit{lq2}$ \\
13 & $\mathit{S2} = \mathit{e} + \mathit{lq2}$ & $\mathit{e} + \mathit{lq2} = \mathit{S2}$ \\
14 & $\mathit{pack\_tth} = \mathit{t} \cdot \mathit{th}$ & $\mathit{t} \cdot \mathit{th} = \mathit{pack\_tth}$ \\
15 & $\mathit{pack\_rhs} = \mathit{V} + \mathit{pack\_tth}$ & $\mathit{V} + \mathit{pack\_tth} = \mathit{pack\_rhs}$ \\
16 & $\mathit{zl} = \mathit{Z} \cdot \mathit{la}$ & $\mathit{Z} \cdot \mathit{la} = \mathit{zl}$ \\
17 & $\mathit{twice\_e} = 2 \cdot \mathit{e}$ & $2 \cdot \mathit{e} = \mathit{twice\_e}$ \\
18 & $\mathit{t2} = \mathit{twice\_e} - \mathit{zl}$ & $\mathit{t2} + \mathit{zl} = \mathit{twice\_e}$ \\
19 & $\mathit{xb5} = \mathit{x} \cdot \mathit{B}$ & $\mathit{x} \cdot \mathit{B} = \mathit{xb5}$ \\
20 & $\mathit{c0} = \mathit{xb5} + 1$ & $\mathit{xb5} + 1 = \mathit{c0}$ \\
21 & $\mathit{C} = \mathit{c0} + \mathit{g}$ & $\mathit{c0} + \mathit{g} = \mathit{C}$ \\
22 & $\mathit{C2} = \mathit{C} \cdot \mathit{C}$ & $\mathit{C} \cdot \mathit{C} = \mathit{C2}$ \\
23 & $\mathit{C4} = \mathit{C2} \cdot \mathit{C2}$ & $\mathit{C2} \cdot \mathit{C2} = \mathit{C4}$ \\
24 & $\mathit{lC4} = \mathit{l} \cdot \mathit{C4}$ & $\mathit{l} \cdot \mathit{C4} = \mathit{lC4}$ \\
25 & $\mathit{bound\_sum} = \mathit{e} + \mathit{lC4}$ & $\mathit{e} + \mathit{lC4} = \mathit{bound\_sum}$ \\
26 & $\mathit{L1} = \mathit{bound\_sum} + \mathit{al}$ & $\mathit{bound\_sum} + \mathit{al} = \mathit{L1}$ \\
27 & $\mathit{P1} = \mathit{t2} \cdot \mathit{C4}$ & $\mathit{t2} \cdot \mathit{C4} = \mathit{P1}$ \\
28 & $\mathit{q4p1} = \mathit{q2} + 1$ & $\mathit{q2} + 1 = \mathit{q4p1}$ \\
29 & $\mathit{P2} = \mathit{Lam} \cdot \mathit{q4p1}$ & $\mathit{Lam} \cdot \mathit{q4p1} = \mathit{P2}$ \\
30 & $\mathit{S3} = \mathit{P1} + \mathit{P2}$ & $\mathit{P1} + \mathit{P2} = \mathit{S3}$ \\
31 & $\mathit{S3q4} = \mathit{S3} \cdot \mathit{q2}$ & $\mathit{S3} \cdot \mathit{q2} = \mathit{S3q4}$ \\
32 & $\mathit{Sin} = \mathit{S2} + \mathit{S3q4}$ & $\mathit{S2} + \mathit{S3q4} = \mathit{Sin}$ \\
33 & $\mathit{Sq3} = \mathit{Sin} \cdot \mathit{q}$ & $\mathit{Sin} \cdot \mathit{q} = \mathit{Sq3}$ \\
34 & $\mathit{S} = \mathit{g} + \mathit{Sq3}$ & $\mathit{g} + \mathit{Sq3} = \mathit{S}$ \\
35 & $\mathit{n2} = \mathit{n} \cdot \mathit{n}$ & $\mathit{n} \cdot \mathit{n} = \mathit{n2}$ \\
36 & $\mathit{n2n} = \mathit{n2} - \mathit{n}$ & $\mathit{n2n} + \mathit{n} = \mathit{n2}$ \\
37 & $\mathit{SA} = \mathit{S} \cdot \mathit{n2n}$ & $\mathit{S} \cdot \mathit{n2n} = \mathit{SA}$ \\
38 & $\mathit{Tcoef} = \mathit{L2} - \mathit{zl}$ & $\mathit{Tcoef} + \mathit{zl} = \mathit{L2}$ \\
39 & $\mathit{Tq3} = \mathit{Tcoef} \cdot \mathit{q}$ & $\mathit{Tcoef} \cdot \mathit{q} = \mathit{Tq3}$ \\
40 & $\mathit{bm1} = \mathit{b} - 1$ & $\mathit{bm1} + 1 = \mathit{b}$ \\
41 & $\mathit{lbm1} = \mathit{l} \cdot \mathit{bm1}$ & $\mathit{l} \cdot \mathit{bm1} = \mathit{lbm1}$ \\
42 & $\mathit{T2} = \mathit{Tq3} - \mathit{lbm1}$ & $\mathit{T2} + \mathit{lbm1} = \mathit{Tq3}$ \\
43 & $\mathit{b5m2} = \mathit{B} - 2$ & $\mathit{b5m2} + 2 = \mathit{B}$ \\
44 & $\mathit{b5m2q8} = \mathit{b5m2} \cdot \mathit{q4}$ & $\mathit{b5m2} \cdot \mathit{q4} = \mathit{b5m2q8}$ \\
45 & $\mathit{T} = \mathit{T2} + \mathit{b5m2q8}$ & $\mathit{T2} + \mathit{b5m2q8} = \mathit{T}$ \\
46 & $\mathit{n2m1} = \mathit{n2} - 1$ & $\mathit{n2m1} + 1 = \mathit{n2}$ \\
47 & $\mathit{TB} = \mathit{T} \cdot \mathit{n2m1}$ & $\mathit{T} \cdot \mathit{n2m1} = \mathit{TB}$ \\
48 & $\mathit{R7} = \mathit{SA} + \mathit{TB}$ & $\mathit{SA} + \mathit{TB} = \mathit{R7}$ \\
49 & $\mathit{wn2} = \mathit{w} \cdot \mathit{n2}$ & $\mathit{w} \cdot \mathit{n2} = \mathit{wn2}$ \\
50 & $\mathit{sn2} = \mathit{s} \cdot \mathit{n2}$ & $\mathit{s} \cdot \mathit{n2} = \mathit{sn2}$ \\
51 & $\mathit{rsn2} = \mathit{r} \cdot \mathit{sn2}$ & $\mathit{r} \cdot \mathit{sn2} = \mathit{rsn2}$ \\
52 & $\mathit{UM} = \mathit{wn2} \cdot \mathit{rsn2}$ & $\mathit{wn2} \cdot \mathit{rsn2} = \mathit{UM}$ \\
53 & $\mathit{R12} = \mathit{UM} + \mathit{rsn2}$ & $\mathit{UM} + \mathit{rsn2} = \mathit{R12}$ \\
54 & $\mathit{wsq} = \mathit{UM} \cdot \mathit{rsn2}$ & $\mathit{UM} \cdot \mathit{rsn2} = \mathit{wsq}$ \\
55 & $\mathit{R8} = 2 \cdot \mathit{wsq}$ & $2 \cdot \mathit{wsq} = \mathit{R8}$ \\
56 & $\mathit{p2} = \mathit{p} \cdot \mathit{p}$ & $\mathit{p} \cdot \mathit{p} = \mathit{p2}$ \\
57 & $\mathit{p2m1} = \mathit{p2} - 1$ & $\mathit{p2m1} + 1 = \mathit{p2}$ \\
58 & $\mathit{k2} = \mathit{k} \cdot \mathit{k}$ & $\mathit{k} \cdot \mathit{k} = \mathit{k2}$ \\
59 & $\mathit{P9} = \mathit{p2m1} \cdot \mathit{k2}$ & $\mathit{p2m1} \cdot \mathit{k2} = \mathit{P9}$ \\
60 & $\mathit{L9} = \mathit{P9} + 1$ & $\mathit{P9} + 1 = \mathit{L9}$ \\
61 & $\mathit{R9} = \mathit{tau} \cdot \mathit{tau}$ & $\mathit{tau} \cdot \mathit{tau} = \mathit{R9}$ \\
62 & $\mathit{ksn2} = \mathit{k} \cdot \mathit{sn2}$ & $\mathit{k} \cdot \mathit{sn2} = \mathit{ksn2}$ \\
63 & $\mathit{R10a} = \mathit{ksn2} + \mathit{eta}$ & $\mathit{ksn2} + \mathit{eta} = \mathit{R10a}$ \\
64 & $\mathit{R10b} = \mathit{eta} + \mathit{zeta}$ & $\mathit{eta} + \mathit{zeta} = \mathit{R10b}$ \\
65 & $\mathit{r1} = \mathit{r} + 1$ & $\mathit{r} + 1 = \mathit{r1}$ \\
66 & $\mathit{pm1} = \mathit{p} - 1$ & $\mathit{pm1} + 1 = \mathit{p}$ \\
67 & $\mathit{hpm1} = \mathit{h} \cdot \mathit{pm1}$ & $\mathit{h} \cdot \mathit{pm1} = \mathit{hpm1}$ \\
68 & $\mathit{R11} = \mathit{r1} + \mathit{hpm1}$ & $\mathit{r1} + \mathit{hpm1} = \mathit{R11}$ \\
69 & $\mathit{tr1} = \mathit{r1} + \mathit{r}$ & $\mathit{r1} + \mathit{r} = \mathit{tr1}$ \\
70 & $\mathit{R13} = \mathit{ka} + \mathit{phi}$ & $\mathit{ka} + \mathit{phi} = \mathit{R13}$ \\
71 & $\mathit{bw} = \mathit{b} \cdot \mathit{w}$ & $\mathit{b} \cdot \mathit{w} = \mathit{bw}$ \\
72 & $\mathit{am2} = \mathit{a} - 2$ & $\mathit{am2} + 2 = \mathit{a}$ \\
73 & $\mathit{cam2} = \mathit{c} \cdot \mathit{am2}$ & $\mathit{c} \cdot \mathit{am2} = \mathit{cam2}$ \\
74 & $\mathit{D1} = \mathit{bw} + \mathit{cam2}$ & $\mathit{bw} + \mathit{cam2} = \mathit{D1}$ \\
75 & $\mathit{a4} = 4 \cdot \mathit{a}$ & $4 \cdot \mathit{a} = \mathit{a4}$ \\
76 & $\mathit{a4m5} = \mathit{a4} - 5$ & $\mathit{a4m5} + 5 = \mathit{a4}$ \\
77 & $\mathit{gam} = \mathit{ga} \cdot \mathit{a4m5}$ & $\mathit{ga} \cdot \mathit{a4m5} = \mathit{gam}$ \\
78 & $\mathit{R14} = \mathit{D1} + \mathit{gam}$ & $\mathit{D1} + \mathit{gam} = \mathit{R14}$ \\
79 & $\mathit{am1} = \mathit{a} - 1$ & $\mathit{am1} + 1 = \mathit{a}$ \\
80 & $\mathit{ap1} = \mathit{a} + 1$ & $\mathit{a} + 1 = \mathit{ap1}$ \\
81 & $\mathit{A} = \mathit{am1} \cdot \mathit{ap1}$ & $\mathit{am1} \cdot \mathit{ap1} = \mathit{A}$ \\
82 & $\mathit{c2} = \mathit{c} \cdot \mathit{c}$ & $\mathit{c} \cdot \mathit{c} = \mathit{c2}$ \\
83 & $\mathit{Ac2} = \mathit{A} \cdot \mathit{c2}$ & $\mathit{A} \cdot \mathit{c2} = \mathit{Ac2}$ \\
84 & $\mathit{R15} = \mathit{Ac2} + 1$ & $\mathit{Ac2} + 1 = \mathit{R15}$ \\
85 & $\mathit{L15} = \mathit{d} \cdot \mathit{d}$ & $\mathit{d} \cdot \mathit{d} = \mathit{L15}$ \\
86 & $\mathit{ic2} = \mathit{i} \cdot \mathit{c2}$ & $\mathit{i} \cdot \mathit{c2} = \mathit{ic2}$ \\
87 & $\mathit{ic22} = \mathit{ic2} \cdot \mathit{ic2}$ & $\mathit{ic2} \cdot \mathit{ic2} = \mathit{ic22}$ \\
88 & $\mathit{AE} = \mathit{A} \cdot \mathit{ic22}$ & $\mathit{A} \cdot \mathit{ic22} = \mathit{AE}$ \\
89 & $\mathit{R16} = \mathit{AE} + 1$ & $\mathit{AE} + 1 = \mathit{R16}$ \\
90 & $\mathit{L16} = \mathit{f} \cdot \mathit{f}$ & $\mathit{f} \cdot \mathit{f} = \mathit{L16}$ \\
91 & $\mathit{of} = \mathit{o} \cdot \mathit{f}$ & $\mathit{o} \cdot \mathit{f} = \mathit{of}$ \\
92 & $\mathit{dof} = \mathit{of} - \mathit{d}$ & $\mathit{dof} + \mathit{d} = \mathit{of}$ \\
93 & $\mathit{L17} = \mathit{dof} \cdot \mathit{dof}$ & $\mathit{dof} \cdot \mathit{dof} = \mathit{L17}$ \\
94 & $\mathit{Gminus1} = \mathit{ap1} \cdot \mathit{AE}$ & $\mathit{ap1} \cdot \mathit{AE} = \mathit{Gminus1}$ \\
95 & $\mathit{Gplus1} = \mathit{Gminus1} + 2$ & $\mathit{Gminus1} + 2 = \mathit{Gplus1}$ \\
96 & $\mathit{G2m1} = \mathit{Gminus1} \cdot \mathit{Gplus1}$ & $\mathit{Gminus1} \cdot \mathit{Gplus1} = \mathit{G2m1}$ \\
97 & $\mathit{jc} = \mathit{j} \cdot \mathit{c}$ & $\mathit{j} \cdot \mathit{c} = \mathit{jc}$ \\
98 & $\mathit{H17} = \mathit{tr1} + \mathit{jc}$ & $\mathit{tr1} + \mathit{jc} = \mathit{H17}$ \\
99 & $\mathit{H2} = \mathit{H17} \cdot \mathit{H17}$ & $\mathit{H17} \cdot \mathit{H17} = \mathit{H2}$ \\
100 & $\mathit{P17} = \mathit{G2m1} \cdot \mathit{H2}$ & $\mathit{G2m1} \cdot \mathit{H2} = \mathit{P17}$ \\
101 & $\mathit{R17} = \mathit{P17} + 1$ & $\mathit{P17} + 1 = \mathit{R17}$ \\
102 & $\mathit{amb} = \mathit{a} - \mathit{B}$ & $\mathit{amb} + \mathit{B} = \mathit{a}$ \\
103 & $\mathit{kamb} = \mathit{ka} \cdot \mathit{amb}$ & $\mathit{ka} \cdot \mathit{amb} = \mathit{kamb}$ \\
104 & $\mathit{qk} = \mathit{q} + \mathit{kamb}$ & $\mathit{q} + \mathit{kamb} = \mathit{qk}$ \\
105 & $\mathit{amb2} = \mathit{amb} \cdot \mathit{amb}$ & $\mathit{amb} \cdot \mathit{amb} = \mathit{amb2}$ \\
106 & $\mathit{Mod} = \mathit{A} - \mathit{amb2}$ & $\mathit{Mod} + \mathit{amb2} = \mathit{A}$ \\
107 & $\mathit{rM} = \mathit{rho} \cdot \mathit{Mod}$ & $\mathit{rho} \cdot \mathit{Mod} = \mathit{rM}$ \\
108 & $\mathit{R18} = \mathit{qk} + \mathit{rM}$ & $\mathit{qk} + \mathit{rM} = \mathit{R18}$ \\
109 & $\mathit{ka2} = \mathit{ka} \cdot \mathit{ka}$ & $\mathit{ka} \cdot \mathit{ka} = \mathit{ka2}$ \\
110 & $\mathit{Aka2} = \mathit{A} \cdot \mathit{ka2}$ & $\mathit{A} \cdot \mathit{ka2} = \mathit{Aka2}$ \\
111 & $\mathit{R19} = \mathit{Aka2} + 1$ & $\mathit{Aka2} + 1 = \mathit{R19}$ \\
112 & $\mathit{L19} = \mathit{mu} \cdot \mathit{mu}$ & $\mathit{mu} \cdot \mathit{mu} = \mathit{L19}$ \\
113 & $\mathit{Dam1} = \mathit{Delta} \cdot \mathit{am1}$ & $\mathit{Delta} \cdot \mathit{am1} = \mathit{Dam1}$ \\
114 & $\mathit{R20} = \mathit{Dam1} + 5^{64}$ & $\mathit{Dam1} + 5^{64} = \mathit{R20}$ \\
\end{longtable}

The twenty-one free equality tests are
\[\begin{array}{@{}l@{}}
\mathit{L1}=\mathit{q},\quad \mathit{b}=\mathit{bx},\quad \mathit{L2}=\mathit{R2},\\
\mathit{L3}=\mathit{B},\quad \mathit{S2}=\mathit{pack\_rhs},\quad \mathit{n}=\mathit{q8},\\
\mathit{r}=\mathit{R7},\quad \mathit{p}=\mathit{R8},\quad \mathit{L9}=\mathit{R9},\\
\mathit{c}=\mathit{R10a},\quad \mathit{k}=\mathit{R10b},\quad \mathit{k}=\mathit{R11},\\
\mathit{a}=\mathit{R12},\quad \mathit{c}=\mathit{R13},\quad \mathit{d}=\mathit{R14},\\
\mathit{L15}=\mathit{R15},\quad \mathit{L16}=\mathit{R16},\quad \mathit{L17}=\mathit{R17},\\
\mathit{mu}=\mathit{R18},\quad \mathit{L19}=\mathit{R19},\quad \mathit{ka}=\mathit{R20}
\end{array}\]


\clearpage
\normalsize
\section{The 112-instruction certificate}\label{app:112}

The fixed admissible index $(Z,V,H)$ is now the direct quadratic index
of Section~\ref{sec:quadratic-masks}. The positive system unknowns are
the 33 listed in Appendix~\ref{app:114} with $p$ removed, leaving 32.
The notation and auxiliary integer domains otherwise remain the same.
In particular, $B=Hb^2$ and the code square $C^2$ are computed explicitly.
The historical name $\mathit{R8}$ now denotes the computed abbreviation
$X=2wn^2(rs n^2)^2$, with no additional equation defining a witness $p$.
The shifted Pell parameter $P=X+1$ is used only in the proof; the table
checks its coefficient as $X(X+2)$ without computing $P$.

The table is generated from \file{round6\_1980\_certificate.py}, which
verifies all 112 primitive statements, the 20 equation residuals, both
triangular corrections, and the finite Sidon and support inequalities.
There are 62 multiplication checks and 50 addition checks. The literal
$5^{16}$ and the other fixed numerals can be generated in seven additional
instructions, giving 119 with that requirement. Theorem~\ref{thm:best112}
supplies the mathematical equivalence for this fixed encoding.

\small
% Generated by round6_1980_schedule_table.py; do not edit.
\begin{longtable}{@{}rll@{}}
\toprule
No. & Assignment & Addition/multiplication check\\
\midrule
\endfirsthead
\toprule
No. & Assignment & Addition/multiplication check\\
\midrule
\endhead
\bottomrule
\endfoot
1 & $\mathit{q2} = \mathit{q} \cdot \mathit{q}$ & $\mathit{q} \cdot \mathit{q} = \mathit{q2}$ \\
2 & $\mathit{q4} = \mathit{q2} \cdot \mathit{q2}$ & $\mathit{q2} \cdot \mathit{q2} = \mathit{q4}$ \\
3 & $\mathit{q8} = \mathit{q4} \cdot \mathit{q4}$ & $\mathit{q4} \cdot \mathit{q4} = \mathit{q8}$ \\
4 & $\mathit{b2} = \mathit{b} \cdot \mathit{b}$ & $\mathit{b} \cdot \mathit{b} = \mathit{b2}$ \\
5 & $\mathit{B} = \mathit{H} \cdot \mathit{b2}$ & $\mathit{H} \cdot \mathit{b2} = \mathit{B}$ \\
6 & $\mathit{bx} = \mathit{x} + \mathit{beta}$ & $\mathit{x} + \mathit{beta} = \mathit{bx}$ \\
7 & $\mathit{L2} = \mathit{la} + \mathit{q2}$ & $\mathit{la} + \mathit{q2} = \mathit{L2}$ \\
8 & $\mathit{Lam} = \mathit{la} \cdot \mathit{B}$ & $\mathit{la} \cdot \mathit{B} = \mathit{Lam}$ \\
9 & $\mathit{R2} = \mathit{Lam} + 1$ & $\mathit{Lam} + 1 = \mathit{R2}$ \\
10 & $\mathit{L3} = \mathit{th} + \mathit{Z}$ & $\mathit{th} + \mathit{Z} = \mathit{L3}$ \\
11 & $\mathit{lq2} = \mathit{l} \cdot \mathit{q}$ & $\mathit{l} \cdot \mathit{q} = \mathit{lq2}$ \\
12 & $\mathit{S2} = \mathit{e} + \mathit{lq2}$ & $\mathit{e} + \mathit{lq2} = \mathit{S2}$ \\
13 & $\mathit{pack\_tth} = \mathit{t} \cdot \mathit{th}$ & $\mathit{t} \cdot \mathit{th} = \mathit{pack\_tth}$ \\
14 & $\mathit{pack\_rhs} = \mathit{V} + \mathit{pack\_tth}$ & $\mathit{V} + \mathit{pack\_tth} = \mathit{pack\_rhs}$ \\
15 & $\mathit{zl} = \mathit{Z} \cdot \mathit{la}$ & $\mathit{Z} \cdot \mathit{la} = \mathit{zl}$ \\
16 & $\mathit{twice\_e} = 2 \cdot \mathit{e}$ & $2 \cdot \mathit{e} = \mathit{twice\_e}$ \\
17 & $\mathit{t2} = \mathit{twice\_e} - \mathit{zl}$ & $\mathit{t2} + \mathit{zl} = \mathit{twice\_e}$ \\
18 & $\mathit{xb5} = \mathit{x} \cdot \mathit{B}$ & $\mathit{x} \cdot \mathit{B} = \mathit{xb5}$ \\
19 & $\mathit{c0} = \mathit{xb5} + 1$ & $\mathit{xb5} + 1 = \mathit{c0}$ \\
20 & $\mathit{C} = \mathit{c0} + \mathit{g}$ & $\mathit{c0} + \mathit{g} = \mathit{C}$ \\
21 & $\mathit{C2} = \mathit{C} \cdot \mathit{C}$ & $\mathit{C} \cdot \mathit{C} = \mathit{C2}$ \\
22 & $\mathit{lC4} = \mathit{l} \cdot \mathit{C2}$ & $\mathit{l} \cdot \mathit{C2} = \mathit{lC4}$ \\
23 & $\mathit{bound\_sum} = \mathit{e} + \mathit{lC4}$ & $\mathit{e} + \mathit{lC4} = \mathit{bound\_sum}$ \\
24 & $\mathit{L1} = \mathit{bound\_sum} + \mathit{al}$ & $\mathit{bound\_sum} + \mathit{al} = \mathit{L1}$ \\
25 & $\mathit{P1} = \mathit{t2} \cdot \mathit{C2}$ & $\mathit{t2} \cdot \mathit{C2} = \mathit{P1}$ \\
26 & $\mathit{q4p1} = \mathit{q2} + 1$ & $\mathit{q2} + 1 = \mathit{q4p1}$ \\
27 & $\mathit{P2} = \mathit{Lam} \cdot \mathit{q4p1}$ & $\mathit{Lam} \cdot \mathit{q4p1} = \mathit{P2}$ \\
28 & $\mathit{S3} = \mathit{P1} + \mathit{P2}$ & $\mathit{P1} + \mathit{P2} = \mathit{S3}$ \\
29 & $\mathit{S3q4} = \mathit{S3} \cdot \mathit{q2}$ & $\mathit{S3} \cdot \mathit{q2} = \mathit{S3q4}$ \\
30 & $\mathit{Sin} = \mathit{S2} + \mathit{S3q4}$ & $\mathit{S2} + \mathit{S3q4} = \mathit{Sin}$ \\
31 & $\mathit{Sq3} = \mathit{Sin} \cdot \mathit{q}$ & $\mathit{Sin} \cdot \mathit{q} = \mathit{Sq3}$ \\
32 & $\mathit{S} = \mathit{g} + \mathit{Sq3}$ & $\mathit{g} + \mathit{Sq3} = \mathit{S}$ \\
33 & $\mathit{n2} = \mathit{n} \cdot \mathit{n}$ & $\mathit{n} \cdot \mathit{n} = \mathit{n2}$ \\
34 & $\mathit{n2n} = \mathit{n2} - \mathit{n}$ & $\mathit{n2n} + \mathit{n} = \mathit{n2}$ \\
35 & $\mathit{SA} = \mathit{S} \cdot \mathit{n2n}$ & $\mathit{S} \cdot \mathit{n2n} = \mathit{SA}$ \\
36 & $\mathit{Tcoef} = \mathit{L2} - \mathit{zl}$ & $\mathit{Tcoef} + \mathit{zl} = \mathit{L2}$ \\
37 & $\mathit{Tq3} = \mathit{Tcoef} \cdot \mathit{q}$ & $\mathit{Tcoef} \cdot \mathit{q} = \mathit{Tq3}$ \\
38 & $\mathit{bm1} = \mathit{b} - 1$ & $\mathit{bm1} + 1 = \mathit{b}$ \\
39 & $\mathit{lbm1} = \mathit{l} \cdot \mathit{bm1}$ & $\mathit{l} \cdot \mathit{bm1} = \mathit{lbm1}$ \\
40 & $\mathit{T2} = \mathit{Tq3} - \mathit{lbm1}$ & $\mathit{T2} + \mathit{lbm1} = \mathit{Tq3}$ \\
41 & $\mathit{b5m2} = \mathit{B} - 2$ & $\mathit{b5m2} + 2 = \mathit{B}$ \\
42 & $\mathit{packed\_target\_shift} = \mathit{lq2} \cdot \mathit{q2}$ & $\mathit{lq2} \cdot \mathit{q2} = \mathit{packed\_target\_shift}$ \\
43 & $\mathit{mask\_term} = \mathit{b5m2} \cdot \mathit{packed\_target\_shift}$ & $\mathit{b5m2} \cdot \mathit{packed\_target\_shift} = \mathit{mask\_term}$ \\
44 & $\mathit{T} = \mathit{T2} + \mathit{mask\_term}$ & $\mathit{T2} + \mathit{mask\_term} = \mathit{T}$ \\
45 & $\mathit{n2m1} = \mathit{n2} - 1$ & $\mathit{n2m1} + 1 = \mathit{n2}$ \\
46 & $\mathit{TB} = \mathit{T} \cdot \mathit{n2m1}$ & $\mathit{T} \cdot \mathit{n2m1} = \mathit{TB}$ \\
47 & $\mathit{R7} = \mathit{SA} + \mathit{TB}$ & $\mathit{SA} + \mathit{TB} = \mathit{R7}$ \\
48 & $\mathit{wn2} = \mathit{w} \cdot \mathit{n2}$ & $\mathit{w} \cdot \mathit{n2} = \mathit{wn2}$ \\
49 & $\mathit{sn2} = \mathit{s} \cdot \mathit{n2}$ & $\mathit{s} \cdot \mathit{n2} = \mathit{sn2}$ \\
50 & $\mathit{rsn2} = \mathit{r} \cdot \mathit{sn2}$ & $\mathit{r} \cdot \mathit{sn2} = \mathit{rsn2}$ \\
51 & $\mathit{UM} = \mathit{wn2} \cdot \mathit{rsn2}$ & $\mathit{wn2} \cdot \mathit{rsn2} = \mathit{UM}$ \\
52 & $\mathit{R12} = \mathit{UM} + \mathit{rsn2}$ & $\mathit{UM} + \mathit{rsn2} = \mathit{R12}$ \\
53 & $\mathit{wsq} = \mathit{UM} \cdot \mathit{rsn2}$ & $\mathit{UM} \cdot \mathit{rsn2} = \mathit{wsq}$ \\
54 & $\mathit{R8} = 2 \cdot \mathit{wsq}$ & $2 \cdot \mathit{wsq} = \mathit{R8}$ \\
55 & $\mathit{Pplus1} = \mathit{R8} + 2$ & $\mathit{R8} + 2 = \mathit{Pplus1}$ \\
56 & $\mathit{p2m1} = \mathit{R8} \cdot \mathit{Pplus1}$ & $\mathit{R8} \cdot \mathit{Pplus1} = \mathit{p2m1}$ \\
57 & $\mathit{k2} = \mathit{k} \cdot \mathit{k}$ & $\mathit{k} \cdot \mathit{k} = \mathit{k2}$ \\
58 & $\mathit{P9} = \mathit{p2m1} \cdot \mathit{k2}$ & $\mathit{p2m1} \cdot \mathit{k2} = \mathit{P9}$ \\
59 & $\mathit{L9} = \mathit{P9} + 1$ & $\mathit{P9} + 1 = \mathit{L9}$ \\
60 & $\mathit{R9} = \mathit{tau} \cdot \mathit{tau}$ & $\mathit{tau} \cdot \mathit{tau} = \mathit{R9}$ \\
61 & $\mathit{ksn2} = \mathit{k} \cdot \mathit{sn2}$ & $\mathit{k} \cdot \mathit{sn2} = \mathit{ksn2}$ \\
62 & $\mathit{R10a} = \mathit{ksn2} + \mathit{eta}$ & $\mathit{ksn2} + \mathit{eta} = \mathit{R10a}$ \\
63 & $\mathit{R10b} = \mathit{eta} + \mathit{zeta}$ & $\mathit{eta} + \mathit{zeta} = \mathit{R10b}$ \\
64 & $\mathit{r1} = \mathit{r} + 1$ & $\mathit{r} + 1 = \mathit{r1}$ \\
65 & $\mathit{hpm1} = \mathit{h} \cdot \mathit{R8}$ & $\mathit{h} \cdot \mathit{R8} = \mathit{hpm1}$ \\
66 & $\mathit{R11} = \mathit{r1} + \mathit{hpm1}$ & $\mathit{r1} + \mathit{hpm1} = \mathit{R11}$ \\
67 & $\mathit{tr1} = \mathit{r1} + \mathit{r}$ & $\mathit{r1} + \mathit{r} = \mathit{tr1}$ \\
68 & $\mathit{R13} = \mathit{ka} + \mathit{phi}$ & $\mathit{ka} + \mathit{phi} = \mathit{R13}$ \\
69 & $\mathit{bw} = \mathit{b} \cdot \mathit{w}$ & $\mathit{b} \cdot \mathit{w} = \mathit{bw}$ \\
70 & $\mathit{am2} = \mathit{a} - 2$ & $\mathit{am2} + 2 = \mathit{a}$ \\
71 & $\mathit{cam2} = \mathit{c} \cdot \mathit{am2}$ & $\mathit{c} \cdot \mathit{am2} = \mathit{cam2}$ \\
72 & $\mathit{D1} = \mathit{bw} + \mathit{cam2}$ & $\mathit{bw} + \mathit{cam2} = \mathit{D1}$ \\
73 & $\mathit{a4} = 4 \cdot \mathit{a}$ & $4 \cdot \mathit{a} = \mathit{a4}$ \\
74 & $\mathit{a4m5} = \mathit{a4} - 5$ & $\mathit{a4m5} + 5 = \mathit{a4}$ \\
75 & $\mathit{gam} = \mathit{ga} \cdot \mathit{a4m5}$ & $\mathit{ga} \cdot \mathit{a4m5} = \mathit{gam}$ \\
76 & $\mathit{R14} = \mathit{D1} + \mathit{gam}$ & $\mathit{D1} + \mathit{gam} = \mathit{R14}$ \\
77 & $\mathit{am1} = \mathit{a} - 1$ & $\mathit{am1} + 1 = \mathit{a}$ \\
78 & $\mathit{ap1} = \mathit{a} + 1$ & $\mathit{a} + 1 = \mathit{ap1}$ \\
79 & $\mathit{A} = \mathit{am1} \cdot \mathit{ap1}$ & $\mathit{am1} \cdot \mathit{ap1} = \mathit{A}$ \\
80 & $\mathit{c2} = \mathit{c} \cdot \mathit{c}$ & $\mathit{c} \cdot \mathit{c} = \mathit{c2}$ \\
81 & $\mathit{Ac2} = \mathit{A} \cdot \mathit{c2}$ & $\mathit{A} \cdot \mathit{c2} = \mathit{Ac2}$ \\
82 & $\mathit{R15} = \mathit{Ac2} + 1$ & $\mathit{Ac2} + 1 = \mathit{R15}$ \\
83 & $\mathit{L15} = \mathit{d} \cdot \mathit{d}$ & $\mathit{d} \cdot \mathit{d} = \mathit{L15}$ \\
84 & $\mathit{ic2} = \mathit{i} \cdot \mathit{c2}$ & $\mathit{i} \cdot \mathit{c2} = \mathit{ic2}$ \\
85 & $\mathit{ic22} = \mathit{ic2} \cdot \mathit{ic2}$ & $\mathit{ic2} \cdot \mathit{ic2} = \mathit{ic22}$ \\
86 & $\mathit{AE} = \mathit{A} \cdot \mathit{ic22}$ & $\mathit{A} \cdot \mathit{ic22} = \mathit{AE}$ \\
87 & $\mathit{R16} = \mathit{AE} + 1$ & $\mathit{AE} + 1 = \mathit{R16}$ \\
88 & $\mathit{L16} = \mathit{f} \cdot \mathit{f}$ & $\mathit{f} \cdot \mathit{f} = \mathit{L16}$ \\
89 & $\mathit{of} = \mathit{o} \cdot \mathit{f}$ & $\mathit{o} \cdot \mathit{f} = \mathit{of}$ \\
90 & $\mathit{dof} = \mathit{of} - \mathit{d}$ & $\mathit{dof} + \mathit{d} = \mathit{of}$ \\
91 & $\mathit{L17} = \mathit{dof} \cdot \mathit{dof}$ & $\mathit{dof} \cdot \mathit{dof} = \mathit{L17}$ \\
92 & $\mathit{Gminus1} = \mathit{ap1} \cdot \mathit{AE}$ & $\mathit{ap1} \cdot \mathit{AE} = \mathit{Gminus1}$ \\
93 & $\mathit{Gplus1} = \mathit{Gminus1} + 2$ & $\mathit{Gminus1} + 2 = \mathit{Gplus1}$ \\
94 & $\mathit{G2m1} = \mathit{Gminus1} \cdot \mathit{Gplus1}$ & $\mathit{Gminus1} \cdot \mathit{Gplus1} = \mathit{G2m1}$ \\
95 & $\mathit{jc} = \mathit{j} \cdot \mathit{c}$ & $\mathit{j} \cdot \mathit{c} = \mathit{jc}$ \\
96 & $\mathit{H17} = \mathit{tr1} + \mathit{jc}$ & $\mathit{tr1} + \mathit{jc} = \mathit{H17}$ \\
97 & $\mathit{H2} = \mathit{H17} \cdot \mathit{H17}$ & $\mathit{H17} \cdot \mathit{H17} = \mathit{H2}$ \\
98 & $\mathit{P17} = \mathit{G2m1} \cdot \mathit{H2}$ & $\mathit{G2m1} \cdot \mathit{H2} = \mathit{P17}$ \\
99 & $\mathit{R17} = \mathit{P17} + 1$ & $\mathit{P17} + 1 = \mathit{R17}$ \\
100 & $\mathit{amb} = \mathit{a} - \mathit{B}$ & $\mathit{amb} + \mathit{B} = \mathit{a}$ \\
101 & $\mathit{kamb} = \mathit{ka} \cdot \mathit{amb}$ & $\mathit{ka} \cdot \mathit{amb} = \mathit{kamb}$ \\
102 & $\mathit{qk} = \mathit{q} + \mathit{kamb}$ & $\mathit{q} + \mathit{kamb} = \mathit{qk}$ \\
103 & $\mathit{amb2} = \mathit{amb} \cdot \mathit{amb}$ & $\mathit{amb} \cdot \mathit{amb} = \mathit{amb2}$ \\
104 & $\mathit{Mod} = \mathit{A} - \mathit{amb2}$ & $\mathit{Mod} + \mathit{amb2} = \mathit{A}$ \\
105 & $\mathit{rM} = \mathit{rho} \cdot \mathit{Mod}$ & $\mathit{rho} \cdot \mathit{Mod} = \mathit{rM}$ \\
106 & $\mathit{R18} = \mathit{qk} + \mathit{rM}$ & $\mathit{qk} + \mathit{rM} = \mathit{R18}$ \\
107 & $\mathit{ka2} = \mathit{ka} \cdot \mathit{ka}$ & $\mathit{ka} \cdot \mathit{ka} = \mathit{ka2}$ \\
108 & $\mathit{Aka2} = \mathit{A} \cdot \mathit{ka2}$ & $\mathit{A} \cdot \mathit{ka2} = \mathit{Aka2}$ \\
109 & $\mathit{R19} = \mathit{Aka2} + 1$ & $\mathit{Aka2} + 1 = \mathit{R19}$ \\
110 & $\mathit{L19} = \mathit{mu} \cdot \mathit{mu}$ & $\mathit{mu} \cdot \mathit{mu} = \mathit{L19}$ \\
111 & $\mathit{Dam1} = \mathit{Delta} \cdot \mathit{am1}$ & $\mathit{Delta} \cdot \mathit{am1} = \mathit{Dam1}$ \\
112 & $\mathit{R20} = \mathit{Dam1} + 5^{16}$ & $\mathit{Dam1} + 5^{16} = \mathit{R20}$ \\
\end{longtable}

The twenty free equality tests are
\[\begin{array}{@{}l@{}}
\mathit{L1}=\mathit{q},\quad \mathit{b}=\mathit{bx},\quad \mathit{L2}=\mathit{R2},\\
\mathit{L3}=\mathit{B},\quad \mathit{S2}=\mathit{pack\_rhs},\quad \mathit{n}=\mathit{q8},\\
\mathit{r}=\mathit{R7},\quad \mathit{L9}=\mathit{R9},\quad \mathit{c}=\mathit{R10a},\\
\mathit{k}=\mathit{R10b},\quad \mathit{k}=\mathit{R11},\quad \mathit{a}=\mathit{R12},\\
\mathit{c}=\mathit{R13},\quad \mathit{d}=\mathit{R14},\quad \mathit{L15}=\mathit{R15},\\
\mathit{L16}=\mathit{R16},\quad \mathit{L17}=\mathit{R17},\quad \mathit{mu}=\mathit{R18},\\
\mathit{L19}=\mathit{R19},\quad \mathit{ka}=\mathit{R20}
\end{array}\]


\clearpage
\normalsize
\section{The 111-instruction certificate}\label{app:111}

The fixed index $(Z,V,H)$ is now that of
Section~\ref{sec:reversed-packing}. Add the positive unknown $\sigma$ to
the 32 witnesses of Appendix~\ref{app:112}; the free equality
$\mathit{S3}=\sigma$ enforces the sign of the computed third block.
All other auxiliary names retain their integer domains. In particular,
positivity is imposed on the 33 declared witnesses, not on every
intermediate expression.

The code register is $\mathit{S2}=\ell+eq$, and the size test is
$\mathit{S2}+C^2+\alpha=q^2$. The table uses the already available $q^2$
and $q^4$ for the new packing shifts. It is generated by
\file{round7\_1980\_schedule\_table.py} from the complete
\file{round7\_1980\_certificate.py} checker. Every primitive and all
21 source residuals are checked exactly, including the changed triangular
identity for (E7). The table contains 61 multiplications and 50 additions;
the same seven numeral instructions give 118 when required.
Theorem~\ref{thm:best111} supplies the encoding equivalence and positive
witness construction. The 112-instruction table remains unchanged above.

\small
% Generated by round7_1980_schedule_table.py; do not edit.
\begin{longtable}{@{}rll@{}}
\toprule
No. & Assignment & Addition/multiplication check\\
\midrule
\endfirsthead
\toprule
No. & Assignment & Addition/multiplication check\\
\midrule
\endhead
\bottomrule
\endfoot
1 & $\mathit{q2} = \mathit{q} \cdot \mathit{q}$ & $\mathit{q} \cdot \mathit{q} = \mathit{q2}$ \\
2 & $\mathit{q4} = \mathit{q2} \cdot \mathit{q2}$ & $\mathit{q2} \cdot \mathit{q2} = \mathit{q4}$ \\
3 & $\mathit{q8} = \mathit{q4} \cdot \mathit{q4}$ & $\mathit{q4} \cdot \mathit{q4} = \mathit{q8}$ \\
4 & $\mathit{b2} = \mathit{b} \cdot \mathit{b}$ & $\mathit{b} \cdot \mathit{b} = \mathit{b2}$ \\
5 & $\mathit{B} = \mathit{H} \cdot \mathit{b2}$ & $\mathit{H} \cdot \mathit{b2} = \mathit{B}$ \\
6 & $\mathit{bx} = \mathit{x} + \mathit{beta}$ & $\mathit{x} + \mathit{beta} = \mathit{bx}$ \\
7 & $\mathit{L2} = \mathit{la} + \mathit{q2}$ & $\mathit{la} + \mathit{q2} = \mathit{L2}$ \\
8 & $\mathit{Lam} = \mathit{la} \cdot \mathit{B}$ & $\mathit{la} \cdot \mathit{B} = \mathit{Lam}$ \\
9 & $\mathit{R2} = \mathit{Lam} + 1$ & $\mathit{Lam} + 1 = \mathit{R2}$ \\
10 & $\mathit{L3} = \mathit{th} + \mathit{Z}$ & $\mathit{th} + \mathit{Z} = \mathit{L3}$ \\
11 & $\mathit{eq} = \mathit{e} \cdot \mathit{q}$ & $\mathit{e} \cdot \mathit{q} = \mathit{eq}$ \\
12 & $\mathit{S2} = \mathit{l} + \mathit{eq}$ & $\mathit{l} + \mathit{eq} = \mathit{S2}$ \\
13 & $\mathit{pack\_tth} = \mathit{t} \cdot \mathit{th}$ & $\mathit{t} \cdot \mathit{th} = \mathit{pack\_tth}$ \\
14 & $\mathit{pack\_rhs} = \mathit{V} + \mathit{pack\_tth}$ & $\mathit{V} + \mathit{pack\_tth} = \mathit{pack\_rhs}$ \\
15 & $\mathit{zl} = \mathit{Z} \cdot \mathit{la}$ & $\mathit{Z} \cdot \mathit{la} = \mathit{zl}$ \\
16 & $\mathit{twice\_e} = 2 \cdot \mathit{e}$ & $2 \cdot \mathit{e} = \mathit{twice\_e}$ \\
17 & $\mathit{t2} = \mathit{twice\_e} - \mathit{zl}$ & $\mathit{t2} + \mathit{zl} = \mathit{twice\_e}$ \\
18 & $\mathit{xb5} = \mathit{x} \cdot \mathit{B}$ & $\mathit{x} \cdot \mathit{B} = \mathit{xb5}$ \\
19 & $\mathit{c0} = \mathit{xb5} + 1$ & $\mathit{xb5} + 1 = \mathit{c0}$ \\
20 & $\mathit{C} = \mathit{c0} + \mathit{g}$ & $\mathit{c0} + \mathit{g} = \mathit{C}$ \\
21 & $\mathit{C2} = \mathit{C} \cdot \mathit{C}$ & $\mathit{C} \cdot \mathit{C} = \mathit{C2}$ \\
22 & $\mathit{bound\_sum} = \mathit{S2} + \mathit{C2}$ & $\mathit{S2} + \mathit{C2} = \mathit{bound\_sum}$ \\
23 & $\mathit{L1} = \mathit{bound\_sum} + \mathit{al}$ & $\mathit{bound\_sum} + \mathit{al} = \mathit{L1}$ \\
24 & $\mathit{P1} = \mathit{t2} \cdot \mathit{C2}$ & $\mathit{t2} \cdot \mathit{C2} = \mathit{P1}$ \\
25 & $\mathit{qp1} = \mathit{q} + 1$ & $\mathit{q} + 1 = \mathit{qp1}$ \\
26 & $\mathit{P2} = \mathit{Lam} \cdot \mathit{qp1}$ & $\mathit{Lam} \cdot \mathit{qp1} = \mathit{P2}$ \\
27 & $\mathit{S3} = \mathit{P1} + \mathit{P2}$ & $\mathit{P1} + \mathit{P2} = \mathit{S3}$ \\
28 & $\mathit{S3q4} = \mathit{S3} \cdot \mathit{q2}$ & $\mathit{S3} \cdot \mathit{q2} = \mathit{S3q4}$ \\
29 & $\mathit{Sin} = \mathit{S2} + \mathit{S3q4}$ & $\mathit{S2} + \mathit{S3q4} = \mathit{Sin}$ \\
30 & $\mathit{Sq4} = \mathit{Sin} \cdot \mathit{q2}$ & $\mathit{Sin} \cdot \mathit{q2} = \mathit{Sq4}$ \\
31 & $\mathit{S} = \mathit{g} + \mathit{Sq4}$ & $\mathit{g} + \mathit{Sq4} = \mathit{S}$ \\
32 & $\mathit{n2} = \mathit{n} \cdot \mathit{n}$ & $\mathit{n} \cdot \mathit{n} = \mathit{n2}$ \\
33 & $\mathit{n2n} = \mathit{n2} - \mathit{n}$ & $\mathit{n2n} + \mathit{n} = \mathit{n2}$ \\
34 & $\mathit{SA} = \mathit{S} \cdot \mathit{n2n}$ & $\mathit{S} \cdot \mathit{n2n} = \mathit{SA}$ \\
35 & $\mathit{Tcoef} = \mathit{L2} - \mathit{zl}$ & $\mathit{Tcoef} + \mathit{zl} = \mathit{L2}$ \\
36 & $\mathit{Tq4} = \mathit{Tcoef} \cdot \mathit{q2}$ & $\mathit{Tcoef} \cdot \mathit{q2} = \mathit{Tq4}$ \\
37 & $\mathit{bm1} = \mathit{b} - 1$ & $\mathit{bm1} + 1 = \mathit{b}$ \\
38 & $\mathit{lbm1} = \mathit{l} \cdot \mathit{bm1}$ & $\mathit{l} \cdot \mathit{bm1} = \mathit{lbm1}$ \\
39 & $\mathit{T2} = \mathit{Tq4} - \mathit{lbm1}$ & $\mathit{T2} + \mathit{lbm1} = \mathit{Tq4}$ \\
40 & $\mathit{b5m2} = \mathit{B} - 2$ & $\mathit{b5m2} + 2 = \mathit{B}$ \\
41 & $\mathit{packed\_target\_shift} = \mathit{l} \cdot \mathit{q4}$ & $\mathit{l} \cdot \mathit{q4} = \mathit{packed\_target\_shift}$ \\
42 & $\mathit{mask\_term} = \mathit{b5m2} \cdot \mathit{packed\_target\_shift}$ & $\mathit{b5m2} \cdot \mathit{packed\_target\_shift} = \mathit{mask\_term}$ \\
43 & $\mathit{T} = \mathit{T2} + \mathit{mask\_term}$ & $\mathit{T2} + \mathit{mask\_term} = \mathit{T}$ \\
44 & $\mathit{n2m1} = \mathit{n2} - 1$ & $\mathit{n2m1} + 1 = \mathit{n2}$ \\
45 & $\mathit{TB} = \mathit{T} \cdot \mathit{n2m1}$ & $\mathit{T} \cdot \mathit{n2m1} = \mathit{TB}$ \\
46 & $\mathit{R7} = \mathit{SA} + \mathit{TB}$ & $\mathit{SA} + \mathit{TB} = \mathit{R7}$ \\
47 & $\mathit{wn2} = \mathit{w} \cdot \mathit{n2}$ & $\mathit{w} \cdot \mathit{n2} = \mathit{wn2}$ \\
48 & $\mathit{sn2} = \mathit{s} \cdot \mathit{n2}$ & $\mathit{s} \cdot \mathit{n2} = \mathit{sn2}$ \\
49 & $\mathit{rsn2} = \mathit{r} \cdot \mathit{sn2}$ & $\mathit{r} \cdot \mathit{sn2} = \mathit{rsn2}$ \\
50 & $\mathit{UM} = \mathit{wn2} \cdot \mathit{rsn2}$ & $\mathit{wn2} \cdot \mathit{rsn2} = \mathit{UM}$ \\
51 & $\mathit{R12} = \mathit{UM} + \mathit{rsn2}$ & $\mathit{UM} + \mathit{rsn2} = \mathit{R12}$ \\
52 & $\mathit{wsq} = \mathit{UM} \cdot \mathit{rsn2}$ & $\mathit{UM} \cdot \mathit{rsn2} = \mathit{wsq}$ \\
53 & $\mathit{R8} = 2 \cdot \mathit{wsq}$ & $2 \cdot \mathit{wsq} = \mathit{R8}$ \\
54 & $\mathit{Pplus1} = \mathit{R8} + 2$ & $\mathit{R8} + 2 = \mathit{Pplus1}$ \\
55 & $\mathit{p2m1} = \mathit{R8} \cdot \mathit{Pplus1}$ & $\mathit{R8} \cdot \mathit{Pplus1} = \mathit{p2m1}$ \\
56 & $\mathit{k2} = \mathit{k} \cdot \mathit{k}$ & $\mathit{k} \cdot \mathit{k} = \mathit{k2}$ \\
57 & $\mathit{P9} = \mathit{p2m1} \cdot \mathit{k2}$ & $\mathit{p2m1} \cdot \mathit{k2} = \mathit{P9}$ \\
58 & $\mathit{L9} = \mathit{P9} + 1$ & $\mathit{P9} + 1 = \mathit{L9}$ \\
59 & $\mathit{R9} = \mathit{tau} \cdot \mathit{tau}$ & $\mathit{tau} \cdot \mathit{tau} = \mathit{R9}$ \\
60 & $\mathit{ksn2} = \mathit{k} \cdot \mathit{sn2}$ & $\mathit{k} \cdot \mathit{sn2} = \mathit{ksn2}$ \\
61 & $\mathit{R10a} = \mathit{ksn2} + \mathit{eta}$ & $\mathit{ksn2} + \mathit{eta} = \mathit{R10a}$ \\
62 & $\mathit{R10b} = \mathit{eta} + \mathit{zeta}$ & $\mathit{eta} + \mathit{zeta} = \mathit{R10b}$ \\
63 & $\mathit{r1} = \mathit{r} + 1$ & $\mathit{r} + 1 = \mathit{r1}$ \\
64 & $\mathit{hpm1} = \mathit{h} \cdot \mathit{R8}$ & $\mathit{h} \cdot \mathit{R8} = \mathit{hpm1}$ \\
65 & $\mathit{R11} = \mathit{r1} + \mathit{hpm1}$ & $\mathit{r1} + \mathit{hpm1} = \mathit{R11}$ \\
66 & $\mathit{tr1} = \mathit{r1} + \mathit{r}$ & $\mathit{r1} + \mathit{r} = \mathit{tr1}$ \\
67 & $\mathit{R13} = \mathit{ka} + \mathit{phi}$ & $\mathit{ka} + \mathit{phi} = \mathit{R13}$ \\
68 & $\mathit{bw} = \mathit{b} \cdot \mathit{w}$ & $\mathit{b} \cdot \mathit{w} = \mathit{bw}$ \\
69 & $\mathit{am2} = \mathit{a} - 2$ & $\mathit{am2} + 2 = \mathit{a}$ \\
70 & $\mathit{cam2} = \mathit{c} \cdot \mathit{am2}$ & $\mathit{c} \cdot \mathit{am2} = \mathit{cam2}$ \\
71 & $\mathit{D1} = \mathit{bw} + \mathit{cam2}$ & $\mathit{bw} + \mathit{cam2} = \mathit{D1}$ \\
72 & $\mathit{a4} = 4 \cdot \mathit{a}$ & $4 \cdot \mathit{a} = \mathit{a4}$ \\
73 & $\mathit{a4m5} = \mathit{a4} - 5$ & $\mathit{a4m5} + 5 = \mathit{a4}$ \\
74 & $\mathit{gam} = \mathit{ga} \cdot \mathit{a4m5}$ & $\mathit{ga} \cdot \mathit{a4m5} = \mathit{gam}$ \\
75 & $\mathit{R14} = \mathit{D1} + \mathit{gam}$ & $\mathit{D1} + \mathit{gam} = \mathit{R14}$ \\
76 & $\mathit{am1} = \mathit{a} - 1$ & $\mathit{am1} + 1 = \mathit{a}$ \\
77 & $\mathit{ap1} = \mathit{a} + 1$ & $\mathit{a} + 1 = \mathit{ap1}$ \\
78 & $\mathit{A} = \mathit{am1} \cdot \mathit{ap1}$ & $\mathit{am1} \cdot \mathit{ap1} = \mathit{A}$ \\
79 & $\mathit{c2} = \mathit{c} \cdot \mathit{c}$ & $\mathit{c} \cdot \mathit{c} = \mathit{c2}$ \\
80 & $\mathit{Ac2} = \mathit{A} \cdot \mathit{c2}$ & $\mathit{A} \cdot \mathit{c2} = \mathit{Ac2}$ \\
81 & $\mathit{R15} = \mathit{Ac2} + 1$ & $\mathit{Ac2} + 1 = \mathit{R15}$ \\
82 & $\mathit{L15} = \mathit{d} \cdot \mathit{d}$ & $\mathit{d} \cdot \mathit{d} = \mathit{L15}$ \\
83 & $\mathit{ic2} = \mathit{i} \cdot \mathit{c2}$ & $\mathit{i} \cdot \mathit{c2} = \mathit{ic2}$ \\
84 & $\mathit{ic22} = \mathit{ic2} \cdot \mathit{ic2}$ & $\mathit{ic2} \cdot \mathit{ic2} = \mathit{ic22}$ \\
85 & $\mathit{AE} = \mathit{A} \cdot \mathit{ic22}$ & $\mathit{A} \cdot \mathit{ic22} = \mathit{AE}$ \\
86 & $\mathit{R16} = \mathit{AE} + 1$ & $\mathit{AE} + 1 = \mathit{R16}$ \\
87 & $\mathit{L16} = \mathit{f} \cdot \mathit{f}$ & $\mathit{f} \cdot \mathit{f} = \mathit{L16}$ \\
88 & $\mathit{of} = \mathit{o} \cdot \mathit{f}$ & $\mathit{o} \cdot \mathit{f} = \mathit{of}$ \\
89 & $\mathit{dof} = \mathit{of} - \mathit{d}$ & $\mathit{dof} + \mathit{d} = \mathit{of}$ \\
90 & $\mathit{L17} = \mathit{dof} \cdot \mathit{dof}$ & $\mathit{dof} \cdot \mathit{dof} = \mathit{L17}$ \\
91 & $\mathit{Gminus1} = \mathit{ap1} \cdot \mathit{AE}$ & $\mathit{ap1} \cdot \mathit{AE} = \mathit{Gminus1}$ \\
92 & $\mathit{Gplus1} = \mathit{Gminus1} + 2$ & $\mathit{Gminus1} + 2 = \mathit{Gplus1}$ \\
93 & $\mathit{G2m1} = \mathit{Gminus1} \cdot \mathit{Gplus1}$ & $\mathit{Gminus1} \cdot \mathit{Gplus1} = \mathit{G2m1}$ \\
94 & $\mathit{jc} = \mathit{j} \cdot \mathit{c}$ & $\mathit{j} \cdot \mathit{c} = \mathit{jc}$ \\
95 & $\mathit{H17} = \mathit{tr1} + \mathit{jc}$ & $\mathit{tr1} + \mathit{jc} = \mathit{H17}$ \\
96 & $\mathit{H2} = \mathit{H17} \cdot \mathit{H17}$ & $\mathit{H17} \cdot \mathit{H17} = \mathit{H2}$ \\
97 & $\mathit{P17} = \mathit{G2m1} \cdot \mathit{H2}$ & $\mathit{G2m1} \cdot \mathit{H2} = \mathit{P17}$ \\
98 & $\mathit{R17} = \mathit{P17} + 1$ & $\mathit{P17} + 1 = \mathit{R17}$ \\
99 & $\mathit{amb} = \mathit{a} - \mathit{B}$ & $\mathit{amb} + \mathit{B} = \mathit{a}$ \\
100 & $\mathit{kamb} = \mathit{ka} \cdot \mathit{amb}$ & $\mathit{ka} \cdot \mathit{amb} = \mathit{kamb}$ \\
101 & $\mathit{qk} = \mathit{q} + \mathit{kamb}$ & $\mathit{q} + \mathit{kamb} = \mathit{qk}$ \\
102 & $\mathit{amb2} = \mathit{amb} \cdot \mathit{amb}$ & $\mathit{amb} \cdot \mathit{amb} = \mathit{amb2}$ \\
103 & $\mathit{Mod} = \mathit{A} - \mathit{amb2}$ & $\mathit{Mod} + \mathit{amb2} = \mathit{A}$ \\
104 & $\mathit{rM} = \mathit{rho} \cdot \mathit{Mod}$ & $\mathit{rho} \cdot \mathit{Mod} = \mathit{rM}$ \\
105 & $\mathit{R18} = \mathit{qk} + \mathit{rM}$ & $\mathit{qk} + \mathit{rM} = \mathit{R18}$ \\
106 & $\mathit{ka2} = \mathit{ka} \cdot \mathit{ka}$ & $\mathit{ka} \cdot \mathit{ka} = \mathit{ka2}$ \\
107 & $\mathit{Aka2} = \mathit{A} \cdot \mathit{ka2}$ & $\mathit{A} \cdot \mathit{ka2} = \mathit{Aka2}$ \\
108 & $\mathit{R19} = \mathit{Aka2} + 1$ & $\mathit{Aka2} + 1 = \mathit{R19}$ \\
109 & $\mathit{L19} = \mathit{mu} \cdot \mathit{mu}$ & $\mathit{mu} \cdot \mathit{mu} = \mathit{L19}$ \\
110 & $\mathit{Dam1} = \mathit{Delta} \cdot \mathit{am1}$ & $\mathit{Delta} \cdot \mathit{am1} = \mathit{Dam1}$ \\
111 & $\mathit{R20} = \mathit{Dam1} + 5^{16}$ & $\mathit{Dam1} + 5^{16} = \mathit{R20}$ \\
\end{longtable}

The twenty-one free equality tests are
\[\begin{array}{@{}l@{}}
\mathit{L1}=\mathit{q2},\quad \mathit{b}=\mathit{bx},\quad \mathit{L2}=\mathit{R2},\\
\mathit{L3}=\mathit{B},\quad \mathit{S2}=\mathit{pack\_rhs},\quad \mathit{n}=\mathit{q8},\\
\mathit{r}=\mathit{R7},\quad \mathit{L9}=\mathit{R9},\quad \mathit{c}=\mathit{R10a},\\
\mathit{k}=\mathit{R10b},\quad \mathit{k}=\mathit{R11},\quad \mathit{a}=\mathit{R12},\\
\mathit{c}=\mathit{R13},\quad \mathit{d}=\mathit{R14},\quad \mathit{L15}=\mathit{R15},\\
\mathit{L16}=\mathit{R16},\quad \mathit{L17}=\mathit{R17},\quad \mathit{mu}=\mathit{R18},\\
\mathit{L19}=\mathit{R19},\quad \mathit{ka}=\mathit{R20},\quad \mathit{S3}=\mathit{sigma}
\end{array}\]


\clearpage
\normalsize
\section{The 110-instruction certificate}\label{app:110}

Use the homogeneous fixed index of Section~\ref{sec:unit-digit}, with
the same 33 positive witnesses and 21 equality tests as in
Appendix~\ref{app:111}. The decoded unit coordinate is part of the
integer code, not an additional system unknown. The code register is
now $C=xB+g$, and the third mask uses $B-4$.

The complete checker \file{round8\_1980\_certificate.py} verifies every
primitive, source residual, and support inequality for this layout.
There are 61 multiplication checks and 49 addition checks; the seven
numeral instructions give 117 when required. Theorem~\ref{thm:best110}
proves the encoding and positive-witness equivalence. The table is
generated by \file{round8\_1980\_schedule\_table.py}.

\small
% Generated by round8_1980_schedule_table.py; do not edit.
\begin{longtable}{@{}rll@{}}
\toprule
No. & Assignment & Addition/multiplication check\\
\midrule
\endfirsthead
\toprule
No. & Assignment & Addition/multiplication check\\
\midrule
\endhead
\bottomrule
\endfoot
1 & $\mathit{q2} = \mathit{q} \cdot \mathit{q}$ & $\mathit{q} \cdot \mathit{q} = \mathit{q2}$ \\
2 & $\mathit{q4} = \mathit{q2} \cdot \mathit{q2}$ & $\mathit{q2} \cdot \mathit{q2} = \mathit{q4}$ \\
3 & $\mathit{q8} = \mathit{q4} \cdot \mathit{q4}$ & $\mathit{q4} \cdot \mathit{q4} = \mathit{q8}$ \\
4 & $\mathit{b2} = \mathit{b} \cdot \mathit{b}$ & $\mathit{b} \cdot \mathit{b} = \mathit{b2}$ \\
5 & $\mathit{B} = \mathit{H} \cdot \mathit{b2}$ & $\mathit{H} \cdot \mathit{b2} = \mathit{B}$ \\
6 & $\mathit{bx} = \mathit{x} + \mathit{beta}$ & $\mathit{x} + \mathit{beta} = \mathit{bx}$ \\
7 & $\mathit{L2} = \mathit{la} + \mathit{q2}$ & $\mathit{la} + \mathit{q2} = \mathit{L2}$ \\
8 & $\mathit{Lam} = \mathit{la} \cdot \mathit{B}$ & $\mathit{la} \cdot \mathit{B} = \mathit{Lam}$ \\
9 & $\mathit{R2} = \mathit{Lam} + 1$ & $\mathit{Lam} + 1 = \mathit{R2}$ \\
10 & $\mathit{L3} = \mathit{th} + \mathit{Z}$ & $\mathit{th} + \mathit{Z} = \mathit{L3}$ \\
11 & $\mathit{eq} = \mathit{e} \cdot \mathit{q}$ & $\mathit{e} \cdot \mathit{q} = \mathit{eq}$ \\
12 & $\mathit{S2} = \mathit{l} + \mathit{eq}$ & $\mathit{l} + \mathit{eq} = \mathit{S2}$ \\
13 & $\mathit{pack\_tth} = \mathit{t} \cdot \mathit{th}$ & $\mathit{t} \cdot \mathit{th} = \mathit{pack\_tth}$ \\
14 & $\mathit{pack\_rhs} = \mathit{V} + \mathit{pack\_tth}$ & $\mathit{V} + \mathit{pack\_tth} = \mathit{pack\_rhs}$ \\
15 & $\mathit{zl} = \mathit{Z} \cdot \mathit{la}$ & $\mathit{Z} \cdot \mathit{la} = \mathit{zl}$ \\
16 & $\mathit{twice\_e} = 2 \cdot \mathit{e}$ & $2 \cdot \mathit{e} = \mathit{twice\_e}$ \\
17 & $\mathit{t2} = \mathit{twice\_e} - \mathit{zl}$ & $\mathit{t2} + \mathit{zl} = \mathit{twice\_e}$ \\
18 & $\mathit{xb5} = \mathit{x} \cdot \mathit{B}$ & $\mathit{x} \cdot \mathit{B} = \mathit{xb5}$ \\
19 & $\mathit{C} = \mathit{xb5} + \mathit{g}$ & $\mathit{xb5} + \mathit{g} = \mathit{C}$ \\
20 & $\mathit{C2} = \mathit{C} \cdot \mathit{C}$ & $\mathit{C} \cdot \mathit{C} = \mathit{C2}$ \\
21 & $\mathit{bound\_sum} = \mathit{S2} + \mathit{C2}$ & $\mathit{S2} + \mathit{C2} = \mathit{bound\_sum}$ \\
22 & $\mathit{L1} = \mathit{bound\_sum} + \mathit{al}$ & $\mathit{bound\_sum} + \mathit{al} = \mathit{L1}$ \\
23 & $\mathit{P1} = \mathit{t2} \cdot \mathit{C2}$ & $\mathit{t2} \cdot \mathit{C2} = \mathit{P1}$ \\
24 & $\mathit{qp1} = \mathit{q} + 1$ & $\mathit{q} + 1 = \mathit{qp1}$ \\
25 & $\mathit{P2} = \mathit{Lam} \cdot \mathit{qp1}$ & $\mathit{Lam} \cdot \mathit{qp1} = \mathit{P2}$ \\
26 & $\mathit{S3} = \mathit{P1} + \mathit{P2}$ & $\mathit{P1} + \mathit{P2} = \mathit{S3}$ \\
27 & $\mathit{S3q4} = \mathit{S3} \cdot \mathit{q2}$ & $\mathit{S3} \cdot \mathit{q2} = \mathit{S3q4}$ \\
28 & $\mathit{Sin} = \mathit{S2} + \mathit{S3q4}$ & $\mathit{S2} + \mathit{S3q4} = \mathit{Sin}$ \\
29 & $\mathit{Sq4} = \mathit{Sin} \cdot \mathit{q2}$ & $\mathit{Sin} \cdot \mathit{q2} = \mathit{Sq4}$ \\
30 & $\mathit{S} = \mathit{g} + \mathit{Sq4}$ & $\mathit{g} + \mathit{Sq4} = \mathit{S}$ \\
31 & $\mathit{n2} = \mathit{n} \cdot \mathit{n}$ & $\mathit{n} \cdot \mathit{n} = \mathit{n2}$ \\
32 & $\mathit{n2n} = \mathit{n2} - \mathit{n}$ & $\mathit{n2n} + \mathit{n} = \mathit{n2}$ \\
33 & $\mathit{SA} = \mathit{S} \cdot \mathit{n2n}$ & $\mathit{S} \cdot \mathit{n2n} = \mathit{SA}$ \\
34 & $\mathit{Tcoef} = \mathit{L2} - \mathit{zl}$ & $\mathit{Tcoef} + \mathit{zl} = \mathit{L2}$ \\
35 & $\mathit{Tq4} = \mathit{Tcoef} \cdot \mathit{q2}$ & $\mathit{Tcoef} \cdot \mathit{q2} = \mathit{Tq4}$ \\
36 & $\mathit{bm1} = \mathit{b} - 1$ & $\mathit{bm1} + 1 = \mathit{b}$ \\
37 & $\mathit{lbm1} = \mathit{l} \cdot \mathit{bm1}$ & $\mathit{l} \cdot \mathit{bm1} = \mathit{lbm1}$ \\
38 & $\mathit{T2} = \mathit{Tq4} - \mathit{lbm1}$ & $\mathit{T2} + \mathit{lbm1} = \mathit{Tq4}$ \\
39 & $\mathit{b5m2} = \mathit{B} - 4$ & $\mathit{b5m2} + 4 = \mathit{B}$ \\
40 & $\mathit{packed\_target\_shift} = \mathit{l} \cdot \mathit{q4}$ & $\mathit{l} \cdot \mathit{q4} = \mathit{packed\_target\_shift}$ \\
41 & $\mathit{mask\_term} = \mathit{b5m2} \cdot \mathit{packed\_target\_shift}$ & $\mathit{b5m2} \cdot \mathit{packed\_target\_shift} = \mathit{mask\_term}$ \\
42 & $\mathit{T} = \mathit{T2} + \mathit{mask\_term}$ & $\mathit{T2} + \mathit{mask\_term} = \mathit{T}$ \\
43 & $\mathit{n2m1} = \mathit{n2} - 1$ & $\mathit{n2m1} + 1 = \mathit{n2}$ \\
44 & $\mathit{TB} = \mathit{T} \cdot \mathit{n2m1}$ & $\mathit{T} \cdot \mathit{n2m1} = \mathit{TB}$ \\
45 & $\mathit{R7} = \mathit{SA} + \mathit{TB}$ & $\mathit{SA} + \mathit{TB} = \mathit{R7}$ \\
46 & $\mathit{wn2} = \mathit{w} \cdot \mathit{n2}$ & $\mathit{w} \cdot \mathit{n2} = \mathit{wn2}$ \\
47 & $\mathit{sn2} = \mathit{s} \cdot \mathit{n2}$ & $\mathit{s} \cdot \mathit{n2} = \mathit{sn2}$ \\
48 & $\mathit{rsn2} = \mathit{r} \cdot \mathit{sn2}$ & $\mathit{r} \cdot \mathit{sn2} = \mathit{rsn2}$ \\
49 & $\mathit{UM} = \mathit{wn2} \cdot \mathit{rsn2}$ & $\mathit{wn2} \cdot \mathit{rsn2} = \mathit{UM}$ \\
50 & $\mathit{R12} = \mathit{UM} + \mathit{rsn2}$ & $\mathit{UM} + \mathit{rsn2} = \mathit{R12}$ \\
51 & $\mathit{wsq} = \mathit{UM} \cdot \mathit{rsn2}$ & $\mathit{UM} \cdot \mathit{rsn2} = \mathit{wsq}$ \\
52 & $\mathit{R8} = 2 \cdot \mathit{wsq}$ & $2 \cdot \mathit{wsq} = \mathit{R8}$ \\
53 & $\mathit{Pplus1} = \mathit{R8} + 2$ & $\mathit{R8} + 2 = \mathit{Pplus1}$ \\
54 & $\mathit{p2m1} = \mathit{R8} \cdot \mathit{Pplus1}$ & $\mathit{R8} \cdot \mathit{Pplus1} = \mathit{p2m1}$ \\
55 & $\mathit{k2} = \mathit{k} \cdot \mathit{k}$ & $\mathit{k} \cdot \mathit{k} = \mathit{k2}$ \\
56 & $\mathit{P9} = \mathit{p2m1} \cdot \mathit{k2}$ & $\mathit{p2m1} \cdot \mathit{k2} = \mathit{P9}$ \\
57 & $\mathit{L9} = \mathit{P9} + 1$ & $\mathit{P9} + 1 = \mathit{L9}$ \\
58 & $\mathit{R9} = \mathit{tau} \cdot \mathit{tau}$ & $\mathit{tau} \cdot \mathit{tau} = \mathit{R9}$ \\
59 & $\mathit{ksn2} = \mathit{k} \cdot \mathit{sn2}$ & $\mathit{k} \cdot \mathit{sn2} = \mathit{ksn2}$ \\
60 & $\mathit{R10a} = \mathit{ksn2} + \mathit{eta}$ & $\mathit{ksn2} + \mathit{eta} = \mathit{R10a}$ \\
61 & $\mathit{R10b} = \mathit{eta} + \mathit{zeta}$ & $\mathit{eta} + \mathit{zeta} = \mathit{R10b}$ \\
62 & $\mathit{r1} = \mathit{r} + 1$ & $\mathit{r} + 1 = \mathit{r1}$ \\
63 & $\mathit{hpm1} = \mathit{h} \cdot \mathit{R8}$ & $\mathit{h} \cdot \mathit{R8} = \mathit{hpm1}$ \\
64 & $\mathit{R11} = \mathit{r1} + \mathit{hpm1}$ & $\mathit{r1} + \mathit{hpm1} = \mathit{R11}$ \\
65 & $\mathit{tr1} = \mathit{r1} + \mathit{r}$ & $\mathit{r1} + \mathit{r} = \mathit{tr1}$ \\
66 & $\mathit{R13} = \mathit{ka} + \mathit{phi}$ & $\mathit{ka} + \mathit{phi} = \mathit{R13}$ \\
67 & $\mathit{bw} = \mathit{b} \cdot \mathit{w}$ & $\mathit{b} \cdot \mathit{w} = \mathit{bw}$ \\
68 & $\mathit{am2} = \mathit{a} - 2$ & $\mathit{am2} + 2 = \mathit{a}$ \\
69 & $\mathit{cam2} = \mathit{c} \cdot \mathit{am2}$ & $\mathit{c} \cdot \mathit{am2} = \mathit{cam2}$ \\
70 & $\mathit{D1} = \mathit{bw} + \mathit{cam2}$ & $\mathit{bw} + \mathit{cam2} = \mathit{D1}$ \\
71 & $\mathit{a4} = 4 \cdot \mathit{a}$ & $4 \cdot \mathit{a} = \mathit{a4}$ \\
72 & $\mathit{a4m5} = \mathit{a4} - 5$ & $\mathit{a4m5} + 5 = \mathit{a4}$ \\
73 & $\mathit{gam} = \mathit{ga} \cdot \mathit{a4m5}$ & $\mathit{ga} \cdot \mathit{a4m5} = \mathit{gam}$ \\
74 & $\mathit{R14} = \mathit{D1} + \mathit{gam}$ & $\mathit{D1} + \mathit{gam} = \mathit{R14}$ \\
75 & $\mathit{am1} = \mathit{a} - 1$ & $\mathit{am1} + 1 = \mathit{a}$ \\
76 & $\mathit{ap1} = \mathit{a} + 1$ & $\mathit{a} + 1 = \mathit{ap1}$ \\
77 & $\mathit{A} = \mathit{am1} \cdot \mathit{ap1}$ & $\mathit{am1} \cdot \mathit{ap1} = \mathit{A}$ \\
78 & $\mathit{c2} = \mathit{c} \cdot \mathit{c}$ & $\mathit{c} \cdot \mathit{c} = \mathit{c2}$ \\
79 & $\mathit{Ac2} = \mathit{A} \cdot \mathit{c2}$ & $\mathit{A} \cdot \mathit{c2} = \mathit{Ac2}$ \\
80 & $\mathit{R15} = \mathit{Ac2} + 1$ & $\mathit{Ac2} + 1 = \mathit{R15}$ \\
81 & $\mathit{L15} = \mathit{d} \cdot \mathit{d}$ & $\mathit{d} \cdot \mathit{d} = \mathit{L15}$ \\
82 & $\mathit{ic2} = \mathit{i} \cdot \mathit{c2}$ & $\mathit{i} \cdot \mathit{c2} = \mathit{ic2}$ \\
83 & $\mathit{ic22} = \mathit{ic2} \cdot \mathit{ic2}$ & $\mathit{ic2} \cdot \mathit{ic2} = \mathit{ic22}$ \\
84 & $\mathit{AE} = \mathit{A} \cdot \mathit{ic22}$ & $\mathit{A} \cdot \mathit{ic22} = \mathit{AE}$ \\
85 & $\mathit{R16} = \mathit{AE} + 1$ & $\mathit{AE} + 1 = \mathit{R16}$ \\
86 & $\mathit{L16} = \mathit{f} \cdot \mathit{f}$ & $\mathit{f} \cdot \mathit{f} = \mathit{L16}$ \\
87 & $\mathit{of} = \mathit{o} \cdot \mathit{f}$ & $\mathit{o} \cdot \mathit{f} = \mathit{of}$ \\
88 & $\mathit{dof} = \mathit{of} - \mathit{d}$ & $\mathit{dof} + \mathit{d} = \mathit{of}$ \\
89 & $\mathit{L17} = \mathit{dof} \cdot \mathit{dof}$ & $\mathit{dof} \cdot \mathit{dof} = \mathit{L17}$ \\
90 & $\mathit{Gminus1} = \mathit{ap1} \cdot \mathit{AE}$ & $\mathit{ap1} \cdot \mathit{AE} = \mathit{Gminus1}$ \\
91 & $\mathit{Gplus1} = \mathit{Gminus1} + 2$ & $\mathit{Gminus1} + 2 = \mathit{Gplus1}$ \\
92 & $\mathit{G2m1} = \mathit{Gminus1} \cdot \mathit{Gplus1}$ & $\mathit{Gminus1} \cdot \mathit{Gplus1} = \mathit{G2m1}$ \\
93 & $\mathit{jc} = \mathit{j} \cdot \mathit{c}$ & $\mathit{j} \cdot \mathit{c} = \mathit{jc}$ \\
94 & $\mathit{H17} = \mathit{tr1} + \mathit{jc}$ & $\mathit{tr1} + \mathit{jc} = \mathit{H17}$ \\
95 & $\mathit{H2} = \mathit{H17} \cdot \mathit{H17}$ & $\mathit{H17} \cdot \mathit{H17} = \mathit{H2}$ \\
96 & $\mathit{P17} = \mathit{G2m1} \cdot \mathit{H2}$ & $\mathit{G2m1} \cdot \mathit{H2} = \mathit{P17}$ \\
97 & $\mathit{R17} = \mathit{P17} + 1$ & $\mathit{P17} + 1 = \mathit{R17}$ \\
98 & $\mathit{amb} = \mathit{a} - \mathit{B}$ & $\mathit{amb} + \mathit{B} = \mathit{a}$ \\
99 & $\mathit{kamb} = \mathit{ka} \cdot \mathit{amb}$ & $\mathit{ka} \cdot \mathit{amb} = \mathit{kamb}$ \\
100 & $\mathit{qk} = \mathit{q} + \mathit{kamb}$ & $\mathit{q} + \mathit{kamb} = \mathit{qk}$ \\
101 & $\mathit{amb2} = \mathit{amb} \cdot \mathit{amb}$ & $\mathit{amb} \cdot \mathit{amb} = \mathit{amb2}$ \\
102 & $\mathit{Mod} = \mathit{A} - \mathit{amb2}$ & $\mathit{Mod} + \mathit{amb2} = \mathit{A}$ \\
103 & $\mathit{rM} = \mathit{rho} \cdot \mathit{Mod}$ & $\mathit{rho} \cdot \mathit{Mod} = \mathit{rM}$ \\
104 & $\mathit{R18} = \mathit{qk} + \mathit{rM}$ & $\mathit{qk} + \mathit{rM} = \mathit{R18}$ \\
105 & $\mathit{ka2} = \mathit{ka} \cdot \mathit{ka}$ & $\mathit{ka} \cdot \mathit{ka} = \mathit{ka2}$ \\
106 & $\mathit{Aka2} = \mathit{A} \cdot \mathit{ka2}$ & $\mathit{A} \cdot \mathit{ka2} = \mathit{Aka2}$ \\
107 & $\mathit{R19} = \mathit{Aka2} + 1$ & $\mathit{Aka2} + 1 = \mathit{R19}$ \\
108 & $\mathit{L19} = \mathit{mu} \cdot \mathit{mu}$ & $\mathit{mu} \cdot \mathit{mu} = \mathit{L19}$ \\
109 & $\mathit{Dam1} = \mathit{Delta} \cdot \mathit{am1}$ & $\mathit{Delta} \cdot \mathit{am1} = \mathit{Dam1}$ \\
110 & $\mathit{R20} = \mathit{Dam1} + 5^{16}$ & $\mathit{Dam1} + 5^{16} = \mathit{R20}$ \\
\end{longtable}

The twenty-one free equality tests are
\[\begin{array}{@{}l@{}}
\mathit{L1}=\mathit{q2},\quad \mathit{b}=\mathit{bx},\quad \mathit{L2}=\mathit{R2},\\
\mathit{L3}=\mathit{B},\quad \mathit{S2}=\mathit{pack\_rhs},\quad \mathit{n}=\mathit{q8},\\
\mathit{r}=\mathit{R7},\quad \mathit{L9}=\mathit{R9},\quad \mathit{c}=\mathit{R10a},\\
\mathit{k}=\mathit{R10b},\quad \mathit{k}=\mathit{R11},\quad \mathit{a}=\mathit{R12},\\
\mathit{c}=\mathit{R13},\quad \mathit{d}=\mathit{R14},\quad \mathit{L15}=\mathit{R15},\\
\mathit{L16}=\mathit{R16},\quad \mathit{L17}=\mathit{R17},\quad \mathit{mu}=\mathit{R18},\\
\mathit{L19}=\mathit{R19},\quad \mathit{ka}=\mathit{R20},\quad \mathit{S3}=\mathit{sigma}
\end{array}\]


\clearpage
\normalsize
\section{The 109-instruction certificate}\label{app:109}

Use the fixed input-unit index of Section~\ref{sec:input-unit}, with
the same positive witnesses and equality tests. The sole arithmetic
change from Appendix~\ref{app:110} is $C=x+g$, which removes the input
multiplication by $B$. The index assigns the input to weight zero and
the decoded unit coordinate to weight one; these are part of its fixed
definition and introduce no extra arithmetic checks.

The complete checker \file{round9\_1980\_certificate.py} verifies the
109 primitive statements and all 21 source residuals, including both
triangular corrections. There are 60 multiplications and 49 additions;
the seven numeral instructions give 116 when required.
Theorem~\ref{thm:best109} proves the relaxed bootstrap and the resulting
encoding equivalence. The table is generated by
\file{round9\_1980\_schedule\_table.py}.

\small
% Generated by round9_1980_schedule_table.py; do not edit.
\begin{longtable}{@{}rll@{}}
\toprule
No. & Assignment & Addition/multiplication check\\
\midrule
\endfirsthead
\toprule
No. & Assignment & Addition/multiplication check\\
\midrule
\endhead
\bottomrule
\endfoot
1 & $\mathit{q2} = \mathit{q} \cdot \mathit{q}$ & $\mathit{q} \cdot \mathit{q} = \mathit{q2}$ \\
2 & $\mathit{q4} = \mathit{q2} \cdot \mathit{q2}$ & $\mathit{q2} \cdot \mathit{q2} = \mathit{q4}$ \\
3 & $\mathit{q8} = \mathit{q4} \cdot \mathit{q4}$ & $\mathit{q4} \cdot \mathit{q4} = \mathit{q8}$ \\
4 & $\mathit{b2} = \mathit{b} \cdot \mathit{b}$ & $\mathit{b} \cdot \mathit{b} = \mathit{b2}$ \\
5 & $\mathit{B} = \mathit{H} \cdot \mathit{b2}$ & $\mathit{H} \cdot \mathit{b2} = \mathit{B}$ \\
6 & $\mathit{bx} = \mathit{x} + \mathit{beta}$ & $\mathit{x} + \mathit{beta} = \mathit{bx}$ \\
7 & $\mathit{L2} = \mathit{la} + \mathit{q2}$ & $\mathit{la} + \mathit{q2} = \mathit{L2}$ \\
8 & $\mathit{Lam} = \mathit{la} \cdot \mathit{B}$ & $\mathit{la} \cdot \mathit{B} = \mathit{Lam}$ \\
9 & $\mathit{R2} = \mathit{Lam} + 1$ & $\mathit{Lam} + 1 = \mathit{R2}$ \\
10 & $\mathit{L3} = \mathit{th} + \mathit{Z}$ & $\mathit{th} + \mathit{Z} = \mathit{L3}$ \\
11 & $\mathit{eq} = \mathit{e} \cdot \mathit{q}$ & $\mathit{e} \cdot \mathit{q} = \mathit{eq}$ \\
12 & $\mathit{S2} = \mathit{l} + \mathit{eq}$ & $\mathit{l} + \mathit{eq} = \mathit{S2}$ \\
13 & $\mathit{pack\_tth} = \mathit{t} \cdot \mathit{th}$ & $\mathit{t} \cdot \mathit{th} = \mathit{pack\_tth}$ \\
14 & $\mathit{pack\_rhs} = \mathit{V} + \mathit{pack\_tth}$ & $\mathit{V} + \mathit{pack\_tth} = \mathit{pack\_rhs}$ \\
15 & $\mathit{zl} = \mathit{Z} \cdot \mathit{la}$ & $\mathit{Z} \cdot \mathit{la} = \mathit{zl}$ \\
16 & $\mathit{twice\_e} = 2 \cdot \mathit{e}$ & $2 \cdot \mathit{e} = \mathit{twice\_e}$ \\
17 & $\mathit{t2} = \mathit{twice\_e} - \mathit{zl}$ & $\mathit{t2} + \mathit{zl} = \mathit{twice\_e}$ \\
18 & $\mathit{C} = \mathit{x} + \mathit{g}$ & $\mathit{x} + \mathit{g} = \mathit{C}$ \\
19 & $\mathit{C2} = \mathit{C} \cdot \mathit{C}$ & $\mathit{C} \cdot \mathit{C} = \mathit{C2}$ \\
20 & $\mathit{bound\_sum} = \mathit{S2} + \mathit{C2}$ & $\mathit{S2} + \mathit{C2} = \mathit{bound\_sum}$ \\
21 & $\mathit{L1} = \mathit{bound\_sum} + \mathit{al}$ & $\mathit{bound\_sum} + \mathit{al} = \mathit{L1}$ \\
22 & $\mathit{P1} = \mathit{t2} \cdot \mathit{C2}$ & $\mathit{t2} \cdot \mathit{C2} = \mathit{P1}$ \\
23 & $\mathit{qp1} = \mathit{q} + 1$ & $\mathit{q} + 1 = \mathit{qp1}$ \\
24 & $\mathit{P2} = \mathit{Lam} \cdot \mathit{qp1}$ & $\mathit{Lam} \cdot \mathit{qp1} = \mathit{P2}$ \\
25 & $\mathit{S3} = \mathit{P1} + \mathit{P2}$ & $\mathit{P1} + \mathit{P2} = \mathit{S3}$ \\
26 & $\mathit{S3q4} = \mathit{S3} \cdot \mathit{q2}$ & $\mathit{S3} \cdot \mathit{q2} = \mathit{S3q4}$ \\
27 & $\mathit{Sin} = \mathit{S2} + \mathit{S3q4}$ & $\mathit{S2} + \mathit{S3q4} = \mathit{Sin}$ \\
28 & $\mathit{Sq4} = \mathit{Sin} \cdot \mathit{q2}$ & $\mathit{Sin} \cdot \mathit{q2} = \mathit{Sq4}$ \\
29 & $\mathit{S} = \mathit{g} + \mathit{Sq4}$ & $\mathit{g} + \mathit{Sq4} = \mathit{S}$ \\
30 & $\mathit{n2} = \mathit{n} \cdot \mathit{n}$ & $\mathit{n} \cdot \mathit{n} = \mathit{n2}$ \\
31 & $\mathit{n2n} = \mathit{n2} - \mathit{n}$ & $\mathit{n2n} + \mathit{n} = \mathit{n2}$ \\
32 & $\mathit{SA} = \mathit{S} \cdot \mathit{n2n}$ & $\mathit{S} \cdot \mathit{n2n} = \mathit{SA}$ \\
33 & $\mathit{Tcoef} = \mathit{L2} - \mathit{zl}$ & $\mathit{Tcoef} + \mathit{zl} = \mathit{L2}$ \\
34 & $\mathit{Tq4} = \mathit{Tcoef} \cdot \mathit{q2}$ & $\mathit{Tcoef} \cdot \mathit{q2} = \mathit{Tq4}$ \\
35 & $\mathit{bm1} = \mathit{b} - 1$ & $\mathit{bm1} + 1 = \mathit{b}$ \\
36 & $\mathit{lbm1} = \mathit{l} \cdot \mathit{bm1}$ & $\mathit{l} \cdot \mathit{bm1} = \mathit{lbm1}$ \\
37 & $\mathit{T2} = \mathit{Tq4} - \mathit{lbm1}$ & $\mathit{T2} + \mathit{lbm1} = \mathit{Tq4}$ \\
38 & $\mathit{b5m2} = \mathit{B} - 4$ & $\mathit{b5m2} + 4 = \mathit{B}$ \\
39 & $\mathit{packed\_target\_shift} = \mathit{l} \cdot \mathit{q4}$ & $\mathit{l} \cdot \mathit{q4} = \mathit{packed\_target\_shift}$ \\
40 & $\mathit{mask\_term} = \mathit{b5m2} \cdot \mathit{packed\_target\_shift}$ & $\mathit{b5m2} \cdot \mathit{packed\_target\_shift} = \mathit{mask\_term}$ \\
41 & $\mathit{T} = \mathit{T2} + \mathit{mask\_term}$ & $\mathit{T2} + \mathit{mask\_term} = \mathit{T}$ \\
42 & $\mathit{n2m1} = \mathit{n2} - 1$ & $\mathit{n2m1} + 1 = \mathit{n2}$ \\
43 & $\mathit{TB} = \mathit{T} \cdot \mathit{n2m1}$ & $\mathit{T} \cdot \mathit{n2m1} = \mathit{TB}$ \\
44 & $\mathit{R7} = \mathit{SA} + \mathit{TB}$ & $\mathit{SA} + \mathit{TB} = \mathit{R7}$ \\
45 & $\mathit{wn2} = \mathit{w} \cdot \mathit{n2}$ & $\mathit{w} \cdot \mathit{n2} = \mathit{wn2}$ \\
46 & $\mathit{sn2} = \mathit{s} \cdot \mathit{n2}$ & $\mathit{s} \cdot \mathit{n2} = \mathit{sn2}$ \\
47 & $\mathit{rsn2} = \mathit{r} \cdot \mathit{sn2}$ & $\mathit{r} \cdot \mathit{sn2} = \mathit{rsn2}$ \\
48 & $\mathit{UM} = \mathit{wn2} \cdot \mathit{rsn2}$ & $\mathit{wn2} \cdot \mathit{rsn2} = \mathit{UM}$ \\
49 & $\mathit{R12} = \mathit{UM} + \mathit{rsn2}$ & $\mathit{UM} + \mathit{rsn2} = \mathit{R12}$ \\
50 & $\mathit{wsq} = \mathit{UM} \cdot \mathit{rsn2}$ & $\mathit{UM} \cdot \mathit{rsn2} = \mathit{wsq}$ \\
51 & $\mathit{R8} = 2 \cdot \mathit{wsq}$ & $2 \cdot \mathit{wsq} = \mathit{R8}$ \\
52 & $\mathit{Pplus1} = \mathit{R8} + 2$ & $\mathit{R8} + 2 = \mathit{Pplus1}$ \\
53 & $\mathit{p2m1} = \mathit{R8} \cdot \mathit{Pplus1}$ & $\mathit{R8} \cdot \mathit{Pplus1} = \mathit{p2m1}$ \\
54 & $\mathit{k2} = \mathit{k} \cdot \mathit{k}$ & $\mathit{k} \cdot \mathit{k} = \mathit{k2}$ \\
55 & $\mathit{P9} = \mathit{p2m1} \cdot \mathit{k2}$ & $\mathit{p2m1} \cdot \mathit{k2} = \mathit{P9}$ \\
56 & $\mathit{L9} = \mathit{P9} + 1$ & $\mathit{P9} + 1 = \mathit{L9}$ \\
57 & $\mathit{R9} = \mathit{tau} \cdot \mathit{tau}$ & $\mathit{tau} \cdot \mathit{tau} = \mathit{R9}$ \\
58 & $\mathit{ksn2} = \mathit{k} \cdot \mathit{sn2}$ & $\mathit{k} \cdot \mathit{sn2} = \mathit{ksn2}$ \\
59 & $\mathit{R10a} = \mathit{ksn2} + \mathit{eta}$ & $\mathit{ksn2} + \mathit{eta} = \mathit{R10a}$ \\
60 & $\mathit{R10b} = \mathit{eta} + \mathit{zeta}$ & $\mathit{eta} + \mathit{zeta} = \mathit{R10b}$ \\
61 & $\mathit{r1} = \mathit{r} + 1$ & $\mathit{r} + 1 = \mathit{r1}$ \\
62 & $\mathit{hpm1} = \mathit{h} \cdot \mathit{R8}$ & $\mathit{h} \cdot \mathit{R8} = \mathit{hpm1}$ \\
63 & $\mathit{R11} = \mathit{r1} + \mathit{hpm1}$ & $\mathit{r1} + \mathit{hpm1} = \mathit{R11}$ \\
64 & $\mathit{tr1} = \mathit{r1} + \mathit{r}$ & $\mathit{r1} + \mathit{r} = \mathit{tr1}$ \\
65 & $\mathit{R13} = \mathit{ka} + \mathit{phi}$ & $\mathit{ka} + \mathit{phi} = \mathit{R13}$ \\
66 & $\mathit{bw} = \mathit{b} \cdot \mathit{w}$ & $\mathit{b} \cdot \mathit{w} = \mathit{bw}$ \\
67 & $\mathit{am2} = \mathit{a} - 2$ & $\mathit{am2} + 2 = \mathit{a}$ \\
68 & $\mathit{cam2} = \mathit{c} \cdot \mathit{am2}$ & $\mathit{c} \cdot \mathit{am2} = \mathit{cam2}$ \\
69 & $\mathit{D1} = \mathit{bw} + \mathit{cam2}$ & $\mathit{bw} + \mathit{cam2} = \mathit{D1}$ \\
70 & $\mathit{a4} = 4 \cdot \mathit{a}$ & $4 \cdot \mathit{a} = \mathit{a4}$ \\
71 & $\mathit{a4m5} = \mathit{a4} - 5$ & $\mathit{a4m5} + 5 = \mathit{a4}$ \\
72 & $\mathit{gam} = \mathit{ga} \cdot \mathit{a4m5}$ & $\mathit{ga} \cdot \mathit{a4m5} = \mathit{gam}$ \\
73 & $\mathit{R14} = \mathit{D1} + \mathit{gam}$ & $\mathit{D1} + \mathit{gam} = \mathit{R14}$ \\
74 & $\mathit{am1} = \mathit{a} - 1$ & $\mathit{am1} + 1 = \mathit{a}$ \\
75 & $\mathit{ap1} = \mathit{a} + 1$ & $\mathit{a} + 1 = \mathit{ap1}$ \\
76 & $\mathit{A} = \mathit{am1} \cdot \mathit{ap1}$ & $\mathit{am1} \cdot \mathit{ap1} = \mathit{A}$ \\
77 & $\mathit{c2} = \mathit{c} \cdot \mathit{c}$ & $\mathit{c} \cdot \mathit{c} = \mathit{c2}$ \\
78 & $\mathit{Ac2} = \mathit{A} \cdot \mathit{c2}$ & $\mathit{A} \cdot \mathit{c2} = \mathit{Ac2}$ \\
79 & $\mathit{R15} = \mathit{Ac2} + 1$ & $\mathit{Ac2} + 1 = \mathit{R15}$ \\
80 & $\mathit{L15} = \mathit{d} \cdot \mathit{d}$ & $\mathit{d} \cdot \mathit{d} = \mathit{L15}$ \\
81 & $\mathit{ic2} = \mathit{i} \cdot \mathit{c2}$ & $\mathit{i} \cdot \mathit{c2} = \mathit{ic2}$ \\
82 & $\mathit{ic22} = \mathit{ic2} \cdot \mathit{ic2}$ & $\mathit{ic2} \cdot \mathit{ic2} = \mathit{ic22}$ \\
83 & $\mathit{AE} = \mathit{A} \cdot \mathit{ic22}$ & $\mathit{A} \cdot \mathit{ic22} = \mathit{AE}$ \\
84 & $\mathit{R16} = \mathit{AE} + 1$ & $\mathit{AE} + 1 = \mathit{R16}$ \\
85 & $\mathit{L16} = \mathit{f} \cdot \mathit{f}$ & $\mathit{f} \cdot \mathit{f} = \mathit{L16}$ \\
86 & $\mathit{of} = \mathit{o} \cdot \mathit{f}$ & $\mathit{o} \cdot \mathit{f} = \mathit{of}$ \\
87 & $\mathit{dof} = \mathit{of} - \mathit{d}$ & $\mathit{dof} + \mathit{d} = \mathit{of}$ \\
88 & $\mathit{L17} = \mathit{dof} \cdot \mathit{dof}$ & $\mathit{dof} \cdot \mathit{dof} = \mathit{L17}$ \\
89 & $\mathit{Gminus1} = \mathit{ap1} \cdot \mathit{AE}$ & $\mathit{ap1} \cdot \mathit{AE} = \mathit{Gminus1}$ \\
90 & $\mathit{Gplus1} = \mathit{Gminus1} + 2$ & $\mathit{Gminus1} + 2 = \mathit{Gplus1}$ \\
91 & $\mathit{G2m1} = \mathit{Gminus1} \cdot \mathit{Gplus1}$ & $\mathit{Gminus1} \cdot \mathit{Gplus1} = \mathit{G2m1}$ \\
92 & $\mathit{jc} = \mathit{j} \cdot \mathit{c}$ & $\mathit{j} \cdot \mathit{c} = \mathit{jc}$ \\
93 & $\mathit{H17} = \mathit{tr1} + \mathit{jc}$ & $\mathit{tr1} + \mathit{jc} = \mathit{H17}$ \\
94 & $\mathit{H2} = \mathit{H17} \cdot \mathit{H17}$ & $\mathit{H17} \cdot \mathit{H17} = \mathit{H2}$ \\
95 & $\mathit{P17} = \mathit{G2m1} \cdot \mathit{H2}$ & $\mathit{G2m1} \cdot \mathit{H2} = \mathit{P17}$ \\
96 & $\mathit{R17} = \mathit{P17} + 1$ & $\mathit{P17} + 1 = \mathit{R17}$ \\
97 & $\mathit{amb} = \mathit{a} - \mathit{B}$ & $\mathit{amb} + \mathit{B} = \mathit{a}$ \\
98 & $\mathit{kamb} = \mathit{ka} \cdot \mathit{amb}$ & $\mathit{ka} \cdot \mathit{amb} = \mathit{kamb}$ \\
99 & $\mathit{qk} = \mathit{q} + \mathit{kamb}$ & $\mathit{q} + \mathit{kamb} = \mathit{qk}$ \\
100 & $\mathit{amb2} = \mathit{amb} \cdot \mathit{amb}$ & $\mathit{amb} \cdot \mathit{amb} = \mathit{amb2}$ \\
101 & $\mathit{Mod} = \mathit{A} - \mathit{amb2}$ & $\mathit{Mod} + \mathit{amb2} = \mathit{A}$ \\
102 & $\mathit{rM} = \mathit{rho} \cdot \mathit{Mod}$ & $\mathit{rho} \cdot \mathit{Mod} = \mathit{rM}$ \\
103 & $\mathit{R18} = \mathit{qk} + \mathit{rM}$ & $\mathit{qk} + \mathit{rM} = \mathit{R18}$ \\
104 & $\mathit{ka2} = \mathit{ka} \cdot \mathit{ka}$ & $\mathit{ka} \cdot \mathit{ka} = \mathit{ka2}$ \\
105 & $\mathit{Aka2} = \mathit{A} \cdot \mathit{ka2}$ & $\mathit{A} \cdot \mathit{ka2} = \mathit{Aka2}$ \\
106 & $\mathit{R19} = \mathit{Aka2} + 1$ & $\mathit{Aka2} + 1 = \mathit{R19}$ \\
107 & $\mathit{L19} = \mathit{mu} \cdot \mathit{mu}$ & $\mathit{mu} \cdot \mathit{mu} = \mathit{L19}$ \\
108 & $\mathit{Dam1} = \mathit{Delta} \cdot \mathit{am1}$ & $\mathit{Delta} \cdot \mathit{am1} = \mathit{Dam1}$ \\
109 & $\mathit{R20} = \mathit{Dam1} + 5^{16}$ & $\mathit{Dam1} + 5^{16} = \mathit{R20}$ \\
\end{longtable}

The twenty-one free equality tests are
\[\begin{array}{@{}l@{}}
\mathit{L1}=\mathit{q2},\quad \mathit{b}=\mathit{bx},\quad \mathit{L2}=\mathit{R2},\\
\mathit{L3}=\mathit{B},\quad \mathit{S2}=\mathit{pack\_rhs},\quad \mathit{n}=\mathit{q8},\\
\mathit{r}=\mathit{R7},\quad \mathit{L9}=\mathit{R9},\quad \mathit{c}=\mathit{R10a},\\
\mathit{k}=\mathit{R10b},\quad \mathit{k}=\mathit{R11},\quad \mathit{a}=\mathit{R12},\\
\mathit{c}=\mathit{R13},\quad \mathit{d}=\mathit{R14},\quad \mathit{L15}=\mathit{R15},\\
\mathit{L16}=\mathit{R16},\quad \mathit{L17}=\mathit{R17},\quad \mathit{mu}=\mathit{R18},\\
\mathit{L19}=\mathit{R19},\quad \mathit{ka}=\mathit{R20},\quad \mathit{S3}=\mathit{sigma}
\end{array}\]


\clearpage
\normalsize
\section{The 107-instruction certificate}\label{app:107}

Use the same fixed input-unit index as in Appendix~\ref{app:109}.
The 34 positive unknowns are the 33 previously listed witnesses,
including $\sigma$, together with $\Omega$. The name
$\mathit{Omega}$ denotes $\Omega$, and the free equality to the
register $\mathit{t2}=Z\lambda-2e$ imposes its positive sign.
The size bound is now $\mathit{S2}+\alpha=q^2$, and the third block is
computed as $\mathit{P2}-\mathit{P1}$.

The first Pell root and index quotient also have new meanings:
$\tau_{\rm old}=2\tau+1$ and $h_{\rm old}=h/2$ in the positive reverse
map proved in Section~\ref{sec:final-two-savings}. The table directly
checks the new equations $\tau(\tau+1)=Q(Q+1)k^2$ and
$k=r+1+hQ$, using the existing register $\mathit{wsq}=Q$.
Neither $2Q$ nor the old root is computed.

The complete checker \file{round12\_1980\_certificate.py} verifies all
107 primitives, all 22 source residuals, the triangular corrections,
and the exact polynomial witness identities. The table is generated by
\file{round12\_1980\_schedule\_table.py}. It has 59 multiplications
and 48 additions, or 114 operations including the seven fixed-numeral
instructions. Theorem~\ref{thm:best107} supplies the composed
positive-domain and encoding proof.

\small
\input{jones1980_theorem5_107_schedule.tex}

\clearpage\normalsize
\section{The 106-instruction certificate}\label{app:106}

Use the modular support of Section~\ref{sec:modular-support}, with
$L=2^{32}$. The positive input named $a$ now denotes $a_0=RY(U+1)$;
the mathematical main Pell parameter is $A=a_0+4$. The base-four
congruence forces $bw=4^{2r+1}$. Theorem~\ref{thm:best106} proves
equivalence with the represented set, including the new ratio estimates
and the reconstruction of all positive Pell witnesses.

The complete checker
\file{round14\_1980\_base\_four\_certificate.py} verifies 59
multiplications and 47 additions, all 22 source residuals, and both
triangular corrections. The table is generated by
\file{round14\_1980\_schedule\_table.py}. Fixed numerals are free.

\small
\input{jones1980_theorem5_106_schedule.tex}

\clearpage\normalsize
\section{The 105-instruction certificate}\label{app:105}

Retain the same fixed modular index and the main parameter $A=a_0+4$.
The positive coordinate named $a$ now denotes
$a_0=(r+1)Y(U+1)$, and $Q=U((r+1)Y)^2$. The first index equation is
$k=h(r+1)$. Theorem~\ref{thm:best105} proves that this forces the exact
Pell index $r+1$ and gives the complete positive-domain equivalence.

The schedule moves the existing register $r1=r+1$ before the scale
product and eliminates the former index-sum register $R11$. The
complete checker
\file{round16\_1980\_index\_multiple\_certificate.py} verifies all
105 primitives: 59 multiplications and 46 additions. It also verifies
all 22 source residuals, including the three changed equations and both
triangular corrections. The table is generated by
\file{round16\_1980\_schedule\_table.py}. There are 34 positive
unknowns, and fixed numerals are free.

\small
% Generated by round16_1980_schedule_table.py; do not edit.
\begin{longtable}{@{}rll@{}}
\toprule
No. & Assignment & Addition/multiplication check\\
\midrule
\endfirsthead
\toprule
No. & Assignment & Addition/multiplication check\\
\midrule
\endhead
\bottomrule
\endfoot
1 & $\mathit{q2} = \mathit{q} \cdot \mathit{q}$ & $\mathit{q} \cdot \mathit{q} = \mathit{q2}$ \\
2 & $\mathit{q4} = \mathit{q2} \cdot \mathit{q2}$ & $\mathit{q2} \cdot \mathit{q2} = \mathit{q4}$ \\
3 & $\mathit{q8} = \mathit{q4} \cdot \mathit{q4}$ & $\mathit{q4} \cdot \mathit{q4} = \mathit{q8}$ \\
4 & $\mathit{b2} = \mathit{b} \cdot \mathit{b}$ & $\mathit{b} \cdot \mathit{b} = \mathit{b2}$ \\
5 & $\mathit{B} = \mathit{H} \cdot \mathit{b2}$ & $\mathit{H} \cdot \mathit{b2} = \mathit{B}$ \\
6 & $\mathit{bx} = \mathit{x} + \mathit{beta}$ & $\mathit{x} + \mathit{beta} = \mathit{bx}$ \\
7 & $\mathit{L2} = \mathit{la} + \mathit{q2}$ & $\mathit{la} + \mathit{q2} = \mathit{L2}$ \\
8 & $\mathit{Lam} = \mathit{la} \cdot \mathit{B}$ & $\mathit{la} \cdot \mathit{B} = \mathit{Lam}$ \\
9 & $\mathit{R2} = \mathit{Lam} + 1$ & $\mathit{Lam} + 1 = \mathit{R2}$ \\
10 & $\mathit{L3} = \mathit{th} + \mathit{Z}$ & $\mathit{th} + \mathit{Z} = \mathit{L3}$ \\
11 & $\mathit{eq} = \mathit{e} \cdot \mathit{q}$ & $\mathit{e} \cdot \mathit{q} = \mathit{eq}$ \\
12 & $\mathit{S2} = \mathit{l} + \mathit{eq}$ & $\mathit{l} + \mathit{eq} = \mathit{S2}$ \\
13 & $\mathit{pack\_tth} = \mathit{t} \cdot \mathit{th}$ & $\mathit{t} \cdot \mathit{th} = \mathit{pack\_tth}$ \\
14 & $\mathit{pack\_rhs} = \mathit{V} + \mathit{pack\_tth}$ & $\mathit{V} + \mathit{pack\_tth} = \mathit{pack\_rhs}$ \\
15 & $\mathit{zl} = \mathit{Z} \cdot \mathit{la}$ & $\mathit{Z} \cdot \mathit{la} = \mathit{zl}$ \\
16 & $\mathit{twice\_e} = 2 \cdot \mathit{e}$ & $2 \cdot \mathit{e} = \mathit{twice\_e}$ \\
17 & $\mathit{t2} = \mathit{zl} - \mathit{twice\_e}$ & $\mathit{t2} + \mathit{twice\_e} = \mathit{zl}$ \\
18 & $\mathit{C} = \mathit{x} + \mathit{g}$ & $\mathit{x} + \mathit{g} = \mathit{C}$ \\
19 & $\mathit{C2} = \mathit{C} \cdot \mathit{C}$ & $\mathit{C} \cdot \mathit{C} = \mathit{C2}$ \\
20 & $\mathit{L1} = \mathit{S2} + \mathit{al}$ & $\mathit{S2} + \mathit{al} = \mathit{L1}$ \\
21 & $\mathit{P1} = \mathit{t2} \cdot \mathit{C2}$ & $\mathit{t2} \cdot \mathit{C2} = \mathit{P1}$ \\
22 & $\mathit{qp1} = \mathit{q} + 1$ & $\mathit{q} + 1 = \mathit{qp1}$ \\
23 & $\mathit{P2} = \mathit{Lam} \cdot \mathit{qp1}$ & $\mathit{Lam} \cdot \mathit{qp1} = \mathit{P2}$ \\
24 & $\mathit{S3} = \mathit{P2} - \mathit{P1}$ & $\mathit{S3} + \mathit{P1} = \mathit{P2}$ \\
25 & $\mathit{S3q4} = \mathit{S3} \cdot \mathit{q2}$ & $\mathit{S3} \cdot \mathit{q2} = \mathit{S3q4}$ \\
26 & $\mathit{Sin} = \mathit{S2} + \mathit{S3q4}$ & $\mathit{S2} + \mathit{S3q4} = \mathit{Sin}$ \\
27 & $\mathit{Sq4} = \mathit{Sin} \cdot \mathit{q2}$ & $\mathit{Sin} \cdot \mathit{q2} = \mathit{Sq4}$ \\
28 & $\mathit{S} = \mathit{g} + \mathit{Sq4}$ & $\mathit{g} + \mathit{Sq4} = \mathit{S}$ \\
29 & $\mathit{n2} = \mathit{n} \cdot \mathit{n}$ & $\mathit{n} \cdot \mathit{n} = \mathit{n2}$ \\
30 & $\mathit{n2n} = \mathit{n2} - \mathit{n}$ & $\mathit{n2n} + \mathit{n} = \mathit{n2}$ \\
31 & $\mathit{SA} = \mathit{S} \cdot \mathit{n2n}$ & $\mathit{S} \cdot \mathit{n2n} = \mathit{SA}$ \\
32 & $\mathit{Tcoef} = \mathit{L2} - \mathit{zl}$ & $\mathit{Tcoef} + \mathit{zl} = \mathit{L2}$ \\
33 & $\mathit{Tq4} = \mathit{Tcoef} \cdot \mathit{q2}$ & $\mathit{Tcoef} \cdot \mathit{q2} = \mathit{Tq4}$ \\
34 & $\mathit{bm1} = \mathit{b} - 1$ & $\mathit{bm1} + 1 = \mathit{b}$ \\
35 & $\mathit{lbm1} = \mathit{l} \cdot \mathit{bm1}$ & $\mathit{l} \cdot \mathit{bm1} = \mathit{lbm1}$ \\
36 & $\mathit{T2} = \mathit{Tq4} - \mathit{lbm1}$ & $\mathit{T2} + \mathit{lbm1} = \mathit{Tq4}$ \\
37 & $\mathit{b5m2} = \mathit{B} - 4$ & $\mathit{b5m2} + 4 = \mathit{B}$ \\
38 & $\mathit{packed\_target\_shift} = \mathit{l} \cdot \mathit{q4}$ & $\mathit{l} \cdot \mathit{q4} = \mathit{packed\_target\_shift}$ \\
39 & $\mathit{mask\_term} = \mathit{b5m2} \cdot \mathit{packed\_target\_shift}$ & $\mathit{b5m2} \cdot \mathit{packed\_target\_shift} = \mathit{mask\_term}$ \\
40 & $\mathit{T} = \mathit{T2} + \mathit{mask\_term}$ & $\mathit{T2} + \mathit{mask\_term} = \mathit{T}$ \\
41 & $\mathit{n2m1} = \mathit{n2} - 1$ & $\mathit{n2m1} + 1 = \mathit{n2}$ \\
42 & $\mathit{TB} = \mathit{T} \cdot \mathit{n2m1}$ & $\mathit{T} \cdot \mathit{n2m1} = \mathit{TB}$ \\
43 & $\mathit{R7} = \mathit{SA} + \mathit{TB}$ & $\mathit{SA} + \mathit{TB} = \mathit{R7}$ \\
44 & $\mathit{wn2} = \mathit{w} \cdot \mathit{n2}$ & $\mathit{w} \cdot \mathit{n2} = \mathit{wn2}$ \\
45 & $\mathit{sn2} = \mathit{s} \cdot \mathit{n2}$ & $\mathit{s} \cdot \mathit{n2} = \mathit{sn2}$ \\
46 & $\mathit{r1} = \mathit{r} + 1$ & $\mathit{r} + 1 = \mathit{r1}$ \\
47 & $\mathit{rsn2} = \mathit{r1} \cdot \mathit{sn2}$ & $\mathit{r1} \cdot \mathit{sn2} = \mathit{rsn2}$ \\
48 & $\mathit{UM} = \mathit{wn2} \cdot \mathit{rsn2}$ & $\mathit{wn2} \cdot \mathit{rsn2} = \mathit{UM}$ \\
49 & $\mathit{R12} = \mathit{UM} + \mathit{rsn2}$ & $\mathit{UM} + \mathit{rsn2} = \mathit{R12}$ \\
50 & $\mathit{wsq} = \mathit{UM} \cdot \mathit{rsn2}$ & $\mathit{UM} \cdot \mathit{rsn2} = \mathit{wsq}$ \\
51 & $\mathit{Qplus1} = \mathit{wsq} + 1$ & $\mathit{wsq} + 1 = \mathit{Qplus1}$ \\
52 & $\mathit{QQplus1} = \mathit{wsq} \cdot \mathit{Qplus1}$ & $\mathit{wsq} \cdot \mathit{Qplus1} = \mathit{QQplus1}$ \\
53 & $\mathit{k2} = \mathit{k} \cdot \mathit{k}$ & $\mathit{k} \cdot \mathit{k} = \mathit{k2}$ \\
54 & $\mathit{L9} = \mathit{QQplus1} \cdot \mathit{k2}$ & $\mathit{QQplus1} \cdot \mathit{k2} = \mathit{L9}$ \\
55 & $\mathit{tauplus1} = \mathit{tau} + 1$ & $\mathit{tau} + 1 = \mathit{tauplus1}$ \\
56 & $\mathit{R9} = \mathit{tau} \cdot \mathit{tauplus1}$ & $\mathit{tau} \cdot \mathit{tauplus1} = \mathit{R9}$ \\
57 & $\mathit{ksn2} = \mathit{k} \cdot \mathit{sn2}$ & $\mathit{k} \cdot \mathit{sn2} = \mathit{ksn2}$ \\
58 & $\mathit{R10a} = \mathit{ksn2} + \mathit{eta}$ & $\mathit{ksn2} + \mathit{eta} = \mathit{R10a}$ \\
59 & $\mathit{R10b} = \mathit{eta} + \mathit{zeta}$ & $\mathit{eta} + \mathit{zeta} = \mathit{R10b}$ \\
60 & $\mathit{hpm1} = \mathit{h} \cdot \mathit{r1}$ & $\mathit{h} \cdot \mathit{r1} = \mathit{hpm1}$ \\
61 & $\mathit{tr1} = \mathit{r1} + \mathit{r}$ & $\mathit{r1} + \mathit{r} = \mathit{tr1}$ \\
62 & $\mathit{R13} = \mathit{ka} + \mathit{phi}$ & $\mathit{ka} + \mathit{phi} = \mathit{R13}$ \\
63 & $\mathit{bw} = \mathit{b} \cdot \mathit{w}$ & $\mathit{b} \cdot \mathit{w} = \mathit{bw}$ \\
64 & $\mathit{cam2} = \mathit{c} \cdot \mathit{a}$ & $\mathit{c} \cdot \mathit{a} = \mathit{cam2}$ \\
65 & $\mathit{D1} = \mathit{bw} + \mathit{cam2}$ & $\mathit{bw} + \mathit{cam2} = \mathit{D1}$ \\
66 & $\mathit{am1} = \mathit{a} + 3$ & $\mathit{a} + 3 = \mathit{am1}$ \\
67 & $\mathit{a4} = 8 \cdot \mathit{a}$ & $8 \cdot \mathit{a} = \mathit{a4}$ \\
68 & $\mathit{a4m5} = \mathit{a4} + 15$ & $\mathit{a4} + 15 = \mathit{a4m5}$ \\
69 & $\mathit{gam} = \mathit{ga} \cdot \mathit{a4m5}$ & $\mathit{ga} \cdot \mathit{a4m5} = \mathit{gam}$ \\
70 & $\mathit{R14} = \mathit{D1} + \mathit{gam}$ & $\mathit{D1} + \mathit{gam} = \mathit{R14}$ \\
71 & $\mathit{ap1} = \mathit{am1} + 2$ & $\mathit{am1} + 2 = \mathit{ap1}$ \\
72 & $\mathit{A} = \mathit{am1} \cdot \mathit{ap1}$ & $\mathit{am1} \cdot \mathit{ap1} = \mathit{A}$ \\
73 & $\mathit{c2} = \mathit{c} \cdot \mathit{c}$ & $\mathit{c} \cdot \mathit{c} = \mathit{c2}$ \\
74 & $\mathit{Ac2} = \mathit{A} \cdot \mathit{c2}$ & $\mathit{A} \cdot \mathit{c2} = \mathit{Ac2}$ \\
75 & $\mathit{R15} = \mathit{Ac2} + 1$ & $\mathit{Ac2} + 1 = \mathit{R15}$ \\
76 & $\mathit{L15} = \mathit{d} \cdot \mathit{d}$ & $\mathit{d} \cdot \mathit{d} = \mathit{L15}$ \\
77 & $\mathit{ic2} = \mathit{i} \cdot \mathit{c2}$ & $\mathit{i} \cdot \mathit{c2} = \mathit{ic2}$ \\
78 & $\mathit{ic22} = \mathit{ic2} \cdot \mathit{ic2}$ & $\mathit{ic2} \cdot \mathit{ic2} = \mathit{ic22}$ \\
79 & $\mathit{AE} = \mathit{A} \cdot \mathit{ic22}$ & $\mathit{A} \cdot \mathit{ic22} = \mathit{AE}$ \\
80 & $\mathit{R16} = \mathit{AE} + 1$ & $\mathit{AE} + 1 = \mathit{R16}$ \\
81 & $\mathit{L16} = \mathit{f} \cdot \mathit{f}$ & $\mathit{f} \cdot \mathit{f} = \mathit{L16}$ \\
82 & $\mathit{of} = \mathit{o} \cdot \mathit{f}$ & $\mathit{o} \cdot \mathit{f} = \mathit{of}$ \\
83 & $\mathit{dof} = \mathit{of} - \mathit{d}$ & $\mathit{dof} + \mathit{d} = \mathit{of}$ \\
84 & $\mathit{L17} = \mathit{dof} \cdot \mathit{dof}$ & $\mathit{dof} \cdot \mathit{dof} = \mathit{L17}$ \\
85 & $\mathit{Gminus1} = \mathit{ap1} \cdot \mathit{AE}$ & $\mathit{ap1} \cdot \mathit{AE} = \mathit{Gminus1}$ \\
86 & $\mathit{Gplus1} = \mathit{Gminus1} + 2$ & $\mathit{Gminus1} + 2 = \mathit{Gplus1}$ \\
87 & $\mathit{G2m1} = \mathit{Gminus1} \cdot \mathit{Gplus1}$ & $\mathit{Gminus1} \cdot \mathit{Gplus1} = \mathit{G2m1}$ \\
88 & $\mathit{jc} = \mathit{j} \cdot \mathit{c}$ & $\mathit{j} \cdot \mathit{c} = \mathit{jc}$ \\
89 & $\mathit{H17} = \mathit{tr1} + \mathit{jc}$ & $\mathit{tr1} + \mathit{jc} = \mathit{H17}$ \\
90 & $\mathit{H2} = \mathit{H17} \cdot \mathit{H17}$ & $\mathit{H17} \cdot \mathit{H17} = \mathit{H2}$ \\
91 & $\mathit{P17} = \mathit{G2m1} \cdot \mathit{H2}$ & $\mathit{G2m1} \cdot \mathit{H2} = \mathit{P17}$ \\
92 & $\mathit{R17} = \mathit{P17} + 1$ & $\mathit{P17} + 1 = \mathit{R17}$ \\
93 & $\mathit{amb} = \mathit{a} - \mathit{b5m2}$ & $\mathit{amb} + \mathit{b5m2} = \mathit{a}$ \\
94 & $\mathit{kamb} = \mathit{ka} \cdot \mathit{amb}$ & $\mathit{ka} \cdot \mathit{amb} = \mathit{kamb}$ \\
95 & $\mathit{qk} = \mathit{q} + \mathit{kamb}$ & $\mathit{q} + \mathit{kamb} = \mathit{qk}$ \\
96 & $\mathit{amb2} = \mathit{amb} \cdot \mathit{amb}$ & $\mathit{amb} \cdot \mathit{amb} = \mathit{amb2}$ \\
97 & $\mathit{Mod} = \mathit{A} - \mathit{amb2}$ & $\mathit{Mod} + \mathit{amb2} = \mathit{A}$ \\
98 & $\mathit{rM} = \mathit{rho} \cdot \mathit{Mod}$ & $\mathit{rho} \cdot \mathit{Mod} = \mathit{rM}$ \\
99 & $\mathit{R18} = \mathit{qk} + \mathit{rM}$ & $\mathit{qk} + \mathit{rM} = \mathit{R18}$ \\
100 & $\mathit{ka2} = \mathit{ka} \cdot \mathit{ka}$ & $\mathit{ka} \cdot \mathit{ka} = \mathit{ka2}$ \\
101 & $\mathit{Aka2} = \mathit{A} \cdot \mathit{ka2}$ & $\mathit{A} \cdot \mathit{ka2} = \mathit{Aka2}$ \\
102 & $\mathit{R19} = \mathit{Aka2} + 1$ & $\mathit{Aka2} + 1 = \mathit{R19}$ \\
103 & $\mathit{L19} = \mathit{mu} \cdot \mathit{mu}$ & $\mathit{mu} \cdot \mathit{mu} = \mathit{L19}$ \\
104 & $\mathit{Dam1} = \mathit{Delta} \cdot \mathit{am1}$ & $\mathit{Delta} \cdot \mathit{am1} = \mathit{Dam1}$ \\
105 & $\mathit{R20} = \mathit{Dam1} + 2^{32}$ & $\mathit{Dam1} + 2^{32} = \mathit{R20}$ \\
\end{longtable}

The twenty-two free equality tests are
\[\begin{array}{@{}l@{}}
\mathit{L1}=\mathit{q2},\quad \mathit{b}=\mathit{bx},\quad \mathit{L2}=\mathit{R2},\\
\mathit{L3}=\mathit{B},\quad \mathit{S2}=\mathit{pack\_rhs},\quad \mathit{n}=\mathit{q8},\\
\mathit{r}=\mathit{R7},\quad \mathit{L9}=\mathit{R9},\quad \mathit{c}=\mathit{R10a},\\
\mathit{k}=\mathit{R10b},\quad \mathit{k}=\mathit{hpm1},\quad \mathit{a}=\mathit{R12},\\
\mathit{c}=\mathit{R13},\quad \mathit{d}=\mathit{R14},\quad \mathit{L15}=\mathit{R15},\\
\mathit{L16}=\mathit{R16},\quad \mathit{L17}=\mathit{R17},\quad \mathit{mu}=\mathit{R18},\\
\mathit{L19}=\mathit{R19},\quad \mathit{ka}=\mathit{R20},\quad \mathit{S3}=\mathit{sigma},\\
\mathit{t2}=\mathit{Omega}
\end{array}\]


\clearpage\normalsize
\section{The 104-instruction certificate}\label{app:104}

Retain the fixed modular index, with $L=2^{32}$, and put
$U=wn^2$, $Y_P=sn^2$, $D=UY_P$. The supplied positive coordinate
$a$ denotes $a_0=D+Y_P$, and the main Pell parameter is $A=a_0+4$.
The first norm and index equation are
\[
 \tau(\tau+1)=(D^2+U)(Y_Pk)^2,\qquad k=r+1+hD.
\]
The products $D$ and $Y_Pk$ are shared with the main parameter and
interval equations. Theorem~\ref{thm:best104} proves the exact index,
positive witness equivalence, and full decoding.

The complete checker
\file{round19\_1980\_shared\_ratio\_product\_certificate.py}
verifies all 104 primitives, all 22 source residuals, both triangular
corrections and the forward witness-map identity. The table is generated
by \file{round19\_1980\_schedule\_table.py}. It contains 57
multiplications and 47 additions. There are 34 positive unknowns;
fixed numerals and equality tests are free.

\small
% Generated by round19_1980_schedule_table.py; do not edit.
\begin{longtable}{@{}rll@{}}
\toprule
No. & Assignment & Addition/multiplication check\\
\midrule
\endfirsthead
\toprule
No. & Assignment & Addition/multiplication check\\
\midrule
\endhead
\bottomrule
\endfoot
1 & $\mathit{q2} = \mathit{q} \cdot \mathit{q}$ & $\mathit{q} \cdot \mathit{q} = \mathit{q2}$ \\
2 & $\mathit{q4} = \mathit{q2} \cdot \mathit{q2}$ & $\mathit{q2} \cdot \mathit{q2} = \mathit{q4}$ \\
3 & $\mathit{q8} = \mathit{q4} \cdot \mathit{q4}$ & $\mathit{q4} \cdot \mathit{q4} = \mathit{q8}$ \\
4 & $\mathit{b2} = \mathit{b} \cdot \mathit{b}$ & $\mathit{b} \cdot \mathit{b} = \mathit{b2}$ \\
5 & $\mathit{B} = \mathit{H} \cdot \mathit{b2}$ & $\mathit{H} \cdot \mathit{b2} = \mathit{B}$ \\
6 & $\mathit{bx} = \mathit{x} + \mathit{beta}$ & $\mathit{x} + \mathit{beta} = \mathit{bx}$ \\
7 & $\mathit{L2} = \mathit{la} + \mathit{q2}$ & $\mathit{la} + \mathit{q2} = \mathit{L2}$ \\
8 & $\mathit{Lam} = \mathit{la} \cdot \mathit{B}$ & $\mathit{la} \cdot \mathit{B} = \mathit{Lam}$ \\
9 & $\mathit{R2} = \mathit{Lam} + 1$ & $\mathit{Lam} + 1 = \mathit{R2}$ \\
10 & $\mathit{L3} = \mathit{th} + \mathit{Z}$ & $\mathit{th} + \mathit{Z} = \mathit{L3}$ \\
11 & $\mathit{eq} = \mathit{e} \cdot \mathit{q}$ & $\mathit{e} \cdot \mathit{q} = \mathit{eq}$ \\
12 & $\mathit{S2} = \mathit{l} + \mathit{eq}$ & $\mathit{l} + \mathit{eq} = \mathit{S2}$ \\
13 & $\mathit{pack\_tth} = \mathit{t} \cdot \mathit{th}$ & $\mathit{t} \cdot \mathit{th} = \mathit{pack\_tth}$ \\
14 & $\mathit{pack\_rhs} = \mathit{V} + \mathit{pack\_tth}$ & $\mathit{V} + \mathit{pack\_tth} = \mathit{pack\_rhs}$ \\
15 & $\mathit{zl} = \mathit{Z} \cdot \mathit{la}$ & $\mathit{Z} \cdot \mathit{la} = \mathit{zl}$ \\
16 & $\mathit{twice\_e} = 2 \cdot \mathit{e}$ & $2 \cdot \mathit{e} = \mathit{twice\_e}$ \\
17 & $\mathit{t2} = \mathit{zl} - \mathit{twice\_e}$ & $\mathit{t2} + \mathit{twice\_e} = \mathit{zl}$ \\
18 & $\mathit{C} = \mathit{x} + \mathit{g}$ & $\mathit{x} + \mathit{g} = \mathit{C}$ \\
19 & $\mathit{C2} = \mathit{C} \cdot \mathit{C}$ & $\mathit{C} \cdot \mathit{C} = \mathit{C2}$ \\
20 & $\mathit{L1} = \mathit{S2} + \mathit{al}$ & $\mathit{S2} + \mathit{al} = \mathit{L1}$ \\
21 & $\mathit{P1} = \mathit{t2} \cdot \mathit{C2}$ & $\mathit{t2} \cdot \mathit{C2} = \mathit{P1}$ \\
22 & $\mathit{qp1} = \mathit{q} + 1$ & $\mathit{q} + 1 = \mathit{qp1}$ \\
23 & $\mathit{P2} = \mathit{Lam} \cdot \mathit{qp1}$ & $\mathit{Lam} \cdot \mathit{qp1} = \mathit{P2}$ \\
24 & $\mathit{S3} = \mathit{P2} - \mathit{P1}$ & $\mathit{S3} + \mathit{P1} = \mathit{P2}$ \\
25 & $\mathit{S3q4} = \mathit{S3} \cdot \mathit{q2}$ & $\mathit{S3} \cdot \mathit{q2} = \mathit{S3q4}$ \\
26 & $\mathit{Sin} = \mathit{S2} + \mathit{S3q4}$ & $\mathit{S2} + \mathit{S3q4} = \mathit{Sin}$ \\
27 & $\mathit{Sq4} = \mathit{Sin} \cdot \mathit{q2}$ & $\mathit{Sin} \cdot \mathit{q2} = \mathit{Sq4}$ \\
28 & $\mathit{S} = \mathit{g} + \mathit{Sq4}$ & $\mathit{g} + \mathit{Sq4} = \mathit{S}$ \\
29 & $\mathit{n2} = \mathit{n} \cdot \mathit{n}$ & $\mathit{n} \cdot \mathit{n} = \mathit{n2}$ \\
30 & $\mathit{n2n} = \mathit{n2} - \mathit{n}$ & $\mathit{n2n} + \mathit{n} = \mathit{n2}$ \\
31 & $\mathit{SA} = \mathit{S} \cdot \mathit{n2n}$ & $\mathit{S} \cdot \mathit{n2n} = \mathit{SA}$ \\
32 & $\mathit{Tcoef} = \mathit{L2} - \mathit{zl}$ & $\mathit{Tcoef} + \mathit{zl} = \mathit{L2}$ \\
33 & $\mathit{Tq4} = \mathit{Tcoef} \cdot \mathit{q2}$ & $\mathit{Tcoef} \cdot \mathit{q2} = \mathit{Tq4}$ \\
34 & $\mathit{bm1} = \mathit{b} - 1$ & $\mathit{bm1} + 1 = \mathit{b}$ \\
35 & $\mathit{lbm1} = \mathit{l} \cdot \mathit{bm1}$ & $\mathit{l} \cdot \mathit{bm1} = \mathit{lbm1}$ \\
36 & $\mathit{T2} = \mathit{Tq4} - \mathit{lbm1}$ & $\mathit{T2} + \mathit{lbm1} = \mathit{Tq4}$ \\
37 & $\mathit{b5m2} = \mathit{B} - 4$ & $\mathit{b5m2} + 4 = \mathit{B}$ \\
38 & $\mathit{packed\_target\_shift} = \mathit{l} \cdot \mathit{q4}$ & $\mathit{l} \cdot \mathit{q4} = \mathit{packed\_target\_shift}$ \\
39 & $\mathit{mask\_term} = \mathit{b5m2} \cdot \mathit{packed\_target\_shift}$ & $\mathit{b5m2} \cdot \mathit{packed\_target\_shift} = \mathit{mask\_term}$ \\
40 & $\mathit{T} = \mathit{T2} + \mathit{mask\_term}$ & $\mathit{T2} + \mathit{mask\_term} = \mathit{T}$ \\
41 & $\mathit{n2m1} = \mathit{n2} - 1$ & $\mathit{n2m1} + 1 = \mathit{n2}$ \\
42 & $\mathit{TB} = \mathit{T} \cdot \mathit{n2m1}$ & $\mathit{T} \cdot \mathit{n2m1} = \mathit{TB}$ \\
43 & $\mathit{R7} = \mathit{SA} + \mathit{TB}$ & $\mathit{SA} + \mathit{TB} = \mathit{R7}$ \\
44 & $\mathit{wn2} = \mathit{w} \cdot \mathit{n2}$ & $\mathit{w} \cdot \mathit{n2} = \mathit{wn2}$ \\
45 & $\mathit{sn2} = \mathit{s} \cdot \mathit{n2}$ & $\mathit{s} \cdot \mathit{n2} = \mathit{sn2}$ \\
46 & $\mathit{UM} = \mathit{wn2} \cdot \mathit{sn2}$ & $\mathit{wn2} \cdot \mathit{sn2} = \mathit{UM}$ \\
47 & $\mathit{R12} = \mathit{UM} + \mathit{sn2}$ & $\mathit{UM} + \mathit{sn2} = \mathit{R12}$ \\
48 & $\mathit{ksn2} = \mathit{k} \cdot \mathit{sn2}$ & $\mathit{k} \cdot \mathit{sn2} = \mathit{ksn2}$ \\
49 & $\mathit{UM2} = \mathit{UM} \cdot \mathit{UM}$ & $\mathit{UM} \cdot \mathit{UM} = \mathit{UM2}$ \\
50 & $\mathit{scaled\_norm\_coefficient} = \mathit{UM2} + \mathit{wn2}$ & $\mathit{UM2} + \mathit{wn2} = \mathit{scaled\_norm\_coefficient}$ \\
51 & $\mathit{ratio\_product2} = \mathit{ksn2} \cdot \mathit{ksn2}$ & $\mathit{ksn2} \cdot \mathit{ksn2} = \mathit{ratio\_product2}$ \\
52 & $\mathit{L9} = \mathit{scaled\_norm\_coefficient} \cdot \mathit{ratio\_product2}$ & $\mathit{scaled\_norm\_coefficient} \cdot \mathit{ratio\_product2} = \mathit{L9}$ \\
53 & $\mathit{tauplus1} = \mathit{tau} + 1$ & $\mathit{tau} + 1 = \mathit{tauplus1}$ \\
54 & $\mathit{R9} = \mathit{tau} \cdot \mathit{tauplus1}$ & $\mathit{tau} \cdot \mathit{tauplus1} = \mathit{R9}$ \\
55 & $\mathit{R10a} = \mathit{ksn2} + \mathit{eta}$ & $\mathit{ksn2} + \mathit{eta} = \mathit{R10a}$ \\
56 & $\mathit{R10b} = \mathit{eta} + \mathit{zeta}$ & $\mathit{eta} + \mathit{zeta} = \mathit{R10b}$ \\
57 & $\mathit{r1} = \mathit{r} + 1$ & $\mathit{r} + 1 = \mathit{r1}$ \\
58 & $\mathit{hpm1} = \mathit{h} \cdot \mathit{UM}$ & $\mathit{h} \cdot \mathit{UM} = \mathit{hpm1}$ \\
59 & $\mathit{R11} = \mathit{r1} + \mathit{hpm1}$ & $\mathit{r1} + \mathit{hpm1} = \mathit{R11}$ \\
60 & $\mathit{tr1} = \mathit{r1} + \mathit{r}$ & $\mathit{r1} + \mathit{r} = \mathit{tr1}$ \\
61 & $\mathit{R13} = \mathit{ka} + \mathit{phi}$ & $\mathit{ka} + \mathit{phi} = \mathit{R13}$ \\
62 & $\mathit{bw} = \mathit{b} \cdot \mathit{w}$ & $\mathit{b} \cdot \mathit{w} = \mathit{bw}$ \\
63 & $\mathit{cam2} = \mathit{c} \cdot \mathit{a}$ & $\mathit{c} \cdot \mathit{a} = \mathit{cam2}$ \\
64 & $\mathit{D1} = \mathit{bw} + \mathit{cam2}$ & $\mathit{bw} + \mathit{cam2} = \mathit{D1}$ \\
65 & $\mathit{am1} = \mathit{a} + 3$ & $\mathit{a} + 3 = \mathit{am1}$ \\
66 & $\mathit{a4} = 8 \cdot \mathit{a}$ & $8 \cdot \mathit{a} = \mathit{a4}$ \\
67 & $\mathit{a4m5} = \mathit{a4} + 15$ & $\mathit{a4} + 15 = \mathit{a4m5}$ \\
68 & $\mathit{gam} = \mathit{ga} \cdot \mathit{a4m5}$ & $\mathit{ga} \cdot \mathit{a4m5} = \mathit{gam}$ \\
69 & $\mathit{R14} = \mathit{D1} + \mathit{gam}$ & $\mathit{D1} + \mathit{gam} = \mathit{R14}$ \\
70 & $\mathit{ap1} = \mathit{am1} + 2$ & $\mathit{am1} + 2 = \mathit{ap1}$ \\
71 & $\mathit{A} = \mathit{am1} \cdot \mathit{ap1}$ & $\mathit{am1} \cdot \mathit{ap1} = \mathit{A}$ \\
72 & $\mathit{c2} = \mathit{c} \cdot \mathit{c}$ & $\mathit{c} \cdot \mathit{c} = \mathit{c2}$ \\
73 & $\mathit{Ac2} = \mathit{A} \cdot \mathit{c2}$ & $\mathit{A} \cdot \mathit{c2} = \mathit{Ac2}$ \\
74 & $\mathit{R15} = \mathit{Ac2} + 1$ & $\mathit{Ac2} + 1 = \mathit{R15}$ \\
75 & $\mathit{L15} = \mathit{d} \cdot \mathit{d}$ & $\mathit{d} \cdot \mathit{d} = \mathit{L15}$ \\
76 & $\mathit{ic2} = \mathit{i} \cdot \mathit{c2}$ & $\mathit{i} \cdot \mathit{c2} = \mathit{ic2}$ \\
77 & $\mathit{ic22} = \mathit{ic2} \cdot \mathit{ic2}$ & $\mathit{ic2} \cdot \mathit{ic2} = \mathit{ic22}$ \\
78 & $\mathit{AE} = \mathit{A} \cdot \mathit{ic22}$ & $\mathit{A} \cdot \mathit{ic22} = \mathit{AE}$ \\
79 & $\mathit{R16} = \mathit{AE} + 1$ & $\mathit{AE} + 1 = \mathit{R16}$ \\
80 & $\mathit{L16} = \mathit{f} \cdot \mathit{f}$ & $\mathit{f} \cdot \mathit{f} = \mathit{L16}$ \\
81 & $\mathit{of} = \mathit{o} \cdot \mathit{f}$ & $\mathit{o} \cdot \mathit{f} = \mathit{of}$ \\
82 & $\mathit{dof} = \mathit{of} - \mathit{d}$ & $\mathit{dof} + \mathit{d} = \mathit{of}$ \\
83 & $\mathit{L17} = \mathit{dof} \cdot \mathit{dof}$ & $\mathit{dof} \cdot \mathit{dof} = \mathit{L17}$ \\
84 & $\mathit{Gminus1} = \mathit{ap1} \cdot \mathit{AE}$ & $\mathit{ap1} \cdot \mathit{AE} = \mathit{Gminus1}$ \\
85 & $\mathit{Gplus1} = \mathit{Gminus1} + 2$ & $\mathit{Gminus1} + 2 = \mathit{Gplus1}$ \\
86 & $\mathit{G2m1} = \mathit{Gminus1} \cdot \mathit{Gplus1}$ & $\mathit{Gminus1} \cdot \mathit{Gplus1} = \mathit{G2m1}$ \\
87 & $\mathit{jc} = \mathit{j} \cdot \mathit{c}$ & $\mathit{j} \cdot \mathit{c} = \mathit{jc}$ \\
88 & $\mathit{H17} = \mathit{tr1} + \mathit{jc}$ & $\mathit{tr1} + \mathit{jc} = \mathit{H17}$ \\
89 & $\mathit{H2} = \mathit{H17} \cdot \mathit{H17}$ & $\mathit{H17} \cdot \mathit{H17} = \mathit{H2}$ \\
90 & $\mathit{P17} = \mathit{G2m1} \cdot \mathit{H2}$ & $\mathit{G2m1} \cdot \mathit{H2} = \mathit{P17}$ \\
91 & $\mathit{R17} = \mathit{P17} + 1$ & $\mathit{P17} + 1 = \mathit{R17}$ \\
92 & $\mathit{amb} = \mathit{a} - \mathit{b5m2}$ & $\mathit{amb} + \mathit{b5m2} = \mathit{a}$ \\
93 & $\mathit{kamb} = \mathit{ka} \cdot \mathit{amb}$ & $\mathit{ka} \cdot \mathit{amb} = \mathit{kamb}$ \\
94 & $\mathit{qk} = \mathit{q} + \mathit{kamb}$ & $\mathit{q} + \mathit{kamb} = \mathit{qk}$ \\
95 & $\mathit{amb2} = \mathit{amb} \cdot \mathit{amb}$ & $\mathit{amb} \cdot \mathit{amb} = \mathit{amb2}$ \\
96 & $\mathit{Mod} = \mathit{A} - \mathit{amb2}$ & $\mathit{Mod} + \mathit{amb2} = \mathit{A}$ \\
97 & $\mathit{rM} = \mathit{rho} \cdot \mathit{Mod}$ & $\mathit{rho} \cdot \mathit{Mod} = \mathit{rM}$ \\
98 & $\mathit{R18} = \mathit{qk} + \mathit{rM}$ & $\mathit{qk} + \mathit{rM} = \mathit{R18}$ \\
99 & $\mathit{ka2} = \mathit{ka} \cdot \mathit{ka}$ & $\mathit{ka} \cdot \mathit{ka} = \mathit{ka2}$ \\
100 & $\mathit{Aka2} = \mathit{A} \cdot \mathit{ka2}$ & $\mathit{A} \cdot \mathit{ka2} = \mathit{Aka2}$ \\
101 & $\mathit{R19} = \mathit{Aka2} + 1$ & $\mathit{Aka2} + 1 = \mathit{R19}$ \\
102 & $\mathit{L19} = \mathit{mu} \cdot \mathit{mu}$ & $\mathit{mu} \cdot \mathit{mu} = \mathit{L19}$ \\
103 & $\mathit{Dam1} = \mathit{Delta} \cdot \mathit{am1}$ & $\mathit{Delta} \cdot \mathit{am1} = \mathit{Dam1}$ \\
104 & $\mathit{R20} = \mathit{Dam1} + 2^{32}$ & $\mathit{Dam1} + 2^{32} = \mathit{R20}$ \\
\end{longtable}

The twenty-two free equality tests are
\[\begin{array}{@{}l@{}}
\mathit{L1}=\mathit{q2},\quad \mathit{b}=\mathit{bx},\quad \mathit{L2}=\mathit{R2},\\
\mathit{L3}=\mathit{B},\quad \mathit{S2}=\mathit{pack\_rhs},\quad \mathit{n}=\mathit{q8},\\
\mathit{r}=\mathit{R7},\quad \mathit{L9}=\mathit{R9},\quad \mathit{c}=\mathit{R10a},\\
\mathit{k}=\mathit{R10b},\quad \mathit{k}=\mathit{R11},\quad \mathit{a}=\mathit{R12},\\
\mathit{c}=\mathit{R13},\quad \mathit{d}=\mathit{R14},\quad \mathit{L15}=\mathit{R15},\\
\mathit{L16}=\mathit{R16},\quad \mathit{L17}=\mathit{R17},\quad \mathit{mu}=\mathit{R18},\\
\mathit{L19}=\mathit{R19},\quad \mathit{ka}=\mathit{R20},\quad \mathit{S3}=\mathit{sigma},\\
\mathit{t2}=\mathit{Omega}
\end{array}\]


\clearpage\normalsize
\section{The 103-instruction coefficient certificate}\label{app:103}

Use the rebuilt fixed coefficient index of Section~\ref{sec:borrow-mask},
with ordinary targets $4F_i^h+\delta^2$ and special target $\delta^2$.
The code positions and $L=2^{32}$ are unchanged, but $Z,V,H$ are
reconstructed from the new coefficients. The third packed block is
\[
 S_3=B\lambda q-(Z\lambda-2e)(x+g)^2.
\]
Theorem~\ref{thm:best103} proves the complete decoding and positive
necessity, including the low carries and the elimination of the
possible quotient ambiguity. The former $q+1$ register is absent.

The checker \file{round20\_1980\_borrow\_mask\_certificate.py}
verifies all 103 primitives, all 22 source residuals and both triangular
corrections.

The table generator is \file{round20\_1980\_schedule\_table.py}.
There are 57 multiplications
and 46 additions, with 34 positive unknowns. Fixed numerals are free.

\small
\input{jones1980_theorem5_103_schedule.tex}

\clearpage\normalsize
\section{The 102-instruction certificate}\label{app:102}

Use the same coefficient index as Appendix~\ref{app:103} and retain
$U=wn^2$, $Y_P=sn^2$, $a_0=Y_P(U+1)$ and $A=a_0+4$.
Use the already calculated $U$ as the first exponent target:
\[
 d=U+c(A-4)+\gamma(8A-17).
\]
Theorem~\ref{thm:best102} proves that this change composes with the
coefficient encoding, including the correct order of the radix
deductions and the positive witness $w=4^{2r+1}/n^2$.
The separate products $bw$ and the low-offset register $q+1$ are absent.

The complete checker \file{round22\_1980\_composed\_certificate.py}
verifies all 102 primitives, all 22 source residuals and both triangular
corrections. It also verifies that the two transformations commute
on the full schedule and every source polynomial.

The table generator is \file{round22\_1980\_schedule\_table.py}.
It contains
56 multiplications and 46 additions, with 34 positive unknowns.
Fixed numerals and equality tests are free.

\small
% Generated by round22_1980_schedule_table.py; do not edit.
\begin{longtable}{@{}rll@{}}
\toprule
No. & Assignment & Addition/multiplication check\\
\midrule
\endfirsthead
\toprule
No. & Assignment & Addition/multiplication check\\
\midrule
\endhead
\bottomrule
\endfoot
1 & $\mathit{q2} = \mathit{q} \cdot \mathit{q}$ & $\mathit{q} \cdot \mathit{q} = \mathit{q2}$ \\
2 & $\mathit{q4} = \mathit{q2} \cdot \mathit{q2}$ & $\mathit{q2} \cdot \mathit{q2} = \mathit{q4}$ \\
3 & $\mathit{q8} = \mathit{q4} \cdot \mathit{q4}$ & $\mathit{q4} \cdot \mathit{q4} = \mathit{q8}$ \\
4 & $\mathit{b2} = \mathit{b} \cdot \mathit{b}$ & $\mathit{b} \cdot \mathit{b} = \mathit{b2}$ \\
5 & $\mathit{B} = \mathit{H} \cdot \mathit{b2}$ & $\mathit{H} \cdot \mathit{b2} = \mathit{B}$ \\
6 & $\mathit{bx} = \mathit{x} + \mathit{beta}$ & $\mathit{x} + \mathit{beta} = \mathit{bx}$ \\
7 & $\mathit{L2} = \mathit{la} + \mathit{q2}$ & $\mathit{la} + \mathit{q2} = \mathit{L2}$ \\
8 & $\mathit{Lam} = \mathit{la} \cdot \mathit{B}$ & $\mathit{la} \cdot \mathit{B} = \mathit{Lam}$ \\
9 & $\mathit{R2} = \mathit{Lam} + 1$ & $\mathit{Lam} + 1 = \mathit{R2}$ \\
10 & $\mathit{L3} = \mathit{th} + \mathit{Z}$ & $\mathit{th} + \mathit{Z} = \mathit{L3}$ \\
11 & $\mathit{eq} = \mathit{e} \cdot \mathit{q}$ & $\mathit{e} \cdot \mathit{q} = \mathit{eq}$ \\
12 & $\mathit{S2} = \mathit{l} + \mathit{eq}$ & $\mathit{l} + \mathit{eq} = \mathit{S2}$ \\
13 & $\mathit{pack\_tth} = \mathit{t} \cdot \mathit{th}$ & $\mathit{t} \cdot \mathit{th} = \mathit{pack\_tth}$ \\
14 & $\mathit{pack\_rhs} = \mathit{V} + \mathit{pack\_tth}$ & $\mathit{V} + \mathit{pack\_tth} = \mathit{pack\_rhs}$ \\
15 & $\mathit{zl} = \mathit{Z} \cdot \mathit{la}$ & $\mathit{Z} \cdot \mathit{la} = \mathit{zl}$ \\
16 & $\mathit{twice\_e} = 2 \cdot \mathit{e}$ & $2 \cdot \mathit{e} = \mathit{twice\_e}$ \\
17 & $\mathit{t2} = \mathit{zl} - \mathit{twice\_e}$ & $\mathit{t2} + \mathit{twice\_e} = \mathit{zl}$ \\
18 & $\mathit{C} = \mathit{x} + \mathit{g}$ & $\mathit{x} + \mathit{g} = \mathit{C}$ \\
19 & $\mathit{C2} = \mathit{C} \cdot \mathit{C}$ & $\mathit{C} \cdot \mathit{C} = \mathit{C2}$ \\
20 & $\mathit{L1} = \mathit{S2} + \mathit{al}$ & $\mathit{S2} + \mathit{al} = \mathit{L1}$ \\
21 & $\mathit{P1} = \mathit{t2} \cdot \mathit{C2}$ & $\mathit{t2} \cdot \mathit{C2} = \mathit{P1}$ \\
22 & $\mathit{P2} = \mathit{Lam} \cdot \mathit{q}$ & $\mathit{Lam} \cdot \mathit{q} = \mathit{P2}$ \\
23 & $\mathit{S3} = \mathit{P2} - \mathit{P1}$ & $\mathit{S3} + \mathit{P1} = \mathit{P2}$ \\
24 & $\mathit{S3q4} = \mathit{S3} \cdot \mathit{q2}$ & $\mathit{S3} \cdot \mathit{q2} = \mathit{S3q4}$ \\
25 & $\mathit{Sin} = \mathit{S2} + \mathit{S3q4}$ & $\mathit{S2} + \mathit{S3q4} = \mathit{Sin}$ \\
26 & $\mathit{Sq4} = \mathit{Sin} \cdot \mathit{q2}$ & $\mathit{Sin} \cdot \mathit{q2} = \mathit{Sq4}$ \\
27 & $\mathit{S} = \mathit{g} + \mathit{Sq4}$ & $\mathit{g} + \mathit{Sq4} = \mathit{S}$ \\
28 & $\mathit{n2} = \mathit{n} \cdot \mathit{n}$ & $\mathit{n} \cdot \mathit{n} = \mathit{n2}$ \\
29 & $\mathit{n2n} = \mathit{n2} - \mathit{n}$ & $\mathit{n2n} + \mathit{n} = \mathit{n2}$ \\
30 & $\mathit{SA} = \mathit{S} \cdot \mathit{n2n}$ & $\mathit{S} \cdot \mathit{n2n} = \mathit{SA}$ \\
31 & $\mathit{Tcoef} = \mathit{L2} - \mathit{zl}$ & $\mathit{Tcoef} + \mathit{zl} = \mathit{L2}$ \\
32 & $\mathit{Tq4} = \mathit{Tcoef} \cdot \mathit{q2}$ & $\mathit{Tcoef} \cdot \mathit{q2} = \mathit{Tq4}$ \\
33 & $\mathit{bm1} = \mathit{b} - 1$ & $\mathit{bm1} + 1 = \mathit{b}$ \\
34 & $\mathit{lbm1} = \mathit{l} \cdot \mathit{bm1}$ & $\mathit{l} \cdot \mathit{bm1} = \mathit{lbm1}$ \\
35 & $\mathit{T2} = \mathit{Tq4} - \mathit{lbm1}$ & $\mathit{T2} + \mathit{lbm1} = \mathit{Tq4}$ \\
36 & $\mathit{b5m2} = \mathit{B} - 4$ & $\mathit{b5m2} + 4 = \mathit{B}$ \\
37 & $\mathit{packed\_target\_shift} = \mathit{l} \cdot \mathit{q4}$ & $\mathit{l} \cdot \mathit{q4} = \mathit{packed\_target\_shift}$ \\
38 & $\mathit{mask\_term} = \mathit{b5m2} \cdot \mathit{packed\_target\_shift}$ & $\mathit{b5m2} \cdot \mathit{packed\_target\_shift} = \mathit{mask\_term}$ \\
39 & $\mathit{T} = \mathit{T2} + \mathit{mask\_term}$ & $\mathit{T2} + \mathit{mask\_term} = \mathit{T}$ \\
40 & $\mathit{n2m1} = \mathit{n2} - 1$ & $\mathit{n2m1} + 1 = \mathit{n2}$ \\
41 & $\mathit{TB} = \mathit{T} \cdot \mathit{n2m1}$ & $\mathit{T} \cdot \mathit{n2m1} = \mathit{TB}$ \\
42 & $\mathit{R7} = \mathit{SA} + \mathit{TB}$ & $\mathit{SA} + \mathit{TB} = \mathit{R7}$ \\
43 & $\mathit{wn2} = \mathit{w} \cdot \mathit{n2}$ & $\mathit{w} \cdot \mathit{n2} = \mathit{wn2}$ \\
44 & $\mathit{sn2} = \mathit{s} \cdot \mathit{n2}$ & $\mathit{s} \cdot \mathit{n2} = \mathit{sn2}$ \\
45 & $\mathit{UM} = \mathit{wn2} \cdot \mathit{sn2}$ & $\mathit{wn2} \cdot \mathit{sn2} = \mathit{UM}$ \\
46 & $\mathit{R12} = \mathit{UM} + \mathit{sn2}$ & $\mathit{UM} + \mathit{sn2} = \mathit{R12}$ \\
47 & $\mathit{ksn2} = \mathit{k} \cdot \mathit{sn2}$ & $\mathit{k} \cdot \mathit{sn2} = \mathit{ksn2}$ \\
48 & $\mathit{UM2} = \mathit{UM} \cdot \mathit{UM}$ & $\mathit{UM} \cdot \mathit{UM} = \mathit{UM2}$ \\
49 & $\mathit{scaled\_norm\_coefficient} = \mathit{UM2} + \mathit{wn2}$ & $\mathit{UM2} + \mathit{wn2} = \mathit{scaled\_norm\_coefficient}$ \\
50 & $\mathit{ratio\_product2} = \mathit{ksn2} \cdot \mathit{ksn2}$ & $\mathit{ksn2} \cdot \mathit{ksn2} = \mathit{ratio\_product2}$ \\
51 & $\mathit{L9} = \mathit{scaled\_norm\_coefficient} \cdot \mathit{ratio\_product2}$ & $\mathit{scaled\_norm\_coefficient} \cdot \mathit{ratio\_product2} = \mathit{L9}$ \\
52 & $\mathit{tauplus1} = \mathit{tau} + 1$ & $\mathit{tau} + 1 = \mathit{tauplus1}$ \\
53 & $\mathit{R9} = \mathit{tau} \cdot \mathit{tauplus1}$ & $\mathit{tau} \cdot \mathit{tauplus1} = \mathit{R9}$ \\
54 & $\mathit{R10a} = \mathit{ksn2} + \mathit{eta}$ & $\mathit{ksn2} + \mathit{eta} = \mathit{R10a}$ \\
55 & $\mathit{R10b} = \mathit{eta} + \mathit{zeta}$ & $\mathit{eta} + \mathit{zeta} = \mathit{R10b}$ \\
56 & $\mathit{r1} = \mathit{r} + 1$ & $\mathit{r} + 1 = \mathit{r1}$ \\
57 & $\mathit{hpm1} = \mathit{h} \cdot \mathit{UM}$ & $\mathit{h} \cdot \mathit{UM} = \mathit{hpm1}$ \\
58 & $\mathit{R11} = \mathit{r1} + \mathit{hpm1}$ & $\mathit{r1} + \mathit{hpm1} = \mathit{R11}$ \\
59 & $\mathit{tr1} = \mathit{r1} + \mathit{r}$ & $\mathit{r1} + \mathit{r} = \mathit{tr1}$ \\
60 & $\mathit{R13} = \mathit{ka} + \mathit{phi}$ & $\mathit{ka} + \mathit{phi} = \mathit{R13}$ \\
61 & $\mathit{cam2} = \mathit{c} \cdot \mathit{a}$ & $\mathit{c} \cdot \mathit{a} = \mathit{cam2}$ \\
62 & $\mathit{D1} = \mathit{wn2} + \mathit{cam2}$ & $\mathit{wn2} + \mathit{cam2} = \mathit{D1}$ \\
63 & $\mathit{am1} = \mathit{a} + 3$ & $\mathit{a} + 3 = \mathit{am1}$ \\
64 & $\mathit{a4} = 8 \cdot \mathit{a}$ & $8 \cdot \mathit{a} = \mathit{a4}$ \\
65 & $\mathit{a4m5} = \mathit{a4} + 15$ & $\mathit{a4} + 15 = \mathit{a4m5}$ \\
66 & $\mathit{gam} = \mathit{ga} \cdot \mathit{a4m5}$ & $\mathit{ga} \cdot \mathit{a4m5} = \mathit{gam}$ \\
67 & $\mathit{R14} = \mathit{D1} + \mathit{gam}$ & $\mathit{D1} + \mathit{gam} = \mathit{R14}$ \\
68 & $\mathit{ap1} = \mathit{am1} + 2$ & $\mathit{am1} + 2 = \mathit{ap1}$ \\
69 & $\mathit{A} = \mathit{am1} \cdot \mathit{ap1}$ & $\mathit{am1} \cdot \mathit{ap1} = \mathit{A}$ \\
70 & $\mathit{c2} = \mathit{c} \cdot \mathit{c}$ & $\mathit{c} \cdot \mathit{c} = \mathit{c2}$ \\
71 & $\mathit{Ac2} = \mathit{A} \cdot \mathit{c2}$ & $\mathit{A} \cdot \mathit{c2} = \mathit{Ac2}$ \\
72 & $\mathit{R15} = \mathit{Ac2} + 1$ & $\mathit{Ac2} + 1 = \mathit{R15}$ \\
73 & $\mathit{L15} = \mathit{d} \cdot \mathit{d}$ & $\mathit{d} \cdot \mathit{d} = \mathit{L15}$ \\
74 & $\mathit{ic2} = \mathit{i} \cdot \mathit{c2}$ & $\mathit{i} \cdot \mathit{c2} = \mathit{ic2}$ \\
75 & $\mathit{ic22} = \mathit{ic2} \cdot \mathit{ic2}$ & $\mathit{ic2} \cdot \mathit{ic2} = \mathit{ic22}$ \\
76 & $\mathit{AE} = \mathit{A} \cdot \mathit{ic22}$ & $\mathit{A} \cdot \mathit{ic22} = \mathit{AE}$ \\
77 & $\mathit{R16} = \mathit{AE} + 1$ & $\mathit{AE} + 1 = \mathit{R16}$ \\
78 & $\mathit{L16} = \mathit{f} \cdot \mathit{f}$ & $\mathit{f} \cdot \mathit{f} = \mathit{L16}$ \\
79 & $\mathit{of} = \mathit{o} \cdot \mathit{f}$ & $\mathit{o} \cdot \mathit{f} = \mathit{of}$ \\
80 & $\mathit{dof} = \mathit{of} - \mathit{d}$ & $\mathit{dof} + \mathit{d} = \mathit{of}$ \\
81 & $\mathit{L17} = \mathit{dof} \cdot \mathit{dof}$ & $\mathit{dof} \cdot \mathit{dof} = \mathit{L17}$ \\
82 & $\mathit{Gminus1} = \mathit{ap1} \cdot \mathit{AE}$ & $\mathit{ap1} \cdot \mathit{AE} = \mathit{Gminus1}$ \\
83 & $\mathit{Gplus1} = \mathit{Gminus1} + 2$ & $\mathit{Gminus1} + 2 = \mathit{Gplus1}$ \\
84 & $\mathit{G2m1} = \mathit{Gminus1} \cdot \mathit{Gplus1}$ & $\mathit{Gminus1} \cdot \mathit{Gplus1} = \mathit{G2m1}$ \\
85 & $\mathit{jc} = \mathit{j} \cdot \mathit{c}$ & $\mathit{j} \cdot \mathit{c} = \mathit{jc}$ \\
86 & $\mathit{H17} = \mathit{tr1} + \mathit{jc}$ & $\mathit{tr1} + \mathit{jc} = \mathit{H17}$ \\
87 & $\mathit{H2} = \mathit{H17} \cdot \mathit{H17}$ & $\mathit{H17} \cdot \mathit{H17} = \mathit{H2}$ \\
88 & $\mathit{P17} = \mathit{G2m1} \cdot \mathit{H2}$ & $\mathit{G2m1} \cdot \mathit{H2} = \mathit{P17}$ \\
89 & $\mathit{R17} = \mathit{P17} + 1$ & $\mathit{P17} + 1 = \mathit{R17}$ \\
90 & $\mathit{amb} = \mathit{a} - \mathit{b5m2}$ & $\mathit{amb} + \mathit{b5m2} = \mathit{a}$ \\
91 & $\mathit{kamb} = \mathit{ka} \cdot \mathit{amb}$ & $\mathit{ka} \cdot \mathit{amb} = \mathit{kamb}$ \\
92 & $\mathit{qk} = \mathit{q} + \mathit{kamb}$ & $\mathit{q} + \mathit{kamb} = \mathit{qk}$ \\
93 & $\mathit{amb2} = \mathit{amb} \cdot \mathit{amb}$ & $\mathit{amb} \cdot \mathit{amb} = \mathit{amb2}$ \\
94 & $\mathit{Mod} = \mathit{A} - \mathit{amb2}$ & $\mathit{Mod} + \mathit{amb2} = \mathit{A}$ \\
95 & $\mathit{rM} = \mathit{rho} \cdot \mathit{Mod}$ & $\mathit{rho} \cdot \mathit{Mod} = \mathit{rM}$ \\
96 & $\mathit{R18} = \mathit{qk} + \mathit{rM}$ & $\mathit{qk} + \mathit{rM} = \mathit{R18}$ \\
97 & $\mathit{ka2} = \mathit{ka} \cdot \mathit{ka}$ & $\mathit{ka} \cdot \mathit{ka} = \mathit{ka2}$ \\
98 & $\mathit{Aka2} = \mathit{A} \cdot \mathit{ka2}$ & $\mathit{A} \cdot \mathit{ka2} = \mathit{Aka2}$ \\
99 & $\mathit{R19} = \mathit{Aka2} + 1$ & $\mathit{Aka2} + 1 = \mathit{R19}$ \\
100 & $\mathit{L19} = \mathit{mu} \cdot \mathit{mu}$ & $\mathit{mu} \cdot \mathit{mu} = \mathit{L19}$ \\
101 & $\mathit{Dam1} = \mathit{Delta} \cdot \mathit{am1}$ & $\mathit{Delta} \cdot \mathit{am1} = \mathit{Dam1}$ \\
102 & $\mathit{R20} = \mathit{Dam1} + 2^{32}$ & $\mathit{Dam1} + 2^{32} = \mathit{R20}$ \\
\end{longtable}

The twenty-two free equality tests are
\[\begin{array}{@{}l@{}}
\mathit{L1}=\mathit{q2},\quad \mathit{b}=\mathit{bx},\quad \mathit{L2}=\mathit{R2},\\
\mathit{L3}=\mathit{B},\quad \mathit{S2}=\mathit{pack\_rhs},\quad \mathit{n}=\mathit{q8},\\
\mathit{r}=\mathit{R7},\quad \mathit{L9}=\mathit{R9},\quad \mathit{c}=\mathit{R10a},\\
\mathit{k}=\mathit{R10b},\quad \mathit{k}=\mathit{R11},\quad \mathit{a}=\mathit{R12},\\
\mathit{c}=\mathit{R13},\quad \mathit{d}=\mathit{R14},\quad \mathit{L15}=\mathit{R15},\\
\mathit{L16}=\mathit{R16},\quad \mathit{L17}=\mathit{R17},\quad \mathit{mu}=\mathit{R18},\\
\mathit{L19}=\mathit{R19},\quad \mathit{ka}=\mathit{R20},\quad \mathit{S3}=\mathit{sigma},\\
\mathit{t2}=\mathit{Omega}
\end{array}\]


\clearpage\normalsize
\section{The 101-instruction certificate}\label{app:101}

Use the primitive circuit index $(V,H,L)$ of Section~\ref{sec:fixed-four}.
Its digit parameter is the universal constant $Z=4$, and the power-of-two
length $L$ is chosen once for the represented set. The supplied parameters
are $x,V,H,L$; the 34 positive unknowns are unchanged. The existing
equation $\theta+4=B$ allows the register $\theta$ to supply both the
third mask and the shift $a-\theta=A-B$ in the second exponent test.
Thus the separate subtraction $B-4$ disappears.

The complete checker \file{round23\_1980\_fixed\_four\_certificate.py}
verifies every primitive, all 22 source residuals and all three triangular
corrections. The new corrections account for reusing $\theta$ in the
packed equation and the second exponential congruence. The full
positive-domain proof is Theorem~\ref{thm:best101}; the arbitrary number
of compiled circuit coordinates changes only the index, not this table.

The generator \file{round23\_1980\_schedule\_table.py} reruns that checker.
There are 56 multiplications and 45 additions. Fixed numerals and
equality tests are free.

\small
\input{jones1980_theorem5_101_schedule.tex}

\clearpage\normalsize
\section{The 100-instruction doubled-index Pell certificate}\label{app:100-pell}

Use the unpaired circuit index of Section~\ref{sec:fixed-four}, replacing
its supplied component $L$ by $T_L=\psi_4(L)$. The doubled-index norm
and the congruence $\kappa=T_L+\Delta a$ are proved in
Section~\ref{sec:doubled-index}. Their joint effect removes the last
uses of $A-1$ and $A+1$, so $A^2-1$ is computed as $a^2+(8a+15)$.
The register named \texttt{A} below holds this norm coefficient;
\texttt{dof} holds the possibly signed $v=of-d^2$;
\texttt{Gminus1} and \texttt{Gplus1} hold $K=(A^2-1)(f^2-1)$ and $K-1$.

The complete checker \file{round24\_1980\_doubled\_pell\_certificate.py}
verifies 100 primitives, 22 source residuals and three corrections.
There are 56 multiplications and 44 additions, with 34 positive unknowns
and supplied parameters $x,V,H,T_L$.

This and the following two tables are generated by
\file{round24\_26\_1980\_schedule\_tables.py}, which reruns each exact
checker. Fixed numerals and equality tests are free.

\small
\input{jones1980_theorem5_100_pell_schedule.tex}

\clearpage\normalsize
\section{The 100-instruction shared-geometric-factor certificate}\label{app:100-encoding}

Use the paired circuit index $(V,H,L)$ of Section~\ref{sec:half-center},
with its first special target and two unit constraints. The coefficient
is $\Omega=2\lambda-e$, and the third block is
$S_3=\theta\lambda q-\Omega\mathcal C^2$. The product
$\theta\lambda$ also enters the geometric equation and the middle mask.
Seven instructions replace eight, saving one multiplication.

The complete checker \file{round25\_1980\_half\_center\_certificate.py}
verifies every primitive and source residual, including the changed
geometric test. Its correction uses the exact equality $\theta+4=B$;
the latter is used first in the proof order, regardless of its position
in the displayed list. There are 55 multiplications and 45 additions,
with 34 positive unknowns and 22 equality tests.

\small
% Generated by round24_26_1980_schedule_tables.py; do not edit.
\begin{longtable}{@{}rll@{}}
\toprule
No. & Assignment & Addition/multiplication check\\
\midrule
\endfirsthead
\toprule
No. & Assignment & Addition/multiplication check\\
\midrule
\endhead
\bottomrule
\endfoot
1 & $\mathit{q2} = \mathit{q} \cdot \mathit{q}$ & $\mathit{q} \cdot \mathit{q} = \mathit{q2}$ \\
2 & $\mathit{q4} = \mathit{q2} \cdot \mathit{q2}$ & $\mathit{q2} \cdot \mathit{q2} = \mathit{q4}$ \\
3 & $\mathit{q8} = \mathit{q4} \cdot \mathit{q4}$ & $\mathit{q4} \cdot \mathit{q4} = \mathit{q8}$ \\
4 & $\mathit{b2} = \mathit{b} \cdot \mathit{b}$ & $\mathit{b} \cdot \mathit{b} = \mathit{b2}$ \\
5 & $\mathit{B} = \mathit{H} \cdot \mathit{b2}$ & $\mathit{H} \cdot \mathit{b2} = \mathit{B}$ \\
6 & $\mathit{bx} = \mathit{x} + \mathit{beta}$ & $\mathit{x} + \mathit{beta} = \mathit{bx}$ \\
7 & $\mathit{theta\_lambda} = \mathit{th} \cdot \mathit{la}$ & $\mathit{th} \cdot \mathit{la} = \mathit{theta\_lambda}$ \\
8 & $\mathit{twice\_lambda} = 2 \cdot \mathit{la}$ & $2 \cdot \mathit{la} = \mathit{twice\_lambda}$ \\
9 & $\mathit{Tcoef} = \mathit{theta\_lambda} + 1$ & $\mathit{theta\_lambda} + 1 = \mathit{Tcoef}$ \\
10 & $\mathit{three\_lambda} = \mathit{twice\_lambda} + \mathit{la}$ & $\mathit{twice\_lambda} + \mathit{la} = \mathit{three\_lambda}$ \\
11 & $\mathit{R2} = \mathit{Tcoef} + \mathit{three\_lambda}$ & $\mathit{Tcoef} + \mathit{three\_lambda} = \mathit{R2}$ \\
12 & $\mathit{t2} = \mathit{twice\_lambda} - \mathit{e}$ & $\mathit{t2} + \mathit{e} = \mathit{twice\_lambda}$ \\
13 & $\mathit{P2} = \mathit{theta\_lambda} \cdot \mathit{q}$ & $\mathit{theta\_lambda} \cdot \mathit{q} = \mathit{P2}$ \\
14 & $\mathit{L3} = \mathit{th} + 4$ & $\mathit{th} + 4 = \mathit{L3}$ \\
15 & $\mathit{eq} = \mathit{e} \cdot \mathit{q}$ & $\mathit{e} \cdot \mathit{q} = \mathit{eq}$ \\
16 & $\mathit{S2} = \mathit{l} + \mathit{eq}$ & $\mathit{l} + \mathit{eq} = \mathit{S2}$ \\
17 & $\mathit{pack\_tth} = \mathit{t} \cdot \mathit{th}$ & $\mathit{t} \cdot \mathit{th} = \mathit{pack\_tth}$ \\
18 & $\mathit{pack\_rhs} = \mathit{V} + \mathit{pack\_tth}$ & $\mathit{V} + \mathit{pack\_tth} = \mathit{pack\_rhs}$ \\
19 & $\mathit{C} = \mathit{x} + \mathit{g}$ & $\mathit{x} + \mathit{g} = \mathit{C}$ \\
20 & $\mathit{C2} = \mathit{C} \cdot \mathit{C}$ & $\mathit{C} \cdot \mathit{C} = \mathit{C2}$ \\
21 & $\mathit{L1} = \mathit{S2} + \mathit{al}$ & $\mathit{S2} + \mathit{al} = \mathit{L1}$ \\
22 & $\mathit{P1} = \mathit{t2} \cdot \mathit{C2}$ & $\mathit{t2} \cdot \mathit{C2} = \mathit{P1}$ \\
23 & $\mathit{S3} = \mathit{P2} - \mathit{P1}$ & $\mathit{S3} + \mathit{P1} = \mathit{P2}$ \\
24 & $\mathit{S3q4} = \mathit{S3} \cdot \mathit{q2}$ & $\mathit{S3} \cdot \mathit{q2} = \mathit{S3q4}$ \\
25 & $\mathit{Sin} = \mathit{S2} + \mathit{S3q4}$ & $\mathit{S2} + \mathit{S3q4} = \mathit{Sin}$ \\
26 & $\mathit{Sq4} = \mathit{Sin} \cdot \mathit{q2}$ & $\mathit{Sin} \cdot \mathit{q2} = \mathit{Sq4}$ \\
27 & $\mathit{S} = \mathit{g} + \mathit{Sq4}$ & $\mathit{g} + \mathit{Sq4} = \mathit{S}$ \\
28 & $\mathit{n2} = \mathit{n} \cdot \mathit{n}$ & $\mathit{n} \cdot \mathit{n} = \mathit{n2}$ \\
29 & $\mathit{n2n} = \mathit{n2} - \mathit{n}$ & $\mathit{n2n} + \mathit{n} = \mathit{n2}$ \\
30 & $\mathit{SA} = \mathit{S} \cdot \mathit{n2n}$ & $\mathit{S} \cdot \mathit{n2n} = \mathit{SA}$ \\
31 & $\mathit{Tq4} = \mathit{Tcoef} \cdot \mathit{q2}$ & $\mathit{Tcoef} \cdot \mathit{q2} = \mathit{Tq4}$ \\
32 & $\mathit{bm1} = \mathit{b} - 1$ & $\mathit{bm1} + 1 = \mathit{b}$ \\
33 & $\mathit{lbm1} = \mathit{l} \cdot \mathit{bm1}$ & $\mathit{l} \cdot \mathit{bm1} = \mathit{lbm1}$ \\
34 & $\mathit{T2} = \mathit{Tq4} - \mathit{lbm1}$ & $\mathit{T2} + \mathit{lbm1} = \mathit{Tq4}$ \\
35 & $\mathit{packed\_target\_shift} = \mathit{l} \cdot \mathit{q4}$ & $\mathit{l} \cdot \mathit{q4} = \mathit{packed\_target\_shift}$ \\
36 & $\mathit{mask\_term} = \mathit{th} \cdot \mathit{packed\_target\_shift}$ & $\mathit{th} \cdot \mathit{packed\_target\_shift} = \mathit{mask\_term}$ \\
37 & $\mathit{T} = \mathit{T2} + \mathit{mask\_term}$ & $\mathit{T2} + \mathit{mask\_term} = \mathit{T}$ \\
38 & $\mathit{n2m1} = \mathit{n2} - 1$ & $\mathit{n2m1} + 1 = \mathit{n2}$ \\
39 & $\mathit{TB} = \mathit{T} \cdot \mathit{n2m1}$ & $\mathit{T} \cdot \mathit{n2m1} = \mathit{TB}$ \\
40 & $\mathit{R7} = \mathit{SA} + \mathit{TB}$ & $\mathit{SA} + \mathit{TB} = \mathit{R7}$ \\
41 & $\mathit{wn2} = \mathit{w} \cdot \mathit{n2}$ & $\mathit{w} \cdot \mathit{n2} = \mathit{wn2}$ \\
42 & $\mathit{sn2} = \mathit{s} \cdot \mathit{n2}$ & $\mathit{s} \cdot \mathit{n2} = \mathit{sn2}$ \\
43 & $\mathit{UM} = \mathit{wn2} \cdot \mathit{sn2}$ & $\mathit{wn2} \cdot \mathit{sn2} = \mathit{UM}$ \\
44 & $\mathit{R12} = \mathit{UM} + \mathit{sn2}$ & $\mathit{UM} + \mathit{sn2} = \mathit{R12}$ \\
45 & $\mathit{ksn2} = \mathit{k} \cdot \mathit{sn2}$ & $\mathit{k} \cdot \mathit{sn2} = \mathit{ksn2}$ \\
46 & $\mathit{UM2} = \mathit{UM} \cdot \mathit{UM}$ & $\mathit{UM} \cdot \mathit{UM} = \mathit{UM2}$ \\
47 & $\mathit{scaled\_norm\_coefficient} = \mathit{UM2} + \mathit{wn2}$ & $\mathit{UM2} + \mathit{wn2} = \mathit{scaled\_norm\_coefficient}$ \\
48 & $\mathit{ratio\_product2} = \mathit{ksn2} \cdot \mathit{ksn2}$ & $\mathit{ksn2} \cdot \mathit{ksn2} = \mathit{ratio\_product2}$ \\
49 & $\mathit{L9} = \mathit{scaled\_norm\_coefficient} \cdot \mathit{ratio\_product2}$ & $\mathit{scaled\_norm\_coefficient} \cdot \mathit{ratio\_product2} = \mathit{L9}$ \\
50 & $\mathit{tauplus1} = \mathit{tau} + 1$ & $\mathit{tau} + 1 = \mathit{tauplus1}$ \\
51 & $\mathit{R9} = \mathit{tau} \cdot \mathit{tauplus1}$ & $\mathit{tau} \cdot \mathit{tauplus1} = \mathit{R9}$ \\
52 & $\mathit{R10a} = \mathit{ksn2} + \mathit{eta}$ & $\mathit{ksn2} + \mathit{eta} = \mathit{R10a}$ \\
53 & $\mathit{R10b} = \mathit{eta} + \mathit{zeta}$ & $\mathit{eta} + \mathit{zeta} = \mathit{R10b}$ \\
54 & $\mathit{r1} = \mathit{r} + 1$ & $\mathit{r} + 1 = \mathit{r1}$ \\
55 & $\mathit{hpm1} = \mathit{h} \cdot \mathit{UM}$ & $\mathit{h} \cdot \mathit{UM} = \mathit{hpm1}$ \\
56 & $\mathit{R11} = \mathit{r1} + \mathit{hpm1}$ & $\mathit{r1} + \mathit{hpm1} = \mathit{R11}$ \\
57 & $\mathit{tr1} = \mathit{r1} + \mathit{r}$ & $\mathit{r1} + \mathit{r} = \mathit{tr1}$ \\
58 & $\mathit{R13} = \mathit{ka} + \mathit{phi}$ & $\mathit{ka} + \mathit{phi} = \mathit{R13}$ \\
59 & $\mathit{cam2} = \mathit{c} \cdot \mathit{a}$ & $\mathit{c} \cdot \mathit{a} = \mathit{cam2}$ \\
60 & $\mathit{D1} = \mathit{wn2} + \mathit{cam2}$ & $\mathit{wn2} + \mathit{cam2} = \mathit{D1}$ \\
61 & $\mathit{am1} = \mathit{a} + 3$ & $\mathit{a} + 3 = \mathit{am1}$ \\
62 & $\mathit{a4} = 8 \cdot \mathit{a}$ & $8 \cdot \mathit{a} = \mathit{a4}$ \\
63 & $\mathit{a4m5} = \mathit{a4} + 15$ & $\mathit{a4} + 15 = \mathit{a4m5}$ \\
64 & $\mathit{gam} = \mathit{ga} \cdot \mathit{a4m5}$ & $\mathit{ga} \cdot \mathit{a4m5} = \mathit{gam}$ \\
65 & $\mathit{R14} = \mathit{D1} + \mathit{gam}$ & $\mathit{D1} + \mathit{gam} = \mathit{R14}$ \\
66 & $\mathit{ap1} = \mathit{am1} + 2$ & $\mathit{am1} + 2 = \mathit{ap1}$ \\
67 & $\mathit{A} = \mathit{am1} \cdot \mathit{ap1}$ & $\mathit{am1} \cdot \mathit{ap1} = \mathit{A}$ \\
68 & $\mathit{c2} = \mathit{c} \cdot \mathit{c}$ & $\mathit{c} \cdot \mathit{c} = \mathit{c2}$ \\
69 & $\mathit{Ac2} = \mathit{A} \cdot \mathit{c2}$ & $\mathit{A} \cdot \mathit{c2} = \mathit{Ac2}$ \\
70 & $\mathit{R15} = \mathit{Ac2} + 1$ & $\mathit{Ac2} + 1 = \mathit{R15}$ \\
71 & $\mathit{L15} = \mathit{d} \cdot \mathit{d}$ & $\mathit{d} \cdot \mathit{d} = \mathit{L15}$ \\
72 & $\mathit{ic2} = \mathit{i} \cdot \mathit{c2}$ & $\mathit{i} \cdot \mathit{c2} = \mathit{ic2}$ \\
73 & $\mathit{ic22} = \mathit{ic2} \cdot \mathit{ic2}$ & $\mathit{ic2} \cdot \mathit{ic2} = \mathit{ic22}$ \\
74 & $\mathit{AE} = \mathit{A} \cdot \mathit{ic22}$ & $\mathit{A} \cdot \mathit{ic22} = \mathit{AE}$ \\
75 & $\mathit{R16} = \mathit{AE} + 1$ & $\mathit{AE} + 1 = \mathit{R16}$ \\
76 & $\mathit{L16} = \mathit{f} \cdot \mathit{f}$ & $\mathit{f} \cdot \mathit{f} = \mathit{L16}$ \\
77 & $\mathit{of} = \mathit{o} \cdot \mathit{f}$ & $\mathit{o} \cdot \mathit{f} = \mathit{of}$ \\
78 & $\mathit{dof} = \mathit{of} - \mathit{d}$ & $\mathit{dof} + \mathit{d} = \mathit{of}$ \\
79 & $\mathit{L17} = \mathit{dof} \cdot \mathit{dof}$ & $\mathit{dof} \cdot \mathit{dof} = \mathit{L17}$ \\
80 & $\mathit{Gminus1} = \mathit{ap1} \cdot \mathit{AE}$ & $\mathit{ap1} \cdot \mathit{AE} = \mathit{Gminus1}$ \\
81 & $\mathit{Gplus1} = \mathit{Gminus1} + 2$ & $\mathit{Gminus1} + 2 = \mathit{Gplus1}$ \\
82 & $\mathit{G2m1} = \mathit{Gminus1} \cdot \mathit{Gplus1}$ & $\mathit{Gminus1} \cdot \mathit{Gplus1} = \mathit{G2m1}$ \\
83 & $\mathit{jc} = \mathit{j} \cdot \mathit{c}$ & $\mathit{j} \cdot \mathit{c} = \mathit{jc}$ \\
84 & $\mathit{H17} = \mathit{tr1} + \mathit{jc}$ & $\mathit{tr1} + \mathit{jc} = \mathit{H17}$ \\
85 & $\mathit{H2} = \mathit{H17} \cdot \mathit{H17}$ & $\mathit{H17} \cdot \mathit{H17} = \mathit{H2}$ \\
86 & $\mathit{P17} = \mathit{G2m1} \cdot \mathit{H2}$ & $\mathit{G2m1} \cdot \mathit{H2} = \mathit{P17}$ \\
87 & $\mathit{R17} = \mathit{P17} + 1$ & $\mathit{P17} + 1 = \mathit{R17}$ \\
88 & $\mathit{amb} = \mathit{a} - \mathit{th}$ & $\mathit{amb} + \mathit{th} = \mathit{a}$ \\
89 & $\mathit{kamb} = \mathit{ka} \cdot \mathit{amb}$ & $\mathit{ka} \cdot \mathit{amb} = \mathit{kamb}$ \\
90 & $\mathit{qk} = \mathit{q} + \mathit{kamb}$ & $\mathit{q} + \mathit{kamb} = \mathit{qk}$ \\
91 & $\mathit{amb2} = \mathit{amb} \cdot \mathit{amb}$ & $\mathit{amb} \cdot \mathit{amb} = \mathit{amb2}$ \\
92 & $\mathit{Mod} = \mathit{A} - \mathit{amb2}$ & $\mathit{Mod} + \mathit{amb2} = \mathit{A}$ \\
93 & $\mathit{rM} = \mathit{rho} \cdot \mathit{Mod}$ & $\mathit{rho} \cdot \mathit{Mod} = \mathit{rM}$ \\
94 & $\mathit{R18} = \mathit{qk} + \mathit{rM}$ & $\mathit{qk} + \mathit{rM} = \mathit{R18}$ \\
95 & $\mathit{ka2} = \mathit{ka} \cdot \mathit{ka}$ & $\mathit{ka} \cdot \mathit{ka} = \mathit{ka2}$ \\
96 & $\mathit{Aka2} = \mathit{A} \cdot \mathit{ka2}$ & $\mathit{A} \cdot \mathit{ka2} = \mathit{Aka2}$ \\
97 & $\mathit{R19} = \mathit{Aka2} + 1$ & $\mathit{Aka2} + 1 = \mathit{R19}$ \\
98 & $\mathit{L19} = \mathit{mu} \cdot \mathit{mu}$ & $\mathit{mu} \cdot \mathit{mu} = \mathit{L19}$ \\
99 & $\mathit{Dam1} = \mathit{Delta} \cdot \mathit{am1}$ & $\mathit{Delta} \cdot \mathit{am1} = \mathit{Dam1}$ \\
100 & $\mathit{R20} = \mathit{Dam1} + \mathit{L}$ & $\mathit{Dam1} + \mathit{L} = \mathit{R20}$ \\
\end{longtable}

The twenty-two free equality tests are
\[\begin{array}{@{}l@{}}
\mathit{L1}=\mathit{q2},\quad \mathit{b}=\mathit{bx},\quad \mathit{q2}=\mathit{R2},\\
\mathit{L3}=\mathit{B},\quad \mathit{S2}=\mathit{pack\_rhs},\quad \mathit{n}=\mathit{q8},\\
\mathit{r}=\mathit{R7},\quad \mathit{L9}=\mathit{R9},\quad \mathit{c}=\mathit{R10a},\\
\mathit{k}=\mathit{R10b},\quad \mathit{k}=\mathit{R11},\quad \mathit{a}=\mathit{R12},\\
\mathit{c}=\mathit{R13},\quad \mathit{d}=\mathit{R14},\quad \mathit{L15}=\mathit{R15},\\
\mathit{L16}=\mathit{R16},\quad \mathit{L17}=\mathit{R17},\quad \mathit{mu}=\mathit{R18},\\
\mathit{L19}=\mathit{R19},\quad \mathit{ka}=\mathit{R20},\quad \mathit{S3}=\mathit{sigma},\\
\mathit{t2}=\mathit{Omega}
\end{array}\]


\clearpage\normalsize
\section{The 99-instruction composed certificate}\label{app:99}

Use the paired, special-first coefficient construction and supply its
underlying length through $T_L=\psi_4(L)$. The supplied parameters are
$x,V,H,T_L$ and the 34 positive unknowns are unchanged. Combine both
changes of the preceding appendices; Theorem~\ref{thm:best99} supplies
the complete positive-domain composition proof.

The standalone checker \file{round26\_1980\_composed\_certificate.py}
verifies all 99 instructions and 22 source residuals, their three
acyclic corrections, and both orders of the full arithmetic and source
transformations. Only the packed equation, auxiliary signed norm,
fixed index equation, positive $S_3$ equation and positive $\Omega$
equation change after renaming the fixed index. The resulting count is
55 multiplications and 44 additions. Fixed numerals and equality tests
are free.

\small
% Generated by round24_26_1980_schedule_tables.py; do not edit.
\begin{longtable}{@{}rll@{}}
\toprule
No. & Assignment & Addition/multiplication check\\
\midrule
\endfirsthead
\toprule
No. & Assignment & Addition/multiplication check\\
\midrule
\endhead
\bottomrule
\endfoot
1 & $\mathit{q2} = \mathit{q} \cdot \mathit{q}$ & $\mathit{q} \cdot \mathit{q} = \mathit{q2}$ \\
2 & $\mathit{q4} = \mathit{q2} \cdot \mathit{q2}$ & $\mathit{q2} \cdot \mathit{q2} = \mathit{q4}$ \\
3 & $\mathit{q8} = \mathit{q4} \cdot \mathit{q4}$ & $\mathit{q4} \cdot \mathit{q4} = \mathit{q8}$ \\
4 & $\mathit{b2} = \mathit{b} \cdot \mathit{b}$ & $\mathit{b} \cdot \mathit{b} = \mathit{b2}$ \\
5 & $\mathit{B} = \mathit{H} \cdot \mathit{b2}$ & $\mathit{H} \cdot \mathit{b2} = \mathit{B}$ \\
6 & $\mathit{bx} = \mathit{x} + \mathit{beta}$ & $\mathit{x} + \mathit{beta} = \mathit{bx}$ \\
7 & $\mathit{theta\_lambda} = \mathit{th} \cdot \mathit{la}$ & $\mathit{th} \cdot \mathit{la} = \mathit{theta\_lambda}$ \\
8 & $\mathit{twice\_lambda} = 2 \cdot \mathit{la}$ & $2 \cdot \mathit{la} = \mathit{twice\_lambda}$ \\
9 & $\mathit{Tcoef} = \mathit{theta\_lambda} + 1$ & $\mathit{theta\_lambda} + 1 = \mathit{Tcoef}$ \\
10 & $\mathit{three\_lambda} = \mathit{twice\_lambda} + \mathit{la}$ & $\mathit{twice\_lambda} + \mathit{la} = \mathit{three\_lambda}$ \\
11 & $\mathit{R2} = \mathit{Tcoef} + \mathit{three\_lambda}$ & $\mathit{Tcoef} + \mathit{three\_lambda} = \mathit{R2}$ \\
12 & $\mathit{t2} = \mathit{twice\_lambda} - \mathit{e}$ & $\mathit{t2} + \mathit{e} = \mathit{twice\_lambda}$ \\
13 & $\mathit{P2} = \mathit{theta\_lambda} \cdot \mathit{q}$ & $\mathit{theta\_lambda} \cdot \mathit{q} = \mathit{P2}$ \\
14 & $\mathit{L3} = \mathit{th} + 4$ & $\mathit{th} + 4 = \mathit{L3}$ \\
15 & $\mathit{eq} = \mathit{e} \cdot \mathit{q}$ & $\mathit{e} \cdot \mathit{q} = \mathit{eq}$ \\
16 & $\mathit{S2} = \mathit{l} + \mathit{eq}$ & $\mathit{l} + \mathit{eq} = \mathit{S2}$ \\
17 & $\mathit{pack\_tth} = \mathit{t} \cdot \mathit{th}$ & $\mathit{t} \cdot \mathit{th} = \mathit{pack\_tth}$ \\
18 & $\mathit{pack\_rhs} = \mathit{V} + \mathit{pack\_tth}$ & $\mathit{V} + \mathit{pack\_tth} = \mathit{pack\_rhs}$ \\
19 & $\mathit{C} = \mathit{x} + \mathit{g}$ & $\mathit{x} + \mathit{g} = \mathit{C}$ \\
20 & $\mathit{C2} = \mathit{C} \cdot \mathit{C}$ & $\mathit{C} \cdot \mathit{C} = \mathit{C2}$ \\
21 & $\mathit{L1} = \mathit{S2} + \mathit{al}$ & $\mathit{S2} + \mathit{al} = \mathit{L1}$ \\
22 & $\mathit{P1} = \mathit{t2} \cdot \mathit{C2}$ & $\mathit{t2} \cdot \mathit{C2} = \mathit{P1}$ \\
23 & $\mathit{S3} = \mathit{P2} - \mathit{P1}$ & $\mathit{S3} + \mathit{P1} = \mathit{P2}$ \\
24 & $\mathit{S3q4} = \mathit{S3} \cdot \mathit{q2}$ & $\mathit{S3} \cdot \mathit{q2} = \mathit{S3q4}$ \\
25 & $\mathit{Sin} = \mathit{S2} + \mathit{S3q4}$ & $\mathit{S2} + \mathit{S3q4} = \mathit{Sin}$ \\
26 & $\mathit{Sq4} = \mathit{Sin} \cdot \mathit{q2}$ & $\mathit{Sin} \cdot \mathit{q2} = \mathit{Sq4}$ \\
27 & $\mathit{S} = \mathit{g} + \mathit{Sq4}$ & $\mathit{g} + \mathit{Sq4} = \mathit{S}$ \\
28 & $\mathit{n2} = \mathit{n} \cdot \mathit{n}$ & $\mathit{n} \cdot \mathit{n} = \mathit{n2}$ \\
29 & $\mathit{n2n} = \mathit{n2} - \mathit{n}$ & $\mathit{n2n} + \mathit{n} = \mathit{n2}$ \\
30 & $\mathit{SA} = \mathit{S} \cdot \mathit{n2n}$ & $\mathit{S} \cdot \mathit{n2n} = \mathit{SA}$ \\
31 & $\mathit{Tq4} = \mathit{Tcoef} \cdot \mathit{q2}$ & $\mathit{Tcoef} \cdot \mathit{q2} = \mathit{Tq4}$ \\
32 & $\mathit{bm1} = \mathit{b} - 1$ & $\mathit{bm1} + 1 = \mathit{b}$ \\
33 & $\mathit{lbm1} = \mathit{l} \cdot \mathit{bm1}$ & $\mathit{l} \cdot \mathit{bm1} = \mathit{lbm1}$ \\
34 & $\mathit{T2} = \mathit{Tq4} - \mathit{lbm1}$ & $\mathit{T2} + \mathit{lbm1} = \mathit{Tq4}$ \\
35 & $\mathit{packed\_target\_shift} = \mathit{l} \cdot \mathit{q4}$ & $\mathit{l} \cdot \mathit{q4} = \mathit{packed\_target\_shift}$ \\
36 & $\mathit{mask\_term} = \mathit{th} \cdot \mathit{packed\_target\_shift}$ & $\mathit{th} \cdot \mathit{packed\_target\_shift} = \mathit{mask\_term}$ \\
37 & $\mathit{T} = \mathit{T2} + \mathit{mask\_term}$ & $\mathit{T2} + \mathit{mask\_term} = \mathit{T}$ \\
38 & $\mathit{n2m1} = \mathit{n2} - 1$ & $\mathit{n2m1} + 1 = \mathit{n2}$ \\
39 & $\mathit{TB} = \mathit{T} \cdot \mathit{n2m1}$ & $\mathit{T} \cdot \mathit{n2m1} = \mathit{TB}$ \\
40 & $\mathit{R7} = \mathit{SA} + \mathit{TB}$ & $\mathit{SA} + \mathit{TB} = \mathit{R7}$ \\
41 & $\mathit{wn2} = \mathit{w} \cdot \mathit{n2}$ & $\mathit{w} \cdot \mathit{n2} = \mathit{wn2}$ \\
42 & $\mathit{sn2} = \mathit{s} \cdot \mathit{n2}$ & $\mathit{s} \cdot \mathit{n2} = \mathit{sn2}$ \\
43 & $\mathit{UM} = \mathit{wn2} \cdot \mathit{sn2}$ & $\mathit{wn2} \cdot \mathit{sn2} = \mathit{UM}$ \\
44 & $\mathit{R12} = \mathit{UM} + \mathit{sn2}$ & $\mathit{UM} + \mathit{sn2} = \mathit{R12}$ \\
45 & $\mathit{ksn2} = \mathit{k} \cdot \mathit{sn2}$ & $\mathit{k} \cdot \mathit{sn2} = \mathit{ksn2}$ \\
46 & $\mathit{UM2} = \mathit{UM} \cdot \mathit{UM}$ & $\mathit{UM} \cdot \mathit{UM} = \mathit{UM2}$ \\
47 & $\mathit{scaled\_norm\_coefficient} = \mathit{UM2} + \mathit{wn2}$ & $\mathit{UM2} + \mathit{wn2} = \mathit{scaled\_norm\_coefficient}$ \\
48 & $\mathit{ratio\_product2} = \mathit{ksn2} \cdot \mathit{ksn2}$ & $\mathit{ksn2} \cdot \mathit{ksn2} = \mathit{ratio\_product2}$ \\
49 & $\mathit{L9} = \mathit{scaled\_norm\_coefficient} \cdot \mathit{ratio\_product2}$ & $\mathit{scaled\_norm\_coefficient} \cdot \mathit{ratio\_product2} = \mathit{L9}$ \\
50 & $\mathit{tauplus1} = \mathit{tau} + 1$ & $\mathit{tau} + 1 = \mathit{tauplus1}$ \\
51 & $\mathit{R9} = \mathit{tau} \cdot \mathit{tauplus1}$ & $\mathit{tau} \cdot \mathit{tauplus1} = \mathit{R9}$ \\
52 & $\mathit{R10a} = \mathit{ksn2} + \mathit{eta}$ & $\mathit{ksn2} + \mathit{eta} = \mathit{R10a}$ \\
53 & $\mathit{R10b} = \mathit{eta} + \mathit{zeta}$ & $\mathit{eta} + \mathit{zeta} = \mathit{R10b}$ \\
54 & $\mathit{r1} = \mathit{r} + 1$ & $\mathit{r} + 1 = \mathit{r1}$ \\
55 & $\mathit{hpm1} = \mathit{h} \cdot \mathit{UM}$ & $\mathit{h} \cdot \mathit{UM} = \mathit{hpm1}$ \\
56 & $\mathit{R11} = \mathit{r1} + \mathit{hpm1}$ & $\mathit{r1} + \mathit{hpm1} = \mathit{R11}$ \\
57 & $\mathit{tr1} = \mathit{r1} + \mathit{r}$ & $\mathit{r1} + \mathit{r} = \mathit{tr1}$ \\
58 & $\mathit{R13} = \mathit{ka} + \mathit{phi}$ & $\mathit{ka} + \mathit{phi} = \mathit{R13}$ \\
59 & $\mathit{cam2} = \mathit{c} \cdot \mathit{a}$ & $\mathit{c} \cdot \mathit{a} = \mathit{cam2}$ \\
60 & $\mathit{D1} = \mathit{wn2} + \mathit{cam2}$ & $\mathit{wn2} + \mathit{cam2} = \mathit{D1}$ \\
61 & $\mathit{a4} = 8 \cdot \mathit{a}$ & $8 \cdot \mathit{a} = \mathit{a4}$ \\
62 & $\mathit{a4m5} = \mathit{a4} + 15$ & $\mathit{a4} + 15 = \mathit{a4m5}$ \\
63 & $\mathit{gam} = \mathit{ga} \cdot \mathit{a4m5}$ & $\mathit{ga} \cdot \mathit{a4m5} = \mathit{gam}$ \\
64 & $\mathit{R14} = \mathit{D1} + \mathit{gam}$ & $\mathit{D1} + \mathit{gam} = \mathit{R14}$ \\
65 & $\mathit{a\_square} = \mathit{a} \cdot \mathit{a}$ & $\mathit{a} \cdot \mathit{a} = \mathit{a\_square}$ \\
66 & $\mathit{A} = \mathit{a\_square} + \mathit{a4m5}$ & $\mathit{a\_square} + \mathit{a4m5} = \mathit{A}$ \\
67 & $\mathit{c2} = \mathit{c} \cdot \mathit{c}$ & $\mathit{c} \cdot \mathit{c} = \mathit{c2}$ \\
68 & $\mathit{Ac2} = \mathit{A} \cdot \mathit{c2}$ & $\mathit{A} \cdot \mathit{c2} = \mathit{Ac2}$ \\
69 & $\mathit{R15} = \mathit{Ac2} + 1$ & $\mathit{Ac2} + 1 = \mathit{R15}$ \\
70 & $\mathit{L15} = \mathit{d} \cdot \mathit{d}$ & $\mathit{d} \cdot \mathit{d} = \mathit{L15}$ \\
71 & $\mathit{ic2} = \mathit{i} \cdot \mathit{c2}$ & $\mathit{i} \cdot \mathit{c2} = \mathit{ic2}$ \\
72 & $\mathit{ic22} = \mathit{ic2} \cdot \mathit{ic2}$ & $\mathit{ic2} \cdot \mathit{ic2} = \mathit{ic22}$ \\
73 & $\mathit{AE} = \mathit{A} \cdot \mathit{ic22}$ & $\mathit{A} \cdot \mathit{ic22} = \mathit{AE}$ \\
74 & $\mathit{R16} = \mathit{AE} + 1$ & $\mathit{AE} + 1 = \mathit{R16}$ \\
75 & $\mathit{L16} = \mathit{f} \cdot \mathit{f}$ & $\mathit{f} \cdot \mathit{f} = \mathit{L16}$ \\
76 & $\mathit{of} = \mathit{o} \cdot \mathit{f}$ & $\mathit{o} \cdot \mathit{f} = \mathit{of}$ \\
77 & $\mathit{dof} = \mathit{of} - \mathit{L15}$ & $\mathit{dof} + \mathit{L15} = \mathit{of}$ \\
78 & $\mathit{vplus1} = \mathit{dof} + 1$ & $\mathit{dof} + 1 = \mathit{vplus1}$ \\
79 & $\mathit{L17} = \mathit{dof} \cdot \mathit{vplus1}$ & $\mathit{dof} \cdot \mathit{vplus1} = \mathit{L17}$ \\
80 & $\mathit{Gminus1} = \mathit{A} \cdot \mathit{AE}$ & $\mathit{A} \cdot \mathit{AE} = \mathit{Gminus1}$ \\
81 & $\mathit{Gplus1} = \mathit{Gminus1} - 1$ & $\mathit{Gplus1} + 1 = \mathit{Gminus1}$ \\
82 & $\mathit{G2m1} = \mathit{Gminus1} \cdot \mathit{Gplus1}$ & $\mathit{Gminus1} \cdot \mathit{Gplus1} = \mathit{G2m1}$ \\
83 & $\mathit{jc} = \mathit{j} \cdot \mathit{c}$ & $\mathit{j} \cdot \mathit{c} = \mathit{jc}$ \\
84 & $\mathit{H17} = \mathit{tr1} + \mathit{jc}$ & $\mathit{tr1} + \mathit{jc} = \mathit{H17}$ \\
85 & $\mathit{H2} = \mathit{H17} \cdot \mathit{H17}$ & $\mathit{H17} \cdot \mathit{H17} = \mathit{H2}$ \\
86 & $\mathit{P17} = \mathit{G2m1} \cdot \mathit{H2}$ & $\mathit{G2m1} \cdot \mathit{H2} = \mathit{P17}$ \\
87 & $\mathit{amb} = \mathit{a} - \mathit{th}$ & $\mathit{amb} + \mathit{th} = \mathit{a}$ \\
88 & $\mathit{kamb} = \mathit{ka} \cdot \mathit{amb}$ & $\mathit{ka} \cdot \mathit{amb} = \mathit{kamb}$ \\
89 & $\mathit{qk} = \mathit{q} + \mathit{kamb}$ & $\mathit{q} + \mathit{kamb} = \mathit{qk}$ \\
90 & $\mathit{amb2} = \mathit{amb} \cdot \mathit{amb}$ & $\mathit{amb} \cdot \mathit{amb} = \mathit{amb2}$ \\
91 & $\mathit{Mod} = \mathit{A} - \mathit{amb2}$ & $\mathit{Mod} + \mathit{amb2} = \mathit{A}$ \\
92 & $\mathit{rM} = \mathit{rho} \cdot \mathit{Mod}$ & $\mathit{rho} \cdot \mathit{Mod} = \mathit{rM}$ \\
93 & $\mathit{R18} = \mathit{qk} + \mathit{rM}$ & $\mathit{qk} + \mathit{rM} = \mathit{R18}$ \\
94 & $\mathit{ka2} = \mathit{ka} \cdot \mathit{ka}$ & $\mathit{ka} \cdot \mathit{ka} = \mathit{ka2}$ \\
95 & $\mathit{Aka2} = \mathit{A} \cdot \mathit{ka2}$ & $\mathit{A} \cdot \mathit{ka2} = \mathit{Aka2}$ \\
96 & $\mathit{R19} = \mathit{Aka2} + 1$ & $\mathit{Aka2} + 1 = \mathit{R19}$ \\
97 & $\mathit{L19} = \mathit{mu} \cdot \mathit{mu}$ & $\mathit{mu} \cdot \mathit{mu} = \mathit{L19}$ \\
98 & $\mathit{Dam1} = \mathit{Delta} \cdot \mathit{a}$ & $\mathit{Delta} \cdot \mathit{a} = \mathit{Dam1}$ \\
99 & $\mathit{R20} = \mathit{Dam1} + \mathit{Tindex}$ & $\mathit{Dam1} + \mathit{Tindex} = \mathit{R20}$ \\
\end{longtable}

The twenty-two free equality tests are
\[\begin{array}{@{}l@{}}
\mathit{L1}=\mathit{q2},\quad \mathit{b}=\mathit{bx},\quad \mathit{q2}=\mathit{R2},\\
\mathit{L3}=\mathit{B},\quad \mathit{S2}=\mathit{pack\_rhs},\quad \mathit{n}=\mathit{q8},\\
\mathit{r}=\mathit{R7},\quad \mathit{L9}=\mathit{R9},\quad \mathit{c}=\mathit{R10a},\\
\mathit{k}=\mathit{R10b},\quad \mathit{k}=\mathit{R11},\quad \mathit{a}=\mathit{R12},\\
\mathit{c}=\mathit{R13},\quad \mathit{d}=\mathit{R14},\quad \mathit{L15}=\mathit{R15},\\
\mathit{L16}=\mathit{R16},\quad \mathit{L17}=\mathit{P17},\quad \mathit{mu}=\mathit{R18},\\
\mathit{L19}=\mathit{R19},\quad \mathit{ka}=\mathit{R20},\quad \mathit{S3}=\mathit{sigma},\\
\mathit{t2}=\mathit{Omega}
\end{array}\]


\clearpage\normalsize
\section{The 98-instruction unit-centered certificate}\label{app:98-encoding}

Use the complete rebuilt index of Section~\ref{sec:unit-center}, including
the copied unit coordinates and the direct guard. The coefficient digits
are centered at one, with $\Omega=\lambda-e>0$ and
$S_3=\theta\lambda q-\Omega\mathcal C^2>0$. Computing $3\lambda$
directly reduces the shared geometric block from seven instructions to six.
The result has 55 multiplications and 43 additions, with 34 positive
unknowns and 22 equality tests. All Pell source equations are retained.

The standalone checker \file{round27\_1980\_unit\_center\_certificate.py}
verifies every primitive and source residual. The positive-domain
equivalence and the rebuilt fixed index are proved in
Theorem~\ref{thm:best98}. This and the following two tables are generated
by \file{round27\_29\_1980\_schedule\_tables.py}, which reruns all three
exact checkers. Fixed numerals and equality tests are free.

\small
\input{jones1980_theorem5_98_encoding_schedule.tex}

\clearpage\normalsize
\section{The 98-instruction relaxed-Pell certificate}\label{app:98-pell}

Retain the index and encoding of Section~\ref{sec:composed-99} and replace
only the auxiliary norm by $(ic^2)^2=D(f^2-1)$, where $D=(a+4)^2-1$.
The existing square $(ic^2)^2$ then supplies the parameter in the signed
auxiliary norm directly. Theorem~\ref{thm:best98-pell} proves recovery
of the required integer Pell sequence and auxiliary index divisibility
from the existing bounds, and constructs all positive witnesses.

The standalone checker \file{round28\_1980\_relaxed\_auxiliary\_certificate.py}
verifies all 98 instructions and all 22 source residuals. Its changed
signed-norm correction is
$-F_{16}(K_{\rm source}+K_{\rm calc}-1)(2r+1+jc)^2$;
$F_{16}$ is checked exactly first. The result has 54 multiplications
and 44 additions. No unknown or equality test is added.

\small
\input{jones1980_theorem5_98_pell_schedule.tex}

\clearpage\normalsize
\section{The 97-instruction composed certificate}\label{app:97}

Combine the unit-centered index and direct guard with the relaxed
auxiliary Pell norm. The supplied parameters are $x,V,H,T_L$, where
$T_L=\psi_4(L)$, and there are 34 positive unknowns and 22 equations.
The complete composition proof is Theorem~\ref{thm:best97}.

The standalone checker \file{round29\_1980\_composed\_certificate.py}
checks all 97 primitives and 22 source residuals, their three acyclic
corrections, and both orders of the full source and schedule changes.
Relative to the 99-operation system, only the packed equation, auxiliary
norm, and positive $S_3$ and $\Omega$ equations change. The resulting
count is 54 multiplications and 43 additions. Fixed numerals and equality
tests are free.

\small
% Generated by round27_29_1980_schedule_tables.py; do not edit.
\begin{longtable}{@{}rll@{}}
\toprule
No. & Assignment & Addition/multiplication check\\
\midrule
\endfirsthead
\toprule
No. & Assignment & Addition/multiplication check\\
\midrule
\endhead
\bottomrule
\endfoot
1 & $\mathit{q2} = \mathit{q} \cdot \mathit{q}$ & $\mathit{q} \cdot \mathit{q} = \mathit{q2}$ \\
2 & $\mathit{q4} = \mathit{q2} \cdot \mathit{q2}$ & $\mathit{q2} \cdot \mathit{q2} = \mathit{q4}$ \\
3 & $\mathit{q8} = \mathit{q4} \cdot \mathit{q4}$ & $\mathit{q4} \cdot \mathit{q4} = \mathit{q8}$ \\
4 & $\mathit{b2} = \mathit{b} \cdot \mathit{b}$ & $\mathit{b} \cdot \mathit{b} = \mathit{b2}$ \\
5 & $\mathit{B} = \mathit{H} \cdot \mathit{b2}$ & $\mathit{H} \cdot \mathit{b2} = \mathit{B}$ \\
6 & $\mathit{bx} = \mathit{x} + \mathit{beta}$ & $\mathit{x} + \mathit{beta} = \mathit{bx}$ \\
7 & $\mathit{theta\_lambda} = \mathit{th} \cdot \mathit{la}$ & $\mathit{th} \cdot \mathit{la} = \mathit{theta\_lambda}$ \\
8 & $\mathit{Tcoef} = \mathit{theta\_lambda} + 1$ & $\mathit{theta\_lambda} + 1 = \mathit{Tcoef}$ \\
9 & $\mathit{three\_lambda} = 3 \cdot \mathit{la}$ & $3 \cdot \mathit{la} = \mathit{three\_lambda}$ \\
10 & $\mathit{R2} = \mathit{Tcoef} + \mathit{three\_lambda}$ & $\mathit{Tcoef} + \mathit{three\_lambda} = \mathit{R2}$ \\
11 & $\mathit{t2} = \mathit{la} - \mathit{e}$ & $\mathit{t2} + \mathit{e} = \mathit{la}$ \\
12 & $\mathit{P2} = \mathit{theta\_lambda} \cdot \mathit{q}$ & $\mathit{theta\_lambda} \cdot \mathit{q} = \mathit{P2}$ \\
13 & $\mathit{L3} = \mathit{th} + 4$ & $\mathit{th} + 4 = \mathit{L3}$ \\
14 & $\mathit{eq} = \mathit{e} \cdot \mathit{q}$ & $\mathit{e} \cdot \mathit{q} = \mathit{eq}$ \\
15 & $\mathit{S2} = \mathit{l} + \mathit{eq}$ & $\mathit{l} + \mathit{eq} = \mathit{S2}$ \\
16 & $\mathit{pack\_tth} = \mathit{t} \cdot \mathit{th}$ & $\mathit{t} \cdot \mathit{th} = \mathit{pack\_tth}$ \\
17 & $\mathit{pack\_rhs} = \mathit{V} + \mathit{pack\_tth}$ & $\mathit{V} + \mathit{pack\_tth} = \mathit{pack\_rhs}$ \\
18 & $\mathit{C} = \mathit{x} + \mathit{g}$ & $\mathit{x} + \mathit{g} = \mathit{C}$ \\
19 & $\mathit{C2} = \mathit{C} \cdot \mathit{C}$ & $\mathit{C} \cdot \mathit{C} = \mathit{C2}$ \\
20 & $\mathit{L1} = \mathit{S2} + \mathit{al}$ & $\mathit{S2} + \mathit{al} = \mathit{L1}$ \\
21 & $\mathit{P1} = \mathit{t2} \cdot \mathit{C2}$ & $\mathit{t2} \cdot \mathit{C2} = \mathit{P1}$ \\
22 & $\mathit{S3} = \mathit{P2} - \mathit{P1}$ & $\mathit{S3} + \mathit{P1} = \mathit{P2}$ \\
23 & $\mathit{S3q4} = \mathit{S3} \cdot \mathit{q2}$ & $\mathit{S3} \cdot \mathit{q2} = \mathit{S3q4}$ \\
24 & $\mathit{Sin} = \mathit{S2} + \mathit{S3q4}$ & $\mathit{S2} + \mathit{S3q4} = \mathit{Sin}$ \\
25 & $\mathit{Sq4} = \mathit{Sin} \cdot \mathit{q2}$ & $\mathit{Sin} \cdot \mathit{q2} = \mathit{Sq4}$ \\
26 & $\mathit{S} = \mathit{g} + \mathit{Sq4}$ & $\mathit{g} + \mathit{Sq4} = \mathit{S}$ \\
27 & $\mathit{n2} = \mathit{n} \cdot \mathit{n}$ & $\mathit{n} \cdot \mathit{n} = \mathit{n2}$ \\
28 & $\mathit{n2n} = \mathit{n2} - \mathit{n}$ & $\mathit{n2n} + \mathit{n} = \mathit{n2}$ \\
29 & $\mathit{SA} = \mathit{S} \cdot \mathit{n2n}$ & $\mathit{S} \cdot \mathit{n2n} = \mathit{SA}$ \\
30 & $\mathit{Tq4} = \mathit{Tcoef} \cdot \mathit{q2}$ & $\mathit{Tcoef} \cdot \mathit{q2} = \mathit{Tq4}$ \\
31 & $\mathit{bm1} = \mathit{b} - 1$ & $\mathit{bm1} + 1 = \mathit{b}$ \\
32 & $\mathit{lbm1} = \mathit{l} \cdot \mathit{bm1}$ & $\mathit{l} \cdot \mathit{bm1} = \mathit{lbm1}$ \\
33 & $\mathit{T2} = \mathit{Tq4} - \mathit{lbm1}$ & $\mathit{T2} + \mathit{lbm1} = \mathit{Tq4}$ \\
34 & $\mathit{packed\_target\_shift} = \mathit{l} \cdot \mathit{q4}$ & $\mathit{l} \cdot \mathit{q4} = \mathit{packed\_target\_shift}$ \\
35 & $\mathit{mask\_term} = \mathit{th} \cdot \mathit{packed\_target\_shift}$ & $\mathit{th} \cdot \mathit{packed\_target\_shift} = \mathit{mask\_term}$ \\
36 & $\mathit{T} = \mathit{T2} + \mathit{mask\_term}$ & $\mathit{T2} + \mathit{mask\_term} = \mathit{T}$ \\
37 & $\mathit{n2m1} = \mathit{n2} - 1$ & $\mathit{n2m1} + 1 = \mathit{n2}$ \\
38 & $\mathit{TB} = \mathit{T} \cdot \mathit{n2m1}$ & $\mathit{T} \cdot \mathit{n2m1} = \mathit{TB}$ \\
39 & $\mathit{R7} = \mathit{SA} + \mathit{TB}$ & $\mathit{SA} + \mathit{TB} = \mathit{R7}$ \\
40 & $\mathit{wn2} = \mathit{w} \cdot \mathit{n2}$ & $\mathit{w} \cdot \mathit{n2} = \mathit{wn2}$ \\
41 & $\mathit{sn2} = \mathit{s} \cdot \mathit{n2}$ & $\mathit{s} \cdot \mathit{n2} = \mathit{sn2}$ \\
42 & $\mathit{UM} = \mathit{wn2} \cdot \mathit{sn2}$ & $\mathit{wn2} \cdot \mathit{sn2} = \mathit{UM}$ \\
43 & $\mathit{R12} = \mathit{UM} + \mathit{sn2}$ & $\mathit{UM} + \mathit{sn2} = \mathit{R12}$ \\
44 & $\mathit{ksn2} = \mathit{k} \cdot \mathit{sn2}$ & $\mathit{k} \cdot \mathit{sn2} = \mathit{ksn2}$ \\
45 & $\mathit{UM2} = \mathit{UM} \cdot \mathit{UM}$ & $\mathit{UM} \cdot \mathit{UM} = \mathit{UM2}$ \\
46 & $\mathit{scaled\_norm\_coefficient} = \mathit{UM2} + \mathit{wn2}$ & $\mathit{UM2} + \mathit{wn2} = \mathit{scaled\_norm\_coefficient}$ \\
47 & $\mathit{ratio\_product2} = \mathit{ksn2} \cdot \mathit{ksn2}$ & $\mathit{ksn2} \cdot \mathit{ksn2} = \mathit{ratio\_product2}$ \\
48 & $\mathit{L9} = \mathit{scaled\_norm\_coefficient} \cdot \mathit{ratio\_product2}$ & $\mathit{scaled\_norm\_coefficient} \cdot \mathit{ratio\_product2} = \mathit{L9}$ \\
49 & $\mathit{tauplus1} = \mathit{tau} + 1$ & $\mathit{tau} + 1 = \mathit{tauplus1}$ \\
50 & $\mathit{R9} = \mathit{tau} \cdot \mathit{tauplus1}$ & $\mathit{tau} \cdot \mathit{tauplus1} = \mathit{R9}$ \\
51 & $\mathit{R10a} = \mathit{ksn2} + \mathit{eta}$ & $\mathit{ksn2} + \mathit{eta} = \mathit{R10a}$ \\
52 & $\mathit{R10b} = \mathit{eta} + \mathit{zeta}$ & $\mathit{eta} + \mathit{zeta} = \mathit{R10b}$ \\
53 & $\mathit{r1} = \mathit{r} + 1$ & $\mathit{r} + 1 = \mathit{r1}$ \\
54 & $\mathit{hpm1} = \mathit{h} \cdot \mathit{UM}$ & $\mathit{h} \cdot \mathit{UM} = \mathit{hpm1}$ \\
55 & $\mathit{R11} = \mathit{r1} + \mathit{hpm1}$ & $\mathit{r1} + \mathit{hpm1} = \mathit{R11}$ \\
56 & $\mathit{tr1} = \mathit{r1} + \mathit{r}$ & $\mathit{r1} + \mathit{r} = \mathit{tr1}$ \\
57 & $\mathit{R13} = \mathit{ka} + \mathit{phi}$ & $\mathit{ka} + \mathit{phi} = \mathit{R13}$ \\
58 & $\mathit{cam2} = \mathit{c} \cdot \mathit{a}$ & $\mathit{c} \cdot \mathit{a} = \mathit{cam2}$ \\
59 & $\mathit{D1} = \mathit{wn2} + \mathit{cam2}$ & $\mathit{wn2} + \mathit{cam2} = \mathit{D1}$ \\
60 & $\mathit{a4} = 8 \cdot \mathit{a}$ & $8 \cdot \mathit{a} = \mathit{a4}$ \\
61 & $\mathit{a4m5} = \mathit{a4} + 15$ & $\mathit{a4} + 15 = \mathit{a4m5}$ \\
62 & $\mathit{gam} = \mathit{ga} \cdot \mathit{a4m5}$ & $\mathit{ga} \cdot \mathit{a4m5} = \mathit{gam}$ \\
63 & $\mathit{R14} = \mathit{D1} + \mathit{gam}$ & $\mathit{D1} + \mathit{gam} = \mathit{R14}$ \\
64 & $\mathit{a\_square} = \mathit{a} \cdot \mathit{a}$ & $\mathit{a} \cdot \mathit{a} = \mathit{a\_square}$ \\
65 & $\mathit{A} = \mathit{a\_square} + \mathit{a4m5}$ & $\mathit{a\_square} + \mathit{a4m5} = \mathit{A}$ \\
66 & $\mathit{c2} = \mathit{c} \cdot \mathit{c}$ & $\mathit{c} \cdot \mathit{c} = \mathit{c2}$ \\
67 & $\mathit{Ac2} = \mathit{A} \cdot \mathit{c2}$ & $\mathit{A} \cdot \mathit{c2} = \mathit{Ac2}$ \\
68 & $\mathit{R15} = \mathit{Ac2} + 1$ & $\mathit{Ac2} + 1 = \mathit{R15}$ \\
69 & $\mathit{L15} = \mathit{d} \cdot \mathit{d}$ & $\mathit{d} \cdot \mathit{d} = \mathit{L15}$ \\
70 & $\mathit{ic2} = \mathit{i} \cdot \mathit{c2}$ & $\mathit{i} \cdot \mathit{c2} = \mathit{ic2}$ \\
71 & $\mathit{ic22} = \mathit{ic2} \cdot \mathit{ic2}$ & $\mathit{ic2} \cdot \mathit{ic2} = \mathit{ic22}$ \\
72 & $\mathit{L16} = \mathit{f} \cdot \mathit{f}$ & $\mathit{f} \cdot \mathit{f} = \mathit{L16}$ \\
73 & $\mathit{f\_square\_minus\_one} = \mathit{L16} - 1$ & $\mathit{f\_square\_minus\_one} + 1 = \mathit{L16}$ \\
74 & $\mathit{R16} = \mathit{A} \cdot \mathit{f\_square\_minus\_one}$ & $\mathit{A} \cdot \mathit{f\_square\_minus\_one} = \mathit{R16}$ \\
75 & $\mathit{of} = \mathit{o} \cdot \mathit{f}$ & $\mathit{o} \cdot \mathit{f} = \mathit{of}$ \\
76 & $\mathit{dof} = \mathit{of} - \mathit{L15}$ & $\mathit{dof} + \mathit{L15} = \mathit{of}$ \\
77 & $\mathit{vplus1} = \mathit{dof} + 1$ & $\mathit{dof} + 1 = \mathit{vplus1}$ \\
78 & $\mathit{L17} = \mathit{dof} \cdot \mathit{vplus1}$ & $\mathit{dof} \cdot \mathit{vplus1} = \mathit{L17}$ \\
79 & $\mathit{Gplus1} = \mathit{ic22} - 1$ & $\mathit{Gplus1} + 1 = \mathit{ic22}$ \\
80 & $\mathit{G2m1} = \mathit{ic22} \cdot \mathit{Gplus1}$ & $\mathit{ic22} \cdot \mathit{Gplus1} = \mathit{G2m1}$ \\
81 & $\mathit{jc} = \mathit{j} \cdot \mathit{c}$ & $\mathit{j} \cdot \mathit{c} = \mathit{jc}$ \\
82 & $\mathit{H17} = \mathit{tr1} + \mathit{jc}$ & $\mathit{tr1} + \mathit{jc} = \mathit{H17}$ \\
83 & $\mathit{H2} = \mathit{H17} \cdot \mathit{H17}$ & $\mathit{H17} \cdot \mathit{H17} = \mathit{H2}$ \\
84 & $\mathit{P17} = \mathit{G2m1} \cdot \mathit{H2}$ & $\mathit{G2m1} \cdot \mathit{H2} = \mathit{P17}$ \\
85 & $\mathit{amb} = \mathit{a} - \mathit{th}$ & $\mathit{amb} + \mathit{th} = \mathit{a}$ \\
86 & $\mathit{kamb} = \mathit{ka} \cdot \mathit{amb}$ & $\mathit{ka} \cdot \mathit{amb} = \mathit{kamb}$ \\
87 & $\mathit{qk} = \mathit{q} + \mathit{kamb}$ & $\mathit{q} + \mathit{kamb} = \mathit{qk}$ \\
88 & $\mathit{amb2} = \mathit{amb} \cdot \mathit{amb}$ & $\mathit{amb} \cdot \mathit{amb} = \mathit{amb2}$ \\
89 & $\mathit{Mod} = \mathit{A} - \mathit{amb2}$ & $\mathit{Mod} + \mathit{amb2} = \mathit{A}$ \\
90 & $\mathit{rM} = \mathit{rho} \cdot \mathit{Mod}$ & $\mathit{rho} \cdot \mathit{Mod} = \mathit{rM}$ \\
91 & $\mathit{R18} = \mathit{qk} + \mathit{rM}$ & $\mathit{qk} + \mathit{rM} = \mathit{R18}$ \\
92 & $\mathit{ka2} = \mathit{ka} \cdot \mathit{ka}$ & $\mathit{ka} \cdot \mathit{ka} = \mathit{ka2}$ \\
93 & $\mathit{Aka2} = \mathit{A} \cdot \mathit{ka2}$ & $\mathit{A} \cdot \mathit{ka2} = \mathit{Aka2}$ \\
94 & $\mathit{R19} = \mathit{Aka2} + 1$ & $\mathit{Aka2} + 1 = \mathit{R19}$ \\
95 & $\mathit{L19} = \mathit{mu} \cdot \mathit{mu}$ & $\mathit{mu} \cdot \mathit{mu} = \mathit{L19}$ \\
96 & $\mathit{Dam1} = \mathit{Delta} \cdot \mathit{a}$ & $\mathit{Delta} \cdot \mathit{a} = \mathit{Dam1}$ \\
97 & $\mathit{R20} = \mathit{Dam1} + \mathit{Tindex}$ & $\mathit{Dam1} + \mathit{Tindex} = \mathit{R20}$ \\
\end{longtable}

The twenty-two free equality tests are
\[\begin{array}{@{}l@{}}
\mathit{L1}=\mathit{q2},\quad \mathit{b}=\mathit{bx},\quad \mathit{q2}=\mathit{R2},\\
\mathit{L3}=\mathit{B},\quad \mathit{S2}=\mathit{pack\_rhs},\quad \mathit{n}=\mathit{q8},\\
\mathit{r}=\mathit{R7},\quad \mathit{L9}=\mathit{R9},\quad \mathit{c}=\mathit{R10a},\\
\mathit{k}=\mathit{R10b},\quad \mathit{k}=\mathit{R11},\quad \mathit{a}=\mathit{R12},\\
\mathit{c}=\mathit{R13},\quad \mathit{d}=\mathit{R14},\quad \mathit{L15}=\mathit{R15},\\
\mathit{ic22}=\mathit{R16},\quad \mathit{L17}=\mathit{P17},\quad \mathit{mu}=\mathit{R18},\\
\mathit{L19}=\mathit{R19},\quad \mathit{ka}=\mathit{R20},\quad \mathit{S3}=\mathit{sigma},\\
\mathit{t2}=\mathit{Omega}
\end{array}\]


\clearpage\normalsize
\section{The 96-instruction linear-radix certificate}\label{app:96}

Use the rebuilt index of Section~\ref{sec:linear-radix}, including the
paired zero targets, both unit tests and their carry-reset padding.
The supplied parameters are $x,V,H,T_L$, with $T_L=\psi_4(L)$.
There are 35 positive unknowns and 23 equations; the additional unknown
is $\alpha_2$, printed as $\mathit{al2}$ in the table.

The radix is $B=Hb$, and the third packed block is
$S_3=(\lambda-e)(q-\mathcal C^2)$ with $\mathcal C=x+g$.
The new positive bound $\ell+\alpha_2=q$ removes the quotient alias.
Theorem~\ref{thm:best96} proves the full encoding and positive-witness
equivalence. Its unit construction is part of the fixed coefficient
index and adds no arithmetic instruction to this table.

The standalone checker \file{round30\_1980\_linear\_radix\_certificate.py}
verifies every primitive, all 23 source residuals and the three acyclic
corrections. The geometric and second-exponent corrections use the
independently exact $\theta+4=B$ equation; the signed auxiliary correction
uses the exact relaxed norm. There are 52 multiplications and 44 additions.
The separate sparse encoding checker verifies finite coefficient and
carry regressions; it does not replace the general proof.

This table is generated by \file{round30\_1980\_schedule\_table.py},
which reruns the complete arithmetic checker. Fixed numerals and
equality tests are free.

\small
% Generated by round30_1980_schedule_table.py; do not edit.
\begin{longtable}{@{}rll@{}}
\toprule
No. & Assignment & Addition/multiplication check\\
\midrule
\endfirsthead
\toprule
No. & Assignment & Addition/multiplication check\\
\midrule
\endhead
\bottomrule
\endfoot
1 & $\mathit{q2} = \mathit{q} \cdot \mathit{q}$ & $\mathit{q} \cdot \mathit{q} = \mathit{q2}$ \\
2 & $\mathit{q4} = \mathit{q2} \cdot \mathit{q2}$ & $\mathit{q2} \cdot \mathit{q2} = \mathit{q4}$ \\
3 & $\mathit{q8} = \mathit{q4} \cdot \mathit{q4}$ & $\mathit{q4} \cdot \mathit{q4} = \mathit{q8}$ \\
4 & $\mathit{B} = \mathit{H} \cdot \mathit{b}$ & $\mathit{H} \cdot \mathit{b} = \mathit{B}$ \\
5 & $\mathit{bx} = \mathit{x} + \mathit{beta}$ & $\mathit{x} + \mathit{beta} = \mathit{bx}$ \\
6 & $\mathit{theta\_lambda} = \mathit{th} \cdot \mathit{la}$ & $\mathit{th} \cdot \mathit{la} = \mathit{theta\_lambda}$ \\
7 & $\mathit{Tcoef} = \mathit{theta\_lambda} + 1$ & $\mathit{theta\_lambda} + 1 = \mathit{Tcoef}$ \\
8 & $\mathit{three\_lambda} = 3 \cdot \mathit{la}$ & $3 \cdot \mathit{la} = \mathit{three\_lambda}$ \\
9 & $\mathit{R2} = \mathit{Tcoef} + \mathit{three\_lambda}$ & $\mathit{Tcoef} + \mathit{three\_lambda} = \mathit{R2}$ \\
10 & $\mathit{t2} = \mathit{la} - \mathit{e}$ & $\mathit{t2} + \mathit{e} = \mathit{la}$ \\
11 & $\mathit{L3} = \mathit{th} + 4$ & $\mathit{th} + 4 = \mathit{L3}$ \\
12 & $\mathit{eq} = \mathit{e} \cdot \mathit{q}$ & $\mathit{e} \cdot \mathit{q} = \mathit{eq}$ \\
13 & $\mathit{S2} = \mathit{l} + \mathit{eq}$ & $\mathit{l} + \mathit{eq} = \mathit{S2}$ \\
14 & $\mathit{indicator\_bound} = \mathit{l} + \mathit{al2}$ & $\mathit{l} + \mathit{al2} = \mathit{indicator\_bound}$ \\
15 & $\mathit{pack\_tth} = \mathit{t} \cdot \mathit{th}$ & $\mathit{t} \cdot \mathit{th} = \mathit{pack\_tth}$ \\
16 & $\mathit{pack\_rhs} = \mathit{V} + \mathit{pack\_tth}$ & $\mathit{V} + \mathit{pack\_tth} = \mathit{pack\_rhs}$ \\
17 & $\mathit{C} = \mathit{x} + \mathit{g}$ & $\mathit{x} + \mathit{g} = \mathit{C}$ \\
18 & $\mathit{C2} = \mathit{C} \cdot \mathit{C}$ & $\mathit{C} \cdot \mathit{C} = \mathit{C2}$ \\
19 & $\mathit{L1} = \mathit{S2} + \mathit{al}$ & $\mathit{S2} + \mathit{al} = \mathit{L1}$ \\
20 & $\mathit{code\_gap} = \mathit{q} - \mathit{C2}$ & $\mathit{code\_gap} + \mathit{C2} = \mathit{q}$ \\
21 & $\mathit{S3} = \mathit{t2} \cdot \mathit{code\_gap}$ & $\mathit{t2} \cdot \mathit{code\_gap} = \mathit{S3}$ \\
22 & $\mathit{S3q4} = \mathit{S3} \cdot \mathit{q2}$ & $\mathit{S3} \cdot \mathit{q2} = \mathit{S3q4}$ \\
23 & $\mathit{Sin} = \mathit{S2} + \mathit{S3q4}$ & $\mathit{S2} + \mathit{S3q4} = \mathit{Sin}$ \\
24 & $\mathit{Sq4} = \mathit{Sin} \cdot \mathit{q2}$ & $\mathit{Sin} \cdot \mathit{q2} = \mathit{Sq4}$ \\
25 & $\mathit{S} = \mathit{g} + \mathit{Sq4}$ & $\mathit{g} + \mathit{Sq4} = \mathit{S}$ \\
26 & $\mathit{n2} = \mathit{n} \cdot \mathit{n}$ & $\mathit{n} \cdot \mathit{n} = \mathit{n2}$ \\
27 & $\mathit{n2n} = \mathit{n2} - \mathit{n}$ & $\mathit{n2n} + \mathit{n} = \mathit{n2}$ \\
28 & $\mathit{SA} = \mathit{S} \cdot \mathit{n2n}$ & $\mathit{S} \cdot \mathit{n2n} = \mathit{SA}$ \\
29 & $\mathit{Tq4} = \mathit{Tcoef} \cdot \mathit{q2}$ & $\mathit{Tcoef} \cdot \mathit{q2} = \mathit{Tq4}$ \\
30 & $\mathit{bm1} = \mathit{b} - 1$ & $\mathit{bm1} + 1 = \mathit{b}$ \\
31 & $\mathit{lbm1} = \mathit{l} \cdot \mathit{bm1}$ & $\mathit{l} \cdot \mathit{bm1} = \mathit{lbm1}$ \\
32 & $\mathit{T2} = \mathit{Tq4} - \mathit{lbm1}$ & $\mathit{T2} + \mathit{lbm1} = \mathit{Tq4}$ \\
33 & $\mathit{packed\_target\_shift} = \mathit{l} \cdot \mathit{q4}$ & $\mathit{l} \cdot \mathit{q4} = \mathit{packed\_target\_shift}$ \\
34 & $\mathit{mask\_term} = \mathit{th} \cdot \mathit{packed\_target\_shift}$ & $\mathit{th} \cdot \mathit{packed\_target\_shift} = \mathit{mask\_term}$ \\
35 & $\mathit{T} = \mathit{T2} + \mathit{mask\_term}$ & $\mathit{T2} + \mathit{mask\_term} = \mathit{T}$ \\
36 & $\mathit{n2m1} = \mathit{n2} - 1$ & $\mathit{n2m1} + 1 = \mathit{n2}$ \\
37 & $\mathit{TB} = \mathit{T} \cdot \mathit{n2m1}$ & $\mathit{T} \cdot \mathit{n2m1} = \mathit{TB}$ \\
38 & $\mathit{R7} = \mathit{SA} + \mathit{TB}$ & $\mathit{SA} + \mathit{TB} = \mathit{R7}$ \\
39 & $\mathit{wn2} = \mathit{w} \cdot \mathit{n2}$ & $\mathit{w} \cdot \mathit{n2} = \mathit{wn2}$ \\
40 & $\mathit{sn2} = \mathit{s} \cdot \mathit{n2}$ & $\mathit{s} \cdot \mathit{n2} = \mathit{sn2}$ \\
41 & $\mathit{UM} = \mathit{wn2} \cdot \mathit{sn2}$ & $\mathit{wn2} \cdot \mathit{sn2} = \mathit{UM}$ \\
42 & $\mathit{R12} = \mathit{UM} + \mathit{sn2}$ & $\mathit{UM} + \mathit{sn2} = \mathit{R12}$ \\
43 & $\mathit{ksn2} = \mathit{k} \cdot \mathit{sn2}$ & $\mathit{k} \cdot \mathit{sn2} = \mathit{ksn2}$ \\
44 & $\mathit{UM2} = \mathit{UM} \cdot \mathit{UM}$ & $\mathit{UM} \cdot \mathit{UM} = \mathit{UM2}$ \\
45 & $\mathit{scaled\_norm\_coefficient} = \mathit{UM2} + \mathit{wn2}$ & $\mathit{UM2} + \mathit{wn2} = \mathit{scaled\_norm\_coefficient}$ \\
46 & $\mathit{ratio\_product2} = \mathit{ksn2} \cdot \mathit{ksn2}$ & $\mathit{ksn2} \cdot \mathit{ksn2} = \mathit{ratio\_product2}$ \\
47 & $\mathit{L9} = \mathit{scaled\_norm\_coefficient} \cdot \mathit{ratio\_product2}$ & $\mathit{scaled\_norm\_coefficient} \cdot \mathit{ratio\_product2} = \mathit{L9}$ \\
48 & $\mathit{tauplus1} = \mathit{tau} + 1$ & $\mathit{tau} + 1 = \mathit{tauplus1}$ \\
49 & $\mathit{R9} = \mathit{tau} \cdot \mathit{tauplus1}$ & $\mathit{tau} \cdot \mathit{tauplus1} = \mathit{R9}$ \\
50 & $\mathit{R10a} = \mathit{ksn2} + \mathit{eta}$ & $\mathit{ksn2} + \mathit{eta} = \mathit{R10a}$ \\
51 & $\mathit{R10b} = \mathit{eta} + \mathit{zeta}$ & $\mathit{eta} + \mathit{zeta} = \mathit{R10b}$ \\
52 & $\mathit{r1} = \mathit{r} + 1$ & $\mathit{r} + 1 = \mathit{r1}$ \\
53 & $\mathit{hpm1} = \mathit{h} \cdot \mathit{UM}$ & $\mathit{h} \cdot \mathit{UM} = \mathit{hpm1}$ \\
54 & $\mathit{R11} = \mathit{r1} + \mathit{hpm1}$ & $\mathit{r1} + \mathit{hpm1} = \mathit{R11}$ \\
55 & $\mathit{tr1} = \mathit{r1} + \mathit{r}$ & $\mathit{r1} + \mathit{r} = \mathit{tr1}$ \\
56 & $\mathit{R13} = \mathit{ka} + \mathit{phi}$ & $\mathit{ka} + \mathit{phi} = \mathit{R13}$ \\
57 & $\mathit{cam2} = \mathit{c} \cdot \mathit{a}$ & $\mathit{c} \cdot \mathit{a} = \mathit{cam2}$ \\
58 & $\mathit{D1} = \mathit{wn2} + \mathit{cam2}$ & $\mathit{wn2} + \mathit{cam2} = \mathit{D1}$ \\
59 & $\mathit{a4} = 8 \cdot \mathit{a}$ & $8 \cdot \mathit{a} = \mathit{a4}$ \\
60 & $\mathit{a4m5} = \mathit{a4} + 15$ & $\mathit{a4} + 15 = \mathit{a4m5}$ \\
61 & $\mathit{gam} = \mathit{ga} \cdot \mathit{a4m5}$ & $\mathit{ga} \cdot \mathit{a4m5} = \mathit{gam}$ \\
62 & $\mathit{R14} = \mathit{D1} + \mathit{gam}$ & $\mathit{D1} + \mathit{gam} = \mathit{R14}$ \\
63 & $\mathit{a\_square} = \mathit{a} \cdot \mathit{a}$ & $\mathit{a} \cdot \mathit{a} = \mathit{a\_square}$ \\
64 & $\mathit{A} = \mathit{a\_square} + \mathit{a4m5}$ & $\mathit{a\_square} + \mathit{a4m5} = \mathit{A}$ \\
65 & $\mathit{c2} = \mathit{c} \cdot \mathit{c}$ & $\mathit{c} \cdot \mathit{c} = \mathit{c2}$ \\
66 & $\mathit{Ac2} = \mathit{A} \cdot \mathit{c2}$ & $\mathit{A} \cdot \mathit{c2} = \mathit{Ac2}$ \\
67 & $\mathit{R15} = \mathit{Ac2} + 1$ & $\mathit{Ac2} + 1 = \mathit{R15}$ \\
68 & $\mathit{L15} = \mathit{d} \cdot \mathit{d}$ & $\mathit{d} \cdot \mathit{d} = \mathit{L15}$ \\
69 & $\mathit{ic2} = \mathit{i} \cdot \mathit{c2}$ & $\mathit{i} \cdot \mathit{c2} = \mathit{ic2}$ \\
70 & $\mathit{ic22} = \mathit{ic2} \cdot \mathit{ic2}$ & $\mathit{ic2} \cdot \mathit{ic2} = \mathit{ic22}$ \\
71 & $\mathit{L16} = \mathit{f} \cdot \mathit{f}$ & $\mathit{f} \cdot \mathit{f} = \mathit{L16}$ \\
72 & $\mathit{f\_square\_minus\_one} = \mathit{L16} - 1$ & $\mathit{f\_square\_minus\_one} + 1 = \mathit{L16}$ \\
73 & $\mathit{R16} = \mathit{A} \cdot \mathit{f\_square\_minus\_one}$ & $\mathit{A} \cdot \mathit{f\_square\_minus\_one} = \mathit{R16}$ \\
74 & $\mathit{of} = \mathit{o} \cdot \mathit{f}$ & $\mathit{o} \cdot \mathit{f} = \mathit{of}$ \\
75 & $\mathit{dof} = \mathit{of} - \mathit{L15}$ & $\mathit{dof} + \mathit{L15} = \mathit{of}$ \\
76 & $\mathit{vplus1} = \mathit{dof} + 1$ & $\mathit{dof} + 1 = \mathit{vplus1}$ \\
77 & $\mathit{L17} = \mathit{dof} \cdot \mathit{vplus1}$ & $\mathit{dof} \cdot \mathit{vplus1} = \mathit{L17}$ \\
78 & $\mathit{Gplus1} = \mathit{ic22} - 1$ & $\mathit{Gplus1} + 1 = \mathit{ic22}$ \\
79 & $\mathit{G2m1} = \mathit{ic22} \cdot \mathit{Gplus1}$ & $\mathit{ic22} \cdot \mathit{Gplus1} = \mathit{G2m1}$ \\
80 & $\mathit{jc} = \mathit{j} \cdot \mathit{c}$ & $\mathit{j} \cdot \mathit{c} = \mathit{jc}$ \\
81 & $\mathit{H17} = \mathit{tr1} + \mathit{jc}$ & $\mathit{tr1} + \mathit{jc} = \mathit{H17}$ \\
82 & $\mathit{H2} = \mathit{H17} \cdot \mathit{H17}$ & $\mathit{H17} \cdot \mathit{H17} = \mathit{H2}$ \\
83 & $\mathit{P17} = \mathit{G2m1} \cdot \mathit{H2}$ & $\mathit{G2m1} \cdot \mathit{H2} = \mathit{P17}$ \\
84 & $\mathit{amb} = \mathit{a} - \mathit{th}$ & $\mathit{amb} + \mathit{th} = \mathit{a}$ \\
85 & $\mathit{kamb} = \mathit{ka} \cdot \mathit{amb}$ & $\mathit{ka} \cdot \mathit{amb} = \mathit{kamb}$ \\
86 & $\mathit{qk} = \mathit{q} + \mathit{kamb}$ & $\mathit{q} + \mathit{kamb} = \mathit{qk}$ \\
87 & $\mathit{amb2} = \mathit{amb} \cdot \mathit{amb}$ & $\mathit{amb} \cdot \mathit{amb} = \mathit{amb2}$ \\
88 & $\mathit{Mod} = \mathit{A} - \mathit{amb2}$ & $\mathit{Mod} + \mathit{amb2} = \mathit{A}$ \\
89 & $\mathit{rM} = \mathit{rho} \cdot \mathit{Mod}$ & $\mathit{rho} \cdot \mathit{Mod} = \mathit{rM}$ \\
90 & $\mathit{R18} = \mathit{qk} + \mathit{rM}$ & $\mathit{qk} + \mathit{rM} = \mathit{R18}$ \\
91 & $\mathit{ka2} = \mathit{ka} \cdot \mathit{ka}$ & $\mathit{ka} \cdot \mathit{ka} = \mathit{ka2}$ \\
92 & $\mathit{Aka2} = \mathit{A} \cdot \mathit{ka2}$ & $\mathit{A} \cdot \mathit{ka2} = \mathit{Aka2}$ \\
93 & $\mathit{R19} = \mathit{Aka2} + 1$ & $\mathit{Aka2} + 1 = \mathit{R19}$ \\
94 & $\mathit{L19} = \mathit{mu} \cdot \mathit{mu}$ & $\mathit{mu} \cdot \mathit{mu} = \mathit{L19}$ \\
95 & $\mathit{Dam1} = \mathit{Delta} \cdot \mathit{a}$ & $\mathit{Delta} \cdot \mathit{a} = \mathit{Dam1}$ \\
96 & $\mathit{R20} = \mathit{Dam1} + \mathit{Tindex}$ & $\mathit{Dam1} + \mathit{Tindex} = \mathit{R20}$ \\
\end{longtable}

The twenty-three free equality tests are
\[\begin{array}{@{}l@{}}
\mathit{L1}=\mathit{q2},\quad \mathit{b}=\mathit{bx},\quad \mathit{q2}=\mathit{R2},\\
\mathit{L3}=\mathit{B},\quad \mathit{S2}=\mathit{pack\_rhs},\quad \mathit{n}=\mathit{q8},\\
\mathit{r}=\mathit{R7},\quad \mathit{L9}=\mathit{R9},\quad \mathit{c}=\mathit{R10a},\\
\mathit{k}=\mathit{R10b},\quad \mathit{k}=\mathit{R11},\quad \mathit{a}=\mathit{R12},\\
\mathit{c}=\mathit{R13},\quad \mathit{d}=\mathit{R14},\quad \mathit{L15}=\mathit{R15},\\
\mathit{ic22}=\mathit{R16},\quad \mathit{L17}=\mathit{P17},\quad \mathit{mu}=\mathit{R18},\\
\mathit{L19}=\mathit{R19},\quad \mathit{ka}=\mathit{R20},\quad \mathit{S3}=\mathit{sigma},\\
\mathit{t2}=\mathit{Omega},\quad \mathit{indicator\_bound}=\mathit{q}
\end{array}\]


\clearpage\normalsize
\section{The 95-instruction affine-radix certificate}\label{app:95}

Use the rebuilt index of Section~\ref{sec:affine-radix}. The fixed
parameter $H$ is now $H_0+1$, where $H_0$ is a sufficiently large
power of two. The supplied parameters are $x,V,H,T_L$, with
$T_L=\psi_4(L)$. There are 34 positive unknowns and 22 equations.

Relative to Appendix~\ref{app:96}, combine the two code bounds as
$\ell+e+\alpha=q$ and remove $\alpha_2$. This still costs two additions.
Replace $B=Hb$ by $B=H+b$, and replace the coefficient $b-1$ in the
first mask by $b$. The resulting first mask forbids exactly the
$H_0$ bit at each coordinate position. The three-way coordinate
representation, three-digit target windows, reset padding, and
independent five- and seven-fold padding carries all belong to the
fixed index. Theorem~\ref{thm:best95} proves the complete equivalence.

The standalone checker \file{round32\_1980\_affine\_radix\_certificate.py}
verifies all 95 primitive instructions and all 22 source residuals,
including their three acyclic corrections. The geometric and
second-exponent corrections use the independently exact
$\theta+4=B$ equation; the signed auxiliary correction uses the
exact relaxed norm. There are 51 multiplications and 44 additions.

The table is generated by \file{round32\_1980\_schedule\_table.py}, which
reruns the complete arithmetic checker. Fixed numerals and equality
tests are free. The register name $\mathit{lbm1}$ is retained from
the preceding schedules; its new value is $\ell b$.

\small
% Generated by round32_1980_schedule_table.py; do not edit.
\begin{longtable}{@{}rll@{}}
\toprule
No. & Assignment & Addition/multiplication check\\
\midrule
\endfirsthead
\toprule
No. & Assignment & Addition/multiplication check\\
\midrule
\endhead
\bottomrule
\endfoot
1 & $\mathit{q2} = \mathit{q} \cdot \mathit{q}$ & $\mathit{q} \cdot \mathit{q} = \mathit{q2}$ \\
2 & $\mathit{q4} = \mathit{q2} \cdot \mathit{q2}$ & $\mathit{q2} \cdot \mathit{q2} = \mathit{q4}$ \\
3 & $\mathit{q8} = \mathit{q4} \cdot \mathit{q4}$ & $\mathit{q4} \cdot \mathit{q4} = \mathit{q8}$ \\
4 & $\mathit{B} = \mathit{H} + \mathit{b}$ & $\mathit{H} + \mathit{b} = \mathit{B}$ \\
5 & $\mathit{bx} = \mathit{x} + \mathit{beta}$ & $\mathit{x} + \mathit{beta} = \mathit{bx}$ \\
6 & $\mathit{theta\_lambda} = \mathit{th} \cdot \mathit{la}$ & $\mathit{th} \cdot \mathit{la} = \mathit{theta\_lambda}$ \\
7 & $\mathit{Tcoef} = \mathit{theta\_lambda} + 1$ & $\mathit{theta\_lambda} + 1 = \mathit{Tcoef}$ \\
8 & $\mathit{three\_lambda} = 3 \cdot \mathit{la}$ & $3 \cdot \mathit{la} = \mathit{three\_lambda}$ \\
9 & $\mathit{R2} = \mathit{Tcoef} + \mathit{three\_lambda}$ & $\mathit{Tcoef} + \mathit{three\_lambda} = \mathit{R2}$ \\
10 & $\mathit{t2} = \mathit{la} - \mathit{e}$ & $\mathit{t2} + \mathit{e} = \mathit{la}$ \\
11 & $\mathit{L3} = \mathit{th} + 4$ & $\mathit{th} + 4 = \mathit{L3}$ \\
12 & $\mathit{eq} = \mathit{e} \cdot \mathit{q}$ & $\mathit{e} \cdot \mathit{q} = \mathit{eq}$ \\
13 & $\mathit{S2} = \mathit{l} + \mathit{eq}$ & $\mathit{l} + \mathit{eq} = \mathit{S2}$ \\
14 & $\mathit{pack\_tth} = \mathit{t} \cdot \mathit{th}$ & $\mathit{t} \cdot \mathit{th} = \mathit{pack\_tth}$ \\
15 & $\mathit{pack\_rhs} = \mathit{V} + \mathit{pack\_tth}$ & $\mathit{V} + \mathit{pack\_tth} = \mathit{pack\_rhs}$ \\
16 & $\mathit{C} = \mathit{x} + \mathit{g}$ & $\mathit{x} + \mathit{g} = \mathit{C}$ \\
17 & $\mathit{C2} = \mathit{C} \cdot \mathit{C}$ & $\mathit{C} \cdot \mathit{C} = \mathit{C2}$ \\
18 & $\mathit{bound\_sum} = \mathit{l} + \mathit{e}$ & $\mathit{l} + \mathit{e} = \mathit{bound\_sum}$ \\
19 & $\mathit{L1} = \mathit{bound\_sum} + \mathit{al}$ & $\mathit{bound\_sum} + \mathit{al} = \mathit{L1}$ \\
20 & $\mathit{code\_gap} = \mathit{q} - \mathit{C2}$ & $\mathit{code\_gap} + \mathit{C2} = \mathit{q}$ \\
21 & $\mathit{S3} = \mathit{t2} \cdot \mathit{code\_gap}$ & $\mathit{t2} \cdot \mathit{code\_gap} = \mathit{S3}$ \\
22 & $\mathit{S3q4} = \mathit{S3} \cdot \mathit{q2}$ & $\mathit{S3} \cdot \mathit{q2} = \mathit{S3q4}$ \\
23 & $\mathit{Sin} = \mathit{S2} + \mathit{S3q4}$ & $\mathit{S2} + \mathit{S3q4} = \mathit{Sin}$ \\
24 & $\mathit{Sq4} = \mathit{Sin} \cdot \mathit{q2}$ & $\mathit{Sin} \cdot \mathit{q2} = \mathit{Sq4}$ \\
25 & $\mathit{S} = \mathit{g} + \mathit{Sq4}$ & $\mathit{g} + \mathit{Sq4} = \mathit{S}$ \\
26 & $\mathit{n2} = \mathit{n} \cdot \mathit{n}$ & $\mathit{n} \cdot \mathit{n} = \mathit{n2}$ \\
27 & $\mathit{n2n} = \mathit{n2} - \mathit{n}$ & $\mathit{n2n} + \mathit{n} = \mathit{n2}$ \\
28 & $\mathit{SA} = \mathit{S} \cdot \mathit{n2n}$ & $\mathit{S} \cdot \mathit{n2n} = \mathit{SA}$ \\
29 & $\mathit{Tq4} = \mathit{Tcoef} \cdot \mathit{q2}$ & $\mathit{Tcoef} \cdot \mathit{q2} = \mathit{Tq4}$ \\
30 & $\mathit{lbm1} = \mathit{l} \cdot \mathit{b}$ & $\mathit{l} \cdot \mathit{b} = \mathit{lbm1}$ \\
31 & $\mathit{T2} = \mathit{Tq4} - \mathit{lbm1}$ & $\mathit{T2} + \mathit{lbm1} = \mathit{Tq4}$ \\
32 & $\mathit{packed\_target\_shift} = \mathit{l} \cdot \mathit{q4}$ & $\mathit{l} \cdot \mathit{q4} = \mathit{packed\_target\_shift}$ \\
33 & $\mathit{mask\_term} = \mathit{th} \cdot \mathit{packed\_target\_shift}$ & $\mathit{th} \cdot \mathit{packed\_target\_shift} = \mathit{mask\_term}$ \\
34 & $\mathit{T} = \mathit{T2} + \mathit{mask\_term}$ & $\mathit{T2} + \mathit{mask\_term} = \mathit{T}$ \\
35 & $\mathit{n2m1} = \mathit{n2} - 1$ & $\mathit{n2m1} + 1 = \mathit{n2}$ \\
36 & $\mathit{TB} = \mathit{T} \cdot \mathit{n2m1}$ & $\mathit{T} \cdot \mathit{n2m1} = \mathit{TB}$ \\
37 & $\mathit{R7} = \mathit{SA} + \mathit{TB}$ & $\mathit{SA} + \mathit{TB} = \mathit{R7}$ \\
38 & $\mathit{wn2} = \mathit{w} \cdot \mathit{n2}$ & $\mathit{w} \cdot \mathit{n2} = \mathit{wn2}$ \\
39 & $\mathit{sn2} = \mathit{s} \cdot \mathit{n2}$ & $\mathit{s} \cdot \mathit{n2} = \mathit{sn2}$ \\
40 & $\mathit{UM} = \mathit{wn2} \cdot \mathit{sn2}$ & $\mathit{wn2} \cdot \mathit{sn2} = \mathit{UM}$ \\
41 & $\mathit{R12} = \mathit{UM} + \mathit{sn2}$ & $\mathit{UM} + \mathit{sn2} = \mathit{R12}$ \\
42 & $\mathit{ksn2} = \mathit{k} \cdot \mathit{sn2}$ & $\mathit{k} \cdot \mathit{sn2} = \mathit{ksn2}$ \\
43 & $\mathit{UM2} = \mathit{UM} \cdot \mathit{UM}$ & $\mathit{UM} \cdot \mathit{UM} = \mathit{UM2}$ \\
44 & $\mathit{scaled\_norm\_coefficient} = \mathit{UM2} + \mathit{wn2}$ & $\mathit{UM2} + \mathit{wn2} = \mathit{scaled\_norm\_coefficient}$ \\
45 & $\mathit{ratio\_product2} = \mathit{ksn2} \cdot \mathit{ksn2}$ & $\mathit{ksn2} \cdot \mathit{ksn2} = \mathit{ratio\_product2}$ \\
46 & $\mathit{L9} = \mathit{scaled\_norm\_coefficient} \cdot \mathit{ratio\_product2}$ & $\mathit{scaled\_norm\_coefficient} \cdot \mathit{ratio\_product2} = \mathit{L9}$ \\
47 & $\mathit{tauplus1} = \mathit{tau} + 1$ & $\mathit{tau} + 1 = \mathit{tauplus1}$ \\
48 & $\mathit{R9} = \mathit{tau} \cdot \mathit{tauplus1}$ & $\mathit{tau} \cdot \mathit{tauplus1} = \mathit{R9}$ \\
49 & $\mathit{R10a} = \mathit{ksn2} + \mathit{eta}$ & $\mathit{ksn2} + \mathit{eta} = \mathit{R10a}$ \\
50 & $\mathit{R10b} = \mathit{eta} + \mathit{zeta}$ & $\mathit{eta} + \mathit{zeta} = \mathit{R10b}$ \\
51 & $\mathit{r1} = \mathit{r} + 1$ & $\mathit{r} + 1 = \mathit{r1}$ \\
52 & $\mathit{hpm1} = \mathit{h} \cdot \mathit{UM}$ & $\mathit{h} \cdot \mathit{UM} = \mathit{hpm1}$ \\
53 & $\mathit{R11} = \mathit{r1} + \mathit{hpm1}$ & $\mathit{r1} + \mathit{hpm1} = \mathit{R11}$ \\
54 & $\mathit{tr1} = \mathit{r1} + \mathit{r}$ & $\mathit{r1} + \mathit{r} = \mathit{tr1}$ \\
55 & $\mathit{R13} = \mathit{ka} + \mathit{phi}$ & $\mathit{ka} + \mathit{phi} = \mathit{R13}$ \\
56 & $\mathit{cam2} = \mathit{c} \cdot \mathit{a}$ & $\mathit{c} \cdot \mathit{a} = \mathit{cam2}$ \\
57 & $\mathit{D1} = \mathit{wn2} + \mathit{cam2}$ & $\mathit{wn2} + \mathit{cam2} = \mathit{D1}$ \\
58 & $\mathit{a4} = 8 \cdot \mathit{a}$ & $8 \cdot \mathit{a} = \mathit{a4}$ \\
59 & $\mathit{a4m5} = \mathit{a4} + 15$ & $\mathit{a4} + 15 = \mathit{a4m5}$ \\
60 & $\mathit{gam} = \mathit{ga} \cdot \mathit{a4m5}$ & $\mathit{ga} \cdot \mathit{a4m5} = \mathit{gam}$ \\
61 & $\mathit{R14} = \mathit{D1} + \mathit{gam}$ & $\mathit{D1} + \mathit{gam} = \mathit{R14}$ \\
62 & $\mathit{a\_square} = \mathit{a} \cdot \mathit{a}$ & $\mathit{a} \cdot \mathit{a} = \mathit{a\_square}$ \\
63 & $\mathit{A} = \mathit{a\_square} + \mathit{a4m5}$ & $\mathit{a\_square} + \mathit{a4m5} = \mathit{A}$ \\
64 & $\mathit{c2} = \mathit{c} \cdot \mathit{c}$ & $\mathit{c} \cdot \mathit{c} = \mathit{c2}$ \\
65 & $\mathit{Ac2} = \mathit{A} \cdot \mathit{c2}$ & $\mathit{A} \cdot \mathit{c2} = \mathit{Ac2}$ \\
66 & $\mathit{R15} = \mathit{Ac2} + 1$ & $\mathit{Ac2} + 1 = \mathit{R15}$ \\
67 & $\mathit{L15} = \mathit{d} \cdot \mathit{d}$ & $\mathit{d} \cdot \mathit{d} = \mathit{L15}$ \\
68 & $\mathit{ic2} = \mathit{i} \cdot \mathit{c2}$ & $\mathit{i} \cdot \mathit{c2} = \mathit{ic2}$ \\
69 & $\mathit{ic22} = \mathit{ic2} \cdot \mathit{ic2}$ & $\mathit{ic2} \cdot \mathit{ic2} = \mathit{ic22}$ \\
70 & $\mathit{L16} = \mathit{f} \cdot \mathit{f}$ & $\mathit{f} \cdot \mathit{f} = \mathit{L16}$ \\
71 & $\mathit{f\_square\_minus\_one} = \mathit{L16} - 1$ & $\mathit{f\_square\_minus\_one} + 1 = \mathit{L16}$ \\
72 & $\mathit{R16} = \mathit{A} \cdot \mathit{f\_square\_minus\_one}$ & $\mathit{A} \cdot \mathit{f\_square\_minus\_one} = \mathit{R16}$ \\
73 & $\mathit{of} = \mathit{o} \cdot \mathit{f}$ & $\mathit{o} \cdot \mathit{f} = \mathit{of}$ \\
74 & $\mathit{dof} = \mathit{of} - \mathit{L15}$ & $\mathit{dof} + \mathit{L15} = \mathit{of}$ \\
75 & $\mathit{vplus1} = \mathit{dof} + 1$ & $\mathit{dof} + 1 = \mathit{vplus1}$ \\
76 & $\mathit{L17} = \mathit{dof} \cdot \mathit{vplus1}$ & $\mathit{dof} \cdot \mathit{vplus1} = \mathit{L17}$ \\
77 & $\mathit{Gplus1} = \mathit{ic22} - 1$ & $\mathit{Gplus1} + 1 = \mathit{ic22}$ \\
78 & $\mathit{G2m1} = \mathit{ic22} \cdot \mathit{Gplus1}$ & $\mathit{ic22} \cdot \mathit{Gplus1} = \mathit{G2m1}$ \\
79 & $\mathit{jc} = \mathit{j} \cdot \mathit{c}$ & $\mathit{j} \cdot \mathit{c} = \mathit{jc}$ \\
80 & $\mathit{H17} = \mathit{tr1} + \mathit{jc}$ & $\mathit{tr1} + \mathit{jc} = \mathit{H17}$ \\
81 & $\mathit{H2} = \mathit{H17} \cdot \mathit{H17}$ & $\mathit{H17} \cdot \mathit{H17} = \mathit{H2}$ \\
82 & $\mathit{P17} = \mathit{G2m1} \cdot \mathit{H2}$ & $\mathit{G2m1} \cdot \mathit{H2} = \mathit{P17}$ \\
83 & $\mathit{amb} = \mathit{a} - \mathit{th}$ & $\mathit{amb} + \mathit{th} = \mathit{a}$ \\
84 & $\mathit{kamb} = \mathit{ka} \cdot \mathit{amb}$ & $\mathit{ka} \cdot \mathit{amb} = \mathit{kamb}$ \\
85 & $\mathit{qk} = \mathit{q} + \mathit{kamb}$ & $\mathit{q} + \mathit{kamb} = \mathit{qk}$ \\
86 & $\mathit{amb2} = \mathit{amb} \cdot \mathit{amb}$ & $\mathit{amb} \cdot \mathit{amb} = \mathit{amb2}$ \\
87 & $\mathit{Mod} = \mathit{A} - \mathit{amb2}$ & $\mathit{Mod} + \mathit{amb2} = \mathit{A}$ \\
88 & $\mathit{rM} = \mathit{rho} \cdot \mathit{Mod}$ & $\mathit{rho} \cdot \mathit{Mod} = \mathit{rM}$ \\
89 & $\mathit{R18} = \mathit{qk} + \mathit{rM}$ & $\mathit{qk} + \mathit{rM} = \mathit{R18}$ \\
90 & $\mathit{ka2} = \mathit{ka} \cdot \mathit{ka}$ & $\mathit{ka} \cdot \mathit{ka} = \mathit{ka2}$ \\
91 & $\mathit{Aka2} = \mathit{A} \cdot \mathit{ka2}$ & $\mathit{A} \cdot \mathit{ka2} = \mathit{Aka2}$ \\
92 & $\mathit{R19} = \mathit{Aka2} + 1$ & $\mathit{Aka2} + 1 = \mathit{R19}$ \\
93 & $\mathit{L19} = \mathit{mu} \cdot \mathit{mu}$ & $\mathit{mu} \cdot \mathit{mu} = \mathit{L19}$ \\
94 & $\mathit{Dam1} = \mathit{Delta} \cdot \mathit{a}$ & $\mathit{Delta} \cdot \mathit{a} = \mathit{Dam1}$ \\
95 & $\mathit{R20} = \mathit{Dam1} + \mathit{Tindex}$ & $\mathit{Dam1} + \mathit{Tindex} = \mathit{R20}$ \\
\end{longtable}

The twenty-two free equality tests are
\[\begin{array}{@{}l@{}}
\mathit{L1}=\mathit{q},\quad \mathit{b}=\mathit{bx},\quad \mathit{q2}=\mathit{R2},\\
\mathit{L3}=\mathit{B},\quad \mathit{S2}=\mathit{pack\_rhs},\quad \mathit{n}=\mathit{q8},\\
\mathit{r}=\mathit{R7},\quad \mathit{L9}=\mathit{R9},\quad \mathit{c}=\mathit{R10a},\\
\mathit{k}=\mathit{R10b},\quad \mathit{k}=\mathit{R11},\quad \mathit{a}=\mathit{R12},\\
\mathit{c}=\mathit{R13},\quad \mathit{d}=\mathit{R14},\quad \mathit{L15}=\mathit{R15},\\
\mathit{ic22}=\mathit{R16},\quad \mathit{L17}=\mathit{P17},\quad \mathit{mu}=\mathit{R18},\\
\mathit{L19}=\mathit{R19},\quad \mathit{ka}=\mathit{R20},\quad \mathit{S3}=\mathit{sigma},\\
\mathit{t2}=\mathit{Omega}
\end{array}\]


\clearpage\normalsize
\section{The 94-instruction shifted-index certificate}\label{app:94}

Translate the fixed index of Appendix~\ref{app:95} by replacing
$H_{\rm old}$ with the supplied positive integer $H=H_{\rm old}-4=H_0-3$.
The actual radix $B=H+b+4$ and all 34 positive unknowns are unchanged.
There are still 22 equations and three fixed index parameters besides
the input. The complete positive-witness bijection is proved in
Theorem~\ref{thm:best94}.

The old additions $B_{\rm reg}=H_{\rm old}+b$ and $L_3=\theta+4$
served only the equality $L_3=B_{\rm reg}$. Retain the addition
$\mathit{theta\_sum}=H+b$ and compare it directly with $\theta$.
The mathematical radix is never separately constructed. Every other
instruction and equality test is unchanged, so the complete table
has 51 multiplications and 43 additions.

The standalone checker \file{round33\_1980\_shifted\_affine\_certificate.py}
verifies every primitive and source residual, the complete fixed-parameter
substitution, and the absence of downstream uses of the deleted register.
The geometric and second-exponent corrections use the exact translated
radix equation; the signed correction uses the exact relaxed norm.
The 95-operation coefficient and carry regression applies without change.

This table is generated by \file{round33\_1980\_schedule\_table.py}.
Fixed numerals and equality tests are free.

\small
% Generated by round33_1980_schedule_table.py; do not edit.
\begin{longtable}{@{}rll@{}}
\toprule
No. & Assignment & Addition/multiplication check\\
\midrule
\endfirsthead
\toprule
No. & Assignment & Addition/multiplication check\\
\midrule
\endhead
\bottomrule
\endfoot
1 & $\mathit{q2} = \mathit{q} \cdot \mathit{q}$ & $\mathit{q} \cdot \mathit{q} = \mathit{q2}$ \\
2 & $\mathit{q4} = \mathit{q2} \cdot \mathit{q2}$ & $\mathit{q2} \cdot \mathit{q2} = \mathit{q4}$ \\
3 & $\mathit{q8} = \mathit{q4} \cdot \mathit{q4}$ & $\mathit{q4} \cdot \mathit{q4} = \mathit{q8}$ \\
4 & $\mathit{theta\_sum} = \mathit{H} + \mathit{b}$ & $\mathit{H} + \mathit{b} = \mathit{theta\_sum}$ \\
5 & $\mathit{bx} = \mathit{x} + \mathit{beta}$ & $\mathit{x} + \mathit{beta} = \mathit{bx}$ \\
6 & $\mathit{theta\_lambda} = \mathit{th} \cdot \mathit{la}$ & $\mathit{th} \cdot \mathit{la} = \mathit{theta\_lambda}$ \\
7 & $\mathit{Tcoef} = \mathit{theta\_lambda} + 1$ & $\mathit{theta\_lambda} + 1 = \mathit{Tcoef}$ \\
8 & $\mathit{three\_lambda} = 3 \cdot \mathit{la}$ & $3 \cdot \mathit{la} = \mathit{three\_lambda}$ \\
9 & $\mathit{R2} = \mathit{Tcoef} + \mathit{three\_lambda}$ & $\mathit{Tcoef} + \mathit{three\_lambda} = \mathit{R2}$ \\
10 & $\mathit{t2} = \mathit{la} - \mathit{e}$ & $\mathit{t2} + \mathit{e} = \mathit{la}$ \\
11 & $\mathit{eq} = \mathit{e} \cdot \mathit{q}$ & $\mathit{e} \cdot \mathit{q} = \mathit{eq}$ \\
12 & $\mathit{S2} = \mathit{l} + \mathit{eq}$ & $\mathit{l} + \mathit{eq} = \mathit{S2}$ \\
13 & $\mathit{pack\_tth} = \mathit{t} \cdot \mathit{th}$ & $\mathit{t} \cdot \mathit{th} = \mathit{pack\_tth}$ \\
14 & $\mathit{pack\_rhs} = \mathit{V} + \mathit{pack\_tth}$ & $\mathit{V} + \mathit{pack\_tth} = \mathit{pack\_rhs}$ \\
15 & $\mathit{C} = \mathit{x} + \mathit{g}$ & $\mathit{x} + \mathit{g} = \mathit{C}$ \\
16 & $\mathit{C2} = \mathit{C} \cdot \mathit{C}$ & $\mathit{C} \cdot \mathit{C} = \mathit{C2}$ \\
17 & $\mathit{bound\_sum} = \mathit{l} + \mathit{e}$ & $\mathit{l} + \mathit{e} = \mathit{bound\_sum}$ \\
18 & $\mathit{L1} = \mathit{bound\_sum} + \mathit{al}$ & $\mathit{bound\_sum} + \mathit{al} = \mathit{L1}$ \\
19 & $\mathit{code\_gap} = \mathit{q} - \mathit{C2}$ & $\mathit{code\_gap} + \mathit{C2} = \mathit{q}$ \\
20 & $\mathit{S3} = \mathit{t2} \cdot \mathit{code\_gap}$ & $\mathit{t2} \cdot \mathit{code\_gap} = \mathit{S3}$ \\
21 & $\mathit{S3q4} = \mathit{S3} \cdot \mathit{q2}$ & $\mathit{S3} \cdot \mathit{q2} = \mathit{S3q4}$ \\
22 & $\mathit{Sin} = \mathit{S2} + \mathit{S3q4}$ & $\mathit{S2} + \mathit{S3q4} = \mathit{Sin}$ \\
23 & $\mathit{Sq4} = \mathit{Sin} \cdot \mathit{q2}$ & $\mathit{Sin} \cdot \mathit{q2} = \mathit{Sq4}$ \\
24 & $\mathit{S} = \mathit{g} + \mathit{Sq4}$ & $\mathit{g} + \mathit{Sq4} = \mathit{S}$ \\
25 & $\mathit{n2} = \mathit{n} \cdot \mathit{n}$ & $\mathit{n} \cdot \mathit{n} = \mathit{n2}$ \\
26 & $\mathit{n2n} = \mathit{n2} - \mathit{n}$ & $\mathit{n2n} + \mathit{n} = \mathit{n2}$ \\
27 & $\mathit{SA} = \mathit{S} \cdot \mathit{n2n}$ & $\mathit{S} \cdot \mathit{n2n} = \mathit{SA}$ \\
28 & $\mathit{Tq4} = \mathit{Tcoef} \cdot \mathit{q2}$ & $\mathit{Tcoef} \cdot \mathit{q2} = \mathit{Tq4}$ \\
29 & $\mathit{lbm1} = \mathit{l} \cdot \mathit{b}$ & $\mathit{l} \cdot \mathit{b} = \mathit{lbm1}$ \\
30 & $\mathit{T2} = \mathit{Tq4} - \mathit{lbm1}$ & $\mathit{T2} + \mathit{lbm1} = \mathit{Tq4}$ \\
31 & $\mathit{packed\_target\_shift} = \mathit{l} \cdot \mathit{q4}$ & $\mathit{l} \cdot \mathit{q4} = \mathit{packed\_target\_shift}$ \\
32 & $\mathit{mask\_term} = \mathit{th} \cdot \mathit{packed\_target\_shift}$ & $\mathit{th} \cdot \mathit{packed\_target\_shift} = \mathit{mask\_term}$ \\
33 & $\mathit{T} = \mathit{T2} + \mathit{mask\_term}$ & $\mathit{T2} + \mathit{mask\_term} = \mathit{T}$ \\
34 & $\mathit{n2m1} = \mathit{n2} - 1$ & $\mathit{n2m1} + 1 = \mathit{n2}$ \\
35 & $\mathit{TB} = \mathit{T} \cdot \mathit{n2m1}$ & $\mathit{T} \cdot \mathit{n2m1} = \mathit{TB}$ \\
36 & $\mathit{R7} = \mathit{SA} + \mathit{TB}$ & $\mathit{SA} + \mathit{TB} = \mathit{R7}$ \\
37 & $\mathit{wn2} = \mathit{w} \cdot \mathit{n2}$ & $\mathit{w} \cdot \mathit{n2} = \mathit{wn2}$ \\
38 & $\mathit{sn2} = \mathit{s} \cdot \mathit{n2}$ & $\mathit{s} \cdot \mathit{n2} = \mathit{sn2}$ \\
39 & $\mathit{UM} = \mathit{wn2} \cdot \mathit{sn2}$ & $\mathit{wn2} \cdot \mathit{sn2} = \mathit{UM}$ \\
40 & $\mathit{R12} = \mathit{UM} + \mathit{sn2}$ & $\mathit{UM} + \mathit{sn2} = \mathit{R12}$ \\
41 & $\mathit{ksn2} = \mathit{k} \cdot \mathit{sn2}$ & $\mathit{k} \cdot \mathit{sn2} = \mathit{ksn2}$ \\
42 & $\mathit{UM2} = \mathit{UM} \cdot \mathit{UM}$ & $\mathit{UM} \cdot \mathit{UM} = \mathit{UM2}$ \\
43 & $\mathit{scaled\_norm\_coefficient} = \mathit{UM2} + \mathit{wn2}$ & $\mathit{UM2} + \mathit{wn2} = \mathit{scaled\_norm\_coefficient}$ \\
44 & $\mathit{ratio\_product2} = \mathit{ksn2} \cdot \mathit{ksn2}$ & $\mathit{ksn2} \cdot \mathit{ksn2} = \mathit{ratio\_product2}$ \\
45 & $\mathit{L9} = \mathit{scaled\_norm\_coefficient} \cdot \mathit{ratio\_product2}$ & $\mathit{scaled\_norm\_coefficient} \cdot \mathit{ratio\_product2} = \mathit{L9}$ \\
46 & $\mathit{tauplus1} = \mathit{tau} + 1$ & $\mathit{tau} + 1 = \mathit{tauplus1}$ \\
47 & $\mathit{R9} = \mathit{tau} \cdot \mathit{tauplus1}$ & $\mathit{tau} \cdot \mathit{tauplus1} = \mathit{R9}$ \\
48 & $\mathit{R10a} = \mathit{ksn2} + \mathit{eta}$ & $\mathit{ksn2} + \mathit{eta} = \mathit{R10a}$ \\
49 & $\mathit{R10b} = \mathit{eta} + \mathit{zeta}$ & $\mathit{eta} + \mathit{zeta} = \mathit{R10b}$ \\
50 & $\mathit{r1} = \mathit{r} + 1$ & $\mathit{r} + 1 = \mathit{r1}$ \\
51 & $\mathit{hpm1} = \mathit{h} \cdot \mathit{UM}$ & $\mathit{h} \cdot \mathit{UM} = \mathit{hpm1}$ \\
52 & $\mathit{R11} = \mathit{r1} + \mathit{hpm1}$ & $\mathit{r1} + \mathit{hpm1} = \mathit{R11}$ \\
53 & $\mathit{tr1} = \mathit{r1} + \mathit{r}$ & $\mathit{r1} + \mathit{r} = \mathit{tr1}$ \\
54 & $\mathit{R13} = \mathit{ka} + \mathit{phi}$ & $\mathit{ka} + \mathit{phi} = \mathit{R13}$ \\
55 & $\mathit{cam2} = \mathit{c} \cdot \mathit{a}$ & $\mathit{c} \cdot \mathit{a} = \mathit{cam2}$ \\
56 & $\mathit{D1} = \mathit{wn2} + \mathit{cam2}$ & $\mathit{wn2} + \mathit{cam2} = \mathit{D1}$ \\
57 & $\mathit{a4} = 8 \cdot \mathit{a}$ & $8 \cdot \mathit{a} = \mathit{a4}$ \\
58 & $\mathit{a4m5} = \mathit{a4} + 15$ & $\mathit{a4} + 15 = \mathit{a4m5}$ \\
59 & $\mathit{gam} = \mathit{ga} \cdot \mathit{a4m5}$ & $\mathit{ga} \cdot \mathit{a4m5} = \mathit{gam}$ \\
60 & $\mathit{R14} = \mathit{D1} + \mathit{gam}$ & $\mathit{D1} + \mathit{gam} = \mathit{R14}$ \\
61 & $\mathit{a\_square} = \mathit{a} \cdot \mathit{a}$ & $\mathit{a} \cdot \mathit{a} = \mathit{a\_square}$ \\
62 & $\mathit{A} = \mathit{a\_square} + \mathit{a4m5}$ & $\mathit{a\_square} + \mathit{a4m5} = \mathit{A}$ \\
63 & $\mathit{c2} = \mathit{c} \cdot \mathit{c}$ & $\mathit{c} \cdot \mathit{c} = \mathit{c2}$ \\
64 & $\mathit{Ac2} = \mathit{A} \cdot \mathit{c2}$ & $\mathit{A} \cdot \mathit{c2} = \mathit{Ac2}$ \\
65 & $\mathit{R15} = \mathit{Ac2} + 1$ & $\mathit{Ac2} + 1 = \mathit{R15}$ \\
66 & $\mathit{L15} = \mathit{d} \cdot \mathit{d}$ & $\mathit{d} \cdot \mathit{d} = \mathit{L15}$ \\
67 & $\mathit{ic2} = \mathit{i} \cdot \mathit{c2}$ & $\mathit{i} \cdot \mathit{c2} = \mathit{ic2}$ \\
68 & $\mathit{ic22} = \mathit{ic2} \cdot \mathit{ic2}$ & $\mathit{ic2} \cdot \mathit{ic2} = \mathit{ic22}$ \\
69 & $\mathit{L16} = \mathit{f} \cdot \mathit{f}$ & $\mathit{f} \cdot \mathit{f} = \mathit{L16}$ \\
70 & $\mathit{f\_square\_minus\_one} = \mathit{L16} - 1$ & $\mathit{f\_square\_minus\_one} + 1 = \mathit{L16}$ \\
71 & $\mathit{R16} = \mathit{A} \cdot \mathit{f\_square\_minus\_one}$ & $\mathit{A} \cdot \mathit{f\_square\_minus\_one} = \mathit{R16}$ \\
72 & $\mathit{of} = \mathit{o} \cdot \mathit{f}$ & $\mathit{o} \cdot \mathit{f} = \mathit{of}$ \\
73 & $\mathit{dof} = \mathit{of} - \mathit{L15}$ & $\mathit{dof} + \mathit{L15} = \mathit{of}$ \\
74 & $\mathit{vplus1} = \mathit{dof} + 1$ & $\mathit{dof} + 1 = \mathit{vplus1}$ \\
75 & $\mathit{L17} = \mathit{dof} \cdot \mathit{vplus1}$ & $\mathit{dof} \cdot \mathit{vplus1} = \mathit{L17}$ \\
76 & $\mathit{Gplus1} = \mathit{ic22} - 1$ & $\mathit{Gplus1} + 1 = \mathit{ic22}$ \\
77 & $\mathit{G2m1} = \mathit{ic22} \cdot \mathit{Gplus1}$ & $\mathit{ic22} \cdot \mathit{Gplus1} = \mathit{G2m1}$ \\
78 & $\mathit{jc} = \mathit{j} \cdot \mathit{c}$ & $\mathit{j} \cdot \mathit{c} = \mathit{jc}$ \\
79 & $\mathit{H17} = \mathit{tr1} + \mathit{jc}$ & $\mathit{tr1} + \mathit{jc} = \mathit{H17}$ \\
80 & $\mathit{H2} = \mathit{H17} \cdot \mathit{H17}$ & $\mathit{H17} \cdot \mathit{H17} = \mathit{H2}$ \\
81 & $\mathit{P17} = \mathit{G2m1} \cdot \mathit{H2}$ & $\mathit{G2m1} \cdot \mathit{H2} = \mathit{P17}$ \\
82 & $\mathit{amb} = \mathit{a} - \mathit{th}$ & $\mathit{amb} + \mathit{th} = \mathit{a}$ \\
83 & $\mathit{kamb} = \mathit{ka} \cdot \mathit{amb}$ & $\mathit{ka} \cdot \mathit{amb} = \mathit{kamb}$ \\
84 & $\mathit{qk} = \mathit{q} + \mathit{kamb}$ & $\mathit{q} + \mathit{kamb} = \mathit{qk}$ \\
85 & $\mathit{amb2} = \mathit{amb} \cdot \mathit{amb}$ & $\mathit{amb} \cdot \mathit{amb} = \mathit{amb2}$ \\
86 & $\mathit{Mod} = \mathit{A} - \mathit{amb2}$ & $\mathit{Mod} + \mathit{amb2} = \mathit{A}$ \\
87 & $\mathit{rM} = \mathit{rho} \cdot \mathit{Mod}$ & $\mathit{rho} \cdot \mathit{Mod} = \mathit{rM}$ \\
88 & $\mathit{R18} = \mathit{qk} + \mathit{rM}$ & $\mathit{qk} + \mathit{rM} = \mathit{R18}$ \\
89 & $\mathit{ka2} = \mathit{ka} \cdot \mathit{ka}$ & $\mathit{ka} \cdot \mathit{ka} = \mathit{ka2}$ \\
90 & $\mathit{Aka2} = \mathit{A} \cdot \mathit{ka2}$ & $\mathit{A} \cdot \mathit{ka2} = \mathit{Aka2}$ \\
91 & $\mathit{R19} = \mathit{Aka2} + 1$ & $\mathit{Aka2} + 1 = \mathit{R19}$ \\
92 & $\mathit{L19} = \mathit{mu} \cdot \mathit{mu}$ & $\mathit{mu} \cdot \mathit{mu} = \mathit{L19}$ \\
93 & $\mathit{Dam1} = \mathit{Delta} \cdot \mathit{a}$ & $\mathit{Delta} \cdot \mathit{a} = \mathit{Dam1}$ \\
94 & $\mathit{R20} = \mathit{Dam1} + \mathit{Tindex}$ & $\mathit{Dam1} + \mathit{Tindex} = \mathit{R20}$ \\
\end{longtable}

The twenty-two free equality tests are
\[\begin{array}{@{}l@{}}
\mathit{L1}=\mathit{q},\quad \mathit{b}=\mathit{bx},\quad \mathit{q2}=\mathit{R2},\\
\mathit{th}=\mathit{theta\_sum},\quad \mathit{S2}=\mathit{pack\_rhs},\quad \mathit{n}=\mathit{q8},\\
\mathit{r}=\mathit{R7},\quad \mathit{L9}=\mathit{R9},\quad \mathit{c}=\mathit{R10a},\\
\mathit{k}=\mathit{R10b},\quad \mathit{k}=\mathit{R11},\quad \mathit{a}=\mathit{R12},\\
\mathit{c}=\mathit{R13},\quad \mathit{d}=\mathit{R14},\quad \mathit{L15}=\mathit{R15},\\
\mathit{ic22}=\mathit{R16},\quad \mathit{L17}=\mathit{P17},\quad \mathit{mu}=\mathit{R18},\\
\mathit{L19}=\mathit{R19},\quad \mathit{ka}=\mathit{R20},\quad \mathit{S3}=\mathit{sigma},\\
\mathit{t2}=\mathit{Omega}
\end{array}\]


\clearpage\normalsize
\section{The 93-instruction factored-mask certificate}\label{app:93}

Use exactly the system and fixed index of Appendix~\ref{app:94}.
Theorem~\ref{thm:best93} replaces its six-instruction packed-mask
block by five instructions, using
\[
 T_+=q^2T_{\rm coef}+\ell(\theta q^4-b).
\]
All other registers and all 22 equality tests have their old values.
There are 50 multiplications and 43 additions, with the same 34
positive unknowns and the same three fixed index parameters besides
the input. The positive-witness map is the identity.

The complete checker \file{round34\_1980\_factored\_mask\_certificate.py}
verifies every primitive, all 22 unchanged source residuals and their
acyclic corrections, equality of every retained register, and the
absence of outside uses for the deleted temporaries. The table is
generated by \file{round34\_1980\_schedule\_table.py}. Its literal
numerals are exactly $1,3,8,15$; fixed numerals and equality tests
are free.

\small
\input{jones1980_theorem5_93_schedule.tex}

\clearpage\normalsize
\section{The 92-instruction smaller-parameter Pell certificate}\label{app:92}

Use the fixed index and coefficient encoding of Appendix~\ref{app:93}.
Theorem~\ref{thm:best92} replaces the signed auxiliary norm with
\[
 u=c+of,\qquad K(u^2-y_{\rm aux}^2)=1-y_{\rm aux}^2,
 \qquad u=2r+1+jc,\quad K=(ic^2)^2.
\]
The computed $u$ is named \texttt{H17} in the schedule. The new
supplied unknown $y_{\rm aux}$ is positive; the auxiliary witnesses
$o,j$ are chosen anew in necessity. Every fixed parameter and all
other canonical witnesses are unchanged. The full proof establishes
the required index before applying the retained exponent and decoding
arguments.

There are 35 positive unknowns, 23 equations, 49 multiplications and
43 additions. The differences $u^2-y_{\rm aux}^2$ and
$1-y_{\rm aux}^2$ are negative integer registers on solutions; their
subtractions are certified by reversed additions.

The complete checker
\file{round35\_1980\_half\_parameter\_pell\_certificate.py}
verifies every primitive, all 23 source residuals, and the retained
geometric and second-exponent corrections. The new auxiliary
correction is $F_{16}(u^2-y_{\rm aux}^2)$, where $F_{16}$ is the
independently exact relaxed norm. The new linear congruence is exact.
The separate focused polynomial, modular and positive-witness
regression corroborates the general proof. The table is generated by
\file{round35\_1980\_schedule\_table.py}. Its literal numerals
remain $1,3,8,15$; fixed numerals and equality tests are free.

\small
\input{jones1980_theorem5_92_schedule.tex}

\clearpage\normalsize
\section{The 91-instruction product-bound certificate}\label{app:91}

Use the coefficient index of Theorem~\ref{thm:best91} and the
half-parameter Pell equations of Appendix~\ref{app:92}. The two changed
coding equations are
\[
 \sigma=(e-\ell)C^2,\qquad \ell+\sigma+\alpha=q,
 \qquad C=x+g.
\]
The computed difference $e-\ell$ is positive because $\sigma>0$ and
$C>0$. No supplied $\Omega$ or separate equality for it is needed.
The new fixed index uses two encoded copies of the input to make all
extra negative-coefficient tests exact. The full proof verifies the
carry bounds, both directions of universality, and the parity needed
by the retained auxiliary Pell construction.

There are 34 positive unknowns, 22 equations, 49 multiplications and
42 additions. Relative to Appendix~\ref{app:92}, compute
\texttt{t2} as $e-\ell$, compute \texttt{bound\_sum} as
$\ell+\sigma$, and multiply \texttt{t2} by the existing $C^2$ to
obtain \texttt{S3}. Delete the subtraction $q-C^2$. The free equality
\texttt{t2}=$\Omega$ disappears with its supplied unknown.

The complete checker
\file{round36\_1980\_product\_bound\_certificate.py}
verifies every primitive and all 22 fresh source residuals, including
the inherited geometric, auxiliary-norm and second-exponent
corrections. A separate sparse compiler regression corroborates the
general proof. The table is generated by
\file{round36\_1980\_schedule\_table.py}. Literal numerals remain
$1,3,8,15$; fixed numerals and equality tests are free.

\small
% Generated by round36_1980_schedule_table.py; do not edit.
\begin{longtable}{@{}rll@{}}
\toprule
No. & Assignment & Addition/multiplication check\\
\midrule
\endfirsthead
\toprule
No. & Assignment & Addition/multiplication check\\
\midrule
\endhead
\bottomrule
\endfoot
1 & $\mathit{q2} = \mathit{q} \cdot \mathit{q}$ & $\mathit{q} \cdot \mathit{q} = \mathit{q2}$ \\
2 & $\mathit{q4} = \mathit{q2} \cdot \mathit{q2}$ & $\mathit{q2} \cdot \mathit{q2} = \mathit{q4}$ \\
3 & $\mathit{q8} = \mathit{q4} \cdot \mathit{q4}$ & $\mathit{q4} \cdot \mathit{q4} = \mathit{q8}$ \\
4 & $\mathit{theta\_sum} = \mathit{H} + \mathit{b}$ & $\mathit{H} + \mathit{b} = \mathit{theta\_sum}$ \\
5 & $\mathit{bx} = \mathit{x} + \mathit{beta}$ & $\mathit{x} + \mathit{beta} = \mathit{bx}$ \\
6 & $\mathit{theta\_lambda} = \mathit{th} \cdot \mathit{la}$ & $\mathit{th} \cdot \mathit{la} = \mathit{theta\_lambda}$ \\
7 & $\mathit{Tcoef} = \mathit{theta\_lambda} + 1$ & $\mathit{theta\_lambda} + 1 = \mathit{Tcoef}$ \\
8 & $\mathit{three\_lambda} = 3 \cdot \mathit{la}$ & $3 \cdot \mathit{la} = \mathit{three\_lambda}$ \\
9 & $\mathit{R2} = \mathit{Tcoef} + \mathit{three\_lambda}$ & $\mathit{Tcoef} + \mathit{three\_lambda} = \mathit{R2}$ \\
10 & $\mathit{t2} = \mathit{e} - \mathit{l}$ & $\mathit{t2} + \mathit{l} = \mathit{e}$ \\
11 & $\mathit{eq} = \mathit{e} \cdot \mathit{q}$ & $\mathit{e} \cdot \mathit{q} = \mathit{eq}$ \\
12 & $\mathit{S2} = \mathit{l} + \mathit{eq}$ & $\mathit{l} + \mathit{eq} = \mathit{S2}$ \\
13 & $\mathit{pack\_tth} = \mathit{t} \cdot \mathit{th}$ & $\mathit{t} \cdot \mathit{th} = \mathit{pack\_tth}$ \\
14 & $\mathit{pack\_rhs} = \mathit{V} + \mathit{pack\_tth}$ & $\mathit{V} + \mathit{pack\_tth} = \mathit{pack\_rhs}$ \\
15 & $\mathit{C} = \mathit{x} + \mathit{g}$ & $\mathit{x} + \mathit{g} = \mathit{C}$ \\
16 & $\mathit{C2} = \mathit{C} \cdot \mathit{C}$ & $\mathit{C} \cdot \mathit{C} = \mathit{C2}$ \\
17 & $\mathit{bound\_sum} = \mathit{l} + \mathit{sigma}$ & $\mathit{l} + \mathit{sigma} = \mathit{bound\_sum}$ \\
18 & $\mathit{L1} = \mathit{bound\_sum} + \mathit{al}$ & $\mathit{bound\_sum} + \mathit{al} = \mathit{L1}$ \\
19 & $\mathit{S3} = \mathit{t2} \cdot \mathit{C2}$ & $\mathit{t2} \cdot \mathit{C2} = \mathit{S3}$ \\
20 & $\mathit{S3q4} = \mathit{S3} \cdot \mathit{q2}$ & $\mathit{S3} \cdot \mathit{q2} = \mathit{S3q4}$ \\
21 & $\mathit{Sin} = \mathit{S2} + \mathit{S3q4}$ & $\mathit{S2} + \mathit{S3q4} = \mathit{Sin}$ \\
22 & $\mathit{Sq4} = \mathit{Sin} \cdot \mathit{q2}$ & $\mathit{Sin} \cdot \mathit{q2} = \mathit{Sq4}$ \\
23 & $\mathit{S} = \mathit{g} + \mathit{Sq4}$ & $\mathit{g} + \mathit{Sq4} = \mathit{S}$ \\
24 & $\mathit{n2} = \mathit{n} \cdot \mathit{n}$ & $\mathit{n} \cdot \mathit{n} = \mathit{n2}$ \\
25 & $\mathit{n2n} = \mathit{n2} - \mathit{n}$ & $\mathit{n2n} + \mathit{n} = \mathit{n2}$ \\
26 & $\mathit{SA} = \mathit{S} \cdot \mathit{n2n}$ & $\mathit{S} \cdot \mathit{n2n} = \mathit{SA}$ \\
27 & $\mathit{Tq4} = \mathit{Tcoef} \cdot \mathit{q2}$ & $\mathit{Tcoef} \cdot \mathit{q2} = \mathit{Tq4}$ \\
28 & $\mathit{theta\_q4} = \mathit{th} \cdot \mathit{q4}$ & $\mathit{th} \cdot \mathit{q4} = \mathit{theta\_q4}$ \\
29 & $\mathit{mask\_gap} = \mathit{theta\_q4} - \mathit{b}$ & $\mathit{mask\_gap} + \mathit{b} = \mathit{theta\_q4}$ \\
30 & $\mathit{indicator\_mask} = \mathit{l} \cdot \mathit{mask\_gap}$ & $\mathit{l} \cdot \mathit{mask\_gap} = \mathit{indicator\_mask}$ \\
31 & $\mathit{T} = \mathit{Tq4} + \mathit{indicator\_mask}$ & $\mathit{Tq4} + \mathit{indicator\_mask} = \mathit{T}$ \\
32 & $\mathit{n2m1} = \mathit{n2} - 1$ & $\mathit{n2m1} + 1 = \mathit{n2}$ \\
33 & $\mathit{TB} = \mathit{T} \cdot \mathit{n2m1}$ & $\mathit{T} \cdot \mathit{n2m1} = \mathit{TB}$ \\
34 & $\mathit{R7} = \mathit{SA} + \mathit{TB}$ & $\mathit{SA} + \mathit{TB} = \mathit{R7}$ \\
35 & $\mathit{wn2} = \mathit{w} \cdot \mathit{n2}$ & $\mathit{w} \cdot \mathit{n2} = \mathit{wn2}$ \\
36 & $\mathit{sn2} = \mathit{s} \cdot \mathit{n2}$ & $\mathit{s} \cdot \mathit{n2} = \mathit{sn2}$ \\
37 & $\mathit{UM} = \mathit{wn2} \cdot \mathit{sn2}$ & $\mathit{wn2} \cdot \mathit{sn2} = \mathit{UM}$ \\
38 & $\mathit{R12} = \mathit{UM} + \mathit{sn2}$ & $\mathit{UM} + \mathit{sn2} = \mathit{R12}$ \\
39 & $\mathit{ksn2} = \mathit{k} \cdot \mathit{sn2}$ & $\mathit{k} \cdot \mathit{sn2} = \mathit{ksn2}$ \\
40 & $\mathit{UM2} = \mathit{UM} \cdot \mathit{UM}$ & $\mathit{UM} \cdot \mathit{UM} = \mathit{UM2}$ \\
41 & $\mathit{scaled\_norm\_coefficient} = \mathit{UM2} + \mathit{wn2}$ & $\mathit{UM2} + \mathit{wn2} = \mathit{scaled\_norm\_coefficient}$ \\
42 & $\mathit{ratio\_product2} = \mathit{ksn2} \cdot \mathit{ksn2}$ & $\mathit{ksn2} \cdot \mathit{ksn2} = \mathit{ratio\_product2}$ \\
43 & $\mathit{L9} = \mathit{scaled\_norm\_coefficient} \cdot \mathit{ratio\_product2}$ & $\mathit{scaled\_norm\_coefficient} \cdot \mathit{ratio\_product2} = \mathit{L9}$ \\
44 & $\mathit{tauplus1} = \mathit{tau} + 1$ & $\mathit{tau} + 1 = \mathit{tauplus1}$ \\
45 & $\mathit{R9} = \mathit{tau} \cdot \mathit{tauplus1}$ & $\mathit{tau} \cdot \mathit{tauplus1} = \mathit{R9}$ \\
46 & $\mathit{R10a} = \mathit{ksn2} + \mathit{eta}$ & $\mathit{ksn2} + \mathit{eta} = \mathit{R10a}$ \\
47 & $\mathit{R10b} = \mathit{eta} + \mathit{zeta}$ & $\mathit{eta} + \mathit{zeta} = \mathit{R10b}$ \\
48 & $\mathit{r1} = \mathit{r} + 1$ & $\mathit{r} + 1 = \mathit{r1}$ \\
49 & $\mathit{hpm1} = \mathit{h} \cdot \mathit{UM}$ & $\mathit{h} \cdot \mathit{UM} = \mathit{hpm1}$ \\
50 & $\mathit{R11} = \mathit{r1} + \mathit{hpm1}$ & $\mathit{r1} + \mathit{hpm1} = \mathit{R11}$ \\
51 & $\mathit{tr1} = \mathit{r1} + \mathit{r}$ & $\mathit{r1} + \mathit{r} = \mathit{tr1}$ \\
52 & $\mathit{R13} = \mathit{ka} + \mathit{phi}$ & $\mathit{ka} + \mathit{phi} = \mathit{R13}$ \\
53 & $\mathit{cam2} = \mathit{c} \cdot \mathit{a}$ & $\mathit{c} \cdot \mathit{a} = \mathit{cam2}$ \\
54 & $\mathit{D1} = \mathit{wn2} + \mathit{cam2}$ & $\mathit{wn2} + \mathit{cam2} = \mathit{D1}$ \\
55 & $\mathit{a4} = 8 \cdot \mathit{a}$ & $8 \cdot \mathit{a} = \mathit{a4}$ \\
56 & $\mathit{a4m5} = \mathit{a4} + 15$ & $\mathit{a4} + 15 = \mathit{a4m5}$ \\
57 & $\mathit{gam} = \mathit{ga} \cdot \mathit{a4m5}$ & $\mathit{ga} \cdot \mathit{a4m5} = \mathit{gam}$ \\
58 & $\mathit{R14} = \mathit{D1} + \mathit{gam}$ & $\mathit{D1} + \mathit{gam} = \mathit{R14}$ \\
59 & $\mathit{a\_square} = \mathit{a} \cdot \mathit{a}$ & $\mathit{a} \cdot \mathit{a} = \mathit{a\_square}$ \\
60 & $\mathit{A} = \mathit{a\_square} + \mathit{a4m5}$ & $\mathit{a\_square} + \mathit{a4m5} = \mathit{A}$ \\
61 & $\mathit{c2} = \mathit{c} \cdot \mathit{c}$ & $\mathit{c} \cdot \mathit{c} = \mathit{c2}$ \\
62 & $\mathit{Ac2} = \mathit{A} \cdot \mathit{c2}$ & $\mathit{A} \cdot \mathit{c2} = \mathit{Ac2}$ \\
63 & $\mathit{R15} = \mathit{Ac2} + 1$ & $\mathit{Ac2} + 1 = \mathit{R15}$ \\
64 & $\mathit{L15} = \mathit{d} \cdot \mathit{d}$ & $\mathit{d} \cdot \mathit{d} = \mathit{L15}$ \\
65 & $\mathit{ic2} = \mathit{i} \cdot \mathit{c2}$ & $\mathit{i} \cdot \mathit{c2} = \mathit{ic2}$ \\
66 & $\mathit{ic22} = \mathit{ic2} \cdot \mathit{ic2}$ & $\mathit{ic2} \cdot \mathit{ic2} = \mathit{ic22}$ \\
67 & $\mathit{L16} = \mathit{f} \cdot \mathit{f}$ & $\mathit{f} \cdot \mathit{f} = \mathit{L16}$ \\
68 & $\mathit{f\_square\_minus\_one} = \mathit{L16} - 1$ & $\mathit{f\_square\_minus\_one} + 1 = \mathit{L16}$ \\
69 & $\mathit{R16} = \mathit{A} \cdot \mathit{f\_square\_minus\_one}$ & $\mathit{A} \cdot \mathit{f\_square\_minus\_one} = \mathit{R16}$ \\
70 & $\mathit{of} = \mathit{o} \cdot \mathit{f}$ & $\mathit{o} \cdot \mathit{f} = \mathit{of}$ \\
71 & $\mathit{aux\_u\_rhs} = \mathit{c} + \mathit{of}$ & $\mathit{c} + \mathit{of} = \mathit{aux\_u\_rhs}$ \\
72 & $\mathit{jc} = \mathit{j} \cdot \mathit{c}$ & $\mathit{j} \cdot \mathit{c} = \mathit{jc}$ \\
73 & $\mathit{H17} = \mathit{tr1} + \mathit{jc}$ & $\mathit{tr1} + \mathit{jc} = \mathit{H17}$ \\
74 & $\mathit{H2} = \mathit{H17} \cdot \mathit{H17}$ & $\mathit{H17} \cdot \mathit{H17} = \mathit{H2}$ \\
75 & $\mathit{aux\_y2} = \mathit{y\_aux} \cdot \mathit{y\_aux}$ & $\mathit{y\_aux} \cdot \mathit{y\_aux} = \mathit{aux\_y2}$ \\
76 & $\mathit{aux\_square\_gap} = \mathit{H2} - \mathit{aux\_y2}$ & $\mathit{aux\_square\_gap} + \mathit{aux\_y2} = \mathit{H2}$ \\
77 & $\mathit{L17} = \mathit{ic22} \cdot \mathit{aux\_square\_gap}$ & $\mathit{ic22} \cdot \mathit{aux\_square\_gap} = \mathit{L17}$ \\
78 & $\mathit{P17} = 1 - \mathit{aux\_y2}$ & $\mathit{P17} + \mathit{aux\_y2} = 1$ \\
79 & $\mathit{amb} = \mathit{a} - \mathit{th}$ & $\mathit{amb} + \mathit{th} = \mathit{a}$ \\
80 & $\mathit{kamb} = \mathit{ka} \cdot \mathit{amb}$ & $\mathit{ka} \cdot \mathit{amb} = \mathit{kamb}$ \\
81 & $\mathit{qk} = \mathit{q} + \mathit{kamb}$ & $\mathit{q} + \mathit{kamb} = \mathit{qk}$ \\
82 & $\mathit{amb2} = \mathit{amb} \cdot \mathit{amb}$ & $\mathit{amb} \cdot \mathit{amb} = \mathit{amb2}$ \\
83 & $\mathit{Mod} = \mathit{A} - \mathit{amb2}$ & $\mathit{Mod} + \mathit{amb2} = \mathit{A}$ \\
84 & $\mathit{rM} = \mathit{rho} \cdot \mathit{Mod}$ & $\mathit{rho} \cdot \mathit{Mod} = \mathit{rM}$ \\
85 & $\mathit{R18} = \mathit{qk} + \mathit{rM}$ & $\mathit{qk} + \mathit{rM} = \mathit{R18}$ \\
86 & $\mathit{ka2} = \mathit{ka} \cdot \mathit{ka}$ & $\mathit{ka} \cdot \mathit{ka} = \mathit{ka2}$ \\
87 & $\mathit{Aka2} = \mathit{A} \cdot \mathit{ka2}$ & $\mathit{A} \cdot \mathit{ka2} = \mathit{Aka2}$ \\
88 & $\mathit{R19} = \mathit{Aka2} + 1$ & $\mathit{Aka2} + 1 = \mathit{R19}$ \\
89 & $\mathit{L19} = \mathit{mu} \cdot \mathit{mu}$ & $\mathit{mu} \cdot \mathit{mu} = \mathit{L19}$ \\
90 & $\mathit{Dam1} = \mathit{Delta} \cdot \mathit{a}$ & $\mathit{Delta} \cdot \mathit{a} = \mathit{Dam1}$ \\
91 & $\mathit{R20} = \mathit{Dam1} + \mathit{Tindex}$ & $\mathit{Dam1} + \mathit{Tindex} = \mathit{R20}$ \\
\end{longtable}

The twenty-two free equality tests are
\[\begin{array}{@{}l@{}}
\mathit{L1}=\mathit{q},\quad \mathit{b}=\mathit{bx},\quad \mathit{q2}=\mathit{R2},\\
\mathit{th}=\mathit{theta\_sum},\quad \mathit{S2}=\mathit{pack\_rhs},\quad \mathit{n}=\mathit{q8},\\
\mathit{r}=\mathit{R7},\quad \mathit{L9}=\mathit{R9},\quad \mathit{c}=\mathit{R10a},\\
\mathit{k}=\mathit{R10b},\quad \mathit{k}=\mathit{R11},\quad \mathit{a}=\mathit{R12},\\
\mathit{c}=\mathit{R13},\quad \mathit{d}=\mathit{R14},\quad \mathit{L15}=\mathit{R15},\\
\mathit{ic22}=\mathit{R16},\quad \mathit{L17}=\mathit{P17},\quad \mathit{mu}=\mathit{R18},\\
\mathit{L19}=\mathit{R19},\quad \mathit{ka}=\mathit{R20},\quad \mathit{S3}=\mathit{sigma},\\
\mathit{H17}=\mathit{aux\_u\_rhs}
\end{array}\]


\clearpage\normalsize
\section{The 90-instruction binary product-bound certificate}\label{app:90}

Use the binary coefficient index of Theorem~\ref{thm:best90}:
$H=H_0-1$, $B=H_0+1+b=H+b+2$, and
$(V,H,T_L)=(\ell_0(2)+2^Le_0(2),H,\psi_2(L))$.
The positive unknowns are those of Appendix~\ref{app:91}, with
the changed Pell base $A=a+2$. There are 34 positive unknowns,
22 equations, 48 multiplications and 42 additions.

The product and bound remain $\sigma=(e-\ell)(x+g)^2$ and
$\ell+\sigma+\alpha=q$. The new reverse helper rows and two unit
tests make the fixed codes binary, permitting $\theta=B-2$.
The geometric identity becomes $q^2=T_{\rm coef}+\lambda$ and
the instruction $3\lambda$ disappears. The first exponent modulus
is $4a+3$ and the discriminant is $a^2+4a+3$; they share the same
three instructions as their predecessors. The complete proof in
Section~\ref{sec:binary-product} includes the changed Pell bootstrap,
exponent and index recovery, exact rounding, and all positive witnesses.

The complete checker
\file{round37\_1980\_binary\_product\_certificate.py}
verifies all 90 primitives and all 22 fresh source residuals, including
the inherited triangular corrections. The table is generated by
\file{round37\_1980\_schedule\_table.py}. Literal numerals are exactly
$1,3,4$; fixed numerals and equality tests are free.

\small
% Generated by round37_1980_schedule_table.py; do not edit.
\begin{longtable}{@{}rll@{}}
\toprule
No. & Assignment & Addition/multiplication check\\
\midrule
\endfirsthead
\toprule
No. & Assignment & Addition/multiplication check\\
\midrule
\endhead
\bottomrule
\endfoot
1 & $\mathit{q2} = \mathit{q} \cdot \mathit{q}$ & $\mathit{q} \cdot \mathit{q} = \mathit{q2}$ \\
2 & $\mathit{q4} = \mathit{q2} \cdot \mathit{q2}$ & $\mathit{q2} \cdot \mathit{q2} = \mathit{q4}$ \\
3 & $\mathit{q8} = \mathit{q4} \cdot \mathit{q4}$ & $\mathit{q4} \cdot \mathit{q4} = \mathit{q8}$ \\
4 & $\mathit{theta\_sum} = \mathit{H} + \mathit{b}$ & $\mathit{H} + \mathit{b} = \mathit{theta\_sum}$ \\
5 & $\mathit{bx} = \mathit{x} + \mathit{beta}$ & $\mathit{x} + \mathit{beta} = \mathit{bx}$ \\
6 & $\mathit{theta\_lambda} = \mathit{th} \cdot \mathit{la}$ & $\mathit{th} \cdot \mathit{la} = \mathit{theta\_lambda}$ \\
7 & $\mathit{Tcoef} = \mathit{theta\_lambda} + 1$ & $\mathit{theta\_lambda} + 1 = \mathit{Tcoef}$ \\
8 & $\mathit{R2} = \mathit{Tcoef} + \mathit{la}$ & $\mathit{Tcoef} + \mathit{la} = \mathit{R2}$ \\
9 & $\mathit{t2} = \mathit{e} - \mathit{l}$ & $\mathit{t2} + \mathit{l} = \mathit{e}$ \\
10 & $\mathit{eq} = \mathit{e} \cdot \mathit{q}$ & $\mathit{e} \cdot \mathit{q} = \mathit{eq}$ \\
11 & $\mathit{S2} = \mathit{l} + \mathit{eq}$ & $\mathit{l} + \mathit{eq} = \mathit{S2}$ \\
12 & $\mathit{pack\_tth} = \mathit{t} \cdot \mathit{th}$ & $\mathit{t} \cdot \mathit{th} = \mathit{pack\_tth}$ \\
13 & $\mathit{pack\_rhs} = \mathit{V} + \mathit{pack\_tth}$ & $\mathit{V} + \mathit{pack\_tth} = \mathit{pack\_rhs}$ \\
14 & $\mathit{C} = \mathit{x} + \mathit{g}$ & $\mathit{x} + \mathit{g} = \mathit{C}$ \\
15 & $\mathit{C2} = \mathit{C} \cdot \mathit{C}$ & $\mathit{C} \cdot \mathit{C} = \mathit{C2}$ \\
16 & $\mathit{bound\_sum} = \mathit{l} + \mathit{sigma}$ & $\mathit{l} + \mathit{sigma} = \mathit{bound\_sum}$ \\
17 & $\mathit{L1} = \mathit{bound\_sum} + \mathit{al}$ & $\mathit{bound\_sum} + \mathit{al} = \mathit{L1}$ \\
18 & $\mathit{S3} = \mathit{t2} \cdot \mathit{C2}$ & $\mathit{t2} \cdot \mathit{C2} = \mathit{S3}$ \\
19 & $\mathit{S3q4} = \mathit{S3} \cdot \mathit{q2}$ & $\mathit{S3} \cdot \mathit{q2} = \mathit{S3q4}$ \\
20 & $\mathit{Sin} = \mathit{S2} + \mathit{S3q4}$ & $\mathit{S2} + \mathit{S3q4} = \mathit{Sin}$ \\
21 & $\mathit{Sq4} = \mathit{Sin} \cdot \mathit{q2}$ & $\mathit{Sin} \cdot \mathit{q2} = \mathit{Sq4}$ \\
22 & $\mathit{S} = \mathit{g} + \mathit{Sq4}$ & $\mathit{g} + \mathit{Sq4} = \mathit{S}$ \\
23 & $\mathit{n2} = \mathit{n} \cdot \mathit{n}$ & $\mathit{n} \cdot \mathit{n} = \mathit{n2}$ \\
24 & $\mathit{n2n} = \mathit{n2} - \mathit{n}$ & $\mathit{n2n} + \mathit{n} = \mathit{n2}$ \\
25 & $\mathit{SA} = \mathit{S} \cdot \mathit{n2n}$ & $\mathit{S} \cdot \mathit{n2n} = \mathit{SA}$ \\
26 & $\mathit{Tq4} = \mathit{Tcoef} \cdot \mathit{q2}$ & $\mathit{Tcoef} \cdot \mathit{q2} = \mathit{Tq4}$ \\
27 & $\mathit{theta\_q4} = \mathit{th} \cdot \mathit{q4}$ & $\mathit{th} \cdot \mathit{q4} = \mathit{theta\_q4}$ \\
28 & $\mathit{mask\_gap} = \mathit{theta\_q4} - \mathit{b}$ & $\mathit{mask\_gap} + \mathit{b} = \mathit{theta\_q4}$ \\
29 & $\mathit{indicator\_mask} = \mathit{l} \cdot \mathit{mask\_gap}$ & $\mathit{l} \cdot \mathit{mask\_gap} = \mathit{indicator\_mask}$ \\
30 & $\mathit{T} = \mathit{Tq4} + \mathit{indicator\_mask}$ & $\mathit{Tq4} + \mathit{indicator\_mask} = \mathit{T}$ \\
31 & $\mathit{n2m1} = \mathit{n2} - 1$ & $\mathit{n2m1} + 1 = \mathit{n2}$ \\
32 & $\mathit{TB} = \mathit{T} \cdot \mathit{n2m1}$ & $\mathit{T} \cdot \mathit{n2m1} = \mathit{TB}$ \\
33 & $\mathit{R7} = \mathit{SA} + \mathit{TB}$ & $\mathit{SA} + \mathit{TB} = \mathit{R7}$ \\
34 & $\mathit{wn2} = \mathit{w} \cdot \mathit{n2}$ & $\mathit{w} \cdot \mathit{n2} = \mathit{wn2}$ \\
35 & $\mathit{sn2} = \mathit{s} \cdot \mathit{n2}$ & $\mathit{s} \cdot \mathit{n2} = \mathit{sn2}$ \\
36 & $\mathit{UM} = \mathit{wn2} \cdot \mathit{sn2}$ & $\mathit{wn2} \cdot \mathit{sn2} = \mathit{UM}$ \\
37 & $\mathit{R12} = \mathit{UM} + \mathit{sn2}$ & $\mathit{UM} + \mathit{sn2} = \mathit{R12}$ \\
38 & $\mathit{ksn2} = \mathit{k} \cdot \mathit{sn2}$ & $\mathit{k} \cdot \mathit{sn2} = \mathit{ksn2}$ \\
39 & $\mathit{UM2} = \mathit{UM} \cdot \mathit{UM}$ & $\mathit{UM} \cdot \mathit{UM} = \mathit{UM2}$ \\
40 & $\mathit{scaled\_norm\_coefficient} = \mathit{UM2} + \mathit{wn2}$ & $\mathit{UM2} + \mathit{wn2} = \mathit{scaled\_norm\_coefficient}$ \\
41 & $\mathit{ratio\_product2} = \mathit{ksn2} \cdot \mathit{ksn2}$ & $\mathit{ksn2} \cdot \mathit{ksn2} = \mathit{ratio\_product2}$ \\
42 & $\mathit{L9} = \mathit{scaled\_norm\_coefficient} \cdot \mathit{ratio\_product2}$ & $\mathit{scaled\_norm\_coefficient} \cdot \mathit{ratio\_product2} = \mathit{L9}$ \\
43 & $\mathit{tauplus1} = \mathit{tau} + 1$ & $\mathit{tau} + 1 = \mathit{tauplus1}$ \\
44 & $\mathit{R9} = \mathit{tau} \cdot \mathit{tauplus1}$ & $\mathit{tau} \cdot \mathit{tauplus1} = \mathit{R9}$ \\
45 & $\mathit{R10a} = \mathit{ksn2} + \mathit{eta}$ & $\mathit{ksn2} + \mathit{eta} = \mathit{R10a}$ \\
46 & $\mathit{R10b} = \mathit{eta} + \mathit{zeta}$ & $\mathit{eta} + \mathit{zeta} = \mathit{R10b}$ \\
47 & $\mathit{r1} = \mathit{r} + 1$ & $\mathit{r} + 1 = \mathit{r1}$ \\
48 & $\mathit{hpm1} = \mathit{h} \cdot \mathit{UM}$ & $\mathit{h} \cdot \mathit{UM} = \mathit{hpm1}$ \\
49 & $\mathit{R11} = \mathit{r1} + \mathit{hpm1}$ & $\mathit{r1} + \mathit{hpm1} = \mathit{R11}$ \\
50 & $\mathit{tr1} = \mathit{r1} + \mathit{r}$ & $\mathit{r1} + \mathit{r} = \mathit{tr1}$ \\
51 & $\mathit{R13} = \mathit{ka} + \mathit{phi}$ & $\mathit{ka} + \mathit{phi} = \mathit{R13}$ \\
52 & $\mathit{cam2} = \mathit{c} \cdot \mathit{a}$ & $\mathit{c} \cdot \mathit{a} = \mathit{cam2}$ \\
53 & $\mathit{D1} = \mathit{wn2} + \mathit{cam2}$ & $\mathit{wn2} + \mathit{cam2} = \mathit{D1}$ \\
54 & $\mathit{a4} = 4 \cdot \mathit{a}$ & $4 \cdot \mathit{a} = \mathit{a4}$ \\
55 & $\mathit{a4m5} = \mathit{a4} + 3$ & $\mathit{a4} + 3 = \mathit{a4m5}$ \\
56 & $\mathit{gam} = \mathit{ga} \cdot \mathit{a4m5}$ & $\mathit{ga} \cdot \mathit{a4m5} = \mathit{gam}$ \\
57 & $\mathit{R14} = \mathit{D1} + \mathit{gam}$ & $\mathit{D1} + \mathit{gam} = \mathit{R14}$ \\
58 & $\mathit{a\_square} = \mathit{a} \cdot \mathit{a}$ & $\mathit{a} \cdot \mathit{a} = \mathit{a\_square}$ \\
59 & $\mathit{A} = \mathit{a\_square} + \mathit{a4m5}$ & $\mathit{a\_square} + \mathit{a4m5} = \mathit{A}$ \\
60 & $\mathit{c2} = \mathit{c} \cdot \mathit{c}$ & $\mathit{c} \cdot \mathit{c} = \mathit{c2}$ \\
61 & $\mathit{Ac2} = \mathit{A} \cdot \mathit{c2}$ & $\mathit{A} \cdot \mathit{c2} = \mathit{Ac2}$ \\
62 & $\mathit{R15} = \mathit{Ac2} + 1$ & $\mathit{Ac2} + 1 = \mathit{R15}$ \\
63 & $\mathit{L15} = \mathit{d} \cdot \mathit{d}$ & $\mathit{d} \cdot \mathit{d} = \mathit{L15}$ \\
64 & $\mathit{ic2} = \mathit{i} \cdot \mathit{c2}$ & $\mathit{i} \cdot \mathit{c2} = \mathit{ic2}$ \\
65 & $\mathit{ic22} = \mathit{ic2} \cdot \mathit{ic2}$ & $\mathit{ic2} \cdot \mathit{ic2} = \mathit{ic22}$ \\
66 & $\mathit{L16} = \mathit{f} \cdot \mathit{f}$ & $\mathit{f} \cdot \mathit{f} = \mathit{L16}$ \\
67 & $\mathit{f\_square\_minus\_one} = \mathit{L16} - 1$ & $\mathit{f\_square\_minus\_one} + 1 = \mathit{L16}$ \\
68 & $\mathit{R16} = \mathit{A} \cdot \mathit{f\_square\_minus\_one}$ & $\mathit{A} \cdot \mathit{f\_square\_minus\_one} = \mathit{R16}$ \\
69 & $\mathit{of} = \mathit{o} \cdot \mathit{f}$ & $\mathit{o} \cdot \mathit{f} = \mathit{of}$ \\
70 & $\mathit{aux\_u\_rhs} = \mathit{c} + \mathit{of}$ & $\mathit{c} + \mathit{of} = \mathit{aux\_u\_rhs}$ \\
71 & $\mathit{jc} = \mathit{j} \cdot \mathit{c}$ & $\mathit{j} \cdot \mathit{c} = \mathit{jc}$ \\
72 & $\mathit{H17} = \mathit{tr1} + \mathit{jc}$ & $\mathit{tr1} + \mathit{jc} = \mathit{H17}$ \\
73 & $\mathit{H2} = \mathit{H17} \cdot \mathit{H17}$ & $\mathit{H17} \cdot \mathit{H17} = \mathit{H2}$ \\
74 & $\mathit{aux\_y2} = \mathit{y\_aux} \cdot \mathit{y\_aux}$ & $\mathit{y\_aux} \cdot \mathit{y\_aux} = \mathit{aux\_y2}$ \\
75 & $\mathit{aux\_square\_gap} = \mathit{H2} - \mathit{aux\_y2}$ & $\mathit{aux\_square\_gap} + \mathit{aux\_y2} = \mathit{H2}$ \\
76 & $\mathit{L17} = \mathit{ic22} \cdot \mathit{aux\_square\_gap}$ & $\mathit{ic22} \cdot \mathit{aux\_square\_gap} = \mathit{L17}$ \\
77 & $\mathit{P17} = 1 - \mathit{aux\_y2}$ & $\mathit{P17} + \mathit{aux\_y2} = 1$ \\
78 & $\mathit{amb} = \mathit{a} - \mathit{th}$ & $\mathit{amb} + \mathit{th} = \mathit{a}$ \\
79 & $\mathit{kamb} = \mathit{ka} \cdot \mathit{amb}$ & $\mathit{ka} \cdot \mathit{amb} = \mathit{kamb}$ \\
80 & $\mathit{qk} = \mathit{q} + \mathit{kamb}$ & $\mathit{q} + \mathit{kamb} = \mathit{qk}$ \\
81 & $\mathit{amb2} = \mathit{amb} \cdot \mathit{amb}$ & $\mathit{amb} \cdot \mathit{amb} = \mathit{amb2}$ \\
82 & $\mathit{Mod} = \mathit{A} - \mathit{amb2}$ & $\mathit{Mod} + \mathit{amb2} = \mathit{A}$ \\
83 & $\mathit{rM} = \mathit{rho} \cdot \mathit{Mod}$ & $\mathit{rho} \cdot \mathit{Mod} = \mathit{rM}$ \\
84 & $\mathit{R18} = \mathit{qk} + \mathit{rM}$ & $\mathit{qk} + \mathit{rM} = \mathit{R18}$ \\
85 & $\mathit{ka2} = \mathit{ka} \cdot \mathit{ka}$ & $\mathit{ka} \cdot \mathit{ka} = \mathit{ka2}$ \\
86 & $\mathit{Aka2} = \mathit{A} \cdot \mathit{ka2}$ & $\mathit{A} \cdot \mathit{ka2} = \mathit{Aka2}$ \\
87 & $\mathit{R19} = \mathit{Aka2} + 1$ & $\mathit{Aka2} + 1 = \mathit{R19}$ \\
88 & $\mathit{L19} = \mathit{mu} \cdot \mathit{mu}$ & $\mathit{mu} \cdot \mathit{mu} = \mathit{L19}$ \\
89 & $\mathit{Dam1} = \mathit{Delta} \cdot \mathit{a}$ & $\mathit{Delta} \cdot \mathit{a} = \mathit{Dam1}$ \\
90 & $\mathit{R20} = \mathit{Dam1} + \mathit{Tindex}$ & $\mathit{Dam1} + \mathit{Tindex} = \mathit{R20}$ \\
\end{longtable}

The twenty-two free equality tests are
\[\begin{array}{@{}l@{}}
\mathit{L1}=\mathit{q},\quad \mathit{b}=\mathit{bx},\quad \mathit{q2}=\mathit{R2},\\
\mathit{th}=\mathit{theta\_sum},\quad \mathit{S2}=\mathit{pack\_rhs},\quad \mathit{n}=\mathit{q8},\\
\mathit{r}=\mathit{R7},\quad \mathit{L9}=\mathit{R9},\quad \mathit{c}=\mathit{R10a},\\
\mathit{k}=\mathit{R10b},\quad \mathit{k}=\mathit{R11},\quad \mathit{a}=\mathit{R12},\\
\mathit{c}=\mathit{R13},\quad \mathit{d}=\mathit{R14},\quad \mathit{L15}=\mathit{R15},\\
\mathit{ic22}=\mathit{R16},\quad \mathit{L17}=\mathit{P17},\quad \mathit{mu}=\mathit{R18},\\
\mathit{L19}=\mathit{R19},\quad \mathit{ka}=\mathit{R20},\quad \mathit{S3}=\mathit{sigma},\\
\mathit{H17}=\mathit{aux\_u\_rhs}
\end{array}\]


\clearpage\normalsize
\section{The 89-instruction narrow binary certificate}\label{app:89}

Use exactly the fixed binary coefficient index and 34 positive unknowns
of Appendix~\ref{app:90}. The new source has 22 equations and uses
47 multiplications and 42 additions. Its three packed fields have widths
$q,q,q^2$, and $n=q^4$. Removing the baseline from the mask permits
this narrower power chain without an additional arithmetic operation.

Section~\ref{sec:narrow-binary-product} proves the wider preliminary
Pell range, recovers the narrower mask bound after the exponent equations,
excludes the possible spill into the fixed-code field, and supplies
fresh positive witnesses at the new packed index. The unchanged raw
query and fixed index still represent the same recursively enumerable set.

The complete checker
\file{round38\_1980\_binary\_product\_certificate.py}
verifies every primitive and all 22 fresh polynomial residuals.
The table is generated by
\file{round38\_1980\_schedule\_table.py}.
Literal numerals remain $1,3,4$; fixed numerals and equality tests are free.

\small
% Generated by round38_1980_schedule_table.py; do not edit.
\begin{longtable}{@{}rll@{}}
\toprule
No. & Assignment & Addition/multiplication check\\
\midrule
\endfirsthead
\toprule
No. & Assignment & Addition/multiplication check\\
\midrule
\endhead
\bottomrule
\endfoot
1 & $\mathit{q2} = \mathit{q} \cdot \mathit{q}$ & $\mathit{q} \cdot \mathit{q} = \mathit{q2}$ \\
2 & $\mathit{q4} = \mathit{q2} \cdot \mathit{q2}$ & $\mathit{q2} \cdot \mathit{q2} = \mathit{q4}$ \\
3 & $\mathit{theta\_sum} = \mathit{H} + \mathit{b}$ & $\mathit{H} + \mathit{b} = \mathit{theta\_sum}$ \\
4 & $\mathit{bx} = \mathit{x} + \mathit{beta}$ & $\mathit{x} + \mathit{beta} = \mathit{bx}$ \\
5 & $\mathit{theta\_lambda} = \mathit{th} \cdot \mathit{la}$ & $\mathit{th} \cdot \mathit{la} = \mathit{theta\_lambda}$ \\
6 & $\mathit{Tcoef} = \mathit{theta\_lambda} + 1$ & $\mathit{theta\_lambda} + 1 = \mathit{Tcoef}$ \\
7 & $\mathit{R2} = \mathit{Tcoef} + \mathit{la}$ & $\mathit{Tcoef} + \mathit{la} = \mathit{R2}$ \\
8 & $\mathit{t2} = \mathit{e} - \mathit{l}$ & $\mathit{t2} + \mathit{l} = \mathit{e}$ \\
9 & $\mathit{eq} = \mathit{e} \cdot \mathit{q}$ & $\mathit{e} \cdot \mathit{q} = \mathit{eq}$ \\
10 & $\mathit{S2} = \mathit{l} + \mathit{eq}$ & $\mathit{l} + \mathit{eq} = \mathit{S2}$ \\
11 & $\mathit{pack\_tth} = \mathit{t} \cdot \mathit{th}$ & $\mathit{t} \cdot \mathit{th} = \mathit{pack\_tth}$ \\
12 & $\mathit{pack\_rhs} = \mathit{V} + \mathit{pack\_tth}$ & $\mathit{V} + \mathit{pack\_tth} = \mathit{pack\_rhs}$ \\
13 & $\mathit{C} = \mathit{x} + \mathit{g}$ & $\mathit{x} + \mathit{g} = \mathit{C}$ \\
14 & $\mathit{C2} = \mathit{C} \cdot \mathit{C}$ & $\mathit{C} \cdot \mathit{C} = \mathit{C2}$ \\
15 & $\mathit{bound\_sum} = \mathit{l} + \mathit{sigma}$ & $\mathit{l} + \mathit{sigma} = \mathit{bound\_sum}$ \\
16 & $\mathit{L1} = \mathit{bound\_sum} + \mathit{al}$ & $\mathit{bound\_sum} + \mathit{al} = \mathit{L1}$ \\
17 & $\mathit{S3} = \mathit{t2} \cdot \mathit{C2}$ & $\mathit{t2} \cdot \mathit{C2} = \mathit{S3}$ \\
18 & $\mathit{S3q4} = \mathit{S2} \cdot \mathit{q}$ & $\mathit{S2} \cdot \mathit{q} = \mathit{S3q4}$ \\
19 & $\mathit{Sin} = \mathit{S3} + \mathit{S3q4}$ & $\mathit{S3} + \mathit{S3q4} = \mathit{Sin}$ \\
20 & $\mathit{Sq4} = \mathit{Sin} \cdot \mathit{q}$ & $\mathit{Sin} \cdot \mathit{q} = \mathit{Sq4}$ \\
21 & $\mathit{S} = \mathit{g} + \mathit{Sq4}$ & $\mathit{g} + \mathit{Sq4} = \mathit{S}$ \\
22 & $\mathit{n2} = \mathit{n} \cdot \mathit{n}$ & $\mathit{n} \cdot \mathit{n} = \mathit{n2}$ \\
23 & $\mathit{n2n} = \mathit{n2} - \mathit{n}$ & $\mathit{n2n} + \mathit{n} = \mathit{n2}$ \\
24 & $\mathit{SA} = \mathit{S} \cdot \mathit{n2n}$ & $\mathit{S} \cdot \mathit{n2n} = \mathit{SA}$ \\
25 & $\mathit{Tq4} = \mathit{theta\_lambda} \cdot \mathit{q2}$ & $\mathit{theta\_lambda} \cdot \mathit{q2} = \mathit{Tq4}$ \\
26 & $\mathit{theta\_q4} = \mathit{th} \cdot \mathit{q}$ & $\mathit{th} \cdot \mathit{q} = \mathit{theta\_q4}$ \\
27 & $\mathit{mask\_gap} = \mathit{theta\_q4} - \mathit{b}$ & $\mathit{mask\_gap} + \mathit{b} = \mathit{theta\_q4}$ \\
28 & $\mathit{indicator\_mask} = \mathit{l} \cdot \mathit{mask\_gap}$ & $\mathit{l} \cdot \mathit{mask\_gap} = \mathit{indicator\_mask}$ \\
29 & $\mathit{T} = \mathit{Tq4} + \mathit{indicator\_mask}$ & $\mathit{Tq4} + \mathit{indicator\_mask} = \mathit{T}$ \\
30 & $\mathit{n2m1} = \mathit{n2} - 1$ & $\mathit{n2m1} + 1 = \mathit{n2}$ \\
31 & $\mathit{TB} = \mathit{T} \cdot \mathit{n2m1}$ & $\mathit{T} \cdot \mathit{n2m1} = \mathit{TB}$ \\
32 & $\mathit{R7} = \mathit{SA} + \mathit{TB}$ & $\mathit{SA} + \mathit{TB} = \mathit{R7}$ \\
33 & $\mathit{wn2} = \mathit{w} \cdot \mathit{n2}$ & $\mathit{w} \cdot \mathit{n2} = \mathit{wn2}$ \\
34 & $\mathit{sn2} = \mathit{s} \cdot \mathit{n2}$ & $\mathit{s} \cdot \mathit{n2} = \mathit{sn2}$ \\
35 & $\mathit{UM} = \mathit{wn2} \cdot \mathit{sn2}$ & $\mathit{wn2} \cdot \mathit{sn2} = \mathit{UM}$ \\
36 & $\mathit{R12} = \mathit{UM} + \mathit{sn2}$ & $\mathit{UM} + \mathit{sn2} = \mathit{R12}$ \\
37 & $\mathit{ksn2} = \mathit{k} \cdot \mathit{sn2}$ & $\mathit{k} \cdot \mathit{sn2} = \mathit{ksn2}$ \\
38 & $\mathit{UM2} = \mathit{UM} \cdot \mathit{UM}$ & $\mathit{UM} \cdot \mathit{UM} = \mathit{UM2}$ \\
39 & $\mathit{scaled\_norm\_coefficient} = \mathit{UM2} + \mathit{wn2}$ & $\mathit{UM2} + \mathit{wn2} = \mathit{scaled\_norm\_coefficient}$ \\
40 & $\mathit{ratio\_product2} = \mathit{ksn2} \cdot \mathit{ksn2}$ & $\mathit{ksn2} \cdot \mathit{ksn2} = \mathit{ratio\_product2}$ \\
41 & $\mathit{L9} = \mathit{scaled\_norm\_coefficient} \cdot \mathit{ratio\_product2}$ & $\mathit{scaled\_norm\_coefficient} \cdot \mathit{ratio\_product2} = \mathit{L9}$ \\
42 & $\mathit{tauplus1} = \mathit{tau} + 1$ & $\mathit{tau} + 1 = \mathit{tauplus1}$ \\
43 & $\mathit{R9} = \mathit{tau} \cdot \mathit{tauplus1}$ & $\mathit{tau} \cdot \mathit{tauplus1} = \mathit{R9}$ \\
44 & $\mathit{R10a} = \mathit{ksn2} + \mathit{eta}$ & $\mathit{ksn2} + \mathit{eta} = \mathit{R10a}$ \\
45 & $\mathit{R10b} = \mathit{eta} + \mathit{zeta}$ & $\mathit{eta} + \mathit{zeta} = \mathit{R10b}$ \\
46 & $\mathit{r1} = \mathit{r} + 1$ & $\mathit{r} + 1 = \mathit{r1}$ \\
47 & $\mathit{hpm1} = \mathit{h} \cdot \mathit{UM}$ & $\mathit{h} \cdot \mathit{UM} = \mathit{hpm1}$ \\
48 & $\mathit{R11} = \mathit{r1} + \mathit{hpm1}$ & $\mathit{r1} + \mathit{hpm1} = \mathit{R11}$ \\
49 & $\mathit{tr1} = \mathit{r1} + \mathit{r}$ & $\mathit{r1} + \mathit{r} = \mathit{tr1}$ \\
50 & $\mathit{R13} = \mathit{ka} + \mathit{phi}$ & $\mathit{ka} + \mathit{phi} = \mathit{R13}$ \\
51 & $\mathit{cam2} = \mathit{c} \cdot \mathit{a}$ & $\mathit{c} \cdot \mathit{a} = \mathit{cam2}$ \\
52 & $\mathit{D1} = \mathit{wn2} + \mathit{cam2}$ & $\mathit{wn2} + \mathit{cam2} = \mathit{D1}$ \\
53 & $\mathit{a4} = 4 \cdot \mathit{a}$ & $4 \cdot \mathit{a} = \mathit{a4}$ \\
54 & $\mathit{a4m5} = \mathit{a4} + 3$ & $\mathit{a4} + 3 = \mathit{a4m5}$ \\
55 & $\mathit{gam} = \mathit{ga} \cdot \mathit{a4m5}$ & $\mathit{ga} \cdot \mathit{a4m5} = \mathit{gam}$ \\
56 & $\mathit{R14} = \mathit{D1} + \mathit{gam}$ & $\mathit{D1} + \mathit{gam} = \mathit{R14}$ \\
57 & $\mathit{a\_square} = \mathit{a} \cdot \mathit{a}$ & $\mathit{a} \cdot \mathit{a} = \mathit{a\_square}$ \\
58 & $\mathit{A} = \mathit{a\_square} + \mathit{a4m5}$ & $\mathit{a\_square} + \mathit{a4m5} = \mathit{A}$ \\
59 & $\mathit{c2} = \mathit{c} \cdot \mathit{c}$ & $\mathit{c} \cdot \mathit{c} = \mathit{c2}$ \\
60 & $\mathit{Ac2} = \mathit{A} \cdot \mathit{c2}$ & $\mathit{A} \cdot \mathit{c2} = \mathit{Ac2}$ \\
61 & $\mathit{R15} = \mathit{Ac2} + 1$ & $\mathit{Ac2} + 1 = \mathit{R15}$ \\
62 & $\mathit{L15} = \mathit{d} \cdot \mathit{d}$ & $\mathit{d} \cdot \mathit{d} = \mathit{L15}$ \\
63 & $\mathit{ic2} = \mathit{i} \cdot \mathit{c2}$ & $\mathit{i} \cdot \mathit{c2} = \mathit{ic2}$ \\
64 & $\mathit{ic22} = \mathit{ic2} \cdot \mathit{ic2}$ & $\mathit{ic2} \cdot \mathit{ic2} = \mathit{ic22}$ \\
65 & $\mathit{L16} = \mathit{f} \cdot \mathit{f}$ & $\mathit{f} \cdot \mathit{f} = \mathit{L16}$ \\
66 & $\mathit{f\_square\_minus\_one} = \mathit{L16} - 1$ & $\mathit{f\_square\_minus\_one} + 1 = \mathit{L16}$ \\
67 & $\mathit{R16} = \mathit{A} \cdot \mathit{f\_square\_minus\_one}$ & $\mathit{A} \cdot \mathit{f\_square\_minus\_one} = \mathit{R16}$ \\
68 & $\mathit{of} = \mathit{o} \cdot \mathit{f}$ & $\mathit{o} \cdot \mathit{f} = \mathit{of}$ \\
69 & $\mathit{aux\_u\_rhs} = \mathit{c} + \mathit{of}$ & $\mathit{c} + \mathit{of} = \mathit{aux\_u\_rhs}$ \\
70 & $\mathit{jc} = \mathit{j} \cdot \mathit{c}$ & $\mathit{j} \cdot \mathit{c} = \mathit{jc}$ \\
71 & $\mathit{H17} = \mathit{tr1} + \mathit{jc}$ & $\mathit{tr1} + \mathit{jc} = \mathit{H17}$ \\
72 & $\mathit{H2} = \mathit{H17} \cdot \mathit{H17}$ & $\mathit{H17} \cdot \mathit{H17} = \mathit{H2}$ \\
73 & $\mathit{aux\_y2} = \mathit{y\_aux} \cdot \mathit{y\_aux}$ & $\mathit{y\_aux} \cdot \mathit{y\_aux} = \mathit{aux\_y2}$ \\
74 & $\mathit{aux\_square\_gap} = \mathit{H2} - \mathit{aux\_y2}$ & $\mathit{aux\_square\_gap} + \mathit{aux\_y2} = \mathit{H2}$ \\
75 & $\mathit{L17} = \mathit{ic22} \cdot \mathit{aux\_square\_gap}$ & $\mathit{ic22} \cdot \mathit{aux\_square\_gap} = \mathit{L17}$ \\
76 & $\mathit{P17} = 1 - \mathit{aux\_y2}$ & $\mathit{P17} + \mathit{aux\_y2} = 1$ \\
77 & $\mathit{amb} = \mathit{a} - \mathit{th}$ & $\mathit{amb} + \mathit{th} = \mathit{a}$ \\
78 & $\mathit{kamb} = \mathit{ka} \cdot \mathit{amb}$ & $\mathit{ka} \cdot \mathit{amb} = \mathit{kamb}$ \\
79 & $\mathit{qk} = \mathit{q} + \mathit{kamb}$ & $\mathit{q} + \mathit{kamb} = \mathit{qk}$ \\
80 & $\mathit{amb2} = \mathit{amb} \cdot \mathit{amb}$ & $\mathit{amb} \cdot \mathit{amb} = \mathit{amb2}$ \\
81 & $\mathit{Mod} = \mathit{A} - \mathit{amb2}$ & $\mathit{Mod} + \mathit{amb2} = \mathit{A}$ \\
82 & $\mathit{rM} = \mathit{rho} \cdot \mathit{Mod}$ & $\mathit{rho} \cdot \mathit{Mod} = \mathit{rM}$ \\
83 & $\mathit{R18} = \mathit{qk} + \mathit{rM}$ & $\mathit{qk} + \mathit{rM} = \mathit{R18}$ \\
84 & $\mathit{ka2} = \mathit{ka} \cdot \mathit{ka}$ & $\mathit{ka} \cdot \mathit{ka} = \mathit{ka2}$ \\
85 & $\mathit{Aka2} = \mathit{A} \cdot \mathit{ka2}$ & $\mathit{A} \cdot \mathit{ka2} = \mathit{Aka2}$ \\
86 & $\mathit{R19} = \mathit{Aka2} + 1$ & $\mathit{Aka2} + 1 = \mathit{R19}$ \\
87 & $\mathit{L19} = \mathit{mu} \cdot \mathit{mu}$ & $\mathit{mu} \cdot \mathit{mu} = \mathit{L19}$ \\
88 & $\mathit{Dam1} = \mathit{Delta} \cdot \mathit{a}$ & $\mathit{Delta} \cdot \mathit{a} = \mathit{Dam1}$ \\
89 & $\mathit{R20} = \mathit{Dam1} + \mathit{Tindex}$ & $\mathit{Dam1} + \mathit{Tindex} = \mathit{R20}$ \\
\end{longtable}

The twenty-two free equality tests are
\[\begin{array}{@{}l@{}}
\mathit{L1}=\mathit{q},\quad \mathit{b}=\mathit{bx},\quad \mathit{q2}=\mathit{R2},\\
\mathit{th}=\mathit{theta\_sum},\quad \mathit{S2}=\mathit{pack\_rhs},\quad \mathit{n}=\mathit{q4},\\
\mathit{r}=\mathit{R7},\quad \mathit{L9}=\mathit{R9},\quad \mathit{c}=\mathit{R10a},\\
\mathit{k}=\mathit{R10b},\quad \mathit{k}=\mathit{R11},\quad \mathit{a}=\mathit{R12},\\
\mathit{c}=\mathit{R13},\quad \mathit{d}=\mathit{R14},\quad \mathit{L15}=\mathit{R15},\\
\mathit{ic22}=\mathit{R16},\quad \mathit{L17}=\mathit{P17},\quad \mathit{mu}=\mathit{R18},\\
\mathit{L19}=\mathit{R19},\quad \mathit{ka}=\mathit{R20},\quad \mathit{S3}=\mathit{sigma},\\
\mathit{H17}=\mathit{aux\_u\_rhs}
\end{array}\]


\end{document}
